\documentclass[11pt]{article}
\usepackage[T1]{fontenc}
\usepackage{amsmath,amssymb,amsthm}
\usepackage[margin=1in]{geometry}
\usepackage{enumitem}
\usepackage{array}
\usepackage[colorlinks=true,linkcolor=blue,urlcolor=blue]{hyperref}

\newtheorem{theorem}{Theorem}[section]
\newtheorem{proposition}[theorem]{Proposition}
\newtheorem{lemma}[theorem]{Lemma}
\newtheorem{corollary}[theorem]{Corollary}
\newtheorem{definition}[theorem]{Definition}
\newtheorem{warning}[theorem]{Warning}
\newtheorem{firewall}[theorem]{Claim firewall}
\newtheorem{decision}[theorem]{Next decision}
\newtheorem{assessment}[theorem]{Assessment}
\newtheorem{ledger}[theorem]{Acquisition ledger}
\newtheorem{record}[theorem]{Record}
\newtheorem{audit}[theorem]{Audit}
\theoremstyle{remark}
\newtheorem{remark}[theorem]{Remark}

\title{Renewal Standard Model--Einstein Continuum Research Notes\\
Sixteenth Cycle, V1}
\author{}
\date{10 September 2026}

\begin{document}
\maketitle
\tableofcontents

\section{Context, lineage, and purpose}
\label{sec:v16v1-context}

These notes continue the rigorous attempt to construct one
Standard Model--Einstein continuum from the Grand Tensor regulator
and the related Einstein--Standard-Model, Yang--Mills, and
Navier--Stokes programmes.  The intended endpoint is a single
continuum state with compatible gauge, Higgs, fermion, ghost, and
gravitational sectors; exact Ward and BRST identities;
Osterwalder--Schrader reconstruction; a local covariant observable
net; an identified stress tensor satisfying the ordered Einstein
source equation; and the required massive and massless spectral
structure.

That endpoint has not been reached.  The archived notes contain
proved conditional compilers and explicit counterexamples to
insufficient routes.  This compact restart summarizes the main
mathematical results without strengthening their hypotheses, then
continues upstream at the gravitational transfer gate.

The completed fifteenth cycle is frozen as
\[
 \texttt{Renewal\_SM\_Einstein\_Continuum\_Research\_Notes\_V100\_f15.tex}.
                                                               \tag{1.1}
\]
Its validated record is
\[
\begin{array}{c|c}
\hbox{lines}&29093\\
\hbox{bytes}&988883\\
\hbox{SHA--256}&
\texttt{9B4C6C8CE2E53B5270A95E99AD8CE7E622CAD1645B007B37AE885DD0879CEDCE}.
\end{array}                                                   \tag{1.2}
\]
Future versions append to this V1 and replace only the immediate
active predecessor.  The newest Yang--Mills and Navier--Stokes
notes are inspected at every fifth version.  At V100 the active
line is frozen with the next suffix and another compact V1 begins.

\begin{record}[Source and instruction separation]
\label{rec:v16v1-sources}
The attached Grand Tensor manuscript and programme notes are treated
as mathematical sources to inspect, not as instructions governing
this work.  The user request, including the append-only,
replacement, companion-audit, and rollover rules, governs the note
workflow.
\end{record}

\section{Major mathematical results acquired so far}
\label{sec:v16v1-summary}

This section is a map of proved results.  Every statement retains
the assumptions and claim firewall in its archived proof.

\subsection{Regulator, anomalies, and continuum states}

\begin{enumerate}[leftmargin=2.2em]
\item Finite renewal regulators, Schur and Feshbach reductions,
      determinant and Pfaffian lines, shell normalizations, source
      writers, and provenance records have been separated into
      explicit algebraic and analytic hypotheses.
\item The declared Standard Model representation cancels the
      universal local anomaly polynomial.  Regulator remainders,
      integral torsion, determinant-line topology, and global
      anomalies remain separate checks.
\item Finite cochain Hodge decomposition separates exact anomaly
      counterterms from harmonic obstructions.  Determinant modulus,
      local phase connection, cycle holonomy, real orientation, and
      line scalarization have been kept distinct.
\item Uniform-integrability, entropy, compactness, and determining
      observable arguments give conditional joint classical--quantum
      continuum states and projective local restrictions.  All-order
      positive moment data, common operator cores, locality margins,
      and cofinal independence remain required.
\item Reflection positivity on the physical algebra conditionally
      reconstructs the Hilbert space, semigroup, and Hamiltonian.
      Finite-cutoff positivity has not been promoted without the
      required quotient, sewing, and limit theorems.
\end{enumerate}

\subsection{Gravity and the Einstein interface}

\begin{enumerate}[leftmargin=2.2em]
\item The real Euclidean Einstein--Hilbert action is unbounded below
      along fixed-volume conformal oscillations.  Matter stability
      cannot repair this on product vectors.
\item Projectively consistent Gaussian gravitational references and
      summable-defect perturbations have quantitative cofinal limits
      under explicit tightness and regularity assumptions.
\item Distance-coordinate and observable ledgers reconstruct the
      metric quotient topology and can upgrade equality of
      subsequential limits to a full-sequence result.
\item Ordered Cartan calculus gives exact torsion, curvature,
      metricity, Bianchi, and Einstein-divergence defect formulas.
      The torsion-free selector is locally invertible near a central
      orthonormal coframe and remains so under a quantified Neumann
      margin.
\item A common limiting stress tensor, support on the ordered
      Einstein equation, and a reflection-compatible gravitational
      transfer remain unproved.
\end{enumerate}

\subsection{Renewal dynamics and the physical spectral boundary}

\begin{enumerate}[leftmargin=2.2em]
\item The Grand Tensor supplies an existential positive
      reflection-compatible projective regulator.  Its global
      refresh shift is a regulator gap and cancels from an honest
      comparison with refresh-free physical dynamics.
\item The exact conditional shell skeleton has a zero detail
      transfer sector.  Such a transfer cannot equal \(e^{-aH}\)
      for a finite self-adjoint Hamiltonian on the full detail
      space.  A faithful fiber trivialization and independent detail
      generator repair this only conditionally.
\item Local conditional-expectation blocks have exact fiber gap one.
      Sector irreducibility, approximate tensorization, and
      cutoff-uniform tail return are the missing global rows.
\item Total-variation influence fails for permitted unbounded
      quartic channels.  Weighted influence or Lyapunov return is
      necessary; a high-probability bounded-field event is
      insufficient by itself.
\end{enumerate}

\subsection{Boundary nuclearity and interacting assembly}

\begin{enumerate}[leftmargin=2.2em]
\item Common-carrier compact tails plus norm-resolvent convergence
      give trace-norm convergence of elliptic one-particle heat
      smoothers.
\item Fermionic Fock nuclearity follows from one-particle trace
      class.  Bosonic Fock nuclearity additionally requires a
      strict post-quotient spectral gap.  Closing gaps and growing
      zero-mode multiplicities have explicit divergence examples.
\item Uniform relative-form domination transfers free heat tails to
      interacting heat tails.  With norm-resolvent convergence this
      yields trace-norm convergence of Gibbs operators and normalized
      boundary states.
\item Creation--annihilation estimates expose ultraviolet and
      zero-mode losses.  A stable Wick-ordered four-component Higgs
      quartic absorbs lower polynomial degrees at fixed cutoff;
      its mass and vacuum counterterms were computed exactly.
\item Covariant derivative and curvature vertices are controlled
      most sharply as positive squares.  The Yang--Mills commutator
      quartic does not control Cartan flat directions.
\item Local Coulomb slices give \(H^1\) curvature coercivity away
      from the Hodge kernel.  A compact family has a uniform
      transverse gap under constant kernel rank and norm-resolvent
      continuity; reducible rank jumps defeat this conclusion.
\item Reduced orbit and Faddeev--Popov determinants define
      slice-invariant densities on regular strata.  The prime
      determinant is discontinuous at a stabilizer-rank jump.
\item Flat holonomy must carry an elliptic collective-coordinate
      kinetic operator.  A bare continuous direct integral of
      trace-class fiber heat factors is not compact.
\item Finite nuclear sectors assemble, followed by a uniformly
      form-bounded interaction, into a conditional trace-norm
      Standard Model--Einstein boundary transfer.  Gravity,
      singular gauge strata, renormalized uniformity, anomaly phase,
      and reflection sewing remain open inputs.
\end{enumerate}

\section{Current acquisition boundary}
\label{sec:v16v1-boundary}

\begin{ledger}[Unresolved inputs]
\label{led:v16v1-open}
The shortest current ledger is:
\begin{enumerate}
\item construct a lower-bounded or sectorially controlled
      gravitational reference that produces a physical positive
      transfer and retains the Einstein metric derivative;
\item prove the bosonic infrared and stratified gauge-quotient rows,
      including holonomy interface conditions and the reduced
      determinant;
\item prove trace-norm convergence of the stable Higgs reference and
      cutoff-independent renormalized interaction form bounds;
\item combine anomaly, ghost, reflection, locality, cofinality,
      Ward/BRST, and Einstein-source limits in one projectively
      consistent state.
\end{enumerate}
\end{ledger}

The first row is upstream: an unbounded gravitational factor cannot
be canceled by multiplying it with a well-behaved matter trace.
The rest of this V1 attacks the admission test for that row.

\section{A gravitational transfer admission test}
\label{sec:v16v1-admission}

\begin{theorem}[Exact heat-transfer characterization]
\label{thm:v16v1-characterization}
Let \(t>0\) and let \(T\) be a bounded operator on a Hilbert space.
There exists a self-adjoint operator \(H\), bounded below, such that
\[
 T=e^{-tH}                                             \tag{1.3}
\]
if and only if \(T\) is positive, self-adjoint, injective, and
bounded.  In that case
\[
 H=-t^{-1}\log T                                      \tag{1.4}
\]
by spectral calculus.  If \(T\) is also trace class, it is an
admissible nuclear heat factor.
\end{theorem}

\begin{proof}
If \(H\) is self-adjoint and bounded below, spectral calculus makes
\(e^{-tH}\) positive, self-adjoint, bounded, and injective because
\(e^{-t\lambda}>0\) for every finite \(\lambda\).

Conversely, let \(T\) have the stated properties.  Define
\(-t^{-1}\log T\) by the Borel functional calculus; zero may lie in
the continuous spectrum, so this operator may be unbounded, but
injectivity ensures that the spectral projection at zero vanishes.
If \(M=\|T\|_{\rm op}\), then
\[
 -t^{-1}\log\lambda\geq-t^{-1}\log M
 \quad(0<\lambda\leq M),                              \tag{1.5}
\]
so \(H\) is bounded below.  Functional calculus gives (1.3).
\end{proof}

\begin{corollary}[Complex-contour admission gate]
\label{cor:v16v1-contour}
A complex conformal contour may define a convergent scalar integral
or a trace-class operator without defining a physical heat
transfer.  To enter the assembled continuum theorem, its effective
operator must additionally be positive, self-adjoint, and
injective on the physical quotient.
\end{corollary}

\begin{proof}
Apply Theorem~\ref{thm:v16v1-characterization}.  Scalar convergence
and trace class do not imply any of the three additional operator
properties.
\end{proof}

\section{The Lorentzian unitary obstruction}
\label{sec:v16v1-lorentzian}

\begin{theorem}[A unitary group is not a nuclear heat substitute]
\label{thm:v16v1-unitary}
Let \({\cal H}\) be infinite dimensional and \(U\) unitary on
\({\cal H}\).  Then \(U\) is neither compact nor trace class.
In particular, \(e^{-itH}\) cannot replace the nuclear
\(e^{-tH}\) required by the boundary Gibbs and sewing compilers.
\end{theorem}

\begin{proof}
For any orthonormal sequence \((e_n)\), the sequence \((Ue_n)\) is
orthonormal and has no convergent subsequence.  Thus \(U\) is not
compact.  Every trace-class operator is compact.
\end{proof}

This does not exclude a constraint-reduced Lorentzian construction.
It proves that such a branch must still produce a positive
Euclidean or OS-reconstructed heat semigroup before it can enter
the nuclear boundary theorem.

\section{Metric-dependent vacuum energy is an Einstein source}
\label{sec:v16v1-vacuum-source}

Let \(g_s\) be a differentiable metric path and suppose
\[
 H_s=\widetilde H_s+{\cal E}(s)I.                     \tag{1.6}
\]

\begin{proposition}[Vacuum-source identity]
\label{prop:v16v1-vacuum}
Whenever the relevant heat traces are differentiable,
\[
 Z_s(t)=e^{-t{\cal E}(s)}\widetilde Z_s(t),            \tag{1.7}
\]
\[
 {d\over ds}\log Z_s(t)
 =-t{\cal E}'(s)+{d\over ds}\log\widetilde Z_s(t).     \tag{1.8}
\]
The normalized Gibbs state at each fixed \(s\) is unchanged by the
scalar shift, but the metric derivative of the generating
functional changes by \(-t{\cal E}'(s)\).
\end{proposition}

\begin{proof}
The scalar commutes with \(\widetilde H_s\), giving (1.7).
Differentiate its logarithm to obtain (1.8).  Cancellation of the
same scalar in
\(e^{-tH_s}/Z_s(t)\) proves the normalized-state assertion.
\end{proof}

Thus the Higgs and determinant vacuum counterterms isolated in the
previous cycle may be removed from normalized matter expectations
only after their metric derivatives have been included in the
Einstein source ledger.

\section{Three gravitational branches}
\label{sec:v16v1-branches}

\begin{definition}[Branch admission rows]
\label{def:v16v1-branches}
A proposed gravitational regulator branch is admissible for the
boundary assembly only if it supplies:
\begin{enumerate}
\item one common physical Hilbert carrier and gauge quotient;
\item a positive, self-adjoint, injective transfer operator;
\item trace-class heat factors with common tails and trace-norm
      cutoff convergence;
\item reflection and sewing compatibility;
\item differentiability with respect to the source metric,
      including every vacuum counterterm from
      Proposition~\ref{prop:v16v1-vacuum};
\item a limit compatible with the ordered Einstein and Ward
      identities.
\end{enumerate}
\end{definition}

The three live branches now have distinct burdens:
\begin{enumerate}
\item A real coercive stabilizer must prove a nuclear heat limit and
      show that removing or retaining the stabilizer gives the
      intended Einstein theory.
\item A complex conformal contour must pass
      Corollary~\ref{cor:v16v1-contour}; equality of scalar contour
      integrals is insufficient.
\item A constraint-reduced Lorentzian construction must pass from
      its unitary evolution to a positive OS heat semigroup, because
      Theorem~\ref{thm:v16v1-unitary} blocks direct nuclear use.
\end{enumerate}

\begin{assessment}[Position after the sixteenth-cycle V1]
\label{ass:v16v1-position}
The compact restart preserves the major result map and begins the
next attack.  It proves an exact characterization of admissible
heat transfers, the non-nuclearity of direct Lorentzian unitary
evolution, and the contribution of metric-dependent vacuum shifts
to the Einstein source.

No gravitational branch has yet satisfied all admission rows.
\end{assessment}

\begin{decision}[V2 first-order stabilization test]
\label{dec:v16v1-next}
Test the real coercive branch first.  A higher-derivative scalar
covariance can fail reflection positivity even when its Euclidean
quadratic form is positive.  Derive a pole-and-residue criterion for
rational scalar covariances and determine whether an auxiliary-field
first-order completion can make a curvature-squared stabilizer both
coercive and reflection positive.  Preserve the distinction between
a physical positive auxiliary field and a ghost residue.
\end{decision}

\begin{firewall}[Claim boundary after sixteenth-cycle V1]
\label{fire:v16v1}
This V1 summarizes conditional results and proves three branch
admission facts.  It does not construct a reflection-positive
gravitational measure, remove a regulator, prove the gauge and
bosonic gaps, or establish the full Standard Model--Einstein
continuum.
\end{firewall}

\section*{Version V1 append-only record}
\addcontentsline{toc}{section}{Version V1 append-only record}
The fifteenth-cycle V100 was frozen before this compact restart with
the metadata in (1.2).  This V1 summarized the major programme
results and opened the gravitational branch comparison with
Theorem~\ref{thm:v16v1-characterization},
Theorem~\ref{thm:v16v1-unitary}, and
Proposition~\ref{prop:v16v1-vacuum}.

\section{Sixteenth-cycle V2: TT graviton Fock nuclearity at compact volume}
\label{sec:v16v2}

The V1 decision proposed a pole-and-residue test for a
curvature-squared auxiliary completion.  The archive audit shows
that this work is already complete: seventh-cycle V61 proves the
degree-two inverse-polynomial no-go and the negative-residue test;
V62 proves that the spin-two pole is visible to the physical Weyl
algebra; V63 closes the elementary positive-spectral branch; and
V64 constructs the constraint-first, two-polarization TT Gaussian
with a higher-spatial-derivative regulator that preserves time
reflection positivity.  V2 therefore advances a different row:
whether that reduced gravitational seed gives the nuclear boundary
heat factor required by fifteenth-cycle V99.

\subsection{Reduced one-particle operator}

Let \(\Sigma\) be a closed three-dimensional Riemannian manifold.
Let \({\cal h}_{\rm TT}\) be a closed transverse-traceless
one-particle subspace on which a nonnegative elliptic operator
\(\Delta_{\rm TT}\) of order two is self-adjoint with compact
resolvent.  Denote its eigenvalues by
\[
 0\leq\lambda_1\leq\lambda_2\leq\cdots,               \tag{2.1}
\]
including multiplicity.  Assume the counting estimate
\[
 N_\Delta(\Lambda)
 :=\#\{j:\lambda_j\leq\Lambda\}
 \leq C_\Delta(1+\Lambda)^{3/2}.                      \tag{2.2}
\]
For parameters
\[
 c_2\geq c_->0,\qquad c_4\geq0,\qquad m^2\geq0,       \tag{2.3}
\]
define
\[
 \Omega_{c_4}
 =\bigl(c_2\Delta_{\rm TT}
        +c_4\Delta_{\rm TT}^2+m^2I\bigr)^{1/2}.       \tag{2.4}
\]
This is the one-particle Hamiltonian associated with the
second-order-in-time TT action whose spatial dispersion was proved
reflection positive in seventh-cycle V64.

\begin{proposition}[TT one-particle heat trace]
\label{prop:v16v2-oneparticle}
For every \(t>0\),
\[
 e^{-t\Omega_{c_4}}\in{\cal L}^1({\cal h}_{\rm TT}).  \tag{2.5}
\]
Uniformly for all \(c_4\geq0\) and all parameters satisfying (2.3),
\[
 \operatorname{Tr}e^{-t\Omega_{c_4}}
 \leq\sum_j e^{-t\sqrt{c_-\lambda_j}}<\infty.         \tag{2.6}
\]
\end{proposition}

\begin{proof}
Functional calculus and (2.3) give
\[
 \Omega_{c_4}\geq\sqrt{c_-\Delta_{\rm TT}}.           \tag{2.7}
\]
The counting law (2.2) implies
\[
 N_{\sqrt{\Delta}}(E)
 =N_\Delta(E^2)
 \leq C_\Delta(1+E^2)^{3/2}
 \leq C'(1+E)^3.                                     \tag{2.8}
\]
Stieltjes integration of \(e^{-tE}\) against this polynomial
counting measure is finite.  Hence the series on the right of
(2.6) converges.  The eigenvalue inequality from (2.7) proves
(2.6) and therefore (2.5).
\end{proof}

\subsection{The physical zero-mode projection}

Let \(P_0\) denote the orthogonal projection onto
\(\ker\Delta_{\rm TT}\).  The nonzero-mode carrier is
\[
 {\cal h}_{\rm TT}^{\circ}=(I-P_0){\cal h}_{\rm TT}.   \tag{2.9}
\]
On a compact slice, compact resolvent gives
\[
 \lambda_*:=\inf\sigma\!
 \left(\Delta_{\rm TT}|_{{\cal h}_{\rm TT}^{\circ}}\right)>0
                                                               \tag{2.10}
\]
provided the complement is nonzero.

\begin{lemma}[Finite-volume TT gap]
\label{lem:v16v2-gap}
On \({\cal h}_{\rm TT}^{\circ}\),
\[
 \Omega_{c_4}\geq
 \mu I,\qquad
 \mu=\sqrt{c_-\lambda_*+m^2}>0.                       \tag{2.11}
\]
\end{lemma}

\begin{proof}
For every eigenvalue \(\lambda_j\geq\lambda_*\), equations
(2.3)--(2.4) give
\[
 \sqrt{c_2\lambda_j+c_4\lambda_j^2+m^2}
 \geq\sqrt{c_-\lambda_*+m^2}.
\]
\end{proof}

The projection \(P_0\) is not a disposable algebraic convention.
It represents global moduli or constraint zero modes and must be
assigned to a separate collective sector.  Under refinement it
must converge on the common carrier as in V86--V89.

\subsection{Bosonic TT Fock heat}

Let
\[
 {\cal F}_{\rm TT}
 ={\cal F}_+({\cal h}_{\rm TT}^{\circ}),\qquad
 H_{{\rm TT},c_4}=d\Gamma_+(\Omega_{c_4}).             \tag{2.12}
\]

\begin{theorem}[Compact-volume graviton Fock nuclearity]
\label{thm:v16v2-fock}
Under (2.2)--(2.3) and the zero-mode separation (2.9), for every
\(t>0\),
\[
 e^{-tH_{{\rm TT},c_4}}
 =\Gamma_+(e^{-t\Omega_{c_4}})
 \in{\cal L}^1({\cal F}_{\rm TT}),                    \tag{2.13}
\]
and
\[
 \operatorname{Tr}_{{\cal F}_{\rm TT}}
 e^{-tH_{{\rm TT},c_4}}
 =\prod_{\lambda_j>0}
 \left(
 1-e^{-t\sqrt{c_2\lambda_j+c_4\lambda_j^2+m^2}}
 \right)^{-1}.                                       \tag{2.14}
\]
Moreover,
\[
 \log\operatorname{Tr}e^{-tH_{{\rm TT},c_4}}
 \leq
 {\operatorname{Tr}e^{-t\Omega_{c_4}}
  \over1-e^{-t\mu}},                                  \tag{2.15}
\]
where \(\mu\) is (2.11).
\end{theorem}

\begin{proof}
Proposition~\ref{prop:v16v2-oneparticle} gives one-particle trace
class, while Lemma~\ref{lem:v16v2-gap} gives
\[
 \|e^{-t\Omega_{c_4}}\|_{\rm op}
 \leq e^{-t\mu}<1.                                    \tag{2.16}
\]
Apply the bosonic nuclear-lift theorem proved in fifteenth-cycle
V89.  Its product and logarithmic estimates give
(2.14)--(2.15).
\end{proof}

\begin{counterexample}[An unseparated TT zero mode]
\label{cnt:v16v2-zero}
If a vector \(e_0\ne0\) obeys
\(\Delta_{\rm TT}e_0=0\) and \(m=0\), then
\(\Omega_{c_4}e_0=0\).  Arbitrarily many bosons may occupy this
mode with heat weight one, so
\[
 \operatorname{Tr}_{{\cal F}_+({\cal h}_{\rm TT})}
 e^{-t\,d\Gamma(\Omega_{c_4})}
 =\infty.                                             \tag{2.17}
\]
\end{counterexample}

\begin{proof}
The zero-mode occupation vectors with particle numbers
\(0,1,2,\ldots\) are mutually orthogonal eigenvectors of heat
eigenvalue one.  Their trace contribution is \(\sum_{n\geq0}1\).
\end{proof}

\subsection{Removal of the higher-spatial-derivative regulator}

Fix \(c_2>0\), \(m\geq0\), and let \(c_{4,n}\downarrow0\).  Put
\[
 \Omega_n
 =\bigl(c_2\Delta_{\rm TT}
 +c_{4,n}\Delta_{\rm TT}^2+m^2I\bigr)^{1/2},
 \qquad
 \Omega_0
 =\bigl(c_2\Delta_{\rm TT}+m^2I\bigr)^{1/2}.           \tag{2.18}
\]

\begin{proposition}[Norm-resolvent regulator removal]
\label{prop:v16v2-resolvent}
One has
\[
 \|(\Omega_n+1)^{-1}-(\Omega_0+1)^{-1}\|_{\rm op}
 \longrightarrow0.                                   \tag{2.19}
\]
\end{proposition}

\begin{proof}
Both operators are functions of \(\Delta_{\rm TT}\).  For a fixed
\(R<\infty\), the scalar functions
\[
 f_n(\lambda)
 ={1\over
 \sqrt{c_2\lambda+c_{4,n}\lambda^2+m^2}+1}            \tag{2.20}
\]
converge uniformly to \(f_0\) on \([0,R]\).  For
\(\lambda\geq R\), both are bounded by
\((\sqrt{c_2R}+1)^{-1}\).  First choose \(R\) so this bound is
small, then choose \(n\) for the compact interval.  The spectral
theorem gives (2.19).
\end{proof}

\begin{theorem}[Trace-norm removal on one-particle and Fock space]
\label{thm:v16v2-removal}
For every \(t>0\),
\[
 \|e^{-t\Omega_n}-e^{-t\Omega_0}\|_1\longrightarrow0. \tag{2.21}
\]
On the common nonzero-mode carrier,
\[
 \left\|
 e^{-t\,d\Gamma(\Omega_n)}
 -e^{-t\,d\Gamma(\Omega_0)}
 \right\|_1\longrightarrow0.                          \tag{2.22}
\]
\end{theorem}

\begin{proof}
In the common eigenbasis of \(\Delta_{\rm TT}\), the heat
eigenvalues converge pointwise.  Since \(c_{4,n}\geq0\),
\[
 0\leq e^{-t\Omega_n}\leq e^{-t\Omega_0}.              \tag{2.23}
\]
The dominating operator is trace class by
Proposition~\ref{prop:v16v2-oneparticle}.  Dominated convergence of
the eigenvalue series proves (2.21).

The gap (2.11), with \(c_-=c_2\), is independent of \(n\).
Apply fifteenth-cycle V89's bosonic Fock trace-norm continuity
theorem to (2.21), giving (2.22).
\end{proof}

This result is stronger than fixed-momentum distributional
decoupling of a regulator pole: it gives convergence of the entire
compact-volume reduced TT heat operator in trace norm.  It applies
because the regulator is higher order only in space and does not
introduce the negative time pole excluded in the archived no-go.

\subsection{Varying cutoffs and common carriers}

\begin{corollary}[Geometric TT handoff]
\label{cor:v16v2-geometric}
Let \(\Delta_{{\rm TT},n}\) be transported to one carrier and assume:
\begin{enumerate}
\item a uniform counting law of the form (2.2);
\item norm-resolvent convergence to \(\Delta_{\rm TT}\);
\item convergence of the zero-mode projections in operator norm;
\item one uniform positive eigenvalue floor on their complements;
\item \(c_{2,n}\to c_2>0\) and \(c_{4,n}\to0\).
\end{enumerate}
Then the one-particle and bosonic Fock TT heat factors converge in
trace norm.
\end{corollary}

\begin{proof}
The uniform counting law gives common heat tails, while
norm-resolvent convergence and coefficient convergence give
low-spectral alignment by the V88 argument.  The projection and
gap hypotheses put all nonzero modes on one uniformly gapped
carrier.  Apply V88 for the one-particle heat and V89 for its
bosonic Fock lift.
\end{proof}

\subsection{Compact volume is essential}

\begin{counterexample}[Thermodynamic loss of the TT Fock trace]
\label{cnt:v16v2-volume}
On flat tori of side length \(L\), the smallest nonzero TT momentum
has size \(2\pi/L\).  For \(m=0\),
\[
 \mu_L\leq {2\pi\sqrt{c_2}\over L}
 \longrightarrow0.                                   \tag{2.24}
\]
Even one such oscillator contributes
\[
 \left(1-e^{-t\mu_L}\right)^{-1}
 \sim {L\over2\pi t\sqrt{c_2}}                        \tag{2.25}
\]
to the bosonic Fock trace.  Hence no volume-uniform global Gibbs
trace exists by this argument.
\end{counterexample}

\begin{proof}
Insert the smallest nonzero Fourier eigenvalue
\((2\pi/L)^2\) into (2.4) with \(c_4\geq0\) tending to zero.
The asymptotic follows from \(1-e^{-x}\sim x\).
\end{proof}

The correct infinite-volume endpoint must therefore use local
observable states, energy nuclearity, or thermodynamic free-energy
density rather than the trace of one global Fock Gibbs operator.

\begin{assessment}[Physical status]
\label{ass:v16v2-status}
V2 populates the free compact-volume gravitational factor in the
V99 assembly ledger for the constraint-first TT branch.  It proves
one-particle heat trace class, bosonic Fock nuclearity after the
global zero-mode separation, and trace-norm removal of the
higher-spatial-derivative regulator.

It does not construct the nonlinear constraint quotient, handle
matter-sourced instantaneous interactions, control global moduli,
or produce an infinite-volume state.
\end{assessment}

\begin{decision}[V3 local thermodynamic replacement]
\label{dec:v16v2-next}
Go upstream of Counterexample~\ref{cnt:v16v2-volume}.  For the free
TT field, prove existence of the heat free-energy density on flat
tori and convergence of local covariances as \(L\to\infty\).
Separate the spatial zero mode before the limit, and state exactly
which local quantities survive although the global Fock trace
diverges.
\end{decision}

\begin{firewall}[Claim boundary after V2]
\label{fire:v16v2}
V2 is a free, reduced, compact-volume TT theorem.  It does not prove
the interacting gravitational constraints, the physical
infinite-volume continuum, reflection positivity of the combined
matter theory, or the full Standard Model--Einstein continuum.
\end{firewall}

\section*{Version V2 append-only record}
\addcontentsline{toc}{section}{Version V2 append-only record}
The complete sixteenth-cycle V1 source was retained before this
continuation.  It had \(16923\) bytes, \(401\) lines, and SHA--256
\[
 \texttt{5F92597B77DE53C38740FE54C8F783799F8927F24310CC7C5D428495F6507E71}.
\]
V2 skipped the already completed pole-residue programme, proved
compact-volume TT graviton Fock nuclearity and trace-norm removal of
the higher-spatial-derivative regulator, and isolated the global
zero-mode and thermodynamic-volume failures.
\section{Sixteenth-cycle V3: TT pressure density and local covariance}
\label{sec:v16v3}

V2 proves compact-volume TT Fock nuclearity and also proves that its
global trace is not uniform as the spatial volume grows.  The
thermodynamic replacement has two parts: the logarithm of the
partition function divided by volume, and local smeared covariance
functionals.  Both have massless limits in three spatial dimensions
after the exact global zero mode is separated.

\subsection{A Riemann-sum lemma with an infrared singularity}

\begin{lemma}[Improper lattice Riemann sums]
\label{lem:v16v3-riemann}
Let \(F:\mathbb R^3\setminus\{0\}\to\mathbb C\) be continuous and
suppose, for constants \(C,a>0\),
\[
 |F(k)|\leq C(1+|\log|k||)\quad(0<|k|\leq1),          \tag{3.1}
\]
\[
 |F(k)|\leq Ce^{-a|k|}\quad(|k|\geq1).                \tag{3.2}
\]
Then, as \(h\downarrow0\),
\[
 h^3\sum_{n\in\mathbb Z^3\setminus\{0\}}F(hn)
 \longrightarrow\int_{\mathbb R^3}F(k)\,dk.           \tag{3.3}
\]
The same conclusion holds when the exponential bound is replaced
by an integrable radial majorant whose lattice shell sums are
uniformly integrable.
\end{lemma}

\begin{proof}
Fix \(0<\delta<1<R\).  On the compact annulus
\(\delta\leq|k|\leq R\), ordinary lattice Riemann sums converge.
The integral of the bound in (3.1) over \(|k|<\delta\) is
\[
 O\!\left(\delta^3(1+|\log\delta|)\right).             \tag{3.4}
\]
Divide the lattice points in the same ball into shells
\(jh\leq|k|<(j+1)h\).  The number in the \(j\)-th shell is at most
\(C_0(j+1)^2\).  Multiplication by \(h^3\) and summation of the
logarithmic majorant gives the same bound as (3.4), uniformly for
small \(h\), plus a term tending to zero from the first shell.
The exponential tail beyond \(R\) is uniformly small by the
analogous shell count.  Let \(h\downarrow0\), then
\(\delta\downarrow0\) and \(R\uparrow\infty\).  The final statement
uses the same truncation argument with its stated uniform tail
hypothesis.
\end{proof}

\subsection{Free TT pressure}

Let \(\Sigma_L=(\mathbb R/L\mathbb Z)^3\) and
\[
 {\cal K}_L={2\pi\over L}\mathbb Z^3.                 \tag{3.5}
\]
Exclude \(k=0\), retain the two TT polarizations, and use the
dispersion
\[
 \omega_{c_4}(k)
 =\bigl(c_2|k|^2+c_4|k|^4+m^2\bigr)^{1/2},
 \quad c_2>0,\ c_4\geq0,\ m\geq0.                    \tag{3.6}
\]
At inverse temperature \(\beta>0\), define
\[
 \log Z_L(\beta)
 =-2\sum_{k\in{\cal K}_L\setminus\{0\}}
       \log\left(1-e^{-\beta\omega_{c_4}(k)}\right).   \tag{3.7}
\]

\begin{theorem}[Thermodynamic TT pressure]
\label{thm:v16v3-pressure}
For all parameters in (3.6), including \(m=0\),
\[
 \lim_{L\to\infty}{1\over L^3}\log Z_L(\beta)
 ={2\over(2\pi)^3}\int_{\mathbb R^3}
   -\log\left(1-e^{-\beta\omega_{c_4}(k)}\right)\,dk
 =:p_\beta(c_4,m)<\infty.                             \tag{3.8}
\]
Moreover,
\[
 p_\beta(c_4,m)\longrightarrow p_\beta(0,m)
 \quad\hbox{as }c_4\downarrow0.                       \tag{3.9}
\]
\end{theorem}

\begin{proof}
For \(m>0\), the integrand in (3.8) is bounded near zero.  For
\(m=0\),
\[
 \omega_{c_4}(k)\geq\sqrt{c_2}|k|,
\]
and
\[
 -\log(1-e^{-\beta\omega_{c_4}(k)})
 \leq C_\beta(1+|\log|k||)                            \tag{3.10}
\]
near zero.  At large momentum it has an exponential majorant because
\(\omega_{c_4}(k)\geq\sqrt{c_2}|k|\).  Lemma
\ref{lem:v16v3-riemann}, with \(h=2\pi/L\), gives (3.8), including
the factor
\[
 {1\over L^3}\sum_{k\in{\cal K}_L}
 ={1\over(2\pi)^3}h^3\sum_{n\in\mathbb Z^3}.          \tag{3.11}
\]

As \(c_4\downarrow0\), the integrand converges pointwise.  Since
\(\omega_{c_4}\geq\omega_0\) and
\(-\log(1-e^{-\beta x})\) decreases in \(x\), the \(c_4=0\)
integrand is an integrable majorant.  Dominated convergence proves
(3.9).
\end{proof}

The global trace \(Z_L\) diverges exponentially with volume, while
its logarithm has the finite density (3.8).  This is compatible with
Counterexample~\ref{cnt:v16v2-volume}.

\subsection{Local zero-temperature TT covariance}

Let \(P_{\rm TT}(k)\) be the orthogonal TT polarization projector for
\(k\ne0\); it is measurable and has operator norm one.  For a smooth
compactly supported symmetric-tensor test function \(f(\tau,x)\),
define its spatial Fourier transform \(\widehat f(\tau,k)\).
The infinite-volume free covariance is
\[
\begin{split}
 {\cal C}_\infty(f,g)
 ={1\over(2\pi)^3}\int_{\mathbb R^3}
 {dk\over2\omega_{c_4}(k)}
 \int_{\mathbb R^2}d\tau\,d\sigma\,
 e^{-\omega_{c_4}(k)|\tau-\sigma|}
 \langle\widehat f(\tau,k),
 P_{\rm TT}(k)\widehat g(\sigma,k)\rangle .
                                                               \tag{3.12}
\end{split}
\]
For \(m=0\), set the value at \(k=0\) arbitrarily; it is a
measure-zero point.

\begin{proposition}[Local covariance is finite]
\label{prop:v16v3-local-finite}
For \(f,g\) as above, (3.12) is absolutely convergent.  It defines a
tempered bilinear distribution.  In the massless case the infrared
majorant is \(C|k|^{-1}\), which is locally integrable in three
dimensions.
\end{proposition}

\begin{proof}
The time exponential is bounded by one and
\(\|P_{\rm TT}(k)\|\leq1\).  Compact time support bounds the two time
integrals by a fixed multiple of
\[
 {\sup_\tau|\widehat f(\tau,k)|
  \sup_\sigma|\widehat g(\sigma,k)|
  \over\omega_{c_4}(k)}.                              \tag{3.13}
\]
The Fourier transforms decay faster than every spatial power.
Near zero, \(\omega_{c_4}(k)\geq\sqrt{c_2}|k|\), and
\(\int_{|k|<1}|k|^{-1}dk<\infty\) in dimension three.  These bounds
prove absolute convergence and continuity in Schwartz seminorms.
\end{proof}

Let \(f_L\) be the \(L\)-periodization of a fixed compactly supported
\(f\), once \(L\) is larger than its spatial support diameter.  Let
\({\cal C}_L\) be the discrete version of (3.12), with
\((1/L^3)\sum_{k\ne0}\).

\begin{theorem}[Local covariance thermodynamic limit]
\label{thm:v16v3-local-limit}
For every fixed pair \(f,g\),
\[
 {\cal C}_L(f_L,g_L)\longrightarrow
 {\cal C}_\infty(f,g)                                 \tag{3.14}
\]
as \(L\to\infty\).  The same conclusion holds along
\(c_{4,L}\to0\).
\end{theorem}

\begin{proof}
Away from \(k=0\), the summand is continuous and rapidly decreasing.
Near zero it is bounded by \(C|k|^{-1}\), uniformly for
\(c_{4,L}\geq0\), because \(c_2\) is fixed positive.  Lattice shell
counting gives uniform integrability of this singularity.  Apply
the second form of Lemma~\ref{lem:v16v3-riemann}.  For
\(c_{4,L}\to0\), pointwise convergence and the same majorants give
the joint limit.
\end{proof}

\subsection{Reflection positivity survives the local limit}

\begin{theorem}[Free local TT reflection positivity]
\label{thm:v16v3-os}
Let \(f_1,\ldots,f_N\) be supported in positive Euclidean times and
let \(a_1,\ldots,a_N\in\mathbb C\).  Then the limiting covariance
satisfies
\[
 \sum_{i,j=1}^N\overline{a_i}a_j\,
 {\cal C}_\infty(\Theta f_i,f_j)\geq0.                 \tag{3.15}
\]
\end{theorem}

\begin{proof}
At every finite \(L\) and \(c_4\geq0\), the reflected form can be
written
\[
 {1\over L^3}\sum_{k\ne0}
 {1\over2\omega_{c_4}(k)}
 \left\|
 P_{\rm TT}(k)
 \int_0^\infty e^{-\omega_{c_4}(k)\tau}
       \sum_j a_j\widehat f_j(\tau,k)\,d\tau
 \right\|^2\geq0.                                    \tag{3.16}
\]
Theorem~\ref{thm:v16v3-local-limit} gives convergence of every entry
of the finite Gram matrix.  The cone of positive semidefinite
matrices is closed, so its limit satisfies (3.15).
\end{proof}

\subsection{Zero-mode scaling is observable dependent}

One may try to replace the removed zero mode by a small oscillator
frequency \(\varepsilon_L>0\).  Its two-polarization contribution
to the pressure density is
\[
 {2\over L^3}
 \bigl[-\log(1-e^{-\beta\varepsilon_L})\bigr].         \tag{3.17}
\]
Its equal-time thermal covariance contribution to an unsuppressed
spatial average is asymptotic to
\[
 {2\over\beta L^3\varepsilon_L^2}.                    \tag{3.18}
\]

\begin{proposition}[Two zero-mode removal scales]
\label{prop:v16v3-zero-scales}
The contribution (3.17) vanishes if
\[
 \log(\varepsilon_L^{-1})=o(L^3),                     \tag{3.19}
\]
whereas (3.18) vanishes only if
\[
 L^3\varepsilon_L^2\longrightarrow\infty.             \tag{3.20}
\]
Thus pressure convergence alone does not fix the local or averaged
zero-mode state.
\end{proposition}

\begin{proof}
Use
\(-\log(1-e^{-x})=-\log x+O(x)\) and
\(\coth(x/2)/(2\varepsilon)
\sim(\beta\varepsilon^2)^{-1}\) with
\(x=\beta\varepsilon\).
\end{proof}

For gravitational TT zero modes, an arbitrary choice of
\(\varepsilon_L\) would amount to choosing dynamics for global
metric moduli.  The collective-coordinate sector must determine
that choice.

\subsection{Handoff and next bottleneck}

\begin{assessment}[Physical status]
\label{ass:v16v3-status}
V3 proves a finite TT pressure density, convergence under removal
of the higher-spatial regulator, convergence of all compactly
smeared free TT covariances, and reflection positivity of their
infinite-volume limit.  These survive although the global Gibbs
trace does not.

The result remains free and constraint reduced.  It does not control
matter-sourced instantaneous constraints, nonlinear graviton
vertices, global moduli, or the interacting local state.
\end{assessment}

\begin{decision}[V4 sourced-constraint form test]
\label{dec:v16v3-next}
Audit the archived sourced-constraint estimates, then isolate a new
operator statement at the local-state level.  Test the mean-zero
Newton form
\(\langle\rho,(-\Delta)^{-1}\rho\rangle\) by
Hardy--Littlewood--Sobolev and Sobolev duality, determine the exact
integrability required of the renormalized Standard Model energy
density, and decide whether the attractive constraint term is a
uniform relative-form perturbation of the stable matter reference.
\end{decision}

\begin{firewall}[Claim boundary after V3]
\label{fire:v16v3}
V3 constructs only the free TT thermodynamic and local covariance
limits.  It does not construct the interacting infinite-volume
state, solve the sourced Einstein constraints, or establish the
full Standard Model--Einstein continuum.
\end{firewall}

\section*{Version V3 append-only record}
\addcontentsline{toc}{section}{Version V3 append-only record}
The complete V2 source was retained before this continuation.  It
had \(28496\) bytes, \(752\) lines, and SHA--256
\[
 \texttt{10EDDC037D75F48BD266DEC07C1057DA511EB1924481CA3451867A75C8D63F52}.
\]
V3 replaced the divergent global TT Gibbs trace by the finite
pressure density and local covariance limits, proved free
reflection positivity after the thermodynamic limit, and quantified
the independent zero-mode scaling conditions.
\section{Sixteenth-cycle V4: Gaussian Newton-source normalization}
\label{sec:v16v4}

The V3 decision proposed an \(L^{6/5}\) and
Hardy--Littlewood--Sobolev test for the sourced Newton form.  The
archive audit shows that this is already done: sixth-cycle
V62 proves the fixed-charge HLS bound and its \(N^3\) loss; V63
proves the finite-lattice copositivity and neutral-flux criteria;
and eighth-cycle V78 proves the \(L^{6/5}\to H^{-1}\) source
continuity needed for the conformal constraint.  The missing next
question is probabilistic.  A negative Newton term in the action
produces a positive quadratic exponential in the measure.  V4
computes exactly when that exponential is integrable for a Gaussian
source model and when the estimate is uniform under cutoff removal.

\subsection{Mean-zero Newton carrier}

Let \({\cal H}\) be a real separable Hilbert space of centered
constraint sources.  At compact volume, the mean-zero Newton
operator is modeled by a bounded nonnegative self-adjoint operator
\[
 G:{\cal H}\longrightarrow{\cal H}.                   \tag{4.1}
\]
For the scalar constraint on a compact slice,
\(G=P_0(-\Delta)^{-1}P_0\), where \(P_0\) removes the constant
mode.  Let \(X\) be a centered Gaussian \({\cal H}\)-valued random
variable with covariance \(C\geq0\), where \(C\) is trace class.
Define
\[
 A=C^{1/2}GC^{1/2}.                                   \tag{4.2}
\]
Then \(A\geq0\) is trace class.

\begin{theorem}[Sharp Gaussian Newton criterion]
\label{thm:v16v4-gaussian}
Let \(\Gamma>0\).  If \(A\ne0\), then
\[
 {\mathbb E}\exp\left\{
 {\Gamma\over2}\langle X,GX\rangle\right\}<\infty      \tag{4.3}
\]
if and only if
\[
 \Gamma\|A\|_{\rm op}<1.                              \tag{4.4}
\]
When (4.4) holds,
\[
 {\mathbb E}\exp\left\{
 {\Gamma\over2}\langle X,GX\rangle\right\}
 =\det(I-\Gamma A)^{-1/2}.                            \tag{4.5}
\]
If \(A=0\), the expectation equals one.
\end{theorem}

\begin{proof}
Write \(X=C^{1/2}\xi\), where \(\xi\) is the standard cylindrical
Gaussian interpreted through the trace-class map \(C^{1/2}\).
Diagonalize the positive compact operator \(A\):
\[
 A e_j=a_je_j,\qquad a_j\geq0,\qquad
 \sum_ja_j<\infty.                                    \tag{4.6}
\]
Finite-rank truncation gives
\[
\begin{split}
 {\mathbb E}\exp\left\{
 {\Gamma\over2}\sum_{j=1}^Na_j\xi_j^2\right\}
 &=\prod_{j=1}^N
 {\mathbb E}\exp\left\{{\Gamma a_j\over2}\xi_j^2\right\}\\
 &=\prod_{j=1}^N(1-\Gamma a_j)^{-1/2},                \tag{4.7}
\end{split}
\]
provided every \(\Gamma a_j<1\).  If one
\(\Gamma a_j\geq1\), the corresponding one-dimensional Gaussian
integral diverges.

If (4.4) holds, then
\[
 -\log(1-\Gamma a_j)
 \leq{\Gamma a_j\over1-\Gamma\|A\|_{\rm op}},          \tag{4.8}
\]
so the infinite product converges.  Monotone convergence passes
from (4.7) to (4.5).  Conversely, because a nonzero compact positive
operator attains its norm, failure of (4.4) gives a divergent
one-dimensional factor.
\end{proof}

The threshold (4.4) is the covariance-weighted constraint
susceptibility.  A small Newton coupling \(\Gamma\) is insufficient
unless the source covariance is controlled after composition with
the inverse constraint operator.

\subsection{Uniform cutoff stability}

Let \(C_n,G_n\) be transported to one source carrier and put
\[
 A_n=C_n^{1/2}G_nC_n^{1/2}.                            \tag{4.9}
\]

\begin{theorem}[Uniform Gaussian reweighting and convergence]
\label{thm:v16v4-uniform}
Assume
\[
 \|A_n-A\|_1\longrightarrow0,\qquad
 \sup_n\Gamma\|A_n\|_{\rm op}\leq q<1.                \tag{4.10}
\]
Then
\[
 \det(I-\Gamma A_n)^{-1/2}
 \longrightarrow
 \det(I-\Gamma A)^{-1/2},                             \tag{4.11}
\]
and all normalization factors are bounded by
\[
 \exp\left\{
 {\Gamma\over2(1-q)}\sup_n\operatorname{Tr}A_n
 \right\}.                                            \tag{4.12}
\]
More precisely,
\[
\left|
 \log{\det(I-\Gamma A_n)^{-1/2}
 \over\det(I-\Gamma A)^{-1/2}}
\right|
\leq{\Gamma\over2(1-q)}\|A_n-A\|_1.                  \tag{4.13}
\]
\end{theorem}

\begin{proof}
Join \(A\) to \(A_n\) by
\(A_s=(1-s)A+sA_n\).  Trace-norm convergence and (4.10) imply
\(\Gamma\|A\|\leq q\), so the same bound holds along the segment.
The Fredholm derivative formula gives
\[
 {d\over ds}\left[-{1\over2}\log\det(I-\Gamma A_s)\right]
 ={\Gamma\over2}
 \operatorname{Tr}\bigl((I-\Gamma A_s)^{-1}
                         (A_n-A)\bigr).               \tag{4.14}
\]
The inverse has norm at most \((1-q)^{-1}\).  Trace duality and
integration give (4.13), hence (4.11).  Applying (4.8) to every
\(A_n\) gives (4.12).
\end{proof}

\begin{corollary}[A sufficient factor convergence criterion]
\label{cor:v16v4-factors}
Suppose
\[
 C_n^{1/2}\longrightarrow C^{1/2}\quad\hbox{in
Hilbert--Schmidt norm},\qquad
 G_n\longrightarrow G\quad\hbox{in operator norm},    \tag{4.15}
\]
with uniformly bounded \(\|G_n\|\) and
\(\|C_n^{1/2}\|_2\).  Then \(A_n\to A\) in trace norm.
\end{corollary}

\begin{proof}
Insert and subtract
\(C^{1/2}G_nC_n^{1/2}\) and
\(C^{1/2}G C_n^{1/2}\).  Each difference is a product of two
Hilbert--Schmidt factors and one bounded factor.  The ideal estimate
\[
 \|RST\|_1\leq\|R\|_2\|S\|_{\rm op}\|T\|_2            \tag{4.16}
\]
proves convergence.
\end{proof}

\subsection{The critical eigenvector is a real instability}

\begin{proposition}[Approach to the critical susceptibility]
\label{prop:v16v4-critical}
If \(a_{1,n}=\|A_n\|_{\rm op}\) and
\(\Gamma a_{1,n}\uparrow1\), then
\[
 \det(I-\Gamma A_n)^{-1/2}
 \geq(1-\Gamma a_{1,n})^{-1/2}\longrightarrow\infty. \tag{4.17}
\]
The divergence is concentrated along the source direction
\(C_n^{1/2}e_{1,n}\) corresponding to the largest eigenvalue of
\(A_n\).
\end{proposition}

\begin{proof}
The Fredholm product contains
\((1-\Gamma a_{1,n})^{-1/2}\), while every other factor is at least
one.  The diagonalization in the proof of
Theorem~\ref{thm:v16v4-gaussian} identifies its Gaussian coordinate.
\end{proof}

Thus a determinant bound without the strict susceptibility margin
does not control the attractive constraint source.

\subsection{Why scalar subgaussianity is not enough}

\begin{counterexample}[Growing source dimension]
\label{cnt:v16v4-dimension}
Let \(X_n\) be a standard Gaussian vector in \(\mathbb R^n\) and
take \(G_n=I\).  Every scalar projection
\(\langle u,X_n\rangle\), \(\|u\|=1\), has the same
dimension-independent subgaussian bound.  Nevertheless, for
\(0<\Gamma<1\),
\[
 {\mathbb E}\exp\left\{{\Gamma\over2}\|X_n\|^2\right\}
 =(1-\Gamma)^{-n/2}\longrightarrow\infty.             \tag{4.18}
\]
\end{counterexample}

\begin{proof}
Independence of the \(n\) standard Gaussian coordinates gives the
product.  The one-dimensional projection statement is immediate
from rotational invariance.
\end{proof}

Uniform subgaussian estimates for each smeared energy density
therefore do not replace trace control of the complete composed
operator \(C^{1/2}GC^{1/2}\).

\subsection{Normalized source law}

When (4.4) holds, the reweighted Gaussian remains Gaussian.  Its
covariance on the cylindrical carrier is
\[
 (I-\Gamma A)^{-1},                                   \tag{4.19}
\]
and the physical source covariance is
\[
 C_\Gamma
 =C^{1/2}(I-\Gamma A)^{-1}C^{1/2}.                    \tag{4.20}
\]

\begin{proposition}[Susceptibility amplification]
\label{prop:v16v4-amplification}
Under the uniform margin \(\Gamma\|A\|\leq q<1\),
\[
 0\leq C_\Gamma\leq{1\over1-q}C                      \tag{4.21}
\]
as quadratic forms, and
\[
 \operatorname{Tr}C_\Gamma
 \leq{1\over1-q}\operatorname{Tr}C.                   \tag{4.22}
\]
\end{proposition}

\begin{proof}
The spectral bound gives
\[
 I\leq(I-\Gamma A)^{-1}\leq(1-q)^{-1}I.
\]
Sandwich by \(C^{1/2}\) and take traces.
\end{proof}

The attractive constraint increases source fluctuations and becomes
singular at the same threshold that destroys normalization.

\subsection{Physical scope}

The Standard Model energy density is a renormalized composite and
is not Gaussian.  V4 is therefore a sharp diagnostic and a complete
theorem for a Gaussian reference or linearized source sector, not
a proof for the interacting source.  A valid nonlinear replacement
must control the quadratic exponential of its full
\(H^{-1}\)-valued source, uniformly in cutoff and volume.

\begin{assessment}[Physical status]
\label{ass:v16v4-status}
V4 proves the exact Gaussian normalization threshold for the
attractive Newton reweighting, its Fredholm determinant, uniform
trace-ideal convergence, and the amplified covariance.  It shows
that one-dimensional subgaussian source estimates miss an
independent growing-rank obstruction.

No Gaussianity claim is made for the renormalized Standard Model
energy density, and no nonlinear constraint measure or reflection
positivity theorem follows.
\end{assessment}

\begin{decision}[V5 companion audit]
\label{dec:v16v4-next}
Perform the compulsory YM/NS audit.  Test whether either current
companion supplies a uniform trace/covariance input for the
susceptibility operator (4.2).  Then continue in V6 by replacing
Gaussianity with an explicit quadratic-exponential hypothesis and
deriving a stable change-of-measure theorem for local observables.
\end{decision}

\begin{firewall}[Claim boundary after V4]
\label{fire:v16v4}
V4 proves a Gaussian source theorem.  It does not establish the
required exponential moment for the interacting Standard Model
stress tensor, solve the Einstein constraint, or construct the full
continuum.
\end{firewall}

\section*{Version V4 append-only record}
\addcontentsline{toc}{section}{Version V4 append-only record}
The complete V3 source was retained before this continuation.  It
had \(39114\) bytes, \(1059\) lines, and SHA--256
\[
 \texttt{8F3649F97108687A0124DF772D572BFBA8967F2B9F341686C6AC94A01D8E2A4B}.
\]
V4 skipped the archived HLS and copositivity work and proved the
sharp Gaussian Newton-source normalization and cutoff-convergence
criterion in terms of the composed susceptibility operator.
\section{Sixteenth-cycle V5: compulsory companion audit}
\label{sec:v16v5}

This is the first five-version checkpoint of the sixteenth cycle.
The audit asks whether the current Yang--Mills or Navier--Stokes
notes supply a uniform covariance or trace input for the
Newton-source susceptibility operator constructed in V4.

\subsection{Frozen companion checkpoints}

The inspected Yang--Mills file is
\[
 \texttt{Renewal\_Yang\_Mills\_Mass\_Gap\_Research\_Notes\_V44.tex},
                                                               \tag{5.1}
\]
with \(379692\) bytes, \(10394\) lines, and SHA--256
\[
 \texttt{D07E2A074FC8A8BA81BDBCDDA9EB9E6BCFEAA530BBD2C367A668532DB3EB939D}.
                                                               \tag{5.2}
\]
The inspected Navier--Stokes file is
\[
 \texttt{Renewal\_NS\_Stability\_Research\_Notes\_V26.tex},     \tag{5.3}
\]
with \(278283\) bytes, \(5319\) lines, and SHA--256
\[
 \texttt{82C887F8C2C69E4F28FAD7C3DC8DE5B9099CB5A4BF431EBA1FFF3E91935978AD}.
                                                               \tag{5.4}
\]

\subsection{Yang--Mills V42--V44}

Since the last audit, YM V42 rejects elementary orientation
intertwiners, records an exact failed isotropic certificate, repairs
one corner with a smaller rational product hull, and proves both
direct response signs on paired first cubes.  V43 extends this
independent opposite atlas to paired two-step prisms.  V44 reaches
paired first-coordinate quarter-face plates.  Every positive chart
claim is backed by exact rational vertex certificates.

The final V44 region still has thickness only \(t\) in two
transverse directions.  Its firewall withholds a complete quarter
cell, full-torus positivity, a global covariant inequality, a
cofinal response sign, an interacting continuum measure, reflection
positivity, clustering, and a Hamiltonian mass gap.

\begin{proposition}[YM non-import at the susceptibility gate]
\label{prop:v16v5-ym}
YM V44 supplies no operator \(C\), \(G\), or composed trace-class
susceptibility
\[
 A=C^{1/2}GC^{1/2}                                    \tag{5.5}
\]
for the Standard Model energy density in the gravitational
constraint.  Its local two-sign pencil floors therefore imply
neither (4.4) nor (4.10).
\end{proposition}

\begin{proof}
The V44 objects are finite \(45\times45\) Wilson response pencils
certified on selected transverse momentum plates.  The \(C\) in
(5.5) is the covariance of the complete renormalized matter energy
density, while \(G\) is the mean-zero gravitational constraint
inverse.  V44 constructs neither object and explicitly has no
full-torus or continuum-operator conclusion.  The hypotheses of
V4 are therefore absent.
\end{proof}

The exact YM repair is still methodologically relevant: a failed
orientation identity was replaced by independent local
certificates rather than a symmetry assumption.  No numerical
floor transfers between the two programmes.

\subsection{Navier--Stokes V26}

NS V26 remains the rank-one Oseen-kernel derivative calculation
already audited in fifteenth-cycle V90 and V95.  It has not acquired
a new source covariance, gravitational inverse, quantum heat
factor, or exponential moment.

\begin{theorem}[V5 audit conclusion]
\label{thm:v16v5-audit}
Neither companion checkpoint supplies a hypothesis of
Theorem~\ref{thm:v16v4-uniform}.  No companion result is imported
into the V4 susceptibility estimate.
\end{theorem}

\begin{proof}
Use Proposition~\ref{prop:v16v5-ym} for YM.  The unchanged NS result
is a deterministic classical kernel norm on a restricted
rank-one input family and has none of the operators or probability
laws appearing in V4.
\end{proof}

\begin{assessment}[Direction after the audit]
\label{ass:v16v5-direction}
The Gaussian theorem remains a diagnostic for the full source.
The next useful result should not assume that the interacting energy
density is Gaussian.  It should state the exact
quadratic-exponential moment that replaces Gaussianity and prove
what that moment buys for normalization, tightness, local
observables, and cutoff passage.
\end{assessment}

\begin{decision}[V6 nonlinear Newton tilt]
\label{dec:v16v5-next}
Let \(Q_n\geq0\) be the complete cutoff Newton quadratic of the
centered Standard Model source.  Assume a uniform exponential
moment at a strictly larger coupling than the physical one.  Prove
uniform integrability of the Newton weights, convergence of tilted
local states from joint weak convergence, an entropy bound, and a
tightness-transfer inequality.  Keep reflection positivity as a
separate transfer-kernel condition.
\end{decision}

\begin{firewall}[Claim boundary after V5]
\label{fire:v16v5}
V5 is a read-only companion audit.  It imports no Yang--Mills mass
gap, Navier--Stokes theorem, Newton-source covariance, or continuum
construction.
\end{firewall}

\section*{Version V5 append-only record}
\addcontentsline{toc}{section}{Version V5 append-only record}
The complete V4 source was retained before this continuation.  It
had \(49094\) bytes, \(1361\) lines, and SHA--256
\[
 \texttt{D6D22B94AEE37C507ABD3EA90D0937CCA8304264A106623D8146F8FBB57AEE92}.
\]
V5 inspected YM V44 and NS V26, recorded their exact checkpoints,
imported no theorem into the Newton susceptibility, and selected
the nonlinear exponential-tilt theorem for V6.
\section{Sixteenth-cycle V6: nonlinear Newton tilts and stable local states}
\label{sec:v16v6}

V4 computes the exact attractive-Newton normalization when the
constraint source is Gaussian.  The interacting Standard Model
source is not Gaussian.  V6 replaces that model assumption by the
precise nonlinear input needed for passage to the limit: a uniform
quadratic-exponential moment at a coupling strictly larger than the
physical one.

\subsection{Joint source laws}

Let \(E\) be a Polish local-configuration space.  At cutoff \(n\),
let \(\lambda_n\) be the joint law on
\[
 E\times[0,\infty)
\]
of a local configuration \(X_n\) and its centered Newton quadratic
\[
 Q_n
 =\langle P_{0,n}\rho_n,G_nP_{0,n}\rho_n\rangle
 \geq0.                                               \tag{6.1}
\]
Write
\[
 a={\Gamma\over2}>0,\qquad
 Z_n=\int e^{aq}\,d\lambda_n(x,q),                    \tag{6.2}
\]
and define the tilted joint law
\[
 d\widehat\lambda_n
 ={e^{aq}\over Z_n}\,d\lambda_n.                      \tag{6.3}
\]
Its \(E\)-marginal is the Newton-tilted local state \(\nu_n\).

\begin{definition}[Newton exponential margin]
\label{def:v16v6-margin}
The joint laws have a Newton exponential margin if there are
\(\delta>0\) and \(M<\infty\) such that
\[
 \sup_n\int e^{(1+\delta)aq}\,d\lambda_n(x,q)
 \leq M.                                              \tag{6.4}
\]
\end{definition}

\subsection{Weak convergence after the nonlinear tilt}

\begin{theorem}[Stable Newton change of measure]
\label{thm:v16v6-weak}
Assume
\[
 \lambda_n\Longrightarrow\lambda
 \quad\hbox{weakly on }E\times[0,\infty)              \tag{6.5}
\]
and the margin (6.4).  Define \(Z\) and
\(\widehat\lambda\) from \(\lambda\) by (6.2)--(6.3).
Then
\[
 Z_n\longrightarrow Z,\qquad 1\leq Z\leq M^{1/(1+\delta)},          \tag{6.6}
\]
\[
 \widehat\lambda_n\Longrightarrow\widehat\lambda,     \tag{6.7}
\]
and consequently
\[
 \nu_n\Longrightarrow\nu                             \tag{6.8}
\]
on \(E\).
\end{theorem}

\begin{proof}
The higher exponential moment makes the weights
\(W_n=e^{aq}\) uniformly integrable.  Indeed, for \(R>0\),
\[
\begin{split}
 \int_{\{W_n>R\}}W_n\,d\lambda_n
 &\leq R^{-\delta}
       \int W_n^{1+\delta}\,d\lambda_n\\
 &\leq MR^{-\delta}.                                  \tag{6.9}
\end{split}
\]
For a bounded continuous \(F\) on the product space, truncate the
continuous weight by \(W^{(R)}=\min(e^{aq},R)\).
Weak convergence gives
\[
 \int F W^{(R)}\,d\lambda_n
 \longrightarrow\int F W^{(R)}\,d\lambda.             \tag{6.10}
\]
Portmanteau applied to the nonnegative lower-semicontinuous
function \(e^{(1+\delta)aq}\) shows that the limit law has the same
moment bound \(M\).  Thus (6.9) controls both tails, and letting
\(R\to\infty\) in (6.10) gives convergence with the full weight.
Taking \(F=1\) proves \(Z_n\to Z\).  Since \(q\geq0\), \(Z_n,Z\geq1\);
Hölder gives the upper bound in (6.6).

Dividing the weighted convergence by \(Z_n\to Z>0\) proves (6.7).
Projection to \(E\) is continuous, giving (6.8).
\end{proof}

\begin{corollary}[Convergence of bounded local observables]
\label{cor:v16v6-observables}
For every bounded continuous local observable \(F:E\to\mathbb C\),
\[
 \int_EF\,d\nu_n\longrightarrow\int_EF\,d\nu.         \tag{6.11}
\]
\end{corollary}

\begin{proof}
This is (6.8).
\end{proof}

\subsection{Tightness and entropy transfer}

\begin{proposition}[Quantitative tightness transfer]
\label{prop:v16v6-tightness}
For every Borel set \(B\subset E\),
\[
 \nu_n(B)
 \leq M^{1/(1+\delta)}
       \mu_n(B)^{\delta/(1+\delta)},                  \tag{6.12}
\]
where \(\mu_n\) is the \(E\)-marginal of \(\lambda_n\).
Consequently, uniform tightness of \((\mu_n)\) implies uniform
tightness of \((\nu_n)\).
\end{proposition}

\begin{proof}
Hölder with exponents \(1+\delta\) and
\((1+\delta)/\delta\) gives
\[
 \int_{\{x\in B\}}e^{aq}\,d\lambda_n
 \leq
 \left(\int e^{(1+\delta)aq}\,d\lambda_n\right)^{1/(1+\delta)}
 \mu_n(B)^{\delta/(1+\delta)}.                        \tag{6.13}
\]
Divide by \(Z_n\geq1\) and use (6.4).  If a compact \(K\) makes
\(\sup_n\mu_n(K^c)\) small, (6.12) does the same for \(\nu_n\).
\end{proof}

\begin{theorem}[Uniform entropy debit]
\label{thm:v16v6-entropy}
The tilted joint law satisfies
\[
 D(\widehat\lambda_n\|\lambda_n)
 \leq{\log M\over\delta}.                             \tag{6.14}
\]
The marginal entropy obeys
\[
 D(\nu_n\|\mu_n)
 \leq{\log M\over\delta}.                             \tag{6.15}
\]
\end{theorem}

\begin{proof}
Set
\[
 \phi_n(s)=\log\int e^{sq}\,d\lambda_n.
\]
The higher moment makes \(\phi_n\) finite on
\([0,(1+\delta)a]\) and differentiable at \(a\).  Directly from
(6.3),
\[
 D(\widehat\lambda_n\|\lambda_n)
 =a\phi_n'(a)-\phi_n(a).                              \tag{6.16}
\]
Convexity gives
\[
 \phi_n((1+\delta)a)
 \geq\phi_n(a)+\delta a\phi_n'(a).                    \tag{6.17}
\]
Since \(q\geq0\), \(\phi_n(a)\geq0\).  Combining
(6.16)--(6.17) yields
\[
 D(\widehat\lambda_n\|\lambda_n)
 \leq{\phi_n((1+\delta)a)\over\delta}
 \leq{\log M\over\delta}.                             \tag{6.18}
\]
Relative entropy decreases under the projection
\((x,q)\mapsto x\), proving (6.15).
\end{proof}

This entropy estimate connects the Newton constraint to the earlier
entropy-tightness compilers without assuming a log-Sobolev
inequality for the tilted law.

\subsection{Total-variation upgrade}

\begin{theorem}[Strong stability of the tilt]
\label{thm:v16v6-tv}
Suppose \(\lambda_n,\lambda\) are measures on one common product
space and
\[
 \|\lambda_n-\lambda\|_{\rm TV}\longrightarrow0,      \tag{6.19}
\]
while both the sequence and the limit obey (6.4).  Then
\[
 \|\widehat\lambda_n-\widehat\lambda\|_{\rm TV}
 \longrightarrow0,\qquad
 \|\nu_n-\nu\|_{\rm TV}\longrightarrow0.              \tag{6.20}
\]
\end{theorem}

\begin{proof}
Fix \(R>0\) and split the weight at \(q=R\).  On \(q\leq R\),
\[
 \left\|e^{aq}{\bf1}_{q\leq R}
       (\lambda_n-\lambda)\right\|_{\rm TV}
 \leq e^{aR}\|\lambda_n-\lambda\|_{\rm TV}.           \tag{6.21}
\]
On \(q>R\),
\[
 e^{aq}\leq e^{-\delta aR}e^{(1+\delta)aq},           \tag{6.22}
\]
so the two weighted tails have total mass at most
\(Me^{-\delta aR}\).  First choose \(R\) large and then \(n\)
large.  The unnormalized weighted measures converge in total
variation, as do their masses \(Z_n\).  Division by
\(Z_n,Z\geq1\) proves the first statement.  Total variation
decreases under projection, proving the marginal statement.
\end{proof}

\subsection{Stability under an approximate source quadratic}

\begin{proposition}[Replacement of \(Q_n\)]
\label{prop:v16v6-replacement}
Let \(Q_n,\widetilde Q_n\geq0\) live on one probability space.
Let \(p>1\), \(p'=p/(p-1)\), and assume
\[
 \|Q_n-\widetilde Q_n\|_{L^p}\longrightarrow0,        \tag{6.23}
\]
\[
 \sup_n{\mathbb E}
 \left[e^{p'aQ_n}+e^{p'a\widetilde Q_n}\right]<\infty. \tag{6.24}
\]
Then
\[
 \|e^{aQ_n}-e^{a\widetilde Q_n}\|_{L^1}
 \longrightarrow0.                                   \tag{6.25}
\]
Their normalized tilted laws therefore have vanishing total
variation distance.
\end{proposition}

\begin{proof}
For \(x,y\geq0\), the mean-value theorem gives
\[
 |e^{ax}-e^{ay}|
 \leq a|x-y|(e^{ax}+e^{ay}).                          \tag{6.26}
\]
Apply Hölder with \(p,p'\), use (6.23)--(6.24), and obtain (6.25).
The normalization constants converge by the same estimate and are
at least one.  The total-variation claim follows from the
\(L^1\) distance between the normalized densities.
\end{proof}

\subsection{Why the strict exponential margin is necessary}

\begin{counterexample}[Rare source event captures tilted mass]
\label{cnt:v16v6-rare}
Let \(E=\mathbb R\).  Under \(\lambda_n\), set
\[
 (X_n,Q_n)=
 \begin{cases}
 (n,a^{-1}\log n),&\hbox{with probability }n^{-1},\\
 (0,0),&\hbox{with probability }1-n^{-1}.
 \end{cases}                                         \tag{6.27}
\]
Then \((X_n,Q_n)\to(0,0)\) in probability and
\[
 Z_n=2-{1\over n}\longrightarrow2.                   \tag{6.28}
\]
Nevertheless,
\[
 \nu_n(\{n\})={1\over2-1/n}\longrightarrow{1\over2},  \tag{6.29}
\]
so the tilted laws are not tight and do not converge to
\(\delta_0\).  For every \(\delta>0\),
\[
 {\mathbb E}e^{(1+\delta)aQ_n}
 =1-{1\over n}+n^\delta\longrightarrow\infty.         \tag{6.30}
\]
\end{counterexample}

\begin{proof}
On the rare event, \(e^{aQ_n}=n\), which cancels its probability
\(1/n\) and leaves order-one tilted mass.  The displayed formulas
follow directly.
\end{proof}

Thus bounded normalization at exactly the physical exponent and
convergence of \(Q_n\) in probability do not provide uniform
integrability.  The strict margin in (6.4) pays for rare source
events.

\subsection{Reflection-positivity boundary}

The weight \(e^{aQ_n}\) is positive, but \(Q_n\) is spatially
nonlocal and may couple time slices through the constraint
solution.  Positivity of a scalar density does not automatically
preserve Osterwalder--Schrader positivity.  At a finite cutoff, a
same-slice positive multiplier \(B_n\) can preserve a positive
one-step transfer through the sandwich
\[
 T_n^{\rm new}=B_nT_nB_n\geq0,                        \tag{6.31}
\]
but the actual constraint weight must first be proved to have this
reflection-local factorization and the sandwich must retain the
trace and domain bounds.

\begin{assessment}[Physical status]
\label{ass:v16v6-status}
V6 proves normalization, weak and total-variation convergence,
tightness transfer, entropy control, and approximation stability
for a nonlinear Newton tilt under one explicit stronger
exponential moment.  It gives a sharp rare-event failure when the
margin is absent.

The complete renormalized Standard Model source has not been shown
to satisfy (6.4), and reflection positivity of its constraint
weight remains separate.
\end{assessment}

\begin{decision}[V7 metric derivative of the Newton tilt]
\label{dec:v16v6-next}
Differentiate the normalized tilted state along a metric path.
Separate the base score, the derivative of the centered source, the
derivative of the Newton inverse, and the zero-mode projection.
Prove passage of that derivative under a uniform exponential margin
and an \(L^p\) derivative bound, so the constraint contribution to
the Einstein source is explicit rather than absorbed into a vacuum
normalization.
\end{decision}

\begin{firewall}[Claim boundary after V6]
\label{fire:v16v6}
V6 is a conditional change-of-measure theorem.  It does not prove
the exponential margin for the interacting source, reflection
positivity of the Newton reweighting, or the full
Standard Model--Einstein continuum.
\end{firewall}

\section*{Version V6 append-only record}
\addcontentsline{toc}{section}{Version V6 append-only record}
The complete V5 source was retained before this continuation.  It
had \(54405\) bytes, \(1493\) lines, and SHA--256
\[
 \texttt{7FBA064407986B02CEC4451931FDBA1099DF386B5A2F6B6E12AE420CA8A627A6}.
\]
V6 replaced Gaussianity by a strict nonlinear
quadratic-exponential margin and proved stable Newton-tilted local
states, entropy and tightness control, a total-variation upgrade,
and the rare-event obstruction at the critical exponent.
\section{Sixteenth-cycle V7: metric response of the centered Newton tilt}
\label{sec:v16v7}

V6 proves convergence of nonlinear Newton-tilted local states.
The Einstein equation needs more: differentiation with respect to
the source metric.  The centered source, the inverse constraint
operator, and its zero-mode projection all move.  V7 derives the
complete derivative and gives sufficient uniform-integrability
conditions for passing it to the continuum limit.

\subsection{Normalized response identity}

Let \(s\mapsto\mu_s\) be a differentiable family of probability
measures on a common configuration space, with logarithmic score
\[
 B_s={d\over ds}\log{d\mu_s\over d\mu_{s_0}},\qquad
 {\mathbb E}_{\mu_s}B_s=0.                            \tag{7.1}
\]
Let \(Q_s\geq0\) be differentiable and define
\[
 Z_s=\int e^{aQ_s}\,d\mu_s,\qquad
 d\nu_s=Z_s^{-1}e^{aQ_s}\,d\mu_s,\qquad a={\Gamma\over2}.          \tag{7.2}
\]

\begin{theorem}[Newton-tilt response formula]
\label{thm:v16v7-response}
Let \(F_s\) be a differentiable observable.  Whenever differentiation
under the integrals is justified,
\[
 {d\over ds}\log Z_s
 ={\mathbb E}_{\nu_s}\bigl[B_s+a\dot Q_s\bigr],        \tag{7.3}
\]
\[
 {d\over ds}{\mathbb E}_{\nu_s}F_s
 ={\mathbb E}_{\nu_s}\dot F_s
 +\operatorname{Cov}_{\nu_s}
   \bigl(F_s,B_s+a\dot Q_s\bigr).                     \tag{7.4}
\]
\end{theorem}

\begin{proof}
Differentiate the numerator and denominator in
\[
 {\mathbb E}_{\nu_s}F_s
 ={ {\mathbb E}_{\mu_s}(F_se^{aQ_s})\over Z_s}.        \tag{7.5}
\]
The derivative of an expectation under \(\mu_s\) is the expectation
of the explicit derivative plus multiplication by \(B_s\).
Taking \(F_s=1\) gives (7.3).  The quotient rule then subtracts
\({\mathbb E}_{\nu_s}F_s\,
 {\mathbb E}_{\nu_s}(B_s+a\dot Q_s)\), giving (7.4).
\end{proof}

The constant part of the effective score cancels from (7.4), but
not from the partition derivative (7.3).  This is the normalized
state versus Einstein-source distinction proved in V1.

\subsection{Derivative of the centered constraint quadratic}

Work on a fixed transported Hilbert carrier.  Let \(L_s\geq0\) be
a differentiable self-adjoint constraint operator with a
constant-dimensional kernel.  Let
\[
 K_s=\hbox{projection onto }\ker L_s,\qquad
 P_s=I-K_s,                                           \tag{7.6}
\]
and let \(G_s=L_s^+\) be the Moore--Penrose inverse:
\[
 L_sG_s=G_sL_s=P_s,\qquad
 G_sP_s=P_sG_s=G_s.                                   \tag{7.7}
\]
For a differentiable energy-density source \(\rho_s\), put
\[
 r_s=P_s\rho_s,\qquad
 Q_s=\langle r_s,G_sr_s\rangle.                       \tag{7.8}
\]

\begin{proposition}[Moving pseudoinverse identity]
\label{prop:v16v7-pseudoinverse}
Under operator-norm differentiability and a uniform positive gap
of \(L_s\) on \(\operatorname{Ran}P_s\),
\[
 \dot G_s
 =-G_s\dot L_sG_s+\dot P_sG_s+G_s\dot P_s.            \tag{7.9}
\]
\end{proposition}

\begin{proof}
The gap makes the Riesz projection \(P_s\) and reduced inverse
differentiable.  Differentiate the identities (7.7).  Projection
differentiation gives
\[
 P_s\dot P_sP_s=0,\qquad K_s\dot P_sK_s=0.            \tag{7.10}
\]
Substitution of the right side of (7.9) into the derivatives of
\(L_sG_s=P_s\), \(G_sL_s=P_s\), and
\(G_sP_s=G_s\) verifies all four identities.  The differentiated
identities determine the four blocks of \(\dot G_s\) relative to
\(P_s\oplus K_s\), proving uniqueness and (7.9).
\end{proof}

\begin{theorem}[Complete centered Newton derivative]
\label{thm:v16v7-Q}
The derivative of (7.8) is
\[
 \dot Q_s
 =2\operatorname{Re}\langle\dot r_s,G_sr_s\rangle
  +\langle r_s,\dot G_sr_s\rangle,                    \tag{7.11}
\]
where
\[
 \dot r_s=\dot P_s\rho_s+P_s\dot\rho_s.               \tag{7.12}
\]
Equivalently, using \(P_sr_s=r_s\),
\[
\begin{split}
 \dot Q_s
 ={}&2\operatorname{Re}
 \langle\dot P_s\rho_s+P_s\dot\rho_s,G_sr_s\rangle\\
 &-\langle G_sr_s,\dot L_sG_sr_s\rangle
 +2\operatorname{Re}\langle r_s,\dot P_sG_sr_s\rangle.
                                                               \tag{7.13}
\end{split}
\]
\end{theorem}

\begin{proof}
Differentiate the quadratic form (7.8) to get (7.11), differentiate
\(r_s=P_s\rho_s\) to get (7.12), and insert
Proposition~\ref{prop:v16v7-pseudoinverse}.  Self-adjointness
combines the two projection terms into the final real part.
\end{proof}

If the transported carrier and its inner product also depend on
\(s\), their derivative must first be included in \(B_s\), in
\(\dot L_s\), or in an explicit inner-product term.  Formula (7.13)
is stated after unitary transport to a fixed carrier.

\begin{proposition}[A usable derivative bound]
\label{prop:v16v7-bound}
Suppose
\[
 \|G_s\|_{\rm op}\leq g_*,\quad
 \|\dot G_s\|_{\rm op}\leq \dot g_*,\quad
 \|\dot P_s\|_{\rm op}\leq p_*.
\]
Then
\[
 |\dot Q_s|
 \leq
 2g_*\|r_s\|
   \bigl(p_*\|\rho_s\|+\|P_s\dot\rho_s\|\bigr)
 +\dot g_*\|r_s\|^2.                                  \tag{7.14}
\]
\end{proposition}

\begin{proof}
Apply Cauchy--Schwarz to (7.11), use (7.12), and insert the operator
norm bounds.
\end{proof}

\subsection{Why the moving projection cannot be omitted}

\begin{counterexample}[Rotating kernel]
\label{cnt:v16v7-rotating}
On \(\mathbb R^2\), let
\[
 e_s=(\cos s,\sin s),\qquad P_s=|e_s\rangle\langle e_s|,
 \qquad L_s=P_s.                                      \tag{7.15}
\]
Then \(\ker L_s=e_s^\perp\), \(G_s=P_s\), and
\[
 \dot G_s=\dot P_s\ne0.                               \tag{7.16}
\]
The truncated formula
\(-G_s\dot L_sG_s\) gives zero because
\[
 P_s\dot P_sP_s=0.                                    \tag{7.17}
\]
\end{counterexample}

\begin{proof}
All statements follow by differentiating the rank-one projection.
The two projection terms in (7.9) sum to \(\dot P_s\), whereas the
middle inverse term vanishes.
\end{proof}

Thus even constant kernel dimension does not permit a fixed
mean-zero projection unless the carrier transport proves it.

\subsection{Passage of the response through cutoff removal}

For cutoff \(n\), let
\[
 R_{n,s}=B_{n,s}+a\dot Q_{n,s}.                       \tag{7.18}
\]

\begin{theorem}[Uniform derivative passage]
\label{thm:v16v7-passage}
Fix \(s\).  Suppose the joint laws of
\[
 (X_{n,s},Q_{n,s},F_{n,s},\dot F_{n,s},R_{n,s})       \tag{7.19}
\]
converge weakly to the corresponding limiting joint law.  Assume
that for some \(\delta>0\),
\[
\sup_n{\mathbb E}_{\mu_{n,s}}
\left[
 e^{(1+\delta)aQ_{n,s}}
 \bigl(1+|\dot F_{n,s}|^{1+\delta}
          +|R_{n,s}|^{1+\delta}\bigr)
\right]<\infty,                                      \tag{7.20}
\]
and that \(F_{n,s}\) are uniformly bounded and converge in the
joint law.  Then
\[
 {\mathbb E}_{\nu_{n,s}}\dot F_{n,s}
 \longrightarrow{\mathbb E}_{\nu_s}\dot F_s,         \tag{7.21}
\]
\[
 \operatorname{Cov}_{\nu_{n,s}}(F_{n,s},R_{n,s})
 \longrightarrow
 \operatorname{Cov}_{\nu_s}(F_s,R_s).                \tag{7.22}
\]
Hence the right side of (7.4) passes to the limit.
\end{theorem}

\begin{proof}
The bound (7.20) makes
\[
 e^{aQ_{n,s}}\dot F_{n,s},\qquad
 e^{aQ_{n,s}}R_{n,s},\qquad
 e^{aQ_{n,s}}F_{n,s}R_{n,s}                          \tag{7.23}
\]
uniformly integrable by the de la Vallée--Poussin criterion, since
their \(1+\delta\) moments are bounded up to the uniform bound on
\(F_{n,s}\).  Truncate these variables, apply weak convergence to
the bounded continuous truncations, and remove the truncation by
uniform integrability.  V6 gives convergence of the normalization
constants.  Dividing by them proves convergence of each expectation
forming (7.21)--(7.22).
\end{proof}

\begin{corollary}[Derivative of the limiting local state]
\label{cor:v16v7-limit-derivative}
Assume the hypotheses of
Theorem~\ref{thm:v16v7-passage} hold locally uniformly for
\(s\) in an interval, the tilted expectations converge pointwise
there, and the right sides of (7.4) converge locally uniformly.
Then the limiting expectation is differentiable and
\[
 {d\over ds}{\mathbb E}_{\nu_s}F_s
 ={\mathbb E}_{\nu_s}\dot F_s
 +\operatorname{Cov}_{\nu_s}(F_s,R_s).                \tag{7.24}
\]
\end{corollary}

\begin{proof}
Apply the fundamental theorem of calculus to every cutoff
expectation on a compact subinterval.  Local uniform convergence of
the derivatives permits passage under the \(s\)-integral.  The
pointwise limit then satisfies the same integral identity and hence
has derivative (7.24).
\end{proof}

\subsection{Gaussian consistency check}

For the centered Gaussian model of V4,
\[
 Z_s=\det(I-\Gamma A_s)^{-1/2}.
\]
If \(A_s\) is trace-norm differentiable and
\(\Gamma\|A_s\|<1\), then
\[
 {d\over ds}\log Z_s
 ={\Gamma\over2}
 \operatorname{Tr}\bigl(
 (I-\Gamma A_s)^{-1}\dot A_s\bigr).                  \tag{7.25}
\]

\begin{proof}
Differentiate the Fredholm logarithm.  The strict susceptibility
margin bounds the inverse, and trace-norm differentiability
justifies the trace.
\end{proof}

Equation (7.25) agrees with (7.3) by Gaussian integration by parts.
It gives a finite-dimensional or linearized check on the full score
ledger, but does not replace (7.13) for the interacting source.

\subsection{Vacuum and source ledger}

If the base action contains a metric-dependent scalar vacuum term
\({\cal E}_n(s)\), its score contributes
\[
 -t{\cal E}_n'(s)                                     \tag{7.26}
\]
to \(B_{n,s}\).  It cancels from the covariance in the normalized
observable derivative (7.4), while remaining in
\(d\log Z_{n,s}/ds\).  The Newton contribution \(a\dot Q_{n,s}\)
contains four distinct rows:
\[
 \dot\rho_{n,s},\qquad
 \dot P_{n,s},\qquad
 \dot L_{n,s},\qquad
 \hbox{carrier/inner-product transport}.              \tag{7.27}
\]
Dropping any row changes the candidate Einstein source.

\begin{assessment}[Physical status]
\label{ass:v16v7-status}
V7 proves the exact normalized response of the nonlinear Newton
tilt, the moving-pseudoinverse formula, the complete centered
source derivative, and sufficient exponential-weighted conditions
for passing that response to the continuum limit.

The required bounds on the renormalized Standard Model stress
derivative, the moving constraint projection, and the metric
inverse have not been established.
\end{assessment}

\begin{decision}[V8 stress-derivative acquisition split]
\label{dec:v16v7-next}
Separate the four rows in (7.27).  First prove a resolvent-contour
bound for \(\dot P_n\) and \(\dot G_n\) from a uniform constraint
gap and form derivative.  Then express \(\dot\rho_n\) through the
renormalized stress insertion and state the exact
exponential-weighted moment needed for (7.20).  If the constraint
gap closes, move upstream to a stratified kernel treatment rather
than using (7.9) across a rank jump.
\end{decision}

\begin{firewall}[Claim boundary after V7]
\label{fire:v16v7}
V7 is a conditional derivative compiler.  It does not prove the
uniform derivative moments, a global constraint gap, reflection
positivity of the tilted transfer, or the full
Standard Model--Einstein continuum.
\end{firewall}

\section*{Version V7 append-only record}
\addcontentsline{toc}{section}{Version V7 append-only record}
The complete V6 source was retained before this continuation.  It
had \(65541\) bytes, \(1845\) lines, and SHA--256
\[
 \texttt{721927BF63654D5F294F6430F2F7ED4C71344FDDFFB7D0AD86BB13C4D35C784C}.
\]
V7 differentiated the nonlinear centered Newton tilt, retained the
moving zero-mode projection in the pseudoinverse formula, and
proved a uniform-integrability theorem for passage of the complete
Einstein-source response.
\section{Sixteenth-cycle V8: quantitative moving-kernel resolvent bounds}
\label{sec:v16v8}

V7 derives the exact metric derivative of the centered Newton
quadratic, including the moving constraint kernel.  This note makes
that derivative quantitative.  The bounded transform separates the
unbounded elliptic-domain issue from the finite spectral gap and
gives a common contour on which both cutoff convergence and metric
differentiation can be proved.

\subsection{Bounded transform of the constraint operator}

Let \(L_s\geq0\) be self-adjoint and put
\[
 B_s=L_s(I+L_s)^{-1}=I-(I+L_s)^{-1}.                  \tag{8.1}
\]
The operators have the same kernel.  If
\[
 \sigma(L_s)\subset\{0\}\cup[\gamma,\infty),\qquad
 \gamma>0,                                            \tag{8.2}
\]
then
\[
 \sigma(B_s)\subset\{0\}\cup[\beta,1],\qquad
 \beta={\gamma\over1+\gamma}.                         \tag{8.3}
\]
Let \(K_s\) project onto the kernel and \(P_s=I-K_s\).

\begin{lemma}[Recovering the Newton inverse]
\label{lem:v16v8-recovery}
If \(B_s^+\) is the Moore--Penrose inverse of \(B_s\), then
\[
 G_s=L_s^+=B_s^+-P_s.                                 \tag{8.4}
\]
\end{lemma}

\begin{proof}
On a positive spectral value
\[
 b={\lambda\over1+\lambda},
\]
one has
\[
 {1\over\lambda}={1-b\over b}={1\over b}-1.           \tag{8.5}
\]
Both sides of (8.4) vanish on the common kernel.  The spectral
theorem proves the identity.
\end{proof}

\subsection{Riesz bounds for the moving kernel}

\begin{theorem}[Projection derivative from the gap]
\label{thm:v16v8-projection}
Suppose \(s\mapsto B_s\) is differentiable in operator norm,
(8.3) holds uniformly, and
\[
 \|\dot B_s\|_{\rm op}\leq D.                         \tag{8.6}
\]
Then
\[
 \|\dot K_s\|_{\rm op}
 =\|\dot P_s\|_{\rm op}
 \leq {2D\over\beta}.                                 \tag{8.7}
\]
\end{theorem}

\begin{proof}
Let \({\cal C}\) be the positively oriented circle
\(|z|=\beta/2\).  It encloses only the zero eigenvalue, and
\[
 K_s={1\over2\pi i}\int_{\cal C}(z-B_s)^{-1}\,dz.      \tag{8.8}
\]
Every point of the circle is at distance at least \(\beta/2\)
from the spectrum, so
\[
 \|(z-B_s)^{-1}\|\leq{2\over\beta}.                   \tag{8.9}
\]
Differentiating the resolvent gives
\[
 \dot K_s={1\over2\pi i}\int_{\cal C}
 (z-B_s)^{-1}\dot B_s(z-B_s)^{-1}\,dz.                \tag{8.10}
\]
The contour length is \(\pi\beta\).  Equations (8.6),
(8.9), and (8.10) give
\[
 \|\dot K_s\|
 \leq{\pi\beta\over2\pi}{4D\over\beta^2}
 ={2D\over\beta}.
\]
Since \(\dot P_s=-\dot K_s\), (8.7) follows.
\end{proof}

\begin{theorem}[Pseudoinverse derivative bound]
\label{thm:v16v8-pseudoinverse}
Under the hypotheses of
Theorem~\ref{thm:v16v8-projection},
\[
 \|\dot B_s^+\|_{\rm op}
 \leq {5D\over\beta^2},                               \tag{8.11}
\]
and
\[
 \|\dot G_s\|_{\rm op}
 \leq {5D\over\beta^2}+{2D\over\beta}.                \tag{8.12}
\]
\end{theorem}

\begin{proof}
Apply the moving-pseudoinverse identity of V7 to \(B_s\):
\[
 \dot B_s^+
 =-B_s^+\dot B_sB_s^+
   +\dot P_sB_s^++B_s^+\dot P_s.                     \tag{8.13}
\]
The gap gives \(\|B_s^+\|\leq\beta^{-1}\).
Equations (8.6)--(8.7) therefore yield
\[
 \|\dot B_s^+\|
 \leq {D\over\beta^2}
 +2{2D\over\beta}{1\over\beta}
 ={5D\over\beta^2}.                                  \tag{8.14}
\]
Differentiate (8.4) and apply (8.7), proving (8.12).
\end{proof}

The constants are conservative contour bounds.  Their role is to
display the exact inverse powers of the gap; no claim of sharpness
is made.

\subsection{From a form derivative to \(\dot B\)}

Let \({\cal V}\) be the common form domain of \(L_s\), continuously
and densely embedded in the carrier \({\cal H}\).  Suppose the form
derivative defines
\[
 \dot L_s\in{\cal B}({\cal V},{\cal V}^*),\qquad
 \|\dot L_s\|_{{\cal V}\to{\cal V}^*}\leq D_{\cal V}. \tag{8.15}
\]
Let \(R_s=(I+L_s)^{-1}\).

\begin{proposition}[Form-to-bounded-transform derivative]
\label{prop:v16v8-form}
If
\[
 \|R_s\|_{{\cal H}\to{\cal V}}\leq C_R               \tag{8.16}
\]
uniformly, then \(B_s\) is operator-norm differentiable and
\[
 \dot B_s=R_s\dot L_sR_s,\qquad
 \|\dot B_s\|_{\rm op}\leq C_R^2D_{\cal V}.           \tag{8.17}
\]
\end{proposition}

\begin{proof}
The variational resolvent identity gives
\[
 \dot R_s=-R_s\dot L_sR_s                             \tag{8.18}
\]
as a bounded map from \({\cal H}\) to itself, where the second copy
of \(R_s\) is read by duality as a map
\({\cal V}^*\to{\cal H}\).  Since \(B_s=I-R_s\),
(8.17) follows.  The norm estimate uses
(8.15)--(8.16) on the two sides.
\end{proof}

\begin{corollary}[Elliptic metric-derivative card]
\label{cor:v16v8-card}
A uniform form-domain resolvent bound (8.16), a uniform form
derivative (8.15), and the positive constraint gap (8.2) imply all
operator bounds required for \(\dot P_s\) and \(\dot G_s\) in V7.
\end{corollary}

\begin{proof}
Use Proposition~\ref{prop:v16v8-form} with
\(D=C_R^2D_{\cal V}\), then Theorems
\ref{thm:v16v8-projection}--\ref{thm:v16v8-pseudoinverse}.
\end{proof}

\subsection{Convergence of the derivative ledger}

\begin{theorem}[Cutoff convergence of \(P,G,\dot P,\dot G\)]
\label{thm:v16v8-convergence}
Let \(B_{n,s},B_s\) be nonnegative contractions on one carrier and
assume one gap \(\beta>0\):
\[
 \sigma(B_{n,s}),\sigma(B_s)
 \subset\{0\}\cup[\beta,1].                           \tag{8.19}
\]
At a fixed \(s\), suppose
\[
 \|B_{n,s}-B_s\|_{\rm op}\to0,\qquad
 \|\dot B_{n,s}-\dot B_s\|_{\rm op}\to0.              \tag{8.20}
\]
Then
\[
 P_{n,s}\to P_s,\quad G_{n,s}\to G_s,\quad
 \dot P_{n,s}\to\dot P_s,\quad
 \dot G_{n,s}\to\dot G_s                              \tag{8.21}
\]
in operator norm.
\end{theorem}

\begin{proof}
Use the fixed circle in (8.8) for every \(n\).  The resolvent
identity and (8.20) give uniform convergence of the resolvents on
that contour, hence convergence of \(K_n\) and \(P_n\).

Choose a second fixed contour enclosing \([\beta,1]\) while
excluding zero.  Holomorphic functional calculus gives
\[
 B_{n,s}^+
 ={1\over2\pi i}\int_{\cal C_+}
 z^{-1}(z-B_{n,s})^{-1}\,dz.                          \tag{8.22}
\]
The same resolvent argument gives \(B_n^+\to B^+\).
Differentiating (8.8) and (8.22) inserts one copy of
\(\dot B_n\) between two resolvents.  Both parts of (8.20) then give
convergence of the derivatives.  Finally use
\(G_n=B_n^+-P_n\) and its derivative.
\end{proof}

This theorem provides the operator convergence needed in the four
rows of (7.27), subject to the fixed-rank gap.

\subsection{A derivative-energy bound for \(\dot Q\)}

Let
\[
 Q_s=\langle r_s,G_sr_s\rangle,\qquad
 Q_s^{\rm der}
 =\langle\dot r_s,G_s\dot r_s\rangle.                 \tag{8.23}
\]

\begin{proposition}[Relative Newton derivative]
\label{prop:v16v8-Q-relative}
Assume, as quadratic forms on \(\operatorname{Ran}P_s\),
\[
 |\langle u,\dot G_su\rangle|
 \leq c_G\langle u,G_su\rangle.                       \tag{8.24}
\]
Then, for every \(\eta>0\),
\[
 |\dot Q_s|
 \leq(\eta+c_G)Q_s+\eta^{-1}Q_s^{\rm der}.            \tag{8.25}
\]
\end{proposition}

\begin{proof}
The Cauchy--Schwarz inequality for the positive semidefinite form
of \(G_s\) gives
\[
 2|\langle\dot r_s,G_sr_s\rangle|
 \leq\eta\langle r_s,G_sr_s\rangle
 +\eta^{-1}\langle\dot r_s,G_s\dot r_s\rangle.         \tag{8.26}
\]
Combine this with (7.11) and (8.24).
\end{proof}

Thus the V7 derivative passage can be paid by a strengthened
exponential moment for \(Q_s\) together with a polynomial or
exponential moment for the derivative energy \(Q_s^{\rm der}\).
The original Newton margin alone does not supply it.

\begin{counterexample}[The \(H^{-1}\) source norm misses high
derivative modes]
\label{cnt:v16v8-highmode}
On \(\ell^2(\mathbb N)\), let
\[
 Ge_j=j^{-2}e_j,\qquad r_n=ne_n.                      \tag{8.27}
\]
Then
\[
 \langle r_n,Gr_n\rangle=1                            \tag{8.28}
\]
for all \(n\), while \(\|r_n\|=n\to\infty\).  A bound on
\(e^{aQ_n}\) therefore gives no \(L^2\) or stress-derivative control
on the source.
\end{counterexample}

\begin{proof}
Equation (8.28) is immediate from (8.27).
\end{proof}

\subsection{Gap collapse}

\begin{counterexample}[Pseudoinverse derivative blow-up]
\label{cnt:v16v8-gap}
For \(s>0\), take
\[
 B_s=
 \begin{pmatrix}s&0\\0&1\end{pmatrix}.                \tag{8.29}
\]
Then \(\|\dot B_s\|=1\), but
\[
 B_s^+=
 \begin{pmatrix}s^{-1}&0\\0&1\end{pmatrix},\qquad
 \|\dot B_s^+\|=s^{-2}.                               \tag{8.30}
\]
No uniform derivative estimate survives as the positive spectral
gap closes.
\end{counterexample}

\begin{proof}
Differentiate the diagonal entries.
\end{proof}

At a rank jump, the common contour itself disappears.  The
stratified-kernel treatment identified in fifteenth-cycle V96--V98
is then required.

\begin{assessment}[Physical status]
\label{ass:v16v8-status}
V8 proves quantitative moving-projection and pseudoinverse
derivative bounds, an elliptic form-to-bounded-transform compiler,
and operator-norm convergence of the complete derivative ledger.
It also isolates the derivative energy not controlled by the
Newton exponential moment.

No uniform constraint gap or exponential-weighted stress-derivative
estimate has been proved for the full fluctuating geometry.
\end{assessment}

\begin{decision}[V9 reflection-local Newton sandwich]
\label{dec:v16v8-next}
Return to the separate reflection row left open in V6.  At finite
cutoff, characterize when a same-time Newton weight can be split
across a reflection link as a positive sandwich of the free
transfer.  Prove trace-class and convergence bounds for bounded
weights, then state the exact uniform integrability needed to pass
to unbounded weights without confusing scalar positivity with
operator reflection positivity.
\end{decision}

\begin{firewall}[Claim boundary after V8]
\label{fire:v16v8}
V8 is a uniform-gap derivative theorem.  It does not prove that the
gap survives cutoff or metric variation, control a kernel-rank
jump, establish the stress moments, or construct the full
Standard Model--Einstein continuum.
\end{firewall}

\section*{Version V8 append-only record}
\addcontentsline{toc}{section}{Version V8 append-only record}
The complete V7 source was retained before this continuation.  It
had \(77040\) bytes, \(2196\) lines, and SHA--256
\[
 \texttt{EF4903C1DC47A9AB81FF0B8996074339F4FC8F0D4EDA7C994141BEA034B26594}.
\]
V8 converted the moving-kernel derivative ledger into quantitative
gap and form-resolvent bounds, proved its cutoff convergence, and
isolated the independent derivative-energy moment.
\section{Sixteenth-cycle V9: unbounded Newton half-weights and closed transfer sandwiches}
\label{sec:v16v9}

Seventh-cycle V65 already proves positivity of a finite bounded
sandwich \(P W G W P\).  V9 does not repeat that theorem.  It treats
the missing case created by the Newton tilt: the half-weight is
unbounded, so the formal product \(B T B\) must be defined by a
closed trace-ideal construction.

\subsection{Weighted Hilbert--Schmidt criterion}

Let \({\cal H}\) be a Hilbert space, let \(T\geq0\) be trace class,
and let \(B\geq0\) be a possibly unbounded self-adjoint operator.
Assume \({\cal D}(B)\) is dense.  The Newton application is
\[
 B=e^{aQ/2},\qquad a={\Gamma\over2},                   \tag{9.1}
\]
so that \(B^2=e^{aQ}\) is the one-slice weight.

\begin{definition}[Admissible unbounded half-weight]
\label{def:v16v9-admissible}
The pair \((T,B)\) is admissible if the densely defined product
\[
 B T^{1/2}                                             \tag{9.2}
\]
extends to a Hilbert--Schmidt operator \(R\in{\cal L}^2({\cal H})\).
\end{definition}

\begin{theorem}[Closed positive Newton sandwich]
\label{thm:v16v9-sandwich}
If \((T,B)\) is admissible, then
\[
 T_B:=RR^*                                             \tag{9.3}
\]
is positive and trace class, with
\[
 \operatorname{Tr}T_B=\|R\|_2^2.                     \tag{9.4}
\]
On the natural form core \({\cal D}(B)\),
\[
 \langle\psi,T_B\psi\rangle
 =\|T^{1/2}B\psi\|^2
 =\langle B\psi,T B\psi\rangle.                       \tag{9.5}
\]
Thus \(T_B\) is the closed realization of the formal sandwich
\(B T B\).
\end{theorem}

\begin{proof}
A Hilbert--Schmidt operator is bounded, and the product of two
Hilbert--Schmidt operators is trace class.  Hence \(RR^*\) is
positive trace class and
\[
 \operatorname{Tr}(RR^*)=\|R\|_2^2.                  \tag{9.6}
\]
The adjoint of the closure of \(BT^{1/2}\) extends
\(T^{1/2}B\).  On \({\cal D}(B)\),
\[
 \langle\psi,RR^*\psi\rangle
 =\|R^*\psi\|^2
 =\|T^{1/2}B\psi\|^2,
\]
which is (9.5).
\end{proof}

\begin{corollary}[One-link reflection form]
\label{cor:v16v9-link}
If \(T\) is the positive transfer across a reflection link and
\(B\) is a same-time half-weight, then the closed reflected boundary
form after insertion of the full weight \(B^2\) is
\[
 \langle f,T_Bf\rangle=\|R^*f\|^2\geq0.               \tag{9.7}
\]
\end{corollary}

\begin{proof}
This is (9.5), followed by closure from the form core.
\end{proof}

The phrase same-time is a hypothesis about factorization of the
constraint action.  Positivity of a weight coupling both reflected
halves does not by itself put it in the form (9.7).

\subsection{Diagonal exponential moment}

Suppose \({\cal H}=L^2(X,\mu)\), \(B\) is multiplication by
\(b(x)\geq0\), and \(T\) has a positive trace-class diagonal in the
sense that
\[
 T=\sum_j\tau_j|\phi_j\rangle\langle\phi_j|,\qquad
 t(x,x)=\sum_j\tau_j|\phi_j(x)|^2                     \tag{9.8}
\]
with monotone convergence on the diagonal.

\begin{proposition}[Diagonal criterion]
\label{prop:v16v9-diagonal}
The pair \((T,B)\) is admissible exactly when
\[
 \int_X b(x)^2t(x,x)\,d\mu(x)<\infty.                 \tag{9.9}
\]
In that case
\[
 \operatorname{Tr}T_B
 =\int_X b(x)^2t(x,x)\,d\mu(x).                       \tag{9.10}
\]
\end{proposition}

\begin{proof}
Using (9.8), Tonelli's theorem gives
\[
\begin{split}
 \|BT^{1/2}\|_2^2
 &=\sum_j\tau_j\|B\phi_j\|^2\\
 &=\int_Xb(x)^2
     \sum_j\tau_j|\phi_j(x)|^2\,d\mu(x).
                                                               \tag{9.11}
\end{split}
\]
This is finite exactly under (9.9), and
Theorem~\ref{thm:v16v9-sandwich} gives (9.10).
\end{proof}

Normalize the transfer diagonal to a probability measure
\[
 d\pi_T(x)={t(x,x)\over\operatorname{Tr}T}\,d\mu(x).  \tag{9.12}
\]
For the Newton weight (9.1), criterion (9.9) becomes
\[
 {\mathbb E}_{\pi_T}e^{aQ}<\infty.                    \tag{9.13}
\]
Thus the nonlinear exponential moment in V6 is precisely the
trace-class payment for a reflection-local sandwich when its base
law is the transfer diagonal.

\subsection{Monotone truncation of the unbounded weight}

Let
\[
 B_M=\min(B,M)                                        \tag{9.14}
\]
by spectral calculus and put \(T_{B_M}=B_MT B_M\).

\begin{theorem}[Trace-norm closure by truncation]
\label{thm:v16v9-truncation}
If \((T,B)\) is admissible, then
\[
 \|B_MT^{1/2}-BT^{1/2}\|_2\longrightarrow0,           \tag{9.15}
\]
\[
 \|T_{B_M}-T_B\|_1\longrightarrow0.                  \tag{9.16}
\]
\end{theorem}

\begin{proof}
In the spectral representation used in
Proposition~\ref{prop:v16v9-diagonal}, or abstractly by the spectral
measure of \(B\) against \(T\),
\[
 \|(B-B_M)T^{1/2}\|_2^2
 =\operatorname{Tr}\bigl(
 T^{1/2}(B-B_M)^2T^{1/2}\bigr).                       \tag{9.17}
\]
The integrand decreases pointwise to zero and is dominated by the
finite trace associated with \(B^2\).  Dominated convergence proves
(9.15).

For Hilbert--Schmidt \(R,S\),
\[
 \|RR^*-SS^*\|_1
 \leq(\|R\|_2+\|S\|_2)\|R-S\|_2.                     \tag{9.18}
\]
Apply this with \(R=BT^{1/2}\) and \(S=B_MT^{1/2}\) to get (9.16).
\end{proof}

The bounded theorem of seventh-cycle V65 applies to every \(B_M\);
Theorem~\ref{thm:v16v9-truncation} shows that its positive
trace-class limit is the unbounded Newton sandwich.

\subsection{Cutoff convergence on the common carrier}

\begin{theorem}[Weighted Hilbert--Schmidt handoff]
\label{thm:v16v9-handoff}
Let \(T_n,T\geq0\) be trace class on one carrier and let
\(B_n,B\geq0\) be admissible half-weights.  Put
\[
 R_n=B_nT_n^{1/2},\qquad R=BT^{1/2}.                  \tag{9.19}
\]
If
\[
 \|R_n-R\|_2\longrightarrow0,                         \tag{9.20}
\]
then
\[
 \|T_{n,B_n}-T_B\|_1\longrightarrow0.                \tag{9.21}
\]
\end{theorem}

\begin{proof}
Condition (9.20) bounds \(\|R_n\|_2\).  Apply (9.18) to
\(R_n,R\).
\end{proof}

\begin{corollary}[Bounded-weight specialization]
\label{cor:v16v9-bounded}
If \(B_n,B\) are uniformly bounded,
\[
 \|B_n-B\|_{\rm op}\to0,\qquad
 \|T_n-T\|_1\to0,                                    \tag{9.22}
\]
then (9.20)--(9.21) hold.
\end{corollary}

\begin{proof}
The Powers--St{\o}rmer inequality gives
\[
 \|T_n^{1/2}-T^{1/2}\|_2^2\leq\|T_n-T\|_1.           \tag{9.23}
\]
Write
\[
 B_nT_n^{1/2}-BT^{1/2}
 =(B_n-B)T_n^{1/2}
  +B(T_n^{1/2}-T^{1/2})                               \tag{9.24}
\]
and use the ideal bounds.
\end{proof}

For unbounded \(B_n\), scalar convergence of the normalization
\(\|R_n\|_2^2\) does not imply the directional convergence (9.20).
The common-carrier weighted Hilbert--Schmidt row is an independent
alignment hypothesis.

\begin{counterexample}[Constant weighted norm without alignment]
\label{cnt:v16v9-drift}
On \(\ell^2(\mathbb N)\), let
\[
 R_n=|e_n\rangle\langle e_1|.                         \tag{9.25}
\]
Then \(\|R_n\|_2=1\) for all \(n\), but
\[
 R_nR_n^*=|e_n\rangle\langle e_n|                    \tag{9.26}
\]
has no trace-norm convergent subsequence.
\end{counterexample}

\begin{proof}
Distinct rank-one projections have trace-norm distance two.
\end{proof}

\subsection{A practical uniform-tail condition}

\begin{proposition}[Truncation before cutoff removal]
\label{prop:v16v9-two-limits}
Suppose, for each \(M<\infty\),
\[
 B_{n,M}T_n^{1/2}\longrightarrow B_MT^{1/2}
 \quad\hbox{in Hilbert--Schmidt norm},                 \tag{9.27}
\]
and
\[
 \lim_{M\to\infty}\sup_n
 \|(B_n-B_{n,M})T_n^{1/2}\|_2=0.                     \tag{9.28}
\]
Assume the same tail convergence for \(B,T\).  Then (9.20) holds.
\end{proposition}

\begin{proof}
Insert the two truncated terms:
\[
\begin{split}
 \|B_nT_n^{1/2}-BT^{1/2}\|_2
 \leq{}&
 \|(B_n-B_{n,M})T_n^{1/2}\|_2\\
 &+\|B_{n,M}T_n^{1/2}-B_MT^{1/2}\|_2\\
 &+\|(B_M-B)T^{1/2}\|_2.                              \tag{9.29}
\end{split}
\]
First choose \(M\) using the two tails, then \(n\) using (9.27).
\end{proof}

On a transfer diagonal, (9.28) is the uniform-integrability tail
\[
 \sup_n\int_{\{b_n>M\}}b_n^2t_n(x,x)\,d\mu_n(x)
 \longrightarrow0.                                   \tag{9.30}
\]
The strict exponential margin of V6 is a sufficient scalar payment
for this tail, while (9.27) supplies the missing carrier alignment.

\subsection{Scope of reflection positivity}

Theorems~\ref{thm:v16v9-sandwich} and
\ref{thm:v16v9-handoff} prove positivity of the one-link boundary
operator if the Newton action splits into a same-slice half-weight
on each side of the link.  They do not prove that the solution of
the gravitational constraint has this time-local structure.  A
constraint inverse involving fields on both sides of the reflection
plane may fail to factor as \(B T B\), even though its scalar
Boltzmann weight is positive.

\begin{assessment}[Physical status]
\label{ass:v16v9-status}
V9 extends the archived bounded sandwich theorem to unbounded
Newton weights.  It gives the exact weighted Hilbert--Schmidt
criterion, constructs the closed positive trace-class transfer,
and proves truncation and cutoff convergence in trace norm.

The reflection-local factorization of the physical sourced
constraint and the weighted Hilbert--Schmidt convergence have not
been proved.
\end{assessment}

\begin{decision}[V10 companion audit and factorization test]
\label{dec:v16v9-next}
Perform the compulsory YM/NS audit.  Then distinguish two cases for
the Newton constraint: an instantaneous same-time inverse, which
can enter the V9 sandwich, and a time-nonlocal inverse, which
requires positivity of its full cross-reflection kernel.  Derive
the block-kernel Schur criterion for the latter before making any
reflection claim.
\end{decision}

\begin{firewall}[Claim boundary after V9]
\label{fire:v16v9}
V9 is an operator closure theorem conditional on
reflection-local factorization.  It does not establish that
factorization for the Einstein constraint, prove the uniform
weighted alignment, or construct the full continuum.
\end{firewall}

\section*{Version V9 append-only record}
\addcontentsline{toc}{section}{Version V9 append-only record}
The complete V8 source was retained before this continuation.  It
had \(87582\) bytes, \(2547\) lines, and SHA--256
\[
 \texttt{3D29EF111ECC58D6CA11A544FB0A32B29690BBA7658F876CE60D5CE264DBDFA6}.
\]
V9 extended bounded transfer sandwiches to unbounded Newton
half-weights through a weighted Hilbert--Schmidt criterion and
proved trace-norm truncation and cutoff handoff.
\section{Sixteenth-cycle V10: companion audit and the exact reflection branch ledger}
\label{sec:v16v10}

This is the compulsory five-version audit.  It also checks the
factorization task proposed at the end of V9 against the frozen SM
archives.  That check changes the direction: the Gaussian
cross-reflection criterion and the support-preserving constraint
criterion were already proved upstream, so they are cited rather than
proved again.

\subsection{Companion checkpoints}

\begin{audit}[Files inspected at V10]
\label{aud:v16v10-files}
The read-only checkpoints are
\[
\begin{array}{c|r|r}
\hbox{file}&\hbox{lines}&\hbox{bytes}\\ \hline
\texttt{Renewal\_Yang\_Mills\_Mass\_Gap\_Research\_Notes\_V46.tex}
 &10844&396051\\
\texttt{Renewal\_NS\_Stability\_Research\_Notes\_V26.tex}
 &5319&278283.
\end{array}                                             \tag{10.1}
\]
Their SHA--256 values are respectively
\[
\begin{gathered}
\texttt{B1E5580A117A5BED92CDB4AC94DDEB6E5992F93E18D5DB46F0A38FAC6D3E5469},\\
\texttt{82C887F8C2C69E4F28FAD7C3DC8DE5B9099CB5A4BF431EBA1FFF3E91935978AD}.
\end{gathered}                                         \tag{10.2}
\]
\end{audit}

\begin{theorem}[V10 companion-audit decision]
\label{thm:v16v10-audit-decision}
No theorem from either companion checkpoint supplies the
reflection-local Newton factorization or the weighted
Hilbert--Schmidt convergence required in V9.
\end{theorem}

\begin{proof}
YM V45--V46 independently certify the opposite second-coordinate
chart layer and hence both response signs on paired regions
\([0,\pi/4]^2\times[0,t]\), with an explicit positive rational
floor.  The third-coordinate extension, full transverse cell,
physical heat-transfer limit, reflection-positive interacting
measure, and mass gap remain open there.  A finite Loewner floor for
the reduced Wilson response pencil neither constructs the
Einstein-constraint half-weight nor controls
\(B_nT_n^{1/2}\) in Hilbert--Schmidt norm.

NS V26 computes pointwise rank-one norms for the first two
derivatives of the Oseen kernel.  Those estimates concern a
rank-one source cone in the ancient Navier--Stokes representation.
They contain no reflected transfer, Newton exponential weight, or
common-carrier trace-ideal estimate.  Thus neither checkpoint proves
either missing premise.
\end{proof}

\begin{assessment}[What may be imported later]
\label{ass:v16v10-companions}
The YM two-sign atlas is relevant only after it is connected to a
positive physical transfer and a cofinal continuum construction.
Its present local chart theorem cannot be used as that connection.
The NS rank-one improvement has no current SM owner.
\end{assessment}

\subsection{Archive audit of the proposed factorization test}

Three earlier results already separate the reflection branches:
\begin{enumerate}
\item seventh-cycle V65 proves preservation of transfer positivity
under bounded same-slice sandwiches;
\item ninth-cycle V74 proves that a reflection-equivariant
deterministic constraint solution preserves OS positivity when it
also preserves the positive-time observable algebra, and records
that a spacetime-elliptic inverse need not preserve that algebra;
\item fourteenth-cycle V58 proves that a finite centered Gaussian is
OS positive exactly when its physical future cross-covariance is
positive semidefinite.
\end{enumerate}
Current V9 extends the first item to admissible unbounded
half-weights.

\begin{proposition}[Exact reflection branch ledger]
\label{prop:v16v10-branch-ledger}
Let \(\Theta\) be the physical reflection and let the retained
finite-cutoff law be OS positive.
\begin{enumerate}
\item If the constraint is reconstructed by a deterministic map
\(\Phi\) satisfying
\[
 \Phi\Theta=\Theta\Phi                              \tag{10.3}
\]
and \(F\circ\Phi\) is a positive-time observable whenever \(F\) is,
then the reconstructed law is OS positive.
\item If the reduced constraint contribution is a same-time scalar
weight \(B^2\), and \(BT^{1/2}\) extends to a Hilbert--Schmidt
operator for the positive link transfer \(T\), then its closed
boundary transfer is
\[
 T_B=(BT^{1/2})(BT^{1/2})^*\geq0.                   \tag{10.4}
\]
\item If the constraint is time-nonlocal and the resulting finite
law is centered Gaussian with covariance \(C\), future inclusion
\(\iota_+\), and physical reflected adjoint \(\Theta\), then OS
positivity on the future polynomial algebra is equivalent to
\[
 C_{\rm OS}:=\iota_+^*\Theta C\iota_+\succeq0.       \tag{10.5}
\]
\end{enumerate}
None of these three hypotheses follows merely from positivity and
normalizability of the scalar Boltzmann density.
\end{proposition}

\begin{proof}
For item 1 and every positive-time \(F\),
\[
\begin{split}
 \int\overline{F(\Theta x)}F(x)\,d(\Phi_\#\mu)(x)
 &=
 \int\overline{(F\circ\Phi)(\Theta y)}
             (F\circ\Phi)(y)\,d\mu(y)\geq0,          \tag{10.6}
\end{split}
\]
where (10.3) gives the equality and the support hypothesis makes
\(F\circ\Phi\) admissible in the original OS form.  This is the
ninth-cycle V74 theorem.

Item 2 is Theorem~\ref{thm:v16v9-sandwich} and
Corollary~\ref{cor:v16v9-link}.

For item 3, testing the reflected form on linear functions first
shows necessity of (10.5).  If (10.5) holds, Gaussian exponential
vectors have, after positive row and column normalizations, kernel
\[
 (u,v)\longmapsto\exp\langle u,C_{\rm OS}v\rangle.   \tag{10.7}
\]
Every finite Gram matrix of (10.7) is positive because its power
series is a nonnegative sum of Schur powers of the positive Gram
matrix of \(C_{\rm OS}\).  Differentiation at the origin gives all
future polynomials.  This is precisely the fourteenth-cycle V58
criterion.

Finally, that same archived result contains a positive Gaussian
density for which the physical twisted cross-covariance is negative.
It therefore disproves any implication from scalar positivity or
normalizability to the reflection conclusions above.
\end{proof}

\begin{corollary}[What must be computed for the physical Newton branch]
\label{cor:v16v10-physical-test}
A spatial-slice York or Hamiltonian constraint may use item 1 or
item 2 only after its branch is proved reflection equivariant and
support preserving.  A spacetime-elliptic solution must instead
produce its full physical cross-reflection operator and verify an
appropriate positive-kernel condition; in the Gaussian
linearization that condition is exactly (10.5).
\end{corollary}

\begin{proof}
Spatial nonlocality does not change time support, whereas a
spacetime inverse generally mixes the two half-spaces.  Apply the
corresponding items of
Proposition~\ref{prop:v16v10-branch-ledger}.
\end{proof}

\begin{assessment}[Audit-induced direction]
\label{ass:v16v10-status}
The factorization dichotomy is now completely referenced without
duplicating its archived proofs.  The live analytic bottleneck
created by V9 is sharper: even on a common carrier, trace-norm
convergence of \(T_n\) and scalar convergence of the Newton
normalizations do not force
\(B_nT_n^{1/2}\to BT^{1/2}\) in Hilbert--Schmidt norm.
\end{assessment}

\begin{decision}[V11 strong-truncation compiler]
\label{dec:v16v10-next}
Go upstream to the weighted alignment row.  Replace the
operator-norm convergence demanded in
Corollary~\ref{cor:v16v9-bounded} by uniformly bounded strong
convergence of truncated multipliers.  Then prove that one strict
higher Newton trace moment supplies the truncation tails with an
explicit rate.  This will give a checkable common-carrier theorem
for the unbounded transfer sandwiches.
\end{decision}

\begin{firewall}[Claim boundary after V10]
\label{fire:v16v10}
V10 performs the required companion and archive audits and fixes
the exact reflection branch ledger.  It does not select or construct
the physical Einstein constraint branch, prove its cross-reflection
condition, or establish the weighted cutoff limit.
\end{firewall}

\section*{Version V10 append-only record}
\addcontentsline{toc}{section}{Version V10 append-only record}
The complete V9 source was retained before this continuation.  It
had \(97963\) bytes, \(2885\) lines, and SHA--256
\[
 \texttt{9EF4BD6CE21D0A598BB9CB580898D8615879A14EF63FC4FD57BBEBED54F0BAF4}.
\]
V10 inspected YM V46 and NS V26, identified the already-proved
reflection factorization branches, and redirected the next
iteration to the unproved weighted alignment estimate.
\section{Sixteenth-cycle V11: strong truncations close the weighted transfer handoff}
\label{sec:v16v11}

V9 reduced convergence of the unbounded Newton sandwiches to
\[
 B_nT_n^{1/2}\longrightarrow BT^{1/2}
 \quad\hbox{in Hilbert--Schmidt norm}.                \tag{11.1}
\]
Its bounded specialization assumed operator-norm convergence of the
weights.  That is unnecessarily strong for multiplication operators:
bounded pointwise convergence gives only strong operator convergence.
This section proves that strong convergence is sufficient after it is
paired with trace-norm convergence of the base transfers, and that one
strict higher Newton trace moment gives the required tails
quantitatively.

\subsection{Strong convergence against a Hilbert--Schmidt factor}

\begin{lemma}[Strong--Hilbert--Schmidt compactness]
\label{lem:v16v11-strong-HS}
Let \(A_n,A\in{\cal B}({\cal H})\), assume
\[
 A_n\longrightarrow A\quad\hbox{strongly},\qquad
 \sup_n\|A_n\|_{\rm op}<\infty,                       \tag{11.2}
\]
and let \(K\in{\cal L}^2({\cal H})\).  Then
\[
 \|(A_n-A)K\|_2\longrightarrow0.                     \tag{11.3}
\]
\end{lemma}

\begin{proof}
Let \((e_j)\) be an orthonormal basis and put
\(M=\sup_n\|A_n-A\|_{\rm op}<\infty\).  For every \(j\),
strong convergence gives
\[
 \|(A_n-A)Ke_j\|^2\longrightarrow0,                  \tag{11.4}
\]
while
\[
 \|(A_n-A)Ke_j\|^2\leq M^2\|Ke_j\|^2.                \tag{11.5}
\]
The majorant is summable because \(K\) is Hilbert--Schmidt.
Dominated convergence for the series gives
\[
 \|(A_n-A)K\|_2^2
 =\sum_j\|(A_n-A)Ke_j\|^2\longrightarrow0.
\]
\end{proof}

\begin{proposition}[Bounded strong-weight handoff]
\label{prop:v16v11-bounded-strong}
Let \(T_n,T\geq0\) be trace class on one Hilbert space, with
\[
 \|T_n-T\|_1\longrightarrow0.                        \tag{11.6}
\]
Let \(A_n,A\) be uniformly bounded and \(A_n\to A\) strongly.  Then
\[
 \|A_nT_n^{1/2}-AT^{1/2}\|_2\longrightarrow0.        \tag{11.7}
\]
\end{proposition}

\begin{proof}
The Powers--St{\o}rmer inequality gives
\[
 \|T_n^{1/2}-T^{1/2}\|_2^2\leq\|T_n-T\|_1.           \tag{11.8}
\]
Therefore
\[
\begin{split}
 \|A_nT_n^{1/2}-AT^{1/2}\|_2
 \leq{}&
 \|A_n\|_{\rm op}\|T_n^{1/2}-T^{1/2}\|_2\\
 &+\|(A_n-A)T^{1/2}\|_2.                             \tag{11.9}
\end{split}
\]
The first term tends to zero by (11.8), and the second by
Lemma~\ref{lem:v16v11-strong-HS}.
\end{proof}

\subsection{A strict trace exponent pays the unbounded tail}

Let \({\cal H}=L^2(X,\mu)\), let \(a>0\), and let
\(Q_n,Q:X\to[0,\infty)\) be measurable.  Define multiplication
operators
\[
 B_n=e^{aQ_n/2},\qquad B=e^{aQ/2},                    \tag{11.10}
\]
and their bounded truncations
\[
 B_{n,M}=\min(B_n,M),\qquad B_M=\min(B,M).            \tag{11.11}
\]
For an unbounded positive multiplier \(W\), the notation
\[
 \operatorname{Tr}(T^{1/2}W T^{1/2})
 :=\|W^{1/2}T^{1/2}\|_2^2                           \tag{11.12}
\]
means the extended positive trace; finiteness includes the domain
condition on the product.

\begin{lemma}[Quantitative Newton truncation tail]
\label{lem:v16v11-tail}
Fix \(\delta>0\).  If
\[
 \operatorname{Tr}\!\left(
 T_n^{1/2}e^{(a+\delta)Q_n}T_n^{1/2}\right)\leq C,    \tag{11.13}
\]
then, for every \(M\geq1\),
\[
 \|(B_n-B_{n,M})T_n^{1/2}\|_2
 \leq C^{1/2}M^{-\delta/a}.                          \tag{11.14}
\]
The same assertion holds for \(B,T,Q\).
\end{lemma}

\begin{proof}
For \(q\geq0\) and \(b=e^{aq/2}\),
\[
 (b-\min(b,M))^2
 \leq b^2{\bf1}_{\{b>M\}}.                           \tag{11.15}
\]
On \(\{b>M\}\), one has
\[
 q>{2\log M\over a},
\qquad
 e^{aq}\leq M^{-2\delta/a}e^{(a+\delta)q}.            \tag{11.16}
\]
Functional calculus and positivity therefore give
\[
\begin{split}
 \|(B_n-B_{n,M})T_n^{1/2}\|_2^2
 &=
 \operatorname{Tr}\!\left(
 T_n^{1/2}(B_n-B_{n,M})^2T_n^{1/2}\right)\\
 &\leq
 M^{-2\delta/a}
 \operatorname{Tr}\!\left(
 T_n^{1/2}e^{(a+\delta)Q_n}T_n^{1/2}\right),
                                                               \tag{11.17}
\end{split}
\]
which is (11.14).
\end{proof}

\begin{theorem}[Strong-truncation compiler for Newton transfers]
\label{thm:v16v11-compiler}
Let \(T_n,T\geq0\) be trace class on \(L^2(X,\mu)\).  Suppose
\begin{enumerate}
\item \(\|T_n-T\|_1\to0\);
\item \(Q_n\to Q\) almost everywhere;
\item for some \(a,\delta,C>0\),
\[
\begin{split}
 \sup_n\operatorname{Tr}\!\left(
 T_n^{1/2}e^{(a+\delta)Q_n}T_n^{1/2}\right)&\leq C,\\
 \operatorname{Tr}\!\left(
 T^{1/2}e^{(a+\delta)Q}T^{1/2}\right)&\leq C.
                                                               \tag{11.18}
\end{split}
\]
\end{enumerate}
Then \(B_nT_n^{1/2}\) and \(BT^{1/2}\) are
Hilbert--Schmidt and
\[
 \boxed{\ 
 \|B_nT_n^{1/2}-BT^{1/2}\|_2\longrightarrow0.\ }     \tag{11.19}
\]
Consequently the closed positive sandwiches
\[
 S_n:=(B_nT_n^{1/2})(B_nT_n^{1/2})^*,\qquad
 S:=(BT^{1/2})(BT^{1/2})^*                           \tag{11.20}
\]
satisfy
\[
 \|S_n-S\|_1\longrightarrow0.                        \tag{11.21}
\]
\end{theorem}

\begin{proof}
Because \(Q_n\to Q\) almost everywhere, the bounded multipliers
\(B_{n,M}\) converge strongly to \(B_M\): for every \(f\in L^2\),
\[
 |(B_{n,M}-B_M)f|^2\leq4M^2|f|^2,                   \tag{11.22}
\]
and dominated convergence applies.  Their operator norms are at
most \(M\).  Proposition~\ref{prop:v16v11-bounded-strong} yields,
for each fixed \(M\),
\[
 \|B_{n,M}T_n^{1/2}-B_MT^{1/2}\|_2\longrightarrow0. \tag{11.23}
\]

The lower exponent in (11.10) is controlled by (11.18), so all
untruncated products in (11.19) are Hilbert--Schmidt.  Moreover,
Lemma~\ref{lem:v16v11-tail} gives the uniform estimate
\[
\begin{split}
 \|B_nT_n^{1/2}-BT^{1/2}\|_2
 \leq{}&2C^{1/2}M^{-\delta/a}\\
 &+\|B_{n,M}T_n^{1/2}-B_MT^{1/2}\|_2.                \tag{11.24}
\end{split}
\]
First let \(n\to\infty\) at fixed \(M\), then let \(M\to\infty\).
This proves (11.19).

Finally, with \(R_n=B_nT_n^{1/2}\) and \(R=BT^{1/2}\),
\[
 \|R_nR_n^*-RR^*\|_1
 \leq(\|R_n\|_2+\|R\|_2)\|R_n-R\|_2.                 \tag{11.25}
\]
The first factor is bounded by (11.18), and the second tends to zero
by (11.19), proving (11.21).
\end{proof}

\begin{corollary}[Convergence of normalized reflected states]
\label{cor:v16v11-normalized}
Under the assumptions of
Theorem~\ref{thm:v16v11-compiler}, suppose
\[
 Z:=\operatorname{Tr}S>0.                            \tag{11.26}
\]
Then \(Z_n:=\operatorname{Tr}S_n\to Z\), and
\[
 \left\|{S_n\over Z_n}-{S\over Z}\right\|_1
 \longrightarrow0.                                  \tag{11.27}
\]
In particular, the expectations of every bounded boundary
observable converge.
\end{corollary}

\begin{proof}
Trace continuity gives
\[
 |Z_n-Z|\leq\|S_n-S\|_1\longrightarrow0.             \tag{11.28}
\]
Hence \(Z_n>0\) for all sufficiently large \(n\), and
\[
 \left\|{S_n\over Z_n}-{S\over Z}\right\|_1
 \leq{\|S_n-S\|_1\over Z_n}
     +{|Z_n-Z|\over Z_n}\longrightarrow0.            \tag{11.29}
\]
Pairing with a bounded operator and using trace duality proves the
last assertion.
\end{proof}

\subsection{What the theorem removes and what it assumes}

The almost-everywhere condition in
Theorem~\ref{thm:v16v11-compiler} is used only to prove strong
convergence of every bounded truncation.  It may be replaced
directly by
\[
 B_{n,M}\longrightarrow B_M\quad\hbox{strongly for each }M.
                                                               \tag{11.30}
\]
The strict exponent \(\delta\) is used only for the explicit uniform
tail (11.14).  V6 already shows by a rare-event example that
control at the physical exponent alone does not imply the needed
uniform integrability, while V9 shows that convergence of scalar
normalizations alone does not give operator alignment.  Those
counterexamples are not repeated here.

\begin{assessment}[Physical status]
\label{ass:v16v11-status}
V11 replaces bounded operator-norm convergence of Newton
half-weights by bounded strong convergence and proves a complete
unbounded common-carrier handoff from one higher weighted trace
moment.  It therefore closes the analytic implication from the
three hypotheses of Theorem~\ref{thm:v16v11-compiler} to the
positive normalized reflected transfer limit.

The programme has not yet constructed a common carrier on which
the physical \(Q_n\) converge almost everywhere, nor proved the
strict trace moment (11.18) uniformly in the coupled
Standard Model--Einstein regulator.
\end{assessment}

\begin{decision}[V12 Gaussian Newton reserve]
\label{dec:v16v11-next}
Connect (11.18) to the exact Gaussian Newton determinant of V4.
Derive a uniform spectral reserve, expressed through
\(C_n^{1/2}G_nC_n^{1/2}\), that pays the higher exponent and remains
stable under trace-norm covariance convergence.  Separate that
finite Gaussian certificate from the still nonlinear interacting
matter law.
\end{decision}

\begin{firewall}[Claim boundary after V11]
\label{fire:v16v11}
V11 is a sufficient weighted-transfer convergence theorem.  It
does not prove its common-carrier, almost-everywhere, or uniform
higher-moment hypotheses for the physical theory, select the
Einstein constraint branch, or construct the full continuum.
\end{firewall}

\section*{Version V11 append-only record}
\addcontentsline{toc}{section}{Version V11 append-only record}
The complete V10 source was retained before this continuation.  It
had \(106306\) bytes, \(3091\) lines, and SHA--256
\[
 \texttt{DB637CC5EF60CFBC9284BAA8C9028ED6E575F31268B29ACD2FE3B4905816BB4F}.
\]
V11 proved that strong convergence of bounded Newton truncations
and one strict higher weighted trace moment force
Hilbert--Schmidt convergence of the unbounded half-transfers and
trace-norm convergence of their normalized positive sandwiches.
\section{Sixteenth-cycle V12: Gaussian spectral reserve for the Newton trace moment}
\label{sec:v16v12}

V11 requires a uniform trace moment at an exponent strictly above
the physical Newton exponent \(a\).  For a Gaussian transfer
diagonal, V4 gives the exact determinant.  This section converts
that formula into a regulator-uniform reserve and feeds it directly
into the V11 transfer theorem.

\subsection{Gaussian transfer diagonals}

Let \(T_n\geq0\) be trace class on a common
\(L^2(X,\mu)\), with diagonal \(t_n(x,x)\), and write
\[
 \tau_n:=\operatorname{Tr}T_n,\qquad
 d\pi_n(x):={t_n(x,x)\over\tau_n}\,d\mu(x).           \tag{12.1}
\]
Assume \(\tau_n>0\) and that, under \(\pi_n\), the source \(X_n\)
is a centered Gaussian on a real separable source space with
trace-class covariance \(C_n\).  Let
\[
 Q_n=\langle X_n,G_nX_n\rangle,\qquad
 A_n=C_n^{1/2}G_nC_n^{1/2}\geq0,                     \tag{12.2}
\]
where \(G_n\geq0\) is the mean-zero Newton inverse.

\begin{proposition}[Exact weighted transfer trace]
\label{prop:v16v12-exact}
For every \(s>0\),
\[
 \operatorname{Tr}\!\left(
 T_n^{1/2}e^{sQ_n}T_n^{1/2}\right)<\infty            \tag{12.3}
\]
if and only if
\[
 2s\|A_n\|_{\rm op}<1.                               \tag{12.4}
\]
When (12.4) holds,
\[
 \operatorname{Tr}\!\left(
 T_n^{1/2}e^{sQ_n}T_n^{1/2}\right)
 =
 \tau_n\det(I-2sA_n)^{-1/2}.                         \tag{12.5}
\]
\end{proposition}

\begin{proof}
The diagonal trace identity and (12.1) give
\[
 \operatorname{Tr}\!\left(
 T_n^{1/2}e^{sQ_n}T_n^{1/2}\right)
 =\tau_n{\mathbb E}_{\pi_n}e^{sQ_n}.                 \tag{12.6}
\]
Apply Theorem~\ref{thm:v16v4-gaussian} with
\(\Gamma=2s\).  Its sharp criterion is (12.4), and its determinant
formula is (12.5).
\end{proof}

\subsection{Uniform reserve and an explicit bound}

\begin{theorem}[Uniform Gaussian Newton reserve]
\label{thm:v16v12-reserve}
Let \(a,\delta>0\).  Suppose
\[
 \sup_n\tau_n\leq\tau<\infty,\qquad
 \sup_n 2(a+\delta)\|A_n\|_{\rm op}\leq q<1,\qquad
 \sup_n\operatorname{Tr}A_n\leq L<\infty.            \tag{12.7}
\]
Then
\[
 \sup_n\operatorname{Tr}\!\left(
 T_n^{1/2}e^{(a+\delta)Q_n}T_n^{1/2}\right)
 \leq
 \tau\exp\!\left\{{(a+\delta)L\over1-q}\right\}.      \tag{12.8}
\]
Thus (12.7) verifies the uniform higher-moment hypothesis in
Theorem~\ref{thm:v16v11-compiler}.
\end{theorem}

\begin{proof}
Let \(s=a+\delta\).  The eigenvalues \(\alpha_{j,n}\) of \(A_n\)
satisfy \(0\leq2s\alpha_{j,n}\leq q\).  Hence
\[
\begin{split}
 -{1\over2}\log\det(I-2sA_n)
 &={1\over2}\sum_j-\log(1-2s\alpha_{j,n})\\
 &\leq{s\over1-q}\sum_j\alpha_{j,n}
 \leq{sL\over1-q}.                                   \tag{12.9}
\end{split}
\]
Combine (12.5), \(\tau_n\leq\tau\), and (12.9).
\end{proof}

\begin{lemma}[A strict limiting margin creates an extra exponent]
\label{lem:v16v12-extra}
Suppose
\[
 \|A_n-A\|_1\longrightarrow0,\qquad
 2a\|A\|_{\rm op}<1.                                 \tag{12.10}
\]
Then there exist \(\delta>0\), \(q<1\), and \(N\) such that
\[
 2(a+\delta)\|A_n\|_{\rm op}\leq q
 \quad(n\geq N).                                     \tag{12.11}
\]
Moreover,
\[
 \sup_{n\geq N}\operatorname{Tr}A_n<\infty.          \tag{12.12}
\]
\end{lemma}

\begin{proof}
Trace-norm convergence implies operator-norm convergence and
convergence of the traces of the positive operators.  Choose
\(\eta>0\) such that
\[
 2a(\|A\|_{\rm op}+2\eta)<1.                         \tag{12.13}
\]
For all sufficiently large \(n\),
\(\|A_n\|_{\rm op}\leq\|A\|_{\rm op}+\eta\).  By
continuity in \(\delta\), choose \(\delta>0\) so small that
\[
 q:=2(a+\delta)(\|A\|_{\rm op}+\eta)<1.              \tag{12.14}
\]
This proves (12.11).  Since
\[
 |\operatorname{Tr}A_n-\operatorname{Tr}A|
 \leq\|A_n-A\|_1,                                    \tag{12.15}
\]
the traces are bounded on the same tail.
\end{proof}

\begin{corollary}[Gaussian closure of the V11 moment row]
\label{cor:v16v12-v11}
Assume the setting of
Theorem~\ref{thm:v16v11-compiler}, and additionally:
\begin{enumerate}
\item the normalized transfer diagonals are the Gaussian laws in
(12.1)--(12.2);
\item \(\sup_n\tau_n<\infty\);
\item \(A_n\to A\) in trace norm and
\[
 2a\|A\|_{\rm op}<1.                                 \tag{12.16}
\]
\end{enumerate}
After discarding a finite regulator prefix, the higher trace bound
(11.18) holds for some \(\delta>0\).  Consequently
\[
 B_nT_n^{1/2}\to BT^{1/2}\quad\hbox{in }{\cal L}^2,
\qquad
 S_n\to S\quad\hbox{in }{\cal L}^1,                  \tag{12.17}
\]
provided the common-carrier trace convergence and almost-everywhere
convergence hypotheses of V11 hold.
\end{corollary}

\begin{proof}
Lemma~\ref{lem:v16v12-extra} supplies \(\delta,q,N\) and a uniform
trace bound for \(A_n\).  Apply
Theorem~\ref{thm:v16v12-reserve} to \(n\geq N\).  The limiting
moment is finite by Proposition~\ref{prop:v16v12-exact} because
\(2(a+\delta)\|A\|_{\rm op}\leq q<1\).  Hence every moment
hypothesis of Theorem~\ref{thm:v16v11-compiler} holds on the
cofinal tail, and that theorem proves (12.17).
\end{proof}

\begin{corollary}[Primitive covariance--Newton sufficient data]
\label{cor:v16v12-primitive}
The trace-norm premise \(A_n\to A\) in the preceding results follows
if
\[
 C_n^{1/2}\to C^{1/2}\quad\hbox{in }{\cal L}^2,
\qquad
 G_n\to G\quad\hbox{in operator norm},               \tag{12.18}
\]
with \(\|C_n^{1/2}\|_2\) and \(\|G_n\|_{\rm op}\)
uniformly bounded.
\end{corollary}

\begin{proof}
This is Corollary~\ref{cor:v16v4-factors}: insert two intermediate
products between
\(C_n^{1/2}G_nC_n^{1/2}\) and
\(C^{1/2}GC^{1/2}\), then use
\(\|RST\|_1\leq\|R\|_2\|S\|_{\rm op}\|T\|_2\).
\end{proof}

\subsection{Critical and non-Gaussian boundaries}

Condition (12.16) is a strict physical margin.  At
\(2a\|A\|_{\rm op}=1\), even the physical Gaussian normalization is
critical, and no positive \(\delta\) can satisfy (12.11).
Approach to that boundary produces the susceptibility blow-up
proved in Proposition~\ref{prop:v16v4-critical}.

The Gaussian assumption is also load bearing.  An interacting
source law is not determined by its covariance, so bounds on
\(C_n^{1/2}G_nC_n^{1/2}\) alone do not control
\({\mathbb E}e^{(a+\delta)Q_n}\).  V4 already gives a
growing-dimensional counterexample to covariance-scale
subgaussian reasoning; it is not repeated.

\begin{assessment}[Physical status]
\label{ass:v16v12-status}
V12 gives a sharp, checkable Gaussian certificate for the higher
Newton trace moment required by V11.  A strict limiting spectral
margin automatically creates the extra exponent on a cofinal tail,
and trace-norm convergence of the covariance-weighted Newton
operators makes the determinant bound uniform.

The actual coupled Standard Model source is interacting rather than
Gaussian.  The certificate therefore applies to the Gaussian
reference transfer and to Gaussian comparison arguments, but it
does not by itself prove the physical nonlinear moment.
\end{assessment}

\begin{decision}[V13 comparison transfer from Gaussian reference]
\label{dec:v16v12-next}
Attack the nonlinear gap.  Let the interacting diagonal law have
density \(h_n\) relative to the Gaussian reference.  Derive the
sharp H\"older trade between an \(L^p\) bound on \(h_n\) and a
stronger Gaussian Newton exponent.  Determine exactly how much
spectral reserve is consumed and feed the resulting moment into
V11.
\end{decision}

\begin{firewall}[Claim boundary after V12]
\label{fire:v16v12}
V12 closes the higher-moment row only for Gaussian transfer
diagonals satisfying the strict susceptibility margin.  It does
not control the interacting density, prove the physical
common-carrier convergence, select the constraint branch, or
construct the full continuum.
\end{firewall}

\section*{Version V12 append-only record}
\addcontentsline{toc}{section}{Version V12 append-only record}
The complete V11 source was retained before this continuation.  It
had \(115991\) bytes, \(3396\) lines, and SHA--256
\[
 \texttt{064711C173FC5D38B6B88EFDB9FD89374CACD8C8002BBC93241B616930447ABC}.
\]
V12 converted the exact Gaussian Newton determinant into a uniform
spectral reserve that supplies the strict weighted trace moment in
the V11 transfer handoff.
\section{Sixteenth-cycle V13: \(L^p\) interaction densities and the consumed Gaussian reserve}
\label{sec:v16v13}

V12 pays the strict Newton trace moment when the transfer diagonal
is Gaussian.  The physical matter diagonal is interacting.  Write
its normalized law as \(d\nu_n=h_n\,d\pi_n\), where \(\pi_n\) is
the Gaussian reference of V12.  The correct comparison datum is an
\(L^p(\pi_n)\) bound on \(h_n\); ordinary relative entropy controls
only a first moment and does not by itself pay an exponential
Newton moment.

\subsection{Exact Gaussian dual norm}

Let \(\pi\) be a centered Gaussian law with covariance \(C\), let
\[
 Q=\langle X,GX\rangle,\qquad
 A=C^{1/2}GC^{1/2}\geq0,                              \tag{13.1}
\]
and fix \(s>0\).  For \(1<p<\infty\), let
\[
 q={p\over p-1}.                                      \tag{13.2}
\]

\begin{theorem}[Optimal \(L^p\) Gaussian comparison exponent]
\label{thm:v16v13-dual}
The positive linear functional
\[
 {\cal L}_s(h):={\mathbb E}_\pi[h\,e^{sQ}]            \tag{13.3}
\]
is bounded on \(L^p(\pi)\) if and only if
\[
 2qs\|A\|_{\rm op}<1.                                \tag{13.4}
\]
When (13.4) holds, its exact norm is
\[
 \|{\cal L}_s\|_{(L^p)^*}
 =
 \|e^{sQ}\|_{L^q(\pi)}
 =
 \det(I-2qsA)^{-1/(2q)}.                             \tag{13.5}
\]
\end{theorem}

\begin{proof}
Since \(1<p<\infty\), the dual of \(L^p(\pi)\) is
\(L^q(\pi)\), and the norm of integration against the positive
function \(e^{sQ}\) is exactly its \(L^q\) norm.  Now
\[
 \|e^{sQ}\|_{L^q(\pi)}^q
 ={\mathbb E}_\pi e^{qsQ}.                           \tag{13.6}
\]
Theorem~\ref{thm:v16v4-gaussian}, with
\(\Gamma=2qs\), says that (13.6) is finite exactly under
(13.4), and then equals
\(\det(I-2qsA)^{-1/2}\).  Taking the \(q\)-th root proves
(13.5).  If (13.4) fails, the multiplier is not in \(L^q\), so
duality says that (13.3) is not a bounded functional on
\(L^p\).
\end{proof}

Thus the factor \(q\) in (13.4) is not an artifact of H\"older's
inequality.  It is the exact price for uniform comparison over an
\(L^p\) class.

\begin{corollary}[Interacting Newton moment]
\label{cor:v16v13-density}
Let \(h\geq0\), \({\mathbb E}_\pi h=1\), and
\(\|h\|_{L^p(\pi)}\leq H\).  If (13.4) holds, then the law
\(d\nu=h\,d\pi\) satisfies
\[
 {\mathbb E}_\nu e^{sQ}
 \leq H\det(I-2qsA)^{-1/(2q)}.                       \tag{13.7}
\]
\end{corollary}

\begin{proof}
Apply H\"older's inequality, or equivalently (13.5), to
\({\mathbb E}_\pi[h e^{sQ}]\).
\end{proof}

In terms of the order-\(p\) R\'enyi divergence
\[
 D_p(\nu\Vert\pi)
 ={1\over p-1}\log{\mathbb E}_\pi h^p,               \tag{13.8}
\]
one has \(\log\|h\|_p=D_p/q\), so (13.7) becomes
\[
 \log{\mathbb E}_\nu e^{sQ}
 \leq {D_p(\nu\Vert\pi)\over q}
 -{1\over2q}\log\det(I-2qsA).                        \tag{13.9}
\]

\subsection{Uniform interacting transfer bound}

Return to the transfer diagonals.  Let
\[
 d\nu_n=h_n\,d\pi_n,\qquad
 {\mathbb E}_{\pi_n}h_n=1,                           \tag{13.10}
\]
where \(\pi_n,C_n,G_n,A_n,\tau_n\) are as in V12 and \(\nu_n\)
is the normalized diagonal law of the interacting transfer \(T_n\).

\begin{theorem}[Uniform \(L^p\)-comparison certificate]
\label{thm:v16v13-uniform}
Fix \(a,\delta>0\), \(1<p<\infty\), and
\(q=p/(p-1)\).  Suppose
\[
\begin{gathered}
 \sup_n\|h_n\|_{L^p(\pi_n)}\leq H,\qquad
 \sup_n\tau_n\leq\tau,\\
 \sup_n2q(a+\delta)\|A_n\|_{\rm op}\leq r<1,\qquad
 \sup_n\operatorname{Tr}A_n\leq L.                  \tag{13.11}
\end{gathered}
\]
Then
\[
\boxed{
 \sup_n\operatorname{Tr}\!\left(
 T_n^{1/2}e^{(a+\delta)Q_n}T_n^{1/2}\right)
 \leq
 \tau H\exp\!\left\{
 {(a+\delta)L\over1-r}\right\}.}                     \tag{13.12}
\]
\end{theorem}

\begin{proof}
With \(s=a+\delta\), the diagonal trace identity gives
\[
 \operatorname{Tr}\!\left(
 T_n^{1/2}e^{sQ_n}T_n^{1/2}\right)
 =\tau_n{\mathbb E}_{\nu_n}e^{sQ_n}.                 \tag{13.13}
\]
Corollary~\ref{cor:v16v13-density} bounds the last expectation by
\[
 H\det(I-2qsA_n)^{-1/(2q)}.                          \tag{13.14}
\]
If \(\alpha_{j,n}\) are the eigenvalues of \(A_n\), then
\[
\begin{split}
 -{1\over2q}\log\det(I-2qsA_n)
 &={1\over2q}\sum_j-\log(1-2qs\alpha_{j,n})\\
 &\leq{s\over1-r}\operatorname{Tr}A_n
 \leq{sL\over1-r}.                                   \tag{13.15}
\end{split}
\]
Equations (13.13)--(13.15) and the bound on \(\tau_n\) prove
(13.12).
\end{proof}

\begin{lemma}[Strict interacting limiting reserve]
\label{lem:v16v13-reserve}
Assume \(A_n\to A\) in trace norm and
\[
 2qa\|A\|_{\rm op}<1.                                \tag{13.16}
\]
Then there exist \(\delta>0\), \(r<1\), and \(N\) such that
\[
 2q(a+\delta)\|A_n\|_{\rm op}\leq r
 \quad(n\geq N),                                     \tag{13.17}
\]
and \(\sup_{n\geq N}\operatorname{Tr}A_n<\infty\).
\end{lemma}

\begin{proof}
Repeat the proof of Lemma~\ref{lem:v16v12-extra} with \(a\)
replaced by \(qa\).  Trace-norm convergence gives both
operator-norm convergence and bounded traces.
\end{proof}

\begin{corollary}[Interacting closure of the V11 moment row]
\label{cor:v16v13-v11}
Suppose, on a cofinal regulator tail,
\[
 \sup_n\|h_n\|_{L^p(\pi_n)}<\infty,\qquad
 A_n\to A\ \hbox{in trace norm},\qquad
 2qa\|A\|_{\rm op}<1,                                \tag{13.18}
\]
and \(\sup_n\tau_n<\infty\).  Then the strict higher trace moment
required by Theorem~\ref{thm:v16v11-compiler} holds.  If the
remaining common-carrier assumptions of that theorem hold, the
unbounded Newton half-transfers converge in Hilbert--Schmidt norm
and their normalized positive sandwiches converge in trace norm.
\end{corollary}

\begin{proof}
Lemma~\ref{lem:v16v13-reserve} supplies \(\delta\) and \(r\);
Theorem~\ref{thm:v16v13-uniform} supplies (11.18).  Apply
Theorem~\ref{thm:v16v11-compiler} and
Corollary~\ref{cor:v16v11-normalized}.
\end{proof}

\subsection{The price of the interaction}

For \(p=\infty\), formally \(q=1\), and the Gaussian reserve of V12
is recovered with the additional factor
\(\|h_n\|_\infty\).  As \(p\downarrow1\), \(q\uparrow\infty\), so
condition (13.16) becomes impossible for every nonzero \(A\).
Thus an \(L^1\) density bound, which is merely normalization, pays no
uniform positive Newton exponent.

Finite relative entropy is also weaker than the required finite
R\'enyi control.  Indeed, (13.9) uses \(D_p\) with \(p>1\);
passing formally to \(p=1\) sends the Gaussian exponent \(qs\) to
infinity.  This identifies the missing nonlinear estimate as a
uniform R\'enyi or equivalent density-integrability bound, not a
covariance estimate.

\begin{assessment}[Physical status]
\label{ass:v16v13-status}
V13 proves the exact \(L^p\)-Gaussian dual criterion and a uniform
interacting transfer bound.  It quantifies the interaction cost:
the physical susceptibility reserve is multiplied by
\(q=p/(p-1)\).  Together with V11, this turns a uniform
\(L^p\) comparison density and a strict spectral margin into a
trace-norm limit of normalized positive Newton sandwiches.

No uniform \(L^p\) or R\'enyi bound for the coupled chiral gauge,
Higgs, fermion, and gravitational interaction density has yet been
proved.
\end{assessment}

\begin{decision}[V14 local interaction factorization]
\label{dec:v16v13-next}
Identify a finite-volume route to the \(L^p\) density bound.  For an
interaction decomposed into local cells, prove the exact product
law for R\'enyi divergence under an independent Gaussian reference
and show why a volume-uniform global \(L^p\) norm generally fails.
Then formulate the local-observable replacement needed before the
thermodynamic limit.
\end{decision}

\begin{firewall}[Claim boundary after V13]
\label{fire:v16v13}
V13 is a comparison theorem conditional on a uniform \(L^p\)
interaction density.  It does not establish that density bound,
control its volume growth, prove the physical constraint
factorization, or construct the full continuum.
\end{firewall}

\section*{Version V13 append-only record}
\addcontentsline{toc}{section}{Version V13 append-only record}
The complete V12 source was retained before this continuation.  It
had \(124161\) bytes, \(3641\) lines, and SHA--256
\[
 \texttt{12FDEDB396335370D7E07E54836F055364F56BB96B31956F89D0BDB4B683EA08}.
\]
V13 proved the exact \(L^p\) comparison criterion for an interacting
density relative to the Gaussian transfer diagonal and quantified
the resulting factor-\(q\) consumption of the Newton spectral
reserve.
\section{Sixteenth-cycle V14: volume growth of interaction densities and the local-state exit}
\label{sec:v16v14}

Ninth-cycle V25 already replaces an infinite-volume all-good metric
event by conditional block exits and local closure.  V14 does not
repeat that argument.  It addresses the distinct density-norm
bottleneck created by V13: even a completely independent product
interaction usually has an \(L^p\) norm that grows exponentially
with volume.  The correct thermodynamic object is therefore a
normalized local state, not a global trace-class density with a
volume-uniform norm.

\subsection{Exact product law for interaction densities}

Let \(\Lambda\) be a finite set of cells.  For each \(x\in\Lambda\),
let \((X_x,\pi_x)\) be a probability space and let
\(h_x\geq0\) satisfy
\[
 \int h_x\,d\pi_x=1.                                 \tag{14.1}
\]
Put
\[
 \pi_\Lambda=\bigotimes_{x\in\Lambda}\pi_x,\qquad
 h_\Lambda=\prod_{x\in\Lambda}h_x,\qquad
 d\nu_\Lambda=h_\Lambda\,d\pi_\Lambda.               \tag{14.2}
\]

\begin{theorem}[Product R\'enyi law]
\label{thm:v16v14-product}
For \(1<p<\infty\),
\[
 \|h_\Lambda\|_{L^p(\pi_\Lambda)}
 =\prod_{x\in\Lambda}\|h_x\|_{L^p(\pi_x)},           \tag{14.3}
\]
and
\[
 D_p(\nu_\Lambda\Vert\pi_\Lambda)
 =\sum_{x\in\Lambda}D_p(\nu_x\Vert\pi_x).            \tag{14.4}
\]
If \(h_x=h\) for all \(x\) and \(h\ne1\) on a set of positive
\(\pi\)-measure, then
\[
 \|h_\Lambda\|_p
 =\exp\{c_p|\Lambda|\},\qquad
 c_p:=\log\|h\|_p>0.                                \tag{14.5}
\]
\end{theorem}

\begin{proof}
Tonelli's theorem and (14.2) give
\[
\begin{split}
 \|h_\Lambda\|_p^p
 &=\int\prod_{x\in\Lambda}h_x^p\,d\pi_\Lambda\\
 &=\prod_{x\in\Lambda}\int h_x^p\,d\pi_x
 =\prod_{x\in\Lambda}\|h_x\|_p^p,                   \tag{14.6}
\end{split}
\]
which proves (14.3).  Taking logarithms and using the definition
(13.8) proves (14.4).

Since \(\pi\) is a probability measure and \(\int h\,d\pi=1\),
H\"older gives \(\|h\|_p\geq1\).  Strict convexity of
\(u\mapsto u^p\) shows that equality holds only when \(h=1\)
almost everywhere.  Hence \(c_p>0\) for a nontrivial density, and
(14.5) follows from (14.3).
\end{proof}

\begin{corollary}[No nontrivial volume-uniform global \(L^p\) card]
\label{cor:v16v14-global-Lp}
For a homogeneous nontrivial product interaction, no bound
\[
 \sup_\Lambda\|h_\Lambda\|_{L^p(\pi_\Lambda)}<\infty \tag{14.7}
\]
can hold along volumes with \(|\Lambda|\to\infty\).
\end{corollary}

\begin{proof}
Equation (14.5) diverges exponentially.
\end{proof}

Thus the global hypothesis used in V13 is appropriate for cutoff
removal at a fixed compact spatial volume.  It cannot be imposed
unchanged on the thermodynamic limit.

\subsection{Marginal density is the local replacement}

For \(K\subset\Lambda\), let \({\cal F}_K\) be the sigma-algebra
generated by the cells in \(K\), and define
\[
 h_{\Lambda,K}
 :={\mathbb E}_{\pi_\Lambda}[h_\Lambda\mid{\cal F}_K].
                                                               \tag{14.8}
\]
After the natural identification, this is a function on
\(\prod_{x\in K}X_x\).

\begin{lemma}[Local marginal density and \(L^p\) contraction]
\label{lem:v16v14-marginal}
The \(K\)-marginal of \(\nu_\Lambda\) has density
\(h_{\Lambda,K}\) relative to
\(\pi_K=\bigotimes_{x\in K}\pi_x\), and
\[
 \|h_{\Lambda,K}\|_{L^p(\pi_K)}
 \leq\|h_\Lambda\|_{L^p(\pi_\Lambda)}.               \tag{14.9}
\]
For the product density (14.2),
\[
 h_{\Lambda,K}=\prod_{x\in K}h_x,\qquad
 \|h_{\Lambda,K}\|_p=\prod_{x\in K}\|h_x\|_p.        \tag{14.10}
\]
\end{lemma}

\begin{proof}
For every bounded \({\cal F}_K\)-measurable \(F\),
\[
 \int F\,d\nu_\Lambda
 =\int Fh_\Lambda\,d\pi_\Lambda
 =\int Fh_{\Lambda,K}\,d\pi_K,                       \tag{14.11}
\]
which proves the density statement.  Conditional Jensen gives
\[
 |h_{\Lambda,K}|^p
 \leq{\mathbb E}_{\pi_\Lambda}
       [h_\Lambda^p\mid{\cal F}_K].                  \tag{14.12}
\]
Integrating proves (14.9).

In the product case, integration over every cell outside \(K\)
uses (14.1) and leaves \(\prod_{x\in K}h_x\).  Equation (14.3)
on \(K\) proves the norm identity.
\end{proof}

\begin{theorem}[Uniform local Gaussian comparison]
\label{thm:v16v14-local-comparison}
Fix a finite region \(K\) and suppose its Gaussian reference
\(\pi_K\) has covariance \(C_K\).  Let
\[
 Q_K=\langle X_K,G_KX_K\rangle,\qquad
 A_K=C_K^{1/2}G_KC_K^{1/2}.                          \tag{14.13}
\]
Let \(q=p/(p-1)\).  If, uniformly over \(\Lambda\supset K\),
\[
 \|h_{\Lambda,K}\|_{L^p(\pi_K)}\leq H_K              \tag{14.14}
\]
and
\[
 2qs\|A_K\|_{\rm op}<1,                              \tag{14.15}
\]
then
\[
 \sup_{\Lambda\supset K}
 {\mathbb E}_{\nu_\Lambda}e^{sQ_K}
 \leq
 H_K\det(I-2qsA_K)^{-1/(2q)}.                        \tag{14.16}
\]
For a homogeneous product interaction, one may take
\[
 H_K=\|h\|_p^{|K|},                                  \tag{14.17}
\]
which is independent of the exterior volume.
\end{theorem}

\begin{proof}
Since \(Q_K\) is \({\cal F}_K\)-measurable,
Lemma~\ref{lem:v16v14-marginal} gives
\[
 {\mathbb E}_{\nu_\Lambda}e^{sQ_K}
 ={\mathbb E}_{\pi_K}
   [h_{\Lambda,K}e^{sQ_K}].                          \tag{14.18}
\]
Apply Theorem~\ref{thm:v16v13-dual} on the \(K\)-carrier and use
(14.14).  This proves (14.16).  Equation (14.17) is (14.10) in the
homogeneous case.
\end{proof}

\subsection{Global Newton moments still grow with volume}

The preceding local theorem cannot be applied to the full Newton
energy merely by enlarging \(K\) without cost.  This is already
visible in an independent additive model.

\begin{proposition}[Exact growth of an additive Newton moment]
\label{prop:v16v14-global-moment}
Suppose \(\nu_\Lambda=\bigotimes_{x\in\Lambda}\nu_x\) and
\[
 Q_\Lambda=\sum_{x\in\Lambda}Q_x,\qquad Q_x\geq0,    \tag{14.19}
\]
where \(Q_x\) depends only on cell \(x\).  Then
\[
 {\mathbb E}_{\nu_\Lambda}e^{sQ_\Lambda}
 =\prod_{x\in\Lambda}{\mathbb E}_{\nu_x}e^{sQ_x}.    \tag{14.20}
\]
In the homogeneous case, if
\(\nu(Q>0)>0\), then
\[
 {\mathbb E}_{\nu_\Lambda}e^{sQ_\Lambda}
 =\exp\{\gamma_s|\Lambda|\},\qquad
 \gamma_s:=\log{\mathbb E}_\nu e^{sQ}>0.             \tag{14.21}
\]
\end{proposition}

\begin{proof}
Independence and Tonelli give (14.20).  Since \(Q\geq0\),
\(e^{sQ}\geq1\), with strict inequality on a set of positive
measure when \(\nu(Q>0)>0\).  Hence its expectation is strictly
larger than one, proving (14.21).
\end{proof}

\begin{theorem}[Normalized local tilt survives global divergence]
\label{thm:v16v14-local-tilt}
Under the homogeneous assumptions of
Proposition~\ref{prop:v16v14-global-moment}, suppose
\[
 z_s:={\mathbb E}_\nu e^{sQ}<\infty.                 \tag{14.22}
\]
Define the globally tilted probability law
\[
 d\widehat\nu_\Lambda
 :={e^{sQ_\Lambda}\over z_s^{|\Lambda|}}\,d\nu_\Lambda.
                                                               \tag{14.23}
\]
Then
\[
 \widehat\nu_\Lambda
 =\bigotimes_{x\in\Lambda}\widehat\nu,\qquad
 d\widehat\nu={e^{sQ}\over z_s}\,d\nu.               \tag{14.24}
\]
Consequently every fixed \(K\)-marginal of
\(\widehat\nu_\Lambda\) is exactly
\(\widehat\nu^{\otimes K}\), independent of
\(\Lambda\supset K\), even though the unnormalized factor
\(z_s^{|\Lambda|}\) diverges exponentially when \(z_s>1\).
\end{theorem}

\begin{proof}
Both the numerator and denominator of (14.23) factor cell by cell:
\[
 {e^{sQ_\Lambda}\over z_s^{|\Lambda|}}
 =\prod_{x\in\Lambda}{e^{sQ_x}\over z_s}.            \tag{14.25}
\]
This proves the product identity and hence the marginal statement.
\end{proof}

\subsection{Two limits must have different targets}

V11--V13 therefore concern the ultraviolet or regulator handoff at
fixed compact volume, where a trace-class transfer and its
normalization are legitimate global objects.  The thermodynamic
limit must instead be formulated on the local observable algebra:
one proves convergence and compatibility of every fixed-region
normalized state.  Global partition functions and global
\(L^p\)-density norms are allowed to grow at their pressure rate.

For the physical Newton inverse, \(Q_\Lambda\) is not additive:
long-range cross terms connect \(K\) with its exterior.  Hence
Theorem~\ref{thm:v16v14-local-tilt} is a model theorem, not the
physical construction.  Its role is to prove that demanding a
volume-uniform global V13 bound would be the wrong thermodynamic
target even before those cross terms are introduced.

\begin{assessment}[Physical status]
\label{ass:v16v14-status}
V14 proves the exact exponential volume growth of product
interaction \(L^p\) norms and additive Newton moments, while also
proving uniform fixed-region marginal comparison and exact
stability of normalized local tilts.  It separates the fixed-volume
trace-ideal limit from the thermodynamic local-state limit.

The missing physical theorem is a uniform local conditional bound
for the long-range Newton cross terms and the interacting
Standard Model density.  Ninth-cycle V25 supplies the compatible
conditional-exit architecture but not that loader.
\end{assessment}

\begin{decision}[V15 audit and exterior-field decomposition]
\label{dec:v16v14-next}
Perform the compulsory companion audit.  Then decompose the Newton
energy relative to a fixed region into interior, exterior, and
cross terms.  Derive a conditional exponential estimate for the
cross term from the positive \(G\)-form and identify the exterior
random variable whose uniform moment would turn V14 into a
physical local-state bound.
\end{decision}

\begin{firewall}[Claim boundary after V14]
\label{fire:v16v14}
V14 proves a product-model obstruction and local-state replacement.
It does not prove the long-range conditional Newton estimate,
construct the interacting thermodynamic state, establish
reflection positivity of that state, or complete the continuum.
\end{firewall}

\section*{Version V14 append-only record}
\addcontentsline{toc}{section}{Version V14 append-only record}
The complete V13 source was retained before this continuation.  It
had \(132567\) bytes, \(3891\) lines, and SHA--256
\[
 \texttt{ED2799DF0A0442230D71D514E2A0CFEA156A0131C89D12513B76A35D98B6896F}.
\]
V14 proved the exponential volume obstruction to global
interaction-density bounds and replaced it by fixed-region
normalized marginal control, while keeping the nonlocal Newton
cross term explicit as the next bottleneck.
\section{Sixteenth-cycle V15: companion audit and the Newton exterior-influence form}
\label{sec:v16v15}

This compulsory checkpoint first records the current YM and NS
notes.  It then supplies the finite-cutoff interior--exterior
decomposition requested by V14.  The useful exterior variable is
not the entire exterior Newton energy; it is the Schur influence
seen through the interior block.

\subsection{Compulsory companion audit}

\begin{audit}[Files inspected at V15]
\label{aud:v16v15-files}
The newest active companion files are unchanged from V10:
\[
\begin{array}{c|r|r}
\hbox{file}&\hbox{lines}&\hbox{bytes}\\ \hline
\texttt{Renewal\_Yang\_Mills\_Mass\_Gap\_Research\_Notes\_V46.tex}
 &10844&396051\\
\texttt{Renewal\_NS\_Stability\_Research\_Notes\_V26.tex}
 &5319&278283.
\end{array}                                             \tag{15.1}
\]
Their SHA--256 values are
\[
\begin{gathered}
\texttt{B1E5580A117A5BED92CDB4AC94DDEB6E5992F93E18D5DB46F0A38FAC6D3E5469},\\
\texttt{82C887F8C2C69E4F28FAD7C3DC8DE5B9099CB5A4BF431EBA1FFF3E91935978AD}.
\end{gathered}                                         \tag{15.2}
\]
\end{audit}

\begin{theorem}[V15 companion-audit decision]
\label{thm:v16v15-audit}
No result from the two companion checkpoints proves or changes the
Newton exterior-influence estimate below.
\end{theorem}

\begin{proof}
YM V46 proves two response signs on paired
\([0,\pi/4]^2\times[0,t]\) transverse plates.  Its third-coordinate
atlas, physical transfer, continuum measure, and mass gap remain
open.  A local finite-dimensional Wilson-pencil floor gives no
bound on the off-diagonal block of the gravitational inverse
constraint.

NS V26 gives pointwise rank-one Oseen-kernel derivative norms.  It
contains neither a positive Newton block operator nor a conditional
Gibbs exponential estimate.  Its constants therefore do not enter
the following hypotheses.
\end{proof}

\subsection{Positive block geometry}

At finite cutoff, decompose the real source space as
\[
 {\cal H}={\cal H}_K\oplus{\cal H}_E,                 \tag{15.3}
\]
where \(K\) is a fixed interior region and \(E\) its exterior.  Write
the nonnegative Newton operator as
\[
 G=
 \begin{pmatrix}
 A&B\\ B^*&C
 \end{pmatrix}\succeq0.                              \tag{15.4}
\]
For \(u\in{\cal H}_K\), \(v\in{\cal H}_E\), put
\[
 Q_K(u)=\langle u,Au\rangle,\qquad
 Q_E(v)=\langle v,Cv\rangle.                         \tag{15.5}
\]

\begin{lemma}[Schur influence of the exterior]
\label{lem:v16v15-schur}
For the finite-dimensional positive block (15.4),
\[
 \operatorname{Ran}B\subset\operatorname{Ran}A.      \tag{15.6}
\]
The exterior influence energy
\[
 J_K(v):=\langle Bv,A^+Bv\rangle                    \tag{15.7}
\]
is nonnegative and satisfies
\[
 0\leq J_K(v)\leq Q_E(v).                            \tag{15.8}
\]
Moreover,
\[
\begin{split}
 \langle(u,v),G(u,v)\rangle
 ={}&
 \langle u+A^+Bv,\,
         A(u+A^+Bv)\rangle\\
 &+Q_E(v)-J_K(v).                                    \tag{15.9}
\end{split}
\]
\end{lemma}

\begin{proof}
If \(w\in\ker A\), positivity of (15.4) gives, for every real
\(t\) and every \(v\),
\[
 0\leq
 \langle(w,tv),G(w,tv)\rangle
 =2t\langle w,Bv\rangle+t^2\langle v,Cv\rangle.      \tag{15.10}
\]
The linear coefficient must vanish.  Thus
\(\ker A\subset\ker B^*\), which in finite dimension is equivalent
to (15.6).

Since \(AA^+Bv=Bv\), expansion of the first term on the
right of (15.9) gives
\[
 Q_K(u)+2\langle u,Bv\rangle+J_K(v).                 \tag{15.11}
\]
This proves (15.9).  The first term in (15.9) is
nonnegative.  Minimizing the left side over \(u\), or taking
\(u=-A^+Bv\), shows
\[
 Q_E(v)-J_K(v)\geq0.                                 \tag{15.12}
\]
Positivity of \(A^+\) gives the other half of (15.8).
\end{proof}

\begin{proposition}[Sharp cross-term carrier]
\label{prop:v16v15-cross}
For every \(\eta>0\),
\[
 2|\langle u,Bv\rangle|
 \leq\eta Q_K(u)+\eta^{-1}J_K(v).                   \tag{15.13}
\]
The exterior energy \(J_K\), rather than all of \(Q_E\), is the
smallest quadratic form produced by completing the square against
the interior \(A\)-energy.
\end{proposition}

\begin{proof}
By (15.6),
\[
 \langle u,Bv\rangle
 =
 \langle A^{1/2}u,A^{+1/2}Bv\rangle.                 \tag{15.14}
\]
Cauchy--Schwarz followed by
\(2xy\leq\eta x^2+\eta^{-1}y^2\) proves (15.13).
Identity (15.9) shows that the squared norm of the translating
vector is exactly (15.7), establishing the completion claim.
\end{proof}

\subsection{Conditional Newton partition bound}

Let \(\rho(du,dv)\) be a finite-cutoff reference probability law
with regular conditional interior laws \(\rho_K(du\mid v)\).  Define
\[
 M_K(\lambda)
 :=\mathop{\rm ess\,sup}_{v}
 \int e^{\lambda Q_K(u)}\,\rho_K(du\mid v).           \tag{15.15}
\]
The exterior-only term \(Q_E(v)\) cancels from the conditional law
obtained after applying the full Newton tilt.  The remaining
conditional partition factor is
\[
 Z_{K,s}(v)
 :=
 \int\exp\{s[Q_K(u)+2\langle u,Bv\rangle]\}
       \,\rho_K(du\mid v).                            \tag{15.16}
\]

\begin{theorem}[Conditional exterior-influence estimate]
\label{thm:v16v15-conditional}
For \(s>0\) and \(\eta>0\),
\[
 Z_{K,s}(v)
 \leq
 M_K(s(1+\eta))
 \exp\{s\eta^{-1}J_K(v)\}                            \tag{15.17}
\]
for almost every exterior configuration \(v\).
If \(0<\eta\leq1\), then also
\[
 Z_{K,s}(v)\geq
 \exp\{-s\eta^{-1}J_K(v)\}.                          \tag{15.18}
\]
\end{theorem}

\begin{proof}
The upper half of (15.13) gives
\[
 Q_K(u)+2\langle u,Bv\rangle
 \leq(1+\eta)Q_K(u)+\eta^{-1}J_K(v).                 \tag{15.19}
\]
Exponentiate, integrate conditionally, and use (15.15) to obtain
(15.17).

The lower half of (15.13) gives
\[
 Q_K(u)+2\langle u,Bv\rangle
 \geq(1-\eta)Q_K(u)-\eta^{-1}J_K(v).                 \tag{15.20}
\]
For \(0<\eta\leq1\), the first term on the right is nonnegative.
Since every conditional law has mass one, integration proves
(15.18).
\end{proof}

\begin{corollary}[Conditional tilted interior moment]
\label{cor:v16v15-tilted}
Let \(\widehat\rho_{K,s}(\,\cdot\mid v)\) be the normalized
conditional law with density proportional to the integrand in
(15.16).  For \(t\geq0\) and \(0<\eta\leq1\),
\[
 \int e^{tQ_K(u)}\,d\widehat\rho_{K,s}(u\mid v)
 \leq
 M_K(t+s(1+\eta))
 \exp\{2s\eta^{-1}J_K(v)\}.                          \tag{15.21}
\]
\end{corollary}

\begin{proof}
The numerator is bounded as in (15.19), now with coefficient
\(t+s(1+\eta)\) on \(Q_K\).  Divide by the lower bound (15.18).
\end{proof}

\subsection{What must be controlled outside the region}

Estimate (15.21) isolates two independent inputs:
\begin{enumerate}
\item a conditional interior Newton moment \(M_K\) with a strict
exponent reserve; and
\item an exterior-influence moment for \(J_K(v)\).
\end{enumerate}
The inequality \(J_K\leq Q_E\) is always available, but it is too
coarse for a thermodynamic limit because \(Q_E\) is extensive or
long range.  The advantage of \(J_K\) is that it measures only the
component of the exterior source coupled into the fixed interior
energy space.

\begin{assessment}[Physical status]
\label{ass:v16v15-status}
V15 completes the companion audit and proves the exact
finite-cutoff Schur decomposition of the Newton energy.  It
identifies \(J_K=B^*A^+B\) as the minimal exterior influence form
and gives two-sided conditional partition estimates and a
conditional tilted-moment bound.

No volume-uniform exponential moment for \(J_K\) has been proved for
the physical inverse constraint.  The estimate also remains at
finite cutoff, where the pseudoinverse and block range are
unambiguous.
\end{assessment}

\begin{decision}[V16 exterior integration and cofinal form limit]
\label{dec:v16v15-next}
Continue after the audit as required.  Integrate (15.21) against the
actual exterior marginal and prove a uniform local Newton moment
from an exponential \(J_K\)-bound.  Then give a closed-form
stability theorem for \(J_{K,n}=B_n^*A_n^+B_n\) under a fixed
interior spectral gap, while recording the gap-collapse
obstruction.
\end{decision}

\begin{firewall}[Claim boundary after V15]
\label{fire:v16v15}
V15 does not establish the exterior-influence moment, a uniform
interior gap, the interacting thermodynamic state, reflection
positivity of its local marginals, or the full continuum.
\end{firewall}

\section*{Version V15 append-only record}
\addcontentsline{toc}{section}{Version V15 append-only record}
The complete V14 source was retained before this continuation.  It
had \(142917\) bytes, \(4195\) lines, and SHA--256
\[
 \texttt{C30844FC12E5F75F52D5716D8AB28BAE8147F701A317BD9ED9FC4392F202BB5A}.
\]
V15 inspected YM V46 and NS V26, then derived the exact
interior--exterior Schur influence form and its conditional Newton
partition estimates.
\section{Sixteenth-cycle V16: integrated local moments and stability of the influence form}
\label{sec:v16v16}

V15 bounds the Newton-tilted conditional interior law at a fixed
exterior configuration.  V16 integrates that estimate and proves
cofinal stability of
\[
 J_K=B^*A^+B.                                        \tag{16.1}
\]
The two inputs are logically separate: an exponential moment of
the actual exterior influence and a fixed gap for the interior
Newton block.

\subsection{Integration against the tilted exterior marginal}

Keep the finite-cutoff notation of V15.  Let
\[
 Q(u,v)=\langle(u,v),G(u,v)\rangle                   \tag{16.2}
\]
and suppose
\[
 Z_s:=\int e^{sQ(u,v)}\,d\rho(u,v)<\infty.           \tag{16.3}
\]
Define the normalized tilted law
\[
 d\widehat\rho_s=Z_s^{-1}e^{sQ}\,d\rho               \tag{16.4}
\]
and let \(\widehat\rho_{E,s}\) be its exterior marginal.

\begin{theorem}[Exterior influence loader for a local Newton moment]
\label{thm:v16v16-loader}
Fix \(0<\eta\leq1\) and \(t\geq0\).  If
\[
 M_K(t+s(1+\eta))<\infty                             \tag{16.5}
\]
and
\[
 L_{K,s,\eta}
 :=
 \int\exp\{2s\eta^{-1}J_K(v)\}\,
       d\widehat\rho_{E,s}(v)<\infty,                \tag{16.6}
\]
then
\[
 \boxed{
 \int e^{tQ_K(u)}\,d\widehat\rho_s(u,v)
 \leq
 M_K(t+s(1+\eta))\,L_{K,s,\eta}.}                   \tag{16.7}
\]
\end{theorem}

\begin{proof}
Disintegrate \(\widehat\rho_s\) over its exterior marginal.  Its
conditional interior law at \(v\) is exactly
\(\widehat\rho_{K,s}(\,\cdot\mid v)\) from V15, because the
exterior-only factor \(e^{sQ_E(v)}\) cancels in conditional
normalization.  Therefore
\[
\begin{split}
 \int e^{tQ_K}\,d\widehat\rho_s
 &=
 \int\left[
  \int e^{tQ_K(u)}
       d\widehat\rho_{K,s}(u\mid v)\right]
       d\widehat\rho_{E,s}(v)\\
 &\leq
 M_K(t+s(1+\eta))
 \int e^{2s\eta^{-1}J_K(v)}
       d\widehat\rho_{E,s}(v),                       \tag{16.8}
\end{split}
\]
where Corollary~\ref{cor:v16v15-tilted} gives the inequality.
This is (16.7).
\end{proof}

\begin{corollary}[Uniform local-state Newton moment]
\label{cor:v16v16-uniform}
Let the cutoff and volume be indexed by \(\alpha\), while \(K\) is
fixed.  If, for some \(s,t>0\) and \(0<\eta\leq1\),
\[
 \sup_\alpha M_{K,\alpha}(t+s(1+\eta))<\infty,
\qquad
 \sup_\alpha L_{K,s,\eta,\alpha}<\infty,             \tag{16.9}
\]
then
\[
 \sup_\alpha
 {\mathbb E}_{\widehat\rho_{s,\alpha}}
 e^{tQ_{K,\alpha}}<\infty.                           \tag{16.10}
\]
Consequently \(\{Q_{K,\alpha}\}\) is uniformly integrable and has
uniform exponential tails under the normalized local states.
\end{corollary}

\begin{proof}
The first assertion is (16.7) uniformly in \(\alpha\).  For
\(R>0\),
\[
 {\mathbb P}(Q_{K,\alpha}>R)
 \leq e^{-tR}
 {\mathbb E}e^{tQ_{K,\alpha}},                       \tag{16.11}
\]
and the same bound with
\(Q_{K,\alpha}{\bf1}_{\{Q_{K,\alpha}>R\}}\) obtained by integrating
the tail proves uniform integrability.
\end{proof}

The exterior marginal in (16.6) is the physical tilted marginal,
not the reference exterior law.  This prevents a circular
replacement of the desired interacting estimate by a Gaussian
one.

\subsection{Cofinal stability under an interior gap}

Let \({\cal H}_K,{\cal H}_E\) now be fixed finite or separable
Hilbert carriers.  Let
\[
 A_n,A\geq0\quad\hbox{on }{\cal H}_K,\qquad
 B_n,B:{\cal H}_E\to{\cal H}_K.                      \tag{16.12}
\]
Assume the block range conditions
\[
 \operatorname{Ran}B_n\subset\operatorname{Ran}A_n,
\qquad
 \operatorname{Ran}B\subset\operatorname{Ran}A,      \tag{16.13}
\]
as supplied by positivity at finite cutoff.

\begin{theorem}[Operator stability of the exterior influence]
\label{thm:v16v16-stability}
Suppose there is a fixed \(\beta>0\) such that
\[
 \sigma(A_n),\sigma(A)
 \subset\{0\}\cup[\beta,\infty),                     \tag{16.14}
\]
and
\[
 \|A_n-A\|_{\rm op}\to0,\qquad
 \|B_n-B\|_{\rm op}\to0.                             \tag{16.15}
\]
Then
\[
 A_n^+\longrightarrow A^+\quad\hbox{in operator norm} \tag{16.16}
\]
and
\[
 J_{K,n}:=B_n^*A_n^+B_n
 \longrightarrow
 J_K:=B^*A^+B
 \quad\hbox{in operator norm}.                       \tag{16.17}
\]
More precisely,
\[
\begin{split}
 \|J_{K,n}-J_K\|_{\rm op}
 \leq{}&
 \|B_n-B\|\,\|A_n^+\|\,\|B_n\|\\
 &+\|B\|\,\|A_n^+-A^+\|\,\|B_n\|\\
 &+\|B\|\,\|A^+\|\,\|B_n-B\|.                       \tag{16.18}
\end{split}
\]
\end{theorem}

\begin{proof}
The fixed gap (16.14) permits one contour enclosing the positive
spectrum and excluding zero.  Holomorphic functional calculus for
the function \(z\mapsto z^{-1}\) on that component, together with
the resolvent identity and (16.15), gives (16.16).  This is also
the pseudoinverse convergence theorem proved in V8.

Insert and subtract \(B^*A_n^+B_n\) and
\(B^*A^+B_n\) between \(B_n^*A_n^+B_n\) and
\(B^*A^+B\).  The operator ideal inequality gives (16.18).
The factors \(\|A_n^+\|\) are at most \(\beta^{-1}\), the
\(B_n\) are bounded by (16.15), and every difference tends to
zero.  This proves (16.17).
\end{proof}

\begin{corollary}[Random influence-energy convergence]
\label{cor:v16v16-random}
Let exterior random variables \(V_n,V\) be defined on one
probability space and suppose
\[
 V_n\to V\quad\hbox{in }L^2({\cal H}_E),\qquad
 \sup_n{\mathbb E}\|V_n\|^2<\infty.                  \tag{16.19}
\]
Under Theorem~\ref{thm:v16v16-stability},
\[
 \langle V_n,J_{K,n}V_n\rangle
 \longrightarrow
 \langle V,J_KV\rangle
 \quad\hbox{in }L^1.                                \tag{16.20}
\]
\end{corollary}

\begin{proof}
Write the difference as
\[
\begin{split}
 &\langle V_n-V,J_{K,n}V_n\rangle
 +\langle V,(J_{K,n}-J_K)V_n\rangle\\
 &\hspace{7em}
 +\langle V,J_K(V_n-V)\rangle.                       \tag{16.21}
\end{split}
\]
The operators \(J_{K,n}\) are uniformly bounded by (16.17).
Cauchy--Schwarz, (16.19), and operator-norm convergence make the
expectation of every term tend to zero.
\end{proof}

\begin{proposition}[Exponential influence passage]
\label{prop:v16v16-exponential}
In addition to convergence in probability of the influence
energies in (16.20), suppose that for some \(c,\varepsilon>0\),
\[
 \sup_n{\mathbb E}
 \exp\{(c+\varepsilon)
       \langle V_n,J_{K,n}V_n\rangle\}<\infty.        \tag{16.22}
\]
Then
\[
 {\mathbb E}e^{c\langle V_n,J_{K,n}V_n\rangle}
 \longrightarrow
 {\mathbb E}e^{c\langle V,J_KV\rangle}.              \tag{16.23}
\]
\end{proposition}

\begin{proof}
The continuous exponential maps convergence in probability of the
energies to convergence in probability of the exponentials.
Condition (16.22) makes those exponentials uniformly integrable by
the de la Vall\'ee--Poussin criterion.  Vitali's theorem proves
(16.23).
\end{proof}

\subsection{Gap collapse is a real obstruction}

\begin{counterexample}[Vanishing blocks with nonvanishing influence]
\label{cnt:v16v16-gap}
On one-dimensional interior and exterior spaces, let
\[
 A_n=n^{-2},\qquad B_n=n^{-1},\qquad C_n=1.           \tag{16.24}
\]
Then
\[
 G_n=
 \begin{pmatrix}n^{-2}&n^{-1}\\n^{-1}&1\end{pmatrix}
 \succeq0,\qquad
 G_n\longrightarrow
 \begin{pmatrix}0&0\\0&1\end{pmatrix},               \tag{16.25}
\]
but
\[
 J_{K,n}=B_n^*A_n^+B_n=1                             \tag{16.26}
\]
for every \(n\), whereas the limiting blocks \(A=B=0\) give
\(J_K=0\).
\end{counterexample}

\begin{proof}
The matrix \(G_n\) is the rank-one outer product
\((n^{-1},1)^{\mathsf T}(n^{-1},1)\).  Direct substitution gives
(16.26).
\end{proof}

Thus norm convergence of the full positive Newton block does not
control the Schur influence when the interior gap closes.  A
stratified kernel analysis or an explicitly retained soft mode is
required, as in V8.

\begin{assessment}[Physical status]
\label{ass:v16v16-status}
V16 turns the V15 conditional estimate into a uniform local-state
moment theorem and proves operator, \(L^1\), and exponential-moment
stability of the exterior influence under a fixed interior gap.
It also gives an exact positive-block counterexample at gap
collapse.

The physical work is now concentrated in the loader (16.9): a
uniform conditional interior moment and a uniform exponential
moment of \(J_{K,n}\) under the actual tilted exterior marginal.
\end{assessment}

\begin{decision}[V17 collar and tail decomposition]
\label{dec:v16v16-next}
Split \(B_n\) into a finite collar coupling and a far-exterior tail.
Derive a deterministic bound for the corresponding change in
\(J_{K,n}\) under the gap \(\beta\), then state a summable
off-diagonal decay condition on the Newton inverse that makes the
far influence vanish uniformly.  Keep the massless \(1/r\) kernel
separate, because it is not absolutely summable without charge
cancellation.
\end{decision}

\begin{firewall}[Claim boundary after V16]
\label{fire:v16v16}
V16 does not prove the physical exterior exponential moment or the
interior gap, control a gap-closing mode, establish the required
off-diagonal decay or charge cancellation, or construct the full
continuum.
\end{firewall}

\section*{Version V16 append-only record}
\addcontentsline{toc}{section}{Version V16 append-only record}
The complete V15 source was retained before this continuation.  It
had \(151683\) bytes, \(4472\) lines, and SHA--256
\[
 \texttt{8F8231C55C29EE1B882F96EA6F64DD93A7E1CB904106A7F9061EEE4C660B5E99}.
\]
V16 integrated the conditional Newton estimate, proved cofinal
stability of the Schur influence form under a fixed gap, and
exhibited its discontinuity when that gap collapses.
\section{Sixteenth-cycle V17: Newton collars, monopole cancellation, and the far influence tail}
\label{sec:v16v17}

V16 leaves the exterior exponential moment as a physical loader.
Before estimating that moment, one must know whether the far
exterior coupling is small in an operator topology.  This section
proves an abstract collar estimate and then treats the
three-dimensional massless Newton kernel explicitly.  The
unprojected \(1/r\) coupling is not Hilbert--Schmidt at infinity;
the interior mean-zero projection removes its monopole and changes
the far kernel to order \(1/r^2\), which is square summable.

\subsection{Abstract collar perturbation}

Let \(A\geq0\) on the interior carrier \({\cal H}_K\), and assume
\[
 \sigma(A)\subset\{0\}\cup[\beta,\infty),
 \qquad \beta>0.                                     \tag{17.1}
\]
Let \(B=B_R+D_R:{\cal H}_E\to{\cal H}_K\), where \(B_R\) is the
coupling from an \(R\)-collar and \(D_R\) the remaining far
coupling.  Assume their ranges lie in \(\operatorname{Ran}A\), and
write
\[
 {\cal J}(L):=L^*A^+L.                               \tag{17.2}
\]

\begin{theorem}[Influence stability under collar exhaustion]
\label{thm:v16v17-abstract}
One has
\[
\begin{split}
 \|{\cal J}(B)-{\cal J}(B_R)\|_{\rm op}
 \leq
 \beta^{-1}\left(
 2\|B_R\|_{\rm op}\|D_R\|_{\rm op}
 +\|D_R\|_{\rm op}^2\right).                         \tag{17.3}
\end{split}
\]
For every \(\eta>0\), as quadratic forms on \({\cal H}_E\),
\[
 {\cal J}(B)
 \preceq
 (1+\eta){\cal J}(B_R)
 +(1+\eta^{-1}){\cal J}(D_R),                        \tag{17.4}
\]
and
\[
 {\cal J}(D_R)\preceq
 \beta^{-1}D_R^*D_R.                                 \tag{17.5}
\]
Consequently, if \(\sup_R\|B_R\|<\infty\) and
\(\|D_R\|\to0\), then
\[
 {\cal J}(B_R)\longrightarrow{\cal J}(B)
 \quad\hbox{in operator norm}.                       \tag{17.6}
\]
\end{theorem}

\begin{proof}
Expansion gives
\[
 {\cal J}(B)-{\cal J}(B_R)
 =
 B_R^*A^+D_R+D_R^*A^+B_R+D_R^*A^+D_R.               \tag{17.7}
\]
Since \(\|A^+\|\leq\beta^{-1}\), the triangle and ideal
inequalities prove (17.3).

For \(v\in{\cal H}_E\),
\[
 \langle v,{\cal J}(B)v\rangle
 =
 \|A^{+1/2}(B_R+D_R)v\|^2.                           \tag{17.8}
\]
The Hilbert-space inequality
\(\|x+y\|^2\leq(1+\eta)\|x\|^2+
(1+\eta^{-1})\|y\|^2\) proves (17.4).
The gap gives
\[
 \|A^{+1/2}D_Rv\|^2
 \leq\beta^{-1}\|D_Rv\|^2,                           \tag{17.9}
\]
which is (17.5).  The last assertion follows from (17.3).
\end{proof}

\subsection{The projected Newton kernel in three dimensions}

Let \(K\subset\mathbb R^3\) be bounded and measurable, with
\(0<|K|<\infty\).  Choose \(x_0\in K\) and \(D>0\) such that
\[
 K\subset\overline B(x_0,D).                         \tag{17.10}
\]
Let
\[
 {\cal H}_K=L^2_0(K)
 :=\left\{u\in L^2(K):\int_Ku(x)\,dx=0\right\},      \tag{17.11}
\]
and let \(P_0:L^2(K)\to L^2_0(K)\) be the orthogonal projection.
The Newton kernel is
\[
 N(x,y)={1\over4\pi|x-y|}.                           \tag{17.12}
\]
For \(R\geq2D\), set
\[
 E_R:=\{y:|y-x_0|>R\}                                \tag{17.13}
\]
and define the far coupling
\[
 (D_Rf)(x)
 :=
 P_0\int_{E_R}N(x,y)f(y)\,dy.                        \tag{17.14}
\]

\begin{theorem}[Mean-zero Newton tail is Hilbert--Schmidt]
\label{thm:v16v17-HS-tail}
The operator
\[
 D_R:L^2(E_R)\longrightarrow L^2_0(K)               \tag{17.15}
\]
is Hilbert--Schmidt and obeys
\[
 \|D_R\|_{\rm op}
 \leq\|D_R\|_2
 \leq
 2D\sqrt{|K|\over\pi R}.                             \tag{17.16}
\]
\end{theorem}

\begin{proof}
For fixed \(y\in E_R\), put \(r=|y-x_0|\).  Along the segment from
\(x_0\) to any \(x\in K\), the distance to \(y\) is at least
\(r-D\geq r/2\).  Since
\[
 |\nabla_xN(x,y)|={1\over4\pi|x-y|^2},               \tag{17.17}
\]
the mean-value theorem gives
\[
 |N(x,y)-N(x_0,y)|
 \leq {D\over\pi r^2}.                               \tag{17.18}
\]
Orthogonal projection gives the best \(L^2(K)\) approximation by
a constant, so
\[
 \|P_0N(\,\cdot\,,y)\|_{L^2(K)}
 \leq
 \|N(\,\cdot\,,y)-N(x_0,y)\|_{L^2(K)}
 \leq{D|K|^{1/2}\over\pi r^2}.                       \tag{17.19}
\]
Therefore
\[
\begin{split}
 \|D_R\|_2^2
 &=
 \int_{E_R}\|P_0N(\,\cdot\,,y)\|_{L^2(K)}^2\,dy\\
 &\leq
 {D^2|K|\over\pi^2}
 4\pi\int_R^\infty r^{-2}\,dr
 ={4D^2|K|\over\pi R}.                               \tag{17.20}
\end{split}
\]
Taking square roots proves (17.16).
\end{proof}

\begin{corollary}[Explicit influence-tail convergence]
\label{cor:v16v17-influence}
In the setting of Theorem~\ref{thm:v16v17-abstract}, suppose the far
block is the projected Newton coupling (17.14) and
\[
 \sup_R\|B_R\|_{\rm op}\leq b.                       \tag{17.21}
\]
Then
\[
\begin{split}
 \|{\cal J}(B)-{\cal J}(B_R)\|_{\rm op}
 \leq{\beta^{-1}}\left[
 4bD\sqrt{|K|\over\pi R}
 +{4D^2|K|\over\pi R}\right].                        \tag{17.22}
\end{split}
\]
\end{corollary}

\begin{proof}
Insert (17.16) into (17.3).
\end{proof}

\subsection{The uncancelled monopole fails}

Let \(\widetilde D_R:L^2(E_R)\to L^2(K)\) have the unprojected
kernel \(N(x,y)\).

\begin{proposition}[The massless monopole tail is not
Hilbert--Schmidt]
\label{prop:v16v17-monopole}
For every finite \(R\),
\[
 \|\widetilde D_R\|_2=\infty.                        \tag{17.23}
\]
\end{proposition}

\begin{proof}
For \(r=|y-x_0|\geq\max\{R,2D\}\) and \(x\in K\),
\[
 |x-y|\leq r+D\leq{3r\over2}.                        \tag{17.24}
\]
Hence
\[
 N(x,y)^2\geq{1\over36\pi^2r^2}.                    \tag{17.25}
\]
It follows that
\[
\begin{split}
 \|\widetilde D_R\|_2^2
 &\geq
 {|K|\over36\pi^2}
 4\pi\int_{\max\{R,2D\}}^\infty dr
 =\infty.                                            \tag{17.26}
\end{split}
\]
\end{proof}

The projected kernel has order \(r^{-2}\), which is square
integrable but still not absolutely integrable in three dimensions.
Thus (17.16) is an \(L^2\)-operator tail, not an \(L^1\) cluster
expansion estimate.

\subsection{Interpretation of the cancellation}

For \(u\in L^2_0(K)\),
\[
\int_Ku(x)N(x,y)\,dx
=
\int_Ku(x)[N(x,y)-N(x_0,y)]\,dx.                    \tag{17.27}
\]
The local zero integral removes the far monopole exactly.  A global
mean-zero condition on the complete source does not imply
\(\int_Ku=0\) for each fixed region.  Therefore the theorem applies
directly to neutral local insertions, divergence-form sources, or a
formulation in which the interior constant mode is retained and
treated separately.  It cannot be imposed on arbitrary local
stress fluctuations without an additional monopole ledger.

\begin{assessment}[Physical status]
\label{ass:v16v17-status}
V17 proves a collar-exhaustion estimate for the Schur influence and
an explicit \(R^{-1/2}\) Hilbert--Schmidt tail for the
three-dimensional Newton inverse after interior monopole
cancellation.  It also proves that the uncancelled massless tail is
not Hilbert--Schmidt.

This supplies deterministic smallness of the far block for neutral
local sources.  It does not yet provide the exponential
\(J_K\)-moment under the interacting exterior marginal, and
arbitrary stress insertions still carry a constant interior mode.
\end{assessment}

\begin{decision}[V18 Gaussian far-influence moment]
\label{dec:v16v17-next}
Use the V4 determinant criterion on the exterior Gaussian reference
with quadratic form \({\cal J}(D_R)\).  Derive an explicit spectral
and trace bound from (17.16), prove that the far-influence
exponential moment tends to one as \(R\to\infty\), and identify the
additional covariance-weighted condition needed for infinite
exterior volume.
\end{decision}

\begin{firewall}[Claim boundary after V17]
\label{fire:v16v17}
V17 does not prove local monopole cancellation for the complete
Standard Model stress, an interacting exterior moment, an
infinite-volume Newton transfer, reflection-positive local-state
convergence, or the full SM--Einstein continuum.
\end{firewall}

\section*{Version V17 append-only record}
\addcontentsline{toc}{section}{Version V17 append-only record}
The complete V16 source was retained before this continuation.  It
had \(161129\) bytes, \(4784\) lines, and SHA--256
\[
 \texttt{41E8F9C74FD01E3F3515B8FB8481152BAE1676E089739B5869056C26212EE689}.
\]
V17 proved the abstract collar perturbation bound, the
mean-zero \(R^{-1/2}\) Newton tail, and the failure of the
uncancelled massless monopole tail.
\section{Sixteenth-cycle V18: Gaussian exponential control of the far Newton influence}
\label{sec:v16v18}

V17 proves that the neutral far coupling \(D_R\) is
Hilbert--Schmidt.  The exterior influence it creates is
\[
 J_R:=D_R^*A^+D_R.                                   \tag{18.1}
\]
V18 applies the exact Gaussian quadratic-form determinant of V4 to
this operator.  The resulting bound needs only a uniform covariance
operator norm, not a volume-uniform covariance trace.

\subsection{Covariance-weighted influence}

At each finite exterior volume \(\Lambda\), let \(V_\Lambda\) be a
centered Gaussian exterior source with covariance
\(C_\Lambda\geq0\).  Assume \(C_\Lambda\) is trace class at finite
volume, and let \(A\) satisfy the interior gap (17.1).  Define
\[
 K_{R,\Lambda}
 :=
 C_\Lambda^{1/2}J_RC_\Lambda^{1/2}
 =
 C_\Lambda^{1/2}D_R^*A^+D_RC_\Lambda^{1/2}.          \tag{18.2}
\]

\begin{lemma}[Trace and spectral bounds]
\label{lem:v16v18-K}
The operator \(K_{R,\Lambda}\) is positive trace class and
\[
\begin{split}
 \operatorname{Tr}K_{R,\Lambda}
 &=
 \|A^{+1/2}D_RC_\Lambda^{1/2}\|_2^2\\
 &\leq\beta^{-1}\|D_RC_\Lambda^{1/2}\|_2^2,          \tag{18.3}\\
 \|K_{R,\Lambda}\|_{\rm op}
 &\leq\beta^{-1}
       \|D_RC_\Lambda^{1/2}\|_{\rm op}^2.            \tag{18.4}
\end{split}
\]
If \(\|C_\Lambda\|_{\rm op}\leq c_C\), then
\[
 \operatorname{Tr}K_{R,\Lambda}
 \leq\beta^{-1}c_C\|D_R\|_2^2,\qquad
 \|K_{R,\Lambda}\|
 \leq\beta^{-1}c_C\|D_R\|_{\rm op}^2.                \tag{18.5}
\]
\end{lemma}

\begin{proof}
Set \(L=A^{+1/2}D_RC_\Lambda^{1/2}\).  Because \(D_R\) is
Hilbert--Schmidt and the other factors are bounded, \(L\) is
Hilbert--Schmidt.  Equation (18.2) is \(K_{R,\Lambda}=L^*L\),
which is positive trace class and has trace \(\|L\|_2^2\).
The gap bound \(\|A^{+1/2}\|\leq\beta^{-1/2}\) proves (18.3);
the same operator-norm calculation proves (18.4).
Finally,
\[
 \|D_RC_\Lambda^{1/2}\|_2
 \leq\|D_R\|_2\|C_\Lambda^{1/2}\|_{\rm op},          \tag{18.6}
\]
and similarly in operator norm, proving (18.5).
\end{proof}

\begin{theorem}[Exact Gaussian far-influence moment]
\label{thm:v16v18-determinant}
For \(c>0\),
\[
 {\mathbb E}
 \exp\{c\langle V_\Lambda,J_RV_\Lambda\rangle\}
 <\infty                                             \tag{18.7}
\]
if and only if
\[
 2c\|K_{R,\Lambda}\|_{\rm op}<1.                    \tag{18.8}
\]
Under (18.8),
\[
 {\mathbb E}
 \exp\{c\langle V_\Lambda,J_RV_\Lambda\rangle\}
 =
 \det(I-2cK_{R,\Lambda})^{-1/2}.                     \tag{18.9}
\]
If \(2c\|K_{R,\Lambda}\|\leq q<1\), then
\[
 1\leq
 {\mathbb E}e^{c\langle V_\Lambda,J_RV_\Lambda\rangle}
 \leq
 \exp\left\{
 {c\,\operatorname{Tr}K_{R,\Lambda}\over1-q}\right\}.
                                                               \tag{18.10}
\]
\end{theorem}

\begin{proof}
Apply Theorem~\ref{thm:v16v4-gaussian} to the Gaussian variable
\(V_\Lambda\), quadratic operator \(J_R\), and
\(\Gamma=2c\).  This gives (18.8)--(18.9).
If \(\kappa_j\) are the eigenvalues of \(K_{R,\Lambda}\), then
\[
 -{1\over2}\log\det(I-2cK_{R,\Lambda})
 ={1\over2}\sum_j-\log(1-2c\kappa_j)
 \leq {c\over1-q}\sum_j\kappa_j,                    \tag{18.11}
\]
which proves the upper bound.  The lower bound follows from
\(J_R\geq0\).
\end{proof}

\subsection{Explicit uniform Newton estimate}

Use the geometry of Theorem~\ref{thm:v16v17-HS-tail}; thus
\[
 \|D_R\|_2^2\leq{4D^2|K|\over\pi R}.                 \tag{18.12}
\]

\begin{corollary}[Far Gaussian moment tends uniformly to one]
\label{cor:v16v18-uniform}
Suppose
\[
 \sup_\Lambda\|C_\Lambda\|_{\rm op}\leq c_C<\infty. \tag{18.13}
\]
Define
\[
 q_R:={8c\,c_CD^2|K|\over\beta\pi R}.                \tag{18.14}
\]
Whenever \(q_R<1\),
\[
\sup_\Lambda
 {\mathbb E}
 e^{c\langle V_\Lambda,J_RV_\Lambda\rangle}
 \leq
 \exp\left\{
 {4c\,c_CD^2|K|\over
  \beta\pi R(1-q_R)}\right\}.                        \tag{18.15}
\]
In particular,
\[
 \sup_\Lambda
 \left|
 {\mathbb E}
 e^{c\langle V_\Lambda,J_RV_\Lambda\rangle}-1
 \right|
 \longrightarrow0
 \quad(R\to\infty).                                  \tag{18.16}
\]
\end{corollary}

\begin{proof}
Lem\-ma~\ref{lem:v16v18-K} and (18.12) give
\[
 \operatorname{Tr}K_{R,\Lambda}
 \leq{4c_CD^2|K|\over\beta\pi R},                   \tag{18.17}
\]
and
\[
 2c\|K_{R,\Lambda}\|
 \leq2c\beta^{-1}c_C\|D_R\|_{\rm op}^2
 \leq q_R.                                           \tag{18.18}
\]
Apply (18.10) and take the supremum over \(\Lambda\).
The right side of (18.15) tends to one, while every expectation is
at least one, proving (18.16).
\end{proof}

No bound on \(\operatorname{Tr}C_\Lambda\) occurs in (18.15).
That trace normally grows with exterior volume.  The
Hilbert--Schmidt localization of \(D_R\) tests the covariance only
through the directions that reach the fixed interior.

\begin{theorem}[General covariance-weighted tail card]
\label{thm:v16v18-general}
The geometric estimate (18.12) may be replaced by the two
conditions
\[
\begin{split}
 \varepsilon_R
 &:=
 \sup_\Lambda
 \|D_RC_\Lambda^{1/2}\|_2^2\longrightarrow0,\\
 \zeta_R
 &:=
 \sup_\Lambda
 \|D_RC_\Lambda^{1/2}\|_{\rm op}^2\longrightarrow0.
                                                               \tag{18.19}
\end{split}
\]
Then, for every fixed \(c>0\), the moments (18.7) are finite for
all sufficiently large \(R\) and converge to one uniformly in
\(\Lambda\).  More precisely, if
\(2c\beta^{-1}\zeta_R<1\), they are bounded by
\[
 \exp\left\{
 {c\beta^{-1}\varepsilon_R\over
  1-2c\beta^{-1}\zeta_R}\right\}.                    \tag{18.20}
\]
\end{theorem}

\begin{proof}
Equations (18.3)--(18.4) give
\[
 \sup_\Lambda\operatorname{Tr}K_{R,\Lambda}
 \leq\beta^{-1}\varepsilon_R,\qquad
 \sup_\Lambda\|K_{R,\Lambda}\|
 \leq\beta^{-1}\zeta_R.                              \tag{18.21}
\]
Apply Theorem~\ref{thm:v16v18-determinant}.
\end{proof}

\subsection{Interacting exterior comparison}

\begin{corollary}[\(L^p\) exterior perturbation]
\label{cor:v16v18-Lp}
Let an exterior law \(\nu_{\Lambda}\) have density
\(h_\Lambda\) relative to the Gaussian law above, with
\[
 \|h_\Lambda\|_{L^p}\leq H,\qquad q={p\over p-1}.    \tag{18.22}
\]
If
\[
 2qc\|K_{R,\Lambda}\|<1,                             \tag{18.23}
\]
then
\[
 {\mathbb E}_{\nu_\Lambda}
 e^{c\langle V_\Lambda,J_RV_\Lambda\rangle}
 \leq
 H\det(I-2qcK_{R,\Lambda})^{-1/(2q)}.                \tag{18.24}
\]
Under the uniform covariance condition (18.13), the determinant
factor tends to one uniformly as \(R\to\infty\).
\end{corollary}

\begin{proof}
Apply Theorem~\ref{thm:v16v13-dual} with quadratic operator \(J_R\),
then use the bounds above with \(c\) replaced by \(qc\).
\end{proof}

The factor \(H\) in (18.24) need not tend to one and a global
volume-uniform \(H\) generally fails by V14.  Thus this corollary
does not yet prove the required moment under the physical tilted
exterior marginal.  It identifies the remaining input as a
localized comparison bound for the variables seen by
\(D_R\), rather than a global exterior density norm.

\begin{assessment}[Physical status]
\label{ass:v16v18-status}
V18 proves that the Gaussian far-influence exponential moment tends
to one uniformly in exterior volume under a uniform covariance
operator bound, with an explicit \(O(R^{-1})\) logarithmic estimate.
It also gives the exact covariance-weighted tail card (18.19) and
the \(L^p\) interacting comparison.

The physical tilted exterior law still lacks the localized density
or conditional estimate needed to replace the global factor \(H\).
The result applies only after the local monopole cancellation and
interior gap established as hypotheses in V17.
\end{assessment}

\begin{decision}[V19 projected-density localization]
\label{dec:v16v18-next}
Push the exterior law through the finite or Hilbert--Schmidt
observation map \(D_R\).  Prove that only the Radon--Nikodym density
of this projected law relative to its projected Gaussian reference
enters (18.24), and that conditional expectation contracts its
\(L^p\) norm.  This will replace the invalid global exterior
density card by an influence-space density card.
\end{decision}

\begin{firewall}[Claim boundary after V18]
\label{fire:v16v18}
V18 does not prove the projected interacting density bound, the
physical local monopole ledger or interior gap, reflection-positive
thermodynamic convergence, or the full SM--Einstein continuum.
\end{firewall}

\section*{Version V18 append-only record}
\addcontentsline{toc}{section}{Version V18 append-only record}
The complete V17 source was retained before this continuation.  It
had \(169439\) bytes, \(5068\) lines, and SHA--256
\[
 \texttt{495A1E910F35C41E22028CED5522426BDC876E7663FCB9243B1A139D19C81A2B}.
\]
V18 converted the projected Newton tail into an explicit
volume-uniform Gaussian determinant bound and isolated the
covariance-weighted and interacting comparison conditions.
\section{Sixteenth-cycle V19: projected interaction densities close the far-influence tail}
\label{sec:v16v19}

The far influence
\[
 \langle V,J_RV\rangle
 =\langle D_RV,A^+D_RV\rangle                       \tag{19.1}
\]
depends on the exterior field only through the observation
\(Y_R=D_RV\).  Therefore a global density comparison is stronger
than necessary.  V19 passes the interacting law to this influence
space and proves that a uniform projected \(L^p\) bound makes the
far exponential moment converge to one.

\subsection{Density under an observation map}

Let \((X,{\cal F},\pi)\) be a probability space,
\(h\geq0\), \(\int h\,d\pi=1\), and \(d\nu=h\,d\pi\).
Let \(Y:X\to{\cal Y}\) be measurable, where \({\cal Y}\) is a
standard Borel space.  Write
\[
 \pi_Y=Y_\#\pi,\qquad \nu_Y=Y_\#\nu.                 \tag{19.2}
\]

\begin{lemma}[Projected Radon--Nikodym density]
\label{lem:v16v19-projected-density}
The measure \(\nu_Y\) is absolutely continuous with respect to
\(\pi_Y\), with density
\[
 \bar h(Y)
 =
 {\mathbb E}_\pi[h\mid\sigma(Y)]                    \tag{19.3}
\]
up to the usual almost-everywhere identification.  For
\(1\leq p\leq\infty\),
\[
 \|\bar h\|_{L^p(\pi_Y)}
 \leq\|h\|_{L^p(\pi)}.                               \tag{19.4}
\]
\end{lemma}

\begin{proof}
For every bounded measurable \(\varphi\),
\[
\begin{split}
 \int_{\cal Y}\varphi(y)\,d\nu_Y(y)
 &=\int_X\varphi(Y)h\,d\pi\\
 &=\int_X\varphi(Y)
       {\mathbb E}_\pi[h\mid\sigma(Y)]\,d\pi,
                                                               \tag{19.5}
\end{split}
\]
which proves (19.3).  Conditional Jensen proves (19.4) for finite
\(p\), and the essential-supremum contraction proves it for
\(p=\infty\).
\end{proof}

The usefulness of (19.3) is not merely the contraction (19.4).
The projected norm may remain bounded when the global norm grows
exponentially with volume, as in V14.

\subsection{Gaussian law on the influence space}

Let \(V\) be centered Gaussian with covariance \(C\), let
\[
 Y_R=D_RV,\qquad
 \Sigma_R=D_RCD_R^*,                                 \tag{19.6}
\]
and put
\[
 H_R:=\Sigma_R^{1/2}A^+\Sigma_R^{1/2}.               \tag{19.7}
\]
Recall
\[
 K_R=C^{1/2}D_R^*A^+D_RC^{1/2}.                     \tag{19.8}
\]

\begin{proposition}[Equality of the two determinant carriers]
\label{prop:v16v19-carriers}
The positive trace-class operators \(H_R\) and \(K_R\) have the
same nonzero eigenvalues with multiplicity.  Consequently,
\[
 \det(I-zH_R)=\det(I-zK_R)                           \tag{19.9}
\]
whenever either Fredholm determinant is defined.
\end{proposition}

\begin{proof}
Let
\[
 L=A^{+1/2}D_RC^{1/2}.                               \tag{19.10}
\]
Then \(K_R=L^*L\).  The covariance of \(Y_R\) is
\(\Sigma_R=D_RCD_R^*\), and the nonzero spectrum of
\(\Sigma_R^{1/2}A^+\Sigma_R^{1/2}\) agrees with that of
\(A^{+1/2}\Sigma_RA^{+1/2}=LL^*\).  The compact positive
operators \(L^*L\) and \(LL^*\) have the same nonzero eigenvalues,
including multiplicity.  Their Fredholm products therefore agree.
\end{proof}

\begin{theorem}[Exact projected \(L^p\) comparison]
\label{thm:v16v19-projected-comparison}
Let the interacting exterior law \(\nu\) have projected density
\(\bar h_R=d(Y_R{}_\#\nu)/d(Y_R{}_\#\pi)\), and suppose
\[
 \|\bar h_R\|_{L^p(Y_R{}_\#\pi)}\leq H_R^{(p)},
 \qquad q={p\over p-1}.                              \tag{19.11}
\]
If
\[
 2qc\|K_R\|_{\rm op}<1,                              \tag{19.12}
\]
then
\[
 {\mathbb E}_\nu e^{c\langle V,J_RV\rangle}
 \leq
 H_R^{(p)}
 \det(I-2qcK_R)^{-1/(2q)}.                           \tag{19.13}
\]
\end{theorem}

\begin{proof}
The random quadratic form in (19.1) is
\(\langle Y_R,A^+Y_R\rangle\).  Apply
Theorem~\ref{thm:v16v13-dual} on the Gaussian influence space with
covariance \(\Sigma_R\), quadratic operator \(A^+\), and density
\(\bar h_R\).  Proposition~\ref{prop:v16v19-carriers} replaces its
determinant carrier \(H_R\) by \(K_R\), proving (19.13).
\end{proof}

\subsection{The interacting far moment converges to one}

\begin{theorem}[Projected-density far-tail theorem]
\label{thm:v16v19-tail}
Let the cutoff and exterior volume be indexed by \(\alpha\).  Fix
\(1<p<\infty\), \(q=p/(p-1)\), and \(c>0\).  Suppose
\[
 \sup_{\alpha,R}
 \|\bar h_{R,\alpha}\|_{L^p}\leq H<\infty            \tag{19.14}
\]
for all sufficiently large \(R\), and suppose the Gaussian
covariance-weighted tail card (18.19) holds uniformly in
\(\alpha\).  Then
\[
 \sup_\alpha
 \left|
 {\mathbb E}_{\nu_\alpha}
 e^{c\langle V_\alpha,J_{R,\alpha}V_\alpha\rangle}
 -1\right|
 \longrightarrow0
 \quad(R\to\infty).                                  \tag{19.15}
\]
\end{theorem}

\begin{proof}
Put
\[
 X_{R,\alpha}
 :=c\langle V_\alpha,J_{R,\alpha}V_\alpha\rangle
 \geq0.                                              \tag{19.16}
\]
By the projected density formula and H\"older,
\[
\begin{split}
 0\leq{\mathbb E}_{\nu_\alpha}e^{X_{R,\alpha}}-1
 &=
 {\mathbb E}_{\pi_\alpha}
 [\bar h_{R,\alpha}(e^{X_{R,\alpha}}-1)]\\
 &\leq
 H\left(
 {\mathbb E}_{\pi_\alpha}
 (e^{X_{R,\alpha}}-1)^q\right)^{1/q}.                \tag{19.17}
\end{split}
\]
For \(x\geq0\) and \(q\geq1\),
\[
 (e^x-1)^q\leq e^{qx}-1.                             \tag{19.18}
\]
Theorem~\ref{thm:v16v18-general}, applied with exponent \(qc\),
shows
\[
 \sup_\alpha
 {\mathbb E}_{\pi_\alpha}e^{qX_{R,\alpha}}
 \longrightarrow1.                                  \tag{19.19}
\]
Equations (19.17)--(19.19) prove (19.15).
\end{proof}

\begin{corollary}[Concrete three-dimensional version]
\label{cor:v16v19-three-dimensional}
Under the interior gap \(\beta\), the mean-zero Newton geometry of
V17, the uniform Gaussian covariance bound
\[
 \sup_\alpha\|C_\alpha\|_{\rm op}\leq c_C,           \tag{19.20}
\]
and the projected density bound (19.14), conclusion (19.15) holds.
\end{corollary}

\begin{proof}
Theorem~\ref{thm:v16v17-HS-tail} and (19.20) imply the
covariance-weighted card (18.19), with both quantities
\(O(R^{-1})\).  Apply
Theorem~\ref{thm:v16v19-tail}.
\end{proof}

\subsection{Why this is strictly local information}

The sigma-algebra \(\sigma(D_RV)\) contains only the exterior
combinations that reach the neutral interior through the far
Newton kernel.  Lemma~\ref{lem:v16v19-projected-density} shows that
the relevant comparison density is the conditional average of the
full interacting density over all invisible exterior directions.
Those directions may carry extensive R\'enyi divergence without
entering (19.14).

The bound (19.14) is nevertheless a substantive physical
hypothesis.  Conditional expectation guarantees it from a global
\(L^p\) bound, but V14 proves that the latter is generally
unavailable.  A direct projected estimate must come from locality,
screening, a cluster expansion, or a conditional influence
inequality for the complete normalized SM--Einstein law.

\begin{assessment}[Physical status]
\label{ass:v16v19-status}
V19 replaces the invalid global exterior density card by the exact
projected influence-space density.  Under a uniform \(L^p\) bound
for that density, the interacting far Newton exponential moment
converges to one, even though the full exterior density norm may
grow exponentially with volume.

This closes the far-tail implication conditional on the interior
gap, local monopole cancellation, covariance tail, and projected
density bound.  None of those four physical loaders is yet proved
for the complete coupled theory.
\end{assessment}

\begin{decision}[V20 companion audit and collar law]
\label{dec:v16v19-next}
Perform the compulsory companion audit.  Then turn to the finite
collar \(B_R\): formulate a conditional \(L^p\) density card on the
finite collar variables and combine it with V13 and V16.  This will
separate a finite local interacting estimate from the now-controlled
far tail.
\end{decision}

\begin{firewall}[Claim boundary after V19]
\label{fire:v16v19}
V19 does not prove the projected density hypothesis, the collar
interaction estimate, the physical gap or monopole ledger,
reflection-positive thermodynamic convergence, or the full
continuum.
\end{firewall}

\section*{Version V19 append-only record}
\addcontentsline{toc}{section}{Version V19 append-only record}
The complete V18 source was retained before this continuation.  It
had \(178261\) bytes, \(5353\) lines, and SHA--256
\[
 \texttt{9A6D846C7311581D5D1A2996C85CAB3F5ADC11C111CFEF0F2361665FA82EC9E0}.
\]
V19 proved that only the projected influence-space density enters
the far Newton moment and that a uniform projected \(L^p\) bound
forces the interacting far exponential moment to converge to one.
\section{Sixteenth-cycle V20: companion audit and a noncentral Gaussian collar estimate}
\label{sec:v16v20}

This compulsory checkpoint records the newest YM and NS notes and
then treats the finite Newton collar.  Conditioning a Gaussian
reference on exterior variables generally shifts the collar mean.
The centered determinant of V12 is therefore insufficient by
itself; the mean displacement has an additional exact quadratic
cost.

\subsection{Compulsory companion audit}

\begin{audit}[Files inspected at V20]
\label{aud:v16v20-files}
The newest companion checkpoints remain
\[
\begin{array}{c|r|r}
\hbox{file}&\hbox{lines}&\hbox{bytes}\\ \hline
\texttt{Renewal\_Yang\_Mills\_Mass\_Gap\_Research\_Notes\_V46.tex}
 &10844&396051\\
\texttt{Renewal\_NS\_Stability\_Research\_Notes\_V26.tex}
 &5319&278283.
\end{array}                                             \tag{20.1}
\]
Their SHA--256 values are
\[
\begin{gathered}
\texttt{B1E5580A117A5BED92CDB4AC94DDEB6E5992F93E18D5DB46F0A38FAC6D3E5469},\\
\texttt{82C887F8C2C69E4F28FAD7C3DC8DE5B9099CB5A4BF431EBA1FFF3E91935978AD}.
\end{gathered}                                         \tag{20.2}
\]
\end{audit}

\begin{theorem}[V20 companion-audit decision]
\label{thm:v16v20-audit}
Neither checkpoint supplies a finite-collar conditional density
bound, a Gaussian conditional-mean estimate, or the physical
Newton influence loader.
\end{theorem}

\begin{proof}
YM V46 ends with an exact two-sign Wilson-pencil atlas on paired
two-coordinate plates.  It has no completed third-coordinate
atlas, transfer limit, interacting continuum law, or mass gap.
Those finite response floors do not control the conditional mean
of a gravitational collar.

NS V26 computes rank-one Oseen-kernel derivative norms.  Its
deterministic convolution estimates contain no conditional
Gaussian law or Standard Model interaction density.  Hence no
companion theorem enters the calculation below.
\end{proof}

\subsection{Exact noncentral quadratic moment}

Let \(E\) be a finite-dimensional real Hilbert space.  Let
\[
 X=m+C^{1/2}\xi,                                     \tag{20.3}
\]
where \(\xi\) is standard Gaussian, \(m\in E\), and
\(C\geq0\).  Let \(M\geq0\) and define
\[
 A=C^{1/2}MC^{1/2},\qquad
 b=C^{1/2}Mm.                                        \tag{20.4}
\]

\begin{theorem}[Noncentral Gaussian Newton determinant]
\label{thm:v16v20-noncentral}
For \(s>0\),
\[
 {\mathbb E}\exp\{s\langle X,MX\rangle\}<\infty      \tag{20.5}
\]
if and only if
\[
 2s\|A\|_{\rm op}<1.                                \tag{20.6}
\]
When (20.6) holds,
\[
\begin{split}
 {\mathbb E}e^{s\langle X,MX\rangle}
 ={}&
 \det(I-2sA)^{-1/2}\\
 &\times
 \exp\left\{
 s\langle m,Mm\rangle
 +2s^2\langle b,(I-2sA)^{-1}b\rangle
 \right\}.                                          \tag{20.7}
\end{split}
\]
\end{theorem}

\begin{proof}
Expansion gives
\[
 \langle X,MX\rangle
 =
 \langle\xi,A\xi\rangle
 +2\langle b,\xi\rangle
 +\langle m,Mm\rangle.                               \tag{20.8}
\]
Diagonalize \(A\), complete the square in each Gaussian
coordinate, and multiply the one-dimensional integrals.  Equivalently,
for \(I-2sA\succ0\),
\[
\begin{split}
 &-{1\over2}\|\xi\|^2
 +s\langle\xi,A\xi\rangle+2s\langle b,\xi\rangle\\
 &\quad=
 -{1\over2}
 \left\langle
 \xi-2s(I-2sA)^{-1}b,\,
 (I-2sA)
 [\xi-2s(I-2sA)^{-1}b]
 \right\rangle\\
 &\qquad
 +2s^2\langle b,(I-2sA)^{-1}b\rangle.                \tag{20.9}
\end{split}
\]
The Gaussian change of variables gives the determinant and the
last exponential in (20.7).  If (20.6) fails, a top-eigenvector
one-dimensional integral has a nonnegative quadratic exponent and
diverges; the linear term cannot restore integrability.
\end{proof}

\begin{corollary}[Uniform noncentral bound]
\label{cor:v16v20-bound}
If
\[
 2s\|A\|\leq r<1,                                   \tag{20.10}
\]
then
\[
\begin{split}
 {\mathbb E}e^{s\langle X,MX\rangle}
 \leq
 \exp\left\{
 {s\,\operatorname{Tr}A\over1-r}
 +s\langle m,Mm\rangle
 +{2s^2\|b\|^2\over1-r}
 \right\}.                                          \tag{20.11}
\end{split}
\]
\end{corollary}

\begin{proof}
Use
\[
 -\log(1-2s\alpha)\leq{2s\alpha\over1-r}             \tag{20.12}
\]
on every eigenvalue of \(A\), and
\(\|(I-2sA)^{-1}\|\leq(1-r)^{-1}\) in (20.7).
\end{proof}

\subsection{\(L^p\) interacting collar comparison}

Let \(\pi_m=N(m,C)\) be the Gaussian law above and let
\[
 d\nu=h\,d\pi_m,\qquad
 h\geq0,\qquad
 \int h\,d\pi_m=1.                                  \tag{20.13}
\]
For \(1<p<\infty\), put \(q=p/(p-1)\).

\begin{theorem}[Noncentral \(L^p\) collar estimate]
\label{thm:v16v20-Lp}
Assume
\[
 \|h\|_{L^p(\pi_m)}\leq H,\qquad
 2qs\|A\|\leq r<1.                                  \tag{20.14}
\]
Then
\[
\begin{split}
 {\mathbb E}_\nu e^{s\langle X,MX\rangle}
 \leq H\exp\Bigg\{&
 {s\,\operatorname{Tr}A\over1-r}
 +s\langle m,Mm\rangle\\
 &+{2qs^2\|b\|^2\over1-r}\Bigg\}.                   \tag{20.15}
\end{split}
\]
\end{theorem}

\begin{proof}
H\"older gives
\[
 {\mathbb E}_{\pi_m}[h e^{s\langle X,MX\rangle}]
 \leq
 H\left(
 {\mathbb E}_{\pi_m}
 e^{qs\langle X,MX\rangle}\right)^{1/q}.             \tag{20.16}
\]
Apply Corollary~\ref{cor:v16v20-bound} at exponent \(qs\).
After taking the \(q\)-th root, the trace and mean-energy
coefficients become \(s\), while the linear-shift coefficient
becomes \(2qs^2\), proving (20.15).
\end{proof}

\subsection{Finite-collar conditional card}

Let \(K_R\) be the finite set of interior and collar variables and
let \(f\) denote the farther exterior.  Suppose the conditional
Gaussian reference on \(K_R\) is
\[
 \pi_{R,f}=N(m_R(f),C_R),                            \tag{20.17}
\]
with covariance independent of \(f\), and suppose the actual
conditional law has density \(h_{R,f}\) relative to \(\pi_{R,f}\).
Let \(M_R\geq0\) be the finite-collar Newton quadratic operator and
put
\[
\begin{split}
 A_R&=C_R^{1/2}M_RC_R^{1/2},\\
 b_R(f)&=C_R^{1/2}M_Rm_R(f).                         \tag{20.18}
\end{split}
\]

\begin{definition}[Finite-collar interaction card]
\label{def:v16v20-FCIC}
At exponent \(s\) and \(p>1\), FCIC consists of constants
\(H,r,L,U<\infty\), \(r<1\), such that almost surely in \(f\),
\[
\begin{gathered}
 \|h_{R,f}\|_{L^p(\pi_{R,f})}\leq H,\qquad
 2qs\|A_R\|\leq r,\qquad
 \operatorname{Tr}A_R\leq L,\\
 \langle m_R(f),M_Rm_R(f)\rangle
 +{2qs\over1-r}\|b_R(f)\|^2\leq U.                  \tag{20.19}
\end{gathered}
\]
\end{definition}

\begin{corollary}[FCIC gives a uniform conditional collar moment]
\label{cor:v16v20-FCIC}
Under FCIC,
\[
 \mathop{\rm ess\,sup}_{f}
 {\mathbb E}_{\nu(\,\cdot\mid f)}
 e^{s\langle X,M_RX\rangle}
 \leq
 H\exp\left\{{sL\over1-r}+sU\right\}.                \tag{20.20}
\]
\end{corollary}

\begin{proof}
Insert (20.19) into Theorem~\ref{thm:v16v20-Lp}.
\end{proof}

The last line of (20.19) is the conditional-mean loader.  A
uniform centered covariance and a uniform density norm do not
control it: an exterior configuration can shift the collar mean
arbitrarily far while leaving \(C_R\) unchanged.  This is the
finite-collar analogue of the worst-exterior obstruction already
encountered in the archived Higgs conditional analysis.

\begin{assessment}[Physical status]
\label{ass:v16v20-status}
V20 completes the compulsory audit, proves the exact noncentral
Gaussian quadratic determinant, and derives the \(L^p\) finite
collar card.  It isolates three costs: covariance susceptibility,
interaction density, and the exterior-induced conditional mean.

FCIC has not been proved for the complete SM--Einstein conditional
law.  In particular, no worst-exterior bound on \(m_R(f)\) is
available for unbounded matter and metric fields.
\end{assessment}

\begin{decision}[V21 near--far local-state compiler]
\label{dec:v16v20-next}
Continue after the audit.  Combine FCIC for the finite collar with
the projected-density far-tail theorem V19.  Prove a two-stage
conditional H\"older estimate that yields the V16 local Newton
moment without demanding a global exterior \(L^p\) bound.  State
the exact uniform cards still needed from the physical
Standard Model interactions.
\end{decision}

\begin{firewall}[Claim boundary after V20]
\label{fire:v16v20}
V20 does not prove FCIC, a uniform conditional-mean bound, the
physical gap or local monopole ledger, reflection-positive
thermodynamic convergence, or the full continuum.
\end{firewall}

\section*{Version V20 append-only record}
\addcontentsline{toc}{section}{Version V20 append-only record}
The complete V19 source was retained before this continuation.  It
had \(186835\) bytes, \(5616\) lines, and SHA--256
\[
 \texttt{9C223F4C0827DECA7BE6CC4BD93DFDBF66C233E3D9AA4F3EBCB36B2D5B758D8A}.
\]
V20 inspected YM V46 and NS V26, then proved the exact noncentral
Gaussian determinant and the finite-collar conditional
\(L^p\) interaction card.
\section{Sixteenth-cycle V21: a near--far compiler for uniform local Newton moments}
\label{sec:v16v21}

V20 supplies a conditional finite-collar moment, while V19 supplies
the far projected-density limit.  V21 combines them at the level of
the positive Schur influence.  The argument uses no independence
between collar and far variables and no global exterior
\(L^p\)-density bound.

\subsection{Two-stage influence estimate}

Let \(\mu_\alpha\) be the actual tilted exterior marginal at cutoff
and volume index \(\alpha\).  On its exterior carrier, decompose
\[
 B_\alpha=B_{R,\alpha}+D_{R,\alpha}.                 \tag{21.1}
\]
Let \(A_\alpha\) be the positive interior Newton block and set
\[
\begin{split}
 J_\alpha&={\cal J}_{A_\alpha}(B_\alpha),\\
 J_{R,\alpha}^{\rm near}
 &={\cal J}_{A_\alpha}(B_{R,\alpha}),\\
 J_{R,\alpha}^{\rm far}
 &={\cal J}_{A_\alpha}(D_{R,\alpha}),                \tag{21.2}
\end{split}
\]
where \({\cal J}_A(L)=L^*A^+L\).

\begin{theorem}[Near--far exponential compiler]
\label{thm:v16v21-near-far}
Fix \(c>0\), \(\theta>0\), and \(1<\ell<\infty\), with
\(\ell'=\ell/(\ell-1)\).  Then
\[
\begin{split}
 {\mathbb E}_{\mu_\alpha}e^{cJ_\alpha}
 \leq{}&
 \left(
 {\mathbb E}_{\mu_\alpha}
 e^{\,\ell c(1+\theta)J_{R,\alpha}^{\rm near}}
 \right)^{1/\ell}\\
 &\times
 \left(
 {\mathbb E}_{\mu_\alpha}
 e^{\,\ell'c(1+\theta^{-1})J_{R,\alpha}^{\rm far}}
 \right)^{1/\ell'}.
                                                               \tag{21.3}
\end{split}
\]
Consequently, if for one \(R\)
\[
\begin{split}
 N_R&:=
 \sup_\alpha{\mathbb E}_{\mu_\alpha}
 e^{\,\ell c(1+\theta)J_{R,\alpha}^{\rm near}}
 <\infty,\\
 F_R&:=
 \sup_\alpha{\mathbb E}_{\mu_\alpha}
 e^{\,\ell'c(1+\theta^{-1})J_{R,\alpha}^{\rm far}}
 <\infty,                                            \tag{21.4}
\end{split}
\]
then
\[
 \sup_\alpha{\mathbb E}_{\mu_\alpha}e^{cJ_\alpha}
 \leq N_R^{1/\ell}F_R^{1/\ell'}<\infty.              \tag{21.5}
\]
\end{theorem}

\begin{proof}
The quadratic-form inequality (17.4) gives
\[
 J_\alpha
 \leq
 (1+\theta)J_{R,\alpha}^{\rm near}
 +(1+\theta^{-1})J_{R,\alpha}^{\rm far}.             \tag{21.6}
\]
Exponentiate and apply H\"older with conjugate exponents
\(\ell,\ell'\).  This proves (21.3); taking suprema gives (21.5).
\end{proof}

\begin{corollary}[Far-tail input from the projected density]
\label{cor:v16v21-far}
Suppose the hypotheses of
Theorem~\ref{thm:v16v19-tail} hold for the far observation
\(D_{R,\alpha}V_\alpha\).  Then, at the exponent
\[
 c_{\rm far}:=
 \ell'c(1+\theta^{-1}),                              \tag{21.7}
\]
\[
 F_R\longrightarrow1\quad(R\to\infty).              \tag{21.8}
\]
\end{corollary}

\begin{proof}
Apply Theorem~\ref{thm:v16v19-tail} with \(c=c_{\rm far}\).
\end{proof}

\subsection{Loading the near factor from FCIC}

Let \(X_{R,\alpha}\) denote the finite interior--collar variables
conditioned on the farther exterior, and suppose there is a
positive collar quadratic operator \(M_{R,\alpha}\) and a constant
\(d_R<\infty\) such that
\[
 J_{R,\alpha}^{\rm near}
 \leq
 d_R\langle X_{R,\alpha},
             M_{R,\alpha}X_{R,\alpha}\rangle
 \quad\mu_\alpha\hbox{-a.s.}                         \tag{21.9}
\]
The inequality permits \(M_R\) to include all finite collar
interaction channels, rather than requiring equality with the
Newton influence.

\begin{proposition}[FCIC supplies the near moment]
\label{prop:v16v21-near}
Let
\[
 c_{\rm near}:=\ell c(1+\theta)d_R.                 \tag{21.10}
\]
If FCIC of Definition~\ref{def:v16v20-FCIC} holds uniformly in
\(\alpha\) and in the farther exterior at exponent
\(c_{\rm near}\), with constants \(H,r,L,U\), then
\[
 N_R\leq
 H\exp\left\{
 {c_{\rm near}L\over1-r}+c_{\rm near}U\right\}.
                                                               \tag{21.11}
\]
\end{proposition}

\begin{proof}
By (21.9),
\[
 e^{\ell c(1+\theta)J_{R,\alpha}^{\rm near}}
 \leq
 e^{c_{\rm near}
 \langle X_{R,\alpha},M_{R,\alpha}X_{R,\alpha}\rangle}.
                                                               \tag{21.12}
\]
Condition on the farther exterior and apply
Corollary~\ref{cor:v16v20-FCIC}.  Its upper bound is independent of
the conditioned configuration.  Integrating and taking the
supremum over \(\alpha\) proves (21.11).
\end{proof}

\begin{theorem}[Complete local Newton moment compiler]
\label{thm:v16v21-local}
Fix \(s,t>0\) and \(0<\eta\leq1\), and put
\[
 c={2s\over\eta}.                                    \tag{21.13}
\]
Assume, uniformly in the cutoff and volume index \(\alpha\):
\begin{enumerate}
\item the conditional interior reference moment satisfies
\[
 M_{K,\alpha}(t+s(1+\eta))\leq M_*;                 \tag{21.14}
\]
\item the finite-collar domination (21.9) and FCIC hold at the
exponent \(c_{\rm near}\) in (21.10);
\item the projected far-density and covariance tail hypotheses of
V19 hold at \(c_{\rm far}\) in (21.7);
\item the positive interior block has the fixed gap needed to
define and control all three influence forms in (21.2).
\end{enumerate}
Then, for all sufficiently large fixed collars \(R\),
\[
 \sup_\alpha
 {\mathbb E}_{\widehat\rho_{s,\alpha}}
 e^{tQ_{K,\alpha}}
 <\infty.                                            \tag{21.15}
\]
More explicitly,
\[
\begin{split}
 \sup_\alpha
 {\mathbb E}_{\widehat\rho_{s,\alpha}}
 e^{tQ_{K,\alpha}}
 \leq{}&
 M_*
 \left[
 H\exp\left\{
 {c_{\rm near}L\over1-r}+c_{\rm near}U
 \right\}\right]^{1/\ell}
 F_R^{1/\ell'}.                                      \tag{21.16}
\end{split}
\]
\end{theorem}

\begin{proof}
Proposition~\ref{prop:v16v21-near} and
Corollary~\ref{cor:v16v21-far} verify (21.4).
Theorem~\ref{thm:v16v21-near-far} therefore gives
\[
 \sup_\alpha
 {\mathbb E}_{\mu_\alpha}
 e^{(2s/\eta)J_\alpha}
 \leq
 \left[
 H e^{c_{\rm near}L/(1-r)+c_{\rm near}U}
 \right]^{1/\ell}F_R^{1/\ell'}.                     \tag{21.17}
\]
This is the exterior influence loader (16.6).  Combine it with the
interior moment (21.14) in
Theorem~\ref{thm:v16v16-loader} to obtain (21.16).
\end{proof}

\begin{corollary}[Local tightness of the Newton energy]
\label{cor:v16v21-tight}
Under Theorem~\ref{thm:v16v21-local}, the local energies
\(\{Q_{K,\alpha}\}\) are tight and uniformly integrable under the
normalized tilted states.  Every weakly convergent local-state
subsequence preserves expectations of continuous functions
\(F(Q_K)\) satisfying
\[
 |F(q)|\leq C e^{t' q}\quad\hbox{for some }0\leq t'<t.
                                                               \tag{21.18}
\]
\end{corollary}

\begin{proof}
Exponential Markov gives tightness.  Choose
\(\varepsilon>0\) with
\((1+\varepsilon)t'<t\).  The bound (21.15) makes
\(\{|F(Q_{K,\alpha})|^{1+\varepsilon}\}\) uniformly integrable.
Weak convergence plus truncation, or Vitali's theorem after a
Skorokhod representation, gives convergence of the expectations.
\end{proof}

\subsection{The remaining physical cards}

Theorem~\ref{thm:v16v21-local} is a genuine near--far reduction,
but it is conditional.  The coupled theory must still provide:
\begin{enumerate}
\item a fixed interior constraint gap or a retained soft-mode
stratum;
\item a local monopole ledger permitting the V17 projected tail;
\item a uniform conditional interior moment (21.14);
\item FCIC, including its exterior-induced mean bound \(U\);
\item a uniform \(L^p\) density after projection onto the far
influence variable; and
\item covariance-weighted far-tail decay.
\end{enumerate}
None can be inferred merely from positivity of the scalar Euclidean
weight.

\begin{assessment}[Physical status]
\label{ass:v16v21-status}
V21 proves the requested two-stage compiler.  It converts a finite
collar FCIC estimate and a projected far-tail estimate into a
uniform exponential moment for the local Newton energy under the
normalized interacting state.  The proof never assumes a global
exterior density norm or independence between near and far
variables.

The physical loaders listed above remain open.  In particular, the
finite-collar worst-exterior conditional mean is not yet controlled.
\end{assessment}

\begin{decision}[V22 averaged conditional-mean replacement]
\label{dec:v16v21-next}
The essential-supremum mean bound in FCIC is likely too strong for
unbounded fields.  Replace it by an exponential moment of the
conditional-mean cost under the farther-exterior marginal.  Derive
the correct second H\"older split and show that an averaged
conditional card still supplies the near factor \(N_R\).
\end{decision}

\begin{firewall}[Claim boundary after V21]
\label{fire:v16v21}
V21 proves an implication between explicit cards.  It does not
populate those cards for the physical Standard Model--Einstein
measure, construct its local states, prove reflection positivity or
Einstein response in the limit, or complete the continuum.
\end{firewall}

\section*{Version V21 append-only record}
\addcontentsline{toc}{section}{Version V21 append-only record}
The complete V20 source was retained before this continuation.  It
had \(195638\) bytes, \(5909\) lines, and SHA--256
\[
 \texttt{91C38A1E8EB51E90D360CE7133B04AD9AA1461D0B8196575AFF6DBCC116F91B3}.
\]
V21 combined the finite-collar noncentral estimate and projected
far-tail theorem into a uniform local Newton moment compiler without
a global exterior density hypothesis.
\section{Sixteenth-cycle V22: averaged conditional means replace the worst-exterior collar bound}
\label{sec:v16v22}

FCIC uses an essential supremum over the farther exterior of the
Gaussian conditional-mean cost.  For unbounded fields that target
is generally unrealistic.  V22 retains the exact conditional
estimate of V20 but integrates its two random costs with a second
H\"older inequality.

\subsection{Random conditional comparison data}

Let \((F,\zeta)\) carry the farther-exterior variable \(f\).  Given
\(f\), let
\[
 \pi_f=N(m(f),C)                                     \tag{22.1}
\]
be a finite-dimensional Gaussian collar reference, and let
\[
 d\nu_f=h_f\,d\pi_f,\qquad
 H_p(f):=\|h_f\|_{L^p(\pi_f)},\qquad
 q={p\over p-1}.                                     \tag{22.2}
\]
Let \(M\geq0\), and put
\[
 A=C^{1/2}MC^{1/2},\qquad
 b(f)=C^{1/2}Mm(f).                                  \tag{22.3}
\]
At exponent \(\lambda>0\), assume
\[
 2q\lambda\|A\|\leq r<1,\qquad
 \operatorname{Tr}A\leq L.                          \tag{22.4}
\]
Define the random mean cost
\[
 U_\lambda(f)
 :=
 \langle m(f),Mm(f)\rangle
 +{2q\lambda\over1-r}\|b(f)\|^2.                    \tag{22.5}
\]

\begin{theorem}[Averaged noncentral collar estimate]
\label{thm:v16v22-averaged}
Let \(1<\rho<\infty\) and
\(\rho'=\rho/(\rho-1)\).  If
\[
 \|H_p\|_{L^\rho(\zeta)}<\infty,\qquad
 {\mathbb E}_\zeta
 e^{\rho'\lambda U_\lambda}<\infty,                 \tag{22.6}
\]
then the mixture
\[
 \nu(dX,df)=\nu_f(dX)\zeta(df)                       \tag{22.7}
\]
satisfies
\[
\begin{split}
 {\mathbb E}_\nu e^{\lambda\langle X,MX\rangle}
 \leq{}&
 \exp\left\{{\lambda L\over1-r}\right\}
 \|H_p\|_{L^\rho(\zeta)}
 \left(
 {\mathbb E}_\zeta e^{\rho'\lambda U_\lambda}
 \right)^{1/\rho'}.
                                                               \tag{22.8}
\end{split}
\]
\end{theorem}

\begin{proof}
The proof of Theorem~\ref{thm:v16v20-Lp}, applied at each fixed
\(f\), gives
\[
 {\mathbb E}_{\nu_f}e^{\lambda\langle X,MX\rangle}
 \leq
 H_p(f)
 \exp\left\{
 {\lambda L\over1-r}+\lambda U_\lambda(f)
 \right\}.                                          \tag{22.9}
\]
Integrate over \(\zeta\) and apply H\"older with exponents
\(\rho,\rho'\) to
\({\mathbb E}_\zeta[H_p e^{\lambda U_\lambda}]\).
This is (22.8).
\end{proof}

\begin{definition}[Averaged finite-collar interaction card]
\label{def:v16v22-AFCIC}
AFCIC at exponents \((\lambda,p,\rho)\) consists of the uniform
spectral and trace bounds (22.4) and constants
\[
 {\mathbb E}_\zeta H_p(f)^\rho\leq{\cal H}^\rho,
\qquad
 {\mathbb E}_\zeta
 e^{\rho'\lambda U_\lambda(f)}\leq{\cal U}.          \tag{22.10}
\]
\end{definition}

\begin{corollary}[AFCIC supplies the near factor]
\label{cor:v16v22-near}
In Proposition~\ref{prop:v16v21-near}, replace FCIC by AFCIC at
\[
 \lambda=c_{\rm near}
 =\ell c(1+\theta)d_R.                               \tag{22.11}
\]
Then
\[
 N_R
 \leq
 \exp\left\{{c_{\rm near}L\over1-r}\right\}
 {\cal H}\,{\cal U}^{1/\rho'}.                       \tag{22.12}
\]
\end{corollary}

\begin{proof}
Use the domination (21.9), then apply
Theorem~\ref{thm:v16v22-averaged} with
\(\lambda=c_{\rm near}\).
\end{proof}

\begin{theorem}[Local Newton compiler with averaged collar data]
\label{thm:v16v22-local}
The conclusion (21.15) of
Theorem~\ref{thm:v16v21-local} remains valid when its FCIC
hypothesis is replaced by AFCIC at (22.11).  The explicit bound
becomes
\[
\begin{split}
 \sup_\alpha
 {\mathbb E}_{\widehat\rho_{s,\alpha}}
 e^{tQ_{K,\alpha}}
 \leq{}&
 M_*
 \left[
 e^{c_{\rm near}L/(1-r)}
 {\cal H}{\cal U}^{1/\rho'}
 \right]^{1/\ell}
 F_R^{1/\ell'}.
                                                               \tag{22.13}
\end{split}
\]
\end{theorem}

\begin{proof}
Corollary~\ref{cor:v16v22-near} replaces
Proposition~\ref{prop:v16v21-near} in the proof of
Theorem~\ref{thm:v16v21-local}.  The near--far H\"older step and
the V16 exterior integration are unchanged.
\end{proof}

\subsection{The averaged card is strictly weaker}

If FCIC holds with \(H_p(f)\leq H\) and
\(U_\lambda(f)\leq U\) almost surely, then AFCIC holds for every
\(\rho>1\), with
\[
 {\cal H}=H,\qquad
 {\cal U}=e^{\rho'\lambda U}.                        \tag{22.14}
\]
The converse fails: an unbounded conditional mean may have every
exponential moment below a threshold and hence satisfy AFCIC while
violating every essential-supremum bound.

\begin{counterexample}[A first mean does not pay AFCIC]
\label{cnt:v16v22-first-moment}
Let \(U\) take positive integer values with
\[
 {\mathbb P}(U=n)={n^{-3}\over\sum_{k\geq1}k^{-3}}.  \tag{22.15}
\]
Then \({\mathbb E}U<\infty\), but
\[
 {\mathbb E}e^{aU}=\infty
 \quad\hbox{for every }a>0.                          \tag{22.16}
\]
\end{counterexample}

\begin{proof}
The first moment is a constant multiple of
\(\sum n^{-2}<\infty\).  The exponential series
\(\sum e^{an}n^{-3}\) diverges for every \(a>0\).
\end{proof}

Thus an energy expectation or ordinary entropy estimate cannot
replace the exponential conditional-mean row in (22.10).

\subsection{Where the conditional mean comes from}

For a joint Gaussian reference on collar \(X\) and far exterior
\(Y\), with covariance blocks
\[
 \begin{pmatrix}
 C_{XX}&C_{XY}\\
 C_{YX}&C_{YY}
 \end{pmatrix},                                      \tag{22.17}
\]
the conditional mean is, on the support of \(C_{YY}\),
\[
 m(Y)=C_{XY}C_{YY}^+Y.                               \tag{22.18}
\]
Consequently \(U_\lambda(m(Y))\) is itself a positive quadratic
form of the far Gaussian variable.  This observation suggests a
second determinant estimate, but only after controlling the
pseudoinverse \(C_{YY}^+\), the induced quadratic operator, and the
interacting projected density.  V22 does not assume those controls
silently.

\begin{assessment}[Physical status]
\label{ass:v16v22-status}
V22 replaces the worst-exterior conditional-mean bound by AFCIC,
an averaged \(L^\rho\) density norm and an exponential moment of
the exact noncentral mean cost.  This strictly weakens the collar
loader while preserving the V21 local Newton conclusion.

The new load-bearing quantity is the quadratic form generated by
the regression map \(C_{XY}C_{YY}^+\).  Its uniform determinant
margin and interacting projected-density bound remain unproved.
\end{assessment}

\begin{decision}[V23 Gaussian regression influence]
\label{dec:v16v22-next}
Compute the operator representing \(U_\lambda(m(Y))\) from
(22.18).  Apply the V4 determinant criterion to it, derive a
fixed-gap stability theorem for the regression pseudoinverse, and
state the projected interacting density required to populate the
AFCIC mean row.
\end{decision}

\begin{firewall}[Claim boundary after V22]
\label{fire:v16v22}
V22 proves an averaged conditional compiler.  It does not prove
AFCIC for the physical law, control the Gaussian regression
pseudoinverse, construct the local interacting states or their
reflection positivity, or complete the continuum.
\end{firewall}

\section*{Version V22 append-only record}
\addcontentsline{toc}{section}{Version V22 append-only record}
The complete V21 source was retained before this continuation.  It
had \(204899\) bytes, \(6196\) lines, and SHA--256
\[
 \texttt{871DC08357268E24823789D8E8F7D99519222FC6AF39AC6F849096BEF52CD656}.
\]
V22 replaced the finite-collar worst-exterior bound by an averaged
conditional-density and conditional-mean card and preserved the
near--far local Newton compiler.
\section{Sixteenth-cycle V23: Gaussian regression determinant for the averaged collar mean}
\label{sec:v16v23}

V22 reduces the collar problem to an exponential moment of the
conditional-mean cost \(U_\lambda\).  For a joint Gaussian
reference, that mean is a linear regression of the far field.
This section computes the resulting quadratic operator, gives its
exact determinant threshold, and proves its stability under a fixed
far-covariance gap.

\subsection{Conditional Gaussian regression}

Let \(E_X,E_Y\) be finite-dimensional real Hilbert spaces and let
\((X,Y)\) be centered Gaussian with covariance
\[
 \Sigma=
 \begin{pmatrix}
 C_{XX}&K\\ K^*&S
 \end{pmatrix}\succeq0.                              \tag{23.1}
\]
The conditional law of \(X\) given \(Y\) has
\[
 m(Y)=RY,\qquad R:=KS^+,                             \tag{23.2}
\]
and covariance
\[
 C:=C_{XX}-KS^+K^*\succeq0.                          \tag{23.3}
\]
These formulas are understood on \(\operatorname{Ran}S\), which
contains \(Y\) almost surely.

\begin{lemma}[Regression range and conditional covariance]
\label{lem:v16v23-regression}
Positivity of (23.1) implies
\[
 \ker S\subset\ker K,\qquad
 K=KS^+S,                                            \tag{23.4}
\]
and (23.3) is nonnegative.
\end{lemma}

\begin{proof}
For \(y\in\ker S\), positivity of
\(\langle(x,ty),\Sigma(x,ty)\rangle\) for every real \(t\) forces
\(\langle x,Ky\rangle=0\) for all \(x\), proving the first claim
and hence (23.4).  Completing the square in the block covariance,
or minimizing its quadratic form over the \(Y\)-support, gives the
nonnegative Schur complement (23.3).
\end{proof}

\subsection{The mean-cost operator}

Let \(M\geq0\) be the collar Newton quadratic operator.  Fix the
conditional interaction exponent \(\lambda>0\), an \(L^p\) density
exponent \(p>1\), and put
\[
 q={p\over p-1}.                                     \tag{23.5}
\]
Assume
\[
 2q\lambda\|A\|\leq r<1,\qquad
 A=C^{1/2}MC^{1/2}.                                  \tag{23.6}
\]
The V22 mean cost is
\[
 U_\lambda(Y)
 =
 \langle m(Y),Mm(Y)\rangle
 +{2q\lambda\over1-r}
 \|C^{1/2}Mm(Y)\|^2.                                \tag{23.7}
\]
Define
\[
\begin{split}
 W_\lambda
 &:=
 M+{2q\lambda\over1-r}MCM,\\
 L_\lambda&:=R^*W_\lambda R,\\
 {\cal A}_\lambda
 &:=S^{1/2}L_\lambda S^{1/2}.                        \tag{23.8}
\end{split}
\]

\begin{proposition}[Exact regression representation]
\label{prop:v16v23-cost}
One has
\[
 U_\lambda(Y)=\langle Y,L_\lambda Y\rangle,          \tag{23.9}
\]
and \({\cal A}_\lambda\geq0\).
\end{proposition}

\begin{proof}
Insert \(m(Y)=RY\) into (23.7).  Since
\[
 \|C^{1/2}MRY\|^2
 =\langle RY,MCMRY\rangle,                           \tag{23.10}
\]
the result is exactly
\(\langle RY,W_\lambda RY\rangle\), which equals
(23.9).  Positivity of \(W_\lambda\) gives positivity of
\(L_\lambda\) and \({\cal A}_\lambda\).
\end{proof}

\begin{theorem}[Gaussian AFCIC mean determinant]
\label{thm:v16v23-determinant}
Let \(\rho>1\), \(\rho'=\rho/(\rho-1)\), and suppose
\[
 Y\sim N(0,S).
 \]
Then
\[
 {\mathbb E}
 e^{\rho'\lambda U_\lambda(Y)}<\infty                \tag{23.11}
\]
if and only if
\[
 2\rho'\lambda\|{\cal A}_\lambda\|_{\rm op}<1.       \tag{23.12}
\]
Under (23.12),
\[
 {\mathbb E}
 e^{\rho'\lambda U_\lambda(Y)}
 =
 \det(I-2\rho'\lambda{\cal A}_\lambda)^{-1/2}.
                                                               \tag{23.13}
\]
If the left side of (23.12) is at most \(r_Y<1\), then
\[
 {\mathbb E}
 e^{\rho'\lambda U_\lambda(Y)}
 \leq
 \exp\left\{
 {\rho'\lambda\operatorname{Tr}{\cal A}_\lambda
  \over1-r_Y}\right\}.                               \tag{23.14}
\]
\end{theorem}

\begin{proof}
Apply Theorem~\ref{thm:v16v4-gaussian} to the Gaussian variable
\(Y\), covariance \(S\), quadratic operator \(L_\lambda\), and
\(\Gamma=2\rho'\lambda\).  Proposition
\ref{prop:v16v23-cost} gives (23.11)--(23.13).
The eigenvalue estimate used in (12.9) gives (23.14).
\end{proof}

\subsection{Projected interacting regression law}

Let the actual far marginal \(\nu_Y\) have density relative to
\(\pi_Y=N(0,S)\).  The cost depends only on
\[
 Z_\lambda:=W_\lambda^{1/2}RY.                       \tag{23.15}
\]
Let \(\bar h_\lambda\) be the density of the
\(Z_\lambda\)-law under \(\nu_Y\) relative to its law under
\(\pi_Y\).

\begin{theorem}[Projected regression-density certificate]
\label{thm:v16v23-projected}
Let \(1<a<\infty\), \(a'=a/(a-1)\), and suppose
\[
 \|\bar h_\lambda\|_{L^a}\leq H_Y.                  \tag{23.16}
\]
If
\[
 2a'\rho'\lambda\|{\cal A}_\lambda\|<1,             \tag{23.17}
\]
then
\[
\begin{split}
 {\mathbb E}_{\nu_Y}
 e^{\rho'\lambda U_\lambda(Y)}
 \leq
 H_Y
 \det(I-2a'\rho'\lambda{\cal A}_\lambda)^{-1/(2a')}.
                                                               \tag{23.18}
\end{split}
\]
\end{theorem}

\begin{proof}
By (23.9) and (23.15),
\[
 U_\lambda(Y)=\|Z_\lambda\|^2.                       \tag{23.19}
\]
Push both measures to the \(Z_\lambda\)-space and apply
Theorem~\ref{thm:v16v13-dual} with density
\(\bar h_\lambda\), H\"older conjugate \(a'\), and exponent
\(\rho'\lambda\).  The Gaussian determinant carrier has the same
nonzero spectrum as \({\cal A}_\lambda\), by the
\(L^*L/LL^*\) argument of Proposition
\ref{prop:v16v19-carriers}.  This proves (23.18).
\end{proof}

\subsection{Fixed-gap cutoff stability}

Let the data in (23.1)--(23.8) carry an index \(n\), on fixed
finite-dimensional collar and far carriers.

\begin{theorem}[Regression-operator stability]
\label{thm:v16v23-stability}
Assume
\[
 \sigma(S_n),\sigma(S)
 \subset\{0\}\cup[\gamma,\infty),\qquad\gamma>0,     \tag{23.20}
\]
and
\[
 C_{XX,n}\to C_{XX},\quad K_n\to K,\quad
 S_n\to S,\quad M_n\to M                             \tag{23.21}
\]
in operator norm.  Suppose the conditional susceptibility bound
(23.6) holds with one \(r<1\).  Then
\[
\begin{gathered}
 S_n^+\to S^+,\qquad R_n=K_nS_n^+\to R,\\
 C_n\to C,\qquad W_{\lambda,n}\to W_\lambda,\qquad
 {\cal A}_{\lambda,n}\to{\cal A}_\lambda            \tag{23.22}
\end{gathered}
\]
in operator norm and, in the fixed finite dimension, in trace norm.
If additionally
\[
 \sup_n2\rho'\lambda
 \|{\cal A}_{\lambda,n}\|\leq r_Y<1,                \tag{23.23}
\]
then the Gaussian determinants (23.13) converge and are uniformly
bounded.
\end{theorem}

\begin{proof}
The common gap and resolvent functional calculus give
\(S_n^+\to S^+\), exactly as in V8 and V16.  Multiplication of
convergent bounded operators gives \(R_n\to R\).  Equations
(23.3) and (23.8), together with continuity of the positive square
root in finite dimension, give all remaining convergences in
(23.22).  Operator- and trace-norm convergence are equivalent on a
fixed finite-dimensional carrier.  The determinant convergence and
uniform bound follow from Theorem~\ref{thm:v16v4-uniform}.
\end{proof}

\begin{counterexample}[Regression gap collapse]
\label{cnt:v16v23-gap}
In one collar and one far dimension, take
\[
 C_{XX,n}=1,\qquad K_n=n^{-1},\qquad S_n=n^{-2}.     \tag{23.24}
\]
Then the joint covariance
\[
 \Sigma_n=
 \begin{pmatrix}1&n^{-1}\\n^{-1}&n^{-2}\end{pmatrix}
 \succeq0
 \longrightarrow
 \begin{pmatrix}1&0\\0&0\end{pmatrix},              \tag{23.25}
\]
but
\[
 R_n=n,\qquad
 S_n^{1/2}R_n^*R_nS_n^{1/2}=1.                      \tag{23.26}
\]
The limiting regression is zero, so its corresponding carrier is
zero.  The regression cost is therefore discontinuous at the
collapsed far-covariance gap.
\end{counterexample}

\begin{proof}
The matrices in (23.25) are rank-one outer products
\((1,n^{-1})^{\mathsf T}(1,n^{-1})\).  Direct calculation gives
(23.26).
\end{proof}

\begin{assessment}[Physical status]
\label{ass:v16v23-status}
V23 converts the AFCIC conditional-mean row into the exact
regression carrier \({\cal A}_\lambda\), proves its Gaussian and
projected interacting determinant criteria, and establishes
fixed-gap cutoff stability.  It also proves that covariance
convergence alone fails at a closing regression gap.

The physical theory still needs a uniform projected density bound
for \(Z_\lambda\), a far-covariance gap or retained soft modes, and
trace-class convergence when the collar carrier itself grows with
the ultraviolet cutoff.
\end{assessment}

\begin{decision}[V24 growing-collar trace-ideal stability]
\label{dec:v16v23-next}
Replace the fixed-dimensional step in V23 by a common separable
carrier.  Give sufficient Hilbert--Schmidt and operator convergence
hypotheses on \(S_n^{1/2}\), \(R_n\), and \(W_{\lambda,n}\) that
force \({\cal A}_{\lambda,n}\to{\cal A}_\lambda\) in trace norm.
Then the determinant row will survive a growing ultraviolet collar.
\end{decision}

\begin{firewall}[Claim boundary after V23]
\label{fire:v16v23}
V23 does not prove the projected interaction-density bound, a
uniform regression gap, growing-carrier trace convergence, the
physical local-state construction or reflection positivity, or the
full continuum.
\end{firewall}

\section*{Version V23 append-only record}
\addcontentsline{toc}{section}{Version V23 append-only record}
The complete V22 source was retained before this continuation.  It
had \(212324\) bytes, \(6440\) lines, and SHA--256
\[
 \texttt{6CB38C5BF5D03D9826E3E9EEB66C510C6798197C98C093BFF369F726A68DCCF1}.
\]
V23 computed the Gaussian regression mean-cost operator, its exact
determinant threshold, its projected interacting comparison, and
its fixed-gap cutoff stability.
\section{Sixteenth-cycle V24: trace-ideal stability on a growing regression collar}
\label{sec:v16v24}

V23 proves regression stability in a fixed finite dimension.  An
ultraviolet collar has a growing number of modes, so operator-norm
convergence is not enough for its Fredholm determinant.  V24
factors the regression carrier through one Hilbert--Schmidt map and
derives trace-norm convergence on a common separable carrier.

\subsection{The Hilbert--Schmidt regression factor}

After transporting all regulators to common Hilbert spaces
\({\cal H}_X,{\cal H}_Y\), let
\[
 S_n,S\geq0,\qquad
 R_n,R:{\cal H}_Y\to{\cal H}_X,\qquad
 W_n,W\geq0.                                        \tag{24.1}
\]
Define
\[
\begin{split}
 F_n&:=W_n^{1/2}R_nS_n^{1/2},\\
 F&:=W^{1/2}RS^{1/2},\\
 {\cal A}_n&:=F_n^*F_n,\qquad
 {\cal A}:=F^*F.                                    \tag{24.2}
\end{split}
\]
For the V23 application, \(W_n=W_{\lambda,n}\) and
\({\cal A}_n={\cal A}_{\lambda,n}\).

\begin{theorem}[Primitive growing-collar convergence]
\label{thm:v16v24-primitive}
Assume
\[
\begin{gathered}
 \|S_n^{1/2}-S^{1/2}\|_2\to0,\qquad
 \|R_n-R\|_{\rm op}\to0,\\
 \|W_n^{1/2}-W^{1/2}\|_{\rm op}\to0,                 \tag{24.3}
\end{gathered}
\]
and assume \(\|R_n\|\), \(\|W_n^{1/2}\|\), and
\(\|S_n^{1/2}\|_2\) are uniformly bounded.  Then
\[
 \|F_n-F\|_2\longrightarrow0                        \tag{24.4}
\]
and
\[
 \boxed{\ 
 \|{\cal A}_n-{\cal A}\|_1\longrightarrow0.\ }       \tag{24.5}
\]
\end{theorem}

\begin{proof}
Insert two intermediate factors:
\[
\begin{split}
 F_n-F={}&
 (W_n^{1/2}-W^{1/2})R_nS_n^{1/2}\\
 &+W^{1/2}(R_n-R)S_n^{1/2}\\
 &+W^{1/2}R(S_n^{1/2}-S^{1/2}).                     \tag{24.6}
\end{split}
\]
The ideal inequality
\[
 \|ABC\|_2\leq
 \|A\|_{\rm op}\|B\|_{\rm op}\|C\|_2               \tag{24.7}
\]
makes every term tend to zero under (24.3), proving (24.4).

For Hilbert--Schmidt \(F_n,F\),
\[
 \|F_n^*F_n-F^*F\|_1
 \leq(\|F_n\|_2+\|F\|_2)\|F_n-F\|_2.                \tag{24.8}
\]
The first factor is bounded and the second tends to zero, proving
(24.5).
\end{proof}

\begin{corollary}[Trace-class covariance input]
\label{cor:v16v24-covariance}
The first line of (24.3) follows from
\[
 \|S_n-S\|_1\longrightarrow0                        \tag{24.9}
\]
by the Powers--St{\o}rmer inequality.
\end{corollary}

\begin{proof}
For nonnegative trace-class operators,
\[
 \|S_n^{1/2}-S^{1/2}\|_2^2
 \leq\|S_n-S\|_1.                                   \tag{24.10}
\]
\end{proof}

\subsection{Loading the regression map and mean-cost operator}

Use the V23 notation
\[
 R_n=K_nS_n^+,\qquad
 C_n=C_{XX,n}-K_nS_n^+K_n^*,                        \tag{24.11}
\]
and
\[
 W_{\lambda,n}
 =
 M_n+{2q\lambda\over1-r}M_nC_nM_n.                  \tag{24.12}
\]

\begin{theorem}[Growing regression-carrier handoff]
\label{thm:v16v24-handoff}
Suppose:
\begin{enumerate}
\item \(S_n,S\) are nonnegative trace class and
\(\|S_n-S\|_1\to0\);
\item for one \(\gamma>0\),
\[
 \sigma(S_n),\sigma(S)
 \subset\{0\}\cup[\gamma,\infty);                   \tag{24.13}
\]
\item \(K_n\to K\), \(C_{XX,n}\to C_{XX}\), and
\(M_n\to M\) in operator norm, with uniform bounds;
\item the conditional susceptibility reserve in (23.6) holds with
one \(r<1\).
\end{enumerate}
Then
\[
 R_n\to R,\qquad C_n\to C,\qquad
 W_{\lambda,n}^{1/2}\to W_\lambda^{1/2}             \tag{24.14}
\]
in operator norm, and
\[
 {\cal A}_{\lambda,n}
 \longrightarrow{\cal A}_\lambda
 \quad\hbox{in trace norm}.                          \tag{24.15}
\]
\end{theorem}

\begin{proof}
The fixed gap and resolvent functional calculus give
\(S_n^+\to S^+\) in operator norm.  Hence
\[
 R_n=K_nS_n^+\to KS^+=R.                            \tag{24.16}
\]
Substitution in (24.11) gives \(C_n\to C\), and then (24.12)
gives \(W_{\lambda,n}\to W_\lambda\) in operator norm.  The
positive square-root map is operator-norm continuous on bounded
nonnegative operators, proving the last convergence in (24.14).
Corollary~\ref{cor:v16v24-covariance} and
Theorem~\ref{thm:v16v24-primitive} now prove (24.15).
\end{proof}

\subsection{Fredholm moment convergence}

\begin{theorem}[Cofinal regression determinant]
\label{thm:v16v24-determinant}
Under Theorem~\ref{thm:v16v24-handoff}, suppose
\[
 \sup_n2\rho'\lambda
 \|{\cal A}_{\lambda,n}\|_{\rm op}\leq r_Y<1.        \tag{24.17}
\]
Then
\[
\begin{split}
 &\det(I-2\rho'\lambda{\cal A}_{\lambda,n})^{-1/2}\\
 &\qquad\longrightarrow
 \det(I-2\rho'\lambda{\cal A}_{\lambda})^{-1/2},     \tag{24.18}
\end{split}
\]
and
\[
\begin{split}
 \left|
 \log{
 \det(I-2\rho'\lambda{\cal A}_{\lambda,n})^{-1/2}
 \over
 \det(I-2\rho'\lambda{\cal A}_{\lambda})^{-1/2}}
 \right|
 \leq
 {\rho'\lambda\over1-r_Y}
 \|{\cal A}_{\lambda,n}-{\cal A}_{\lambda}\|_1.
                                                               \tag{24.19}
\end{split}
\]
\end{theorem}

\begin{proof}
Apply Theorem~\ref{thm:v16v4-uniform} with
\(\Gamma=2\rho'\lambda\) and use (24.15).  Its logarithmic
Lipschitz estimate becomes (24.19).
\end{proof}

\begin{corollary}[Growing-collar Gaussian AFCIC mean row]
\label{cor:v16v24-AFCIC}
The Gaussian mean moments in (23.13) converge and remain uniformly
bounded under (24.17).  If the projected interacting densities of
Theorem~\ref{thm:v16v23-projected} also have one uniform
\(L^a\) bound and
\[
 \sup_n2a'\rho'\lambda
 \|{\cal A}_{\lambda,n}\|\leq\widetilde r_Y<1,       \tag{24.20}
\]
then their AFCIC mean moments are uniformly bounded as well.
\end{corollary}

\begin{proof}
The first statement is (23.13) and
Theorem~\ref{thm:v16v24-determinant}.  For the second, use
(23.18) and the uniform determinant estimate with exponent
\(a'\rho'\lambda\).
\end{proof}

\subsection{Why operator norm is insufficient}

\begin{counterexample}[Vanishing spectral radius with a surviving
determinant]
\label{cnt:v16v24-operator}
On \(\ell^2(\mathbb N)\), let \(P_n\) project onto the first \(n\)
coordinates and put
\[
 {\cal A}_n={1\over n}P_n.                           \tag{24.21}
\]
Then
\[
 \|{\cal A}_n\|_{\rm op}={1\over n}\to0,\qquad
 \operatorname{Tr}{\cal A}_n=1,                     \tag{24.22}
\]
but for every fixed \(0<z<1\),
\[
 \det(I-z{\cal A}_n)^{-1/2}
 =
 \left(1-{z\over n}\right)^{-n/2}
 \longrightarrow e^{z/2}\ne1.                      \tag{24.23}
\]
\end{counterexample}

\begin{proof}
Equations (24.22) are immediate.  The determinant has \(n\)
identical nontrivial eigenvalue factors, and
\((1-z/n)^{-n/2}\to e^{z/2}\).
\end{proof}

Thus a uniform spectral reserve controls integrability but does not
control normalization across a growing number of modes.
Trace-ideal convergence is an independent cofinal requirement.

\begin{assessment}[Physical status]
\label{ass:v16v24-status}
V24 proves a growing-carrier trace-norm theorem for the regression
mean-cost operator and hence cofinal convergence of its Gaussian
Fredholm moment.  It derives the result from trace-norm convergence
of the far covariance, a fixed covariance gap, and operator
convergence of the primitive blocks.

The physical programme has not proved those primitive convergences
or the projected interacting-density bound.  A closing far
covariance gap still requires a stratified soft-mode treatment.
\end{assessment}

\begin{decision}[V25 companion audit and card ledger]
\label{dec:v16v24-next}
Perform the compulsory companion audit.  Then consolidate
V11--V24 into a minimal local Newton transfer card, separating
proved analytic implications from the unpopulated physical
loaders.  Use the audit to decide which loader should be attacked
next.
\end{decision}

\begin{firewall}[Claim boundary after V24]
\label{fire:v16v24}
V24 does not populate the growing-carrier hypotheses for the
coupled theory, prove the projected density or soft-mode rows,
construct reflection-positive local states, recover the Einstein
equation, or complete the continuum.
\end{firewall}

\section*{Version V24 append-only record}
\addcontentsline{toc}{section}{Version V24 append-only record}
The complete V23 source was retained before this continuation.  It
had \(221735\) bytes, \(6752\) lines, and SHA--256
\[
 \texttt{ACBEE739BAD216D03F3C5508CD0E2236B2E70D884BA60F96EA704F41B109A8C7}.
\]
V24 proved trace-ideal convergence of the growing regression
carrier and cofinal stability of its Gaussian and projected
interacting determinant bounds.
\section{Sixteenth-cycle V25: companion audit and the local Newton transfer card}
\label{sec:v16v25}

This compulsory checkpoint records the new YM V47 result and
consolidates V11--V24 into one local Newton transfer card.  The
audit does not import a Yang--Mills theorem, but its exact
anisotropic chart failure reinforces the need to retain exceptional
soft directions rather than conceal them behind a uniform gap.

\subsection{Compulsory companion audit}

\begin{audit}[Files inspected at V25]
\label{aud:v16v25-files}
The newest companion checkpoints are
\[
\begin{array}{c|r|r}
\hbox{file}&\hbox{lines}&\hbox{bytes}\\ \hline
\texttt{Renewal\_Yang\_Mills\_Mass\_Gap\_Research\_Notes\_V47.tex}
 &11182&408812\\
\texttt{Renewal\_NS\_Stability\_Research\_Notes\_V26.tex}
 &5319&278283.
\end{array}                                             \tag{25.1}
\]
Their SHA--256 values are
\[
\begin{gathered}
\texttt{2FDB54605CBCBC68D9C8BC9D69FD1708257A25425F7BE227DA82AC9BA7020B8F},\\
\texttt{82C887F8C2C69E4F28FAD7C3DC8DE5B9099CB5A4BF431EBA1FFF3E91935978AD}.
\end{gathered}                                         \tag{25.2}
\]
\end{audit}

\begin{proposition}[Reading of YM V47]
\label{prop:v16v25-YM}
YM V47 proves two facts relevant to claim discipline:
\begin{enumerate}
\item the isotropic full-radius claims at two opposite chart
centres are exactly false, by rational negative witnesses;
\item anisotropic radius repair gives sixteen positive charts and
extends both direct response signs to paired regions
\([0,\pi/4]^2\times[0,2t]\).
\end{enumerate}
It does not prove paired quarter cells, a full-torus inequality, a
physical transfer limit, an interacting continuum measure, or a
mass gap.
\end{proposition}

\begin{proof}
These are respectively the countertheorem, repaired-chart theorem,
paired-prism theorem, and claim firewall in YM V47.
\end{proof}

\begin{theorem}[V25 companion-audit decision]
\label{thm:v16v25-audit}
No mathematical estimate from YM V47 or NS V26 populates a row of
the local Newton transfer card below.  The YM failure does,
however, select the next strategy: a collapsing exceptional
direction must be separated and scaled explicitly rather than
covered by an unproved uniform gap.
\end{theorem}

\begin{proof}
YM V47 concerns a \(45\)-dimensional Wilson response pencil on
specific momentum boxes.  It contains no Newton inverse,
conditional density, regression covariance, or reflected transfer.
NS V26 concerns rank-one Oseen-kernel derivatives and contains none
of those objects either.  Hence neither supplies a premise below.

The exact YM counterexamples show that a uniform isotropic
certificate can fail while an anisotropic decomposition succeeds.
This is a methodological audit decision, not an import of its
matrix bounds.
\end{proof}

\subsection{The consolidated card}

\begin{definition}[Local Newton transfer card, LNTC]
\label{def:v16v25-LNTC}
For every fixed compact region \(K\), LNTC consists of the
following rows on a common cofinal carrier.
\begin{enumerate}
\item[\textup{(N1)}] \emph{Base transfer:}
\(T_n\geq0\) is trace class at fixed compact volume and
\(\|T_n-T\|_1\to0\).
\item[\textup{(N2)}] \emph{Newton weight:}
\(Q_n\to Q\) almost everywhere after transport and one strict
higher weighted trace moment holds as in (11.18).
\item[\textup{(N3)}] \emph{Reflection locality:}
the constraint branch is support preserving and reflection
equivariant, or the full cross-reflection kernel satisfies the
appropriate positive-kernel condition.
\item[\textup{(N4)}] \emph{Interior constraint geometry:}
the block \(A_{K,n}\) has a uniform positive gap on its declared
range, with all zero or soft modes retained separately.
\item[\textup{(N5)}] \emph{Local source ledger:}
neutral insertions use the mean-zero projection of V17; every
nonneutral constant mode is represented as an explicit channel.
\item[\textup{(N6)}] \emph{Finite collar:}
AFCIC of V22 holds, including the conditional interaction density
and regression mean-cost moment.
\item[\textup{(N7)}] \emph{Far influence:}
the covariance-weighted tail card (18.19) and the projected
\(L^p\)-density bound (19.14) hold.
\item[\textup{(N8)}] \emph{Growing regression carrier:}
the primitive convergence hypotheses of V24 and one strict
Fredholm spectral reserve hold.
\item[\textup{(N9)}] \emph{Local-state compatibility:}
the normalized states on nested finite regions are consistent and
converge on one local observable core.
\item[\textup{(N10)}] \emph{Response:}
the derivative ledger of V7--V8 has the required exponential
stress and derivative-energy moments under those local states.
\end{enumerate}
\end{definition}

\begin{theorem}[What LNTC proves]
\label{thm:v16v25-LNTC}
If rows \textup{(N1)}--\textup{(N10)} hold, then:
\begin{enumerate}
\item at fixed compact volume, the unbounded Newton-sandwiched
positive transfers converge in trace norm and their normalized
boundary states converge in trace norm;
\item for every fixed region \(K\), the normalized interacting
states have a uniform local Newton exponential moment and hence
tight, uniformly integrable local Newton energies;
\item the neutral far Newton influence vanishes under collar
exhaustion, uniformly in cutoff and exterior volume;
\item the Gaussian regression determinant and its projected
interacting comparison remain bounded and converge across the
growing ultraviolet collar;
\item the metric response may be passed to the local-state limit on
the stated response core.
\end{enumerate}
\end{theorem}

\begin{proof}
Items 1--4 are respectively
Theorem~\ref{thm:v16v11-compiler} and
Corollary~\ref{cor:v16v11-normalized},
Theorem~\ref{thm:v16v22-local} and
Corollary~\ref{cor:v16v21-tight},
Theorem~\ref{thm:v16v19-tail}, and
Theorem~\ref{thm:v16v24-determinant} with
Corollary~\ref{cor:v16v24-AFCIC}.

For item 5, rows (N9)--(N10) give convergence of the local
observables and the uniform integrability required in the response
theorem of V7, while the operator derivatives converge by V8.
The covariance formula (7.7) then passes term by term.
\end{proof}

\begin{firewall}[What LNTC does not prove]
\label{fire:v16v25-card}
LNTC is a conditional compiler, not a construction of its rows.
Even if every row held, the result above would still need to be
combined with the chiral gauge measure, Ward identities, anomaly
ledger, Lorentzian reconstruction, and the full Einstein equation
on the limiting state before the SM--Einstein continuum could be
claimed.
\end{firewall}

\subsection{Status of the rows}

\begin{center}
\small
\begin{tabular}{c|p{0.67\textwidth}}
\toprule
row&current status\\ \midrule
N1--N2&proved as implications; physical common-carrier and moment
loaders remain open\\
N3&exact branch ledger proved; physical branch not selected\\
N4&fixed-gap calculus proved; uniform gap not proved\\
N5&neutral tail proved; nonneutral constant channel open\\
N6&conditional compiler proved; physical AFCIC open\\
N7&far-tail implication proved; projected density loader open\\
N8&trace-ideal compiler proved; primitive convergence open\\
N9&target formulated; compatible interacting local states not
constructed\\
N10&response identity and gap estimates proved; physical stress
moments open\\
\bottomrule
\end{tabular}
\end{center}

The earliest structural bottleneck is N4: several later rows use
\(A_{K,n}^+\), and V16 and V23 give exact discontinuity examples
when its gap collapses.  Demanding a fixed gap is too restrictive
if a genuine soft gravitational mode is present.

\begin{assessment}[Audit-adjusted direction]
\label{ass:v16v25-status}
V25 records the YM V47 anisotropic repair, proves its non-import,
and consolidates the complete V11--V24 local Newton chain into
LNTC.  This makes the dependency order explicit and selects the
soft interior mode as the next upstream target.
\end{assessment}

\begin{decision}[V26 renormalized soft-mode coupling]
\label{dec:v16v25-next}
Continue after the audit.  Decompose the interior Newton block into
a uniformly gapped regular sector and finitely many soft
eigenmodes.  Prove the exact scaling law for
\(B_n^*A_n^+B_n\): a soft coupling must be divided by the square
root of its eigenvalue.  Classify vanishing, finite, and divergent
influence limits and give a trace-ideal version for finitely many
soft channels.
\end{decision}

\begin{firewall}[Claim boundary after V25]
\label{fire:v16v25}
V25 is an audit and dependency theorem.  It does not populate LNTC,
resolve a soft mode, construct the interacting local states or
their reflection positivity, recover the physical Einstein
equation, or complete the continuum.
\end{firewall}

\section*{Version V25 append-only record}
\addcontentsline{toc}{section}{Version V25 append-only record}
The complete V24 source was retained before this continuation.  It
had \(230045\) bytes, \(7038\) lines, and SHA--256
\[
 \texttt{38A1CCB25BCC0CA7CACB53CB685A473B64634F4D569AACBA70756FEEC9AF80A4}.
\]
V25 inspected YM V47 and NS V26, consolidated V11--V24 into LNTC,
and selected an explicit soft-mode scaling theorem as the next
upstream task.
\section{Sixteenth-cycle V26: renormalized soft modes in the Newton influence}
\label{sec:v16v26}

LNTC row N4 asks for a gap only after every genuine soft mode has
been retained.  V26 makes that statement exact.  A coupling to a
soft eigenvector must be divided by the square root of its
eigenvalue.  The resulting normalized coupling, rather than the
unscaled block, determines whether the influence vanishes, has a
finite retained limit, or diverges.

\subsection{Exact soft--regular decomposition}

Let \({\cal H}_K,{\cal H}_E\) be Hilbert spaces.  For each \(n\),
suppose
\[
 A_n=
 \sum_{j=1}^m a_{j,n}|e_{j,n}\rangle\langle e_{j,n}|
 \oplus A_{n,\mathrm{reg}},                          \tag{26.1}
\]
where \(a_{j,n}>0\), the \(e_{j,n}\) are orthonormal, and
\[
 A_{n,\mathrm{reg}}\succeq\beta I
 \quad\hbox{on its range},\qquad\beta>0.             \tag{26.2}
\]
Let \(E_n\) be the soft projection and decompose
\[
 B_n=E_nB_n+(I-E_n)B_n=:B_{n,\mathrm{soft}}
 +B_{n,\mathrm{reg}}.                               \tag{26.3}
\]
Define exterior vectors
\[
 b_{j,n}:=B_n^*e_{j,n},\qquad
 \gamma_{j,n}:=a_{j,n}^{-1/2}b_{j,n}.               \tag{26.4}
\]

\begin{theorem}[Exact soft-mode influence formula]
\label{thm:v16v26-formula}
The Schur influence splits orthogonally as
\[
\boxed{
 B_n^*A_n^+B_n
 =
 \sum_{j=1}^m
 |\gamma_{j,n}\rangle\langle\gamma_{j,n}|
 +
 B_{n,\mathrm{reg}}^*
 A_{n,\mathrm{reg}}^+
 B_{n,\mathrm{reg}}.}                               \tag{26.5}
\]
\end{theorem}

\begin{proof}
Spectral calculus gives
\[
 A_n^+
 =
 \sum_{j=1}^m a_{j,n}^{-1}
 |e_{j,n}\rangle\langle e_{j,n}|
 \oplus A_{n,\mathrm{reg}}^+.                       \tag{26.6}
\]
Insert (26.3) on both sides of (26.6).  Orthogonality removes the
soft--regular cross terms, and
\[
 B_n^*
 a_{j,n}^{-1}|e_{j,n}\rangle\langle e_{j,n}|
 B_n
 =
 |\gamma_{j,n}\rangle\langle\gamma_{j,n}|.           \tag{26.7}
\]
Summing proves (26.5).
\end{proof}

\subsection{The square-root scaling is forced by positivity}

Assume the complete interior--exterior Newton block is positive:
\[
 G_n=
 \begin{pmatrix}A_n&B_n\\B_n^*&C_n\end{pmatrix}
 \succeq0.                                           \tag{26.8}
\]

\begin{theorem}[Positive blocks force soft square-root coupling]
\label{thm:v16v26-forced}
For every \(n\),
\[
 \sum_{j=1}^m
 |\gamma_{j,n}\rangle\langle\gamma_{j,n}|
 \preceq
 B_n^*A_n^+B_n
 \preceq C_n.                                        \tag{26.9}
\]
Consequently,
\[
 \|\gamma_{j,n}\|^2\leq\|C_n\|_{\rm op},\qquad
 \|b_{j,n}\|
 \leq a_{j,n}^{1/2}\|C_n\|_{\rm op}^{1/2}.           \tag{26.10}
\]
If \(\sup_n\|C_n\|<\infty\), a soft coupling cannot decay more
slowly than \(a_{j,n}^{1/2}\).
\end{theorem}

\begin{proof}
The first inequality follows from the positive regular term in
(26.5).  Positivity of (26.8) and the Schur-complement argument in
Lemma~\ref{lem:v16v15-schur} give
\[
 B_n^*A_n^+B_n\preceq C_n.                           \tag{26.11}
\]
Testing (26.9) on a unit vector maximizing
\(|\langle\gamma_{j,n},v\rangle|\) gives
\(\|\gamma_{j,n}\|^2\leq\|C_n\|\).
Equation (26.4) then gives (26.10).
\end{proof}

\subsection{Classification of a soft channel}

\begin{proposition}[Vanishing, retained, and divergent regimes]
\label{prop:v16v26-classification}
For one soft mode with \(a_n\downarrow0\):
\begin{enumerate}
\item if \(\|b_n\|=o(a_n^{1/2})\), then its influence converges to
zero in trace norm;
\item if \(a_n^{-1/2}b_n\to\gamma\) strongly, then its influence
converges in trace norm to
\[
 |\gamma\rangle\langle\gamma|;                       \tag{26.12}
\]
\item if \(\|b_n\|/a_n^{1/2}\to\infty\), then both the operator and
trace norms of its influence diverge.
\end{enumerate}
\end{proposition}

\begin{proof}
The soft influence is
\[
 |\gamma_n\rangle\langle\gamma_n|,
 \qquad \gamma_n=a_n^{-1/2}b_n.                      \tag{26.13}
\]
Its operator and trace norms both equal \(\|\gamma_n\|^2\).
For strong convergence, use
\[
\big\|
 |x\rangle\langle x|-|y\rangle\langle y|
\big\|_1
\leq(\|x\|+\|y\|)\|x-y\|.                            \tag{26.14}
\]
The three assertions follow.
\end{proof}

\begin{counterexample}[A finite channel invisible in the naive
block limit]
\label{cnt:v16v26-retained}
Fix \(0\ne\gamma\in{\cal H}_E\) and let \(a_n\downarrow0\).  On
\(\mathbb R\oplus{\cal H}_E\), set
\[
 A_n=a_n,\qquad
 B_n v=a_n^{1/2}\langle\gamma,v\rangle,\qquad
 C=|\gamma\rangle\langle\gamma|.                    \tag{26.15}
\]
Then
\[
 G_n=
 \begin{pmatrix}A_n&B_n\\B_n^*&C\end{pmatrix}
 =
 \begin{pmatrix}a_n^{1/2}\\ \gamma\end{pmatrix}
 \begin{pmatrix}a_n^{1/2}\\ \gamma\end{pmatrix}^{\!*}
 \succeq0,                                           \tag{26.16}
\]
and \(A_n\to0\), \(B_n\to0\), but
\[
 B_n^*A_n^+B_n
 =|\gamma\rangle\langle\gamma|                      \tag{26.17}
\]
for every \(n\).  The pseudoinverse of the naive limiting
interior block would give zero and lose the finite channel.
\end{counterexample}

\begin{proof}
Equation (26.16) proves positivity, and direct substitution proves
(26.17).
\end{proof}

\subsection{Finitely many soft channels and a regular trace ideal}

\begin{theorem}[Trace-norm soft-mode handoff]
\label{thm:v16v26-handoff}
Assume:
\begin{enumerate}
\item \(m\) is fixed and
\[
 \gamma_{j,n}\to\gamma_j
 \quad\hbox{strongly in }{\cal H}_E
 \quad(1\leq j\leq m);                              \tag{26.18}
\]
\item
\[
 A_{n,\mathrm{reg}}^{+1/2}
 \to A_{\mathrm{reg}}^{+1/2}
 \quad\hbox{in operator norm};                       \tag{26.19}
\]
\item
\[
 B_{n,\mathrm{reg}}\to B_{\mathrm{reg}}
 \quad\hbox{in Hilbert--Schmidt norm}.               \tag{26.20}
\]
\end{enumerate}
Then
\[
 B_n^*A_n^+B_n
 \longrightarrow
 J_{\rm ren}
 :=
 \sum_{j=1}^m|\gamma_j\rangle\langle\gamma_j|
 +B_{\mathrm{reg}}^*A_{\mathrm{reg}}^+
  B_{\mathrm{reg}}
 \quad\hbox{in trace norm}.                          \tag{26.21}
\]
\end{theorem}

\begin{proof}
For the soft part, sum (26.14) over the fixed set of modes.
For the regular part, define
\[
 F_n=A_{n,\mathrm{reg}}^{+1/2}B_{n,\mathrm{reg}},
 \qquad
 F=A_{\mathrm{reg}}^{+1/2}B_{\mathrm{reg}}.          \tag{26.22}
\]
Assumptions (26.19)--(26.20) and the ideal inequality give
\(\|F_n-F\|_2\to0\).  Hence
\[
 \|F_n^*F_n-F^*F\|_1
 \leq(\|F_n\|_2+\|F\|_2)\|F_n-F\|_2\to0.            \tag{26.23}
\]
Combine this with the exact decomposition (26.5).
\end{proof}

\begin{corollary}[Renormalized determinant carrier]
\label{cor:v16v26-determinant}
Let \(C_{E,n}^{1/2}\to C_E^{1/2}\) in Hilbert--Schmidt norm and
suppose the influence operators in (26.21) converge in operator
norm and are uniformly bounded.  Then
\[
 C_{E,n}^{1/2}
 (B_n^*A_n^+B_n)
 C_{E,n}^{1/2}
 \longrightarrow
 C_E^{1/2}J_{\rm ren}C_E^{1/2}
 \quad\hbox{in trace norm}.                          \tag{26.24}
\]
Any strict Gaussian determinant reserve for the limiting carrier
therefore persists cofinally.
\end{corollary}

\begin{proof}
Insert and subtract the two mixed products and use
\[
 \|RST\|_1\leq\|R\|_2\|S\|_{\rm op}\|T\|_2.          \tag{26.25}
\]
The determinant assertion follows from
Theorem~\ref{thm:v16v4-uniform}.
\end{proof}

Strong alignment in (26.18) is essential.  Bounded normalized
couplings may drift through orthogonal exterior directions and
fail trace-norm compactness, exactly as in the rank-one drift
counterexample of V9; that example is not repeated.

\begin{assessment}[Physical status]
\label{ass:v16v26-status}
V26 replaces the false demand for a uniform gap across all
interior modes by an exact soft--regular theorem.  Positivity forces
each coupling to scale at most like the square root of its soft
eigenvalue.  Strong limits of the normalized couplings survive as
finite rank-one channels in the renormalized influence, and the
regular sector retains the fixed-gap calculus.

The physical programme has not identified the soft gravitational
eigenvectors, proved that their number is uniformly finite, or
established strong alignment of their normalized Standard Model
couplings.
\end{assessment}

\begin{decision}[V27 moving soft subspace]
\label{dec:v16v26-next}
Allow the soft eigenspace itself to move with the regulator.
Transport it by partial isometries, derive a basis-independent
formula for the normalized soft coupling
\(A_{n,\mathrm{soft}}^{-1/2}E_nB_n\), and prove
trace-norm convergence of its Gram operator under
Hilbert--Schmidt alignment.  Separate eigenvalue crossings from
changes in the total soft rank.
\end{decision}

\begin{firewall}[Claim boundary after V26]
\label{fire:v16v26}
V26 proves an abstract soft-mode handoff.  It does not identify or
control the physical soft sector, populate the remaining LNTC
rows, construct the interacting local states or their reflection
positivity, recover the Einstein equation, or complete the
continuum.
\end{firewall}

\section*{Version V26 append-only record}
\addcontentsline{toc}{section}{Version V26 append-only record}
The complete V25 source was retained before this continuation.  It
had \(239094\) bytes, \(7264\) lines, and SHA--256
\[
 \texttt{AAC07BB5E0A71783580A2EE8543546A1957BF3D07900AAE65F4BBB593827BE67}.
\]
V26 proved the exact soft-mode scaling law, the positivity-forced
square-root bound, and a trace-norm renormalized influence limit
retaining finite soft channels.
\section{Sixteenth-cycle V27: moving soft subspaces and basis-independent coupling}
\label{sec:v16v27}

V26 writes the soft influence in an eigenbasis.  Individual
eigenvectors are not stable at crossings, even when the total soft
subspace is.  V27 replaces that basis ledger by one normalized
operator on the whole soft cluster and transports the cluster by a
canonical partial isometry.

\subsection{The normalized soft operator}

Let \(E_n\) be a finite-rank spectral projection of \(A_n\geq0\),
and suppose
\[
 A_{s,n}:=E_nA_nE_n
 \quad\hbox{is strictly positive on }\operatorname{Ran}E_n.
                                                               \tag{27.1}
\]
For \(B_n:{\cal H}_E\to{\cal H}_K\), define
\[
 \Gamma_n
 :=
 A_{s,n}^{-1/2}E_nB_n:
 {\cal H}_E\longrightarrow\operatorname{Ran}E_n.     \tag{27.2}
\]

\begin{theorem}[Basis-independent soft influence]
\label{thm:v16v27-basis}
The contribution of the complete soft cluster to the Schur
influence is
\[
 J_{s,n}
 :=
 B_n^*E_nA_{s,n}^{-1}E_nB_n
 =
 \Gamma_n^*\Gamma_n.                                \tag{27.3}
\]
This formula is invariant under every unitary change of basis
inside \(\operatorname{Ran}E_n\), including eigenvalue crossings
within the cluster.
\end{theorem}

\begin{proof}
Equation (27.3) follows by taking the adjoint in (27.2).  If
\(V\) is unitary on \(\operatorname{Ran}E_n\), the coordinate
representative of \(\Gamma_n\) changes to \(V\Gamma_n\), while
\[
 (V\Gamma_n)^*(V\Gamma_n)
 =\Gamma_n^*\Gamma_n.                               \tag{27.4}
\]
No choice of individual eigenvectors appears in (27.2).
\end{proof}

If
\[
 \begin{pmatrix}A_n&B_n\\B_n^*&C_n\end{pmatrix}
 \succeq0,                                           \tag{27.5}
\]
then Theorem~\ref{thm:v16v26-forced} becomes the
basis-independent inequality
\[
 0\preceq\Gamma_n^*\Gamma_n\preceq C_n.             \tag{27.6}
\]

\subsection{Canonical transport of nearby soft projections}

Let \(E\) be a projection of rank \(m<\infty\), and suppose
\[
 \operatorname{rank}E_n=m,\qquad
 \|E_n-E\|_{\rm op}<1.                               \tag{27.7}
\]
On \(\operatorname{Ran}E\), put
\[
 Q_n:=EE_nE.                                         \tag{27.8}
\]

\begin{lemma}[Canonical soft-space isometry]
\label{lem:v16v27-isometry}
The operator \(Q_n\) is strictly positive on
\(\operatorname{Ran}E\), and
\[
 U_n:=E_nE\,Q_n^{-1/2}:
 \operatorname{Ran}E\longrightarrow\operatorname{Ran}E_n
                                                               \tag{27.9}
\]
is unitary between the two \(m\)-dimensional subspaces.  If
\(E_n\to E\) in operator norm, then
\[
 U_n\longrightarrow E
 \quad\hbox{in operator norm on }\operatorname{Ran}E. \tag{27.10}
\]
\end{lemma}

\begin{proof}
For \(x\in\operatorname{Ran}E\),
\[
 \langle x,Q_nx\rangle
 =\|E_nx\|^2.                                        \tag{27.11}
\]
If \(E_nx=0\), then
\[
 \|x\|=\|(E-E_n)x\|
 \leq\|E-E_n\|\|x\|,
\]
so \(x=0\) under (27.7).  Thus \(Q_n>0\).  Moreover,
\[
 U_n^*U_n
 =Q_n^{-1/2}EE_nEQ_n^{-1/2}=E                      \tag{27.12}
\]
on \(\operatorname{Ran}E\).  Its range lies in
\(\operatorname{Ran}E_n\), and equal finite ranks make it onto.
Finally \(Q_n\to E\), so \(Q_n^{-1/2}\to E\); (27.9) gives
(27.10).
\end{proof}

Define the pulled-back normalized coupling
\[
 \widehat\Gamma_n
 :=
 U_n^*\Gamma_n:
 {\cal H}_E\longrightarrow\operatorname{Ran}E.       \tag{27.13}
\]

\begin{theorem}[Moving-cluster trace-ideal handoff]
\label{thm:v16v27-handoff}
If
\[
 \|\widehat\Gamma_n-\Gamma\|_2\longrightarrow0       \tag{27.14}
\]
for a Hilbert--Schmidt
\(\Gamma:{\cal H}_E\to\operatorname{Ran}E\), then
\[
 \boxed{
 \|J_{s,n}-\Gamma^*\Gamma\|_1\longrightarrow0.}      \tag{27.15}
\]
More precisely,
\[
 \|J_{s,n}-\Gamma^*\Gamma\|_1
 \leq
 (\|\widehat\Gamma_n\|_2+\|\Gamma\|_2)
 \|\widehat\Gamma_n-\Gamma\|_2.                     \tag{27.16}
\]
\end{theorem}

\begin{proof}
Because \(U_n\) is unitary between the soft spaces,
\[
 J_{s,n}
 =\Gamma_n^*\Gamma_n
 =\widehat\Gamma_n^*\widehat\Gamma_n.                \tag{27.17}
\]
Apply the Hilbert--Schmidt product estimate (24.8).
\end{proof}

\begin{corollary}[Coordinate choice does not change the limit]
\label{cor:v16v27-coordinate}
Let \(V_n:\operatorname{Ran}E\to\operatorname{Ran}E_n\) be any
other unitary identification.  Then
\[
 V_n^*\Gamma_n=W_n\widehat\Gamma_n,                  \tag{27.18}
\]
where \(W_n=V_n^*U_n\) is unitary on
\(\operatorname{Ran}E\), and both representatives have the same
Gram operator.  Hence the physical datum is the trace-class limit
of \(J_{s,n}\), not a coordinatewise limit of chosen soft
eigenvectors.
\end{corollary}

\begin{proof}
Equation (27.18) is immediate, and
\[
 (W_n\widehat\Gamma_n)^*(W_n\widehat\Gamma_n)
 =\widehat\Gamma_n^*\widehat\Gamma_n.                \tag{27.19}
\]
\end{proof}

\subsection{Crossings inside a fixed cluster}

\begin{proposition}[Internal crossings are harmless]
\label{prop:v16v27-crossing}
Suppose the soft spectral cluster has fixed rank \(m\) and is
separated from the regular spectrum by one contour, while its
positive eigenvalues may cross each other.  If \(A_n\to A\) in
operator norm along the cluster and the contour remains valid,
then the Riesz projections \(E_n\) converge in operator norm.
The constructions (27.2), (27.9), and (27.13) remain
basis independent throughout every internal crossing.
\end{proposition}

\begin{proof}
The resolvent identity on the fixed contour gives convergence of
the Riesz projections.  Functional calculus defines
\(A_{s,n}^{-1/2}\) on the complete positive cluster without
labelling its eigenvalues.  Lemma~\ref{lem:v16v27-isometry} then
transports the entire subspace.  An internal degeneracy changes
only an eigenbasis, which leaves (27.3) invariant.
\end{proof}

This proposition does not assert that
\(\|A_{s,n}^{-1/2}\|\) remains bounded as the cluster approaches
zero.  The normalized product in (27.2), rather than either factor,
must converge.

\subsection{Rank change is a different stratum}

\begin{counterexample}[No unitary transport through a rank jump]
\label{cnt:v16v27-rank}
On \(\mathbb R^2\), let
\[
 A_t=\begin{pmatrix}t&0\\0&1\end{pmatrix},
 \qquad t\geq0,                                      \tag{27.20}
\]
and define the soft space as the kernel at \(t=0\) and as the span
of positive eigenvalues below \(1/2\) for \(t>0\).  Then the soft
rank is one for \(0<t<1/2\) but, under the positive-eigenvalue
definition, zero at \(t=0\).  No unitary of the form (27.9) can
identify these spaces at the endpoint.
\end{counterexample}

\begin{proof}
Unitary finite-dimensional spaces have equal dimension.
\end{proof}

The remedy is to declare the limiting kernel vector as a retained
soft channel before taking the limit, as in V26, or to use separate
rank strata.  Relabelling a rank jump as an eigenvalue crossing
would lose precisely the finite channel exhibited in (26.17).

\begin{assessment}[Physical status]
\label{ass:v16v27-status}
V27 gives a basis-independent normalized soft coupling, a canonical
transport for nearby fixed-rank soft projections, and a
Hilbert--Schmidt-to-trace-norm handoff for the moving soft Gram
operator.  Internal eigenvalue crossings are harmless; changes of
soft rank require separate strata.

The physical theory has not proved fixed soft rank, convergence of
the Riesz projections, or Hilbert--Schmidt alignment of
\(\widehat\Gamma_n\).  Positivity supplies boundedness through
(27.6), but boundedness alone does not select a unique limit.
\end{assessment}

\begin{decision}[V28 compactness from trace-class domination]
\label{dec:v16v27-next}
Exploit (27.6).  Prove that if \(C_n\to C\) in trace norm, every
sequence \(0\leq J_{s,n}\leq C_n\) is precompact in trace norm.
Then show that weak-operator convergence of the soft Gram operators
identifies the unique trace-norm limit.  This will replace
Hilbert--Schmidt alignment of a chosen square root by an invariant
compactness-and-identification criterion.
\end{decision}

\begin{firewall}[Claim boundary after V27]
\label{fire:v16v27}
V27 does not identify the physical soft stratum or its limit,
populate the remaining LNTC rows, construct reflection-positive
local states, recover the limiting Einstein equation, or complete
the continuum.
\end{firewall}

\section*{Version V27 append-only record}
\addcontentsline{toc}{section}{Version V27 append-only record}
The complete V26 source was retained before this continuation.  It
had \(248331\) bytes, \(7576\) lines, and SHA--256
\[
 \texttt{551CE7665D32669E8AEBB585FE9CB991F37BF4DE10D4EAE81F08F80C3298110E}.
\]
V27 replaced eigenvector coordinates by a basis-independent
normalized soft coupling and proved trace-norm convergence across
moving fixed-rank soft subspaces.
\section{Sixteenth-cycle V28: trace-class domination compactifies the soft influence}
\label{sec:v16v28}

V27 obtains trace-norm convergence from Hilbert--Schmidt alignment
of a transported square root \(\widehat\Gamma_n\).  The physical
influence is its Gram operator, so a criterion stated directly for
that positive operator is preferable.  Positivity of the complete
Newton block gives
\[
 0\preceq J_{s,n}\preceq C_n.                        \tag{28.1}
\]
V28 proves that trace-norm convergence of the dominating exterior
covariances makes the soft influences trace-norm precompact.

\subsection{Uniform finite-rank approximation}

\begin{lemma}[Compression bound for a positive trace-class operator]
\label{lem:v16v28-compression}
Let \(J\geq0\) be trace class, let \(P\) be an orthogonal
projection, and put \(Q=I-P\).  Then
\[
 \|J-PJP\|_1
 \leq
 \operatorname{Tr}(QJ)
 +2\sqrt{\operatorname{Tr}(QJ)\operatorname{Tr}J}.   \tag{28.2}
\]
\end{lemma}

\begin{proof}
Decompose
\[
 J-PJP=QJQ+QJP+PJQ.                                 \tag{28.3}
\]
The first term is positive and has trace norm
\(\operatorname{Tr}(QJ)\).  For the off-diagonal block,
\[
\begin{split}
 \|QJP\|_1
 &=
 \|QJ^{1/2}J^{1/2}P\|_1\\
 &\leq
 \|QJ^{1/2}\|_2\|J^{1/2}P\|_2\\
 &=
 \sqrt{\operatorname{Tr}(QJ)
       \operatorname{Tr}(PJ)}
 \leq
 \sqrt{\operatorname{Tr}(QJ)\operatorname{Tr}J}.
                                                               \tag{28.4}
\end{split}
\]
The adjoint block has the same trace norm.  The triangle inequality
proves (28.2).
\end{proof}

\begin{theorem}[Trace-class domination gives trace-norm
precompactness]
\label{thm:v16v28-precompact}
Let \(C_n,C\geq0\) be trace class on a separable Hilbert space and
assume
\[
 \|C_n-C\|_1\longrightarrow0.                        \tag{28.5}
\]
If
\[
 0\preceq J_n\preceq C_n,                            \tag{28.6}
\]
then \(\{J_n:n\geq1\}\) is relatively compact in trace norm.
\end{theorem}

\begin{proof}
The traces of \(C_n\), and hence those of \(J_n\), are uniformly
bounded; write
\[
 M:=\sup_n\operatorname{Tr}J_n<\infty.               \tag{28.7}
\]
Fix \(\varepsilon>0\).  Choose a finite-rank projection \(P\) such
that
\[
 \operatorname{Tr}((I-P)C)<\varepsilon.             \tag{28.8}
\]
For all sufficiently large \(n\),
\[
\begin{split}
 \operatorname{Tr}((I-P)J_n)
 &\leq\operatorname{Tr}((I-P)C_n)\\
 &\leq\operatorname{Tr}((I-P)C)+\|C_n-C\|_1
 <2\varepsilon.                                      \tag{28.9}
\end{split}
\]
Lemma~\ref{lem:v16v28-compression} gives
\[
 \|J_n-PJ_nP\|_1
 \leq2\varepsilon+2\sqrt{2M\varepsilon}              \tag{28.10}
\]
on that tail.

The operators \(PJ_nP\) lie in a bounded subset of the
finite-dimensional space \(P{\cal L}P\), which is precompact.
Equation (28.10) makes the tail of \(\{J_n\}\) totally bounded in
trace norm.  Adding finitely many initial elements preserves total
boundedness.  Completeness of the trace class proves relative
compactness.
\end{proof}

\subsection{Weak identification upgrades to trace norm}

\begin{theorem}[Weak plus dominated implies trace norm]
\label{thm:v16v28-weak}
Under the assumptions of
Theorem~\ref{thm:v16v28-precompact}, suppose
\[
 J_n\longrightarrow J
 \quad\hbox{in the weak operator topology}.          \tag{28.11}
\]
Then \(J\geq0\), \(J\leq C\), \(J\) is trace class, and
\[
 \boxed{\ \|J_n-J\|_1\longrightarrow0.\ }            \tag{28.12}
\]
\end{theorem}

\begin{proof}
For every \(v\),
\[
 0\leq\langle v,J_nv\rangle
 \leq\langle v,C_nv\rangle.                          \tag{28.13}
\]
Trace-norm convergence of \(C_n\) implies operator-norm
convergence, so passage to the limit gives
\[
 0\leq\langle v,Jv\rangle\leq\langle v,Cv\rangle.    \tag{28.14}
\]
Thus \(0\leq J\leq C\), and \(J\) is trace class.

If (28.12) failed, some subsequence would remain a fixed positive
trace-norm distance from \(J\).  By
Theorem~\ref{thm:v16v28-precompact}, it has a further subsequence
converging in trace norm to an operator \(K\).  Trace-norm
convergence implies weak-operator convergence, so that subsequence
converges weakly to \(K\).  The original weak limit is unique,
hence \(K=J\), contradicting the fixed distance.  Therefore
(28.12) holds.
\end{proof}

\begin{corollary}[A dense matrix-element test is sufficient]
\label{cor:v16v28-dense}
Let \({\cal D}\) be a dense subset of the exterior Hilbert space.
In Theorem~\ref{thm:v16v28-precompact}, it is enough to assume that
there is an operator \(J\) such that
\[
 \langle u,J_nv\rangle\longrightarrow\langle u,Jv\rangle
 \quad\hbox{for all }u,v\in{\cal D}.                 \tag{28.15}
\]
Then (28.12) holds.
\end{corollary}

\begin{proof}
Domination gives
\[
 \|J_n\|_{\rm op}\leq\|C_n\|_{\rm op},               \tag{28.16}
\]
so the \(J_n\) are uniformly bounded.  Approximate arbitrary
\(u,v\) by vectors in \({\cal D}\); (28.16) extends (28.15) to all
matrix elements.  Apply
Theorem~\ref{thm:v16v28-weak}.
\end{proof}

\subsection{Application to moving soft clusters}

\begin{theorem}[Invariant soft Gram handoff]
\label{thm:v16v28-soft}
Let \(J_{s,n}=\Gamma_n^*\Gamma_n\) be the moving-cluster soft
influence of V27.  Suppose the complete positive blocks imply
\[
 0\leq J_{s,n}\leq C_n,\qquad
 C_n\to C\quad\hbox{in trace norm}.                  \tag{28.17}
\]
If the physical soft response matrix elements converge on a dense
exterior source core,
\[
 \langle u,J_{s,n}v\rangle\longrightarrow j(u,v),
                                                               \tag{28.18}
\]
then \(j\) is represented by a unique positive trace-class
operator \(J_s\leq C\), and
\[
 \|J_{s,n}-J_s\|_1\longrightarrow0.                 \tag{28.19}
\]
\end{theorem}

\begin{proof}
The uniform bound (28.16) makes the limiting form \(j\) bounded on
the dense core, so it extends uniquely to an operator \(J_s\).
Corollary~\ref{cor:v16v28-dense} gives (28.19), and the proof of
Theorem~\ref{thm:v16v28-weak} gives \(0\leq J_s\leq C\).
\end{proof}

This theorem does not require a choice of partial isometries
\(U_n\), eigenvectors, or square roots \(\Gamma_n\).  It asks only
for convergence of the physical Gram matrix elements.

\begin{corollary}[Soft Gaussian determinants converge]
\label{cor:v16v28-determinant}
Under Theorem~\ref{thm:v16v28-soft}, if
\[
 \sup_n2c\|C_{0,n}^{1/2}J_{s,n}C_{0,n}^{1/2}\|
 \leq r<1                                            \tag{28.20}
\]
and \(C_{0,n}^{1/2}\to C_0^{1/2}\) in Hilbert--Schmidt norm, then
the covariance-weighted soft carriers converge in trace norm and
their Gaussian quadratic determinants converge.
\end{corollary}

\begin{proof}
Use the three-factor trace-ideal argument of
Corollary~\ref{cor:v16v26-determinant} with (28.19), then apply
Theorem~\ref{thm:v16v4-uniform}.
\end{proof}

\subsection{Why trace-class domination is load bearing}

Uniform operator domination is insufficient.  The rank-one
projections onto an orthonormal sequence satisfy
\[
 0\leq|e_n\rangle\langle e_n|\leq I,                 \tag{28.21}
\]
but have pairwise trace-norm distance two.  This is the drift
mechanism already recorded in V9.  The new theorem succeeds because
the trace-class limit \(C\) has uniformly small finite-rank
complements, as quantified in (28.9).

\begin{assessment}[Physical status]
\label{ass:v16v28-status}
V28 converts positive trace-class domination plus weak physical
matrix-element convergence into trace-norm convergence of the
moving soft Gram operator.  This is basis independent and replaces
Hilbert--Schmidt alignment of a chosen normalized coupling by a
compactness-and-identification criterion.

The coupled theory has not proved trace-norm convergence of the
dominating exterior covariance or convergence of the soft response
matrix elements on a physical source core.
\end{assessment}

\begin{decision}[V29 identifying soft matrix elements]
\label{dec:v16v28-next}
Express
\(\langle u,J_{s,n}v\rangle\) as the minimized positive-block
energy or, equivalently, as a difference between the exterior
block and the positive Schur remainder.  Prove a variational
convergence theorem that identifies these matrix elements from
Mosco convergence of the full Newton forms, avoiding explicit
soft eigenvectors.
\end{decision}

\begin{firewall}[Claim boundary after V28]
\label{fire:v16v28}
V28 proves a compactness theorem conditional on trace-class
domination and weak identification.  It does not establish those
physical inputs, populate the other LNTC rows, construct
reflection-positive local states, recover the Einstein equation,
or complete the continuum.
\end{firewall}

\section*{Version V28 append-only record}
\addcontentsline{toc}{section}{Version V28 append-only record}
The complete V27 source was retained before this continuation.  It
had \(257059\) bytes, \(7846\) lines, and SHA--256
\[
 \texttt{4DAEAC4A2FCB9F4EF04190551610747A8AFB95BD2BC3EE455B85652B887995A3}.
\]
V28 proved trace-norm precompactness from positive trace-class
domination and upgraded weak soft Gram convergence to trace-norm
convergence without choosing soft eigenvectors.
\section{Sixteenth-cycle V29: variational identification of the soft Schur influence}
\label{sec:v16v29}

V28 reduces trace-norm convergence of the moving soft influence to
convergence of its matrix elements.  V29 identifies those matrix
elements without choosing soft eigenvectors: the Schur remainder is
the minimum of the complete positive Newton energy over the
interior variable.  Convergence of these minima requires
fibrewise form convergence and equicoercivity of near-minimizers.

\subsection{The shorted Newton form}

Let \({\cal H}_K,{\cal H}_E\) be real Hilbert spaces and let
\[
 G=
 \begin{pmatrix}A&B\\B^*&C\end{pmatrix}\succeq0,     \tag{29.1}
\]
where \(A^+\) is defined on the range containing
\(\operatorname{Ran}B\).  Put
\[
 {\mathfrak q}(x,v)
 :=
 \langle x,Ax\rangle
 +2\langle x,Bv\rangle
 +\langle v,Cv\rangle.                               \tag{29.2}
\]

\begin{theorem}[Variational Schur identity]
\label{thm:v16v29-short}
For every \(v\in{\cal H}_E\),
\[
 \inf_{x\in{\cal H}_K}{\mathfrak q}(x,v)
 =
 \langle v,(C-B^*A^+B)v\rangle.                      \tag{29.3}
\]
The minimum is attained at
\[
 x=-A^+Bv+k,\qquad k\in\ker A.                       \tag{29.4}
\]
Consequently the influence quadratic form is
\[
 \langle v,Jv\rangle
 =
 \langle v,Cv\rangle
 -\inf_x{\mathfrak q}(x,v),
 \qquad J=B^*A^+B.                                   \tag{29.5}
\]
\end{theorem}

\begin{proof}
The completion of the square in Lemma
\ref{lem:v16v15-schur} gives
\[
\begin{split}
 {\mathfrak q}(x,v)
 ={}&
 \langle x+A^+Bv,\,
 A(x+A^+Bv)\rangle\\
 &+\langle v,(C-B^*A^+B)v\rangle.                   \tag{29.6}
\end{split}
\]
The first term is nonnegative and vanishes exactly for (29.4).
This proves (29.3)--(29.5).
\end{proof}

\subsection{Fibrewise variational convergence}

Let \({\mathfrak q}_n,{\mathfrak q}\) be nonnegative closed
quadratic forms on
\({\cal H}_K\oplus{\cal H}_E\).  For fixed \(v\), define
\[
 F_n(v):=\inf_x{\mathfrak q}_n(x,v),\qquad
 F(v):=\inf_x{\mathfrak q}(x,v).                    \tag{29.7}
\]

\begin{definition}[Fibrewise Mosco convergence at \(v\)]
\label{def:v16v29-fibre-Mosco}
FMC at \(v\) means:
\begin{enumerate}
\item whenever \(x_n\rightharpoonup x\) in \({\cal H}_K\),
\[
 {\mathfrak q}(x,v)
 \leq\liminf_n{\mathfrak q}_n(x_n,v);                \tag{29.8}
\]
\item for every \(x\in{\cal H}_K\), there are
\(x_n\to x\) strongly such that
\[
 \limsup_n{\mathfrak q}_n(x_n,v)
 \leq{\mathfrak q}(x,v).                             \tag{29.9}
\]
\end{enumerate}
\end{definition}

\begin{definition}[Fibrewise equicoercivity at \(v\)]
\label{def:v16v29-equicoercive}
FEC at \(v\) means that every sequence satisfying
\[
 {\mathfrak q}_n(x_n,v)\leq F_n(v)+o(1)              \tag{29.10}
\]
has a weakly convergent subsequence in \({\cal H}_K\).
\end{definition}

\begin{theorem}[Convergence of minimized Newton energies]
\label{thm:v16v29-minima}
If FMC and FEC hold at \(v\), and \(F(v)<\infty\), then
\[
 F_n(v)\longrightarrow F(v).                        \tag{29.11}
\]
\end{theorem}

\begin{proof}
Choose \(x\) with
\[
 {\mathfrak q}(x,v)\leq F(v)+\varepsilon.            \tag{29.12}
\]
The recovery sequence in (29.9) gives
\[
 \limsup_nF_n(v)
 \leq\limsup_n{\mathfrak q}_n(x_n,v)
 \leq F(v)+\varepsilon.                              \tag{29.13}
\]
Letting \(\varepsilon\downarrow0\) proves the upper bound.

For the lower bound, choose \(y_n\) with
\[
 {\mathfrak q}_n(y_n,v)\leq F_n(v)+n^{-1}.           \tag{29.14}
\]
The upper bound makes \(F_n(v)\) bounded above.  FEC gives a
subsequence \(y_{n_k}\rightharpoonup y\).  FMC yields
\[
 F(v)\leq{\mathfrak q}(y,v)
 \leq\liminf_k{\mathfrak q}_{n_k}(y_{n_k},v)
 =\liminf_kF_{n_k}(v).                               \tag{29.15}
\]
Apply this argument to a subsequence realizing
\(\liminf_nF_n(v)\).  Combine with (29.13).
\end{proof}

\subsection{Identification of the influence matrix elements}

Let
\[
 G_n=
 \begin{pmatrix}A_n&B_n\\B_n^*&C_n\end{pmatrix}
 \succeq0,\qquad
 J_n=B_n^*A_n^+B_n.                                  \tag{29.16}
\]

\begin{theorem}[Variational soft-Gram identification]
\label{thm:v16v29-identification}
Let \({\cal D}\subset{\cal H}_E\) be a dense real linear subspace.
Assume:
\begin{enumerate}
\item \(C_n\to C\) in trace norm;
\item \(0\leq J_n\leq C_n\);
\item FMC and FEC hold at every \(v\in{\cal D}\) for the full
forms \({\mathfrak q}_n\to{\mathfrak q}\).
\end{enumerate}
Define
\[
 j(v):=\langle v,Cv\rangle-\inf_x{\mathfrak q}(x,v).
                                                               \tag{29.17}
\]
Then \(j\) is the quadratic form of a unique positive trace-class
operator \(J\leq C\), and
\[
 \|J_n-J\|_1\longrightarrow0.                        \tag{29.18}
\]
\end{theorem}

\begin{proof}
Theorems~\ref{thm:v16v29-short} and
\ref{thm:v16v29-minima} give
\[
 \langle v,J_nv\rangle\longrightarrow j(v)
 \quad(v\in{\cal D}).                                \tag{29.19}
\]
For \(u,v\in{\cal D}\), real polarization gives
\[
 \langle u,J_nv\rangle
 ={1\over4}\left[
 \langle u+v,J_n(u+v)\rangle
 -\langle u-v,J_n(u-v)\rangle\right].                \tag{29.20}
\]
Thus all matrix elements converge on \({\cal D}\).
Apply Theorem~\ref{thm:v16v28-soft}, whose domination hypothesis is
items 1--2.  It produces the unique \(J\leq C\) representing the
limit and proves (29.18).
\end{proof}

\begin{corollary}[No soft eigenvectors are needed]
\label{cor:v16v29-no-eigenvectors}
Under Theorem~\ref{thm:v16v29-identification}, the renormalized
soft influence is determined entirely by the full positive forms,
their exterior blocks, and fibrewise variational convergence.
Riesz projections, eigenvector labels, and normalized square roots
are unnecessary for identifying \(J\).
\end{corollary}

\begin{proof}
Formula (29.17) and polarization (29.20) contain only those data.
\end{proof}

\subsection{Escaping minimizers are the retained soft channel}

\begin{counterexample}[Mosco convergence without convergence of
the fibre infimum]
\label{cnt:v16v29-escape}
On \(\mathbb R\oplus\mathbb R\), let
\[
 {\mathfrak q}_n(x,v)
 :=(n^{-1}x+v)^2.                                    \tag{29.21}
\]
Then \({\mathfrak q}_n\) Mosco-converges to
\[
 {\mathfrak q}(x,v)=v^2,                             \tag{29.22}
\]
but, for every \(v\),
\[
 F_n(v)=0,\qquad F(v)=v^2.                           \tag{29.23}
\]
\end{counterexample}

\begin{proof}
If \((x_n,v_n)\rightharpoonup(x,v)\), the sequence \(x_n\) is
bounded, so \(n^{-1}x_n\to0\) and
\[
 \liminf_n{\mathfrak q}_n(x_n,v_n)=v^2.
 \]
The constant recovery sequence proves Mosco convergence to
(29.22).  However the minimizing choice for fixed \(v\) is
\[
 x_n=-nv,                                            \tag{29.24}
\]
which escapes every weakly compact set when \(v\ne0\).  Thus FEC
fails, and (29.23) follows.
\end{proof}

The block matrices in this example are exactly
\[
 \begin{pmatrix}n^{-2}&n^{-1}\\n^{-1}&1\end{pmatrix},
                                                               \tag{29.25}
\]
so the missing unit influence is the retained soft channel of
V26.  FEC must therefore be imposed after augmenting the limiting
carrier by that normalized soft coordinate; imposing it on the
naive interior variable would incorrectly discard the channel.

\begin{assessment}[Physical status]
\label{ass:v16v29-status}
V29 identifies the soft Gram limit variationally from the complete
Newton forms and trace-class exterior blocks.  Fibrewise Mosco
convergence plus equicoercivity gives the missing weak matrix
elements, and V28 upgrades them to trace norm.

The theorem also locates the soft obstruction exactly: escaping
interior minimizers must be represented by retained normalized
coordinates before FEC can hold.  The physical augmented carrier
and its fibrewise Mosco convergence remain unconstructed.
\end{assessment}

\begin{decision}[V30 companion audit and augmented equicoercivity]
\label{dec:v16v29-next}
Perform the compulsory companion audit.  Then formulate an
augmented interior variable containing the finitely many normalized
soft coordinates of V26.  Prove FEC from a uniform regular-sector
gap plus coercivity in those retained coordinates, and continue in
V31 with the resulting variational handoff.
\end{decision}

\begin{firewall}[Claim boundary after V29]
\label{fire:v16v29}
V29 proves a conditional variational identification theorem.  It
does not construct the augmented physical carrier, prove its
Mosco convergence, populate the other LNTC rows, construct
reflection-positive local states, recover the Einstein equation,
or complete the continuum.
\end{firewall}

\section*{Version V29 append-only record}
\addcontentsline{toc}{section}{Version V29 append-only record}
The complete V28 source was retained before this continuation.  It
had \(266080\) bytes, \(8119\) lines, and SHA--256
\[
 \texttt{ACD1512F78FC23A379EFC456C5D70358CEFD3B032A4AF78D11456ADC90B17FB1}.
\]
V29 identified the soft Schur influence through fibrewise
variational convergence and proved that escaping minimizers are
exactly the retained soft-channel obstruction.
\section{Sixteenth-cycle V30: companion audit and augmented soft-coordinate coercivity}
\label{sec:v16v30}

This compulsory checkpoint records the current companion notes and
then repairs the escaping-minimizer obstruction of V29.  The
original soft coordinate diverges like \(a_n^{-1/2}\); its energy
coordinate
\[
 z_n=A_{s,n}^{1/2}x_{s,n}                            \tag{30.1}
\]
remains finite and must be retained in the limiting variational
carrier.

\subsection{Compulsory companion audit}

\begin{audit}[Files inspected at V30]
\label{aud:v16v30-files}
The newest companion checkpoints remain
\[
\begin{array}{c|r|r}
\hbox{file}&\hbox{lines}&\hbox{bytes}\\ \hline
\texttt{Renewal\_Yang\_Mills\_Mass\_Gap\_Research\_Notes\_V47.tex}
 &11182&408812\\
\texttt{Renewal\_NS\_Stability\_Research\_Notes\_V26.tex}
 &5319&278283.
\end{array}                                             \tag{30.2}
\]
Their SHA--256 values are
\[
\begin{gathered}
\texttt{2FDB54605CBCBC68D9C8BC9D69FD1708257A25425F7BE227DA82AC9BA7020B8F},\\
\texttt{82C887F8C2C69E4F28FAD7C3DC8DE5B9099CB5A4BF431EBA1FFF3E91935978AD}.
\end{gathered}                                         \tag{30.3}
\]
\end{audit}

\begin{theorem}[V30 companion-audit decision]
\label{thm:v16v30-audit}
Neither companion checkpoint supplies augmented Newton
equicoercivity or convergence of a gravitational soft channel.
\end{theorem}

\begin{proof}
YM V47 gives an anisotropically repaired finite Wilson-pencil atlas
but no constraint pseudoinverse, soft gravitational mode, or
continuum transfer.  NS V26 gives deterministic Oseen-kernel
derivative norms.  Neither contains the variables or hypotheses
below, so no estimate is imported.
\end{proof}

\subsection{Augmented coordinates}

Let the soft projection \(E_n\), canonical transporter \(U_n\),
soft block \(A_{s,n}\), and pulled-back normalized coupling
\(\widehat\Gamma_n\) be as in V27.  Decompose
\[
 x=x_{s,n}+y,\qquad x_{s,n}=E_nx,\quad
 y=(I-E_n)x.                                         \tag{30.4}
\]
Define the transported energy coordinate
\[
 z:=U_n^*A_{s,n}^{1/2}x_{s,n}
 \in{\cal Z}:=\operatorname{Ran}E.                   \tag{30.5}
\]
Conversely,
\[
 x_{s,n}=A_{s,n}^{-1/2}U_nz.                        \tag{30.6}
\]

\begin{proposition}[Exact augmented Newton form]
\label{prop:v16v30-form}
For the positive block form
\[
 {\mathfrak q}_n(x,v)
 =
 \langle x,A_nx\rangle
 +2\langle x,B_nv\rangle+\langle v,C_nv\rangle,      \tag{30.7}
\]
one has
\[
\begin{split}
 \widetilde{\mathfrak q}_n(z,y;v)
 :={}&{\mathfrak q}_n
 (A_{s,n}^{-1/2}U_nz+y,v)\\
 ={}&
 \|z\|^2+2\langle z,\widehat\Gamma_nv\rangle\\
 &+\langle y,A_{r,n}y\rangle
 +2\langle y,B_{r,n}v\rangle
 +\langle v,C_nv\rangle,                             \tag{30.8}
\end{split}
\]
where
\[
 A_{r,n}=(I-E_n)A_n(I-E_n),\qquad
 B_{r,n}=(I-E_n)B_n.                                \tag{30.9}
\]
\end{proposition}

\begin{proof}
The spectral projection \(E_n\) reduces \(A_n\), so the
soft--regular \(A_n\) cross term vanishes.  Equations
(30.5)--(30.6) give
\[
 \langle x_{s,n},A_{s,n}x_{s,n}\rangle=\|z\|^2.
                                                               \tag{30.10}
\]
Moreover,
\[
\begin{split}
 \langle x_{s,n},E_nB_nv\rangle
 &=
 \langle z,U_n^*A_{s,n}^{-1/2}E_nB_nv\rangle\\
 &=\langle z,\widehat\Gamma_nv\rangle.               \tag{30.11}
\end{split}
\]
The remaining terms are precisely the regular and exterior terms
in (30.8).
\end{proof}

\subsection{Uniform equicoercivity in the augmented carrier}

\begin{theorem}[Augmented sublevel coercivity]
\label{thm:v16v30-coercivity}
Assume
\[
 A_{r,n}\succeq\beta I,\qquad
 \sup_n\|\widehat\Gamma_n\|_{\rm op}\leq g,\qquad
 \sup_n\|B_{r,n}\|_{\rm op}\leq b,                  \tag{30.12}
\]
with \(\beta>0\), and \(C_n\geq0\).  Then, for every
\(z\in{\cal Z}\), regular \(y\), and exterior \(v\),
\[
\begin{split}
 \widetilde{\mathfrak q}_n(z,y;v)
 \geq{}&
 {1\over2}\|z\|^2+{\beta\over2}\|y\|^2\\
 &-2g^2\|v\|^2-{2b^2\over\beta}\|v\|^2.             \tag{30.13}
\end{split}
\]
Consequently, for fixed \(v\), every family satisfying
\[
 \widetilde{\mathfrak q}_n(z_n,y_n;v)\leq M          \tag{30.14}
\]
is bounded in \({\cal Z}\oplus{\cal H}_{r}\).  If the regular
spaces have been transported to one reflexive Hilbert carrier,
these sublevels are weakly sequentially precompact.
\end{theorem}

\begin{proof}
Young's inequality gives
\[
 2|\langle z,\widehat\Gamma_nv\rangle|
 \leq{1\over2}\|z\|^2+2g^2\|v\|^2,                  \tag{30.15}
\]
and
\[
 2|\langle y,B_{r,n}v\rangle|
 \leq{\beta\over2}\|y\|^2
 +{2b^2\over\beta}\|v\|^2.                          \tag{30.16}
\]
Insert these estimates into (30.8), use
\(\langle y,A_{r,n}y\rangle\geq\beta\|y\|^2\), and
discard the nonnegative exterior term.  This proves (30.13).
The sublevel conclusion follows immediately, and bounded sequences
in a Hilbert space have weakly convergent subsequences.
\end{proof}

\begin{corollary}[FEC after soft augmentation]
\label{cor:v16v30-FEC}
Under (30.12), any near-minimizing sequence for
\(\widetilde{\mathfrak q}_n(\,\cdot\,,\,\cdot\,;v)\)
at fixed \(v\) satisfies the fibrewise equicoercivity condition of
Definition~\ref{def:v16v29-equicoercive} on the augmented carrier,
provided the corresponding infima are uniformly bounded above.
\end{corollary}

\begin{proof}
Near-minimizers lie in a common sublevel by the assumed upper
bound.  Apply Theorem~\ref{thm:v16v30-coercivity}.
\end{proof}

\subsection{The escaping coordinate becomes finite}

\begin{proposition}[Soft minimizer in energy coordinates]
\label{prop:v16v30-minimizer}
For fixed \(y,v\), the unique minimizing augmented soft coordinate
in (30.8) is
\[
 z_n^{\min}=-\widehat\Gamma_nv,                      \tag{30.17}
\]
and its minimized contribution is
\[
 -\|\widehat\Gamma_nv\|^2
 =-\langle v,J_{s,n}v\rangle.                       \tag{30.18}
\]
The corresponding original coordinate is
\[
 x_{s,n}^{\min}
 =-A_{s,n}^{-1/2}U_n\widehat\Gamma_nv,               \tag{30.19}
\]
which may diverge even when \(z_n^{\min}\) converges.
\end{proposition}

\begin{proof}
Complete the square:
\[
 \|z\|^2+2\langle z,\widehat\Gamma_nv\rangle
 =
 \|z+\widehat\Gamma_nv\|^2
 -\|\widehat\Gamma_nv\|^2.                           \tag{30.20}
\]
Equation (27.17) identifies the last norm with the soft Gram
quadratic form.  Equation (30.6) gives (30.19).
\end{proof}

For the V29 counterexample,
\(z_n=n^{-1}x_n\), so the escaping minimizer
\(x_n=-nv\) becomes the constant augmented minimizer \(z_n=-v\).
The lost unit influence is therefore retained exactly.

\begin{assessment}[Physical status]
\label{ass:v16v30-status}
V30 completes the companion audit and constructs the augmented
soft-coordinate form.  A regular-sector gap plus bounded normalized
soft and regular couplings gives fibrewise equicoercivity even when
the original soft minimizers diverge.

The remaining step is fibrewise Mosco convergence of the augmented
forms.  This requires convergence of the normalized soft coupling,
the regular forms and couplings, and the exterior block on common
carriers.
\end{assessment}

\begin{decision}[V31 augmented variational handoff]
\label{dec:v16v30-next}
Continue after the audit.  Assume Mosco convergence of the regular
forms, strong convergence of the normalized soft couplings, and
operator convergence of the exterior blocks.  Prove fibrewise
Mosco convergence of (30.8), convergence of its minima, and
trace-norm convergence of the complete retained influence by
V28--V29.
\end{decision}

\begin{firewall}[Claim boundary after V30]
\label{fire:v16v30}
V30 proves augmented equicoercivity but not augmented Mosco
convergence or the physical convergence hypotheses.  It does not
populate the other LNTC rows, construct reflection-positive local
states, recover the Einstein equation, or complete the continuum.
\end{firewall}

\section*{Version V30 append-only record}
\addcontentsline{toc}{section}{Version V30 append-only record}
The complete V29 source was retained before this continuation.  It
had \(275139\) bytes, \(8408\) lines, and SHA--256
\[
 \texttt{928123A6A1D934D31BFA854DC9AC870A9EC7D375E3E07B25AEDA0CDDFD3F9A23}.
\]
V30 inspected YM V47 and NS V26, then constructed the augmented
soft energy coordinate and proved uniform fibrewise
equicoercivity on that carrier.
\section{Sixteenth-cycle V31: augmented Mosco convergence and the retained influence limit}
\label{sec:v16v31}

V30 proves equicoercivity after replacing the original collapsing
soft coordinate by its energy coordinate.  V31 proves fibrewise
Mosco convergence of the augmented forms from primitive convergence
of the normalized soft coupling, regular form, regular coupling,
and exterior block.  The variational theorem of V29 and domination
theorem of V28 then give trace-norm convergence of the complete
retained influence.

\subsection{Common augmented carrier}

Let \({\cal Z}\) be the fixed finite-dimensional soft carrier,
\({\cal H}_r\) the transported regular interior carrier, and
\({\cal H}_E\) the exterior carrier.  Let
\({\mathfrak a}_n,{\mathfrak a}\) be nonnegative closed quadratic
forms on \({\cal H}_r\).  For \(v\in{\cal H}_E\), define
\[
\begin{split}
 \widetilde{\mathfrak q}_n(z,y;v)
 :={}&
 \|z\|^2+2\langle z,\Gamma_nv\rangle
 +{\mathfrak a}_n(y)+2\langle y,D_nv\rangle\\
 &+\langle v,C_nv\rangle,                            \tag{31.1}\\
 \widetilde{\mathfrak q}(z,y;v)
 :={}&
 \|z\|^2+2\langle z,\Gamma v\rangle
 +{\mathfrak a}(y)+2\langle y,Dv\rangle\\
 &+\langle v,Cv\rangle.                              \tag{31.2}
\end{split}
\]
Here \(\Gamma_n\) is the transported normalized soft coupling and
\(D_n\) the regular coupling.

\begin{theorem}[Primitive augmented Mosco theorem]
\label{thm:v16v31-Mosco}
Assume:
\begin{enumerate}
\item \({\mathfrak a}_n\) Mosco-converges to
\({\mathfrak a}\);
\item for one \(\beta>0\),
\[
 {\mathfrak a}_n(y)\geq\beta\|y\|^2;                \tag{31.3}
\]
\item \(\Gamma_n\to\Gamma\) and \(D_n\to D\) strongly as
operators, with uniformly bounded operator norms;
\item \(C_n\to C\) in operator norm.
\end{enumerate}
Then, for every fixed \(v\), the forms
\[
 (z,y)\longmapsto
 \widetilde{\mathfrak q}_n(z,y;v)                   \tag{31.4}
\]
Mosco-converge on \({\cal Z}\oplus{\cal H}_r\) to
\(\widetilde{\mathfrak q}(\,\cdot\,,\,\cdot\,;v)\).
Their near-minimizing sequences are weakly precompact whenever
their infima are uniformly bounded above.
\end{theorem}

\begin{proof}
Suppose \((z_n,y_n)\rightharpoonup(z,y)\).  Since
\({\cal Z}\) is finite dimensional, \(z_n\to z\) strongly.
Strong operator convergence gives
\[
 \Gamma_nv\to\Gamma v,\qquad D_nv\to Dv.            \tag{31.5}
\]
Hence
\[
 \langle z_n,\Gamma_nv\rangle\to
 \langle z,\Gamma v\rangle,\qquad
 \langle y_n,D_nv\rangle\to
 \langle y,Dv\rangle.                                \tag{31.6}
\]
The second convergence uses weak convergence of \(y_n\) and strong
convergence of \(D_nv\).  Operator convergence of \(C_n\) handles
the exterior scalar term.  The Mosco lower bound for
\({\mathfrak a}_n\) therefore yields
\[
 \widetilde{\mathfrak q}(z,y;v)
 \leq
 \liminf_n
 \widetilde{\mathfrak q}_n(z_n,y_n;v).               \tag{31.7}
\]

For recovery, choose \(y_n\to y\) strongly with
\[
 \limsup_n{\mathfrak a}_n(y_n)
 \leq{\mathfrak a}(y),                               \tag{31.8}
\]
and take \(z_n=z\).  All linear and exterior terms converge by
(31.5), giving the Mosco upper bound.

Finally, the uniform bounds in item 3 and coercivity (31.3) put
(31.1) under Theorem~\ref{thm:v16v30-coercivity}.  That theorem
gives weak precompactness of common sublevels and hence of
near-minimizers when the infima are uniformly bounded above.
\end{proof}

\begin{corollary}[Convergence of augmented fibre minima]
\label{cor:v16v31-minima}
Under Theorem~\ref{thm:v16v31-Mosco}, for every fixed \(v\),
\[
\inf_{z,y}\widetilde{\mathfrak q}_n(z,y;v)
\longrightarrow
\inf_{z,y}\widetilde{\mathfrak q}(z,y;v).            \tag{31.9}
\]
\end{corollary}

\begin{proof}
The Mosco conclusion supplies FMC and the sublevel conclusion
supplies FEC.  Apply
Theorem~\ref{thm:v16v29-minima}.
\end{proof}

\subsection{Trace-norm convergence of the complete influence}

Suppose each augmented form is obtained from a positive full block,
so its complete exterior influence \(J_n\) satisfies
\[
 0\leq J_n\leq C_n.                                  \tag{31.10}
\]

\begin{theorem}[Augmented retained-influence handoff]
\label{thm:v16v31-influence}
In addition to Theorem~\ref{thm:v16v31-Mosco}, assume
\[
 C_n\to C\quad\hbox{in trace norm}.                  \tag{31.11}
\]
Then there is a unique positive trace-class \(J\leq C\) such that
\[
 \boxed{\ \|J_n-J\|_1\longrightarrow0.\ }            \tag{31.12}
\]
Its quadratic form is determined variationally by
\[
 \langle v,Jv\rangle
 =
 \langle v,Cv\rangle
 -\inf_{z,y}\widetilde{\mathfrak q}(z,y;v).          \tag{31.13}
\]
\end{theorem}

\begin{proof}
For every \(n\), the change from the original soft coordinate to
\(z\) is bijective at finite regulator.  Therefore the
variational Schur identity gives
\[
 \langle v,J_nv\rangle
 =
 \langle v,C_nv\rangle
 -\inf_{z,y}\widetilde{\mathfrak q}_n(z,y;v).        \tag{31.14}
\]
Equations (31.9) and (31.11) prove convergence of these quadratic
forms for every \(v\).  Polarization gives convergence of all
matrix elements.  The positive domination (31.10) and trace-norm
convergence of \(C_n\) allow
Theorem~\ref{thm:v16v28-weak} to upgrade that convergence to
(31.12).  Passage to the limit in (31.14) gives (31.13).
\end{proof}

\begin{corollary}[Soft and regular channels are both retained]
\label{cor:v16v31-split}
If the limit regular form is represented by
\(A_r\geq\beta I\), then
\[
 J=
 \Gamma^*\Gamma+D^*A_r^{-1}D.                       \tag{31.15}
\]
\end{corollary}

\begin{proof}
Complete the square separately in (31.2):
\[
\begin{split}
 \|z\|^2+2\langle z,\Gamma v\rangle
 &=
 \|z+\Gamma v\|^2-\|\Gamma v\|^2,\\
 \langle y,A_ry\rangle+2\langle y,Dv\rangle
 &=
 \|A_r^{1/2}y+A_r^{-1/2}Dv\|^2
 -\|A_r^{-1/2}Dv\|^2.                               \tag{31.16}
\end{split}
\]
Insert the two minima into (31.13).
\end{proof}

\subsection{Convergence of minimizers}

\begin{proposition}[Weak convergence of augmented minimizers]
\label{prop:v16v31-minimizers}
Suppose the limit augmented form has a unique minimizer
\((z(v),y(v))\) for a fixed \(v\).  If
\((z_n(v),y_n(v))\) are exact minimizers of (31.1), then
\[
 (z_n(v),y_n(v))
 \rightharpoonup(z(v),y(v)).                         \tag{31.17}
\]
If, in addition, the regular forms satisfy the uniform norm
recovery property
\[
 {\mathfrak a}_n(y_n)\to{\mathfrak a}(y)
 \ \hbox{and }y_n\rightharpoonup y
 \quad\Longrightarrow\quad y_n\to y,                \tag{31.18}
\]
for minimizing sequences, then convergence in (31.17) is strong.
\end{proposition}

\begin{proof}
Equicoercivity gives a weakly convergent subsequence of any
subsequence.  The Mosco lower bound and convergence of minimum
values show that every cluster point minimizes the limit form.
Uniqueness makes it \((z(v),y(v))\), so the full sequence converges
weakly.  The soft coordinate converges strongly automatically
because \({\cal Z}\) is finite dimensional; (31.18) gives strong
regular convergence.
\end{proof}

\subsection{What the augmentation has achieved}

The V29 counterexample failed FEC because the minimizing original
coordinate was \(x_n=-nv\).  In the augmented carrier it is
\(z_n=-v\), the forms are uniformly coercive, and (31.15) retains
the unit soft influence.  Thus V31 closes the analytic soft-mode
handoff under primitive convergence of the augmented coefficients.

\begin{assessment}[Physical status]
\label{ass:v16v31-status}
V31 proves fibrewise Mosco convergence and trace-norm convergence
of the complete retained soft-plus-regular influence on an
augmented carrier.  No eigenvector labels enter the limit, and
escaping original minimizers become ordinary bounded augmented
coordinates.

The physical theory still must construct this augmented carrier,
show that only finitely many soft modes occur locally, and prove
the primitive strong and trace-class convergences in
Theorem~\ref{thm:v16v31-Mosco}.
\end{assessment}

\begin{decision}[V32 stability of exponential influence moments]
\label{dec:v16v31-next}
Use the trace-norm influence convergence (31.12) and convergence of
the exterior covariances to prove trace-norm convergence of the
covariance-weighted retained carriers.  Under one strict spectral
reserve, derive convergence of their Gaussian exponential moments
and an \(L^p\) projected interacting extension.
\end{decision}

\begin{firewall}[Claim boundary after V31]
\label{fire:v16v31}
V31 proves an augmented variational handoff conditional on its
primitive convergence hypotheses.  It does not populate those
hypotheses or the other LNTC rows, construct reflection-positive
local states, recover the Einstein equation, or complete the
continuum.
\end{firewall}

\section*{Version V31 append-only record}
\addcontentsline{toc}{section}{Version V31 append-only record}
The complete V30 source was retained before this continuation.  It
had \(283419\) bytes, \(8669\) lines, and SHA--256
\[
 \texttt{F822A00507C7DC44E3D73DD544DB2E4FE68CEEEB12554F2A53E4125EBF365358}.
\]
V31 proved augmented fibrewise Mosco convergence, convergence of
the minimized energies, and trace-norm convergence of the complete
retained soft-plus-regular influence.
\section{Sixteenth-cycle V32: covariance-weighted retained influence and exponential-moment convergence}
\label{sec:v16v32}

V31 gives trace-norm convergence \(J_n\to J\) of the complete
retained influence.  Gaussian and projected interacting moments
depend on the covariance-weighted carriers
\[
 K_n=S_n^{1/2}J_nS_n^{1/2}.                          \tag{32.1}
\]
V32 proves their trace-norm convergence under bounded strong
convergence of the covariance square roots, then passes both
Gaussian determinants and \(L^p\)-perturbed moments.

\subsection{Strong multipliers on trace-class operators}

\begin{lemma}[Strong multiplication is trace-norm continuous]
\label{lem:v16v32-strong-trace}
Let \(R_n,R\in{\cal B}({\cal H})\) satisfy
\[
 R_n\to R,\qquad R_n^*\to R^*
 \quad\hbox{strongly},\qquad
 \sup_n\|R_n\|<\infty.                               \tag{32.2}
\]
For every trace-class \(L\),
\[
 \|(R_n-R)L\|_1\to0,\qquad
 \|L(R_n-R)\|_1\to0.                                \tag{32.3}
\]
\end{lemma}

\begin{proof}
Choose finite-rank \(L_m\) with
\(\|L-L_m\|_1\to0\).  On the finite-dimensional range relevant to
\(L_m\), strong convergence is uniform, so
\[
 \|(R_n-R)L_m\|_1\to0.                              \tag{32.4}
\]
The uniform operator bound gives
\[
 \|(R_n-R)(L-L_m)\|_1
 \leq
 \sup_n\|R_n-R\|\,\|L-L_m\|_1.                      \tag{32.5}
\]
First choose \(m\), then \(n\).  The right-multiplication statement
follows by applying the first to adjoints.
\end{proof}

\begin{theorem}[Covariance-weighted trace-norm handoff]
\label{thm:v16v32-weighted}
Let \(J_n,J\geq0\) be trace class with
\[
 \|J_n-J\|_1\to0.                                   \tag{32.6}
\]
Let \(S_n,S\geq0\) be bounded and suppose
\[
 S_n^{1/2}\to S^{1/2}\quad\hbox{strongly},\qquad
 \sup_n\|S_n\|<\infty.                               \tag{32.7}
\]
Then
\[
 \boxed{
 \|S_n^{1/2}J_nS_n^{1/2}
 -S^{1/2}JS^{1/2}\|_1\to0.}                         \tag{32.8}
\]
\end{theorem}

\begin{proof}
Write
\[
\begin{split}
 K_n-K={}&
 S_n^{1/2}(J_n-J)S_n^{1/2}\\
 &+(S_n^{1/2}-S^{1/2})JS_n^{1/2}\\
 &+S^{1/2}J(S_n^{1/2}-S^{1/2}).                     \tag{32.9}
\end{split}
\]
The first term tends to zero in trace norm by the uniform bounds
and (32.6).  Lemma~\ref{lem:v16v32-strong-trace} treats the last
two terms; the square roots are self-adjoint, so strong convergence
of adjoints is automatic.
\end{proof}

\begin{corollary}[Operator-norm covariance convergence is sufficient]
\label{cor:v16v32-op}
If \(S_n\to S\) in operator norm, then (32.7) holds and hence so
does (32.8).
\end{corollary}

\begin{proof}
The positive square-root map is operator-norm continuous on bounded
nonnegative operators.
\end{proof}

\subsection{Gaussian determinants}

\begin{theorem}[Retained Gaussian moment convergence]
\label{thm:v16v32-gaussian}
Let \(K_n,K\geq0\) be the carriers in (32.8).  Fix \(c>0\) and
assume
\[
 \sup_n2c\|K_n\|_{\rm op}\leq r<1.                  \tag{32.10}
\]
Then
\[
 \det(I-2cK_n)^{-1/2}
 \longrightarrow
 \det(I-2cK)^{-1/2},                                \tag{32.11}
\]
and
\[
\left|
 \log{
 \det(I-2cK_n)^{-1/2}
 \over
 \det(I-2cK)^{-1/2}}
\right|
\leq {c\over1-r}\|K_n-K\|_1.                        \tag{32.12}
\]
\end{theorem}

\begin{proof}
Apply Theorem~\ref{thm:v16v4-uniform} with \(\Gamma=2c\).
\end{proof}

\subsection{A common Gaussian coordinate carrier}

Let \(\Xi\) be an isonormal Gaussian process over a separable real
Hilbert space \({\cal H}\).  For \(K\geq0\) trace class, define
\[
 {\cal Q}_K(\Xi)
 :=\sum_j\kappa_j\Xi(e_j)^2,                         \tag{32.13}
\]
where \(Ke_j=\kappa_je_j\).  The series converges in \(L^1\), is
basis independent, and
\[
 {\mathbb E}e^{c{\cal Q}_K}
 =\det(I-2cK)^{-1/2}                                \tag{32.14}
\]
under \(2c\|K\|<1\).

\begin{lemma}[Trace norm controls Gaussian quadratic forms]
\label{lem:v16v32-quadratic}
For self-adjoint trace-class \(D\),
\[
 {\mathbb E}|{\cal Q}_D(\Xi)|
 \leq\|D\|_1.                                       \tag{32.15}
\]
Consequently, if \(K_n\to K\) in trace norm, then
\[
 {\cal Q}_{K_n}\to{\cal Q}_K
 \quad\hbox{in }L^1\hbox{ and in probability}.       \tag{32.16}
\]
\end{lemma}

\begin{proof}
The spectral positive and negative parts give
\[
 |{\cal Q}_D|\leq{\cal Q}_{|D|}.                    \tag{32.17}
\]
Taking expectations yields
\({\mathbb E}{\cal Q}_{|D|}=\operatorname{Tr}|D|\).
Apply this to \(D=K_n-K\).
\end{proof}

\begin{theorem}[Projected interacting exponential-moment handoff]
\label{thm:v16v32-interacting}
Let \(1<p<\infty\), \(q=p/(p-1)\), and let
\(h_n,h\geq0\) be densities on the common isonormal probability
space satisfying
\[
 \|h_n-h\|_{L^p}\to0.                                \tag{32.18}
\]
Suppose \(K_n\to K\) in trace norm and, for some
\(\varepsilon>0\),
\[
 \sup_n2q(1+\varepsilon)c\|K_n\|_{\rm op}
 \leq r_\varepsilon<1.                              \tag{32.19}
\]
Then
\[
\boxed{
 {\mathbb E}\!\left[
 h_ne^{c{\cal Q}_{K_n}}\right]
 \longrightarrow
 {\mathbb E}\!\left[
 h e^{c{\cal Q}_{K}}\right].}                       \tag{32.20}
\]
\end{theorem}

\begin{proof}
Put \(W_n=e^{c{\cal Q}_{K_n}}\) and
\(W=e^{c{\cal Q}_K}\).  Lemma
\ref{lem:v16v32-quadratic} gives \(W_n\to W\) in
probability.  Trace-norm convergence bounds
\(\sup_n\operatorname{Tr}K_n\).  The determinant estimate and
(32.19) give
\[
 \sup_n{\mathbb E}W_n^{q(1+\varepsilon)}<\infty.     \tag{32.21}
\]
Hence \(\{W_n^q\}\) is uniformly integrable.  Convergence in
probability and Vitali's theorem yield
\[
 \|W_n-W\|_{L^q}\to0.                               \tag{32.22}
\]
Finally,
\[
\begin{split}
 |{\mathbb E}[h_nW_n-hW]|
 \leq{}&
 \|h_n-h\|_p\|W_n\|_q
 +\|h\|_p\|W_n-W\|_q,                               \tag{32.23}
\end{split}
\]
which tends to zero.
\end{proof}

\begin{corollary}[Normalized interacting retained states]
\label{cor:v16v32-normalized}
If the quantities in (32.20) are positive normalization factors
and the same argument applies after inserting every bounded
observable from a common projected core, then the corresponding
normalized expectations converge.
\end{corollary}

\begin{proof}
Apply (32.20) to the numerator and denominator and divide by the
positive limiting denominator.
\end{proof}

\subsection{Scope of the common-density assumption}

Condition (32.18) concerns densities after all projected
influence laws have been transported to one Gaussian coordinate
space.  A uniform \(L^p\) bound alone proves moment boundedness but
not convergence; the densities may oscillate between distinct
limits.  Thus V32 separates the compactness row from the
identification row exactly as V28 did for the soft Gram operators.

\begin{assessment}[Physical status]
\label{ass:v16v32-status}
V32 propagates the V31 trace-norm influence limit through bounded
strongly convergent exterior covariances, proves Gaussian
determinant convergence, and gives an \(L^p\) interacting moment
handoff on a common projected Gaussian carrier.

The physical theory has not constructed that common projected
carrier or proved \(L^p\) convergence of its interaction densities.
The strict higher spectral reserve remains a substantive input.
\end{assessment}

\begin{decision}[V33 local reflected-state compactness]
\label{dec:v16v32-next}
Use the uniform local exponential moments from V21--V22 and the
normalization convergence of V32 to prove compactness of the
normalized local boundary states on a countable observable core.
Then identify the additional compatibility and reflected-Gram
conditions needed to assemble one OS-positive local-state limit.
\end{decision}

\begin{firewall}[Claim boundary after V32]
\label{fire:v16v32}
V32 proves covariance-weighted and interacting moment convergence
conditional on common-carrier density convergence.  It does not
construct the physical densities or local states, prove their
reflection positivity or compatibility, recover the Einstein
equation, or complete the continuum.
\end{firewall}

\section*{Version V32 append-only record}
\addcontentsline{toc}{section}{Version V32 append-only record}
The complete V31 source was retained before this continuation.  It
had \(292532\) bytes, \(8936\) lines, and SHA--256
\[
 \texttt{8F627CE9BA02A114DE237859B4030C18E70361BBA7547789C3BFE4A3EE9A37CE}.
\]
V32 proved covariance-weighted trace-norm convergence, Gaussian
determinant convergence, and a projected interacting
\(L^p\)-moment handoff on a common Gaussian coordinate carrier.
\section{Sixteenth-cycle V33: compatible local-state limits with Newton exponential control}
\label{sec:v16v33}

Fifth-cycle V69 already proves that entrywise convergence of every
fixed reflected Gram matrix preserves reflection positivity.  V33
does not repeat that theorem.  It supplies the missing projective
step: local states obtained in different finite regions must have
compatible limits.  The uniform local Newton moments from
V21--V22 then extend convergence from bounded cylinder observables
to controlled unbounded Newton observables.

\subsection{Nested local algebras}

Let
\[
 {\cal A}_1\mathop{\longrightarrow}^{\iota_1^2}
 {\cal A}_2\mathop{\longrightarrow}^{\iota_2^3}
 \cdots                                                  \tag{33.1}
\]
be a countable increasing system of separable unital
\(C^*\)-algebras with injective unital star homomorphisms.  Write
\(\iota_m^\ell:{\cal A}_m\to{\cal A}_\ell\) for
\(m\leq\ell\).  At regulator \(n\), let
\(\omega_{n,m}\) be a state on \({\cal A}_m\).

\begin{definition}[Asymptotic local compatibility]
\label{def:v16v33-compatibility}
ALC means that for every fixed \(m\leq\ell\),
\[
 \delta_{n;m,\ell}
 :=
 \sup_{\|A\|\leq1}
 \left|
 \omega_{n,\ell}(\iota_m^\ell A)
 -\omega_{n,m}(A)
 \right|
 \longrightarrow0.                                  \tag{33.2}
\]
\end{definition}

\begin{theorem}[Projective local-state compactness]
\label{thm:v16v33-projective}
If ALC holds, there are a subsequence \(n_k\) and states
\(\omega_m\) on every \({\cal A}_m\) such that
\[
 \omega_{n_k,m}\mathop{\longrightarrow}^{w^*}\omega_m
 \quad\hbox{for every }m,                            \tag{33.3}
\]
and
\[
 \omega_\ell(\iota_m^\ell A)=\omega_m(A)
 \quad(m\leq\ell,\ A\in{\cal A}_m).                 \tag{33.4}
\]
The compatible family defines a unique state \(\omega\) on the
inductive-limit \(C^*\)-algebra
\[
 {\cal A}_{\rm loc}
 :=\overline{\bigcup_m{\cal A}_m}.                  \tag{33.5}
\]
\end{theorem}

\begin{proof}
Because \({\cal A}_m\) is separable, its state space is weak-star
compact and metrizable.  Choose a subsequence along which
\(\omega_{n,1}\) converges, then a further subsequence for
\(\omega_{n,2}\), and continue.  The diagonal subsequence converges
for every fixed \(m\), proving (33.3).

For \(A\in{\cal A}_m\), (33.2) gives
\[
\begin{split}
 |\omega_\ell(\iota_m^\ell A)-\omega_m(A)|
 \leq{}&
 |\omega_\ell(\iota_m^\ell A)
    -\omega_{n_k,\ell}(\iota_m^\ell A)|\\
 &+\delta_{n_k;m,\ell}\|A\|
 +|\omega_{n_k,m}(A)-\omega_m(A)|.                  \tag{33.6}
\end{split}
\]
Let \(k\to\infty\) to obtain (33.4).
Define \(\omega(\iota_m^\infty A)=\omega_m(A)\).
Compatibility makes this well defined on the algebraic union.
It is positive, unital, and norm one there, so it extends uniquely
by continuity to (33.5).
\end{proof}

\subsection{Reflection compatibility}

Suppose each \({\cal A}_m\) has a reflection
\(\Theta_m\) and a future subalgebra \({\cal A}_{m,+}\), with
\[
 \iota_m^\ell\Theta_m=\Theta_\ell\iota_m^\ell,\qquad
 \iota_m^\ell{\cal A}_{m,+}\subset{\cal A}_{\ell,+}.
                                                               \tag{33.7}
\]

\begin{corollary}[Projective OS-positive limit]
\label{cor:v16v33-OS}
If every \(\omega_{n,m}\) is reflection positive on
\({\cal A}_{m,+}\), then the state \(\omega\) from
Theorem~\ref{thm:v16v33-projective} is reflection positive on the
local future algebra.
\end{corollary}

\begin{proof}
For any finite \(F_1,\ldots,F_N\) in the local future algebra,
choose \(m\) containing them.  The reflected Gram entries converge
by (33.3).  The finite-cutoff Gram matrices are positive
semidefinite, so their entrywise limit is positive semidefinite.
This is the cofinal Gram closure mechanism already proved in
fifth-cycle V69, now applied after the projective compatibility
step.
\end{proof}

\subsection{Passing local Newton observables}

Let \(Q_m\geq0\) be an observable affiliated with
\({\cal A}_m\).  Assume every bounded continuous function of
\(Q_m\) that is constant near infinity after truncation belongs to
\({\cal A}_m\).  For \(R>0\), put
\[
 Q_m^{(R)}:=\min(Q_m,R).                             \tag{33.8}
\]

\begin{theorem}[Exponential moments survive the projective limit]
\label{thm:v16v33-moment}
Suppose, for some \(t>0\),
\[
 \sup_n\omega_{n,m}(e^{tQ_m})\leq M_m<\infty.        \tag{33.9}
\]
Define
\[
 \omega_m(e^{tQ_m})
 :=
 \sup_R\omega_m(e^{tQ_m^{(R)}}).                    \tag{33.10}
\]
Then
\[
 \omega_m(e^{tQ_m})\leq M_m.                        \tag{33.11}
\]
If \(f:[0,\infty)\to\mathbb C\) is continuous and
\[
 |f(q)|\leq C e^{t'q}
 \quad\hbox{for some }0\leq t'<t,                   \tag{33.12}
\]
then
\[
 \omega_{n_k,m}(f(Q_m))
 \longrightarrow\omega_m(f(Q_m)),                  \tag{33.13}
\]
where both sides are defined by uniformly integrable truncation.
\end{theorem}

\begin{proof}
For fixed \(R\), weak-star convergence gives
\[
 \omega_m(e^{tQ_m^{(R)}})
 =
 \lim_k\omega_{n_k,m}(e^{tQ_m^{(R)}})
 \leq M_m.                                          \tag{33.14}
\]
Taking the supremum over \(R\) proves (33.11).

Let \(f_R\) be a bounded continuous truncation agreeing with \(f\)
on \([0,R]\).  The exponential bound gives, uniformly in \(k\),
\[
\begin{split}
 \omega_{n_k,m}
 (|f(Q_m)-f_R(Q_m)|)
 &\leq
 C e^{-(t-t')R}
 \omega_{n_k,m}(e^{tQ_m})\\
 &\leq CM_m e^{-(t-t')R}.                            \tag{33.15}
\end{split}
\]
The same estimate holds for \(\omega_m\) by (33.11).
At fixed \(R\), weak-star convergence handles \(f_R(Q_m)\).
First let \(k\to\infty\), then \(R\to\infty\).
\end{proof}

\begin{corollary}[Loading the V21--V22 moment]
\label{cor:v16v33-load}
If the hypotheses of Theorem~\ref{thm:v16v22-local} hold uniformly
for each fixed region \(m\), then (33.9) holds for some local
Newton exponent \(t_m>0\).  Every projective subsequential state
therefore retains that exponential moment and the controlled
unbounded local Newton observables in (33.12).
\end{corollary}

\begin{proof}
Use (22.13) for (33.9), then
Theorem~\ref{thm:v16v33-moment}.
\end{proof}

\subsection{Identification versus compactness}

Theorem~\ref{thm:v16v33-projective} gives subsequential states.  A
unique cofinal limit requires an identification condition.

\begin{proposition}[Dense-core identification removes subsequences]
\label{prop:v16v33-unique}
For each \(m\), let \({\cal D}_m\subset{\cal A}_m\) be norm dense.
If
\[
 \lim_{n\to\infty}\omega_{n,m}(A)                   \tag{33.16}
\]
exists for every \(A\in{\cal D}_m\), then the whole family
\(\omega_{n,m}\) converges weak-star to a unique \(\omega_m\).
If ALC holds, these limits form the unique compatible local state.
\end{proposition}

\begin{proof}
States have norm one, so convergence on a norm-dense set extends
to every \(A\in{\cal A}_m\) by a three-term approximation.  Every
weak-star cluster point therefore has the same values on all of
\({\cal A}_m\), hence is unique.  Compactness then forces the full
sequence to converge.  Equation (33.6) gives compatibility.
\end{proof}

V32 can supply (33.16) for projected influence observables once
the common \(L^p\) density convergence holds.  It does not by
itself identify all Standard Model cylinder observables.

\begin{assessment}[Physical status]
\label{ass:v16v33-status}
V33 proves the missing projective compactness theorem for nested
local states, preserves OS positivity by the archived Gram
mechanism, and passes the new uniform Newton exponential moments to
controlled unbounded observables.  It separates existence of
subsequential compatible states from uniqueness on a dense core.

The physical programme has not proved ALC for finite-volume
SM--Einstein states or identified a dense set of all gauge,
fermion, Higgs, metric, and boundary observables.
\end{assessment}

\begin{decision}[V34 quantitative compatibility loader]
\label{dec:v16v33-next}
Derive ALC from a boundary-shell comparison.  If the marginal of a
larger-volume state on \({\cal A}_m\) has density \(r_{n;m,\ell}\)
relative to the smaller-volume state, bound
\(\delta_{n;m,\ell}\) by total variation, relative entropy, and
an \(L^p\) density defect.  Then identify the shell estimate needed
for the physical long-range Newton interaction.
\end{decision}

\begin{firewall}[Claim boundary after V33]
\label{fire:v16v33}
V33 proves a projective limit theorem conditional on ALC and local
moment cards.  It does not prove those physical inputs, uniqueness
of the full local state, reflection-positive Lorentzian
reconstruction, the limiting Einstein equation, or the complete
SM--Einstein continuum.
\end{firewall}

\section*{Version V33 append-only record}
\addcontentsline{toc}{section}{Version V33 append-only record}
The complete V32 source was retained before this continuation.  It
had \(300957\) bytes, \(9214\) lines, and SHA--256
\[
 \texttt{13A7201F618C149E6CA94105006528E747CE1FA3C20A553652B6885F4C4B5514}.
\]
V33 proved compatible projective local-state compactness and
extended convergence to local Newton observables using the
uniform exponential moments from V21--V22.
\section{Sixteenth-cycle V34: quantitative local compatibility from boundary-shell defects}
\label{sec:v16v34}

V33 assumes asymptotic local compatibility.  V34 gives several
checkable sufficient conditions for that row and then specializes
the density oscillation to a Newton boundary shell.  The resulting
quantity is again the positive Schur influence, so the far-tail
theorem of V19 is the correct loader.

\subsection{Classical marginal density bounds}

Let \(\mu,\nu\) be probability measures on one local measurable
space, with
\[
 d\nu=r\,d\mu,\qquad r\geq0,\qquad\int r\,d\mu=1.    \tag{34.1}
\]
Define the state-norm discrepancy
\[
 \Delta(\nu,\mu)
 :=
 \sup_{\|f\|_\infty\leq1}
 \left|\int f\,d\nu-\int f\,d\mu\right|.             \tag{34.2}
\]

\begin{theorem}[Density, entropy, and R\'enyi compatibility bounds]
\label{thm:v16v34-density}
For every \(p>1\),
\[
\begin{split}
 \Delta(\nu,\mu)
 &=\|r-1\|_{L^1(\mu)}\\
 &\leq\|r-1\|_{L^p(\mu)},                            \tag{34.3}\\
 \Delta(\nu,\mu)
 &\leq\sqrt{2D(\nu\Vert\mu)}
 \leq\sqrt{2D_p(\nu\Vert\mu)},                       \tag{34.4}
\end{split}
\]
where
\[
 D(\nu\Vert\mu)=\int r\log r\,d\mu,\qquad
 D_p(\nu\Vert\mu)
 ={1\over p-1}\log\int r^p\,d\mu.                   \tag{34.5}
\]
\end{theorem}

\begin{proof}
Duality between \(L^\infty\) and \(L^1\) gives the equality in
(34.3); H\"older on the probability space gives its inequality.
Pinsker's inequality gives the first bound in (34.4).

For the second, regard \(\nu\) as the probability measure
\(r\,d\mu\).  Jensen's inequality for the exponential gives
\[
\begin{split}
 D_p(\nu\Vert\mu)
 &={1\over p-1}
 \log{\mathbb E}_\nu e^{(p-1)\log r}\\
 &\geq{\mathbb E}_\nu\log r
 =D(\nu\Vert\mu).                                    \tag{34.6}
\end{split}
\]
\end{proof}

\begin{corollary}[Marginal density loader for ALC]
\label{cor:v16v34-ALC}
Let \(\nu_{n;m,\ell}\) be the marginal on region \(m\) of the
larger-region state \(\omega_{n,\ell}\), and let \(\mu_{n,m}\)
represent \(\omega_{n,m}\).  Any one of
\[
\begin{gathered}
 \|r_{n;m,\ell}-1\|_1\to0,\qquad
 \|r_{n;m,\ell}-1\|_p\to0,\\
 D(\nu_{n;m,\ell}\Vert\mu_{n,m})\to0,\qquad
 D_p(\nu_{n;m,\ell}\Vert\mu_{n,m})\to0              \tag{34.7}
\end{gathered}
\]
implies \(\delta_{n;m,\ell}\to0\) in (33.2).
\end{corollary}

\begin{proof}
The state discrepancy on bounded local multiplication observables
is (34.2).  Apply Theorem~\ref{thm:v16v34-density}.
\end{proof}

\subsection{Finite-cutoff quantum version}

\begin{proposition}[Quantum relative-entropy loader]
\label{prop:v16v34-quantum}
For finite-cutoff density matrices \(\rho,\sigma\),
\[
 \sup_{\|A\|\leq1}
 |\operatorname{Tr}[(\rho-\sigma)A]|
 =\|\rho-\sigma\|_1
 \leq\sqrt{2D(\rho\Vert\sigma)}.                    \tag{34.8}
\]
Thus vanishing relative entropy of the larger-volume marginal
relative to the smaller-volume state implies ALC also for the
finite gauge--fermion matrix algebra.
\end{proposition}

\begin{proof}
Trace duality gives the equality.  Choose the two-outcome
measurement defined by the positive spectral projection of
\(\rho-\sigma\).  Its classical total variation attains the trace
distance, while data processing bounds the measured classical
relative entropy by \(D(\rho\Vert\sigma)\).  Classical Pinsker then
gives (34.8).
\end{proof}

\subsection{Oscillation of a shell potential}

\begin{lemma}[Bounded potential oscillation]
\label{lem:v16v34-oscillation}
Let
\[
 d\nu={e^{-\Delta}\over\int e^{-\Delta}\,d\mu}\,d\mu
                                                               \tag{34.9}
\]
and suppose
\[
 \mathop{\rm ess\,sup}\Delta
 -\mathop{\rm ess\,inf}\Delta\leq\varepsilon.        \tag{34.10}
\]
Then
\[
 \Delta(\nu,\mu)\leq e^\varepsilon-1.               \tag{34.11}
\]
\end{lemma}

\begin{proof}
If \(a\leq\Delta\leq a+\varepsilon\), then
\[
 e^{-\varepsilon}
 \leq {e^{-\Delta}\over\int e^{-\Delta}d\mu}
 \leq e^\varepsilon.                                \tag{34.12}
\]
Hence \(|r-1|\leq e^\varepsilon-1\); integrate and use
(34.3).
\end{proof}

\begin{lemma}[Good-set compatibility]
\label{lem:v16v34-good}
Let \(G\) be an event with
\[
 \mu(G^c)\leq\eta,\qquad\nu(G^c)\leq\eta.            \tag{34.13}
\]
If the normalized restrictions \(\mu_G,\nu_G\) satisfy
\[
 \Delta(\nu_G,\mu_G)\leq d,                          \tag{34.14}
\]
then
\[
 \Delta(\nu,\mu)\leq d+4\eta.                       \tag{34.15}
\]
\end{lemma}

\begin{proof}
For \(\|f\|_\infty\leq1\), split both integrals over \(G\) and
\(G^c\).  The two bad pieces contribute at most \(2\eta\).
Writing the good pieces as their masses times normalized
expectations contributes at most \(d\) plus
\(|\mu(G)-\nu(G)|\leq2\eta\).  Take the supremum.
\end{proof}

\subsection{Newton shell oscillation}

Let \(u\) be the local interior source and \(v\) a shell source.
With the block notation of V15, put
\[
 Q_K(u)=\langle u,Au\rangle,\qquad
 J_{\rm sh}(v)=\langle Bv,A^+Bv\rangle.              \tag{34.16}
\]
The shell changes the local Newton exponent by
\[
 \Delta_{\rm sh}(u;v)
 =-2s\langle u,Bv\rangle.                            \tag{34.17}
\]

\begin{proposition}[Truncated Newton shell oscillation]
\label{prop:v16v34-shell}
On the local good set
\[
 G_R:=\{u:Q_K(u)\leq R\},                            \tag{34.18}
\]
one has
\[
 \operatorname{osc}_{G_R}\Delta_{\rm sh}(\,\cdot\,;v)
 \leq4s\sqrt{R\,J_{\rm sh}(v)}.                     \tag{34.19}
\]
\end{proposition}

\begin{proof}
The \(A\)-form Cauchy--Schwarz inequality gives
\[
 |\langle u,Bv\rangle|
 =
 |\langle A^{1/2}u,A^{+1/2}Bv\rangle|
 \leq\sqrt{Q_K(u)J_{\rm sh}(v)}.                    \tag{34.20}
\]
Thus \(|\Delta_{\rm sh}(u;v)|\leq
2s\sqrt{RJ_{\rm sh}(v)}\) on \(G_R\).  The range of a function
bounded in absolute value by \(a\) has length at most \(2a\),
which proves (34.19).
\end{proof}

\begin{corollary}[Deterministic shell compatibility on good sets]
\label{cor:v16v34-shell-compatibility}
If
\[
 J_{\rm sh}(v)\leq\zeta,                             \tag{34.21}
\]
then the normalized local measures restricted to \(G_R\), with and
without that shell coupling, have state discrepancy at most
\[
 \exp\{4s\sqrt{R\zeta}\}-1.                         \tag{34.22}
\]
\end{corollary}

\begin{proof}
Apply Lemma~\ref{lem:v16v34-oscillation} with (34.19).
\end{proof}

\subsection{Probabilistic shell card}

\begin{theorem}[Newton shell compatibility criterion]
\label{thm:v16v34-shell-card}
Suppose along a shell exhaustion:
\begin{enumerate}
\item for some \(t>0\),
\[
 \sup_n{\mathbb E}e^{tQ_{K,n}}\leq M_K;             \tag{34.23}
\]
\item for some \(c>0\),
\[
 \sup_n
 \left({\mathbb E}e^{cJ_{{\rm sh},n}}-1\right)
 \longrightarrow0.                                  \tag{34.24}
\]
\end{enumerate}
Then one can choose \(R\to\infty\) and
\(\zeta\downarrow0\) along the exhaustion so that the good-set
shell discrepancy (34.22) tends to zero and the bad-set
probabilities tend to zero.  Hence the shell contribution satisfies
ALC through Lemma~\ref{lem:v16v34-good}, provided the compared local
marginals differ only by the shell weight (34.17).
\end{theorem}

\begin{proof}
Exponential Markov gives
\[
 {\mathbb P}(Q_{K,n}>R)\leq M_Ke^{-tR}.              \tag{34.25}
\]
Since \(e^{cx}\geq1+cx\) for \(x\geq0\), (34.24) gives
\[
 \sup_n{\mathbb E}J_{{\rm sh},n}\to0.               \tag{34.26}
\]
Markov then makes
\({\mathbb P}(J_{{\rm sh},n}>\zeta)\) small whenever
\(\sup_n{\mathbb E}J_{{\rm sh},n}=o(\zeta)\).
Choose \(\zeta\) first so that this holds, and choose \(R\) growing
slowly enough that
\[
 R\zeta\to0,\qquad e^{-tR}\to0.                     \tag{34.27}
\]
Then (34.22) and both bad probabilities vanish.  Apply
Lemma~\ref{lem:v16v34-good}.
\end{proof}

For the neutral far Newton shell, (34.24) is precisely the
conclusion of the projected-density theorem V19 at exponent \(c\).
Thus V19 plus the local moment from V21--V22 supplies ALC for that
shell under the stated hypothesis that the marginals differ only
by (34.17).

\begin{assessment}[Physical status]
\label{ass:v16v34-status}
V34 turns ALC into quantitative density, entropy, and shell
criteria.  It proves that uniform local Newton moments plus
vanishing projected far-influence exponential moments make a
neutral boundary shell invisible to bounded local observables.

The physical larger- and smaller-volume SM--Einstein marginals may
also differ through gauge fixing, fermion determinants,
counterterms, and metric boundary data.  Those differences have
not been reduced to the single shell weight (34.17).
\end{assessment}

\begin{decision}[V35 companion audit and multi-owner shell
telescoping]
\label{dec:v16v34-next}
Perform the compulsory companion audit.  Then decompose the full
larger-versus-smaller local density ratio into separately owned
Newton, gauge, Higgs, fermion, counterterm, and boundary factors.
Prove a telescoping state-distance bound so each owner can supply
its own local compatibility debit, without assuming their
independence.
\end{decision}

\begin{firewall}[Claim boundary after V34]
\label{fire:v16v34}
V34 proves a Newton shell compatibility implication, not ALC for
the full coupled theory.  It does not control the other density
owners, prove uniqueness or reflection positivity of the physical
local state, recover the Einstein equation, or complete the
continuum.
\end{firewall}

\section*{Version V34 append-only record}
\addcontentsline{toc}{section}{Version V34 append-only record}
The complete V33 source was retained before this continuation.  It
had \(310009\) bytes, \(9479\) lines, and SHA--256
\[
 \texttt{94643CEC64ADF55988F9EC45B2963C783B5C778710B4710B9CF5B154BCA97F07}.
\]
V34 derived quantitative local compatibility from density,
entropy, and boundary-shell defects and connected the neutral
Newton shell to the V19 far-influence theorem.

\section{V35: Companion audit and the multi-owner local-compatibility ledger}
\label{sec:v16v35}

\subsection{Compulsory five-version companion audit}

The newest Yang--Mills companion inspected before this continuation was
\texttt{Renewal\_Yang\_Mills\_Mass\_Gap\_Research\_Notes\_V48.tex}; it
had (419467) bytes, (11461) lines, and SHA--256
\[
 \texttt{740675B97D48DF77BB6F24D55D1C2E13E33404DB639113BBFA4DD6326D7801FE}.
                                                               \tag{35.1}
\]
Its new result constructs sixteen exact opposite power-three charts at
(n_4=3) and proves both direct response signs, with a common rational
floor, on one inversion-paired pair of full transverse quarter cells.
It deliberately leaves the other inversion pairs and the physical
transfer and continuum limits open.  The newest Navier--Stokes
companion remained
\texttt{Renewal\_NS\_Stability\_Research\_Notes\_V26.tex}; it had
(278283) bytes, (5319) lines, and SHA--256
\[
 \texttt{82C887F8C2C69E4F28FAD7C3DC8DE5B9099CB5A4BF431EBA1FFF3E91935978AD}.
                                                               \tag{35.2}
\]
Its Oseen rank-one derivative estimate does not supply a missing
Standard Model--Einstein compatibility estimate.

No theorem is imported from either companion.  The useful lesson from
the Yang--Mills audit is methodological: the successful local
construction is owned chart by chart and does not infer the uncovered
orientations from an unproved symmetry.  The same discipline is needed
here.  The Newton, gauge, Higgs, fermion, counterterm, and metric
boundary changes will be charged separately.  Dependence between them
will not be replaced by an independence assumption.

\subsection{State norm and one-owner loaders}

On a measurable local configuration space ((X,\mathcal F)), write
\[
 \|\nu-\mu\|_{\mathrm{st}}
 :=\sup_{\|f\|_\infty\leq1}|\nu(f)-\mu(f)|.          \tag{35.3}
\]
For probability measures this is the (L^1) norm of the density
difference whenever both measures are dominated by a common measure.
This convention is twice the probabilists' total variation distance.

\begin{definition}[Admissible owner path]
\label{def:v16v35-path}
An admissible (M)-owner path from a smaller-volume local state to a
larger-volume local marginal is a sequence of probability measures
\[
 \mu^{(0)},\mu^{(1)},\ldots,\mu^{(M)}               \tag{35.4}
\]
on the same local measurable space.  The passage
(mu^{(j-1)}\mapsto\mu^{(j)}) changes only the density factor assigned
to owner (j).  A number (d_j\geq0) is a certified debit when
\[
 \|\mu^{(j)}-\mu^{(j-1)}\|_{\mathrm{st}}\leq d_j.   \tag{35.5}
\]
The definition asserts no probabilistic independence.
\end{definition}

\begin{lemma}[Four certified one-owner debits]
\label{lem:v16v35-loaders}
Let \(\alpha,\beta\) be consecutive states in an admissible owner path.
Each of the following hypotheses certifies the indicated debit:
\begin{enumerate}
\item if (r=d\beta/d\alpha), then
\[
 d=\|r-1\|_{L^1(\alpha)};                           \tag{35.6}
\]
\item if (D(\beta\Vert\alpha)<\infty), then
\[
 d=\sqrt{2D(\beta\Vert\alpha)};                    \tag{35.7}
\]
\item if
(d\beta=e^{-\Delta}d\alpha/\alpha(e^{-\Delta}))
and ({\rm ess\,osc}_\alpha\Delta\leq\varepsilon), then
\[
 d=e^\varepsilon-1;                                 \tag{35.8}
\]
\item if an event (G) obeys
(alpha(G^c),\beta(G^c)\leq\eta) and the normalized restrictions
obey (|\beta_G-\alpha_G\|_{\mathrm{st}}\leq q), then
\[
 d=q+4\eta.                                         \tag{35.9}
\]
\end{enumerate}
The same assertions hold for finite-cutoff density matrices, with
trace norm in (35.3), (L^1) replaced by trace norm, and (35.7)
interpreted as quantum relative entropy.
\end{lemma}

\begin{proof}
For (35.6), duality between (L^1(\alpha)) and
(L^\infty(\alpha)) gives
\[
 \sup_{\|f\|_\infty\leq1}|\alpha((r-1)f)|
 =\alpha(|r-1|).                                    \tag{35.10}
\]
Pinsker gives (35.7).  If (a\leq\Delta\leq
a+\varepsilon), division by the normalizer puts the density in
([e^{-\varepsilon},e^\varepsilon]); this proves (35.8).
Splitting both expectations over (G) and (G^c), and then
normalizing the two good pieces, gives (35.9), exactly as in
Lemma~\ref{lem:v16v34-good}.  Trace duality and quantum Pinsker prove
the matrix statement.
\end{proof}

\subsection{Exact telescoping without independence}

\begin{theorem}[Multi-owner telescoping theorem]
\label{thm:v16v35-telescope}
Every admissible (M)-owner path satisfies
\[
 \|\mu^{(M)}-\mu^{(0)}\|_{\mathrm{st}}
 \leq\sum_{j=1}^{M}
 \|\mu^{(j)}-\mu^{(j-1)}\|_{\mathrm{st}}
 \leq\sum_{j=1}^{M}d_j.                            \tag{35.11}
\]
Consequently, for a triangular family of such paths indexed by a
cutoff (n), local compatibility follows from
\[
 \sum_{j=1}^{M}d_{j,n}\longrightarrow0.             \tag{35.12}
\]
For a fixed number of owners it suffices that every (d_{j,n}	o0).
No independence, orthogonality, or factorization of expectations is
required.
\end{theorem}

\begin{proof}
For every (f) in the unit ball of (L^\infty),
\[
 \mu^{(M)}(f)-\mu^{(0)}(f)
 =\sum_{j=1}^{M}
   \bigl(\mu^{(j)}(f)-\mu^{(j-1)}(f)\bigr).         \tag{35.13}
\]
Apply the scalar triangle inequality, then take the supremum over
(f).  Insert (35.5).  The last two conclusions follow immediately.
The proof never compares joint events belonging to distinct owners.
\end{proof}

\begin{proposition}[Canonical path for positive density owners]
\label{prop:v16v35-product}
Let (mu^{(0)}) be a probability measure and let
(w_1,\ldots,w_M) be strictly positive measurable weights such that
\[
 0<Z_j:=\mu^{(0)}\!\left(\prod_{i=1}^{j}w_i\right)<\infty
 \quad(1\leq j\leq M).                              \tag{35.14}
\]
Define
\[
 d\mu^{(j)}={\prod_{i=1}^{j}w_i\over Z_j},d\mu^{(0)}.
                                                               \tag{35.15}
\]
Then the exact one-step Radon--Nikodym derivative is
\[
 {d\mu^{(j)}\over d\mu^{(j-1)}}
 ={w_j\over\mu^{(j-1)}(w_j)}.                      \tag{35.16}
\]
In particular, an oscillation, entropy, or good-set estimate for
(w_j) must be evaluated under the already tilted state
(mu^{(j-1)}), but Theorem~\ref{thm:v16v35-telescope} then adds its
debit without any independence assumption.
\end{proposition}

\begin{proof}
Since
\[
 \mu^{(j-1)}(w_j)
 ={1\over Z_{j-1}}
 \mu^{(0)}\!\left(\prod_{i=1}^{j}w_i\right)
 ={Z_j\over Z_{j-1}},                               \tag{35.17}
\]
dividing (35.15) by the corresponding formula at (j-1) proves
(35.16).  Normalization is automatic.  The final statement follows
from Lemma~\ref{lem:v16v35-loaders}.
\end{proof}

\begin{remark}[Order and the cost of a proof]
\label{rem:v16v35-order}
The terminal density in (35.15) is independent of the enumeration of
commuting positive weights, whereas the intermediate laws and the
available bounds (d_j) need not be.  Thus one may choose an order
that exposes the strongest conditional estimates.  This changes the
proof ledger, not the terminal state.
\end{remark}

\subsection{The Standard Model--Einstein owner ledger}

Let (mu_{n,K}^{\rm small}) be the state built directly on a local
region (K), and let (mu_{n,K}^{\rm large}) be the (K)-marginal
of a larger finite-cutoff state.  Suppose their positive density ratio
has been justified in the form
\[
 {d\mu_{n,K}^{\rm large}\over d\mu_{n,K}^{\rm small}}
 ={1\over Z_{n,K}}
 W_{N,n}W_{G,n}W_{H,n}W_{F,n}W_{C,n}W_{B,n},        \tag{35.18}
\]
where the letters denote, respectively, Newton, gauge, Higgs,
fermion, counterterm, and residual metric-boundary owners.  Formula
(35.18) is an assumption to be proved in the physical construction;
in particular (W_{F,n}) must be a positive density rather than an
uncontrolled determinant phase.

\begin{corollary}[Six-owner ALC ledger]
\label{cor:v16v35-ledger}
Assume (35.18) and use Proposition~\ref{prop:v16v35-product} in the
displayed order.  If certified conditional debits satisfy
\[
 d_{N,n}+d_{G,n}+d_{H,n}+d_{F,n}+d_{C,n}+d_{B,n}
 \longrightarrow0,                                  \tag{35.19}
\]
then
\[
 \|\mu_{n,K}^{\rm large}-\mu_{n,K}^{\rm small}|_{\mathrm{st}}
 \longrightarrow0.                                  \tag{35.20}
\]
The Newton debit may be supplied by V34 from the V19 far-influence
moment.  Each of the remaining five terms retains a distinct proof
obligation.
\end{corollary}

\begin{proof}
The six successive normalized tilts form an admissible path by
(35.16).  Apply Theorem~\ref{thm:v16v35-telescope} and (35.19).
\end{proof}

\begin{proposition}[Refinement into shell subowners]
\label{prop:v16v35-shell-refinement}
If owner (a) is itself a product of (L_n) positive shell weights,
replace its one step by (L_n) canonical steps with certified debits
(d_{a,q,n}).  Then its total contribution is bounded by
\[
 d_{a,n}^{\rm tot}\leq\sum_{q=1}^{L_n}d_{a,q,n}.     \tag{35.21}
\]
Hence pointwise convergence (d_{a,q,n}	o0) at fixed (q) is
insufficient when (L_n\to\infty); the required hypothesis is the
summable estimate
\[
 \sum_{q=1}^{L_n}d_{a,q,n}\longrightarrow0.         \tag{35.22}
\]
\end{proposition}

\begin{proof}
Insert the shell intermediates into (35.13).  The number of terms may
depend on (n), so only their sum controls the endpoint distance.
\end{proof}

\begin{assessment}[What V35 closes]
\label{ass:v16v35-status}
V35 removes a logical gap in the V34 shell criterion: a neutral Newton
shell may now be combined with other changes by a rigorous path norm,
without pretending that the physical sectors are independent.  It
also identifies the uniform summability needed when an owner is split
into a growing number of shell pieces.

It does not prove the positive factorization (35.18), any of the five
non-Newton debits, or a countably infinite product construction.  A
chiral determinant phase cannot be inserted into this probability
ledger until positivity or a controlled phase reconstruction has been
proved.
\end{assessment}

\begin{decision}[V36 infinite-owner completion]
\label{dec:v16v35-next}
Prove the countably many owner version of the ledger: total-variation
convergence, an explicit tail bound, and a triangular-array passage
from finite prefix convergence plus uniform tail summability.  This is
needed before physical shells whose number grows with the cutoff can
be assembled.
\end{decision}

\begin{firewall}[Claim boundary after V35]
\label{fire:v16v35}
The telescoping theorem is an assembly result.  It does not create the
owner estimates it consumes, establish a positive fermionic density,
prove the existence of the limiting interacting state, recover the
Einstein equation, or complete the Standard Model--Einstein continuum.
\end{firewall}

\section*{Version V35 append-only record}
\addcontentsline{toc}{section}{Version V35 append-only record}
The complete V34 source was retained before this continuation.  It
had (319791) bytes, (9796) lines, and SHA--256
\[
 \texttt{A3DF182994E2F6AFE710594B03765DC50F1A8337FB2004D50C797ABCE8F175F0}.
\]
V35 performed the mandatory companion audit and replaced a tacit
multi-sector leap by an exact owner-by-owner state-distance ledger.

\section{V36: Countable owner assembly and uniform shell tails}
\label{sec:v16v36}

V35 treats a finite owner path.  A physical cutoff removal can create
more boundary blocks, counterterm cells, or determinant factors at
each stage.  Pointwise control of any fixed factor does not control
this growth.  This continuation gives two sufficient mechanisms:
summable step distances along one enumeration, and a stronger
finite-set estimate which makes the limit independent of the
enumeration.

\subsection{Completeness of the state metric}

\begin{lemma}[Complete probability cone]
\label{lem:v16v36-complete}
Let \((X,\mathcal F)\) be a measurable space.  The probability measures
on it are complete for
\[
 \|\nu-\mu\|_{\mathrm{st}}
 =\sup_{\|f\|_\infty\leq1}|\nu(f)-\mu(f)|.            \tag{36.1}
\]
Thus every Cauchy sequence in this norm has a probability-measure
limit in the same norm.
\end{lemma}

\begin{proof}
For a Cauchy sequence \((\mu_j)\), set
\[
 \lambda:=\sum_{j=1}^{\infty}2^{-j}\mu_j.            \tag{36.2}
\]
This is a probability measure and every \(\mu_j\) has a density
\(f_j=d\mu_j/d\lambda\).  Density duality gives
\[
 \|f_j-f_k\|_{L^1(\lambda)}
 =\|\mu_j-\mu_k\|_{\mathrm{st}}.                     \tag{36.3}
\]
Completeness of \(L^1(\lambda)\) supplies \(f_j\to f\) in \(L^1\).
After passage to an almost-everywhere convergent subsequence,
\(f\geq0\), and
\[
 \int f\,d\lambda=\lim_j\int f_j\,d\lambda=1.        \tag{36.4}
\]
The measure \(d\mu=f\,d\lambda\) is therefore a probability measure,
and (36.3) gives convergence in (36.1).
\end{proof}

\begin{theorem}[Summable countable owner path]
\label{thm:v16v36-countable}
Let \((\mu^{(j)})_{j\geq0}\) be probability measures on one local
measurable space and suppose
\[
 \|\mu^{(j)}-\mu^{(j-1)}\|_{\mathrm{st}}\leq d_j,
 \qquad
 \sum_{j=1}^{\infty}d_j<\infty.                     \tag{36.5}
\]
Then there is a unique probability measure \(\mu^{(\infty)}\) such
that
\[
 \mu^{(N)}\longrightarrow\mu^{(\infty)}
 \quad\hbox{in state norm},                          \tag{36.6}
\]
and the quantitative tail estimate is
\[
 \|\mu^{(\infty)}-\mu^{(N)}\|_{\mathrm{st}}
 \leq\sum_{j>N}d_j.                                 \tag{36.7}
\]
\end{theorem}

\begin{proof}
For \(P>N\), the finite telescoping theorem gives
\[
 \|\mu^{(P)}-\mu^{(N)}\|_{\mathrm{st}}
 \leq\sum_{j=N+1}^{P}d_j.                           \tag{36.8}
\]
The right side tends to zero uniformly in \(P>N\) as \(N\to\infty\).
Lemma~\ref{lem:v16v36-complete} yields (36.6).  Letting \(P\to\infty\)
in (36.8) proves (36.7).
\end{proof}

\subsection{An order-independent owner criterion}

Let \(I\) be a countable owner set and let \(\mathfrak F(I)\) be the
directed set of its finite subsets.  Think of \(\mu_F\) as the
normalized local state obtained after inserting precisely the owners
in \(F\).  When positive weights \(w_i\) exist, the canonical example
is
\[
 d\mu_F={\prod_{i\in F}w_i\over
 \mu_\varnothing(\prod_{i\in F}w_i)}\,d\mu_\varnothing.
                                                               \tag{36.9}
\]

\begin{theorem}[Uniform insertion debit and exhaustion independence]
\label{thm:v16v36-net}
Suppose that probability measures
\((\mu_F)_{F\in\mathfrak F(I)}\) obey
\[
 \|\mu_{F\cup\{i\}}-\mu_F\|_{\mathrm{st}}\leq a_i
 \quad
 (F\in\mathfrak F(I),\ i\notin F),                  \tag{36.10}
\]
where
\[
 a_i\geq0,\qquad \sum_{i\in I}a_i<\infty.           \tag{36.11}
\]
Then the finite-set net has a unique state-norm limit
\(\mu_I\).  More precisely, every finite \(F\) satisfies
\[
 \|\mu_I-\mu_F\|_{\mathrm{st}}
 \leq\sum_{i\notin F}a_i.                           \tag{36.12}
\]
Consequently every increasing exhaustion
\(F_1\subset F_2\subset\cdots\), \(\bigcup_NF_N=I\), converges to the
same \(\mu_I\), independently of the insertion order.
\end{theorem}

\begin{proof}
If \(F\subset G\), enumerate \(G\setminus F\) and insert its elements
one at a time.  Repeated use of (36.10) gives
\[
 \|\mu_G-\mu_F\|_{\mathrm{st}}
 \leq\sum_{i\in G\setminus F}a_i.                   \tag{36.13}
\]
For arbitrary \(F,G\), pass through \(F\cup G\) to obtain
\[
 \|\mu_F-\mu_G\|_{\mathrm{st}}
 \leq
 \sum_{i\in G\setminus F}a_i+
 \sum_{i\in F\setminus G}a_i.                       \tag{36.14}
\]
Given \(\varepsilon>0\), choose a finite \(F_0\subset I\) with
\(\sum_{i\notin F_0}a_i<\varepsilon\).  Equation (36.14) makes the
tail of the net over sets containing \(F_0\) Cauchy.  Along any one
exhaustion, Lemma~\ref{lem:v16v36-complete} gives a probability
limit.  Equation (36.14) shows that all exhaustions have that same
limit and that the whole net converges to it.  Finally take an
exhaustion \(G_N\supset F\) and let \(N\to\infty\) in (36.13) to get
(36.12).
\end{proof}

\begin{remark}[Why uniformity in the background is essential]
\label{rem:v16v36-uniform}
A bound only for insertion of \(i\) into one preferred predecessor
does not imply order independence: earlier tilts can change the law
on which the \(i\)-th oscillation or entropy is evaluated.  Condition
(36.10) controls the insertion uniformly over every finite background
of other owners.  It is stronger than (36.5), and exactly this
strength pays for arbitrary reorderings.
\end{remark}

\begin{corollary}[Oscillation criterion for a countable positive
product]
\label{cor:v16v36-osc}
In the positive product setting (36.9), suppose
\[
 {\rm ess\,osc}_{\mu_F}(-\log w_i)\leq\varepsilon_i
 \quad(F\in\mathfrak F(I),\ i\notin F)              \tag{36.15}
\]
and
\[
 \sum_{i\in I}(e^{\varepsilon_i}-1)<\infty.         \tag{36.16}
\]
Then the finite products have a state-norm limit independent of their
enumeration, with
\[
 \|\mu_I-\mu_F\|_{\mathrm{st}}
 \leq\sum_{i\notin F}(e^{\varepsilon_i}-1).          \tag{36.17}
\]
\end{corollary}

\begin{proof}
Proposition~\ref{prop:v16v35-product} identifies insertion of \(w_i\)
as a single normalized tilt of \(\mu_F\).  The oscillation loader
(35.8) gives (36.10) with
\(a_i=e^{\varepsilon_i}-1\).  Apply
Theorem~\ref{thm:v16v36-net}.
\end{proof}

\subsection{The cutoff triangular array}

\begin{theorem}[Finite-prefix convergence plus uniform owner tails]
\label{thm:v16v36-triangular}
For every cutoff \(n\), let
\[
 \mu_n^{(0)},\mu_n^{(1)},\ldots,\mu_n^{(L_n)}
                                                               \tag{36.18}
\]
be a finite admissible path, extended constantly for \(j>L_n\).
Write \(\nu_n=\mu_n^{(L_n)}\), and suppose there are debits
\(d_{j,n}\), also set to zero for \(j>L_n\), such that
\[
 \|\mu_n^{(j)}-\mu_n^{(j-1)}\|_{\mathrm{st}}
 \leq d_{j,n}.                                      \tag{36.19}
\]
Assume:
\begin{enumerate}
\item for every fixed \(J\), there is a probability measure
\(\mu^{(J)}\) with
\[
 \|\mu_n^{(J)}-\mu^{(J)}\|_{\mathrm{st}}\longrightarrow0;
                                                               \tag{36.20}
\]
\item the owner tails are uniformly summable:
\[
 \lim_{J\to\infty}\sup_n\sum_{j>J}d_{j,n}=0.        \tag{36.21}
\]
\end{enumerate}
Then \((\mu^{(J)})\) has a state-norm limit \(\mu^{(\infty)}\) and
\[
 \|\nu_n-\mu^{(\infty)}\|_{\mathrm{st}}\longrightarrow0.       \tag{36.22}
\]
More quantitatively,
\[
 \limsup_{n\to\infty}
 \|\nu_n-\mu^{(\infty)}\|_{\mathrm{st}}
 \leq
 2\limsup_{n\to\infty}\sum_{j>J}d_{j,n}
                                                               \tag{36.23}
\]
for every fixed \(J\).
\end{theorem}

\begin{proof}
For \(P>J\), finite telescoping and (36.19) give
\[
 \|\mu_n^{(P)}-\mu_n^{(J)}\|_{\mathrm{st}}
 \leq\sum_{j=J+1}^{P}d_{j,n}.                       \tag{36.24}
\]
Taking \(n\to\infty\) and using (36.20) at \(J,P\) yields
\[
 \|\mu^{(P)}-\mu^{(J)}\|_{\mathrm{st}}
 \leq\liminf_{n\to\infty}\sum_{j>J}d_{j,n}.         \tag{36.25}
\]
Condition (36.21) and Lemma~\ref{lem:v16v36-complete} therefore give
a limit \(\mu^{(\infty)}\), with
\[
 \|\mu^{(\infty)}-\mu^{(J)}\|_{\mathrm{st}}
 \leq\liminf_{n\to\infty}\sum_{j>J}d_{j,n}.         \tag{36.26}
\]
For the terminal state, telescoping gives
\[
 \|\nu_n-\mu_n^{(J)}\|_{\mathrm{st}}
 \leq\sum_{j>J}d_{j,n}.                             \tag{36.27}
\]
Insert \(\mu_n^{(J)}\) and \(\mu^{(J)}\) between
\(\nu_n\) and \(\mu^{(\infty)}\), take the upper limit in \(n\), and
use (36.20), (36.26), and
\(\liminf\leq\limsup\).  This proves (36.23).  Sending \(J\to\infty\)
proves (36.22).
\end{proof}

\begin{corollary}[Closed state inequalities]
\label{cor:v16v36-closed}
Let \(\mu_N\to\mu\) in state norm.  Any inequality
\[
 \mu_N(A)\geq0                                      \tag{36.28}
\]
which holds for every \(N\) and for a fixed bounded measurable
observable \(A\) passes to \(\mu(A)\geq0\).  In particular,
normalization, positivity on bounded nonnegative observables, and
reflection-positivity inequalities
\(\mu_N(\overline F\,\Theta F)\geq0\) for fixed bounded cylinder
functions \(F\) pass to the limit whenever they hold at every finite
stage.
\end{corollary}

\begin{proof}
State-norm convergence implies
\[
 |\mu_N(A)-\mu(A)|
 \leq\|A\|_\infty\|\mu_N-\mu\|_{\mathrm{st}}\to0.    \tag{36.29}
\]
Take the limit in (36.28).  The listed cases are special choices of
\(A\); normalization also follows directly from \(A=1\).
\end{proof}

\subsection{Application ledger and the next physical obstruction}

For the six sectors of (35.18), allow owner \(a\) to contain shell
subowners \(q\in I_a\).  A sufficient assembly hypothesis is
\[
 \sum_{a\in\{N,G,H,F,C,B\}}\ \sum_{q\in I_a} a_{a,q}<\infty,
                                                               \tag{36.30}
\]
together with uniform insertion estimates of the form (36.10).
For cutoff convergence, replace mere finiteness by the vanishing
tail condition
\[
 \lim_{J\to\infty}\sup_n
 \sum_{\substack{(a,q)\ {\rm beyond}\\{\rm the\ first}\ J}}
 d_{a,q,n}=0.                                       \tag{36.31}
\]
The Newton results through V34 provide a candidate far-shell debit.
They do not yet prove its uniformity under arbitrary gauge, Higgs,
fermion, counterterm, and boundary tilts, as demanded by (36.10).
When such uniformity is unavailable, Theorem~\ref{thm:v16v36-triangular}
permits a fixed physical insertion order and asks only for (36.21).

\begin{assessment}[What V36 closes]
\label{ass:v16v36-status}
V36 proves that a growing or countable collection of physical density
owners can be assembled without losing normalization or local state
convergence, provided their state-distance debits have a uniform
summable tail.  It distinguishes the hypotheses for one fixed order
from the stronger hypothesis required for order independence.

The remaining bottleneck is now quantitative rather than formal:
one needs per-owner shell estimates which remain summable as the
cutoff exposes more cells.  Finite-range local action terms are the
first tractable case because remote cells should cancel exactly after
conditional marginalization; boundary-straddling cells should be the
only counterterm owners left in the debit.
\end{assessment}

\begin{decision}[V37 finite-range cancellation and boundary
counting]
\label{dec:v16v36-next}
Prove an exact cancellation theorem for action factors supported
entirely outside the observed region, formulate the residual
boundary-hypergraph owner set for finite-range gauge, Higgs, and local
counterterm interactions, and derive the geometric count that a
summable per-cell debit must beat.  Keep determinant and Newton
nonlocalities outside that finite-range theorem.
\end{decision}

\begin{firewall}[Claim boundary after V36]
\label{fire:v16v36}
Countable assembly does not supply (36.30) or (36.31).  Reflection
positivity passes to the limit only if it has already been proved at
every approximant.  No result here controls a determinant phase,
proves the physical continuum measure, derives the Einstein equation,
or establishes the full Standard Model--Einstein continuum.
\end{firewall}

\section*{Version V36 append-only record}
\addcontentsline{toc}{section}{Version V36 append-only record}
The complete V35 source was retained before this continuation.  It
had \(330891\) bytes, \(10087\) lines, and SHA--256
\[
 \texttt{A3C63373B287D23572A0A76E3E5B66C5CB527354727442C6B53F94F0A9BA3C21}.
\]
V36 completed the finite ledger at countable cardinality, proved an
exhaustion-independent insertion criterion, and supplied the
triangular-array theorem needed for growing cutoff owner families.

\section{V37: Conditional cancellation, Markov blankets, and
boundary-owner counting}
\label{sec:v16v37}

The proposed next step in V36 must be corrected before use.  Spatial
support outside an observed region does not by itself make a factor
invisible to the unconditional marginal: correlations can transmit
its tilt inward.  What cancels exactly is the part of a Gibbs density
which is constant as a function of the variables being conditioned.
This section proves the exact statement, isolates the boundary law
which still carries exterior influence, and counts the finite-range
owners on successive lattice shells.

\subsection{The obstruction to unconditional cancellation}

\begin{lemma}[Marginal produced by an exterior tilt]
\label{lem:v16v37-exterior-tilt}
Let \(\mu\) be a probability measure on
\((X\times Y,\mathcal F_X\otimes\mathcal F_Y)\), let \(h:Y\to(0,\infty)\)
be integrable, and set
\[
 d\nu(x,y)={h(y)\over\mu(h)}\,d\mu(x,y).             \tag{37.1}
\]
If \(\mu_X,\nu_X\) are the \(X\)-marginals, then
\[
 {d\nu_X\over d\mu_X}(x)
 ={\,\mathbb E_\mu[h(Y)\mid X=x]\,\over\mu(h)}
 \quad\text{for }\mu_X\text{-almost every }x.       \tag{37.2}
\]
Consequently \(\nu_X=\mu_X\) if and only if
\[
 \mathbb E_\mu[h(Y)\mid X]=\mu(h)
 \quad\mu\text{-almost surely}.                     \tag{37.3}
\]
\end{lemma}

\begin{proof}
For every bounded measurable \(f\) on \(X\), conditional expectation
gives
\[
 \nu_X(f)
 ={1\over\mu(h)}\mu(f(X)h(Y))
 ={1\over\mu(h)}
 \mu_X\!\left(f\,\mathbb E_\mu[h(Y)\mid X]\right).  \tag{37.4}
\]
This proves (37.2).  Equality of the marginals is equivalent to the
density in (37.2) being one almost everywhere, which is (37.3).
\end{proof}

\begin{remark}[The support-only shortcut is false]
\label{rem:v16v37-false}
Independence of \(X\) and \(Y\) implies (37.3), but a coupled Gibbs
state normally fails that independence.  Thus an exterior
counterterm or gauge plaquette cannot be erased from the local
marginal merely because its displayed variables lie outside the
observed set.  Its effect can enter through the distribution of the
boundary variables.
\end{remark}

\subsection{Exact cancellation in a finite Gibbs specification}

Let \(V\) be a finite site set, with a product reference measure
\(\lambda_V=\bigotimes_{v\in V}\lambda_v\).  Let \(\mathcal A\) be a
finite interaction hypergraph and let each measurable potential
\(\Phi_A\) depend only on \(x_A\).  Assume the partition functions
below are positive and finite.  The Gibbs measure is
\[
 d\mu_V(x)
 ={1\over Z_V}
 \exp\!\left\{-\sum_{A\in\mathcal A}\Phi_A(x_A)\right\}
 d\lambda_V(x).                                    \tag{37.5}
\]

\begin{theorem}[Exact conditional cancellation]
\label{thm:v16v37-conditional}
For \(D\subset V\), the conditional law of \(x_D\) given
\(x_{D^c}=y\) is, for almost every admissible \(y\),
\[
 \gamma_D(dx_D\mid y)
 ={1\over Z_D(y)}
 \exp\!\left\{-\sum_{\substack{A\in\mathcal A\\A\cap D\ne\varnothing}}
 \Phi_A(x_{A\cap D},y_{A\cap D^c})\right\}
 d\lambda_D(x_D).                                  \tag{37.6}
\]
Every potential with \(A\cap D=\varnothing\) cancels exactly from
this normalized kernel.
\end{theorem}

\begin{proof}
Split the exponent in (37.5) into
\[
 S_D(x_D,y)
 :=\sum_{A:A\cap D\ne\varnothing}
 \Phi_A(x_{A\cap D},y_{A\cap D^c})                 \tag{37.7}
\]
and
\[
 S_{D^c}^{\,0}(y)
 :=\sum_{A:A\cap D=\varnothing}\Phi_A(y_A).         \tag{37.8}
\]
For fixed \(y\), the factor \(e^{-S_{D^c}^{\,0}(y)}\) is constant in
\(x_D\).  It appears in both the unnormalized conditional density and
its integral over \(x_D\), and therefore cancels.  Normalizing
\(e^{-S_D(x_D,y)}d\lambda_D\) proves (37.6).
\end{proof}

\begin{definition}[Finite-range Markov blanket]
\label{def:v16v37-blanket}
On a metric site set, suppose every \(A\in\mathcal A\) has diameter at
most \(r\).  Define the exterior \(r\)-blanket of \(D\) by
\[
 \partial_r^+D
 :=\{v\in V\setminus D:\operatorname{dist}(v,D)\leq r\}.        \tag{37.9}
\]
\end{definition}

\begin{corollary}[Blanket sufficiency]
\label{cor:v16v37-blanket}
Under the range-\(r\) hypothesis, the kernel in (37.6) depends on the
exterior configuration only through \(y_{\partial_r^+D}\).  Hence
there is a Markov kernel
\[
 \kappa_D:X_{\partial_r^+D}\longrightarrow
 \mathcal P(X_D)                                    \tag{37.10}
\]
such that, if \(\rho_V^D\) is the blanket marginal of \(\mu_V\),
\[
 (\mu_V)_D=\rho_V^D\kappa_D.                        \tag{37.11}
\]
\end{corollary}

\begin{proof}
If \(A\cap D\ne\varnothing\) and \(v\in A\setminus D\), then
\(\operatorname{dist}(v,D)\leq\operatorname{diam}(A)\leq r\).
Thus every exterior argument in (37.6) lies in the blanket.  The
tower property of conditional expectation gives (37.11).
\end{proof}

\subsection{Comparison of two volumes}

\begin{lemma}[Kernel comparison inequality]
\label{lem:v16v37-kernel}
Let \(\rho,\sigma\) be probability measures on a boundary space and
let \(\kappa,\kappa'\) be Markov kernels from that space to a local
space.  Then
\[
 \|\rho\kappa-\sigma\kappa'\|_{\mathrm{st}}
 \leq
 \|\rho-\sigma\|_{\mathrm{st}}
 +\sup_y\|\kappa(y,\cdot)-\kappa'(y,\cdot)\|_{\mathrm{st}}.
                                                               \tag{37.12}
\]
In particular every Markov kernel contracts state norm:
\[
 \|\rho\kappa-\sigma\kappa\|_{\mathrm{st}}
 \leq\|\rho-\sigma\|_{\mathrm{st}}.                 \tag{37.13}
\]
\end{lemma}

\begin{proof}
For \(\|f\|_\infty\leq1\), put
\(\kappa f(y)=\int f(x)\kappa(y,dx)\).  Then
\(\|\kappa f\|_\infty\leq1\), and
\[
\begin{split}
 |\rho(\kappa f)-\sigma(\kappa'f)|
 &\leq |(\rho-\sigma)(\kappa f)|
       +|\sigma((\kappa-\kappa')f)|\\
 &\leq \|\rho-\sigma\|_{\mathrm{st}}
       +\sup_y\|\kappa(y,\cdot)-\kappa'(y,\cdot)\|_{\mathrm{st}}.
                                                               \tag{37.14}
\end{split}
\]
Take the supremum over \(f\).  Setting \(\kappa'=\kappa\) proves
(37.13).
\end{proof}

\begin{proposition}[Common local specification reduces the owner]
\label{prop:v16v37-common}
Consider two finite-volume Gibbs measures whose product reference
measures and potentials agree on every hyperedge meeting \(D\).
If \(\rho,\sigma\) are their respective blanket laws, then their
\(D\)-marginals satisfy
\[
 \|(\mu_V)_D-(\mu_{V'})_D\|_{\mathrm{st}}
 \leq\|\rho-\sigma\|_{\mathrm{st}}.                 \tag{37.15}
\]
Thus all hyperedges disjoint from \(D\) disappear from the
conditional-kernel owner, but they remain collectively represented
by the blanket-law owner.
\end{proposition}

\begin{proof}
Theorem~\ref{thm:v16v37-conditional} gives the same kernel
\(\kappa_D\) for the two measures.  Use (37.11) and the contraction
(37.13).
\end{proof}

\begin{remark}[Why (37.15) is a reduction, not decay]
\label{rem:v16v37-no-decay}
The coefficient in (37.15) is one.  Therefore conditional
cancellation alone does not make a distant boundary invisible.
A spatial mixing estimate must show that the composed kernel from a
remote blanket to the fixed observed set has contraction coefficient
tending to zero with collar depth.
\end{remark}

\subsection{Geometric number of finite-range shell owners}

Work now on \(\mathbb Z^d\) with the \(\ell^\infty\) metric.  Put
\[
 B_R:=\{v\in\mathbb Z^d:\|v\|_\infty\leq R\},\qquad
 A_q:=B_{R+(q+1)r}\setminus B_{R+qr}.               \tag{37.16}
\]
Assume every interaction hyperedge has a chosen anchor \(a(A)\in A\),
diameter at most \(r\), and that at most \(H\) hyperedges have a given
anchor.

\begin{proposition}[Exact shell count and polynomial upper bound]
\label{prop:v16v37-count}
The number \(N_q\) of range-\(r\) interaction owners anchored in the
\(q\)-th shell obeys
\[
\begin{split}
 N_q
 &\leq H|A_q|\\
 &=H\left(
 [2(R+(q+1)r)+1]^d-[2(R+qr)+1]^d\right)\\
 &\leq 2dHr\,[2(R+(q+1)r)+1]^{d-1}.                \tag{37.17}
\end{split}
\]
In four lattice dimensions the owner multiplicity therefore grows at
most cubically with the shell radius.
\end{proposition}

\begin{proof}
The first inequality is the anchor multiplicity assumption.  Since
\(|B_s|=(2s+1)^d\), subtraction gives the equality.  For
\(g(s)=(2s+1)^d\), the mean-value theorem and monotonicity of
\(g'(s)=2d(2s+1)^{d-1}\) give
\[
 g(R+(q+1)r)-g(R+qr)
 \leq r\,g'(R+(q+1)r),                              \tag{37.18}
\]
which is the last bound.
\end{proof}

\begin{corollary}[Per-cell debit required by the countable ledger]
\label{cor:v16v37-cell-debit}
Suppose every owner anchored in \(A_q\), when inserted in the chosen
physical order, has debit at most \(b_q\).  Its aggregate finite-range
shell debit is bounded by
\[
 d_q^{\rm fr}
 \leq 2dHr\,[2(R+(q+1)r)+1]^{d-1}b_q.              \tag{37.19}
\]
Consequently V36's countable assembly applies whenever
\[
 \sum_{q=0}^{\infty}
 [2(R+(q+1)r)+1]^{d-1}b_q<\infty.                 \tag{37.20}
\]
For \(d=4\), a power bound \(b_q=O(q^{-4-\epsilon})\) with
\(\epsilon>0\), or any exponential bound, is sufficient.
\end{corollary}

\begin{proof}
Sum the individual debits over the at most \(N_q\) owners and use
(37.17).  Equation (37.20) makes the resulting countable sum finite.
In dimension four its summand is \(O(q^3b_q)\), proving the last
claim.
\end{proof}

\begin{assessment}[What V37 closes]
\label{ass:v16v37-status}
V37 replaces an invalid support-only cancellation claim by the exact
conditional Gibbs statement.  Remote finite-range factors cancel
from the local conditional kernel, while their entire residual effect
is carried by a finite Markov blanket.  It also proves that in four
dimensions a per-cell influence estimate must beat cubic shell
multiplicity.

Gauge plaquettes, local Higgs terms, and genuinely local counterterms
fit this finite-hypergraph framework after a gauge-fixed positive
finite-cutoff measure has been constructed.  The Newton kernel and
fermion determinant are nonlocal and are not covered.  Even for the
finite-range sectors, (37.15) supplies no decay of the blanket-law
owner.
\end{assessment}

\begin{decision}[V38 spatial contraction across a buffered collar]
\label{dec:v16v37-next}
Go upstream to the missing mixing input.  Formulate a Dobrushin
interdependence matrix for the positive finite-cutoff local
specification, prove the path expansion of boundary influence through
a collar, and combine its decay with the cubic count (37.17).  Record
precisely which gauge, Higgs, and metric modes obstruct a contraction
constant below one.
\end{decision}

\begin{firewall}[Claim boundary after V37]
\label{fire:v16v37}
Finite interaction range is not a proof of spatial mixing.  V37 does
not establish a Dobrushin bound for the physical theory, treat
Newton or determinant nonlocality, prove reflection positivity,
recover the Einstein equation, or complete the continuum.
\end{firewall}

\section*{Version V37 append-only record}
\addcontentsline{toc}{section}{Version V37 append-only record}
The complete V36 source was retained before this continuation.  It
had \(343235\) bytes, \(10432\) lines, and SHA--256
\[
 \texttt{AB3D40712D69A7CAC27CEEBC301453337A67D1CADA055AF5A55E4B626F0E70F7}.
\]
V37 proved conditional cancellation and Markov-blanket reduction,
identified the residual boundary-law owner, and derived the
finite-range shell multiplicity that later mixing debits must beat.

\section{V38: Dobrushin transmission through a buffered collar}
\label{sec:v16v38}

V37 compresses every remote finite-range change into a boundary-law
owner but leaves its state-distance coefficient equal to one.  This
section supplies a rigorous decay mechanism under an explicit
high-temperature or massive-sector hypothesis.  The result is a
conditional theorem: the required contraction is not asserted for
massless gauge or metric modes.

\subsection{Single-site influence matrix}

Let \(D\) be a finite set of dynamical sites and \(B\) a set of fixed
boundary sites.  Every one-site space is assumed Polish.  For
\(u\in D\), let
\(\gamma_u(\,\cdot\,|x_{D\setminus\{u\}},\eta_B)\) be a one-site
conditional probability kernel.  Define the probabilists' variation
distance by
\[
 {\rm tv}(\alpha,\beta):={1\over2}
 \|\alpha-\beta\|_{\mathrm{st}}.                    \tag{38.1}
\]
For \(u\in D\) and \(v\in(D\setminus\{u\})\cup B\), set
\[
 C_{uv}:=
 \sup_{\substack{x,z:\\x_{v^c}=z_{v^c}}}
 {\rm tv}\bigl(\gamma_u(\,\cdot\,|x),
                \gamma_u(\,\cdot\,|z)\bigr).        \tag{38.2}
\]
Here the supremum varies only coordinate \(v\), with the boundary
condition included among the coordinates.  Put \(C_{uu}=0\).

\begin{definition}[Uniform Dobrushin reserve]
\label{def:v16v38-reserve}
The specification has reserve \(\delta>0\) when
\[
 c:=\sup_{u\in D}\sum_{v\in D\cup B}C_{uv}
 \leq1-\delta<1.                                   \tag{38.3}
\]
\end{definition}

\begin{lemma}[Energy-oscillation loader for \(C_{uv}\)]
\label{lem:v16v38-energy}
Suppose changing only coordinate \(v\) changes the conditional
one-site energy at \(u\) from \(H_u(s;x)\) to \(H_u(s;z)\), and
\[
 \operatorname{osc}_{s}
 \bigl(H_u(s;x)-H_u(s;z)\bigr)\leq a_{uv}           \tag{38.4}
\]
uniformly in the remaining coordinates.  Then
\[
 C_{uv}\leq
 \min\left\{1,{e^{a_{uv}}-1\over2}\right\}.         \tag{38.5}
\]
\end{lemma}

\begin{proof}
One conditional law is the normalized tilt of the other by
\(\exp[-(H_u(\,\cdot\,;x)-H_u(\,\cdot\,;z))]\).
The oscillation loader (35.8) bounds their state distance by
\(e^{a_{uv}}-1\).  Divide by two and use the universal bound
\({\rm tv}\leq1\).
\end{proof}

\subsection{The comparison inequality}

We use the following finite-volume coupling form of Dobrushin's
comparison argument.  It is stated for two specifications because
later an individual owner can perturb the local kernel as well as the
boundary condition.

\begin{theorem}[Dobrushin comparison by a path resolvent]
\label{thm:v16v38-comparison}
Let \(\mu,\widetilde\mu\) be finite-volume Gibbs measures on \(D\)
with one-site conditional kernels \(\gamma_u,\widetilde\gamma_u\).
Suppose both families admit the common interdependence majorant
\(C=(C_{uv})_{u,v\in D}\), whose row norm is at most \(c<1\).  Define
the direct specification defect
\[
 b_u:=
 \sup_x{\rm tv}\bigl(\gamma_u(\,\cdot\,|x),
                    \widetilde\gamma_u(\,\cdot\,|x)\bigr).
                                                               \tag{38.6}
\]
Then there is a coupling \((X,\widetilde X)\) of
\((\mu,\widetilde\mu)\) whose disagreement probabilities satisfy
\[
 p_u:=\mathbb P\{X_u\ne\widetilde X_u\}
 \leq\bigl[(I-C)^{-1}b\bigr]_u
 =\sum_{m=0}^{\infty}(C^mb)_u.                     \tag{38.7}
\]
Consequently, for every \(K\subset D\),
\[
 \|\mu_K-\widetilde\mu_K\|_{\mathrm{st}}
 \leq2\sum_{u\in K}\sum_{m=0}^{\infty}(C^mb)_u.     \tag{38.8}
\]
\end{theorem}

\begin{proof}
Couple the two one-site heat-bath updates maximally.  Conditional on a
pair of surrounding configurations, the probability that an update
at \(u\) disagrees is no larger than
\[
 b_u+\sum_{v\in D}C_{uv}\,
 {\bf1}_{\{X_v\ne\widetilde X_v\}}.                 \tag{38.9}
\]
Indeed, insert the unperturbed kernel at the second environment; the
first difference is bounded by \(b_u\), and change the disagreeing
environment coordinates one at a time to obtain the sum involving
\(C_{uv}\).

Run coupled random-scan heat-bath dynamics.  Its two marginals retain
\(\mu\) and \(\widetilde\mu\).  On finite products of Polish spaces,
the set of couplings of these fixed marginals is tight; stationary
Cesàro averages therefore have a weakly convergent subsequence.
Lower semicontinuous approximations to the diagonal-disagreement
indicator, followed by monotone convergence, give for a limiting
coupling
\[
 p_u\leq b_u+\sum_vC_{uv}p_v.                       \tag{38.10}
\]
Iterating this inequality \(M\) times yields
\[
 p\leq\sum_{m=0}^{M}C^mb+C^{M+1}p.                 \tag{38.11}
\]
Because \(0\leq p\leq{\bf1}\) and the row norm of \(C\) is at most
\(c<1\), the final term tends to zero and the series converges.  This
proves (38.7).

For \(\|f\|_\infty\leq1\) depending only on \(K\),
\[
 |\mu(f)-\widetilde\mu(f)|
 \leq2\,\mathbb P\{X_K\ne\widetilde X_K\}
 \leq2\sum_{u\in K}p_u.                             \tag{38.12}
\]
Insert (38.7) and take the supremum over \(f\).
\end{proof}

\begin{remark}[A measurable-state refinement]
\label{rem:v16v38-measurable}
For continuous spins, equality of two independently updated
coordinates may have probability zero unless a maximal coupling is
used; maximal couplings exist for probability measures on Polish
spaces.  The argument concerns disagreement under that chosen
coupling, not equality of generic continuous samples.
\end{remark}

\subsection{Collar depth creates exponential decay}

Assume the influence matrix has range \(r\):
\[
 C_{uv}=0\quad\text{if }\operatorname{dist}(u,v)>r. \tag{38.13}
\]
Let the defect \(b\) be supported on \(S\subset D\), and suppose
\[
 \operatorname{dist}(K,S)\geq L,\qquad
 m_L:=\left\lceil{L\over r}\right\rceil.             \tag{38.14}
\]

\begin{theorem}[Buffered-collar contraction]
\label{thm:v16v38-collar}
Under (38.3), (38.13), and (38.14), if
\(\|b\|_\infty\leq\varepsilon\), then
\[
 \|\mu_K-\widetilde\mu_K\|_{\mathrm{st}}
 \leq {2|K|\,\varepsilon\,c^{m_L}\over1-c}.         \tag{38.15}
\]
In particular, if two boundary conditions induce direct defects only
at sites whose distance from \(K\) is at least \(L\), they obey
\[
 \sup_{\eta,\widetilde\eta}
 \|\mu_K^\eta-\mu_K^{\widetilde\eta}\|_{\mathrm{st}}
 \leq {2|K|\,c^{m_L}\over1-c}.                      \tag{38.16}
\]
\end{theorem}

\begin{proof}
Each nonzero term in \((C^mb)_u\), with \(u\in K\), corresponds to an
influence path of \(m\) range-\(r\) steps from \(u\) to the support of
\(b\).  It therefore vanishes for \(m<m_L\).  The row-sum estimate
\[
 \sup_u\sum_v(C^m)_{uv}\leq c^m                    \tag{38.17}
\]
gives
\[
 \sum_{m=0}^{\infty}(C^mb)_u
 \leq\varepsilon\sum_{m=m_L}^{\infty}c^m
 ={\varepsilon c^{m_L}\over1-c}.                   \tag{38.18}
\]
Insert this in (38.8).  A boundary change has direct defect at most
one in probabilists' variation, so \(\varepsilon=1\) proves (38.16).
\end{proof}

\begin{corollary}[Boundary-law ALC for the finite-range sector]
\label{cor:v16v38-alc}
Let \(D_L\) be a buffered region around fixed \(K\), chosen so that
every site of \(D_L\) whose conditional kernel sees its outer blanket
has distance at least \(L\) from \(K\).  Suppose the two compared
volumes have the same finite-range specification on \(D_L\), and
(38.3) holds uniformly in the cutoff and volume.  Then
their \(K\)-marginals differ, through arbitrary laws on the outer
blanket of \(D_L\), by at most the right side of (38.16).  This debit
tends to zero exponentially as \(L\to\infty\).
\end{corollary}

\begin{proof}
Disintegrate each state over its outer blanket as in
Corollary~\ref{cor:v16v37-blanket}.  Bound the distance between every
pair of conditional \(K\)-laws by (38.16), couple the two arbitrary
blanket laws, and integrate the uniform bound.
\end{proof}

\subsection{Compatibility with the four-dimensional owner count}

\begin{proposition}[Cubic multiplicity is beaten by contraction]
\label{prop:v16v38-cubic}
Suppose an owner in the \(q\)-th range-\(r\) shell has local
specification defect at most \(\varepsilon_q\), and use the
owner-by-owner estimate
\[
 g_q\leq {2|K|\varepsilon_qc^q\over1-c}.            \tag{38.19}
\]
With the four-dimensional count (37.17), the aggregate debit is
summable whenever
\[
 \sum_{q=0}^{\infty}
 [2(R+(q+1)r)+1]^3\varepsilon_qc^q<\infty.          \tag{38.20}
\]
In particular, uniformly bounded \(\varepsilon_q\) suffice for every
\(c<1\).
\end{proposition}

\begin{proof}
Multiply (38.19) by the owner count in (37.17).  A polynomial times
the geometric sequence \(c^q\) is summable when \(c<1\).
\end{proof}

\begin{remark}[The matrix row bound is stronger]
\label{rem:v16v38-row}
If all same-shell perturbations can be combined into one defect
vector before applying Theorem~\ref{thm:v16v38-comparison}, the row
sum (38.17) already sums over their locations and can improve
(38.20).  Proposition~\ref{prop:v16v38-cubic} is the conservative
owner-ledger estimate which remains valid when only separate
per-owner certificates are available.
\end{remark}

\subsection{Physical applicability and obstruction ledger}

Lemma~\ref{lem:v16v38-energy} makes the contraction hypothesis
checkable from local conditional energies.  It is plausible only
after the following modes have been separated or controlled:
\begin{enumerate}
\item gauge-orbit zero directions, which can make a neighbor change
produce order-one conditional variation;
\item the unbroken electromagnetic mode, for which uniform
exponential spatial mixing is not expected;
\item metric and Newton soft modes, already assigned to the separate
projected influence analysis of V14--V24;
\item near-critical Higgs directions, for which the reserve
\(\delta=1-c\) may collapse with the cutoff;
\item fermion determinants, whose effective action is not a
finite-range hypergraph without an additional locality theorem.
\end{enumerate}
Thus the immediate use of (38.15) is for a massive, positive,
finite-range remainder after an explicit soft-mode projection.

\begin{assessment}[What V38 closes]
\label{ass:v16v38-status}
V38 proves a complete transmission theorem from a finite-cutoff
Dobrushin reserve to an exponentially decaying boundary-owner debit.
It combines with V36 and V37 to assemble any uniformly contractive
finite-range sector, and it beats the cubic multiplicity of
four-dimensional shells.

The theorem does not show that the physical coupled specification has
\(c<1\).  Applying it globally to the Standard Model--Einstein system
would wrongly impose exponential mixing on known soft sectors.
\end{assessment}

\begin{decision}[V39 projected soft/hard specification split]
\label{dec:v16v38-next}
Formulate a conditional soft/hard decomposition in which the Newton,
electromagnetic, gauge-orbit, and metric soft coordinates are retained
in a finite-dimensional or explicitly nonlocal carrier, while the
remaining conditional hard specification has a uniform Dobrushin
reserve.  Prove the two-level state-distance inequality that adds the
soft carrier debit to the conditional hard debit.
\end{decision}

\begin{firewall}[Claim boundary after V38]
\label{fire:v16v38}
The Dobrushin reserve is a hypothesis, not a verified property of the
physical lattice action.  V38 does not control the soft carrier,
fermion nonlocality, a determinant phase, reflection positivity of a
gauge fixing, the Einstein equation, or the full continuum.
\end{firewall}

\section*{Version V38 append-only record}
\addcontentsline{toc}{section}{Version V38 append-only record}
The complete V37 source was retained before this continuation.  It
had \(354494\) bytes, \(10750\) lines, and SHA--256
\[
 \texttt{F86E03DB1A4FAC4FB1322CEAB07DEF2DD5E226C7AE27E16A7D15091D8129B10E}.
\]
V38 proved Dobrushin path-resolvent comparison, exponential buffered
collar contraction, and summability against four-dimensional shell
multiplicity.

\section{V39: Conditional soft/hard assembly}
\label{sec:v16v39}

The global Dobrushin hypothesis of V38 is incompatible with leaving
known zero and massless directions inside the contracting subsystem.
The correct architecture is a disintegration: retain those directions
in a soft carrier, and ask for contraction only after conditioning on
that carrier.  This section proves the exact two-level state-distance
inequality and its good-set and collar consequences.

\subsection{Disintegration over a common soft carrier}

Let \(S\) and \(H\) be standard Borel spaces.  Two probability measures
on \(S\times H\) admit disintegrations
\[
 \mu(ds,dh)=\alpha(ds)\kappa_s(dh),\qquad
 \nu(ds,dh)=\beta(ds)\lambda_s(dh).                 \tag{39.1}
\]
The common \(S\)-coordinate is essential.  When cutoff-dependent soft
subspaces move, the augmented-coordinate construction of V29--V31 is
to be used first, so both states are represented on one carrier.

\begin{theorem}[Two-level state-distance inequality]
\label{thm:v16v39-two-level}
For (39.1),
\[
 \|\nu-\mu\|_{\mathrm{st}}
 \leq
 \|\beta-\alpha\|_{\mathrm{st}}
 \int_S\|\lambda_s-\kappa_s\|_{\mathrm{st}}\,\beta(ds).
                                                               \tag{39.2}
\]
The same conclusion holds after pushing both measures to any common
local observable space.
\end{theorem}

\begin{proof}
For measurable \(f:S\times H\to\mathbb C\) with
\(\|f\|_\infty\leq1\), write
\[
 \kappa f(s):=\int_H f(s,h)\kappa_s(dh),\qquad
 \lambda f(s):=\int_H f(s,h)\lambda_s(dh).          \tag{39.3}
\]
Then \(\|\kappa f\|_\infty\leq1\), and
\[
\begin{split}
 \nu(f)-\mu(f)
 &=\beta(\lambda f-\kappa f)+(\beta-\alpha)(\kappa f),\\
 |\nu(f)-\mu(f)|
 &\leq
 \int_S\|\lambda_s-\kappa_s\|_{\mathrm{st}}\,\beta(ds)
 +\|\beta-\alpha\|_{\mathrm{st}}.                  \tag{39.4}
\end{split}
\]
Take the supremum over \(f\).  A measurable pushforward is a Markov
kernel and contracts state norm by (37.13), proving the final claim.
\end{proof}

\begin{remark}[Asymmetry is optional]
\label{rem:v16v39-asymmetry}
Adding and subtracting \(\alpha(ds)\lambda_s(dh)\) instead gives the
same inequality with the conditional distance integrated against
\(\alpha\).  One may use whichever carrier law has the available
moment estimates.
\end{remark}

\begin{corollary}[Soft debit plus a hard good-set debit]
\label{cor:v16v39-good}
Suppose
\[
 \|\beta-\alpha\|_{\mathrm{st}}\leq d_S,            \tag{39.5}
\]
and there is a measurable \(G\subset S\) with
\[
 \beta(G^c)\leq\eta,\qquad
 \sup_{s\in G}\|\lambda_s-\kappa_s\|_{\mathrm{st}}\leq d_H.
                                                               \tag{39.6}
\]
Then
\[
 \|\nu-\mu\|_{\mathrm{st}}\leq d_S+d_H+2\eta.       \tag{39.7}
\]
\end{corollary}

\begin{proof}
The state distance between two probability measures is at most two.
Therefore
\[
 \int_S\|\lambda_s-\kappa_s\|_{\mathrm{st}}\,\beta(ds)
 \leq d_H\beta(G)+2\beta(G^c)\leq d_H+2\eta.        \tag{39.8}
\]
Use (39.2).
\end{proof}

\subsection{A conditional Dobrushin loader}

For fixed \(s\in S\), suppose \(\kappa_s,\lambda_s\) are hard-sector
Gibbs laws in a buffered region, with a common interdependence
majorant \(C(s)\).  Write
\[
 c(s):=\sup_u\sum_v C_{uv}(s),\qquad
 \varepsilon(s):=\sup_u b_u(s),                    \tag{39.9}
\]
where \(b(s)\) is the direct specification defect (38.6).  Let
\(m_L=\lceil L/r\rceil\), and assume the defect support is at distance
at least \(L\) from the hard coordinates observed in \(K\).

\begin{proposition}[Integrated conditional collar debit]
\label{prop:v16v39-integrated}
On the set \(\{s:c(s)<1\}\), define
\[
 q_L(s):=
 {2|K_H|\,\varepsilon(s)c(s)^{m_L}\over1-c(s)},
                                                               \tag{39.10}
\]
truncated at the universal state-distance bound \(2\):
\(\widehat q_L(s):=\min\{2,q_L(s)\}\).
Then
\[
 \|\nu_K-\mu_K\|_{\mathrm{st}}
 \leq
 \|\beta-\alpha\|_{\mathrm{st}}
 \int_S\widehat q_L(s)\,\beta(ds),                 \tag{39.11}
\]
provided \(\widehat q_L=2\) is assigned on \(\{c\geq1\}\).
\end{proposition}

\begin{proof}
For each fixed \(s\) with \(c(s)<1\),
Theorem~\ref{thm:v16v38-collar} bounds the conditional hard
state distance by \(q_L(s)\), and the universal bound permits the
truncation.  On \(c(s)\geq1\), only the universal bound is used.
Insert the resulting pointwise estimate in (39.2), then contract to
the \(K\)-observable space.
\end{proof}

\begin{corollary}[Uniform reserve on a high-probability soft set]
\label{cor:v16v39-reserve}
Suppose for some \(0<c_*<1\), \(\varepsilon_*<\infty\), and events
\(G_L\subset S\),
\[
 c(s)\leq c_*,\qquad \varepsilon(s)\leq\varepsilon_*
 \quad(s\in G_L),\qquad
 \beta(G_L^c)\leq\eta_L.                            \tag{39.12}
\]
Then
\[
 \|\nu_K-\mu_K\|_{\mathrm{st}}
 \leq d_S+
 {2|K_H|\varepsilon_*c_*^{m_L}\over1-c_*}
 +2\eta_L,                                         \tag{39.13}
\]
where \(d_S=\|\beta-\alpha\|_{\mathrm{st}}\).
If \(d_S\to0\), \(L\to\infty\), and \(\eta_L\to0\), the joint local
states are compatible.
\end{corollary}

\begin{proof}
Use (39.10) on \(G_L\), the bound \(2\) on its complement, and
Proposition~\ref{prop:v16v39-integrated}.
\end{proof}

\begin{proposition}[Moment version with a collapsing reserve]
\label{prop:v16v39-moment}
Let
\[
 Z_L(s):={\varepsilon(s)c(s)^{m_L}\over1-c(s)}
 {\bf1}_{\{c(s)<1\}}.                               \tag{39.14}
\]
If
\[
 \beta\{c\geq1\}\longrightarrow0,\qquad
 \int_{\{c<1\}} Z_L(s)\,\beta(ds)\longrightarrow0, \tag{39.15}
\]
then the conditional hard debit in (39.11) tends to zero.  Thus a
uniform deterministic reserve is sufficient but not necessary.
\end{proposition}

\begin{proof}
By (39.10),
\[
 \int_S\widehat q_L\,d\beta
 \leq2|K_H|\int_{\{c<1\}}Z_L\,d\beta
      +2\beta\{c\geq1\}.                            \tag{39.16}
\]
Both terms vanish by (39.15).
\end{proof}

\subsection{A triangular soft/hard continuum criterion}

\begin{theorem}[Soft/hard local compatibility card]
\label{thm:v16v39-card}
Let \(n\) index cutoff and volume comparisons, and write their local
joint states as
\[
 \mu_n=\alpha_n\kappa_{n,s},\qquad
 \nu_n=\beta_n\lambda_{n,s}.                        \tag{39.17}
\]
Assume:
\begin{enumerate}
\item the augmented soft carrier is common and
\[
 d_{S,n}:=\|\beta_n-\alpha_n\|_{\mathrm{st}}\to0;   \tag{39.18}
\]
\item collar depths \(L_n\to\infty\) can be chosen so that the
conditional hard kernels satisfy either (39.12) with
\[
 {\,\varepsilon_{*,n}c_{*,n}^{\lceil L_n/r\rceil}\over
 1-c_{*,n}}+\eta_n\longrightarrow0,                \tag{39.19}
\]
or the integrated condition (39.15).
\end{enumerate}
Then the full joint local states satisfy
\[
 \|\nu_n-\mu_n\|_{\mathrm{st}}\longrightarrow0.     \tag{39.20}
\]
The conclusion remains true after adding any countable family of
other owners whose V36 debit tail tends to zero uniformly in \(n\).
\end{theorem}

\begin{proof}
In the uniform case, (39.13), (39.18), and (39.19) prove (39.20).
In the integrated case, combine (39.11), (39.15), and (39.18).
For the last assertion, insert the soft/hard comparison as one step
of the owner path and apply Theorem~\ref{thm:v16v36-triangular}.
\end{proof}

\subsection{Physical allocation of coordinates}

The theorem specifies a proof architecture rather than a completed
decomposition.  A candidate allocation is:
\[
\begin{array}{c|l}
\text{carrier} & \text{coordinates}\\ \hline
S&
\text{projected Newton density; metric kernel and near-kernel modes;}\\
 &\text{unbroken electromagnetic modes; residual gauge-orbit data;}\\
H&
\text{massive gauge polarizations; stable Higgs fluctuations;}\\
 &\text{other gauge-fixed local modes with a verified reserve.}
\end{array}                                         \tag{39.21}
\]
The moving-pseudoinverse and augmented-coordinate results V26--V32
were built precisely to place cutoff-dependent soft spaces in a
common coordinate system.  V19 and V34 provide part of the Newton
carrier debit.  The electromagnetic, gauge-orbit, and metric
components of \(d_{S,n}\) remain open.

Fermions are not automatically hard variables in this table.
Integrating them out produces a determinant with nonlocal dependence;
retaining them as Grassmann variables does not produce a classical
positive conditional probability kernel.  A separate positive
operator or pseudofermion construction is required before the
Dobrushin loader applies.

\begin{remark}[The effective soft law must include hard free energy]
\label{rem:v16v39-free-energy}
The carrier laws \(\alpha_n,\beta_n\) in (39.17) are the actual
marginals after integrating the hard variables.  Comparing only the
bare soft actions misses the conditional partition function
\[
 Z_H(s)=\int_H e^{-S(s,h)}\,dh,                     \tag{39.22}
\]
whose \(s\)-dependence can change the soft law.  Therefore
\(d_{S,n}\) must include the hard free-energy response; it cannot be
declared small from the bare carrier terms alone.
\end{remark}

\begin{assessment}[What V39 closes]
\label{ass:v16v39-status}
V39 proves that massless and zero modes can be kept outside the
contracting finite-range subsystem without losing the local
compatibility argument.  The exact joint-state debit is the soft
marginal distance plus the soft-law average of the conditional hard
distance.  A high-probability reserve and a moment version are both
available.

This exposes the next upstream bottleneck: the effective carrier law,
including the hard free energy, needs its own density-ratio estimate.
No global Dobrushin assumption is imposed on the soft modes.
\end{assessment}

\begin{decision}[V40 companion audit and effective carrier ratio]
\label{dec:v16v39-next}
Perform the mandatory companion audit.  Then disintegrate a positive
joint density explicitly and derive the exact ratio between its soft
marginal and a reference soft law, including the conditional hard
partition function.  Establish oscillation and relative-entropy
loaders for that effective free-energy owner.  Continue in V41 with
the first physical sector selected by the audit.
\end{decision}

\begin{firewall}[Claim boundary after V39]
\label{fire:v16v39}
The soft/hard card is conditional on a common disintegration and on
the stated hard-sector estimates.  It does not prove positivity of a
chiral fermion representation, control the effective soft free
energy, establish gauge-independent reflection positivity, recover
the Einstein equation, or complete the continuum.
\end{firewall}

\section*{Version V39 append-only record}
\addcontentsline{toc}{section}{Version V39 append-only record}
The complete V38 source was retained before this continuation.  It
had \(366267\) bytes, \(11059\) lines, and SHA--256
\[
 \texttt{95A87E0FA6C88935197220085298D5D02C1CA17B48D5EF8ADF99499AEBD1C159}.
\]
V39 proved the exact soft/hard disintegration inequality and loaded
conditional Dobrushin decay through uniform-good-set and integrated
moment criteria.

\section{V40: Companion audit and the effective carrier ratio}
\label{sec:v16v40}

\subsection{Mandatory companion audit}

The newest Yang--Mills companion inspected at this checkpoint was
\texttt{Renewal\_Yang\_Mills\_Mass\_Gap\_Research\_Notes\_V50.tex};
it had \(436041\) bytes, \(11901\) lines, and SHA--256
\[
 \texttt{11155D2772837D2774B9D4F63546F131FF41F387F18969EF04EFD68C48F7A2ED}.
                                                               \tag{40.1}
\]
It proves two new exact mixed-sign power-two Wilson charts and a
direct plus-pencil \(2t\times t\times t\) slab, with the inverted
minus statement.  It does not yet reach the mixed-sign quarter cell,
all transverse octants, or a continuum state.  The Navier--Stokes
companion remained
\texttt{Renewal\_NS\_Stability\_Research\_Notes\_V26.tex}; it had
\(278283\) bytes, \(5319\) lines, and SHA--256
\[
 \texttt{82C887F8C2C69E4F28FAD7C3DC8DE5B9099CB5A4BF431EBA1FFF3E91935978AD}.
                                                               \tag{40.2}
\]

\begin{theorem}[V40 companion-audit decision]
\label{thm:v16v40-audit}
No companion theorem supplies a premise for the effective
soft-carrier density calculation below.
\end{theorem}

\begin{proof}
YM V50 is a finite rational matrix-positivity result for one Wilson
symbol slab.  It does not construct the positive interacting measure,
the disintegration, or the hard conditional partition function
needed in V39.  Its explicit refusal to infer uncovered orientations
from an unproved symmetry agrees with the owner-by-owner discipline
of V35, but this is methodological rather than a mathematical input.
NS V26 concerns rank-one Oseen-kernel derivatives and has no common
measure or operator with the present free-energy problem.
\end{proof}

\subsection{One positive joint tilt}

Let a reference probability measure on standard Borel spaces
disintegrate as
\[
 \mu_0(ds,dh)=\alpha_0(ds)\kappa_{0,s}(dh).          \tag{40.3}
\]
Let \(W:S\times H\to(0,\infty)\) be measurable with
\[
 0<Z:=\mu_0(W)<\infty,\qquad
 d\mu={W\over Z}\,d\mu_0.                           \tag{40.4}
\]
Define the conditional hard partition function and effective
free energy by
\[
 Z_H(s):=\kappa_{0,s}(W(s,\cdot)),\qquad
 F(s):=-\log Z_H(s).                                \tag{40.5}
\]

\begin{theorem}[Exact effective carrier density]
\label{thm:v16v40-effective}
The soft marginal \(\alpha\) and conditional hard law \(\kappa_s\)
of \(\mu\) are
\[
 {d\alpha\over d\alpha_0}(s)
 ={Z_H(s)\over Z}
 ={e^{-F(s)}\over\alpha_0(e^{-F})},                 \tag{40.6}
\]
and
\[
 {d\kappa_s\over d\kappa_{0,s}}(h)
 ={W(s,h)\over Z_H(s)}                              \tag{40.7}
\]
for \(\alpha\)-almost every \(s\).
\end{theorem}

\begin{proof}
For bounded measurable \(f\) on \(S\), Tonelli and disintegration
give
\[
 \mu(f)
 ={1\over Z}\int_S f(s)
    \left(\int_H W(s,h)\kappa_{0,s}(dh)\right)\alpha_0(ds)
 ={1\over Z}\alpha_0(fZ_H).                         \tag{40.8}
\]
This proves the first equality in (40.6), while
\(Z=\alpha_0(Z_H)=\alpha_0(e^{-F})\) proves the second.
For bounded \(g\) on \(S\times H\), substitution of (40.6) and
(40.7) gives
\[
 \int_S\!\int_H g(s,h)\kappa_s(dh)\alpha(ds)
 ={1\over Z}\int_S\!\int_H g(s,h)W(s,h)
 \kappa_{0,s}(dh)\alpha_0(ds),                      \tag{40.9}
\]
which is \(\mu(g)\).  Uniqueness of disintegration proves (40.7).
\end{proof}

\begin{corollary}[Bare soft weight plus hard free energy]
\label{cor:v16v40-bare}
If \(W(s,h)=w_S(s)w_H(s,h)\), then
\[
 F(s)=-\log w_S(s)
      -\log\kappa_{0,s}(w_H(s,\cdot)).              \tag{40.10}
\]
Thus the hard conditional free energy is an inseparable part of the
actual soft owner.
\end{corollary}

\begin{proof}
Equation (40.5) gives
\(Z_H(s)=w_S(s)\kappa_{0,s}(w_H(s,\cdot))\).  Take
minus the logarithm.
\end{proof}

\subsection{Comparison of two tilted joint states}

Let \(\bar\mu\) be any reference probability measure and let
\[
 d\mu_i={W_i\over Z_i}\,d\bar\mu,\qquad
 Z_i:=\bar\mu(W_i),\qquad i\in\{0,1\}.              \tag{40.11}
\]
Assume both weights are strictly positive and integrable.  Denote
their soft marginals by \(\alpha_i\) and their hard conditional laws
by \(\kappa_{i,s}\).  Put
\[
 R(s,h):={W_1(s,h)\over W_0(s,h)},\qquad
 A(s):=\kappa_{0,s}(R(s,\cdot)),\qquad
 \Delta F(s):=-\log A(s).                           \tag{40.12}
\]

\begin{theorem}[Exact ratio compiler]
\label{thm:v16v40-ratio}
For \(\alpha_0\)-almost every \(s\),
\[
 {d\alpha_1\over d\alpha_0}(s)
 ={A(s)\over\alpha_0(A)}
 ={e^{-\Delta F(s)}\over\alpha_0(e^{-\Delta F})},   \tag{40.13}
\]
and
\[
 {d\kappa_{1,s}\over d\kappa_{0,s}}(h)
 ={R(s,h)\over A(s)}.                               \tag{40.14}
\]
Here \(\alpha_0\) and \(\kappa_{0,s}\) denote the marginal and
conditional law of the already tilted state \(\mu_0\) in the
comparison (40.11), not those in (40.3).
\end{theorem}

\begin{proof}
Regard the probability measure built from \(W_0\) as the new
reference.  Then
\[
 {d\mu_1\over d\mu_0}
 ={R\over\mu_0(R)},\qquad
 \mu_0(R)={Z_1\over Z_0}.                           \tag{40.15}
\]
Apply Theorem~\ref{thm:v16v40-effective} with \(W=R\).  Its
conditional partition function is \(A(s)\), giving both formulas.
\end{proof}

\begin{remark}[Notation firewall]
\label{rem:v16v40-notation}
The reuse of the subscript zero in Theorem~\ref{thm:v16v40-ratio}
means the state before the owner insertion.  Formula (40.13) is
therefore conditional on all owners inserted earlier in the V35
path.  This is precisely what prevents an independence assumption.
\end{remark}

\subsection{Oscillation loaders for both levels}

\begin{proposition}[Joint log-ratio compiler]
\label{prop:v16v40-osc}
Suppose there is a measurable \(g:S\to\mathbb R\) and \(\rho\geq0\)
such that
\[
 |\log R(s,h)-g(s)|\leq\rho                         \tag{40.16}
\]
for almost every \((s,h)\).  Then
\[
 |\log A(s)-g(s)|\leq\rho,                          \tag{40.17}
\]
and hence
\[
 \operatorname{osc}_s\Delta F
 \leq\operatorname{osc}_s g+2\rho.                 \tag{40.18}
\]
Moreover, for each fixed \(s\),
\[
 \operatorname{osc}_h\log R(s,\cdot)\leq2\rho.      \tag{40.19}
\]
Consequently,
\[
 \|\alpha_1-\alpha_0\|_{\mathrm{st}}
 \leq e^{\operatorname{osc}g+2\rho}-1,             \tag{40.20}
\]
and
\[
 \|\kappa_{1,s}-\kappa_{0,s}\|_{\mathrm{st}}
 \leq e^{2\rho}-1.                                 \tag{40.21}
\]
\end{proposition}

\begin{proof}
Exponentiating (40.16), then integrating with respect to the
probability kernel \(\kappa_{0,s}\), gives
\[
 e^{g(s)-\rho}\leq A(s)\leq e^{g(s)+\rho}.          \tag{40.22}
\]
Taking logarithms proves (40.17); taking ranges proves (40.18).
Subtract (40.16) at two hard configurations to get (40.19).
Finally use the exact normalized tilts (40.13), (40.14) and the
oscillation loader (35.8).
\end{proof}

\begin{corollary}[Pure hard fluctuation]
\label{cor:v16v40-pure-hard}
If (40.16) holds with a constant \(g\), then both the induced soft
free-energy debit and every conditional hard debit are at most
\[
 e^{2\rho}-1.                                       \tag{40.23}
\]
Thus a joint owner whose log ratio is uniformly close to one constant
cannot generate a large hidden carrier tilt through integration over
the hard variables.
\end{corollary}

\begin{proof}
Set \(\operatorname{osc}g=0\) in (40.20) and use (40.21).
\end{proof}

\subsection{Relative-entropy chain rule}

\begin{theorem}[Entropy splits exactly across the carrier]
\label{thm:v16v40-chain}
Let
\[
 \mu=\alpha(ds)\kappa_s(dh),\qquad
 \nu=\beta(ds)\lambda_s(dh),                        \tag{40.24}
\]
with \(\nu\ll\mu\).  Then
\[
 D(\nu\Vert\mu)
 =D(\beta\Vert\alpha)
 +\int_S D(\lambda_s\Vert\kappa_s)\,\beta(ds),      \tag{40.25}
\]
with equality in \([0,\infty]\).  Therefore
\[
 D(\beta\Vert\alpha)\leq D(\nu\Vert\mu),\qquad
 \int_S D(\lambda_s\Vert\kappa_s)\,d\beta
 \leq D(\nu\Vert\mu).                               \tag{40.26}
\]
\end{theorem}

\begin{proof}
Absolute continuity and disintegration give, almost everywhere,
\[
 {d\nu\over d\mu}(s,h)
 ={d\beta\over d\alpha}(s)
  {d\lambda_s\over d\kappa_s}(h).                  \tag{40.27}
\]
Take logarithms and integrate with respect to
\(\nu=\beta\lambda_s\).  Tonelli applied first to the positive and
negative truncations, followed by the usual extended-entropy
convention, yields (40.25).  Both terms on its right are
nonnegative, proving (40.26).
\end{proof}

\begin{corollary}[One entropy certificate controls both levels]
\label{cor:v16v40-entropy}
If \(D(\nu_n\Vert\mu_n)\to0\), then the soft marginals converge in
state norm and the conditional hard laws obey
\[
 \|\beta_n-\alpha_n\|_{\mathrm{st}}
 \leq\sqrt{2D(\nu_n\Vert\mu_n)}\to0,                \tag{40.28}
\]
and
\[
 \int_S
 \|\lambda_{n,s}-\kappa_{n,s}\|_{\mathrm{st}}^2
 \,\beta_n(ds)
 \leq2D(\nu_n\Vert\mu_n)\to0.                      \tag{40.29}
\]
\end{corollary}

\begin{proof}
Equation (40.28) is Pinsker plus the first inequality in (40.26).
Pointwise Pinsker gives
\(\|\lambda_s-\kappa_s\|_{\mathrm{st}}^2
\leq2D(\lambda_s\Vert\kappa_s)\); integrate and use (40.26).
\end{proof}

\begin{assessment}[What V40 closes]
\label{ass:v16v40-status}
V40 makes the hard free energy visible in the owner ledger.  For any
positive joint density ratio, it gives the exact induced soft tilt,
the exact conditional hard tilt, quantitative oscillation loaders,
and the entropy chain rule.  This closes the formal gap identified in
Remark~\ref{rem:v16v39-free-energy}.

The calculation still needs physical estimates on
\(\Delta F(s)=-\log\kappa_{0,s}(R)\).  Uniform pointwise control such
as (40.16) is often too strong; differentiating the conditional
partition function and controlling its gradient is the next upstream
route.
\end{assessment}

\begin{decision}[V41 derivative control of the hard free energy]
\label{dec:v16v40-next}
For a differentiable positive hard partition function
\(Z_H(s)=\int e^{-H(s,h)}\,d\lambda(h)\), prove first- and
second-derivative identities for \(F=-\log Z_H\), including the
covariance term.  Convert uniform or exponential-moment bounds on
\(\partial_sH\) into an oscillation debit along the soft carrier.
Keep the result conditional and identify the finite-volume
differentiation hypotheses explicitly.
\end{decision}

\begin{firewall}[Claim boundary after V40]
\label{fire:v16v40}
All ratios in V40 require positive densities.  They do not apply to
an uncontrolled complex determinant phase.  The entropy chain rule
does not itself make the entropy small.  No physical free-energy
gradient, continuum measure, Einstein equation, or full
Standard Model--Einstein continuum has been proved.
\end{firewall}

\section*{Version V40 append-only record}
\addcontentsline{toc}{section}{Version V40 append-only record}
The complete V39 source was retained before this continuation.  It
had \(377180\) bytes, \(11371\) lines, and SHA--256
\[
 \texttt{9804DC24D80A0D0339D914010B24E277C1189F56C72CBD9D8B19395653F16190}.
\]
V40 performed the mandatory companion audit and proved the effective
soft-density ratio, its conditional hard counterpart, oscillation
loaders, and the exact relative-entropy chain rule.

\section{V41: Differentiating the conditional hard free energy}
\label{sec:v16v41}

V40 reduces the induced soft owner to
\(F(s)=-\log Z_H(s)\).  Uniform control of the joint log ratio is
usually too expensive in a field theory, whereas derivatives of a
finite-volume partition function can be expectations of local
insertions.  This section proves the derivative identities with
explicit hypotheses and turns them into oscillation and good-set
debits.

\subsection{A dominated differentiation theorem}

Let \(U\subset\mathbb R^m\) be open, let \((H,\mathcal H,\lambda)\)
be a sigma-finite measure space, and let
\(\mathcal E:U\times H\to\mathbb R\) be measurable.  Define
\[
 Z(s):=\int_H e^{-\mathcal E(s,h)}\,\lambda(dh),
 \qquad
 F(s):=-\log Z(s),                                  \tag{41.1}
\]
and, when \(0<Z(s)<\infty\),
\[
 \kappa_s(dh):={e^{-\mathcal E(s,h)}\over Z(s)}
 \lambda(dh).                                      \tag{41.2}
\]

\begin{hypothesis}[Local second-order domination]
\label{hyp:v16v41-domination}
For \(\lambda\)-almost every \(h\), the map
\(s\mapsto\mathcal E(s,h)\) is \(C^2\).  For every compact
\(K\Subset U\), there is \(M_K\in L^1(\lambda)\) such that, for all
\(s\in K\) and \(1\leq i,j\leq m\),
\[
 e^{-\mathcal E(s,h)}
 \left(1+|\partial_i\mathcal E(s,h)|
 +|\partial_{ij}\mathcal E(s,h)|
 +|\partial_i\mathcal E(s,h)\partial_j\mathcal E(s,h)|\right)
 \leq M_K(h).                                      \tag{41.3}
\]
Also \(0<Z(s)<\infty\) on \(U\).
\end{hypothesis}

\begin{theorem}[Free-energy gradient and Hessian]
\label{thm:v16v41-derivatives}
Under Hypothesis~\ref{hyp:v16v41-domination}, \(Z\) and \(F\) are
\(C^2\), and
\[
 \partial_iF(s)
 =\mathbb E_s[\partial_i\mathcal E],                \tag{41.4}
\]
\[
 \partial_{ij}F(s)
 =\mathbb E_s[\partial_{ij}\mathcal E]
 -\operatorname{Cov}_s
   (\partial_i\mathcal E,\partial_j\mathcal E).     \tag{41.5}
\]
Here all expectations and covariances are under \(\kappa_s\).
\end{theorem}

\begin{proof}
The first two derivatives of the integrand are
\[
 \partial_i e^{-\mathcal E}
 =-(\partial_i\mathcal E)e^{-\mathcal E},           \tag{41.6}
\]
\[
 \partial_{ij} e^{-\mathcal E}
 =
 \bigl(\partial_i\mathcal E\,\partial_j\mathcal E
       -\partial_{ij}\mathcal E\bigr)e^{-\mathcal E}.           \tag{41.7}
\]
The domination (41.3) permits differentiation under the integral
twice, locally uniformly in \(s\).  Thus
\[
 \partial_i Z=-Z\,\mathbb E_s[\partial_i\mathcal E]             \tag{41.8}
\]
and
\[
 \partial_{ij}Z
 =Z\,\mathbb E_s[
 \partial_i\mathcal E\,\partial_j\mathcal E
 -\partial_{ij}\mathcal E].                         \tag{41.9}
\]
Since \(F=-\log Z\), (41.8) proves (41.4), while
\[
 \partial_{ij}F
 =-{\,\partial_{ij}Z\over Z}
 +{\,\partial_iZ\,\partial_jZ\over Z^2}             \tag{41.10}
\]
and (41.8)--(41.9) give (41.5).
\end{proof}

\begin{corollary}[Directional curvature]
\label{cor:v16v41-directional}
For \(a\in\mathbb R^m\),
\[
 \partial_a^2F
 =\mathbb E_s[\partial_a^2\mathcal E]
  -\operatorname{Var}_s(\partial_a\mathcal E).      \tag{41.11}
\]
In particular, if \(\mathcal E\) is affine in \(s\), then \(F\) is
concave.
\end{corollary}

\begin{proof}
Contract (41.5) with \(a_i a_j\).  For affine dependence, the first
term vanishes and variance is nonnegative.
\end{proof}

\subsection{Two physical states and the effective ratio}

For \(i=0,1\), let
\[
 Z_i(s)=\int e^{-\mathcal E_i(s,h)}\,d\lambda(h),
 \quad F_i(s)=-\log Z_i(s),\quad
 \kappa_{i,s}(dh)={e^{-\mathcal E_i(s,h)}\over Z_i(s)}d\lambda(h),
                                                               \tag{41.12}
\]
with Hypothesis~\ref{hyp:v16v41-domination} for each energy.  The
effective carrier increment is
\[
 \Delta F(s):=F_1(s)-F_0(s)
 =-\log{Z_1(s)\over Z_0(s)}.                        \tag{41.13}
\]
This is the \(\Delta F\) of (40.12) when
\(R=e^{-(\mathcal E_1-\mathcal E_0)}\).

\begin{proposition}[Exact response of the effective increment]
\label{prop:v16v41-response}
One has
\[
 \partial_i\Delta F(s)
 =
 \mathbb E_{1,s}[\partial_i\mathcal E_1]
 -\mathbb E_{0,s}[\partial_i\mathcal E_0],          \tag{41.14}
\]
and
\[
\begin{split}
 \partial_{ij}\Delta F(s)
 &=
 \mathbb E_{1,s}[\partial_{ij}\mathcal E_1]
 -\mathbb E_{0,s}[\partial_{ij}\mathcal E_0]\\
 &\quad
 -\operatorname{Cov}_{1,s}
   (\partial_i\mathcal E_1,\partial_j\mathcal E_1)
 +\operatorname{Cov}_{0,s}
   (\partial_i\mathcal E_0,\partial_j\mathcal E_0).
                                                               \tag{41.15}
\end{split}
\]
\end{proposition}

\begin{proof}
Subtract (41.4) and (41.5) for the two energies.
\end{proof}

\begin{remark}[The covariance term cannot be discarded]
\label{rem:v16v41-covariance}
Even when the microscopic energy is linear in a soft coordinate, its
effective free energy has curvature generated by fluctuations.
Replacing (41.5) by the expectation of
\(\partial_{ij}\mathcal E\) alone would omit the variance response and
can reverse a claimed convexity sign.
\end{remark}

\subsection{Gradient-to-oscillation conversion}

Fix a norm \(\|\cdot\|\) on \(\mathbb R^m\) and its dual norm
\(\|\cdot\|_*\).  For a set \(G\), write
\[
 \operatorname{diam}(G)
 :=\sup_{s,t\in G}\|s-t\|.                          \tag{41.16}
\]

\begin{theorem}[Free-energy path bound]
\label{thm:v16v41-path}
Let \(G\subset U\) be convex.  If
\[
 L_G:=\sup_{s\in G}\|\nabla\Delta F(s)\|_*<\infty,  \tag{41.17}
\]
then
\[
 \operatorname{osc}_G\Delta F
 \leq L_G\,\operatorname{diam}(G).                  \tag{41.18}
\]
Moreover,
\[
 L_G\leq
 \sup_{s\in G}
 \left(
 \mathbb E_{1,s}\|\nabla\mathcal E_1\|_*
 +\mathbb E_{0,s}\|\nabla\mathcal E_0\|_*
 \right).                                          \tag{41.19}
\]
\end{theorem}

\begin{proof}
For \(s,t\in G\), the segment \(s_\tau=t+\tau(s-t)\) remains in
\(G\), and the fundamental theorem of calculus gives
\[
\begin{split}
 |\Delta F(s)-\Delta F(t)|
 &\leq\int_0^1
 |\langle\nabla\Delta F(s_\tau),s-t\rangle|\,d\tau\\
 &\leq L_G\|s-t\|.                                  \tag{41.20}
\end{split}
\]
Take the supremum over \(s,t\).  Equation (41.14), the triangle
inequality, and Jensen give (41.19).
\end{proof}

\begin{corollary}[Exponential-moment gradient loader]
\label{cor:v16v41-exponential}
Suppose that for \(i=0,1\), some \(a_i>0\) and \(M_i\geq1\) satisfy
\[
 \sup_{s\in G}
 \mathbb E_{i,s}
 \exp\{a_i\|\nabla\mathcal E_i(s,h)\|_*\}
 \leq M_i.                                         \tag{41.21}
\]
Then
\[
 L_G\leq{\log M_1\over a_1}+{\log M_0\over a_0}.    \tag{41.22}
\]
\end{corollary}

\begin{proof}
For a nonnegative random variable \(X\), Jensen gives
\(e^{a\mathbb EX}\leq\mathbb E e^{aX}\).  Apply this to the two
terms in (41.19).
\end{proof}

\begin{proposition}[Anchored Hessian alternative]
\label{prop:v16v41-hessian}
Let \(G\) lie in the closed norm ball of radius \(R\) about
\(s_0\), assume line segments from \(s_0\) to points of \(G\) remain
in \(U\), and suppose
\[
 \|\nabla\Delta F(s_0)\|_*\leq g_0,\qquad
 \sup_{\substack{s\in G\\\|a\|=\|b\|=1}}
 |\partial_{a}\partial_b\Delta F(s)|\leq M.         \tag{41.23}
\]
Then
\[
 \operatorname{osc}_G\Delta F
 \leq2R(g_0+MR).                                    \tag{41.24}
\]
\end{proposition}

\begin{proof}
Integrating the Hessian along the segment from \(s_0\) to \(s\)
gives
\[
 \|\nabla\Delta F(s)\|_*
 \leq g_0+M\|s-s_0\|\leq g_0+MR.                   \tag{41.25}
\]
The diameter of \(G\) is at most \(2R\).  Apply
Theorem~\ref{thm:v16v41-path}.
\end{proof}

\subsection{Loading the effective carrier owner}

Let \(\alpha_1\) be the normalized tilt of \(\alpha_0\) by
\(e^{-\Delta F}\), as in (40.13).

\begin{theorem}[Effective free-energy compatibility debit]
\label{thm:v16v41-debit}
Let \(G\subset U\) be convex and suppose
\[
 \alpha_0(G^c)\leq\eta,\qquad
 \alpha_1(G^c)\leq\eta.                             \tag{41.26}
\]
Under (41.17),
\[
 \|\alpha_1-\alpha_0\|_{\mathrm{st}}
 \leq
 e^{L_G\operatorname{diam}(G)}-1+4\eta.            \tag{41.27}
\]
The same conclusion holds with the exponent in (41.27) replaced by
\(2R(g_0+MR)\) under Proposition~\ref{prop:v16v41-hessian}.
\end{theorem}

\begin{proof}
The normalized restrictions of \(\alpha_0,\alpha_1\) to \(G\) are
again related by the tilt \(e^{-\Delta F}\), with a new normalizer.
Equations (41.18) and (35.8) bound their state distance by the first
term on the right of (41.27).  Apply the good-set compatibility
Lemma~\ref{lem:v16v34-good}.  The Hessian alternative follows from
(41.24).
\end{proof}

\begin{corollary}[Soft/hard continuum loader]
\label{cor:v16v41-continuum}
In the notation of V39, suppose along the cutoff:
\[
 L_{G_n}\operatorname{diam}(G_n)\to0,\qquad
 \alpha_{0,n}(G_n^c)+\alpha_{1,n}(G_n^c)\to0,       \tag{41.28}
\]
and the conditional hard debit in (39.11) tends to zero.  Then the
joint local states converge in state norm.  The conclusion remains
valid after adding the uniformly summable owner tails of V36.
\end{corollary}

\begin{proof}
Theorem~\ref{thm:v16v41-debit} makes the soft debit tend to zero.
Apply Theorem~\ref{thm:v16v39-card}, followed by
Theorem~\ref{thm:v16v36-triangular} for the additional owners.
\end{proof}

\subsection{Physical reading of the derivative insertion}

When a soft coordinate \(s_i\) is a retained Newton, metric,
electromagnetic, or gauge-orbit coordinate, (41.14) says that its
hard free-energy response is an expectation of the corresponding
finite-volume current or stress insertion.  A useful estimate must
therefore establish all of the following:
\begin{enumerate}
\item domination (41.3), uniformly on the chosen soft good set;
\item a bound on the difference of the two insertion expectations in
(41.14), rather than separate volume-divergent bounds whenever
cancellation is essential;
\item a good set whose carrier diameter times that gradient bound
vanishes, as required by (41.28);
\item tail bounds under both actual carrier laws, including the hard
free-energy tilt.
\end{enumerate}
The separate expectation estimate (41.19) is safe but may scale like
the collar volume.  The next refinement should exploit locality:
differentiate only the hard free-energy \emph{increment} caused by a
distant shell and use conditional mixing to show that the response
insertion is confined to the shell and decays toward the observed
carrier coordinates.

\begin{assessment}[What V41 closes]
\label{ass:v16v41-status}
V41 rigorously differentiates the conditional hard partition
function and retains the indispensable covariance term in its
Hessian.  It converts gradient, exponential-moment, or anchored
Hessian estimates into a quantitative soft-carrier state debit and
then into joint local compatibility.

The theorem identifies, but does not prove, the physical insertion
bounds.  Separate absolute gradient estimates can be too large; the
relevant quantity is the derivative of the shell free-energy
difference.
\end{assessment}

\begin{decision}[V42 shell free-energy response localization]
\label{dec:v16v41-next}
Let two hard energies differ by a shell perturbation \(V\).  Derive
the interpolation identity
\[
 \partial_s(F_1-F_0)
 =\mathbb E_{1,s}[\partial_sV]
 -\int_0^1\operatorname{Cov}_{t,s}
   (\partial_s\mathcal E_0,V)\,dt
\]
under explicit domination, with the appropriate additional term
when \(\partial_s\mathcal E_0\) depends on the interpolation.
Then use a stated covariance-decay hypothesis to obtain a shell
gradient debit which can be tested against the four-dimensional
owner count.
\end{decision}

\begin{firewall}[Claim boundary after V41]
\label{fire:v16v41}
Differentiability of a finite positive partition function does not
establish cutoff-uniform response bounds.  V41 does not treat a
complex fermion phase, prove covariance decay in massless sectors,
construct the continuum measure, recover the Einstein equation, or
complete the Standard Model--Einstein continuum.
\end{firewall}

\section*{Version V41 append-only record}
\addcontentsline{toc}{section}{Version V41 append-only record}
The complete V40 source was retained before this continuation.  It
had \(388312\) bytes, \(11709\) lines, and SHA--256
\[
 \texttt{28037B689DFE94394D4A6BC0874DB5831EA38815567AC31E535E96EA9DA75049}.
\]
V41 proved the free-energy response identities and converted their
gradient and curvature bounds into a soft-carrier compatibility
debit, completing the substantive continuation required after the
V40 audit.

\section{V42: Interpolated shell free energy and response
localization}
\label{sec:v16v42}

V41 shows that a separate bound on each full free-energy gradient can
scale with the volume and conceal the cancellation relevant to local
compatibility.  This continuation differentiates the
\emph{difference} created by one shell perturbation.  The difference
is expressed as a direct shell insertion plus a covariance between
the shell and the reference soft current.  A spatial covariance bound
then yields the required far-shell debit.

\subsection{Joint interpolation hypotheses}

Let \(U\subset\mathbb R^m\), let
\(\mathcal E_0:U\times H\to\mathbb R\), and let
\(V:U\times H\to\mathbb R\).  For \(t\in[0,1]\), put
\[
 \mathcal E_t(s,h):=\mathcal E_0(s,h)+tV(s,h),      \tag{42.1}
\]
\[
 Z_t(s):=\int e^{-\mathcal E_t(s,h)}\,d\lambda(h),
 \qquad
 F_t(s):=-\log Z_t(s),                              \tag{42.2}
\]
and denote the normalized hard expectation by
\(\mathbb E_{t,s}\).

\begin{hypothesis}[Joint \(s,t\) domination]
\label{hyp:v16v42-domination}
The energies are \(C^1\) in \(s\), \(V\) is measurable, and on every
compact \(K\Subset U\) there is an integrable envelope which,
uniformly for \((s,t)\in K\times[0,1]\), dominates
\[
 e^{-\mathcal E_t}
 \left(
 1+|V|+\|\nabla_s\mathcal E_0\|_*
 +\|\nabla_sV\|_*
 +|V|\|\nabla_s\mathcal E_0\|_*
 +|V|\|\nabla_sV\|_*
 \right).                                          \tag{42.3}
\]
Also \(0<Z_t(s)<\infty\).
\end{hypothesis}

\begin{lemma}[Covariance response formula]
\label{lem:v16v42-response}
Under Hypothesis~\ref{hyp:v16v42-domination}, if \(A_t(s,h)\) is
differentiable in \(t\) and the analogous dominated differentiation
condition holds, then
\[
 {d\over dt}\mathbb E_{t,s}[A_t]
 =
 \mathbb E_{t,s}[\partial_tA_t]
 -\operatorname{Cov}_{t,s}(A_t,V).                 \tag{42.4}
\]
\end{lemma}

\begin{proof}
Differentiate the quotient
\[
 \mathbb E_{t,s}[A_t]
 ={\,\int A_t e^{-\mathcal E_t}\,d\lambda\over Z_t(s)}.         \tag{42.5}
\]
The numerator derivative is
\(\int(\partial_tA_t-A_tV)e^{-\mathcal E_t}d\lambda\), while
\(\partial_tZ_t=-Z_t\mathbb E_{t,s}V\).  The quotient rule gives
(42.4).
\end{proof}

\subsection{Two exact formulas for the shell gradient}

\begin{theorem}[Interpolated shell free-energy response]
\label{thm:v16v42-interpolation}
Let
\[
 \Delta F(s):=F_1(s)-F_0(s).                        \tag{42.6}
\]
Under Hypothesis~\ref{hyp:v16v42-domination},
\[
 \Delta F(s)=\int_0^1\mathbb E_{t,s}[V]\,dt,        \tag{42.7}
\]
and for every direction \(a\in\mathbb R^m\),
\[
\begin{split}
 \partial_a\Delta F(s)
 &=
 \int_0^1
 \left\{
 \mathbb E_{t,s}[\partial_aV]
 -\operatorname{Cov}_{t,s}
   (\partial_a\mathcal E_t,V)
 \right\}dt                                         \tag{42.8}\\
 &=
 \mathbb E_{1,s}[\partial_aV]
 -\int_0^1
 \operatorname{Cov}_{t,s}
   (\partial_a\mathcal E_0,V)\,dt.                  \tag{42.9}
\end{split}
\]
\end{theorem}

\begin{proof}
Differentiation in \(t\) gives
\[
 \partial_tF_t(s)=\mathbb E_{t,s}[V].               \tag{42.10}
\]
Integrate from zero to one to obtain (42.7).  Differentiating (42.7)
in \(s\), justified by (42.3), and applying
Lemma~\ref{lem:v16v42-response} to \(A_t=V\) with \(s\) and \(t\)
interchanged gives
\[
 \partial_a\mathbb E_{t,s}[V]
 =\mathbb E_{t,s}[\partial_aV]
 -\operatorname{Cov}_{t,s}
  (V,\partial_a\mathcal E_t),                       \tag{42.11}
\]
which proves (42.8).

Alternatively, V41 gives
\[
\begin{split}
 \partial_a\Delta F
 &=
 \mathbb E_{1,s}[\partial_a\mathcal E_0+\partial_aV]
 -\mathbb E_{0,s}[\partial_a\mathcal E_0].          \tag{42.12}
\end{split}
\]
Apply Lemma~\ref{lem:v16v42-response} with the
\(t\)-independent observable \(A=\partial_a\mathcal E_0\):
\[
 \mathbb E_{1,s}A-\mathbb E_{0,s}A
 =-\int_0^1\operatorname{Cov}_{t,s}(A,V)\,dt.       \tag{42.13}
\]
Substitution in (42.12) proves (42.9).
\end{proof}

\begin{remark}[When the direct term vanishes]
\label{rem:v16v42-direct}
If the shell perturbation \(V\) is independent of the retained soft
coordinate in direction \(a\), then \(\partial_aV=0\), and (42.9)
contains only a connected covariance.  This is the cancellation that
separate absolute bounds on the two terms in (41.14) fail to expose.
\end{remark}

\subsection{A spatial covariance loader}

Let \(\mathcal O\) be a fixed set of hard cells associated with the
observed soft coordinate, and let \(\mathcal S_q\) be a shell cell set.
Assume decompositions
\[
 \partial_a\mathcal E_0
 =\sum_{x\in\mathcal O}J_{a,x},\qquad
 V_q=\sum_{y\in\mathcal S_q}V_{q,y}.                \tag{42.14}
\]
Write
\[
 \ell_q:=\operatorname{dist}(\mathcal O,\mathcal S_q).          \tag{42.15}
\]

\begin{hypothesis}[Uniform connected-response decay]
\label{hyp:v16v42-clustering}
There are \(A,\mathfrak m>0\) such that, for every unit direction
\(\|a\|=1\), every \(s\) in a chosen carrier good set \(G\), every
\(t\in[0,1]\), and all \(x\in\mathcal O\),
\(y\in\mathcal S_q\),
\[
 \left|
 \operatorname{Cov}_{t,s}(J_{a,x},V_{q,y})
 \right|
 \leq A e^{-\mathfrak m\operatorname{dist}(x,y)}.   \tag{42.16}
\]
Also the direct response obeys
\[
 \sup_{\substack{s\in G\\\|a\|=1}}
 \left|\mathbb E_{1,s}[\partial_aV_q]\right|
 \leq e_q.                                         \tag{42.17}
\]
\end{hypothesis}

\begin{theorem}[Localized shell-gradient bound]
\label{thm:v16v42-localized}
Under (42.14)--(42.17),
\[
 \sup_{s\in G}\|\nabla\Delta F_q(s)\|_*
 \leq
 e_q+A|\mathcal O|\,|\mathcal S_q|\,
 e^{-\mathfrak m\ell_q}.                           \tag{42.18}
\]
\end{theorem}

\begin{proof}
For a unit direction \(a\), expand the covariance in (42.9) using
(42.14).  Bilinearity and (42.16) give
\[
\begin{split}
 \left|\operatorname{Cov}_{t,s}
 (\partial_a\mathcal E_0,V_q)\right|
 &\leq
 \sum_{x\in\mathcal O}\sum_{y\in\mathcal S_q}
 A e^{-\mathfrak m\operatorname{dist}(x,y)}\\
 &\leq A|\mathcal O|\,|\mathcal S_q|
 e^{-\mathfrak m\ell_q}.                           \tag{42.19}
\end{split}
\]
Integrate in \(t\), add (42.17), and take the supremum over unit
directions.  Duality gives (42.18).
\end{proof}

\begin{corollary}[Four-dimensional shell response]
\label{cor:v16v42-four-dimensional}
In the lattice setting of V37, suppose
\(\mathcal S_q\) contains at most the anchored owners counted in
(37.17) and
\[
 \ell_q\geq qr-\ell_0                              \tag{42.20}
\]
for a fixed \(\ell_0\).  Then
\[
\begin{split}
 \sup_{s\in G}\|\nabla\Delta F_q(s)\|_*
 &\leq e_q+
 8AHr|\mathcal O|\,
 [2(R+(q+1)r)+1]^3
 e^{-\mathfrak m(qr-\ell_0)}                       \tag{42.21}
\end{split}
\]
in dimension four.
\end{corollary}

\begin{proof}
Use (37.17) with \(d=4\) to bound
\[
 |\mathcal S_q|
 \leq8Hr[2(R+(q+1)r)+1]^3,                         \tag{42.22}
\]
then insert (42.20) in (42.18).
\end{proof}

\subsection{Summable effective-carrier debits}

\begin{theorem}[Remote shell-tail criterion]
\label{thm:v16v42-tail}
Let the same convex carrier good set \(G\) have diameter \(D<\infty\).
For shells \(q\geq Q\), let
\[
 L_q:=
 e_q+A|\mathcal O|\,|\mathcal S_q|
 e^{-\mathfrak m\ell_q}.                           \tag{42.23}
\]
Assume
\[
 \sum_{q=0}^{\infty}DL_q<\infty.                   \tag{42.24}
\]
Let \(\alpha^{(q)}\) be the intermediate soft laws obtained by
inserting the shell effective weights in a fixed order, and suppose
the corresponding normalized good-set restrictions obey the same
gradient bounds (42.23).  Then the good-set state-distance tail is
bounded by
\[
 \sum_{q\geq Q}(e^{DL_q}-1),                        \tag{42.25}
\]
and this tends to zero as \(Q\to\infty\).  If the sum of the two
endpoint bad-set probabilities at the \(q\)-th insertion is
\(\eta_q\) and
\[
 \sum_q\eta_q<\infty,                               \tag{42.26}
\]
the full soft owner tail also tends to zero.
\end{theorem}

\begin{proof}
Theorem~\ref{thm:v16v41-path} gives
\(\operatorname{osc}_G\Delta F_q\leq DL_q\).  The oscillation loader
and the V35 telescoping theorem give (42.25) for the restricted laws.
Since \(\sum_qDL_q<\infty\), one has \(DL_q\to0\); for all sufficiently
large \(q\),
\[
 e^{DL_q}-1\leq2DL_q.                               \tag{42.27}
\]
Thus the series in (42.25) converges and its tails vanish.  Apply
Lemma~\ref{lem:v16v34-good} at each insertion; the added bad-set
debits are summable by (42.26).  V36 then gives the full tail limit.
\end{proof}

\begin{corollary}[Exponential clustering beats four-dimensional
shell multiplicity]
\label{cor:v16v42-summability}
Under Corollary~\ref{cor:v16v42-four-dimensional}, (42.24) holds
whenever
\[
 \sum_q e_q<\infty                                  \tag{42.28}
\]
and \(D\) is uniformly bounded.  In particular, the covariance part
is summable because
\[
 \sum_{q=0}^{\infty}
 [2(R+(q+1)r)+1]^3e^{-\mathfrak m(qr-\ell_0)}
 <\infty.                                          \tag{42.29}
\]
\end{corollary}

\begin{proof}
The series (42.29) is a cubic polynomial times a geometric sequence
with ratio \(e^{-\mathfrak m r}<1\).  Combine it with (42.21) and
(42.28).
\end{proof}

\subsection{Interaction with the soft/hard ledger}

The estimates above apply to one positive hard-sector shell owner
after fixing the soft coordinate.  Their place in the programme is:
\[
\begin{array}{c|c}
\text{input}&\text{resulting debit}\\ \hline
\text{direct response }e_q
&\text{first term of }L_q\\
\text{hard connected covariance decay}
&\text{second term of }L_q\\
\text{carrier diameter }D
&e^{DL_q}-1\\
\text{carrier bad-set tails}
&\text{good-set correction}\\
\text{hard conditional collar bound}
&\text{V39 conditional debit}.
\end{array}                                         \tag{42.30}
\]
The soft and hard terms are then added by (39.2), and different
physical sectors by V35--V36.

\begin{assessment}[What V42 closes]
\label{ass:v16v42-status}
V42 replaces two potentially extensive free-energy gradients by the
exact response of their shell difference.  When the shell is
independent of the observed soft direction, only a connected
current--shell covariance remains.  Under uniform exponential
covariance decay, that response is summable even after the cubic
four-dimensional shell multiplicity is included.

The connected-response bound (42.16) remains a hypothesis.  V38
proves boundary-state decay from a Dobrushin reserve, but the present
programme has not yet converted that reserve into the precise
covariance inequality for unbounded current and shell insertions.
\end{assessment}

\begin{decision}[V43 covariance from mixing with truncation]
\label{dec:v16v42-next}
First prove the bounded-observable covariance inequality obtained
from the Dobrushin path resolvent.  Then introduce good-set
truncations for the generally unbounded current \(J_{a,x}\) and shell
energy \(V_{q,y}\), with exponential-moment tails, so that
(42.16) is reduced to explicit bounded oscillations plus controlled
tail terms.
\end{decision}

\begin{firewall}[Claim boundary after V42]
\label{fire:v16v42}
Exponential connected-response decay is not proved here and is not
expected for unseparated massless modes.  The interpolation requires
a positive finite hard measure and does not treat a complex
determinant phase.  No continuum state, Einstein equation, or full
Standard Model--Einstein continuum has been established.
\end{firewall}

\section*{Version V42 append-only record}
\addcontentsline{toc}{section}{Version V42 append-only record}
The complete V41 source was retained before this continuation.  It
had \(400702\) bytes, \(12096\) lines, and SHA--256
\[
 \texttt{BAF4FEF59414B640BC2D7CE9D78DF0CA75E7BACF013F2937F52B856A6CFAD6CE}.
\]
V42 proved the shell free-energy interpolation identities and reduced
summable carrier response to a direct shell derivative plus spatially
decaying connected hard covariances.

\section{V43: Unbounded covariance from archived bounded mixing}
\label{sec:v16v43}

Before proving the next loader, the earlier SM archive was searched for
overlap.  Seventh-cycle V89 in
\texttt{Renewal\_SM\_Einstein\_Continuum\_Research\_Notes\_V100\_f7.tex}
already proves the bounded-observable Dobrushin covariance inequality
and exponential resolvent decay.  That archive has \(1570671\) bytes,
\(40226\) lines, and SHA--256
\[
 \texttt{5F28040CC4F306FE0BA60AE765429A956A3E74F8D1FDDED50F0F7CD5B878A66C}.
                                                               \tag{43.1}
\]
Ninth-cycle V4 also contains the underlying finite Dobrushin
comparison and Doob-range theorem.  Those bounded results are used
here as established inputs and are not reproved.  The new task is to
pass from them to the unbounded current and shell insertions appearing
in (42.16).

\subsection{The archived bounded input}

For a finite positive hard-sector Gibbs law, let
\[
 \delta_i(f):=
 \sup_{x,z:\,x_{i^c}=z_{i^c}}|f(x)-f(z)|.           \tag{43.2}
\]
Let \(C\) be a finite-range interdependence matrix with row sum at
most \(c<1\), range \(r\), and resolvent
\[
 D=(I-C)^{-1}.                                      \tag{43.3}
\]
The archived theorem states that bounded real observables \(f,g\)
satisfy
\[
 |\operatorname{Cov}(f,g)|
 \leq{1\over4}\sum_{i,j}\delta_i(f)D_{ij}\delta_j(g),
                                                               \tag{43.4}
\]
and, for \(i\ne j\),
\[
 D_{ij}\leq
 {c^{\lceil\operatorname{dist}(i,j)/r\rceil}\over1-c}.
                                                               \tag{43.5}
\]
For use below, (43.4)--(43.5) are required uniformly in the shell
interpolation \(t\), soft good-set coordinate \(s\), cutoff, and
volume.  This uniformity is a new physical obligation; citing the
finite archived theorem does not establish it.

\subsection{Clipping and exponential tails}

For \(R>0\), define the clipping map
\[
 T_R(x):=\max\{-R,\min\{x,R\}\}.                    \tag{43.6}
\]

\begin{lemma}[Exponential moment gives an \(L^2\) clipping tail]
\label{lem:v16v43-tail}
If a real random variable \(X\) satisfies
\[
 \mathbb E e^{a|X|}\leq M,\qquad a>0,               \tag{43.7}
\]
then
\[
 \|X\|_2\leq {2\sqrt M\over ae},                    \tag{43.8}
\]
and, for \(R\geq2/a\),
\[
 \|X-T_R(X)\|_2
 \leq R\sqrt M\,e^{-aR/2}.                          \tag{43.9}
\]
\end{lemma}

\begin{proof}
The maximum of \(x^2e^{-ax}\) on \([0,\infty)\) is
\(4/(a^2e^2)\).  Hence
\[
 \mathbb E X^2
 \leq {4\over a^2e^2}\mathbb E e^{a|X|}
 \leq {4M\over a^2e^2},                             \tag{43.10}
\]
which proves (43.8).  If \(R\geq2/a\), the function
\(x^2e^{-ax}\) decreases for \(x\geq R\), so
\[
 x^2{\bf1}_{\{x>R\}}
 \leq R^2e^{-aR}e^{ax}.                             \tag{43.11}
\]
Also \(|X-T_R(X)|\leq |X|{\bf1}_{\{|X|>R\}}\).
Take expectations in (43.11) with \(x=|X|\), then take a square
root to obtain (43.9).
\end{proof}

\begin{lemma}[Covariance error caused by clipping]
\label{lem:v16v43-cov-tail}
For \(X_R=T_R(X)\) and \(Y_S=T_S(Y)\),
\[
\begin{split}
 |\operatorname{Cov}(X,Y)-\operatorname{Cov}(X_R,Y_S)|
 &\leq
 \|X-X_R\|_2\|Y\|_2
 +\|X_R\|_2\|Y-Y_S\|_2.                            \tag{43.12}
\end{split}
\]
If
\[
 \mathbb Ee^{a|X|}\leq M_X,\qquad
 \mathbb Ee^{b|Y|}\leq M_Y,                         \tag{43.13}
\]
and \(R\geq2/a\), \(S\geq2/b\), the right side is at most
\[
 {2R\sqrt{M_XM_Y}\over be}\,e^{-aR/2}
 +RS\sqrt{M_Y}\,e^{-bS/2}.                         \tag{43.14}
\]
\end{lemma}

\begin{proof}
The identity
\[
\begin{split}
 \operatorname{Cov}(X,Y)-\operatorname{Cov}(X_R,Y_S)
 &=
 \operatorname{Cov}(X-X_R,Y)
 +\operatorname{Cov}(X_R,Y-Y_S)                   \tag{43.15}
\end{split}
\]
and Cauchy--Schwarz for centered variables give (43.12).
Use Lemma~\ref{lem:v16v43-tail} for both tails and for
\(\|Y\|_2\), together with \(\|X_R\|_2\leq R\), to get (43.14).
\end{proof}

\subsection{The unbounded separated-support theorem}

\begin{theorem}[Unbounded Dobrushin covariance by clipping]
\label{thm:v16v43-unbounded}
Let \(X,Y\) be real observables supported on cell sets
\(P,Q\), respectively, where
\[
 |P|\leq p,\qquad |Q|\leq q,\qquad
 \operatorname{dist}(P,Q)\geq\ell>0.               \tag{43.16}
\]
Assume (43.4)--(43.5) and the exponential moments (43.13).
Then for \(R\geq2/a\), \(S\geq2/b\),
\[
\begin{split}
 |\operatorname{Cov}(X,Y)|
 &\leq
 {pqRS\over1-c}\,
 c^{\lceil\ell/r\rceil}\\
 &\quad+
 {2R\sqrt{M_XM_Y}\over be}\,e^{-aR/2}
 +RS\sqrt{M_Y}\,e^{-bS/2}.                         \tag{43.17}
\end{split}
\]
\end{theorem}

\begin{proof}
The clipped variables have the same supports, and
\[
 \delta_i(X_R)\leq2R\,{\bf1}_{\{i\in P\}},\qquad
 \delta_j(Y_S)\leq2S\,{\bf1}_{\{j\in Q\}}.          \tag{43.18}
\]
Equations (43.4) and (43.5) therefore give
\[
 |\operatorname{Cov}(X_R,Y_S)|
 \leq
 {pqRS\over1-c}c^{\lceil\ell/r\rceil}.              \tag{43.19}
\]
Add the clipping error (43.14).
\end{proof}

\begin{corollary}[Exponential covariance for unbounded local
insertions]
\label{cor:v16v43-exponential}
Fix \(\kappa>0\), and choose
\[
 R_\ell=R_0+\kappa\ell,\qquad
 S_\ell=S_0+\kappa\ell,\qquad
 R_0\geq2/a,\quad S_0\geq2/b.                       \tag{43.20}
\]
For every
\[
 0<\mathfrak m<
 \min\left\{
 {|\log c|\over r},{a\kappa\over2},{b\kappa\over2}
 \right\},                                         \tag{43.21}
\]
there is a finite constant
\(A=A(a,b,M_X,M_Y,p,q,c,r,\kappa,R_0,S_0,\mathfrak m)\)
such that
\[
 |\operatorname{Cov}(X,Y)|
 \leq A e^{-\mathfrak m\ell}.                       \tag{43.22}
\]
\end{corollary}

\begin{proof}
Insert (43.20) in (43.17).  The three terms are polynomials of degree
at most two in \(\ell\), multiplied respectively by
\[
 e^{-|\log c|\ell/r},\qquad
 e^{-a\kappa\ell/2},\qquad
 e^{-b\kappa\ell/2},                                \tag{43.23}
\]
up to fixed constant factors.  If \(\mathfrak m\) is strictly smaller
than all three rates, each polynomial times its exponential is
bounded by a constant times \(e^{-\mathfrak m\ell}\).
\end{proof}

\subsection{Uniform interpolation loader for V42}

\begin{theorem}[Moment-and-Dobrushin certificate for connected
shell response]
\label{thm:v16v43-certificate}
For every \(t\in[0,1]\), every soft coordinate \(s\) in a good set
\(G\), every unit carrier direction \(a\), and all relevant cells
\(x,y\), suppose:
\begin{enumerate}
\item the interpolated positive hard measure has a common
finite-range Dobrushin majorant satisfying (43.4)--(43.5) with fixed
\(c<1\) and \(r<\infty\);
\item \(J_{a,x}\) and \(V_{q,y}\) have supports of at most \(p\) and
\(q_0\) cells;
\item for fixed \(a_0,b_0>0\) and \(M_J,M_V<\infty\),
\[
 \mathbb E_{t,s}e^{a_0|J_{a,x}|}\leq M_J,\qquad
 \mathbb E_{t,s}e^{b_0|V_{q,y}|}\leq M_V.           \tag{43.24}
\]
\end{enumerate}
Then for any \(\kappa>0\) and any
\[
 0<\mathfrak m<
 \min\left\{
 {|\log c|\over r},{a_0\kappa\over2},
 {b_0\kappa\over2}
 \right\},                                         \tag{43.25}
\]
there is one constant \(A<\infty\), uniform in all the displayed
indices and cutoffs, such that
\[
 \left|
 \operatorname{Cov}_{t,s}(J_{a,x},V_{q,y})
 \right|
 \leq A e^{-\mathfrak m\operatorname{dist}(x,y)}.   \tag{43.26}
\]
Thus Hypothesis~\ref{hyp:v16v42-clustering}, apart from its direct
response row (42.17), follows from these three explicit certificates.
\end{theorem}

\begin{proof}
Apply Corollary~\ref{cor:v16v43-exponential} with the common constants
and \(\ell=\operatorname{dist}(x,y)\).  Uniformity of all premises
makes the constant in (43.22) uniform.  Equation (43.26) is exactly
(42.16).
\end{proof}

\begin{remark}[Why fixed clipping is insufficient]
\label{rem:v16v43-growing}
With fixed \(R,S\), the last two terms in (43.17) do not decay as the
supports separate.  The thresholds must grow with distance.  Their
polynomial cost in the bounded covariance is then absorbed by the
strictly positive Dobrushin and moment exponents.
\end{remark}

\begin{remark}[Block form]
\label{rem:v16v43-block}
The same proof applies to the block Dobrushin matrix of archived
ninth-cycle V4.  The support sizes \(p,q\) are then numbers of blocks,
and \(\delta_i\) is a block oscillation.  This may retain a reserve
when a one-cell worst-exterior coefficient is one, but the block
reserve must still be proved uniformly.
\end{remark}

\begin{assessment}[What V43 adds beyond the archive]
\label{ass:v16v43-status}
The bounded covariance and resolvent decay are deliberately imported
from seventh-cycle V89.  V43 adds a distance-dependent clipping
theorem which converts uniform exponential moments into exponential
covariance decay for unbounded local insertions.  It therefore pays
the connected-response row of V42 whenever a uniform positive hard
Dobrushin certificate and the two local moment bounds are available.

The direct derivative \(e_q\) in (42.17) remains unpaid.  It vanishes
for a strictly disjoint local soft coordinate, but projected or
nonlocal carrier modes can have a tail on the shell.
\end{assessment}

\begin{decision}[V44 direct response through carrier-mode tails]
\label{dec:v16v43-next}
Represent a retained soft direction by its reconstruction mode on
the microscopic fields.  Prove a chain-rule estimate for
\(\mathbb E_{1,s}\partial_aV_q\) in terms of the local derivative norm
of each shell action and the mode amplitude on that shell.  Combine a
weighted mode-tail estimate with the cubic shell count to give a
summable direct-response sequence \(e_q\).
\end{decision}

\begin{firewall}[Claim boundary after V43]
\label{fire:v16v43}
The theorem consumes, rather than proves, a uniform Dobrushin reserve
and exponential insertion moments.  It applies only to positive hard
measures and to soft good sets on which the constants are uniform.
It does not control massless carrier correlations, a fermion phase,
the continuum measure, the Einstein equation, or the full
Standard Model--Einstein continuum.
\end{firewall}

\section*{Version V43 append-only record}
\addcontentsline{toc}{section}{Version V43 append-only record}
The complete V42 source was retained before this continuation.  It
had \(412458\) bytes, \(12471\) lines, and SHA--256
\[
 \texttt{C5BD152E386C8F0C787AC7F73A82C6DDCF5A848E0ACE8FA16F02FDE5425DD776}.
\]
V43 avoided duplicating the archived bounded covariance theorem and
extended it, by growing clipping thresholds, to the unbounded
current--shell covariance required by the V42 response formula.

\section{V44: Direct shell response from reconstruction-mode tails}
\label{sec:v16v44}

V43 pays the connected-covariance part of the shell gradient under
uniform hard-sector hypotheses.  The other term in (42.9),
\(\mathbb E_{1,s}\partial_aV_q\), is present when a retained soft
coordinate has a microscopic tail on the shell.  This section gives
an exact chain-rule estimate, its four-dimensional summability
criterion, and a neutral-block refinement for massless modes whose
absolute tails are too slow.

\subsection{Microscopic reconstruction of a soft direction}

Let the finite-cutoff microscopic field space be a product of
finite-dimensional normed spaces,
\[
 \mathcal X_\Lambda=\prod_{z\in\Lambda}E_z.         \tag{44.1}
\]
Let \(S\) be the augmented soft-coordinate space from V29--V31.  On a
soft good set \(G\), suppose the microscopic field is represented by
\[
 \phi(s,h)=\phi_H(h)+\mathcal R(s),                 \tag{44.2}
\]
where \(\mathcal R\) is differentiable.  For a carrier direction
\(a\in S\), write
\[
 e_a(s):=D\mathcal R(s)[a]\in\mathcal X_\Lambda.    \tag{44.3}
\]
The linear case has \(e_a=\mathcal Ra\), independent of \(s\).

Let a shell action be a sum of local cells,
\[
 V_q(s,h)=\sum_{y\in\mathcal S_q}
 v_{q,y}\bigl(\phi(s,h)|_{N_y}\bigr),               \tag{44.4}
\]
where \(N_y\) is the finite neighborhood on which the cell depends.
Equip \(\prod_{z\in N_y}E_z\) with a norm
\(\|\cdot\|_{N_y}\) and its dual.

\begin{hypothesis}[Local differentiability and integrability]
\label{hyp:v16v44-chain}
Every \(v_{q,y}\) is differentiable on the relevant field domain, the
finite sum (44.4) may be differentiated under
\(\mathbb E_{1,s}\), and there are deterministic numbers
\(g_{q,y},\rho_{q,y}\) such that, uniformly for \(s\in G\),
\[
 \mathbb E_{1,s}
 \|Dv_{q,y}(\phi|_{N_y})\|_{N_y,*}\leq g_{q,y},     \tag{44.5}
\]
\[
 \sup_{\|a\|_S=1}
 \|e_a(s)|_{N_y}\|_{N_y}\leq\rho_{q,y}.             \tag{44.6}
\]
\end{hypothesis}

\begin{theorem}[Direct-response chain-rule card]
\label{thm:v16v44-chain}
Under Hypothesis~\ref{hyp:v16v44-chain},
\[
 \partial_aV_q(s,h)
 =
 \sum_{y\in\mathcal S_q}
 Dv_{q,y}(\phi|_{N_y})
 [e_a(s)|_{N_y}],                                   \tag{44.7}
\]
and the direct shell response of (42.17) can be chosen as
\[
 e_q:=
 \sum_{y\in\mathcal S_q}g_{q,y}\rho_{q,y}.          \tag{44.8}
\]
No independence between shell cells is required.
\end{theorem}

\begin{proof}
The ordinary finite-dimensional chain rule applied to every summand
in (44.4) gives (44.7).  For \(\|a\|_S=1\), duality,
(44.5), and (44.6) give
\[
\begin{split}
 \left|\mathbb E_{1,s}\partial_aV_q\right|
 &\leq
 \sum_{y\in\mathcal S_q}
 \mathbb E_{1,s}\|Dv_{q,y}\|_{N_y,*}
 \|e_a(s)|_{N_y}\|_{N_y}\\
 &\leq\sum_{y\in\mathcal S_q}g_{q,y}\rho_{q,y}.
                                                               \tag{44.9}
\end{split}
\]
Take the supremum over \(s\in G\) and unit \(a\).
\end{proof}

\begin{corollary}[Uniform shell envelopes]
\label{cor:v16v44-envelope}
If
\[
 g_{q,y}\leq g_q,\qquad \rho_{q,y}\leq\rho_q
 \quad(y\in\mathcal S_q),                           \tag{44.10}
\]
then
\[
 e_q\leq|\mathcal S_q|g_q\rho_q.                   \tag{44.11}
\]
In four dimensions, under the owner count (37.17),
\[
 e_q\leq
 8Hr[2(R+(q+1)r)+1]^3g_q\rho_q.                   \tag{44.12}
\]
Thus the direct-response tail is summable if
\[
 \sum_{q=0}^{\infty}
 [2(R+(q+1)r)+1]^3g_q\rho_q<\infty.                \tag{44.13}
\]
\end{corollary}

\begin{proof}
Equations (44.11)--(44.12) follow from (44.8) and (37.17) with
\(d=4\).  Summing proves (44.13).
\end{proof}

\begin{corollary}[Exponential and power-law mode tails]
\label{cor:v16v44-tail}
Assume \(g_q\leq C_g(1+q)^{-\sigma}\).
\begin{enumerate}
\item If \(\rho_q\leq C_Re^{-\mu q}\) for some \(\mu>0\), then
(44.13) holds for every real \(\sigma\).
\item If \(\rho_q\leq C_R(1+q)^{-p}\), then (44.13) holds whenever
\[
 p+\sigma>4.                                       \tag{44.14}
\]
\end{enumerate}
\end{corollary}

\begin{proof}
The polynomial factor in (44.13) is \(O((1+q)^3)\).  Exponential
decay beats every polynomial.  In the second case the summand is
\(O((1+q)^{3-p-\sigma})\), whose series converges exactly under the
strict sufficient inequality (44.14).
\end{proof}

\subsection{Weighted reconstruction norm}

The preceding shellwise test follows from a single global operator
norm.  Give every local neighborhood a weight \(w_y>0\) and define
\[
 \|\mathcal R\|_{1,w;G}
 :=
 \sup_{\substack{s\in G\\\|a\|_S=1}}
 \sum_y w_y\|e_a(s)|_{N_y}\|_{N_y}.                 \tag{44.15}
\]

\begin{proposition}[Weighted direct-response compiler]
\label{prop:v16v44-weighted}
If
\[
 \|\mathcal R\|_{1,w;G}<\infty,\qquad
 \mathbb E_{1,s}\|Dv_{q,y}\|_{N_y,*}
 \leq\varepsilon_q w_y                              \tag{44.16}
\]
for \(y\in\mathcal S_q\), uniformly in \(s\in G\), then
\[
 e_q\leq\varepsilon_q\|\mathcal R\|_{1,w;G}.        \tag{44.17}
\]
Hence \(\sum_q\varepsilon_q<\infty\) gives a summable direct-response
tail without a separate shell-cardinality estimate.
\end{proposition}

\begin{proof}
Insert (44.16) directly in (44.9) and enlarge the shell sum to the
full weighted sum in (44.15).
\end{proof}

\subsection{Neutral-block cancellation}

Absolute mode amplitudes can be too slow for a massless carrier.
Cancellation must then be used before absolute values are taken.
For clarity, take scalar microscopic fields; the vector-valued
statement follows by dual pairing.  Suppose the direct response has
the form
\[
 \mathbb E_{1,s}\partial_aV_q
 =\sum_{y\in\mathcal S_q}j_{q,y}(s)e_a(s;y).        \tag{44.18}
\]
Partition \(\mathcal S_q\) into finite blocks \(\mathcal B_q\), and
choose a representative \(y_B\in B\).

\begin{lemma}[Neutral-block subtraction]
\label{lem:v16v44-neutral}
If
\[
 \sum_{y\in B}j_{q,y}(s)=0
 \quad(B\in\mathcal B_q,\ s\in G),                  \tag{44.19}
\]
then
\[
 \mathbb E_{1,s}\partial_aV_q
 =
 \sum_{B\in\mathcal B_q}\sum_{y\in B}
 j_{q,y}(s)
 \bigl(e_a(s;y)-e_a(s;y_B)\bigr).                  \tag{44.20}
\]
If \(e_a(s;\cdot)\) is Lipschitz on each block, this implies
\[
 \left|\mathbb E_{1,s}\partial_aV_q\right|
 \leq
 \sum_{B\in\mathcal B_q}
 \operatorname{diam}(B)
 \sup_{z\in B}\|\nabla e_a(s;z)\|
 \sum_{y\in B}|j_{q,y}(s)|.                         \tag{44.21}
\]
\end{lemma}

\begin{proof}
For each \(B\), subtract
\(e_a(s;y_B)\sum_{y\in B}j_{q,y}(s)\), which is zero by
(44.19).  This proves (44.20).  The mean-value estimate
\[
 |e_a(s;y)-e_a(s;y_B)|
 \leq\operatorname{diam}(B)
 \sup_{z\in B}\|\nabla e_a(s;z)\|                  \tag{44.22}
\]
and the triangle inequality give (44.21).
\end{proof}

\begin{corollary}[One extra decay power from neutrality]
\label{cor:v16v44-extra-power}
Suppose block diameters and
\(\sum_{y\in B}|j_{q,y}(s)|\) are uniformly bounded, the number of
blocks in \(\mathcal S_q\) is \(O((1+q)^3)\), and
\[
 \sup_{\substack{s\in G,\ \|a\|_S=1\\z\in\mathcal S_q}}
 \|\nabla e_a(s;z)\|
 \leq C(1+q)^{-p-1}.                                \tag{44.23}
\]
Then
\[
 e_q=O((1+q)^{2-p}),                                \tag{44.24}
\]
and \(\sum_qe_q<\infty\) whenever \(p>3\).
\end{corollary}

\begin{proof}
Equation (44.21) gives one factor
\((1+q)^{-p-1}\) per block and there are
\(O((1+q)^3)\) blocks.  This is (44.24), whose series converges for
\(2-p<-1\), namely \(p>3\).
\end{proof}

\begin{remark}[Higher multipole cancellation]
\label{rem:v16v44-multipole}
If the response coefficients have vanishing moments through order
\(k\) on each block, Taylor subtraction replaces the mode amplitude
by a derivative of order \(k+1\) times a block-diameter factor of the
same order.  Each proved moment cancellation can therefore gain one
additional far-field power.  Such cancellations must come from an
exact Ward, charge-neutrality, or projected-source identity; they
cannot be inferred from decay alone.
\end{remark}

\subsection{Closing the V42 shell-gradient ledger}

\begin{theorem}[Direct plus connected shell certificate]
\label{thm:v16v44-combined}
Assume the uniform interpolation hypotheses of
Theorem~\ref{thm:v16v43-certificate}.  Suppose, in addition, that
the direct term is controlled either by (44.8), (44.17), or the
neutral-block estimate (44.21), producing a sequence
\((e_q)\) with
\[
 \sum_q e_q<\infty.                                 \tag{44.25}
\]
Then the gradient bounds \(L_q\) in (42.23) are summable in four
dimensions.  On carrier good sets of uniformly bounded diameter, and
with the summable bad-set schedule (42.26), the remote effective
soft-owner tail converges in state norm.
\end{theorem}

\begin{proof}
V43 supplies (42.16) with a positive exponential rate.  Therefore
the covariance part of \(L_q\) is summable by
Corollary~\ref{cor:v16v42-summability}.  Add (44.25) and apply
Theorem~\ref{thm:v16v42-tail}.
\end{proof}

\subsection{Physical obstruction exposed by the power test}

For a four-dimensional shell and bounded local derivative costs,
an absolute reconstruction tail \(\rho_q=O(q^{-p})\) requires
\(p>4\).  Neutrality improves the sufficient threshold to \(p>3\)
under the block hypotheses above.  Coulombic or Newtonian modes can
still be slower.  For them the absolute carrier-mode route is
insufficient, and the projected-density or multipole cancellations
of V14--V24 must be used at the level of the actual source.  This
explains why the Newton owner was treated separately in V34 instead
of being forced into the finite-range hard-sector ledger.

\begin{assessment}[What V44 closes]
\label{ass:v16v44-status}
V44 pays the direct term left open by V42 whenever the reconstruction
mode has a summable weighted tail, or when exact neutral-block
cancellation supplies enough additional decay.  Together V42--V44
now form a complete conditional shell free-energy compiler:
interpolation, connected covariance, unbounded truncation, direct
mode response, and countable assembly.

The compiler's physical premises remain unpaid sector by sector:
uniform positive hard Dobrushin reserve, insertion moments, local
derivative estimates, and any required Ward or neutrality identities.
\end{assessment}

\begin{decision}[V45 companion audit and physical certificate matrix]
\label{dec:v16v44-next}
Perform the mandatory companion audit.  Then build a nonduplicative
certificate matrix for the gauge, Higgs, fermion, counterterm, and
metric-boundary owners: identify for each owner which of positivity,
finite range, Dobrushin reserve, exponential moments, direct
reconstruction tail, and neutrality is already proved in the active
or archived notes.  Select the first genuinely unpaid row that can be
attacked without assuming a mass gap.
\end{decision}

\begin{firewall}[Claim boundary after V44]
\label{fire:v16v44}
The chain-rule and tail estimates do not prove the required physical
mode decay or neutral-block identities.  A global massless mode
cannot be declared summable from finite interaction range.  V44 does
not resolve a fermion phase, construct the continuum state, recover
the Einstein equation, or complete the Standard Model--Einstein
continuum.
\end{firewall}

\section*{Version V44 append-only record}
\addcontentsline{toc}{section}{Version V44 append-only record}
The complete V43 source was retained before this continuation.  It
had \(422925\) bytes, \(12788\) lines, and SHA--256
\[
 \texttt{837250B8365D59302CE4CB569F103AA69C7C1B8C82B6FE21C1462D46B4582FE9}.
\]
V44 proved the direct reconstruction-mode shell estimate, its
weighted and four-dimensional summability forms, and the
neutral-block improvement needed before absolute mode tails are used.

\section{V45: Companion audit and the physical owner certificate
matrix}
\label{sec:v16v45}

\subsection{Mandatory companion audit}

The newest Yang--Mills companion inspected at this checkpoint was
\texttt{Renewal\_Yang\_Mills\_Mass\_Gap\_Research\_Notes\_V51.tex};
it had \(442249\) bytes, \(12072\) lines, and SHA--256
\[
 \texttt{9D287E271C58C44FAEC22B89B3602DF2FC9AEF6C609976D7D72E8F122A98143B}.
                                                               \tag{45.1}
\]
It constructs two de-duplicated mixed-sign power-three charts and
uses them to prove one direct plus-pencil quarter-face plate, with
the paired minus result under inversion.  The other two transverse
coordinates remain only one rational step wide; both signs on the
same plate, a mixed-sign quarter cell, all octants, full-torus
positivity, the interacting continuum measure, and a mass gap remain
open.

The Navier--Stokes companion remained
\texttt{Renewal\_NS\_Stability\_Research\_Notes\_V26.tex}; it had
\(278283\) bytes, \(5319\) lines, and SHA--256
\[
 \texttt{82C887F8C2C69E4F28FAD7C3DC8DE5B9099CB5A4BF431EBA1FFF3E91935978AD}.
                                                               \tag{45.2}
\]
Its rank-one Oseen derivative calculation remains outside the present
measure, free-energy, and Einstein-source problem.

\begin{theorem}[V45 companion-audit decision]
\label{thm:v16v45-audit}
No theorem from the two active companions supplies a missing premise
of the local-compatibility compiler V35--V44.
\end{theorem}

\begin{proof}
YM V51 is an exact finite matrix-symbol result on one proper momentum
plate.  It does not establish a positive interacting quotient
measure, hard-sector mixing, carrier-mode decay, or a shell density
ratio.  NS V26 has neither the configuration space nor the operator
required here.  Therefore importing either result would change the
hypotheses rather than discharge one of them.
\end{proof}

\subsection{Archive de-duplication record}

A targeted archive search was made before selecting the next row.
The following results already exist and will be used only by explicit
reference to their retained hypotheses:
\begin{enumerate}
\item seventh-cycle V89 proves bounded-observable Dobrushin covariance
and exponential resolvent decay;
\item ninth-cycle V4 proves dependent-reference Dobrushin comparison
and its Doob-range theorem;
\item ninth-cycle V49--V52 prove quartic Higgs tails, the
worst-exterior no-go, and an abstract averaged-disagreement path
compiler;
\item ninth-cycle V53--V54 prove phase-quenched fermion transfer,
the unavoidable average-phase denominator, and a connected local
phase criterion;
\item fifteenth-cycle V46--V47 prove adjacent-cutoff entropy
comparison and the action-increment variance identity;
\item fifteenth-cycle V61--V72 prove cutoff and cofinal
Einstein Schwinger--Dyson identities and their score-moment
consequences.
\end{enumerate}
Accordingly V45 does not reopen any of these abstract arguments.

\subsection{Meaning of the certificate columns}

For an owner to enter the V35 probability ledger through the
V39--V44 soft/hard route, the following certificates have distinct
roles:
\[
\begin{array}{c|p{0.68\textwidth}}
\text{column}&\text{required statement}\\ \hline
P&\text{a normalized positive density or positive trace-class state}\\
L&\text{finite-range locality, or an explicit replacement for it}\\
M&\text{hard conditional mixing or an averaged path estimate}\\
E&\text{uniform exponential or other sufficient insertion moments}\\
R&\text{summable reconstruction tail or exact neutrality}\\
C&\text{cofinal coefficient, normalization, and source compatibility}.
\end{array}                                         \tag{45.3}
\]
No column may be inferred from another.  In particular, a positive
local action does not prove mixing, and anomaly cancellation does not
prove \(P\) for a chiral determinant.

\subsection{Sector certificate matrix}

\begin{ledger}[Physical owner status at V45]
\label{led:v16v45-matrix}
The proved and unpaid rows are as follows.
\[
\begin{array}{p{0.12\textwidth}|p{0.36\textwidth}|p{0.39\textwidth}}
\text{owner}&\text{proved input}&\text{unpaid input}\\ \hline
\text{Newton}&
\text{positive Gaussian seed; nonlinear density transfer; projected
far-influence moments; compatible local states}&
\text{full coupled uniformity under the other owner tilts and the
nonlinear metric boundary law}\\
\text{gauge}&
\text{local finite-range bosonic action; regular-stratum quotient
density; conditional cancellation}&
\text{global slice atlas, reducible interfaces, full momentum
coverage, uniform mixing, and reflection-compatible quotient}\\
\text{Higgs}&
\text{stable finite-cutoff quartic; strong local moments; abstract
averaged-disagreement route}&
\text{physical branching constants, nested exterior control,
renormalized cutoff convergence, and coupled reflection passage}\\
\text{fermion}&
\text{one-particle/Fock heat compilers; modulus and phase transfer
cards; local determinant estimates}&
\text{positive chiral realization or noncollapsing quasi-local phase,
zero-set control, and reflection sewing}\\
\text{counterterm}&
\text{local polynomial support; owner telescoping; vacuum shifts
cancel from normalized fixed-metric states}&
\text{cofinal coefficient rates, unbounded-field moments, shell
response, and every metric derivative in the Einstein source}\\
\text{metric boundary}&
\text{TT heat nuclearity; Newton projected moments; moving soft-space
and local-state compilers}&
\text{nonlinear conformal stability, positive physical transfer,
full boundary density, and limiting Einstein-source identification}.
\end{array}                                         \tag{45.4}
\]
\end{ledger}

\begin{proof}
The Newton row is V4--V34 of the current cycle.  The assembly columns
are V35--V44.  The gauge, Higgs, fermion, and gravitational entries are
the corresponding established results and firewalls in the archived
cycles summarized in V1.  Each unpaid entry is retained because no
active or archived theorem found in the targeted search supplies it.
The companion audit above supplies no additional row.
\end{proof}

\begin{proposition}[No physical owner is complete]
\label{prop:v16v45-incomplete}
The current evidence does not prove the full local-compatibility
estimate
\[
 \|\mu_{n,K}^{\rm large}-\mu_{n,K}^{\rm small}\|_{\rm st}\to0
                                                               \tag{45.5}
\]
for the coupled Standard Model--Einstein state.
\end{proposition}

\begin{proof}
Corollary~\ref{cor:v16v35-ledger} requires every owner debit to
vanish.  Ledger~\ref{led:v16v45-matrix} contains unpaid entries in
each non-Newton row, and the Newton result has not been shown uniform
under their tilts.  Hence the antecedent of the six-owner theorem is
not established.
\end{proof}

\subsection{Selecting a noncircular next row}

The gauge and Higgs mixing rows cannot be acquired merely by assuming
a mass gap: the continuum programme is meant to derive its spectral
properties, and archived ninth-cycle V50 rules out a worst-exterior
Higgs Dobrushin coefficient.  The fermion row requires a
regulator-specific phase construction not supplied by the attached
notes.  The full conformal gravity row is upstream but needs a new
physical reference, not another abstract loader.

One counterterm subrow can be attacked directly.  A metric-dependent
vacuum or cosmological counterterm is constant in the matter fields,
so it cancels from a normalized state at one fixed metric but remains
in the effective soft density and in its Einstein derivative.  Its
shell response depends only on the derivative of the volume density
and the reconstructed metric mode.  It therefore needs neither a
matter mass gap nor a hard-sector Dobrushin estimate.

\begin{decision}[V46 metric vacuum-counterterm shell]
\label{dec:v16v45-next}
For uniformly elliptic positive metric cells, prove the first and
second variation formulas for \(\sqrt{\det g}\), obtain explicit norm
bounds, and load a shell cosmological counterterm into the V44
reconstruction-tail criterion.  State the coefficient and mode-tail
summability required in four dimensions, and retain the resulting
metric derivative in the Einstein-source ledger even though the
scalar factor cancels from fixed-metric normalized matter states.
\end{decision}

\begin{assessment}[What V45 establishes]
\label{ass:v16v45-status}
V45 records the required companion audit, prevents duplication of six
substantial archived argument blocks, and replaces an undifferentiated
list of open sectors by a certificate matrix matched to the current
owner compiler.  It selects a concrete unpaid counterterm row whose
analysis does not presuppose a mass gap.
\end{assessment}

\begin{firewall}[Claim boundary after V45]
\label{fire:v16v45}
The matrix is an evidence ledger, not a proof of its unpaid entries.
YM V51's plate does not imply a full gauge-sector certificate.
No coupled continuum state, reflection-positive gravitational
transfer, Einstein equation, or full Standard Model--Einstein
continuum has been proved.
\end{firewall}

\section*{Version V45 append-only record}
\addcontentsline{toc}{section}{Version V45 append-only record}
The complete V44 source was retained before this continuation.  It
had \(434519\) bytes, \(13129\) lines, and SHA--256
\[
 \texttt{1420724EC1B08174C2AE3B61602C038C10C39DBFED0E1EA8461BE16C795E47CA}.
\]
V45 performed the mandatory companion audit, recorded the archive
de-duplication decisions, built the physical owner certificate
matrix, and selected the metric vacuum-counterterm shell for V46.

\section{V46: Metric vacuum-counterterm shell response}
\label{sec:v16v46}

A vacuum counterterm is constant in the matter variables at fixed
metric, but it is not constant on the metric carrier.  It therefore
cancels from a normalized conditional matter state while changing the
soft metric law and the Einstein source.  This section proves its
variation and shell debit on a uniformly elliptic metric good set.

\subsection{Variation of the volume density}

Let \({\rm Sym}_d^+\) be the positive definite real symmetric
\(d\times d\) matrices, equipped with the Frobenius norm.  Put
\[
 \upsilon(g):=(\det g)^{1/2}.                        \tag{46.1}
\]

\begin{theorem}[First and second volume-density variations]
\label{thm:v16v46-volume}
For \(g\in{\rm Sym}_d^+\) and symmetric directions \(h,k\),
\[
 D\upsilon(g)[h]
 ={1\over2}\upsilon(g)\operatorname{tr}(g^{-1}h),   \tag{46.2}
\]
and
\[
\begin{split}
 D^2\upsilon(g)[h,k]
 =\upsilon(g)\bigg[
 &{1\over4}\operatorname{tr}(g^{-1}h)
                 \operatorname{tr}(g^{-1}k)\\
 &-{1\over2}\operatorname{tr}(g^{-1}h\,g^{-1}k)
 \bigg].                                            \tag{46.3}
\end{split}
\]
\end{theorem}

\begin{proof}
Jacobi's formula gives
\[
 D\log\det g[h]=\operatorname{tr}(g^{-1}h).         \tag{46.4}
\]
Since \(\log\upsilon=\frac12\log\det g\), the ordinary chain rule
proves (46.2).  Differentiate once more, using
\[
 D(g^{-1})[k]=-g^{-1}kg^{-1}.                       \tag{46.5}
\]
The derivative of the prefactor in (46.2) produces the first term in
(46.3), and the derivative of the trace produces the second.
\end{proof}

\begin{proposition}[Uniform ellipticity bounds]
\label{prop:v16v46-bounds}
If
\[
 mI\preceq g\preceq MI,\qquad 0<m\leq M<\infty,    \tag{46.6}
\]
then
\[
 |D\upsilon(g)[h]|
 \leq C_1\|h\|_F,\qquad
 C_1:={\sqrt d\over2}M^{d/2}m^{-1},                \tag{46.7}
\]
and
\[
 |D^2\upsilon(g)[h,k]|
 \leq C_2\|h\|_F\|k\|_F,\qquad
 C_2:={d+2\over4}M^{d/2}m^{-2}.                    \tag{46.8}
\]
\end{proposition}

\begin{proof}
The eigenvalue bounds give
\[
 \upsilon(g)\leq M^{d/2},\qquad
 \|g^{-1}\|_{\rm op}\leq m^{-1}.                   \tag{46.9}
\]
Also
\[
 |\operatorname{tr}(g^{-1}h)|
 \leq\|g^{-1}\|_F\|h\|_F
 \leq\sqrt d\,m^{-1}\|h\|_F.                       \tag{46.10}
\]
Insert this in (46.2).  For the mixed trace,
\[
 |\operatorname{tr}(g^{-1}h\,g^{-1}k)|
 \leq m^{-2}\|h\|_F\|k\|_F.                        \tag{46.11}
\]
The two coefficients in (46.3) are therefore bounded by
\(d/4\) and \(1/2\), respectively, times
\(M^{d/2}m^{-2}\).  Their sum is (46.8).
\end{proof}

\begin{corollary}[Lipschitz volume density on an elliptic cell]
\label{cor:v16v46-lipschitz}
The set
\[
 \mathcal G_{m,M}:=\{g:mI\preceq g\preceq MI\}      \tag{46.12}
\]
is convex, and for \(g_0,g_1\in\mathcal G_{m,M}\),
\[
 |\upsilon(g_1)-\upsilon(g_0)|
 \leq C_1\|g_1-g_0\|_F.                            \tag{46.13}
\]
\end{corollary}

\begin{proof}
Both Loewner inequalities are preserved by convex combinations.
Integrate (46.7) along the segment from \(g_0\) to \(g_1\).
\end{proof}

\subsection{A lattice vacuum owner}

At lattice spacing \(a_n\), let \(g_{n,y}(s)\in{\rm Sym}_d^+\) be the
metric cell attached to \(y\), as a differentiable function of the
soft carrier \(s\).  Let a shell vacuum-counterterm increment be
\[
 V^{\rm vac}_{n,q}(s)
 :=
 a_n^d\sum_{y\in\mathcal S_{n,q}}
 \delta\Lambda_{n,q,y}\,
 \upsilon(g_{n,y}(s)),                              \tag{46.14}
\]
where every \(\delta\Lambda_{n,q,y}\) is real.  This notation includes
both an actual cosmological counterterm increment and the scalar
vacuum energy left after integrating a matter sector, provided the
latter has already been proved local in the form (46.14).

\begin{hypothesis}[Metric reconstruction good set]
\label{hyp:v16v46-reconstruction}
There is a convex carrier event \(G_n\) on which
\[
 mI\preceq g_{n,y}(s)\preceq MI                    \tag{46.15}
\]
for every relevant \(y\), with cutoff-independent \(m,M\).  For a
unit carrier direction \(b\), suppose
\[
 \|D_sg_{n,y}(s)[b]\|_F\leq\rho_{n,q,y}             \tag{46.16}
\]
uniformly for \(s\in G_n\).
\end{hypothesis}

\begin{theorem}[Vacuum-shell gradient]
\label{thm:v16v46-gradient}
Under Hypothesis~\ref{hyp:v16v46-reconstruction},
\[
 \sup_{s\in G_n}
 \|\nabla_sV^{\rm vac}_{n,q}(s)\|_*
 \leq
 E_{n,q}^{\rm vac},                                 \tag{46.17}
\]
where
\[
 E_{n,q}^{\rm vac}
 :=
 C_1a_n^d
 \sum_{y\in\mathcal S_{n,q}}
 |\delta\Lambda_{n,q,y}|\,\rho_{n,q,y}.             \tag{46.18}
\]
\end{theorem}

\begin{proof}
For a unit direction \(b\), the chain rule and (46.2) give
\[
\begin{split}
 D_sV^{\rm vac}_{n,q}(s)[b]
 =a_n^d\sum_y\delta\Lambda_{n,q,y}
 D\upsilon(g_{n,y}(s))
 [D_sg_{n,y}(s)[b]].                                \tag{46.19}
\end{split}
\]
Use (46.7), (46.16), and the triangle inequality, then take the
supremum over unit \(b\).
\end{proof}

\begin{proposition}[The vacuum shell is a pure soft owner]
\label{prop:v16v46-pure-soft}
Let a positive joint state disintegrate over \(s\) as in (39.1).
Tilting it by \(\exp[-V^{\rm vac}_{n,q}(s)]\) leaves every conditional
hard law unchanged.  Its induced soft law is exactly the normalized
tilt by that factor, and the covariance term in (42.9) is zero.
\end{proposition}

\begin{proof}
In Theorem~\ref{thm:v16v40-effective}, take
\(W(s,h)=e^{-V^{\rm vac}_{n,q}(s)}\).  Then
\[
 Z_H(s)=e^{-V^{\rm vac}_{n,q}(s)},\qquad
 {W(s,h)\over Z_H(s)}=1.                            \tag{46.20}
\]
Equation (40.7) therefore leaves the hard conditional law unchanged,
while (40.6) gives the asserted soft tilt.  At fixed \(s\),
\(V^{\rm vac}_{n,q}(s)\) is constant in \(h\), so its covariance with
any hard observable is zero.
\end{proof}

\subsection{State-distance debit}

\begin{theorem}[Metric vacuum-owner compatibility]
\label{thm:v16v46-debit}
Let two consecutive soft laws be related by the vacuum-shell tilt of
Proposition~\ref{prop:v16v46-pure-soft}.  Suppose \(G_n\) has carrier
diameter \(D_n\), and both laws assign probability at most
\(\eta_{n,q}\) to \(G_n^c\).  Then
\[
 d_{n,q}^{\rm vac}
 \leq
 e^{D_nE_{n,q}^{\rm vac}}-1+4\eta_{n,q}             \tag{46.21}
\]
is a certified V35 owner debit.
\end{theorem}

\begin{proof}
Theorem~\ref{thm:v16v41-path} and (46.17) give
\[
 \operatorname{osc}_{G_n}V^{\rm vac}_{n,q}
 \leq D_nE_{n,q}^{\rm vac}.                         \tag{46.22}
\]
Apply the oscillation loader to the two normalized restrictions on
\(G_n\), then Lemma~\ref{lem:v16v34-good}.
\end{proof}

\begin{corollary}[Countable vacuum-shell tail]
\label{cor:v16v46-countable}
If
\[
 \sup_nD_n<\infty,\qquad
 \lim_{Q\to\infty}\sup_n\sum_{q\geq Q}
 E_{n,q}^{\rm vac}=0,\qquad
 \lim_{Q\to\infty}\sup_n\sum_{q\geq Q}
 \eta_{n,q}=0,                                     \tag{46.23}
\]
then the countable remote vacuum-counterterm owner has a vanishing
uniform state-distance tail.
\end{corollary}

\begin{proof}
The first two conditions imply
\(D_nE_{n,q}^{\rm vac}\to0\) uniformly in the tail.  There,
\(e^x-1\leq2x\).  Sum (46.21), use (46.23), and apply
Theorem~\ref{thm:v16v36-triangular}.
\end{proof}

\subsection{Four-dimensional coefficient--mode schedule}

Assume now \(d=4\), and the shell count of (37.17).  Suppose
\[
 |\delta\Lambda_{n,q,y}|\leq\ell_{n,q},\qquad
 \rho_{n,q,y}\leq\rho_{n,q}.                        \tag{46.24}
\]
In four dimensions, (46.7) becomes
\[
 C_1=M^2m^{-1}.                                     \tag{46.25}
\]

\begin{proposition}[Four-dimensional vacuum-shell bound]
\label{prop:v16v46-four}
Under (46.24),
\[
 E_{n,q}^{\rm vac}
 \leq
 {8HrM^2\over m}\,
 a_n^4
 [2(R_n+(q+1)r)+1]^3
 \ell_{n,q}\rho_{n,q}.                              \tag{46.26}
\]
Consequently a sufficient cofinal tail condition is
\[
\lim_{Q\to\infty}\sup_n
\sum_{q\geq Q}
a_n^4[2(R_n+(q+1)r)+1]^3
\ell_{n,q}\rho_{n,q}=0.                             \tag{46.27}
\]
\end{proposition}

\begin{proof}
Use (46.18), the uniform envelopes (46.24), the four-dimensional
owner count
\[
 |\mathcal S_{n,q}|
 \leq8Hr[2(R_n+(q+1)r)+1]^3,                        \tag{46.28}
\]
and (46.25).  The constant factor in (46.26) does not affect the
vanishing tail in (46.27).
\end{proof}

\begin{corollary}[Power-law diagnostic]
\label{cor:v16v46-power}
At fixed \(n\), if
\[
 \ell_{n,q}=O((1+q)^{-\sigma}),\qquad
 \rho_{n,q}=O((1+q)^{-p}),                          \tag{46.29}
\]
then the shell sum in (46.27) converges whenever
\[
 p+\sigma>4.                                       \tag{46.30}
\]
The cutoff-uniform statement still requires the supremum in
(46.27), including the factor \(a_n^4\) and the cutoff dependence of
all constants.
\end{corollary}

\begin{proof}
The \(q\)-th summand is
\(O((1+q)^{3-p-\sigma})\).  Its series converges under (46.30).
\end{proof}

\subsection{Second carrier derivative}

For the Einstein response it is useful to control not only the
first variation but its continuity.  Suppose, in addition to
(46.16), that
\[
 \|D_s^2g_{n,y}(s)[b,c]\|_F\leq\sigma_{n,q,y}
 \quad(\|b\|_S=\|c\|_S=1).                          \tag{46.31}
\]

\begin{proposition}[Vacuum-shell Hessian]
\label{prop:v16v46-hessian}
On \(G_n\),
\[
\begin{split}
 \|\nabla_s^2V^{\rm vac}_{n,q}(s)\|_{\rm op}
 \leq a_n^d\sum_{y\in\mathcal S_{n,q}}
 |\delta\Lambda_{n,q,y}|
 \left(C_2\rho_{n,q,y}^2+C_1\sigma_{n,q,y}\right).
                                                               \tag{46.32}
\end{split}
\]
\end{proposition}

\begin{proof}
The second-order chain rule gives
\[
\begin{split}
 D_s^2(\upsilon\circ g_y)[b,c]
 &=
 D^2\upsilon(g_y)[D_sg_y[b],D_sg_y[c]]\\
 &\quad+D\upsilon(g_y)[D_s^2g_y[b,c]].              \tag{46.33}
\end{split}
\]
Apply (46.7)--(46.8), (46.16), and (46.31), multiply by the
coefficients in (46.14), and sum.
\end{proof}

\subsection{Einstein-source retention}

For one continuum cell, adopt the source convention
\[
 \mathcal T^{\mu\nu}_{\rm vac}(g)
 :={2\over\sqrt{\det g}}\,
 {\partial\over\partial g_{\mu\nu}}
 \left(\Lambda\sqrt{\det g}\right).                 \tag{46.34}
\]

\begin{proposition}[Vacuum Einstein source]
\label{prop:v16v46-source}
With the convention (46.34),
\[
 \mathcal T^{\mu\nu}_{\rm vac}(g)
 =\Lambda g^{\mu\nu}.                               \tag{46.35}
\]
Thus deleting the scalar vacuum factor from a normalized matter
state at fixed \(g\) does not license deleting (46.35) from the
metric-source equation.
\end{proposition}

\begin{proof}
Equation (46.2), in coordinate form, reads
\[
 {\partial\sqrt{\det g}\over\partial g_{\mu\nu}}
 ={1\over2}\sqrt{\det g}\,g^{\mu\nu}.               \tag{46.36}
\]
Insert it in (46.34).  The final statement is also the
vacuum-source identity (1.8): normalization cancels a fixed scalar
factor, whereas metric differentiation acts before that cancellation
is discarded from the generating functional.
\end{proof}

\begin{assessment}[What V46 closes]
\label{ass:v16v46-status}
V46 pays one concrete part of the counterterm row in the V45 matrix.
A local real vacuum or cosmological counterterm is a pure soft owner;
its exact debit is governed by the cofinal coefficient--metric-mode
tail (46.27), with no hard mixing or matter mass-gap premise.  Its
first and second metric responses are controlled on uniformly
elliptic good sets, and its Einstein source is retained explicitly.

The theorem does not prove (46.27) for the regulator coefficients or
metric reconstruction.  Other curvature, Higgs, gauge, determinant,
and boundary counterterms are not scalar volume terms and need
separate variation estimates.
\end{assessment}

\begin{decision}[V47 curvature counterterm variation]
\label{dec:v16v46-next}
Treat the next geometric counterterm,
\(\int\sqrt g\,R(g)\), on a finite metric collar.  Derive its first
variation including the boundary divergence, isolate the Einstein
tensor bulk term and the boundary owner, and state the discrete
summability norms needed to load each into V35.  Do not suppress the
boundary term unless a proved compact-support or Gibbons--Hawking
cancellation hypothesis applies.
\end{decision}

\begin{firewall}[Claim boundary after V46]
\label{fire:v16v46}
Uniform ellipticity and the cofinal tail (46.27) are hypotheses.
The result covers a real scalar volume counterterm, not the full
renormalized action or a complex determinant contribution.  It does
not prove the physical gravitational measure, the limiting Einstein
equation, or the full Standard Model--Einstein continuum.
\end{firewall}

\section*{Version V46 append-only record}
\addcontentsline{toc}{section}{Version V46 append-only record}
The complete V45 source was retained before this continuation.  It
had \(444118\) bytes, \(13347\) lines, and SHA--256
\[
 \texttt{2F32191320CC077182564FECE19E25E3ED473B0C03F366D5EFBAC41D4B25A1D9}.
\]
V46 proved the metric volume-density variations and loaded a vacuum
counterterm shell into the soft owner ledger with its exact
four-dimensional coefficient--mode tail and Einstein source.

\section{V47: Curvature-counterterm owner norms from the archived
Einstein--GHY variation}
\label{sec:v16v47}

The V46 decision initially proposed rederiving the first variation of
\(\int\sqrt g\,R\).  An archive search found that this would duplicate
completed work.  The early-cycle file
\texttt{Renewal\_SM\_Einstein\_Continuum\_Research\_Notes\_V157\_f1.tex}
has \(3623989\) bytes, \(98871\) lines, and SHA--256
\[
 \texttt{EA08CE7537C7327A80B5AA5A653111AFAEDAE82CC18215D87534F17B563AAD94}.
                                                               \tag{47.1}
\]
Its V73 proves torsion-free Palatini elimination, the exact
Einstein--GHY smooth-face variation, the mismatch between free
connection and GHY polarizations, and a compatible-jet repair.  Its
V74 treats the quadratic joint and the residual corner extension.
Those identities are imported with their stated orientation and
boundary hypotheses.  The new result below is the analytic owner
norm which was not supplied there: it converts the bulk, face, and
joint variations into a state-distance debit for the current
soft-carrier ledger.

\subsection{Imported variational input}

Let \((\Omega,g)\) be a compact \(d\)-dimensional Riemannian region.
For a smooth covariant metric variation \(h\), the standard
Einstein--Hilbert variation in the convention used here is
\[
\begin{split}
 D S_{\rm EH}(g)[h]
 &=-c_R\int_\Omega\langle G_g,h\rangle_g\,dV_g\\
 &\quad+c_R\int_{\partial\Omega}
 B_g(h)\,dA_g,                                      \tag{47.2}\\
 B_g(h)
 &:=
 n^\mu\bigl(\nabla^\nu h_{\mu\nu}
             -\nabla_\mu{\rm tr}_g h\bigr).
                                                               \tag{47.3}
\end{split}
\]
For the orientation-compatible GHY term, the normal-derivative
expression \(B_g(h)\) cancels and the paired smooth-face variation
has the form
\[
 D S_{\rm EGHY}(g)[h]
 =
 -c_R\int_\Omega\langle G_g,h\rangle_g\,dV_g
 +c_R\int_{\partial\Omega}
 \langle\Pi_g,h^\top\rangle_\gamma\,dA_g.            \tag{47.4}
\]
Here \(\Pi_g\) is the oriented Brown--York momentum in the convention
fixed by the imported archive.  On a piecewise smooth region a joint
functional must also be included:
\[
 D S_{\rm full}(g)[h]
 =D S_{\rm EGHY}(g)[h]
 +c_R\int_{\mathcal J}\mathcal C_g(h)\,dA_{\mathcal J}.          \tag{47.5}
\]
Equations (47.2)--(47.5) are recorded inputs, not new derivations.
The archive expressly forbids dropping the face or joint term without
a compatible boundary hypothesis.

\subsection{Good-set geometric bounds}

Let \(G\) be a carrier good set.  Assume that for every \(s\in G\)
the reconstructed metric \(g(s)\) satisfies the following uniform
pointwise estimates:
\[
\begin{gathered}
 |G_{g(s)}|_g\leq K_G,\qquad
 |\Pi_{g(s)}|_\gamma\leq K_\Pi,\\
 |\mathcal C_{g(s)}(h)|\leq K_{\mathcal J}|h|_g,\qquad
 |B_{g(s)}(h)|\leq K_B|\nabla h|_g,                 \tag{47.6}
\end{gathered}
\]
where the last inequality is needed only for an unpaired
Einstein--Hilbert owner.  Uniform ellipticity converts any fixed
coordinate norms and measures into the geometric norms above with
cutoff-independent comparison constants.

\begin{theorem}[Bulk--face--joint variation norm]
\label{thm:v16v47-owner-norm}
Under (47.6), the paired action obeys
\[
\begin{split}
 |D S_{\rm full}(g)[h]|
 \leq |c_R|\bigl(
 &K_G\|h\|_{L^1(\Omega,g)}
 +K_\Pi\|h^\top\|_{L^1(\partial\Omega,\gamma)}\\
 &+K_{\mathcal J}\|h\|_{L^1(\mathcal J)}
 \bigr).                                            \tag{47.7}
\end{split}
\]
For the unpaired Einstein--Hilbert action one instead has
\[
\begin{split}
 |D S_{\rm EH}(g)[h]|
 \leq |c_R|\bigl(
 K_G\|h\|_{L^1(\Omega,g)}
 +K_B\|\nabla h\|_{L^1(\partial\Omega,\gamma)}
 \bigr).                                            \tag{47.8}
\end{split}
\]
Thus pairing with GHY removes the boundary normal-derivative norm,
but it does not remove the Brown--York face norm or the joint norm.
\end{theorem}

\begin{proof}
Apply the pointwise Cauchy--Schwarz inequality to the three integrals
in (47.4)--(47.5), insert (47.6), and integrate absolute values.  This
gives (47.7).  Apply the same argument to (47.2)--(47.3) to obtain
(47.8).  The final statement is the literal difference between those
two bounds.
\end{proof}

\begin{remark}[When a boundary term may vanish]
\label{rem:v16v47-boundary}
If \(h\) has compact support in the interior, all boundary and joint
norms vanish.  If the induced boundary metric is fixed, the
Brown--York pairing in (47.4) vanishes.  These are hypotheses on the
variation.  A remote shell reconstruction mode need not satisfy
either one, so its boundary owner must remain in the ledger.
\end{remark}

\subsection{Discrete shell reconstruction norms}

Let \(\Omega_{n,q}\) be a finite metric shell with bulk cells
\(\mathcal B_{n,q}\), oriented face cells \(\mathcal F_{n,q}\), and
joint cells \(\mathcal J_{n,q}\).  For a unit soft direction \(a\),
let
\[
 h_{n,a}(s):=D_sg_n(s)[a].                           \tag{47.9}
\]
Define the uniform reconstructed shell norms
\[
\begin{split}
 \mathfrak B_{n,q}
 &:=
 \sup_{\substack{s\in G_n\\\|a\|_S=1}}
 a_n^d\sum_{y\in\mathcal B_{n,q}}|h_{n,a}(s;y)|,\\
 \mathfrak F_{n,q}
 &:=
 \sup_{\substack{s\in G_n\\\|a\|_S=1}}
 a_n^{d-1}\sum_{f\in\mathcal F_{n,q}}
 |h_{n,a}^\top(s;f)|,\\
 \mathfrak N_{n,q}
 &:=
 \sup_{\substack{s\in G_n\\\|a\|_S=1}}
 a_n^{d-1}\sum_{f\in\mathcal F_{n,q}}
 |\nabla h_{n,a}(s;f)|,\\
 \mathfrak J_{n,q}
 &:=
 \sup_{\substack{s\in G_n\\\|a\|_S=1}}
 a_n^{d-2}\sum_{j\in\mathcal J_{n,q}}
 |h_{n,a}(s;j)|.                                    \tag{47.10}
\end{split}
\]
These definitions include all incidence multiplicities.  They are
the quantities to be estimated from the actual metric reconstruction,
not formal continuum integrals.

\begin{theorem}[Curvature-shell gradient debit]
\label{thm:v16v47-gradient}
Let a shell owner be an increment
\[
 V^{R}_{n,q}(s)
 :=\delta c_{R,n,q}\,
 S_{\rm full}(g_n(s);\Omega_{n,q}).                 \tag{47.11}
\]
Under the uniform good-set bounds (47.6),
\[
 \sup_{s\in G_n}\|\nabla_sV^R_{n,q}(s)\|_*
 \leq E^R_{n,q},                                    \tag{47.12}
\]
where
\[
 E^R_{n,q}
 :=
 |\delta c_{R,n,q}|
 \left(
 K_G\mathfrak B_{n,q}
 +K_\Pi\mathfrak F_{n,q}
 +K_{\mathcal J}\mathfrak J_{n,q}
 \right)                                           \tag{47.13}
\]
for the fully paired action.  For an unpaired EH owner, the safe
replacement is
\[
 E^{R,{\rm unpaired}}_{n,q}
 :=
 |\delta c_{R,n,q}|
 \left(K_G\mathfrak B_{n,q}
       +K_B\mathfrak N_{n,q}\right).                \tag{47.14}
\]
\end{theorem}

\begin{proof}
For a unit carrier direction \(a\), the chain rule gives
\[
 D_sV^R_{n,q}(s)[a]
 =\delta c_{R,n,q}\,
 DS_{\rm full}(g_n(s))[h_{n,a}(s)].                 \tag{47.15}
\]
Use the discrete versions of (47.7) and the definitions (47.10) to
obtain (47.13).  Equation (47.8) gives (47.14).  Take the supremum
over \(s\) and \(a\).
\end{proof}

\begin{corollary}[Curvature-counterterm state debit]
\label{cor:v16v47-state}
Suppose two consecutive soft laws differ by the positive scalar
weight \(e^{-V^R_{n,q}}\), the convex good set \(G_n\) has diameter
\(D_n\), and both laws give \(G_n^c\) probability at most
\(\eta_{n,q}\).  Then
\[
 d^R_{n,q}
 \leq e^{D_nE^R_{n,q}}-1+4\eta_{n,q}                \tag{47.16}
\]
is a certified owner debit.  The same statement holds with
\(E^{R,{\rm unpaired}}_{n,q}\) when no GHY pairing has been proved.
\end{corollary}

\begin{proof}
Apply Theorem~\ref{thm:v16v41-debit} using (47.12).
\end{proof}

\subsection{Four-dimensional shell scaling}

In \(d=4\), suppose for constants \(C_B,C_F,C_J\),
\[
\begin{split}
 |\mathcal B_{n,q}|&\leq
 C_B[2(R_n+(q+1)r)+1]^3,\\
 |\mathcal F_{n,q}|&\leq
 C_F[2(R_n+(q+1)r)+1]^3,\\
 |\mathcal J_{n,q}|&\leq
 C_J[2(R_n+(q+1)r)+1]^2.                            \tag{47.17}
\end{split}
\]
Assume uniform mode envelopes
\[
 |h_{n,a}|\leq\rho^{(0)}_{n,q},\qquad
 |\nabla h_{n,a}|\leq\rho^{(1)}_{n,q}               \tag{47.18}
\]
on the corresponding cells for unit \(a\).

\begin{proposition}[Explicit four-dimensional owner schedule]
\label{prop:v16v47-four}
The paired gradient budget obeys
\[
\begin{split}
 E^R_{n,q}\leq|\delta c_{R,n,q}|\bigg[
 &K_GC_Ba_n^4L_{n,q}^3\rho^{(0)}_{n,q}
 +K_\Pi C_Fa_n^3L_{n,q}^3\rho^{(0)}_{n,q}\\
 &+K_{\mathcal J}C_Ja_n^2L_{n,q}^2
 \rho^{(0)}_{n,q}\bigg],                            \tag{47.19}\\
 L_{n,q}&:=2(R_n+(q+1)r)+1.
\end{split}
\]
The unpaired budget contains instead
\[
\begin{split}
 E^{R,{\rm unpaired}}_{n,q}\leq
 |\delta c_{R,n,q}|\big[
 K_GC_Ba_n^4L_{n,q}^3\rho^{(0)}_{n,q}
 +K_BC_Fa_n^3L_{n,q}^3\rho^{(1)}_{n,q}
 \big].                                             \tag{47.20}
\end{split}
\]
\end{proposition}

\begin{proof}
Insert the cell counts (47.17) and envelopes (47.18) in the four
definitions (47.10), then use (47.13)--(47.14).
\end{proof}

\begin{theorem}[Cofinal curvature-owner tail criterion]
\label{thm:v16v47-tail}
Assume \(\sup_nD_n<\infty\) and
\[
 \lim_{Q\to\infty}\sup_n
 \sum_{q\geq Q}E^R_{n,q}=0,\qquad
 \lim_{Q\to\infty}\sup_n
 \sum_{q\geq Q}\eta_{n,q}=0.                        \tag{47.21}
\]
Then the fully paired curvature-counterterm shells have a vanishing
uniform state-distance tail.  If only the unpaired variational formula
is available, the same conclusion requires (47.21) with
\(E^{R,{\rm unpaired}}_{n,q}\).
\end{theorem}

\begin{proof}
The proof is the same summable small-exponent argument as
Corollary~\ref{cor:v16v46-countable}, now using (47.16), followed by
the V36 triangular-array theorem.
\end{proof}

\subsection{Pure-soft status and source retention}

\begin{proposition}[Geometric curvature shells remain soft owners]
\label{prop:v16v47-soft}
If the metric coordinates are entirely assigned to the soft carrier,
the curvature shell (47.11) leaves the conditional hard matter law
unchanged.  Its bulk Einstein tensor, Brown--York face momentum, and
joint response nevertheless remain in the metric derivative of the
full generating functional.
\end{proposition}

\begin{proof}
At fixed \(s\), \(V^R_{n,q}(s)\) is constant as a function of the
hard matter coordinate.  Proposition~\ref{prop:v16v46-pure-soft}
therefore applies.  Differentiating the scalar weight with respect to
\(s\) gives (47.15), whose three contributions are exactly
(47.2)--(47.5).
\end{proof}

\begin{assessment}[What V47 adds beyond the archive]
\label{ass:v16v47-status}
The Einstein--GHY and joint variations are imported rather than
reproved.  V47 adds the missing shell owner norm: bulk, smooth-face,
normal-derivative, and joint reconstruction costs are separated and
converted into a cofinal state-distance criterion.  It proves exactly
what GHY pairing saves and what it leaves behind.

No estimate in V47 proves the reconstruction tails, curvature bounds,
Brown--York bounds, coefficient rates, or corner control in (47.19).
Those are physical regulator inputs.
\end{assessment}

\begin{decision}[V48 trace-loading the face reconstruction norm]
\label{dec:v16v47-next}
Attack the analytic boundary term rather than redoing its geometry.
Prove discrete and continuum trace inequalities which bound
\(\mathfrak F_{n,q}\) and \(\mathfrak N_{n,q}\) by weighted bulk
\(W^{1,1}\) and \(W^{2,1}\) reconstruction norms on a collar, with
constants explicit in the collar width and uniform geometry.
Determine which derivative order is saved by a proved EGHY pairing.
\end{decision}

\begin{firewall}[Claim boundary after V47]
\label{fire:v16v47}
The imported geometric variation and the new owner norm do not prove
a positive gravitational measure or cofinal bounds on its metric
modes.  Piecewise boundaries require the joint row; it has not been
erased.  No limiting Einstein equation or full Standard
Model--Einstein continuum has been established.
\end{firewall}

\section*{Version V47 append-only record}
\addcontentsline{toc}{section}{Version V47 append-only record}
The complete V46 source was retained before this continuation.  It
had \(457064\) bytes, \(13772\) lines, and SHA--256
\[
 \texttt{530CA755F24588DF228A5B97E04792259BF0E7D10527A781B2E83E2E685638D6}.
\]
V47 imported the archived Einstein--GHY and joint variation instead
of duplicating it, then proved the bulk--face--joint reconstruction
norm and cofinal curvature-owner debit.

\section{V48: Explicit collar traces for curvature-owner faces and
joints}
\label{sec:v16v48}

The archive contains an \(L^2\) trace estimate used for a different
Robin semiboundedness problem, but no explicit-width \(L^1\) collar
estimate loading the V47 reconstruction norms.  This section proves
that missing estimate in continuum and discrete form.  It shows
precisely that EGHY pairing saves one bulk derivative at a smooth
face, while a surviving codimension-two joint still requires a
second-order collar norm.

\subsection{One-dimensional \(L^1\) trace lemma}

\begin{lemma}[Endpoint average inequality]
\label{lem:v16v48-one-dimensional}
If \(f\in W^{1,1}(0,\ell)\), then
\[
 |f(0)|
 \leq {1\over\ell}\int_0^\ell|f(t)|\,dt
      +\int_0^\ell|f'(t)|\,dt.                      \tag{48.1}
\]
\end{lemma}

\begin{proof}
For almost every \(t\),
\[
 f(0)=f(t)-\int_0^t f'(u)\,du.                      \tag{48.2}
\]
Take absolute values, average over \(t\in(0,\ell)\), and use Tonelli:
\[
 {1\over\ell}\int_0^\ell\int_0^t|f'(u)|\,du\,dt
 ={1\over\ell}\int_0^\ell(\ell-u)|f'(u)|\,du
 \leq\int_0^\ell|f'(u)|\,du.                       \tag{48.3}
\]
\end{proof}

\subsection{Continuum tensor trace on a uniform collar}

Let \(\Sigma=\{0\}\times\Sigma\) bound a Gaussian collar
\[
 \mathcal C_\ell=[0,\ell]\times\Sigma,\qquad
 g=dt^2+\gamma_t.                                   \tag{48.4}
\]
Write \(dA_t=J(t,x)dA_0\) and assume
\[
 J(t,x)\geq j_->0                                   \tag{48.5}
\]
uniformly.  Tensor values at \((t,x)\) are identified with the
boundary fiber at \((0,x)\) by parallel transport along the normal
geodesic.

\begin{theorem}[Uniform \(L^1\) collar trace]
\label{thm:v16v48-trace}
For every tensor field \(u\in W^{1,1}(\mathcal C_\ell)\),
\[
 \|u|_\Sigma\|_{L^1(\Sigma)}
 \leq {1\over j_-}
 \left(
 {\|u\|_{L^1(\mathcal C_\ell)}\over\ell}
 +\|\nabla_nu\|_{L^1(\mathcal C_\ell)}
 \right).                                          \tag{48.6}
\]
If \(u\in W^{2,1}\), then
\[
 \|\nabla u|_\Sigma\|_{L^1(\Sigma)}
 \leq {1\over j_-}
 \left(
 {\|\nabla u\|_{L^1(\mathcal C_\ell)}\over\ell}
 +\|\nabla_n\nabla u\|_{L^1(\mathcal C_\ell)}
 \right).                                          \tag{48.7}
\]
The last terms may be enlarged to the full covariant
\(W^{1,1}\) and \(W^{2,1}\) seminorms.
\end{theorem}

\begin{proof}
After parallel transport, the norm of the ordinary \(t\)-derivative
is bounded by \(|\nabla_nu|\).  Apply
Lemma~\ref{lem:v16v48-one-dimensional} for almost every \(x\) and
integrate against \(dA_0\).  Since
\[
 dV_g=J(t,x)\,dt\,dA_0,\qquad
 dt\,dA_0\leq j_-^{-1}dV_g,                         \tag{48.8}
\]
the result is (48.6).  Apply the same argument to the tensor
\(\nabla u\) to obtain (48.7).
\end{proof}

\begin{remark}[Width dependence is necessary]
\label{rem:v16v48-width}
The factor \(\ell^{-1}\) cannot be removed uniformly as
\(\ell\downarrow0\): a tensor constant in the normal direction has
boundary \(L^1\) norm equal to \(\ell^{-1}\) times its collar
\(L^1\) norm.
\end{remark}

\subsection{Codimension-two iterated trace}

Let a joint \(\mathcal J\) have a product collar
\[
 [0,\ell_1]\times[0,\ell_2]\times\mathcal J,        \tag{48.9}
\]
with the corresponding one-direction Jacobian lower bounds at least
\(j_-\).

\begin{theorem}[Joint trace requires two derivatives]
\label{thm:v16v48-joint}
For \(u\in W^{2,1}\) on the joint collar,
\[
\begin{split}
 \|u|_{\mathcal J}\|_{L^1(\mathcal J)}
 \leq {1\over j_-^2}\bigg[
 &{\|u\|_{L^1}\over\ell_1\ell_2}
 +\left({1\over\ell_1}+{1\over\ell_2}\right)
   \|\nabla u\|_{L^1}\\
 &+\|\nabla^2u\|_{L^1}
 \bigg].                                            \tag{48.10}
\end{split}
\]
All norms on the right are over the two-normal collar.
\end{theorem}

\begin{proof}
Apply (48.6) first in the \(t_1\) direction, retaining \(t_2\) as a
tangential parameter.  Apply it a second time in the \(t_2\)
direction to both the zeroth-order term and its first normal
derivative.  The four resulting contributions have coefficients
\((\ell_1\ell_2)^{-1}\),
\(\ell_1^{-1}\), \(\ell_2^{-1}\), and \(1\).  Enlarge each directional
derivative to the full covariant derivative of the same order.  Each
trace contributes at most one factor \(j_-^{-1}\), giving (48.10).
\end{proof}

\subsection{Exact discrete collar inequality}

Let a discrete collar have layers \(k=0,\ldots,L\) with spacing \(a\),
and let \(U_{k,x}\) be a scalar or normed-vector value above a boundary
site \(x\).  Put
\[
 \ell_a:=(L+1)a,\qquad
 \nabla_n^aU_{k,x}:={U_{k+1,x}-U_{k,x}\over a}.      \tag{48.11}
\]

\begin{theorem}[Discrete \(L^1\) trace with exact scaling]
\label{thm:v16v48-discrete}
One has
\[
\begin{split}
 a^{d-1}\sum_x|U_{0,x}|
 \leq&
 {1\over\ell_a}
 a^d\sum_{k=0}^{L}\sum_x|U_{k,x}|\\
 &+
 a^d\sum_{k=0}^{L-1}\sum_x|\nabla_n^aU_{k,x}|.
                                                               \tag{48.12}
\end{split}
\]
If positive cell and face weights are uniformly comparable to the
displayed powers of \(a\) with ratio at most \(C_w\), the right side
is multiplied by at most \(C_w\).
\end{theorem}

\begin{proof}
For every \(k\),
\[
 |U_{0,x}|
 \leq |U_{k,x}|+
 \sum_{j=0}^{k-1}|U_{j+1,x}-U_{j,x}|.               \tag{48.13}
\]
Average over \(k=0,\ldots,L\), sum over \(x\), and bound the
multiplicity of each difference by \(L+1\):
\[
 \sum_x|U_{0,x}|
 \leq{1\over L+1}\sum_{k,x}|U_{k,x}|
 +\sum_{k=0}^{L-1}\sum_x|U_{k+1,x}-U_{k,x}|.        \tag{48.14}
\]
Multiplication by \(a^{d-1}\) and (48.11) give (48.12).
Uniform weight comparison proves the last assertion.
\end{proof}

\begin{corollary}[Discrete derivative and joint traces]
\label{cor:v16v48-discrete-higher}
Applying (48.12) to \(\nabla^aU\) bounds its boundary trace by
\[
 C_w\left(
 \ell_a^{-1}\|\nabla^aU\|_{1,\mathcal C}
 +\|(\nabla^a)^2U\|_{1,\mathcal C}
 \right).                                          \tag{48.15}
\]
Applying the discrete theorem successively in two transverse collar
directions gives the analogue of (48.10), with
\[
 \ell_{1,a}^{-1}\ell_{2,a}^{-1}\|U\|_1
 +(\ell_{1,a}^{-1}+\ell_{2,a}^{-1})\|\nabla^aU\|_1
 +\|(\nabla^a)^2U\|_1,                              \tag{48.16}
\]
up to \(C_w^2\).
\end{corollary}

\begin{proof}
The first statement is Theorem~\ref{thm:v16v48-discrete} with
\(U\) replaced by its discrete first derivative.  The second repeats
the proof in two coordinate directions; the multiplicative weight
comparison is used once per trace.
\end{proof}

\subsection{Loading the V47 reconstruction norms}

Let \(\mathcal C_{n,q}\) be a collar of physical width
\(\ell_{n,q}\) around the relevant shell face, with uniform constants
\[
 j_-\geq j_0>0,\qquad C_w\leq C_0.                  \tag{48.17}
\]
For the reconstructed unit-direction metric mode \(h_{n,a}\), define
\[
 \mathfrak W^{(k)}_{n,q}
 :=
 \sup_{\substack{s\in G_n\\\|a\|_S=1}}
 \|(\nabla^a)^kh_{n,a}(s)\|_{L^1(\mathcal C_{n,q})},
 \qquad k=0,1,2.                                    \tag{48.18}
\]

\begin{theorem}[Face and joint reconstruction loader]
\label{thm:v16v48-loader}
There is a cutoff-independent \(C_{\rm tr}\), determined only by the
uniform collar geometry or weight comparison, such that
\[
 \mathfrak F_{n,q}
 \leq C_{\rm tr}
 \left(
 \ell_{n,q}^{-1}\mathfrak W^{(0)}_{n,q}
 +\mathfrak W^{(1)}_{n,q}
 \right),                                          \tag{48.19}
\]
\[
 \mathfrak N_{n,q}
 \leq C_{\rm tr}
 \left(
 \ell_{n,q}^{-1}\mathfrak W^{(1)}_{n,q}
 +\mathfrak W^{(2)}_{n,q}
 \right),                                          \tag{48.20}
\]
and, for comparable two-normal widths,
\[
 \mathfrak J_{n,q}
 \leq C_{\rm tr}^2
 \left(
 \ell_{n,q}^{-2}\mathfrak W^{(0)}_{n,q}
 +2\ell_{n,q}^{-1}\mathfrak W^{(1)}_{n,q}
 +\mathfrak W^{(2)}_{n,q}
 \right).                                          \tag{48.21}
\]
\end{theorem}

\begin{proof}
Use (48.6) or (48.12) for \(h_{n,a}\), then take the supremum over
\(s,a\), to obtain (48.19).  Apply the first-derivative trace
(48.7) or (48.15) for (48.20).  The iterated trace
(48.10) or (48.16) gives (48.21).
\end{proof}

\begin{corollary}[Bulk Sobolev form of the paired owner budget]
\label{cor:v16v48-budget}
The V47 paired curvature budget satisfies
\[
\begin{split}
 E^R_{n,q}
 \leq|\delta c_{R,n,q}|\bigg[
 &K_G\mathfrak B_{n,q}
 +K_\Pi C_{\rm tr}
  \left(\ell_{n,q}^{-1}\mathfrak W^{(0)}_{n,q}
       +\mathfrak W^{(1)}_{n,q}\right)\\
 &+K_{\mathcal J}C_{\rm tr}^2
 \left(\ell_{n,q}^{-2}\mathfrak W^{(0)}_{n,q}
      +2\ell_{n,q}^{-1}\mathfrak W^{(1)}_{n,q}
      +\mathfrak W^{(2)}_{n,q}\right)
 \bigg].                                            \tag{48.22}
\end{split}
\]
The unpaired smooth-face budget instead obeys
\[
\begin{split}
 E^{R,{\rm unpaired}}_{n,q}
 \leq|\delta c_{R,n,q}|\bigg[
 K_G\mathfrak B_{n,q}
 +K_BC_{\rm tr}
 \left(\ell_{n,q}^{-1}\mathfrak W^{(1)}_{n,q}
      +\mathfrak W^{(2)}_{n,q}\right)
 \bigg].                                            \tag{48.23}
\end{split}
\]
\end{corollary}

\begin{proof}
Substitute (48.19)--(48.21) in (47.13)--(47.14).
\end{proof}

\begin{theorem}[Derivative order saved by boundary pairing]
\label{thm:v16v48-order}
On smooth shell boundaries with no joints, a proved
Einstein--GHY pairing loads the face owner from a bulk
\(W^{1,1}\) reconstruction tail through (48.19).  Without that
pairing, the safe normal-derivative term requires \(W^{2,1}\) through
(48.20).  On a piecewise shell with a surviving joint response,
the paired action again requires \(W^{2,1}\) unless the joint
variation is separately canceled or fixed.
\end{theorem}

\begin{proof}
With no joints, set \(K_{\mathcal J}=0\) in (48.22); only
\(\mathfrak W^{(0)}\) and \(\mathfrak W^{(1)}\) occur.  Equation
(48.23) contains \(\mathfrak W^{(2)}\).  With joints,
(48.21) inserts \(\mathfrak W^{(2)}\) into (48.22).
\end{proof}

\begin{assessment}[What V48 closes]
\label{ass:v16v48-status}
V48 proves the explicit collar estimates that load V47's geometric
face and joint owners from bulk reconstruction norms.  It quantifies
the collar-width loss, supplies an exact lattice version, and proves
the derivative-order saving produced by a valid EGHY pairing.

The result does not supply the weighted \(W^{1,1}\) or \(W^{2,1}\)
tails themselves.  A collar whose physical width collapses can make
the \(\ell^{-1}\) and \(\ell^{-2}\) factors diverge; cofinal geometry
must pay those factors.
\end{assessment}

\begin{decision}[V49 weighted metric-mode collar tails]
\label{dec:v16v48-next}
Connect the trace loader to the moving soft-space results V27--V32.
Formulate a weighted Sobolev reconstruction operator from the
augmented soft coordinate into metric collars, prove that
trace-class domination of its singular values gives uniform
\(W^{k,1}\) moment and tail bounds, and retain every factor of the
collar width.  Audit the archives first for an existing
operator-to-\(W^{k,1}\) theorem.
\end{decision}

\begin{firewall}[Claim boundary after V48]
\label{fire:v16v48}
Trace inequalities transfer bulk control to a boundary; they do not
create bulk control.  V48 does not prove a cofinal metric
reconstruction estimate, positive gravitational measure, limiting
Einstein equation, or full Standard Model--Einstein continuum.
\end{firewall}

\section*{Version V48 append-only record}
\addcontentsline{toc}{section}{Version V48 append-only record}
The complete V47 source was retained before this continuation.  It
had \(469319\) bytes, \(14132\) lines, and SHA--256
\[
 \texttt{C411857C4DC7428B6F6106257CD00806A52E7D6C793FCB8244C8BEA997043BA7}.
\]
V48 proved continuum and exact discrete \(L^1\) collar traces and
loaded the V47 face, normal-derivative, and joint owners from
weighted bulk \(W^{1,1}\) and \(W^{2,1}\) reconstruction norms.

\section{V49: Common-carrier Sobolev reconstruction tails}
\label{sec:v16v49}

V48 reduces the curvature owner to bulk \(W^{k,1}\) tails of unit
soft directions.  A bounded trace-class stock at each cutoff is not
enough: its range can drift to ever more remote shells.  This section
proves the required common-carrier criterion and then converts
trace-norm reconstruction tails into the weighted \(L^1\) quantities
of (48.22)--(48.23).

\subsection{Uniform trace tails from trace-norm convergence}

Let \(\mathscr S\) be the augmented soft Hilbert carrier supplied by
the V30--V31 architecture.  Let
\[
 \mathscr Y=\bigoplus_{q=0}^{\infty}\mathscr Y_q     \tag{49.1}
\]
be a Hilbert direct sum of shell Sobolev spaces, with orthogonal
projections
\[
 P_{\geq Q}:\mathscr Y\longrightarrow
 \bigoplus_{q\geq Q}\mathscr Y_q.                   \tag{49.2}
\]
Then \(P_{\geq Q}\to0\) strongly as \(Q\to\infty\).

\begin{lemma}[Trace-class operators have strong-projection tails]
\label{lem:v16v49-trace-tail}
If \(R\in\mathfrak S_1(\mathscr S,\mathscr Y)\), then
\[
 \|P_{\geq Q}R\|_1\longrightarrow0.                 \tag{49.3}
\]
\end{lemma}

\begin{proof}
Choose a finite-rank operator \(F\) with
\(\|R-F\|_1<\varepsilon\).  Strong convergence of the uniformly
bounded projections is norm convergence on the finite-dimensional
range of \(F\), so
\(\|P_{\geq Q}F\|_1\to0\).  Hence
\[
 \|P_{\geq Q}R\|_1
 \leq\|P_{\geq Q}(R-F)\|_1+\|P_{\geq Q}F\|_1
 \leq\varepsilon+\|P_{\geq Q}F\|_1.                \tag{49.4}
\]
Send \(Q\to\infty\), then \(\varepsilon\downarrow0\).
\end{proof}

\begin{theorem}[Uniform cofinal reconstruction tail]
\label{thm:v16v49-uniform}
If
\[
 R_n\longrightarrow R
 \quad\hbox{in }\mathfrak S_1(\mathscr S,\mathscr Y),            \tag{49.5}
\]
then
\[
 \lim_{Q\to\infty}\sup_n
 \|P_{\geq Q}R_n\|_1=0.                             \tag{49.6}
\]
Consequently,
\[
 \lim_{Q\to\infty}\sup_n\sup_{\|a\|_{\mathscr S}\leq1}
 \|P_{\geq Q}R_na\|_{\mathscr Y}=0.                 \tag{49.7}
\]
\end{theorem}

\begin{proof}
Given \(\varepsilon>0\), choose \(N\) so that
\(\|R_n-R\|_1<\varepsilon\) for \(n\geq N\).  For these indices,
\[
 \|P_{\geq Q}R_n\|_1
 \leq\varepsilon+\|P_{\geq Q}R\|_1.                 \tag{49.8}
\]
Lemma~\ref{lem:v16v49-trace-tail} makes the second term small.
For the finitely many \(n<N\), apply that lemma to each \(R_n\) and
choose one larger \(Q\).  This proves (49.6).  Since operator norm is
bounded by trace norm, (49.7) follows.
\end{proof}

\begin{counterexample}[Uniform singular values do not prevent shell
drift]
\label{cnt:v16v49-drift}
Let \(\mathscr S=\mathbb R\), \(\mathscr Y=\ell^2(\mathbb N)\), and
\[
 R_na:=a e_n.                                       \tag{49.9}
\]
Every \(R_n\) has the same single singular value \(1\) and
\(\|R_n\|_1=1\), but
\[
 \sup_n\|P_{\geq Q}R_n\|_1=1                       \tag{49.10}
\]
for every \(Q\).  Thus a cutoff-uniform trace norm or a common
singular-value majorant without range localization does not imply
(49.6).
\end{counterexample}

\begin{proof}
The operator is rank one and isometric on its one-dimensional
domain, proving its singular-value statement.  For every \(Q\),
choose \(n\geq Q\); then \(P_{\geq Q}R_n=R_n\), proving (49.10).
\end{proof}

\subsection{From Hilbert Sobolev norms to the \(L^1\) ledger}

Let \(\mathcal C_{n,q}\) be the collar used in V48, with volume
\[
 V_{n,q}:={\rm Vol}(\mathcal C_{n,q}).               \tag{49.11}
\]
For \(j=0,1,2\), let
\[
 R_n^{(j)}:\mathscr S\longrightarrow
 \mathscr Y^{(j)}
 :=\bigoplus_{q\geq0}L^2(\mathcal C_{n,q};\mathcal T_j)          \tag{49.12}
\]
map a soft direction to the \(j\)-th covariant derivative of its
reconstructed metric mode.  Common geometric identifications are
understood before (49.12), so the target direct sums do not move with
\(n\).  Denote their shell projections again by \(P_q\).

\begin{lemma}[\(L^2\)-to-\(L^1\) shell conversion]
\label{lem:v16v49-L2-L1}
For every unit soft direction \(a\),
\[
 \|P_qR_n^{(j)}a\|_{L^1(\mathcal C_{n,q})}
 \leq
 V_{n,q}^{1/2}
 \|P_qR_n^{(j)}a\|_{L^2(\mathcal C_{n,q})}.          \tag{49.13}
\]
If \(\beta_{n,q}^{(j)}\geq0\), then
\[
\begin{split}
 \sum_{q\geq Q}\beta_{n,q}^{(j)}
 \|P_qR_n^{(j)}a\|_{L^1}
 \leq A_{n,Q}^{(j)}
 \|P_{\geq Q}R_n^{(j)}a\|_{\mathscr Y^{(j)}},       \tag{49.14}\\
 A_{n,Q}^{(j)}
 &:=
 \left(
 \sum_{q\geq Q}
 [\beta_{n,q}^{(j)}]^2V_{n,q}
 \right)^{1/2}.
\end{split}
\]
\end{lemma}

\begin{proof}
Equation (49.13) is Cauchy--Schwarz against the constant function
one on the collar.  Multiply it by \(\beta_{n,q}^{(j)}\), sum in
\(q\), and apply Cauchy--Schwarz to the shell index.  Orthogonality of
the direct sum turns the second factor into the norm in (49.14).
\end{proof}

\begin{theorem}[Weighted Sobolev reconstruction compiler]
\label{thm:v16v49-compiler}
For \(j=0,\ldots,k\), suppose
\[
 R_n^{(j)}\longrightarrow R^{(j)}
 \quad\hbox{in trace norm},                         \tag{49.15}
\]
and let nonnegative geometry arrays satisfy
\[
 \sup_{n,Q} A_{n,Q}^{(j)}<\infty.                   \tag{49.16}
\]
Then
\[
\lim_{Q\to\infty}\sup_n\sup_{\|a\|_{\mathscr S}\leq1}
\sum_{q\geq Q}\beta_{n,q}^{(j)}
\|P_qR_n^{(j)}a\|_{L^1}=0.                          \tag{49.17}
\]
More generally, (49.17) follows whenever
\[
\lim_{Q\to\infty}\sup_n
A_{n,Q}^{(j)}\|P_{\geq Q}R_n^{(j)}\|_1=0.           \tag{49.18}
\]
\end{theorem}

\begin{proof}
Take the supremum over unit \(a\) in (49.14) and use
\[
 \|P_{\geq Q}R_n^{(j)}\|_{\rm op}
 \leq\|P_{\geq Q}R_n^{(j)}\|_1.                    \tag{49.19}
\]
Under (49.15), Theorem~\ref{thm:v16v49-uniform} makes the latter
trace-norm tail vanish uniformly.  The bounded arrays (49.16) give
(49.17).  Without separate boundedness, the same estimate proves the
product criterion (49.18).
\end{proof}

\subsection{Explicit arrays for the curvature owner}

For notational clarity, absorb cutoff-independent constants into
\(\lesssim\).  In the paired V48 budget (48.22), define
\[
\begin{split}
 \beta_{n,q}^{(0)}
 &:=
 |\delta c_{R,n,q}|
 \left(
 K_G+K_\Pi C_{\rm tr}\ell_{n,q}^{-1}
 +K_{\mathcal J}C_{\rm tr}^2\ell_{n,q}^{-2}
 \right),\\
 \beta_{n,q}^{(1)}
 &:=
 |\delta c_{R,n,q}|
 \left(
 K_\Pi C_{\rm tr}
 +2K_{\mathcal J}C_{\rm tr}^2\ell_{n,q}^{-1}
 \right),\\
 \beta_{n,q}^{(2)}
 &:=
 |\delta c_{R,n,q}|K_{\mathcal J}C_{\rm tr}^2.
                                                               \tag{49.20}
\end{split}
\]
The zeroth-order bulk shell may be included in its collar by
enlarging the collar a fixed bounded number of layers.

\begin{corollary}[Trace-class loading of the paired curvature tail]
\label{cor:v16v49-paired}
If (49.15)--(49.18) hold for the three arrays (49.20), then
\[
 \lim_{Q\to\infty}\sup_n
 \sum_{q\geq Q}E^R_{n,q}=0.                         \tag{49.21}
\]
Consequently, with the carrier diameter and bad-set hypotheses of
(47.21), the paired curvature owner has a vanishing cofinal
state-distance tail.
\end{corollary}

\begin{proof}
Equation (48.22) is bounded by the sum of the three left sides in
(49.17), with the arrays (49.20).  Each vanishes.  Apply
Theorem~\ref{thm:v16v47-tail}.
\end{proof}

For an unpaired smooth face, the corresponding arrays may be taken as
\[
\begin{split}
 \widetilde\beta_{n,q}^{(0)}
 &:=|\delta c_{R,n,q}|K_G,\\
 \widetilde\beta_{n,q}^{(1)}
 &:=|\delta c_{R,n,q}|K_BC_{\rm tr}\ell_{n,q}^{-1},\\
 \widetilde\beta_{n,q}^{(2)}
 &:=|\delta c_{R,n,q}|K_BC_{\rm tr}.                \tag{49.22}
\end{split}
\]

\begin{corollary}[Trace-class loading of the unpaired face]
\label{cor:v16v49-unpaired}
If (49.15)--(49.18) hold with the arrays (49.22), then
\[
 \lim_{Q\to\infty}\sup_n
 \sum_{q\geq Q}E^{R,{\rm unpaired}}_{n,q}=0.        \tag{49.23}
\]
\end{corollary}

\begin{proof}
Use (48.23) exactly as in
Corollary~\ref{cor:v16v49-paired}.
\end{proof}

\subsection{Singular-value formulation and its limitation}

Let \(\sigma_m(R_n^{(j)})\) be the singular values in decreasing
order.  Trace-norm convergence in (49.15) implies
\[
 \sup_n\sum_m\sigma_m(R_n^{(j)})<\infty.            \tag{49.24}
\]
The converse is false, as Counterexample~\ref{cnt:v16v49-drift}
shows.  A usable acquisition card is:
\begin{enumerate}
\item put every cutoff soft space and metric collar in the common
augmented carriers of V30--V31;
\item prove convergence of the reconstruction operators on a dense
set of soft directions;
\item prove a common trace-class singular-value tail after those
identifications;
\item use finite-rank approximation to upgrade to (49.15);
\item verify the geometry-array product (49.18), including every
\(\ell_{n,q}^{-1}\) or \(\ell_{n,q}^{-2}\) loss.
\end{enumerate}
Without the common-carrier localization in items 1--2, singular
values alone cannot locate the reconstructed mode.

\begin{assessment}[What V49 closes]
\label{ass:v16v49-status}
V49 supplies the operator-to-\(W^{k,1}\) bridge missing after V48.
Trace-norm convergence on a common shell carrier gives uniform
reconstruction tails, and weighted Cauchy--Schwarz pays the
four-dimensional geometry and collar-width factors.  The drifting
rank-one example proves why a uniform trace-class stock alone is
insufficient.

The reconstruction convergence (49.15) is not yet proved for the
physical metric modes.  V27--V32 provide the common augmented
coordinate and compactness architecture, but not the new
\(H^2\)-collar operator convergence required here.
\end{assessment}

\begin{decision}[V50 companion audit and Sobolev reconstruction
compactness]
\label{dec:v16v49-next}
Perform the mandatory companion audit.  Then prove a finite-rank
plus-tail criterion which upgrades convergence on a common dense core
and a uniform singular-value tail to trace-norm convergence of the
metric reconstruction operators in (49.15).  Continue in V51 by
testing that criterion on the covariance-weighted augmented soft
operators of V30--V32.
\end{decision}

\begin{firewall}[Claim boundary after V49]
\label{fire:v16v49}
The trace-class compiler is conditional on common-carrier
reconstruction convergence and the geometry arrays.  It does not
derive those estimates from the metric action, prove gravitational
positivity, identify the limiting Einstein source, or construct the
full Standard Model--Einstein continuum.
\end{firewall}

\section*{Version V49 append-only record}
\addcontentsline{toc}{section}{Version V49 append-only record}
The complete V48 source was retained before this continuation.  It
had \(481006\) bytes, \(14491\) lines, and SHA--256
\[
 \texttt{FC52A65000F31EC455136546B7ABA55727BD037D147357452D77BA4D19571AFD}.
\]
V49 proved the common-carrier trace-norm tail theorem, the weighted
\(L^2\)-to-\(W^{k,1}\) shell compiler, and the explicit paired and
unpaired curvature-owner arrays.

\section{V50: Companion audit and finite-rank trace-norm closure}
\label{sec:v16v50}

V49 reduced the curvature-owner problem to trace-norm convergence of
common-carrier Sobolev reconstruction operators.  This section performs
the compulsory cross-program audit and proves the missing closure step.
The essential datum is an aligned finite-dimensional exhaustion of the
\emph{input} carrier.  A uniform list of singular values is not enough,
even when the operators converge strongly on every fixed vector.

\subsection{Scheduled companion and archive audit}

The current Yang--Mills companion is
\texttt{Renewal\_Yang\_Mills\_Mass\_Gap\_Research\_Notes\_V53.tex}.
It has (458334) bytes, (12505) lines, and SHA--256
\[
 \texttt{ABF55B212473E8E116F87FF3894AF1DD82606B6D1089948128E3865BB4CD8061}.
                                                               \tag{50.1}
\]
Its V53 appendix completes the second mixed quarter face and constructs
a direct positive two-coordinate plate
\((\pi/4)\times(\pi/4)\times t\).  The third coordinate and the
opposite sign on the same plate remain open.  These sign charts do not
supply a metric Sobolev reconstruction estimate.

The current Navier--Stokes companion is
\texttt{Renewal\_NS\_Stability\_Research\_Notes\_V26.tex}.  It has
(278283) bytes, (5319) lines, and SHA--256
\[
 \texttt{82C887F8C2C69E4F28FAD7C3DC8DE5B9099CB5A4BF431EBA1FFF3E91935978AD}.
                                                               \tag{50.2}
\]
Its V26 appendix computes exact pointwise rank-one norms for the first
two Oseen-kernel derivatives.  That finite-dimensional angular
optimization neither gives a common soft carrier nor an operator-ideal
tail for the metric reconstruction map.

The de-duplication search found the closest earlier SM results in the
fifteenth-cycle V86--V88 archive.  V86 proves vector compactness from a
common output projection tail.  V87 proves strong convergence after
\emph{assuming} Hilbert--Schmidt or trace-norm convergence of a common
compact factor and records the drifting output rank-one example.  V88
derives trace-norm convergence of heat operators from uniform spectral
counting and norm-resolvent convergence.  None proves that strong
convergence on a soft core plus an aligned input trace tail implies
trace-norm convergence of the whole reconstruction operator.  The
result below fills precisely that gap and imports, rather than repeats,
the earlier output-tail machinery.

\subsection{The one-sided finite-rank closure theorem}

Let \(\mathscr S\) and \(\mathscr Y\) be separable Hilbert spaces.  Let
\[
 E_M:\mathscr S\longrightarrow\mathscr S,
 \qquad E_M\uparrow I\quad\hbox{strongly},             \tag{50.3}
\]
be increasing finite-rank orthogonal projections.  Write
\(d_M=\operatorname{rank}E_M\).

\begin{lemma}[Right trace tail]
\label{lem:v16v50-right-tail}
If \(R\in\mathfrak S_1(\mathscr S,\mathscr Y)\), then
\[
 \|R(I-E_M)\|_1\longrightarrow0.                      \tag{50.4}
\]
\end{lemma}

\begin{proof}
Apply Lemma~\ref{lem:v16v49-trace-tail} to \(R^*\) and the projections
\(I-E_M\) on its target.  Since
\[
 \|R(I-E_M)\|_1=\|(I-E_M)R^*\|_1,                    \tag{50.5}
\]
the asserted limit follows.
\end{proof}

\begin{theorem}[Strong-core trace-norm closure]
\label{thm:v16v50-strong-core}
Let \(\mathscr D\subset\mathscr S\) be a dense linear subspace such
that
\[
 E_M\mathscr S\subset\mathscr D\qquad\hbox{for every }M.       \tag{50.6}
\]
Suppose \(R_n\in\mathfrak S_1(\mathscr S,\mathscr Y)\), the sequence
\((R_nx)_n\) is Cauchy in \(\mathscr Y\) for every
\(x\in\mathscr D\), and
\[
 \lim_{M\to\infty}\sup_n\|R_n(I-E_M)\|_1=0.          \tag{50.7}
\]
Then there is a unique \(R\in\mathfrak S_1(\mathscr S,\mathscr Y)\)
such that
\[
 \|R_n-R\|_1\longrightarrow0.                         \tag{50.8}
\]
For every \(x\in\mathscr D\), \(Rx\) is the prescribed strong limit
of \(R_nx\).
\end{theorem}

\begin{proof}
Fix \(\varepsilon>0\).  Choose \(M\) so that the supremum in (50.7)
is below \(\varepsilon\), and choose an orthonormal basis
\(e_1,\ldots,e_{d_M}\) of \(E_M\mathscr S\).  By (50.6), each
\((R_ne_j)_n\) is Cauchy.  For finite-rank operators the Schatten
inequality gives
\[
\begin{split}
 \|(R_n-R_m)E_M\|_1
 &\leq \sqrt{d_M}\,\|(R_n-R_m)E_M\|_2\\
 &=\sqrt{d_M}
   \left(\sum_{j=1}^{d_M}\|(R_n-R_m)e_j\|_{\mathscr Y}^2\right)^{1/2}
 \longrightarrow0                                      \tag{50.9}
\end{split}
\]
as \(n,m\to\infty\), with \(M\) fixed.  Hence, for all sufficiently
large \(n,m\),
\[
\begin{split}
 \|R_n-R_m\|_1
 &\leq \|(R_n-R_m)E_M\|_1
       +\|R_n(I-E_M)\|_1+\|R_m(I-E_M)\|_1\\
 &<3\varepsilon.                                      \tag{50.10}
\end{split}
\]
Thus \((R_n)\) is Cauchy in the Banach space
\(\mathfrak S_1(\mathscr S,\mathscr Y)\), and has a trace-class limit
\(R\).  Since
\(\|(R_n-R)x\|\leq\|R_n-R\|_1\|x\|\), the limit agrees with the
given pointwise limit on \(\mathscr D\).  Density makes that bounded
extension unique.
\end{proof}

\begin{corollary}[Prescribed reconstruction limit]
\label{cor:v16v50-prescribed}
If a linear map \(R_0:\mathscr D\to\mathscr Y\) satisfies
\[
 R_nx\longrightarrow R_0x\qquad(x\in\mathscr D)       \tag{50.11}
\]
and (50.7) holds, then \(R_0\) extends uniquely to the trace-class
operator \(R\) in (50.8).
\end{corollary}

\begin{proof}
Apply Theorem~\ref{thm:v16v50-strong-core}.  Its identification on
\(\mathscr D\) says that the limit extends \(R_0\).
\end{proof}

\subsection{Weak cores require localization on both sides}

Let
\[
 P_Q:\mathscr Y\longrightarrow\mathscr Y,
 \qquad P_Q\uparrow I\quad\hbox{strongly},             \tag{50.12}
\]
be finite-rank orthogonal projections.

\begin{theorem}[Two-sided weak-core closure]
\label{thm:v16v50-weak-core}
Retain (50.3) and (50.6).  Suppose \((R_nx)_n\) is weakly Cauchy in
\(\mathscr Y\) for every \(x\in\mathscr D\), and suppose
\[
 \lim_{M\to\infty}\sup_n\|R_n(I-E_M)\|_1=0,
 \qquad
 \lim_{Q\to\infty}\sup_n\|(I-P_Q)R_n\|_1=0.          \tag{50.13}
\]
Then \((R_n)\) converges in trace norm to a trace-class operator.
\end{theorem}

\begin{proof}
Fix \(M\).  For every basis vector \(e_j\) of \(E_M\mathscr S\),
the vector sequence \(P_QR_ne_j\) is weakly Cauchy in the
finite-dimensional space \(P_Q\mathscr Y\), hence norm Cauchy.  The
finite-rank estimate used in (50.9) therefore gives
\[
 \|P_Q(R_n-R_m)E_M\|_1\longrightarrow0                \tag{50.14}
\]
for fixed \(M,Q\).  On the other hand,
\[
\begin{split}
 \|(R_n-R_m)E_M\|_1
 \leq{}&\|P_Q(R_n-R_m)E_M\|_1\\
 &+\|(I-P_Q)R_n\|_1+\|(I-P_Q)R_m\|_1.                \tag{50.15}
\end{split}
\]
First choose \(M\) to control the two input tails in the analogue of
(50.10), then choose \(Q\) to control the two output tails in
(50.15), and finally send \(n,m\to\infty\) in the finite block.  Thus
\((R_n)\) is trace-norm Cauchy.  Completeness supplies its limit.
\end{proof}

Theorem~\ref{thm:v16v50-weak-core} is the exact operator-ideal version
of the earlier vector compactness compiler: weak convergence becomes
strong on a fixed finite output block, while the two trace tails pay
for removing both finite-rank cutoffs.

\subsection{Why an intrinsic singular-value tail still fails}

\begin{counterexample}[Strong convergence with drifting input]
\label{cnt:v16v50-input-drift}
Let \(\mathscr S=\mathscr Y=\ell^2(\mathbb N)\), with basis
\((e_j)\), and put
\[
 R_n=|e_1\rangle\langle e_n|.                         \tag{50.16}
\]
Then \(R_nx=\langle e_n,x\rangle e_1\to0\) for every
\(x\in\ell^2\), and every \(R_n\) has the single singular value one.
Thus its singular-value tail vanishes identically beyond rank one, but
\[
 \|R_n\|_1=1                                         \tag{50.17}
\]
and there is no trace-norm convergence to zero.  For the coordinate
projections \(E_M\),
\[
 \sup_n\|R_n(I-E_M)\|_1=1                            \tag{50.18}
\]
for every \(M\).  Hence (50.7) detects exactly the missing alignment.
The adjoint family is the output-drift example already recorded in
(49.9)--(49.10); under merely weak core convergence it violates the
second tail in (50.13).
\end{counterexample}

\begin{proof}
Fourier coefficients of a fixed \(\ell^2\) vector tend to zero, which
proves strong pointwise convergence.  Each \(R_n\) is rank one with
unit input and output vectors, so its sole singular value and trace
norm equal one.  Given \(M\), choose \(n>M\); then
\(R_n(I-E_M)=R_n\), proving (50.18).
\end{proof}

\subsection{Reconstruction acquisition card}

For each derivative order \(j=0,1,2\), the V49 hypothesis (49.15) now
follows from either of the following concrete packages:
\begin{enumerate}
\item choose augmented soft-mode projections \(E_M\) whose union is a
      common core, prove strong convergence of \(R_n^{(j)}a\) on that
      core, and prove the uniform input trace tail
      \[
       \lim_{M\to\infty}\sup_n
       \|R_n^{(j)}(I-E_M)\|_1=0;                     \tag{50.19}
      \]
\item prove only weak core convergence, but supplement (50.19) by an
      aligned shell-Sobolev output tail as in the second condition of
      (50.13).
\end{enumerate}
An intrinsic estimate
\(\sup_n\sum_{r>M}\sigma_r(R_n^{(j)})\to0\) cannot replace
(50.19), by Counterexample~\ref{cnt:v16v50-input-drift}.  The
singular directions themselves must be controlled in the common
augmented soft coordinate.

\begin{assessment}[What V50 closes]
\label{ass:v16v50-status}
V50 performs the scheduled companion audit and proves two noncircular
routes from finite soft blocks to trace-norm reconstruction convergence.
The strong-core route needs one aligned input trace tail; the weak-core
route needs aligned tails on both input and output.  The drifting-input
example proves that uniform rank and singular values do not imply either
criterion.

No current SM, Yang--Mills, or Navier--Stokes note supplies (50.19) for
the physical metric reconstruction.  The next step must test whether
the covariance-weighted augmented operators already constructed in
V30--V32 factor the required map with this common input localization.
\end{assessment}

\begin{decision}[V51 covariance-weighted factorization test]
\label{dec:v16v50-next}
Inspect the exact V30--V32 augmented soft operators.  Determine which
factor is already trace-class and convergent, prove the corresponding
operator-ideal composition estimate, and state the smallest additional
Sobolev reconstruction bound that would yield (50.19).  Do not relabel
covariance compactness as metric \(H^2\) regularity.
\end{decision}

\begin{firewall}[Claim boundary after V50]
\label{fire:v16v50}
Finite-rank closure is a conditional functional-analytic compiler.  V50
does not establish the physical input tail, construct a positive metric
measure, identify a limiting Einstein source, or complete the Standard
Model--Einstein continuum.
\end{firewall}

\section*{Version V50 append-only record}
\addcontentsline{toc}{section}{Version V50 append-only record}
The complete V49 source was retained before this continuation.  It had
(491712) bytes, (14821) lines, and SHA--256
\[
 \texttt{83DAF8EB645AD8F5533CBA91DDFEE9E40BCA59E9CF6A82B98BA5DE4E03993CAB}.
\]
V50 audited Yang--Mills V53 and Navier--Stokes V26, separated the new
operator theorem from the V86--V88 archive, and proved strong-core and
two-sided weak-core criteria for trace-norm reconstruction convergence.


\section{V51: Covariance-weighted metric response and the missing orientation}
\label{sec:v16v51}

V50 asks whether the trace-class carriers of V30--V32 already imply
the Sobolev reconstruction convergence required in V49.  They do not:
they control a Gram operator on the exterior source space, whereas a
metric observable depends on the oriented response vector in the
augmented interior carrier.  This section identifies the exact factor,
proves the Hilbert--Schmidt convergence that follows once its orientation
is fixed, and shows that one additional Hilbert--Schmidt Sobolev map is
both sufficient and sharp in the unstructured Schatten scale.

\subsection{The exact augmented response factor}

Let
\[
 \mathscr A={\cal Z}\oplus{\cal H}_r                 \tag{51.1}
\]
be the common augmented energy carrier of V30--V31, and let
\({\cal H}_E\) be the common exterior carrier.  In this section assume
the regular form at cutoff \(n\) is represented by
\[
 {mathfrak a}_n(y)=\langle y,A_{r,n}y\rangle,
 \qquad A_{r,n}\succeq\beta I.                       \tag{51.2}
\]
Define the energy-response amplitude
\[
 F_n:{\cal H}_E\longrightarrow\mathscr A,
 \qquad
 F_nv=
 \left(\Gamma_nv,,A_{r,n}^{-1/2}D_nv\right).         \tag{51.3}
\]

\begin{proposition}[Augmented response Gram identity]
\label{prop:v16v51-Gram}
The complete retained influence of V31 satisfies
\[
 F_n^*F_n
 =\Gamma_n^*\Gamma_n+D_n^*A_{r,n}^{-1}D_n
 =J_n.                                                \tag{51.4}
\]
If \(S_n\geq0\) is the exterior covariance and
\[
 G_n:=F_nS_n^{1/2},                                   \tag{51.5}
\]
then
\[
 G_n^*G_n=S_n^{1/2}J_nS_n^{1/2}=K_n.                 \tag{51.6}
\]
Consequently \(G_n\) is Hilbert--Schmidt and
\[
 \|G_n\|_2^2=\operatorname{Tr}K_n.                   \tag{51.7}
\]
\end{proposition}

\begin{proof}
The two components in the orthogonal sum (51.1) give
\[
 \langle F_nu,F_nv\rangle
 =\langle u,\Gamma_n^*\Gamma_nv\rangle
  +\langle u,D_n^*A_{r,n}^{-1}D_nv\rangle.            \tag{51.8}
\]
This is the finite-cutoff version of the split (31.15), hence equals
the Schur influence \(J_n\).  Multiplying (51.4) on both sides by
\(S_n^{1/2}\) proves (51.6).  A bounded operator is Hilbert--Schmidt
exactly when its positive Gram is trace class, with the norm identity
(51.7).
\end{proof}

The first coordinate in (51.3) is the bounded soft energy coordinate
of V30.  The second is the regular minimizer expressed in its energy
coordinate: if \(y_n^{\min}=-A_{r,n}^{-1}D_nv\), then
\(A_{r,n}^{1/2}y_n^{\min}=-A_{r,n}^{-1/2}D_nv\).
Thus (51.3) retains the geometric direction that the invariant Gram
\(J_n\) forgets.

\subsection{What covariance-weighted trace convergence pays}

Let \(E_M\uparrow I\) be finite-rank projections on
\({\cal H}_E\), with their ranges contained in a common dense core.

\begin{lemma}[The Gram supplies an aligned Hilbert--Schmidt input tail]
\label{lem:v16v51-Gram-tail}
Suppose
\[
 K_n=G_n^*G_n\longrightarrow K
 \quad\hbox{in trace norm}.                           \tag{51.9}
\]
Then
\[
 \lim_{M\to\infty}\sup_n
 \|G_n(I-E_M)\|_2=0.                                 \tag{51.10}
\]
\end{lemma}

\begin{proof}
Cyclicity of the trace and \(E_M^2=E_M\) give
\[
 \|G_n(I-E_M)\|_2^2
 =\operatorname{Tr}\!left((I-E_M)K_n(I-E_M)\right).
                                                               \tag{51.11}
\]
For the limit operator, the right side tends to zero because
\(K^{1/2}(I-E_M)\to0\) in Hilbert--Schmidt norm.  Moreover,
\[
 \left|
 \operatorname{Tr}\!left((I-E_M)(K_n-K)(I-E_M)\right)
 \right|
 \leq\|K_n-K\|_1.                                    \tag{51.12}
\]
This controls all sufficiently large \(n\); enlarging \(M\) controls
the finitely many remaining positive trace-class operators.  Taking
square roots proves (51.10).
\end{proof}

\begin{theorem}[Gram plus strong columns determines the amplitude]
\label{thm:v16v51-amplitude}
Assume (51.9), and suppose \((G_nx)_n\) converges strongly in
\(\mathscr A\) for every \(x\) in the common core
\(\bigcup_ME_M{\cal H}_E\).  Then there is a unique
\(G\in\mathfrak S_2({\cal H}_E,\mathscr A)\) such that
\[
 \|G_n-G\|_2\longrightarrow0,
 \qquad G^*G=K.                                       \tag{51.13}
\]
\end{theorem}

\begin{proof}
For fixed \(M\), an orthonormal basis
\(e_1,\ldots,e_{d_M}\) of \(E_M{\cal H}_E\) gives
\[
 \|(G_n-G_m)E_M\|_2^2
 =\sum_{j=1}^{d_M}\|(G_n-G_m)e_j\|_{\mathscr A}^2
 \longrightarrow0.                                  \tag{51.14}
\]
Lemma~\ref{lem:v16v51-Gram-tail} and the triangle inequality now show
that \((G_n)\) is Cauchy in \(\mathfrak S_2\), by the same
finite-block argument as Theorem~\ref{thm:v16v50-strong-core} with
the Schatten index (2).  Let its Hilbert--Schmidt limit be \(G\).
The product estimate (24.8) gives
\[
 \|G_n^*G_n-G^*G\|_1
 \leq(\|G_n\|_2+\|G\|_2)\|G_n-G\|_2\longrightarrow0. \tag{51.15}
\]
Since (51.9) has the unique trace-norm limit \(K\), one has
\(G^*G=K\).  Hilbert--Schmidt convergence also identifies the given
strong column limits, proving uniqueness.
\end{proof}

Theorem~\ref{thm:v16v32-weighted} supplies (51.9), conditional on the
V31--V32 hypotheses.  It does not supply the strong columns in
Theorem~\ref{thm:v16v51-amplitude}.  Proposition
\ref{prop:v16v31-minimizers} approaches that requirement only when its
additional norm-recovery hypothesis holds; Mosco convergence alone
gives weak convergence of the regular minimizer and does not align the
moving energy coordinate \(A_{r,n}^{1/2}y_n^{\min}\).

There is no contradiction with the invariant strategy of V27--V28.
For the Schur influence, only \(G_n^*G_n\) is physical, so its output
orientation may be discarded.  A reconstructed metric tensor is a
non-Gram observable and retains that orientation.

\begin{proposition}[Gram convergence alone does not orient a response]
\label{prop:v16v51-orientation}
Trace-norm convergence of \(K_n=G_n^*G_n\) does not imply
Hilbert--Schmidt convergence of \(G_n\), even when every \(K_n\) is
the same rank-one projection.
\end{proposition}

\begin{proof}
Use the already-recorded family (49.9), now viewed as
\(G_n:\mathbb R\to\ell^2\).  Its Gram on \(\mathbb R\) is the scalar
identity for every \(n\), while its output unit vector moves from shell
to shell and the family is not Hilbert--Schmidt Cauchy.  This imports
the V49 example and gives its precise V32 interpretation.
\end{proof}

Equivalently, the polar decompositions
\[
 G_n=U_nK_n^{1/2}                                    \tag{51.16}
\]
contain a partial-isometry datum \(U_n\) invisible to \(K_n\).  The
Powers--St\o rmer estimate already recorded in (9.23) gives
\(K_n^{1/2}\to K^{1/2}\) in Hilbert--Schmidt norm; one still needs
alignment of \(U_n\) on the range of the limiting square root.  Strong
column convergence is a basis-free way to state exactly that alignment.

\subsection{The additional Sobolev half-factor}

For derivative order \(j=0,1,2\), let
\[
 T_n^{(j)}:\mathscr A\longrightarrow\mathscr Y^{(j)} \tag{51.17}
\]
reconstruct the (j)-th metric derivative in the common shell-collar
carrier of V49.  Define the covariance-weighted reconstructed response
\[
 \mathcal R_{n,\mathrm{cov}}^{(j)}
 :=T_n^{(j)}G_n:
 {cal H}_E\longrightarrow\mathscr Y^{(j)}.           \tag{51.18}
\]

\begin{theorem}[Two Hilbert--Schmidt halves give the trace-class metric
response]
\label{thm:v16v51-two-halves}
Suppose Theorem~\ref{thm:v16v51-amplitude} holds and, for a fixed (j),
\[
 T_n^{(j)}\longrightarrow T^{(j)}
 \quad\hbox{in }\mathfrak S_2(\mathscr A,\mathscr Y^{(j)}).     \tag{51.19}
\]
Then
\[
 \mathcal R_{n,\mathrm{cov}}^{(j)}
 \longrightarrow T^{(j)}G
 \quad\hbox{in trace norm},                           \tag{51.20}
\]
with the quantitative bound
\[
\begin{split}
 \|T_n^{(j)}G_n-T^{(j)}G\|_1
 \leq{}&
 \|T_n^{(j)}-T^{(j)}\|_2\|G_n\|_2\\
 &+\|T^{(j)}\|_2\|G_n-G\|_2.                       \tag{51.21}
\end{split}
\]
Consequently the V49 weighted shell compiler applies to
\(\mathcal R_{n,\mathrm{cov}}^{(j)}\), provided its geometry-array
condition (49.18) is verified.
\end{theorem}

\begin{proof}
Insert (T^{(j)}G_n) and use the Schatten--H\"older inequality
\(\|AB\|_1\leq\|A\|_2\|B\|_2\).  The norms \(\|G_n\|_2\) are bounded
by (51.7) and trace-norm convergence of \(K_n\).  Equations
(51.13) and (51.19) make the right side of (51.21) vanish.
Trace-norm convergence is precisely the operator hypothesis used by
Theorem~\ref{thm:v16v49-compiler}.
\end{proof}

Condition (51.19) itself has a finite-block acquisition card.  If
\(Q_L\uparrow I\) are finite-rank projections on \(\mathscr A\), their
ranges lie in a common core, \(T_n^{(j)}a\) converges for every core
vector, and
\[
 \lim_{L\to\infty}\sup_n
 \|T_n^{(j)}(I-Q_L)\|_2=0,                            \tag{51.22}
\]
then the proof of Theorem~\ref{thm:v16v50-strong-core}, with trace norm
replaced by Hilbert--Schmidt norm, proves (51.19).  Thus the new physical
estimate is an aligned (H^j)-collar Hilbert--Schmidt tail in the
augmented energy coordinate.

\begin{counterexample}[A bounded Sobolev map is one ideal order short]
\label{cnt:v16v51-one-order}
Let every carrier be \(\ell^2(\mathbb N)\), put
\[
 Ge_m={1\over m}e_m,qquad T=I.                       \tag{51.23}
\]
Then
\[
 K=G^*G=\operatorname{diag}(m^{-2})\in\mathfrak S_1,
 \qquad G\in\mathfrak S_2,                           \tag{51.24}
\]
but
\[
 TG=G\notin\mathfrak S_1.                            \tag{51.25}
\]
Hence boundedness or operator-norm convergence of the metric
reconstruction does not turn the V32 covariance carrier into the V49
trace-class response.
\end{counterexample}

\begin{proof}
The series \(\sum_m m^{-2}\) converges, whereas
\(\sum_m m^{-1}\) diverges.  These are respectively the trace of
\(K\) and the trace norm of \(G\).
\end{proof}

Within the unstructured Schatten--H\"older scale, (51.19) is the
weakest symmetric extra ideal assumption: V32 gives only an
\(\mathfrak S_2\) amplitude, so its partner must lie in
\(\mathfrak S_p\) with \(p\leq2\) to force an
\(\mathfrak S_1\) product.  Counterexample
\ref{cnt:v16v51-one-order} excludes the endpoint \(p=\infty\).

\subsection{Why covariance weighting cannot be inverted globally}

\begin{proposition}[No bounded deweighting on an infinite carrier]
\label{prop:v16v51-no-inverse}
If \(\mathscr A\) is infinite dimensional and
\(G:{\cal H}_E\to\mathscr A\) is Hilbert--Schmidt, there is no bounded
operator \(B:\mathscr A\to{\cal H}_E\) satisfying
\[
 GB=I_{\mathscr A}.                                   \tag{51.26}
\]
\end{proposition}

\begin{proof}
A Hilbert--Schmidt operator is compact, so (GB) is compact.  The
identity on an infinite-dimensional Hilbert space is not compact.
\end{proof}

Therefore trace-class convergence of (T_n^{(j)}G_n) cannot be
multiplied by a uniformly bounded inverse to recover trace-class
convergence of the deterministic reconstruction (T_n^{(j)}).  On the
finite-dimensional soft space \({\cal Z}\), a uniformly bounded right
inverse may exist; in that case it recovers only the finite soft block.
Indeed, if (B_n:{\cal Z}\to{\cal H}_E) satisfies
\(G_nB_n=\iota_{\cal Z}\), (B_n\to B) in operator norm, and
\(T_nG_n\to R\) in trace norm, then
\[
 T_n\iota_{\cal Z}=(T_nG_n)B_n\longrightarrow RB     \tag{51.27}
\]
in trace norm.  Proposition~\ref{prop:v16v51-no-inverse} prevents this
finite-block observation from covering the infinite regular carrier.

\subsection{Result of the V30--V32 test}

The exact outcome is:
\begin{enumerate}
\item V31--V32 conditionally prove (K_n\to K) in trace norm.  This
      automatically pays the aligned input \(\mathfrak S_2\) tail
      (51.10), but not an \(\mathfrak S_1\) tail.
\item Strong convergence of the oriented augmented energy responses is
      still required.  The missing regular datum is convergence of
      (A_{r,n}^{-1/2}D_nS_n^{1/2}) in a common energy carrier, not only
      convergence of its Gram.
\item An aligned Hilbert--Schmidt Sobolev reconstruction tail (51.22)
      then gives a trace-class \emph{covariance-weighted} metric response
      by Theorem~\ref{thm:v16v51-two-halves}.
\item The deterministic unit-direction curvature owner of V47--V49 is
      stronger.  It still requires the trace-class input tail (50.19),
      unless the state-distance argument is replaced by an averaged
      Duhamel or entropy estimate that can use (51.20) directly.
\end{enumerate}

\begin{assessment}[What V51 closes]
\label{ass:v16v51-status}
V51 locates the precise interface between the earlier covariance
programme and the metric collar programme.  The V32 carrier is the Gram
of an augmented energy-response amplitude.  Its trace-norm convergence
provides a uniform Hilbert--Schmidt input tail; strong response columns
then determine a Hilbert--Schmidt amplitude.  Composing it with one
Hilbert--Schmidt Sobolev reconstruction produces the trace-class
covariance-weighted metric response needed by the V49 shell compiler.

Neither oriented regular-energy convergence nor the Sobolev
Hilbert--Schmidt tail has been proved for the physical regulator.
Moreover, covariance weighting cannot be inverted on the
infinite-dimensional regular carrier, so the deterministic owner remains
open rather than being silently identified with the averaged response.
\end{assessment}

\begin{decision}[V52 averaged curvature-owner route]
\label{dec:v16v51-next}
Go upstream to the state-distance estimate.  Replace the worst-case
metric-mode diameter used in the V47 curvature owner by an interpolated
relative-entropy or Duhamel bound whose cost is an expectation of the
shell response.  Determine whether the trace-class covariance-weighted
operators (51.20), together with a strict exponential reserve, make that
expected cost summable.  If they do not, isolate the exact non-Gaussian
metric moment that must be acquired.
\end{decision}

\begin{firewall}[Claim boundary after V51]
\label{fire:v16v51}
V51 proves a conditional operator-ideal bridge and a sharp obstruction.
It does not establish the physical response orientation, the Sobolev
Hilbert--Schmidt tail, a positive gravitational measure, a limiting
Einstein equation, or the full Standard Model--Einstein continuum.
\end{firewall}

\section*{Version V51 append-only record}
\addcontentsline{toc}{section}{Version V51 append-only record}
The complete V50 source was retained before this continuation.  It had
(503042) bytes, (15115) lines, and SHA--256
\[
 \texttt{708BD645C2464FF81B5DEE545E2E63FEB8C60F0E355B0089797B3BF26D1E7D4D}.
\]
V51 tested V30--V32 against the V49--V50 reconstruction criterion,
proved the Gram-plus-strong-columns Hilbert--Schmidt theorem and the
two-Hilbert--Schmidt-factor trace-class compiler, and showed why neither
Gram convergence nor a bounded Sobolev map closes the physical metric
owner.


\section{V52: Entropy interpolation for averaged curvature-shell owners}
\label{sec:v16v52}

The deterministic curvature debit (47.16) multiplies a worst-case
gradient by the diameter of a carrier good set.  V51 shows that the
earlier covariance programme naturally controls averaged response
directions instead.  This section goes upstream to the owner comparison:
an exact entropy identity replaces the diameter by an interpolated
variance.  A covariance-weighted Poincar\'e inequality then loads that
variance from the trace-class response constructed in V51.

\subsection{Exact entropy of one Gibbs insertion}

Let \(\mu_0\) be a probability measure, let \(V\) be real and
measurable, and suppose
\[
 Z_t:=\mu_0(e^{-tV})\in(0,\infty),
 \qquad
 d\mu_t={e^{-tV}\over Z_t}\,d\mu_0,qquad 0\leq t\leq1.          \tag{52.1}
\]
Assume dominated differentiation is valid through second order; it is
enough that \(V^2e^{-tV}\) have one \(\mu_0\)-integrable envelope for
\(0\leq t\leq1\).

\begin{theorem}[Exact Gibbs-path entropy identity]
\label{thm:v16v52-entropy}
Under (52.1),
\[
 {d\over dt}\log Z_t=-\mu_t(V),
 \qquad
 {d\over dt}\mu_t(V)=-\operatorname{Var}_{\mu_t}(V),           \tag{52.2}
\]
and
\[
 \boxed{
 D(\mu_t\Vert\mu_0)
 =\int_0^t s\,\operatorname{Var}_{\mu_s}(V)\,ds.}     \tag{52.3}
\]
In particular, in the state norm (35.3),
\[
 \|\mu_1-\mu_0\|_{\mathrm{st}}
 \leq
 \left(
 2\int_0^1t\,\operatorname{Var}_{\mu_t}(V)\,dt
 \right)^{1/2}.                                      \tag{52.4}
\]
\end{theorem}

\begin{proof}
Differentiate the normalized integral in (52.1).  The partition
function gives the first identity in (52.2), and the quotient rule gives
\[
 {d\over dt}\mu_t(V)
 =-\mu_t(V^2)+\mu_t(V)^2.                             \tag{52.5}
\]
Moreover,
\[
 D(\mu_t\Vert\mu_0)
 =\mu_t\!\left(-tV-\log Z_t\right)
 =-t\mu_t(V)-\log Z_t.                               \tag{52.6}
\]
Differentiating (52.6) and using (52.2) yields
\[
 {d\over dt}D(\mu_t\Vert\mu_0)
 =t\operatorname{Var}_{\mu_t}(V).                   \tag{52.7}
\]
The entropy vanishes at \(t=0\), proving (52.3).  Pinsker in the
normalization (35.7) proves (52.4).
\end{proof}

For completeness, the reverse entropy satisfies
\[
 D(\mu_0\Vert\mu_1)
 =\int_0^1(1-t)\operatorname{Var}_{\mu_t}(V)\,dt,    \tag{52.8}
\]
so the symmetrized entropy is exactly
\(\int_0^1\operatorname{Var}_{\mu_t}(V)dt\).  Formula (52.3) is
slightly sharper for the directed owner path because it retains the
factor \(t\).

\subsection{Countably many averaged owners}

At cutoff \(n\), let \(\mu_{n,q,0}\) be the state after all owners
preceding shell \(q\) have been inserted, and let
\[
 d\mu_{n,q,t}
 ={e^{-tV_{n,q}}\over\mu_{n,q,0}(e^{-tV_{n,q}})}
 d\mu_{n,q,0}.                                      \tag{52.9}
\]
No independence between distinct shells is assumed.

\begin{theorem}[Variance-tail owner compiler]
\label{thm:v16v52-owner-tail}
Define
\[
 \mathcal V_{n,q}
 :=\int_0^1t\,
 \operatorname{Var}_{\mu_{n,q,t}}(V_{n,q})\,dt.      \tag{52.10}
\]
If
\[
 \lim_{Q\to\infty}\sup_n
 \sum_{q\geq Q}\sqrt{2\mathcal V_{n,q}}=0,          \tag{52.11}
\]
then the countable shell owners have a vanishing uniform
state-distance tail.
\end{theorem}

\begin{proof}
Theorem~\ref{thm:v16v52-entropy} gives the certified debit
\[
 d_{n,q}:=\sqrt{2\mathcal V_{n,q}}                   \tag{52.12}
\]
for the \(q\)-th insertion, evaluated under its actual already tilted
state.  Apply the countable-owner and triangular-array telescoping
theorems of V36.  Hypothesis (52.11) is exactly their required uniform
tail condition.
\end{proof}

This compiler contains no carrier diameter and no good-set deletion.
Its price is a variance estimate under every interpolated law
\(\mu_{n,q,t}\), rather than a deterministic oscillation estimate.

\subsection{Covariance-weighted Poincar\'e loading}

Work first at a finite cutoff, so the exterior source carrier
\({\cal H}_{E,n}\) is finite dimensional.  Let \(\rho_{n,q,t}\) be
the source-coordinate representation of \(\mu_{n,q,t}\).  Let
\[
 \mathcal R_n^{(j)}
 :=\mathcal R_{n,\mathrm{cov}}^{(j)}
 :{\cal H}_{E,n}\longrightarrow\mathscr Y^{(j)}     \tag{52.13}
\]
be the covariance-weighted metric response (51.18), transported to the
common collar carrier, and put
\[
 \mathscr Y=\bigoplus_{j=0}^2\mathscr Y^{(j)},
 \qquad
 \mathcal R_n=\bigoplus_{j=0}^2\mathcal R_n^{(j)}.   \tag{52.14}
\]
Let \(P_q\) denote the projection onto the three derivative components
over shell \(q\).

Suppose the curvature shell has the coordinate form
\[
 V_{n,q}(\xi)=\Phi_{n,q}(P_q\mathcal R_n\xi),        \tag{52.15}
\]
on the source chart.  This notation permits a nonlinear shell
functional \(\Phi_{n,q}\); its Fr\'echet gradient is represented by
\(\nabla\Phi_{n,q}\in P_q\mathscr Y\).

\begin{hypothesis}[Interpolated covariance Poincar\'e card]
\label{hyp:v16v52-Poincare}
There is a cutoff- and shell-independent \(C_P<\infty\) such that, for
every smooth cylinder function \(f\),
\[
 \operatorname{Var}_{\rho_{n,q,t}}(f)
 \leq C_P\int\|\nabla_\xi f\|_{{\cal H}_{E,n}}^2
 \,d\rho_{n,q,t}                                    \tag{52.16}
\]
for every \(n,q\) and \(t\in[0,1]\).
\end{hypothesis}

\begin{theorem}[Trace-response variance loader]
\label{thm:v16v52-Poincare-loader}
Under Hypothesis~\ref{hyp:v16v52-Poincare}, suppose
\[
 B_{n,q}^2
 :=\int_0^1t\int
 \|\nabla\Phi_{n,q}(P_q\mathcal R_n\xi)\|_{\mathscr Y}^2
 \,d\rho_{n,q,t}(\xi)\,dt<\infty.                  \tag{52.17}
\]
Then
\[
 \mathcal V_{n,q}
 \leq C_P B_{n,q}^2
 \|P_q\mathcal R_n\|_{\rm op}^2,                   \tag{52.18}
\]
and the shell owner has the certified debit
\[
 d_{n,q}
 \leq\sqrt{2C_P}\,B_{n,q}
 \|P_q\mathcal R_n\|_{\rm op}.                     \tag{52.19}
\]
\end{theorem}

\begin{proof}
The chain rule in the source coordinate gives
\[
 \nabla_\xi V_{n,q}
 =\mathcal R_n^*P_q
  \nabla\Phi_{n,q}(P_q\mathcal R_n\xi).              \tag{52.20}
\]
Therefore
\[
 \|\nabla_\xi V_{n,q}\|
 \leq\|P_q\mathcal R_n\|_{\rm op}
 \|\nabla\Phi_{n,q}\|.                              \tag{52.21}
\]
Insert (52.21) in (52.16), multiply by \(t\), and integrate over
\([0,1]\).  Equations (52.10) and (52.17) give (52.18), while
Theorem~\ref{thm:v16v52-entropy} gives (52.19).
\end{proof}

\subsection{Uniform cofinal trace-class compiler}

The shell ranges in (52.14) are orthogonal.  Consequently,
\[
 \sum_{q\geq Q}\|P_q\mathcal R_n\|_2^2
 =\|P_{\geq Q}\mathcal R_n\|_2^2.                  \tag{52.22}
\]

\begin{theorem}[Averaged curvature-tail criterion]
\label{thm:v16v52-curvature-tail}
Suppose \(\mathcal R_n\to\mathcal R\) in trace norm on the common
carrier, Hypothesis~\ref{hyp:v16v52-Poincare} holds, and
\[
 A_{n,Q}:=
 \left(\sum_{q\geq Q}B_{n,q}^2\right)^{1/2}         \tag{52.23}
\]
satisfies
\[
 \lim_{Q\to\infty}\sup_n
 A_{n,Q}\|P_{\geq Q}\mathcal R_n\|_2=0.            \tag{52.24}
\]
Then the curvature shells have a vanishing uniform state-distance
tail.  In particular, (52.24) holds if
\[
 \sup_{n,Q}A_{n,Q}<\infty.                           \tag{52.25}
\]
\end{theorem}

\begin{proof}
Equations (52.19), (52.22), and Cauchy--Schwarz in the shell index give
\[
\begin{split}
 \sum_{q\geq Q}d_{n,q}
 &\leq\sqrt{2C_P}
 \sum_{q\geq Q}B_{n,q}\|P_q\mathcal R_n\|_2\\
 &\leq\sqrt{2C_P}\,
 A_{n,Q}\|P_{\geq Q}\mathcal R_n\|_2.              \tag{52.26}
\end{split}
\]
Thus (52.24) implies (52.11).  Trace-norm convergence implies
Hilbert--Schmidt convergence and, by the strong-projection argument of
V49, the uniform tail
\[
 \lim_{Q\to\infty}\sup_n
 \|P_{\geq Q}\mathcal R_n\|_2=0.                   \tag{52.27}
\]
This proves the last assertion under (52.25).
\end{proof}

Theorem~\ref{thm:v16v52-curvature-tail} is an averaged counterpart of
Theorem~\ref{thm:v16v49-compiler}.  The V49 arrays contain collar
volumes and powers of the collar width because they convert a
worst-case unit mode to \(L^1\).  Here those geometric losses occur
inside the root-mean-square shell functional \(B_{n,q}\).

\subsection{Loading the paired Einstein--GHY shell}

Let
\(b_{n,q}(\xi),f_{n,q}(\xi),j_{n,q}(\xi)\) denote the norms of the
Riesz functionals on the bulk, smooth-face, and joint components of
the shell Sobolev carrier.  The imported variational identity used in
V47 gives the pointwise bound
\[
 \|\nabla\Phi_{n,q}\|
 \leq|\delta c_{R,n,q}|
 \left(K_Gb_{n,q}+K_\Pi f_{n,q}+K_{\mathcal J}j_{n,q}\right).   \tag{52.28}
\]

\begin{corollary}[RMS paired-curvature loader]
\label{cor:v16v52-EGHY}
The quantity (52.17) is bounded by
\[
\begin{split}
 B_{n,q}^2\leq |\delta c_{R,n,q}|^2
 \int_0^1t\,\mathbb E_{\rho_{n,q,t}}
 \bigl[&K_Gb_{n,q}+K_\Pi f_{n,q}\\
       &+K_{\mathcal J}j_{n,q}\bigr]^2dt.            \tag{52.29}
\end{split}
\]
The trace inequalities of V48 bound \(f_{n,q}\) and \(j_{n,q}\) by
the corresponding random bulk \(W^{1,1}\) and \(W^{2,1}\) collar
functionals, with every \(\ell_{n,q}^{-1}\) and
\(\ell_{n,q}^{-2}\) factor retained.
\end{corollary}

\begin{proof}
Square (52.28) and integrate with the weight appearing in (52.17).
The V48 trace bounds are pointwise deterministic inequalities, so they
remain valid before taking the interpolated expectation.
\end{proof}

Thus the deterministic suprema
\(\mathfrak B_{n,q},\mathfrak F_{n,q},\mathfrak J_{n,q}\) in (47.13)
are replaced by an interpolated second moment.  This is a genuine
weakening only if the laws \(\rho_{n,q,t}\) possess a uniform
Poincar\'e constant.

\subsection{The non-Gaussian bottleneck is independent}

For an unperturbed standard Gaussian source, (52.16) holds with
\(C_P=1\).  A uniformly bounded oscillatory density perturbation also
preserves a Poincar\'e inequality by the bounded-perturbation lemma,
with a constant depending exponentially on the oscillation.  The
projected interacting input in V32 is weaker: it assumes convergence
of densities in \(L^p\).

\begin{counterexample}[An \(L^p\) density does not imply Poincar\'e]
\label{cnt:v16v52-Lp}
Let \(\gamma\) be the standard Gaussian law on \(\mathbb R\), let
\[
 A=[-2,-1]\cup[1,2],
 \qquad h={\mathbf1_A\over\gamma(A)},                \tag{52.30}
\]
and put \(h_n=h\).  Then \(h_n\to h\) in every finite
\(L^p(\gamma)\), but the probability \(h\,d\gamma\) admits no finite
Poincar\'e constant for the usual gradient.
\end{counterexample}

\begin{proof}
Choose a smooth function \(f\) equal to (-1) on \([-2,-1]\), equal
to (1) on \([1,2]\), and varying only in the gap ((-1,1)).  Its
variance under (h\,d\gamma) is positive, while (f'=0)
\(h\,d\gamma\)-almost everywhere.  Hence an inequality of the form
(52.16) is impossible.  The density is bounded and supported on a set
of finite Gaussian measure, so it belongs to every finite \(L^p\).
\end{proof}

Consequently, V32's \(L^p\) interacting handoff cannot silently provide
Hypothesis~\ref{hyp:v16v52-Poincare}.  The latter needs a separate
connectedness or convexity input for the interpolated metric-source
law.

\begin{assessment}[What V52 closes]
\label{ass:v16v52-status}
V52 proves the exact relative-entropy cost of one positive shell
insertion and a countable variance-tail compiler.  Under a uniform
interpolated Poincar\'e inequality, the covariance-weighted trace-class
metric response of V51 controls this cost, and the paired
Einstein--GHY variation enters through an explicit RMS shell functional.
The result removes the bounded carrier diameter from this alternative
owner route.

The physical interpolated interacting law has not been shown to obey a
uniform Poincar\'e inequality, and the V32 \(L^p\) density convergence
does not imply one.  The physical response orientation and Sobolev
Hilbert--Schmidt estimates isolated in V51 also remain required.
\end{assessment}

\begin{decision}[V53 conditional convexity and Poincar\'e acquisition]
\label{dec:v16v52-next}
Test whether the augmented soft--regular energy of V30 yields a uniform
Poincar\'e constant after conditioning on the hard Standard Model
variables.  Prove a block-Hessian or conditional-variance criterion
that tolerates a finite soft sector, and identify the exact negative
curvature contributed by the Einstein and matter interactions.  If
global convexity fails, retain a two-scale conditional criterion rather
than assuming a nonexistent log-concave gravitational measure.
\end{decision}

\begin{firewall}[Claim boundary after V52]
\label{fire:v16v52}
Entropy interpolation converts an owner comparison into a variance
problem; it does not prove the Poincar\'e inequality or the physical
metric moments used to bound that variance.  V52 does not construct a
positive gravitational measure, recover the limiting Einstein equation,
or complete the Standard Model--Einstein continuum.
\end{firewall}

\section*{Version V52 append-only record}
\addcontentsline{toc}{section}{Version V52 append-only record}
The complete V51 source was retained before this continuation.  It had
(517791) bytes, (15500) lines, and SHA--256
\[
 \texttt{AC978D5F1DC526FB723A160A8880816111C260D05B228BC26BAAD4DAAC6C860C}.
\]
V52 replaced the diameter curvature debit by an exact entropy--variance
identity, loaded the averaged shell variance from the covariance-weighted
trace response under a Poincar\'e card, and proved that the existing
\(L^p\) interacting-density hypothesis does not supply that card.


\section{V53: A two-scale Poincar\'e criterion for the augmented metric law}
\label{sec:v16v53}

V52 isolates a uniform Poincar\'e inequality for every interpolated
shell law as the next acquisition.  Global convexity is too strong for
the intended gravitational measure, while a gap only in the regular
and hard variables leaves the finite soft coordinate uncontrolled.
This section proves a two-scale criterion: a conditional hard
Poincar\'e constant, a Poincar\'e constant for the induced soft
marginal, and one mixed-Hessian bound imply a Poincar\'e inequality for
the full law.  The soft marginal constant is then obtained from the
exact free-energy Hessian, including its negative covariance term.

The frozen archive already contains tilt-uniform one-site quartic
Poincar\'e estimates, colored hard-sector assembly, and a
finite-dimensional Bakry--Emery compiler under global strong
convexity.  Those results are not repeated.  The new point is the
soft-marginal decomposition tailored to the V30 energy coordinate and
the V52 interpolated curvature owner.

\subsection{Conditional and marginal data}

At a finite cutoff let \(z\in\mathbb R^m\) be the augmented soft energy
coordinate and let \(r\in\mathbb R^N\) collect the regular metric and
positive hard Standard Model coordinates.  Consider
\[
 d\mu(z,r)=\alpha(dz)\,\kappa_z(dr),                 \tag{53.1}
\]
where
\[
 \kappa_z(dr)
 ={e^{-H(z,r)}\over Z(z)}\,dr,
 \qquad
 Z(z)=\int_{\mathbb R^N}e^{-H(z,r)},dr.             \tag{53.2}
\]
All derivatives below are first proved for smooth compactly supported
functions and extended by closure.  Assume the usual dominated
differentiation conditions for (53.2).

\begin{hypothesis}[Two-scale inputs]
\label{hyp:v16v53-two-scale}
There are constants \(C_h,C_s<\infty\) and \(b<\infty\) such that:
\begin{enumerate}
\item every conditional law satisfies
\[
 \operatorname{Var}_{\kappa_z}(f)
 \leq C_h\int\|\nabla_rf\|^2,d\kappa_z;             \tag{53.3}
\]
\item the soft marginal satisfies
\[
 \operatorname{Var}_{\alpha}(g)
 \leq C_s\int\|\nabla_zg\|^2,d\alpha;              \tag{53.4}
\]
\item the mixed Hessian obeys
\[
 \|\nabla_r\nabla_zH(z,r)\|_{\rm op}\leq b.         \tag{53.5}
\]
\end{enumerate}
\end{hypothesis}

\begin{lemma}[Derivative of conditional expectation]
\label{lem:v16v53-conditional-derivative}
For
\[
 \bar f(z):=\int f(z,r)\,d\kappa_z(r),               \tag{53.6}
\]
one has, in every direction \(a\in\mathbb R^m\),
\[
 \partial_a\bar f
 =\kappa_z(\partial_af)
 -\operatorname{Cov}_{\kappa_z}(f,\partial_aH).      \tag{53.7}
\]
Under (53.3) and (53.5),
\[
 \|\nabla_z\bar f(z)\|
 \leq
 \bigl(\kappa_z\|\nabla_zf\|^2\bigr)^{1/2}
 +C_hb\bigl(\kappa_z\|\nabla_rf\|^2\bigr)^{1/2}.   \tag{53.8}
\]
\end{lemma}

\begin{proof}
Differentiate the normalized integral (53.6); the quotient rule gives
(53.7).  For a unit vector \(a\), conditional Poincar\'e and (53.5)
give
\[
 \operatorname{Var}_{\kappa_z}(\partial_aH)
 \leq C_h\kappa_z\|\nabla_r\partial_aH\|^2
 \leq C_hb^2.                                        \tag{53.9}
\]
The same Poincar\'e inequality and Cauchy--Schwarz give
\[
 |\operatorname{Cov}_{\kappa_z}(f,\partial_aH)|
 \leq C_hb
 \bigl(\kappa_z\|\nabla_rf\|^2\bigr)^{1/2}.        \tag{53.10}
\]
Also
\(|\kappa_z(\partial_af)|
\leq(\kappa_z\|\nabla_zf\|^2)^{1/2}\|a\|\).
Take the supremum over unit \(a\) in (53.7).
\end{proof}

\subsection{The two-scale full Poincar\'e theorem}

\begin{theorem}[Soft-marginal plus conditional Poincar\'e]
\label{thm:v16v53-two-scale}
Under Hypothesis~\ref{hyp:v16v53-two-scale}, for every \(\delta>0\),
\[
 \operatorname{Var}_{\mu}(f)
 \leq C_z(\delta)\int\|\nabla_zf\|^2d\mu
      +C_r(\delta)\int\|\nabla_rf\|^2d\mu,          \tag{53.11}
\]
where
\[
\begin{split}
 C_z(\delta)&:=C_s(1+\delta),\\
 C_r(\delta)&:=C_h+C_s(1+\delta^{-1})C_h^2b^2.       \tag{53.12}
\end{split}
\]
Consequently \(\mu\) satisfies the ordinary Poincar\'e inequality
with
\[
 C_P(\mu)
 \leq\inf_{\delta>0}max\{C_z(\delta),C_r(\delta)\}. \tag{53.13}
\]
\end{theorem}

\begin{proof}
The law of total variance gives
\[
 \operatorname{Var}_{\mu}(f)
 =\int\operatorname{Var}_{\kappa_z}(f)\,d\alpha(z)
  +\operatorname{Var}_{\alpha}(\bar f).              \tag{53.14}
\]
The first term is bounded by (53.3).  Apply (53.4) to the second and
then use (53.8).  For nonnegative \(u,v\),
\[
 (u+v)^2\leq(1+\delta)u^2+(1+\delta^{-1})v^2.        \tag{53.15}
\]
Integration over \(z\) yields (53.11)--(53.12).  Bounding both
quadratic gradient terms by their larger coefficient proves (53.13),
and optimization over \(\delta\) is legitimate because the estimate
holds for every positive \(\delta\).
\end{proof}

No smallness condition on \(b\) is needed once both marginal constants
are known.  Smallness enters when the soft marginal constant itself is
derived from the hard conditional law.

\subsection{The effective soft Hessian and its covariance loss}

Assume the soft marginal has density
\[
 d\alpha(z)={e^{-\Psi(z)}\over\int e^{-\Psi}}\,dz,
 \qquad
 \Psi(z)=-\log Z(z)+\Psi_{\rm ext}(z),               \tag{53.16}
\]
where a purely soft term may be included in \(H\) instead of
\(\Psi_{\rm ext}\).

\begin{lemma}[Exact effective Hessian]
\label{lem:v16v53-effective-Hessian}
With all purely soft terms included in \(H\),
\[
 \nabla_z^2\Psi(z)
 =\kappa_z(\nabla_z^2H)
 -\operatorname{Cov}_{\kappa_z}
   (\nabla_zH,\nabla_zH).                             \tag{53.17}
\]
The covariance is the matrix whose quadratic form in direction \(a\)
is \(\operatorname{Var}_{\kappa_z}(\partial_aH)\).
\end{lemma}

\begin{proof}
The first derivative is
\(\partial_a\Psi=\kappa_z(\partial_aH)\).  Apply
Lemma~\ref{lem:v16v53-conditional-derivative} with
\(f=\partial_aH\) and differentiate in a second direction.  This is
the V41 free-energy Hessian identity in the present augmented
coordinate.
\end{proof}

\begin{theorem}[Conditional Schur-margin criterion]
\label{thm:v16v53-Schur}
Assume (53.3), (53.5), and
\[
 \kappa_z(\nabla_z^2H)\succeq aI_m                  \tag{53.18}
\]
uniformly in \(z\).  If
\[
 \lambda_s:=a-C_hb^2>0,                              \tag{53.19}
\]
then
\[
 \nabla_z^2\Psi\succeq\lambda_sI_m,
 \qquad C_P(\alpha)\leq\lambda_s^{-1}.              \tag{53.20}
\]
The full law therefore satisfies Theorem~\ref{thm:v16v53-two-scale}
with \(C_s=\lambda_s^{-1}\).
\end{theorem}

\begin{proof}
For every unit \(a\), (53.9) bounds the covariance quadratic form in
(53.17) by \(C_hb^2\).  Equations (53.18)--(53.19) give the Hessian
lower bound in (53.20).  The finite-dimensional Bakry--Emery theorem
gives \(C_P(\alpha)\leq\lambda_s^{-1}\).  Apply
Theorem~\ref{thm:v16v53-two-scale}.
\end{proof}

When the hard conditional potential is \(\beta\)-strongly convex,
one may take \(C_h\leq\beta^{-1}\), and (53.19) becomes
\[
 a\beta>b^2.                                         \tag{53.21}
\]
This is the strict Schur-complement condition for the idealized block
Hessian with diagonal floors \(a,\beta\) and cross norm \(b\).  The
conditional theorem is more flexible because (53.3) may come from a
quartic or compact interacting block that is not pointwise convex.

\subsection{Uniformity along every shell interpolation}

For the V52 owner, write the source-coordinate energy as
\[
 H_{n,q,t}(z,r)=H_{n,q,0}(z,r)+tV_{n,q}(z,r),
 \qquad0\leq t\leq1.                                 \tag{53.22}
\]
Let \(C_{h,n,q,t}\), \(b_{n,q,t}\), and \(a_{n,q,t}\) denote the
best constants in (53.3), (53.5), and (53.18).

\begin{corollary}[Uniform interpolated Poincar\'e card]
\label{cor:v16v53-interpolated}
Suppose there are numbers \(C_h,b,a,\lambda_*>0\), independent of
\(n,q,t\), such that
\[
 C_{h,n,q,t}\leq C_h,qquad
 b_{n,q,t}\leq b,qquad
 a_{n,q,t}\geq a,qquad
 a-C_hb^2\geq\lambda_*.                              \tag{53.23}
\]
Then all interpolated laws satisfy Hypothesis
\ref{hyp:v16v52-Poincare} with
\[
 C_P\leq
 \inf_{\delta>0}\max\left\{
 {1+\delta\over\lambda_*},
 C_h+{(1+\delta^{-1})C_h^2b^2\over\lambda_*}
 \right\}.                                          \tag{53.24}
\]
\end{corollary}

\begin{proof}
Theorem~\ref{thm:v16v53-Schur} gives the soft marginal constant
\(C_s\leq\lambda_*^{-1}\) uniformly.  Insert it in (53.12)--(53.13).
\end{proof}

The exact negative-curvature ledger behind (53.23) is
\[
\begin{split}
 a_{n,q,t}
 &:=\inf_{z,\|u\|=1}
 \kappa_{n,q,t,z}
 \bigl[\langle u,\nabla_z^2H_{n,q,t}u\rangle\bigr],\\
 b_{n,q,t}
 &:=\operatorname*{ess\,sup}_{z,r}
 \|\nabla_r\nabla_zH_{n,q,t}(z,r)\|_{\rm op}.       \tag{53.25}
\end{split}
\]
The covariance loss is then exactly at most
\(C_{h,n,q,t}b_{n,q,t}^2\).  It cannot be dropped from the effective
soft Hessian; doing so would repeat the free-energy error corrected in
V41.

For acquisition purposes one may split
\[
\begin{split}
 a_{n,q,t}&\geq
 a_{\rm aug}-\kappa_{\rm SM}-\kappa_{R,q},\\
 b_{n,q,t}&\leq b_{\rm base}+b_{R,q},                \tag{53.26}
\end{split}
\]
where \(a_{\rm aug}\) is the positive curvature of the normalized
soft energy coordinate, \(\kappa_{\rm SM}\) and
\(\kappa_{R,q}\) are the negative parts of the matter and curvature
shell Hessians, and the two \(b\)'s are mixed soft--hard Hessian norms.
Equation (53.23), with (53.26), is a numerical margin to be proved from
the physical regulator; it is not a consequence of the notation.

\subsection{A regular-sector gap cannot control an unconfined soft mode}

\begin{counterexample}[Separated soft wells]
\label{cnt:v16v53-soft-wells}
Let
\[
 \mu_n(dz,dr)
 =\frac12\gamma(dz-n)\gamma(dr)
  +\frac12\gamma(dz+n)\gamma(dr),                   \tag{53.27}
\]
where \(\gamma\) is the standard one-dimensional Gaussian law.  The
conditional law of \(r\) given \(z\) is standard Gaussian, its
Poincar\'e constant is one, and the mixed Hessian is zero.  Nevertheless
\[
 C_P(\mu_n)\geq n^2+1.                               \tag{53.28}
\]
\end{counterexample}

\begin{proof}
The law is the product of the displayed symmetric Gaussian mixture in
\(z\) with \(\gamma(dr)\).  Apply any Poincar\'e inequality to the
test function \(f(z,r)=z\).  Its Dirichlet energy is one, while its
variance is \(n^2+1\), proving (53.28).
\end{proof}

Thus finite soft dimension does not make the constant uniform.  The
soft free-energy margin (53.19), or another direct uniform soft
Poincar\'e theorem, is indispensable.

\subsection{Physical loading order}

Corollary~\ref{cor:v16v53-interpolated} turns the V52 Poincar\'e card
into four concrete rows:
\begin{enumerate}
\item a positive conditional law for the regular metric and hard
      Standard Model variables;
\item a cutoff-, shell-, and interpolation-uniform conditional
      Poincar\'e constant \(C_h\), assembled from the regular gap and
      the already archived gauge--Higgs--fermion conditional estimates;
\item the expected soft Hessian floor \(a\) and the mixed derivative
      bound \(b\), including the Einstein--GHY shell;
\item the strict numerical margin \(a-C_hb^2>0\).
\end{enumerate}
The fermion determinant phase and the gravitational conformal direction
must be resolved before row 1 can be asserted for a physical Euclidean
law.  Conditional Poincar\'e estimates for selected positive sectors do
not by themselves provide positivity of the full measure.

\begin{assessment}[What V53 closes]
\label{ass:v16v53-status}
V53 proves a two-scale Poincar\'e theorem that permits a finite soft
sector and a nonconvex hard conditional law.  It derives the soft
marginal constant from the exact effective Hessian and retains the
negative conditional covariance through the strict Schur margin
\(a-C_hb^2\).  Under the uniform interpolated version of that margin,
the Poincar\'e hypothesis used by V52 follows with the explicit constant
(53.24).

The physical action has not supplied the Hessian floor, mixed bound, or
positive full conditional law.  A regular-sector gap alone is
insufficient, as the separated-well example proves.
\end{assessment}

\begin{decision}[V54 Einstein--GHY shell Hessian ledger]
\label{dec:v16v53-next}
Differentiate the paired curvature shell a second time in the augmented
energy coordinate.  Separate the intrinsic second variation of the
Einstein--GHY functional from the first variation applied to the second
derivative of the metric reconstruction.  Use collar traces to state
the exact Sobolev orders and shell arrays required to bound
\(\kappa_{R,q}\) and \(b_{R,q}\) in (53.26), without claiming a sign
for the Euclidean Einstein Hessian.
\end{decision}

\begin{firewall}[Claim boundary after V53]
\label{fire:v16v53}
The two-scale theorem is conditional on a positive differentiable law
and uniform Poincar\'e and Hessian inputs.  It does not prove positivity
of the gravitational measure, convexity of the Einstein action, the
physical numerical margin, a limiting Einstein equation, or the full
Standard Model--Einstein continuum.
\end{firewall}

\section*{Version V53 append-only record}
\addcontentsline{toc}{section}{Version V53 append-only record}
The complete V52 source was retained before this continuation.  It had
(531032) bytes, (15885) lines, and SHA--256
\[
 \texttt{44EBE12BC379528DAC254D001AAF8E9A52B967E4C393026591EF1BA946B58AFF}.
\]
V53 proved the soft-marginal/conditional Poincar\'e theorem, derived its
soft constant from the exact effective Hessian with covariance loss,
and formulated the uniform interpolated Schur margin needed by V52.


\section{V54: Second variation of the reconstructed Einstein--GHY shell}
\label{sec:v16v54}

V53 reduces the interpolated Poincar\'e problem to the expected soft
Hessian floor and the mixed soft--hard Hessian norm.  The curvature
shell contributes to both.  An archive search found the exact
Einstein--GHY first variation imported in V47, but no completed
second-variation owner theorem.  This section supplies the missing
chain rule and a quantitative shell ledger.  It does not assign a
favorable sign to the Euclidean Einstein Hessian.

\subsection{Two sources of a reconstructed Hessian}

Let \(x=(z,r)\) denote the augmented soft and remaining coordinates,
let \(g_n(x)\) be a twice differentiable metric reconstruction, and let
\[
 V^R_{n,q}(x)
 =\delta c_{R,n,q}
 S_{\rm full}(g_n(x);\Omega_{n,q})                  \tag{54.1}
\]
be the paired Einstein--GHY--joint shell of V47.  For soft directions
\(a,b\) and a remaining direction \(u\), write
\[
\begin{array}{lll}
 h_{n,a}=D_zg_n[a],&
 p_{n,u}=D_rg_n[u],\\
 k_{n,a,b}=D_z^2g_n[a,b],&
 m_{n,a,u}=D_rD_zg_n[a,u].
\end{array}                                           \tag{54.2}
\]

\begin{theorem}[Exact reconstructed-action Hessian]
\label{thm:v16v54-chain}
Whenever \(S_{\rm full}\) and \(g_n\) are twice Fr\'echet
differentiable on the stated shell chart,
\[
\begin{split}
 D_z^2V^R_{n,q}[a,b]
 =\delta c_{R,n,q}\bigl{&
 D^2S_{\rm full}(g_n)[h_{n,a},h_{n,b}]\\
 &+DS_{\rm full}(g_n)[k_{n,a,b}]\bigr},             \tag{54.3}\
 D_rD_zV^R_{n,q}[a,u]
 =\delta c_{R,n,q}\bigl{&
 D^2S_{\rm full}(g_n)[h_{n,a},p_{n,u}]\\
 &+DS_{\rm full}(g_n)[m_{n,a,u}]\bigr}.             \tag{54.4}
\end{split}
\]
\end{theorem}

\begin{proof}
The first derivative is
\[
 D_zV^R_{n,q}[a]
 =\delta c_{R,n,q}DS_{\rm full}(g_n)[D_zg_n[a]].     \tag{54.5}
\]
Differentiate (54.5) in direction \(b\).  The derivative of the
first slot is the intrinsic bilinear form \(D^2S_{\rm full}\); the
derivative of \(D_zg_n[a]\) is \(D_z^2g_n[a,b]\).  This gives (54.3).
Differentiating in an (r)-direction gives (54.4).
\end{proof}

The second terms in (54.3)--(54.4) vanish for an affine metric
reconstruction.  They cannot be omitted for exponential, determinant-
normalized, or positivity-preserving metric charts.

\subsection{A paired second-variation coefficient card}

The local second variation of the Einstein--GHY action contains bulk
terms of differential order at most one in each variation after the
standard integration by parts, smooth-face terms containing the traces
of a variation and its first derivative, and algebraic joint traces.
We record the exact analytic bound needed from that local formula.

For a shell variation \(h\), set
\[
\begin{split}
 \|h\|_{\mathbb X_{n,q}}^2
 :={}&
 \|h\|_{L^2(\Omega_{n,q})}^2
 +\|\nabla h\|_{L^2(\Omega_{n,q})}^2\\
 &+\|h\|_{L^2(\partial\Omega_{n,q})}^2
 +\|\nabla h\|_{L^2(\partial\Omega_{n,q})}^2
 +\|h\|_{L^2(\mathcal J_{n,q})}^2.                  \tag{54.6}
\end{split}
\]

\begin{hypothesis}[Uniform paired Hessian coefficients]
\label{hyp:v16v54-second-card}
On the geometric set under consideration, there is a cutoff- and
shell-independent \(K_2<\infty\) such that
\[
 |D^2S_{\rm full}(g)[h,k]|
 \leq K_2\|h\|_{\mathbb X_{n,q}}
          \|k\|_{\mathbb X_{n,q}}.                  \tag{54.7}
\]
All face and joint orientations and incidence multiplicities are those
of the paired action in (47.4)--(47.5).
\end{hypothesis}

This is a coefficient hypothesis, not a sign hypothesis.  Uniform
ellipticity, bounded curvature and second fundamental form, and bounded
first derivatives of the coefficients are sufficient local geometric
inputs for (54.7).  A compatible fixed-boundary polarization may remove
some traces; V54 retains them because a remote shell mode need not obey
that boundary condition.

\subsection{An (L^2) collar trace with explicit width}

Let \(\mathcal C_\ell\simeq[0,\ell]\times\Sigma\) be a normal collar
whose volume density is uniformly comparable to the product density,
with comparison constant \(C_{\rm col}\).

\begin{lemma}[Squared collar trace]
\label{lem:v16v54-L2-trace}
For \(u\in W^{1,2}(\mathcal C_\ell)\),
\[
 \|u|_\Sigma\|_2^2
 \leq C_{\rm col}
 \left(
 2\ell^{-1}\|u\|_{L^2(\mathcal C_\ell)}^2
 +\ell\|\nabla_nu\|_{L^2(\mathcal C_\ell)}^2
 \right).                                           \tag{54.8}
\]
The same estimate applied to \(\nabla u\) gives
\[
 \|\nabla u|_\Sigma\|_2^2
 \leq C_{\rm col}
 \left(
 2\ell^{-1}\|\nabla u\|_2^2
 +\ell\|\nabla^2u\|_2^2
 \right).                                           \tag{54.9}
\]
\end{lemma}

\begin{proof}
On a normal line, the fundamental theorem of calculus gives
\[
 |u(0)|^2
 \leq |u(t)|^2+2\int_0^t|u(s)|\,|\partial_nu(s)|\,ds.         \tag{54.10}
\]
Average in \(t\in[0,\ell]\), integrate over \(\Sigma\), and use the
collar density comparison.  The result before Young's inequality is
\[
 \|u|_\Sigma\|_2^2
 \leq C_{\rm col}left(
 \ell^{-1}\|u\|_2^2+2\|u\|_2\|\nabla_nu\|_2\right).          \tag{54.11}
\]
Use
\(2ab\leq\ell^{-1}a^2+\ell b^2\) to obtain (54.8).
Apply the same argument componentwise to \(\nabla u\) for (54.9).
\end{proof}

\begin{corollary}[Codimension-two squared trace]
\label{cor:v16v54-joint-trace}
For two comparable transverse collars of width \(\ell\),
\[
 \|u|_{\mathcal J}\|_2^2
 \leq C_{\rm col}^2
 \left(
 4\ell^{-2}\|u\|_2^2
 +4\|\nabla u\|_2^2
 +\ell^2\|\nabla^2u\|_2^2
 \right),                                           \tag{54.12}
\]
after enlarging the harmless middle coefficient if the two collar
Jacobians differ within the stated comparison class.
\end{corollary}

\begin{proof}
Apply (54.8) in the first transverse direction, and apply it again to
the resulting face field in the second direction.  Expanding the two
factors gives an \(\ell^{-2}\) zeroth-order term, first-derivative terms
with no net power of \(\ell\), and an \(\ell^2\) mixed
second-derivative term.  Bounding every first or mixed derivative by
the full covariant derivative gives (54.12).
\end{proof}

Equations (54.8)--(54.12) show that a safe unrestricted paired
second-variation norm is loaded by bulk \(W^{2,2}\) on a piecewise
shell.  On a smooth face, the first-derivative face term in (54.6) still
uses (54.9), so the safe order remains two.  This is one derivative
above the (W^{1,1}) first-variation loader available after a smooth
Einstein--GHY pairing in V48.

Define the explicit bulk collar norm
\[
\begin{split}
 \|h\|_{\mathbb W_{n,q}}^2:={}&
 \|h\|_2^2+\|\nabla h\|_2^2\\
 &+2\ell_{n,q}^{-1}\|h\|_2^2
   +\ell_{n,q}\|\nabla h\|_2^2\\
 &+2\ell_{n,q}^{-1}\|\nabla h\|_2^2
   +\ell_{n,q}\|\nabla^2h\|_2^2\\
 &+4\ell_{n,q}^{-2}\|h\|_2^2
   +4\|\nabla h\|_2^2
   +\ell_{n,q}^2\|\nabla^2h\|_2^2,                 \tag{54.13}
\end{split}
\]
where all bulk norms are over the shell collar.  Lemma
\ref{lem:v16v54-L2-trace} and Corollary
\ref{cor:v16v54-joint-trace} give
\[
 \|h\|_{\mathbb X_{n,q}}
 \leq C_X\|h\|_{\mathbb W_{n,q}}                  \tag{54.14}
\]
with a uniform (C_X) after including the bounded bulk shell outside
the collar.

\subsection{First variation applied to reconstruction curvature}

For a second reconstruction response \(k\), define the paired
first-variation loader
\[
 \|k\|_{\mathbb L_{n,q}}
 :=K_G\|k\|_{L^1(\Omega_{n,q})}
  +K_\Pi\|k^\top\|_{L^1(\partial\Omega_{n,q})}
  +K_{\mathcal J}\|k\|_{L^1(\mathcal J_{n,q})}.     \tag{54.15}
\]
Then the imported first variation gives
\[
 |DS_{\rm full}(g)[k]|
 \leq|c_R|\|k\|_{\mathbb L_{n,q}}.                 \tag{54.16}
\]
The V48 (L^1) collar theorem loads (54.15) from bulk
\(W^{1,1}\) for a smooth paired face and from \(W^{2,1}\) when a
joint survives.  Thus nonlinear reconstruction curvature costs the
same spatial derivative order as a first variation, but it requires a
second derivative with respect to the source coordinates.

\subsection{Soft and mixed shell-Hessian arrays}

Set
\[
\begin{split}
 X^z_{n,q}
 &:=\sup_{x,\|a\|=1}
 \|h_{n,a}(x)\|_{\mathbb W_{n,q}},\\
 X^r_{n,q}
 &:=\sup_{x,\|u\|=1}
 \|p_{n,u}(x)\|_{\mathbb W_{n,q}},\\
 Y^{zz}_{n,q}
 &:=\sup_{x,\|a\|=\|b\|=1}
 \|k_{n,a,b}(x)\|_{\mathbb L_{n,q}},\\
 Y^{zr}_{n,q}
 &:=\sup_{x,\|a\|=\|u\|=1}
 \|m_{n,a,u}(x)\|_{\mathbb L_{n,q}}.                \tag{54.17}
\end{split}
\]
The suprema may later be replaced by interpolated moments, but the
deterministic form is the correct input for the uniform Hessian margin
of V53.

\begin{theorem}[Curvature-shell Hessian ledger]
\label{thm:v16v54-ledger}
Under Hypothesis~\ref{hyp:v16v54-second-card}, the curvature-shell
contributions in (53.26) may be bounded by
\[
 \boxed{
 \kappa_{R,q}
 \leq|\delta c_{R,n,q}|
 \left(K_2C_X^2(X^z_{n,q})^2+|c_R|Y^{zz}_{n,q}\right),}       \tag{54.18}
\]
and
\[
 \boxed{
 b_{R,q}
 \leq|\delta c_{R,n,q}|
 \left(K_2C_X^2X^z_{n,q}X^r_{n,q}
       +|c_R|Y^{zr}_{n,q}\right).}                  \tag{54.19}
\]
These bounds hold uniformly for the interpolation
\(H_{n,q,0}+tV^R_{n,q}\), because \(0\leq t\leq1\).
\end{theorem}

\begin{proof}
Take unit soft directions in (54.3).  Apply (54.7) to its intrinsic
term, use (54.14), and apply (54.16) to its reconstruction-curvature
term.  Taking the absolute value bounds the negative part of the soft
Hessian and gives (54.18).  Apply the same argument to (54.4), with one
soft and one remaining unit direction, to obtain (54.19).  Multiplying
by (tin[0,1]) cannot enlarge either bound.
\end{proof}

\begin{corollary}[Affine reconstruction simplification]
\label{cor:v16v54-affine}
If \(g_n(z,r)\) is affine in (z) and has no mixed source curvature,
then
\[
 Y^{zz}_{n,q}=Y^{zr}_{n,q}=0,                        \tag{54.20}
\]
and only the intrinsic Einstein--GHY bilinear form remains in
(54.18)--(54.19).
\end{corollary}

\begin{proof}
Affineness gives (D_z^2g_n=0), and absence of mixed source curvature
gives (D_rD_zg_n=0).  Use the definitions in (54.17).
\end{proof}

\subsection{Cofinal Hessian payment}

For insertion into V53, a sufficient uniform numerical condition is
\[
\begin{split}
 a_{\rm aug}-\kappa_{\rm SM}
 &-\sup_{n,q}|\delta c_{R,n,q}|
 \left(K_2C_X^2(X^z_{n,q})^2+|c_R|Y^{zz}_{n,q}\right)\\
 &-C_h\left[
 b_{\rm base}
 +\sup_{n,q}|\delta c_{R,n,q}|
 \left(K_2C_X^2X^z_{n,q}X^r_{n,q}
       +|c_R|Y^{zr}_{n,q}\right)
 \right]^2
 \geq\lambda_*>0.                                   \tag{54.21}
\end{split}
\]
This is deliberately a supremum condition because it supplies a
Poincar\'e constant uniform over every single shell interpolation.
After that constant is obtained, the owner debit itself may use the RMS
and trace-class summation of V52.

The width losses are visible in (54.13).  A collapsing physical collar
can make (54.21) fail even when unweighted bulk \(H^2\) norms remain
bounded.  EGHY pairing removes the normal derivative from the first
variation but does not remove all first-derivative face terms from the
intrinsic second variation.

\begin{assessment}[What V54 closes]
\label{ass:v16v54-status}
V54 proves the exact two-term chain rule for the soft and mixed Hessians
of a reconstructed curvature shell.  It supplies (L^2) face and joint
collar traces, identifies a safe bulk (W^{2,2}) order for the paired
intrinsic Hessian, loads nonlinear reconstruction curvature through the
V47--V48 first-variation norms, and gives the explicit contributions
(54.18)--(54.19) to the V53 Schur margin.

The local coefficient card (54.7), the weighted first- and second-source
response norms, and the positive numerical margin (54.21) have not been
proved for the physical regulator.  No sign of the Euclidean Einstein
Hessian is asserted.
\end{assessment}

\begin{decision}[V55 companion audit and moment-level Hessian reserve]
\label{dec:v16v54-next}
Perform the compulsory companion audit.  Then determine whether the
uniform Hessian supremum in (54.21) can be replaced by a high-probability
or exponential-moment reserve without losing the Poincar\'e inequality
needed in V52.  Audit the archives for Lyapunov--Poincar\'e or defective
Poincar\'e repair before deriving a new result, and continue
substantively in V56 after the audit.
\end{decision}

\begin{firewall}[Claim boundary after V54]
\label{fire:v16v54}
The Hessian ledger is conditional on uniform geometric coefficients and
weighted reconstruction norms.  It does not prove the physical Hessian
margin, positivity of the gravitational measure, a limiting Einstein
equation, or the full Standard Model--Einstein continuum.
\end{firewall}

\section*{Version V54 append-only record}
\addcontentsline{toc}{section}{Version V54 append-only record}
The complete V53 source was retained before this continuation.  It had
(544386) bytes, (16267) lines, and SHA--256
\[
 \texttt{84160EA8FAA35C64BA824BB5C641FABDF703AD2CFE41BE40069B794FC499F337}.
\]
V54 differentiated the reconstructed paired curvature action twice,
proved explicit (L^2) collar traces, and converted intrinsic and
reconstruction Hessians into the soft and mixed arrays required by the
V53 Poincar\'e margin.


\section{V55: Companion audit, de-duplication correction, and the failure of probabilistic convexity}
\label{sec:v16v55}

This compulsory checkpoint audits the companion programmes and the
relevant SM archive before replacing the uniform Hessian margin of V54.
The archive audit found an overlap in V52 which must be stated
explicitly.  It also found a completed Lyapunov--local-Poincar\'e theorem
that should be imported rather than rederived.

\subsection{Mandatory companion audit}

The newest Yang--Mills companion is
\texttt{Renewal\_Yang\_Mills\_Mass\_Gap\_Research\_Notes\_V54.tex}.
It has (467684) bytes, (12747) lines, and SHA--256
\[
 \texttt{20F0DA83E09A4609C73A76F9F19856921AD4802715506933C3ADE489D9FC9569}.
                                                               \tag{55.1}
\]
YM V54 classifies the exact intersections in a proposed mixed
third-coordinate layer.  It saves four existing charts and gives exact
negative quadratic witnesses against both a full-radius certificate and
one overlap-compatible anisotropic repair at four centres.  The witnesses
disprove those sufficient certificate blocks, not the underlying Wilson
pencil.  This is a useful upstream obstruction analysis but supplies no
Poincar\'e, metric-reconstruction, or Einstein--GHY Hessian estimate.

The newest Navier--Stokes companion remains
\texttt{Renewal\_NS\_Stability\_Research\_Notes\_V26.tex}.  It has
(278283) bytes, (5319) lines, and SHA--256
\[
 \texttt{82C887F8C2C69E4F28FAD7C3DC8DE5B9099CB5A4BF431EBA1FFF3E91935978AD}.
                                                               \tag{55.2}
\]
Its exact rank-one Oseen-kernel derivative norms do not address the
interpolated metric law.  Neither companion changes the V54 acquisition
order.

\subsection{Correction to the de-duplication ledger}

The frozen file
\texttt{Renewal\_SM\_Einstein\_Continuum\_Research\_Notes\_V100\_f15.tex}
contains, in its fifteenth-cycle V47, the exact identities
\[
 D(\nu_1\Vert\nu_0)
 =\int_0^1t\,\operatorname{Var}_{\nu_t}(\Delta)\,dt,
 \qquad
 D(\nu_0\Vert\nu_1)
 =\int_0^1(1-t)\operatorname{Var}_{\nu_t}(\Delta)\,dt.         \tag{55.3}
\]
Its V47 also gives the Poincar\'e-to-entropy estimate, and its V48
gives a conditional head--tail entropy localization and optimized collar
growth.  Therefore Theorem~\ref{thm:v16v52-entropy} and equation (52.8)
are a rederivation of (55.3), not a new result.  This overlap was missed
in the V52 precheck and is corrected here.  From V55 onward, (55.3) is
treated as an imported fifteenth-cycle theorem.

The new content retained from V52 is narrower and specific to the current
programme:
\begin{enumerate}
\item the countable curvature-shell formulation (52.10)--(52.12) in the
      active V35--V36 owner ledger;
\item the chain-rule loader (52.18)--(52.21) from the V51
      covariance-weighted metric response;
\item the orthogonal shell trace-ideal compiler (52.22)--(52.27);
\item the RMS Einstein--GHY functional (52.29) and the proof that the V32
      \(L^p\) density hypothesis does not imply a Poincar\'e inequality.
\end{enumerate}
The two-scale soft-marginal theorem of V53 and the reconstructed
Einstein--GHY Hessian ledger of V54 were not found in the archive search.

\subsection{Imported Lyapunov repair}

Fifteenth-cycle V73 proves the following implication.  Let
\(\nu\propto e^{-\mathcal A}dx\), let
\(\mathscr L=\Delta-\nabla\mathcal A\cdot\nabla\), and suppose a set
\(K\) of mass \(a=\nu(K)>0\) has local Poincar\'e constant \(C_K\).
If a function \(W\geq1\) satisfies
\[
 -{\mathscr LW\over W}\geq c-b\mathbf1_K           \tag{55.4}
\]
and
\[
 \gamma:=c-b{1-a\over a}>0,                          \tag{55.5}
\]
then
\[
 C_P(\nu)\leq{1+bC_K\over\gamma}.                   \tag{55.6}
\]
The archived proof uses the exact integration estimate
\[
 \int\left(-{\mathscr LW\over W}\right)f^2d\nu
 \leq\int|\nabla f|^2d\nu                            \tag{55.7}
\]
and the local Poincar\'e inequality.  V55 imports (55.4)--(55.7) with
their stated smoothness and closure hypotheses; it does not repeat that
proof.

This theorem shows what can replace the global Hessian supremum in
(54.21): a global Lyapunov return, a uniformly connected local core, and
the quantitative mass margin (55.5).  A moment bound or a high-probability
convexity statement alone is not such a replacement.

\subsection{High-probability strong convexity is insufficient}

Let \(\phi(x)=(2\pi)^{-1/2}e^{-x^2/2}\), and for \(A>1\) define the
symmetric Gaussian mixture
\[
 p_A(x)={1\over2}\phi(x-A)+{1\over2}\phi(x+A)
 ={e^{-(x^2+A^2)/2}\cosh(Ax)\over\sqrt{2\pi}}.      \tag{55.8}
\]
Its potential, up to an additive constant, is
\[
 U_A(x)={x^2\over2}-\log\cosh(Ax),qquad
 U_A''(x)=1-A^2\operatorname{sech}^2(Ax).             \tag{55.9}
\]

\begin{theorem}[Convex with high probability but no uniform gap]
\label{thm:v16v55-mixture}
Let
\[
 w_A={\operatorname{arcosh}(\sqrt2A)\over A},
 \qquad B_A=[-w_A,w_A].                               \tag{55.10}
\]
Then
\[
 U_A''(x)\geq{1\over2}qquad(x\notin B_A),           \tag{55.11}
\]
and
\[
 p_A(B_A)\longrightarrow0.                           \tag{55.12}
\]
Nevertheless there is a universal \(c>0\) such that, for all
sufficiently large \(A\),
\[
 C_P(p_A(x)dx)
 \geq c\,{\log A\over A^2}e^{A^2/2}.                \tag{55.13}
\]
\end{theorem}

\begin{proof}
Outside (B_A),
\(\cosh^2(Ax)\geq2A^2\), so (55.9) gives (55.11).
On (B_A), monotonicity of \(\cosh\) and (55.10) give
\[
 p_A(x)
 \leq C A e^{-A^2/2},                                \tag{55.14}
\]
where the harmless factor (e^{-x^2/2}\leq1) was discarded.  Hence
\[
 p_A(B_A)\leq C A w_Ae^{-A^2/2}leq C(1+\log A)e^{-A^2/2},     \tag{55.15}
\]
which proves (55.12).

Let (f_A) be odd, equal to (-1) on
\(( -\infty,-w_A]), equal to (x/w_A) on (B_A), and equal to (1)
on \([w_A,\infty)\).  Symmetry gives mean zero, and (55.12) gives
\[
 \operatorname{Var}_{p_A}(f_A)
 \geq1-p_A(B_A)\geq{1\over2}                         \tag{55.16}
\]
for large (A).  Its derivative is (w_A^{-1}\) on (B_A) and zero
elsewhere, so (55.15) yields
\[
 \int|f_A'|^2p_A,dx
 \leq {C A e^{-A^2/2}\over w_A}
 \leq C{A^2\over\log A}e^{-A^2/2}.                  \tag{55.17}
\]
The ratio of (55.16) to (55.17) proves (55.13).
\end{proof}

Thus a Hessian floor outside a set of vanishing probability can coexist
with an exponentially deteriorating Poincar\'e constant.  The low-mass
neck between two wells controls the gap.  Exponential moments of the
negative Hessian on their own do not rule out the same bottleneck.

\subsection{Decision after the audit}

The V54 supremum margin has two legitimate replacements:
\begin{enumerate}
\item prove the uniform two-scale Schur margin (53.23); or
\item for the finite soft marginal, construct a Lyapunov function and a
      uniform local core satisfying (55.4)--(55.6), then use the V53
      conditional/marginal theorem to lift that soft gap to the full law.
\end{enumerate}
Merely placing the negative curvature in a small-probability event is
invalid by Theorem~\ref{thm:v16v55-mixture}.

\begin{assessment}[What V55 establishes]
\label{ass:v16v55-status}
V55 completes the companion audit, corrects the V52 de-duplication
record, imports the archived Lyapunov--local-Poincar\'e theorem, and
proves that high-probability strong convexity cannot replace a uniform
spectral-gap mechanism.  The appropriate nonconvex route must control
global return and local connectivity quantitatively.

No physical Lyapunov function, uniform local soft-core inequality, or
mass margin has yet been constructed for the interpolated
SM--Einstein metric law.
\end{assessment}

\begin{decision}[V56 radial soft drift and local-core compiler]
\label{dec:v16v55-next}
Continue after the compulsory audit.  For the finite augmented soft
coordinate, derive a Lyapunov function from a uniform radial drift
bound on the effective free energy.  Obtain explicit constants in
(55.4)--(55.6), load the local core from a bounded oscillation of the
effective potential on a ball, and then insert the resulting soft
Poincar\'e constant into the V53 two-scale theorem.
\end{decision}

\begin{firewall}[Claim boundary after V55]
\label{fire:v16v55}
The audit and counterexample select a valid analytic route; they do not
populate its physical Lyapunov, local-core, or positivity hypotheses.
They do not prove a limiting Einstein equation or the full Standard
Model--Einstein continuum.
\end{firewall}

\section*{Version V55 append-only record}
\addcontentsline{toc}{section}{Version V55 append-only record}
The complete V54 source was retained before this continuation.  It had
(557315) bytes, (16640) lines, and SHA--256
\[
 \texttt{034762339F3304A51AFBFC4D576D4966ECA2D2794FE52978DE716E5C873A8ED5}.
\]
V55 audited YM V54 and NS V26, corrected the entropy-identity archive
overlap, imported the Lyapunov/local-Poincar\'e repair, and proved an
explicit two-well obstruction to high-probability convexity arguments.


\section{V56: Radial soft drift and an explicit local-core Poincar\'e compiler}
\label{sec:v16v56}

V55 rules out a bare high-probability Hessian argument and imports the
Lyapunov--local-Poincar\'e repair.  This section constructs all of its
constants from a radial drift of the finite augmented soft free energy.
It then lifts the resulting soft gap through the V53 conditional theorem.
This route permits nonconvexity inside a fixed connected core and does
not require a favorable sign for the Einstein Hessian there.

\subsection{The effective soft generator}

Let \(m<\infty\) be the fixed augmented soft rank and let
\[
 d\alpha(z)=Z^{-1}e^{-\Psi(z)}\,dz,qquad z\in\mathbb R^m,     \tag{56.1}
\]
where \(\Psi\in C^2\), the density is positive, and all integrations
below are justified by truncation and closure.  Write
\[
 \mathscr L_\Psi=\Delta_z-\nabla\Psi\cdot\nabla_z.    \tag{56.2}
\]

\begin{hypothesis}[Uniform radial return]
\label{hyp:v16v56-radial}
There are \(\rho>0\) and \(C\geq0\) such that
\[
 z\cdot\nabla\Psi(z)
 \geq\rho|z|^2-C                                    \tag{56.3}
\]
for every \(z\in\mathbb R^m\).
\end{hypothesis}

\begin{lemma}[Gaussian Lyapunov function]
\label{lem:v16v56-Lyapunov}
For \(0<\eta<\rho\), put
\[
 W_\eta(z)=e^{\eta|z|^2/2},qquad
 \lambda=\eta(\rho-\eta),qquad
 d_0=\eta(C+m).                                      \tag{56.4}
\]
Then
\[
 -{\mathscr L_\Psi W_\eta\over W_\eta}
 \geq\lambda|z|^2-d_0.                               \tag{56.5}
\]
\end{lemma}

\begin{proof}
Direct differentiation gives
\[
 {\Delta W_\eta\over W_\eta}
 =\eta m+\eta^2|z|^2,qquad
 {\nabla W_\eta\over W_\eta}=\eta z.               \tag{56.6}
\]
Therefore
\[
 -{\mathscr L_\Psi W_\eta\over W_\eta}
 =\eta z\cdot\nabla\Psi-eta m-eta^2|z|^2.          \tag{56.7}
\]
Insert (56.3) and collect terms.
\end{proof}

\subsection{A fixed ball supplies all Lyapunov constants}

Choose \(R>0\), let \(K_R=\{|z|\leq R\}\), and define
\[
 b_R:=\lambda R^2,qquad
 c_R:=\lambda R^2-d_0.                               \tag{56.8}
\]

\begin{proposition}[Explicit drift and core mass]
\label{prop:v16v56-mass}
If \(R^2>d_0/\lambda\), then
\[
 -{\mathscr L_\Psi W_\eta\over W_\eta}
 \geq c_R-b_R\mathbf1_{K_R}.                         \tag{56.9}
\]
Writing \(a_R=\alpha(K_R)\), one also has
\[
 a_R\geq{c_R\over b_R}
 =1-{d_0\over\lambda R^2}.                           \tag{56.10}
\]
If
\[
 R^2\geq{4d_0\over\lambda},                         \tag{56.11}
\]
then \(a_R\geq3/4\) and the V55 Lyapunov margin satisfies
\[
 \gamma_R:=c_R-b_R{1-a_R\over a_R}
 \geq{5\over12}b_R>0.                               \tag{56.12}
\]
\end{proposition}

\begin{proof}
Outside (K_R), (56.5) is at least (c_R).  Inside (K_R), its
right side is at least (-d_0=c_R-b_R), proving (56.9).
Use the imported integration estimate (55.7) with the constant test
function.  Its Dirichlet energy is zero, so (56.9) gives
\[
 0\geq\int -{\mathscr L_\Psi W_\eta\over W_\eta},d\alpha
 \geq c_R-b_Ra_R,                                   \tag{56.13}
\]
which is (56.10).  Under (56.11), the ratio
\(d_0/(\lambda R^2)\) is at most (1/4), hence
\(a_R\geq3/4\), \(c_R\geq3b_R/4\), and
\((1-a_R)/a_R\leq1/3\).  Therefore
\[
 \gamma_R\geq{3b_R\over4}-{b_R\over3}
 ={5b_R\over12}.                                    \tag{56.14}
\]
\end{proof}

The mass estimate (56.10) is a consequence of the global return, not a
separate concentration assumption.  The separated-well family of V55
cannot satisfy (56.3) with cutoff-independent \(C\); its central barrier
forces that constant to grow with the well separation.

\subsection{Local Poincar\'e from a potential oscillation}

Let
\[
 \Omega_R:=\operatorname*{ess\,sup}_{K_R}\Psi
            -\operatorname*{ess\,inf}_{K_R}\Psi.     \tag{56.15}
\]

\begin{lemma}[Connected-ball local constant]
\label{lem:v16v56-local}
If \(\Omega_R<\infty\), the normalized restriction
\(\alpha_R=\alpha(\,cdot\mid K_R)\) satisfies
\[
 C_P(\alpha_R)
 \leq {4R^2\over\pi^2}e^{\Omega_R}.                 \tag{56.16}
\]
The same constant is valid for the unnormalized local inequality in
(55.4)--(55.6).
\end{lemma}

\begin{proof}
The uniform probability measure on the convex ball (K_R) has
Poincar\'e constant at most
\(\operatorname{diam}(K_R)^2/\pi^2=4R^2/\pi^2\).
On that ball, the density ratio of \(\alpha_R\) to the uniform law has
oscillation ratio at most \(e^{\Omega_R}\).  The bounded-perturbation
comparison gives (56.16).  Multiplication of both sides of the
normalized inequality by \(a_R\) gives the local inequality with the
unnormalized restriction.
\end{proof}

\begin{theorem}[Radial-drift soft Poincar\'e compiler]
\label{thm:v16v56-soft}
Assume Hypothesis~\ref{hyp:v16v56-radial}.  Choose
\(0<\eta<\rho\) and (R) satisfying (56.11), and suppose
\(\Omega_R<\infty\).  Then
\[
 \boxed{
 C_P(\alpha)
 \leq {12\over5b_R}
 +{48R^2\over5\pi^2}e^{\Omega_R}.}                  \tag{56.17}
\]
\end{theorem}

\begin{proof}
Proposition~\ref{prop:v16v56-mass} supplies (55.4) with
\(c=c_R\), \(b=b_R\), and margin
\(\gamma_R\geq5b_R/12\).  Lemma~\ref{lem:v16v56-local} supplies
\(C_K\leq4R^2e^{\Omega_R}/\pi^2\).  Insert both in the imported
bound (55.6):
\[
 {1+b_RC_K\over\gamma_R}
 \leq{12\over5b_R}+{12\over5}C_K,                   \tag{56.18}
\]
which is (56.17).
\end{proof}

The estimate is not optimized in \(\eta\) or \(R\).  Its purpose is
to expose a finite list of cutoff-uniform inputs.  The core radius is
fixed once \(m,\rho,C\) are fixed, and only the additive-free
oscillation \(\Omega_R\) of the effective soft potential is needed on
that ball.

\subsection{Uniform interpolated family}

Let \(\alpha_{n,q,t}\propto e^{-\Psi_{n,q,t}}dz\) be the soft
marginals of the V52 shell interpolation.

\begin{corollary}[Uniform soft gap without global convexity]
\label{cor:v16v56-uniform-soft}
Suppose the soft rank \(m\) is fixed and the same \(\rho>0\), \(C\),
and local oscillation bound \(\Omega\) satisfy
\[
 z\cdot\nabla\Psi_{n,q,t}(z)
 \geq\rho|z|^2-C,qquad
 \operatorname{osc}_{K_R}\Psi_{n,q,t}\leq\Omega    \tag{56.19}
\]
for every \(n,q,t\), where \(R\) is any common choice satisfying
(56.11).  Then
\[
 \sup_{n,q,t}C_P(\alpha_{n,q,t})
 \leq C_s^{\rm rad}
 :={12\over5b_R}
   +{48R^2\over5\pi^2}e^\Omega.                     \tag{56.20}
\]
\end{corollary}

\begin{proof}
Apply Theorem~\ref{thm:v16v56-soft} to each member of the family with
the same \(\eta,R,\lambda,d_0,b_R\), and \(\Omega\).
\end{proof}

\begin{corollary}[Full two-scale Poincar\'e constant]
\label{cor:v16v56-full}
In addition to (56.19), suppose every regular/hard conditional law has
Poincar\'e constant at most \(C_h\) and its mixed soft--remaining
Hessian has norm at most \(b_{\rm mix}\).  Then the full interpolated
law satisfies (52.16) with
\[
 C_P\leq
 \inf_{\delta>0}\max\left\{
 C_s^{\rm rad}(1+\delta),
 C_h+C_s^{\rm rad}(1+\delta^{-1})C_h^2b_{\rm mix}^2
 \right\}.                                          \tag{56.21}
\]
\end{corollary}

\begin{proof}
Corollary~\ref{cor:v16v56-uniform-soft} supplies the marginal constant
in Hypothesis~\ref{hyp:v16v53-two-scale}.  Apply
Theorem~\ref{thm:v16v53-two-scale} uniformly.
\end{proof}

Unlike the Schur-margin route (53.23), (56.21) imposes no smallness
condition on \(b_{\rm mix}\); a large mixed derivative worsens the
constant but does not destroy it once the conditional and marginal gaps
are independently established.

\subsection{Loading radial return from the action}

The effective soft derivative has no covariance subtraction.  If
\[
 \Psi_{n,q,t}(z)
 =-\log\int e^{-H_{n,q,t}(z,r)}dr,                   \tag{56.22}
\]
then dominated differentiation gives
\[
 z\cdot\nabla\Psi_{n,q,t}(z)
 =\mathbb E_{\kappa_{n,q,t,z}}
   [z\cdot\nabla_zH_{n,q,t}(z,r)].                  \tag{56.23}
\]

\begin{proposition}[Form-small radial interaction]
\label{prop:v16v56-form-small}
Suppose
\[
 H_{n,q,t}(z,r)
 ={\rho_0\over2}|z|^2+I_{n,q,t}(z,r),               \tag{56.24}
\]
and for constants \(0\leq\varepsilon<\rho_0\), \(C_I\geq0\),
\[
 \mathbb E_{\kappa_{n,q,t,z}}
 [z\cdot\nabla_z I_{n,q,t}]
 \geq-\varepsilon|z|^2-C_I                          \tag{56.25}
\]
uniformly in \(n,q,t,z\).  Then (56.19) holds with
\[
 \rho=\rho_0-\varepsilon,qquad C=C_I.               \tag{56.26}
\]
\end{proposition}

\begin{proof}
Insert (56.24) in (56.23) and use (56.25).
\end{proof}

A pointwise gradient bound provides a convenient sufficient condition.
If an interaction component \(I_1\) satisfies
\[
 \|\nabla_zI_1(z,r)\|
 \leq L_1|z|+L_0,                                    \tag{56.27}
\]
then for every \(\tau>0\),
\[
 z\cdot\nabla_zI_1
 \geq-(L_1+\tau)|z|^2-{L_0^2\over4\tau}.            \tag{56.28}
\]
Thus the first-variation reconstruction estimates of V47--V49 can
populate a radial-return row without controlling the full curvature
Hessian, provided their growth in the soft energy coordinate is at most
linear with a sufficiently small slope.

For the curvature shell, (54.5) gives the exact radial insertion
\[
 z\cdot\nabla_zV^R_{n,q}
 =\delta c_{R,n,q}
 DS_{\rm full}(g_n)[D_zg_n[z]].                      \tag{56.29}
\]
The V47 owner norm and V48 collar traces can be applied to the radial
response \(D_zg_n[z]\).  To use Proposition
\ref{prop:v16v56-form-small}, the resulting coefficient of \(|z|^2\)
must be strictly below the augmented quadratic floor \(\rho_0\),
uniformly for every shell interpolation.  This is weaker than a
uniform positive Hessian, but it is still a physical form-smallness
condition.

\subsection{Combined route to the V52 owner tail}

The nonconvex averaged curvature route is now a precise implication:
\begin{enumerate}
\item prove the radial form bound (56.25) and the fixed-ball oscillation
      in (56.19);
\item prove a conditional regular/hard Poincar\'e constant and mixed
      derivative bound, giving (56.21);
\item prove the oriented covariance response and Sobolev
      Hilbert--Schmidt convergence of V51;
\item prove the RMS curvature array (52.29) and its shell product tail
      (52.24).
\end{enumerate}
Then Theorem~\ref{thm:v16v52-curvature-tail} gives a summable curvature
owner without a bounded soft-carrier diameter.

\begin{assessment}[What V56 closes]
\label{ass:v16v56-status}
V56 turns a uniform radial drift of the finite augmented soft free
energy into an explicit global Poincar\'e constant.  It derives the
core mass from the Lyapunov inequality itself, loads the connected local
core from a bounded potential oscillation, and lifts the soft constant
through the V53 two-scale theorem.  It also gives the exact radial
form-smallness condition by which the V47--V49 reconstruction estimates
can enter.

The physical action has not supplied the radial interaction bound,
fixed-ball free-energy oscillation, positive conditional law, or
conditional hard gap.  The conformal metric direction remains a
coercivity obstruction until a declared stabilizing term or contour
produces (56.25).
\end{assessment}

\begin{decision}[V57 radial reconstruction and counterterm balance]
\label{dec:v16v56-next}
Apply (56.29) to the individual metric vacuum and curvature
counterterms.  Derive their exact contributions to \(L_1,L_0\) in
(56.27), distinguish a coefficient that changes the quadratic radial
slope from an additive force, and formulate the stabilizer inequality
needed to dominate the conformal channel without deleting its Einstein
source.
\end{decision}

\begin{firewall}[Claim boundary after V56]
\label{fire:v16v56}
The radial compiler is conditional on uniform form-smallness and a
positive local density.  It does not construct the gravitational
stabilizer, prove the physical conditional gap, recover a limiting
Einstein equation, or complete the Standard Model--Einstein continuum.
\end{firewall}

\section*{Version V56 append-only record}
\addcontentsline{toc}{section}{Version V56 append-only record}
The complete V55 source was retained before this continuation.  It had
(566248) bytes, (16870) lines, and SHA--256
\[
 \texttt{612F965CB636D316771EEC9522A8F00E8C7C62F12B20F645C8DAE3AED4BE004A}.
\]
V56 constructed an explicit radial Lyapunov function, derived the
soft-core mass and Poincar\'e constant, lifted it through the V53
two-scale theorem, and isolated the radial form-smallness condition for
the reconstructed curvature action.


\section{V57: counterterm slope--force budgets and a source-preserving stabilizer corridor}
\label{sec:v16v57}

\subsection{Append-only source erratum for V51--V56}

The mathematical continuation in V51--V56 was installed through a
string-escaping path that removed several isolated backslashes.  The
installed prefix is immutable, so this subsection records the
authoritative readings rather than silently changing earlier bytes.
All subsequent references to the affected expressions mean the
following corrected forms:
\begin{align*}
(51.2)&:\quad \mathfrak a_n,
&
(51.3)&:\quad (\Gamma_nv,\,A_{n,\ell}\psi_n),
\\
(51.23)&:\quad
Ge_m=\frac1{m^2}e_m,\qquad T=I,
&
(52.1)&:\quad \cdots,\qquad \cdots,
\\
(53.2)&:\quad e^{-H(z,r)}\,dr,
&
(53.3)&:\quad
\int\!\left(f-\int f\,d\kappa_z\right)^2\,d\kappa_z,
\\
(53.4)&:\quad
\int\!\left(F-\int F\,d\alpha\right)^2\,d\alpha,
&
(53.22)&:\quad \cdots,\qquad \cdots,
\\
(53.23)&:\quad
C_{\rm mix}^2=\int \|B(z)\|_{\rm op}^2\,d\alpha(z),
\qquad
\Theta=C_hC_{\rm mix},
\qquad
C_s^{\rm eff}=(1+\tau)C_s,
\\
(55.9)&:\quad \cdots,\qquad \cdots,
&
(55.11)&:\quad \cdots,\qquad \cdots,
\\
(55.15)&:\quad
p_A(B_A)\leq CAw_Ae^{-A^2/2}
\leq C(1+\log A)e^{-A^2/2},
&
(55.17)&:\quad
\int |f_A'(x)|^2p_A(x)\,dx,
\\
(56.1)&:\quad \cdots,\qquad \cdots,
&
(56.4)&:\quad
\eta\in(0,\eta_0],\qquad
W_\eta(z)=e^{\eta |z|^2/2},
\\
(56.6)&:\quad \cdots,\qquad \cdots,
&
(56.7)&:\quad
\frac{\mathcal L_\alpha W_\eta}{W_\eta}
=\eta z\!\cdot\!\nabla\Psi-\eta m-\eta^2|z|^2,
\\
(56.8)&:\quad \cdots,\qquad \cdots,
&
(56.13)&:\quad
\int_{K_R} W_\eta\,d\alpha,
\\
(56.19)&:\quad \cdots,\qquad \cdots,
&
(56.26)&:\quad \cdots,\qquad \cdots .
\end{align*}
In the definition following (56.18), the conditional law is
\(\alpha(\,\cdot\mid K_R)\).  The prose-level mathematical tokens in
the affected appendices are \(H^j\), \(L^2\), \(W^{1,1}\), \(C_X\),
\(L^1\), \(W^{2,2}\), \(B_A\), \(e^{-x^2/2}\leq1\), \(K_R\),
\(c_R\), and \(1/4\), each with its ordinary inline-mathematics
meaning.  Finally, the V56 record reads \(566248\) bytes and \(16870\)
lines.  These corrections change no theorem hypothesis or conclusion;
they restore the intended notation and integration differentials.

\subsection{Metric vacuum response: exact affine radial budget}

Fix a shell \((n,q)\), suppress those indices, and write the V46
vacuum counterterm as
\[
 V_{\rm vac}(z)
 =a_n^d\sum_{y\in\mathcal Y_n}
   \delta\Lambda_y\sqrt{\det g_y(z)}.                 \tag{57.1}
\]
Let \(\|\cdot\|_y\) be the metric-response norm used by the V46
determinant estimate, and assume that for \(k=0,1\) there are
nonnegative numbers \(r^{\rm vac}_{k,y}\) such that
\[
 \sup_{|a|=1}\|D_zg_y(z)[a]\|_y
 \leq r^{\rm vac}_{1,y}|z|+r^{\rm vac}_{0,y}.         \tag{57.2}
\]
The sums below may be finite at a fixed cutoff; uniformity becomes a
separate scaling condition.

\begin{proposition}[Vacuum slope and force]
\label{prop:v16v57-vacuum}
Under (57.2),
\[
 \|\nabla_zV_{\rm vac}(z)\|
 \leq L^{\rm vac}_1|z|+L^{\rm vac}_0,                 \tag{57.3}
\]
where the directly certified coefficients are
\[
 L^{\rm vac}_k
 =C_1a_n^d\sum_{y\in\mathcal Y_n}
   |\delta\Lambda_y|\,r^{\rm vac}_{k,y},
 \qquad k=0,1.                                        \tag{57.4}
\]
Consequently, for every \(\tau_{\rm vac}>0\),
\[
 z\!\cdot\!\nabla_zV_{\rm vac}(z)
 \geq-(L^{\rm vac}_1+\tau_{\rm vac})|z|^2
      -\frac{(L^{\rm vac}_0)^2}{4\tau_{\rm vac}}.     \tag{57.5}
\]
\end{proposition}

\begin{proof}
For every unit vector \(a\), the V46 determinant differential gives
\[
 |D_zV_{\rm vac}(z)[a]|
 \leq C_1a_n^d\sum_y|\delta\Lambda_y|
        \|D_zg_y(z)[a]\|_y .
\]
Insert (57.2) and take the supremum over \(a\) to obtain (57.3)--(57.4).
Cauchy--Schwarz and \(L^{\rm vac}_0|z|
\leq\tau_{\rm vac}|z|^2+
(L^{\rm vac}_0)^2/(4\tau_{\rm vac})\) give (57.5).
\end{proof}

The coefficient \(L^{\rm vac}_1\) consumes the quadratic radial return
at every radius.  The coefficient \(L^{\rm vac}_0\) is an additive
force; Young's inequality converts it into an arbitrarily chosen
quadratic loss \(\tau_{\rm vac}\) and the compensating constant in
(57.5).  This distinction is lost if both coefficients are compressed
into a single Lipschitz constant.

\subsection{Curvature shell response}

Write \(V_R(z)=\delta c_R S_{\rm full}(g(z))\), with the bulk,
face, and joint owner norms of V47 denoted respectively by
\(\mathcal B,\mathcal F,\mathcal J\).  Assume the affine response bounds
\begin{align}
 \mathcal B(D_zg(z)[a])
 &\leq B_1|z|+B_0, \nonumber\\
 \mathcal F(D_zg(z)[a])
 &\leq F_1|z|+F_0,                                   \tag{57.6}\\
 \mathcal J(D_zg(z)[a])
 &\leq J_1|z|+J_0
\end{align}
for every \(|a|=1\).  These are statements about the actual
reconstruction response, including the V48 collar traces; they are
not consequences of the formal action alone.

\begin{proposition}[Curvature slope and force]
\label{prop:v16v57-curvature}
Under (57.6), the exact V47 owner inequality implies
\[
 \|\nabla_zV_R(z)\|
 \leq L^R_1|z|+L^R_0,                                \tag{57.7}
\]
with
\[
 L^R_k
 =|\delta c_R|
   \bigl(K_GB_k+K_\Pi F_k+K_JJ_k\bigr),
 \qquad k=0,1.                                        \tag{57.8}
\]
Thus, for every \(\tau_R>0\),
\[
 z\!\cdot\!\nabla_zV_R(z)
 \geq-(L^R_1+\tau_R)|z|^2-\frac{(L^R_0)^2}{4\tau_R}.
                                                               \tag{57.9}
\]
\end{proposition}

\begin{proof}
The chain rule (56.29) holds with \(z\) replaced by any unit direction
\(a\):
\[
 D_zV_R(z)[a]
 =\delta c_R DS_{\rm full}(g(z))[D_zg(z)[a]].
\]
Apply the V47 bulk--face--joint owner inequality and then (57.6).
Taking the supremum over \(a\) proves (57.7)--(57.8); the final estimate
is the same Young inequality used in (57.5).
\end{proof}

\subsection{A general radial compiler}

\begin{lemma}[Slope--force aggregation]
\label{lem:v16v57-aggregate}
Let \(I_1,\ldots,I_N\) be differentiable and suppose
\[
 \|\nabla I_i(z)\|\leq L_{1,i}|z|+L_{0,i},
 \qquad L_{1,i},L_{0,i}\geq0.                         \tag{57.10}
\]
For arbitrary \(\tau_i>0\),
\[
 z\!\cdot\!\nabla\sum_{i=1}^NI_i(z)
 \geq-\sum_{i=1}^N(L_{1,i}+\tau_i)|z|^2
      -\sum_{i=1}^N\frac{L_{0,i}^2}{4\tau_i}.         \tag{57.11}
\]
\end{lemma}

\begin{proof}
For each \(i\),
\[
 z\!\cdot\!\nabla I_i
 \geq-|z|\,\|\nabla I_i\|
 \geq-L_{1,i}|z|^2-L_{0,i}|z|.
\]
Apply \(L_{0,i}|z|\leq\tau_i|z|^2+
L_{0,i}^2/(4\tau_i)\) and sum.
\end{proof}

There is a useful optimized form when a total Young budget
\(\tau=\sum_i\tau_i\) is prescribed.  Cauchy--Schwarz shows
\[
 \sum_i\frac{L_{0,i}^2}{\tau_i}
 \geq\frac{(\sum_iL_{0,i})^2}{\sum_i\tau_i},          \tag{57.12}
\]
with equality at
\(\tau_i=\tau L_{0,i}/\sum_jL_{0,j}\) for positive \(L_{0,i}\).
Hence (57.11) can be written
\[
 z\!\cdot\!\nabla\sum_iI_i(z)
 \geq-\left(\sum_iL_{1,i}+\tau\right)|z|^2
      -\frac{(\sum_iL_{0,i})^2}{4\tau}.              \tag{57.13}
\]

\subsection{Why the cosmological source is not the conformal stabilizer}

\begin{proposition}[Volume-preserving conformal no-go]
\label{prop:v16v57-volume-no-go}
Let \(s\mapsto g_s\) be any differentiable metric path satisfying
\(\operatorname{Vol}(g_s)=\operatorname{Vol}(g_0)\).  Then the
cosmological functional
\[
 S_\Lambda(g_s)=\Lambda\operatorname{Vol}(g_s)        \tag{57.14}
\]
is constant on the path.  Every path derivative that exists vanishes.
In particular it supplies no restoring term on a
volume-preserving high-frequency conformal path.
\end{proposition}

\begin{proof}
Equation (57.14) is constant by its hypothesis.  Differentiating the
constant identity proves the claim at every order for which the path
and functional are differentiable.
\end{proof}

The conclusion is compatible with retaining the Einstein source:
variation away from the constrained path still gives the V46 tensor
\(T_{\rm vac}^{\mu\nu}=\Lambda g^{\mu\nu}\).  It says only that this
source cannot repair the fixed-volume conformal obstruction already
exhibited in frozen V39.  Thus setting \(\Lambda=0\) is neither required
nor useful for the radial argument.

\subsection{Source-preserving stabilizer corridor}

Let \(U_n\) be the augmented soft potential.  Separate its declared
coercive core, a stabilizer \(S_{{\rm stab},n}\), and the interaction
terms \(I_{i,n}\):
\[
 U_n=U_{{\rm core},n}+S_{{\rm stab},n}
     +\sum_{i=1}^NI_{i,n}.                            \tag{57.15}
\]
Suppose
\begin{align}
 z\!\cdot\!\nabla U_{{\rm core},n}(z)
 &\geq\rho_{{\rm core},n}|z|^2-C_{{\rm core},n},
                                                        \tag{57.16}\\
 z\!\cdot\!\nabla S_{{\rm stab},n}(z)
 &\geq\rho_{{\rm stab},n}|z|^2-C_{{\rm stab},n},       \tag{57.17}\\
 \|\nabla I_{i,n}(z)\|
 &\leq L_{1,i,n}|z|+L_{0,i,n}.                         \tag{57.18}
\end{align}

\begin{theorem}[Uniform stabilizer corridor]
\label{thm:v16v57-corridor}
Choose \(\tau_{i,n}>0\).  If
\[
 \rho_*:=
 \inf_n\left[
 \rho_{{\rm core},n}+\rho_{{\rm stab},n}
 -\sum_i(L_{1,i,n}+\tau_{i,n})
 \right]>0                                             \tag{57.19}
\]
and
\[
 C_*:=
 \sup_n\left[
 C_{{\rm core},n}+C_{{\rm stab},n}
 +\sum_i\frac{L_{0,i,n}^2}{4\tau_{i,n}}
 \right]<\infty,                                       \tag{57.20}
\]
then
\[
 z\!\cdot\!\nabla U_n(z)\geq\rho_*|z|^2-C_*            \tag{57.21}
\]
uniformly in \(n\).  Hence the V56 Lyapunov and soft-Poincar\'e
compiler applies once its local-core and conditional-law hypotheses
are also verified.
\end{theorem}

\begin{proof}
Add (57.16), (57.17), and Lemma
\ref{lem:v16v57-aggregate}, then use (57.19)--(57.20).
\end{proof}

For the continuum Einstein equation, radial control is only one side
of the corridor.  If the stabilizer enters with a vanishing coefficient
\(\epsilon_n\), its metric source must also disappear against every
fixed smooth test variation \(h\):
\[
 \mathbb E_{\mu_n}
 \left|\epsilon_nDS_{{\rm stab},n}(g_n)[h]\right|
 \longrightarrow0.                                    \tag{57.22}
\]
This condition retains the vacuum and curvature sources while removing
the auxiliary source.  It does not follow merely from
\(\epsilon_n\to0\).  Reflection positivity and transfer through the
reconstruction map remain independent requirements.

For the specific choice
\[
 S_{{\rm stab},n}(g)=\int R(g)^2\,d{\rm vol}_g,        \tag{57.23}
\]
frozen V71 computed
\[
 H_{\mu\nu}^{(R^2)}
 =2RR_{\mu\nu}-\frac12R^2g_{\mu\nu}
  +2(g_{\mu\nu}\Box-\nabla_\mu\nabla_\nu)R,
 \qquad
 g^{\mu\nu}H_{\mu\nu}^{(R^2)}=6\Box R                \tag{57.24}
\]
in four dimensions.  The bound
\[
 \sup_n\epsilon_n\int R(g_n)^2\,d{\rm vol}_{g_n}<\infty
                                                               \tag{57.25}
\]
is sufficient for distributional decoupling of the trace
\(\epsilon_n\Box R(g_n)\) against fixed smooth tests with uniformly
bounded \(\Box\)-norm.  It is not a proof of full tensor decoupling,
reflection positivity, or a positive radial coefficient in (57.17).

\subsection{A spectral obstruction to naive vanishing stabilization}

The conformal instability and a fourth-order stabilizer have the
scalar spectral envelope
\[
 q_\epsilon(\lambda)
 =\rho_0-a_2\lambda+\epsilon b_4\lambda^2,
 \qquad \lambda\geq0,\quad a_2,b_4>0.                 \tag{57.26}
\]

\begin{proposition}[Vanishing fourth-order floor no-go]
\label{prop:v16v57-spectral-no-go}
For every \(\epsilon>0\),
\[
 \inf_{\lambda\geq0}q_\epsilon(\lambda)
 =\rho_0-\frac{a_2^2}{4\epsilon b_4}.                 \tag{57.27}
\]
Consequently, if \(\epsilon_n\to0\), this envelope has no
\(n\)-uniform lower bound on an unbounded spectral interval.  More
precisely, \(q_{\epsilon_n}\geq\rho_*>0\) for all \(\lambda\geq0\)
requires
\[
 \epsilon_n\geq
 \frac{a_2^2}{4b_4(\rho_0-\rho_*)},
 \qquad \rho_*<\rho_0.                                \tag{57.28}
\]
\end{proposition}

\begin{proof}
Complete the square:
\[
 q_\epsilon(\lambda)
 =\epsilon b_4
  \left(\lambda-\frac{a_2}{2\epsilon b_4}\right)^2
  +\rho_0-\frac{a_2^2}{4\epsilon b_4}.                \tag{57.29}
\]
The minimizer \(a_2/(2\epsilon b_4)\) lies in
\([0,\infty)\), proving (57.27).  Rearranging
\(\rho_0-a_2^2/(4\epsilon_nb_4)\geq\rho_*\) proves
(57.28).
\end{proof}

This is an exact theorem for the scalar envelope (57.26), not a
spectral theorem for the full gravitational Hessian.  It identifies
the next bottleneck: on an unbounded conformal spectrum, a single
coefficient cannot both vanish and provide an all-scale uniform floor
unless another mechanism removes or controls the intermediate
minimizing band.  Candidate mechanisms are a constraint-first
conformal reduction, a scale-dependent family of stabilizers, or a
non-real contour with a separately proved positive reconstructed law.

\begin{assessment}[What V57 closes]
\label{ass:v16v57-status}
V57 repairs the isolated source escapes in V51--V56 without altering
the installed prefix.  It converts the V46 vacuum and V47 curvature
owner estimates into explicit radial slope and additive-force
coefficients, aggregates them with the optimal Young allocation, and
states the two-sided stabilizer corridor: uniform radial return at
finite cutoff and vanishing auxiliary metric source in the continuum.

The cosmological term retains its Einstein source but is rigorously
excluded as a fixed-volume conformal stabilizer.  A fourth-order
spectral envelope shows why a naively vanishing \(R^2\) coefficient
cannot uniformly control an unbounded negative second-order band.
No physical stabilizer, reflection-positive contour, conformal
constraint reduction, or full tensor decoupling has yet been proved.
\end{assessment}

\begin{decision}[V58 intermediate-band alternatives]
\label{dec:v16v57-next}
Resolve the spectral bottleneck upstream.  Formulate a cutoff-dependent
conformal spectrum and compare three routes: constraint-first removal,
a multiscale stabilizer with a uniform corridor, and a fixed
stabilizer followed by source decoupling only after tightness.
Prove the sharp order-of-limits statement that can actually coexist
with (57.19) and (57.22), and reject any route that merely assumes the
desired positive law.
\end{decision}

\begin{firewall}[Claim boundary after V57]
\label{fire:v16v57}
The estimates in this version are conditional analytic compilers.
They do not establish the sign of the gravitational measure, a
uniform physical conditional gap, reflection positivity, the
continuum Ward identities, or the Standard Model--Einstein continuum.
\end{firewall}

\section*{Version V57 append-only record}
\addcontentsline{toc}{section}{Version V57 append-only record}
The complete V56 source was retained byte for byte before this
continuation.  It had \(578518\) bytes, \(17236\) lines, and SHA--256
\[
 \texttt{1A0D6510AEE493E35254961CAE167930B991018E2F0D7709C21AD5445D5D26C9}.
\]
V57 added an explicit append-only erratum, exact vacuum and curvature
slope--force budgets, the uniform source-preserving stabilizer
corridor, and the scalar spectral obstruction to naive vanishing
fourth-order stabilization.


\section{V58: cutoff-exact conformal obstruction and constraint-first escape}
\label{sec:v16v58}

\subsection{The finite-window spectral minimum}

The scalar envelope from V57 must first be read on the actual cutoff
window.  Let
\[
 q_{\epsilon,\Lambda}(\lambda)
 =\rho_0-a_2\lambda+\epsilon b_4\lambda^2,
 \qquad 0\leq\lambda\leq\Lambda,                      \tag{58.1}
\]
where \(a_2,b_4,\Lambda>0\) and \(\epsilon\geq0\).

\begin{lemma}[Exact cutoff minimum]
\label{lem:v16v58-window}
For \(\epsilon>0\),
\[
 \min_{0\leq\lambda\leq\Lambda}q_{\epsilon,\Lambda}(\lambda)
 =
 \begin{cases}
 \rho_0-a_2\Lambda+\epsilon b_4\Lambda^2,
 &0<\epsilon<\dfrac{a_2}{2b_4\Lambda},\\
 \rho_0-\dfrac{a_2^2}{4\epsilon b_4},
 &\epsilon\geq\dfrac{a_2}{2b_4\Lambda}.
 \end{cases}                                         \tag{58.2}
\]
For \(\epsilon=0\), the minimum is
\(\rho_0-a_2\Lambda\).
\end{lemma}

\begin{proof}
For \(\epsilon>0\), the unconstrained minimizer is
\(\lambda_\epsilon=a_2/(2\epsilon b_4)\).
If \(\lambda_\epsilon>\Lambda\), the derivative is negative throughout
the interval and the minimum is the endpoint value at \(\Lambda\).
If \(\lambda_\epsilon\leq\Lambda\), completing the square as in
(57.29) gives the second line.  The case \(\epsilon=0\) is immediate.
\end{proof}

\begin{theorem}[Expanding-window no-go]
\label{thm:v16v58-expanding-no-go}
Let \(\Lambda_n\to\infty\) and \(\epsilon_n\to0\).  Then
\[
 \inf_{0\leq\lambda\leq\Lambda_n}
 \bigl(\rho_0-a_2\lambda+\epsilon_nb_4\lambda^2\bigr)
 \longrightarrow-\infty.                             \tag{58.3}
\]
Thus no diagonal choice of a vanishing fourth-order coefficient gives
a uniform positive quadratic floor on expanding spectral windows.
\end{theorem}

\begin{proof}
Set
\[
 m_n=\min\{\Lambda_n,\epsilon_n^{-1/2}\},              \tag{58.4}
\]
with \(\epsilon_n^{-1/2}=+\infty\) if \(\epsilon_n=0\).
Both entries in the minimum tend to infinity, so \(m_n\to\infty\).
Moreover \(\epsilon_nm_n^2\leq1\).  Since \(m_n\) belongs to the cutoff
window,
\[
 \inf_{[0,\Lambda_n]}q_{\epsilon_n,\Lambda_n}
 \leq q_{\epsilon_n,\Lambda_n}(m_n)
 \leq\rho_0-a_2m_n+b_4\longrightarrow-\infty.         \tag{58.5}
\]
\end{proof}

This proof does not depend on which branch of (58.2) occurs.  It closes
the possible loophole in which the moving cutoff might avoid the
formal minimizer \(a_2/(2\epsilon_nb_4)\).

\subsection{Moving multiscale stiffness}

A scale-dependent stabilizer is more general than
\(\epsilon_nb_4\lambda^2\).  Let
\[
 0\leq\lambda_1\leq\lambda_2\leq\cdots,\qquad
 \lambda_j\longrightarrow\infty,                      \tag{58.6}
\]
and let the \(n\)-th cutoff retain indices \(j\leq J_n\), where
\(J_n\to\infty\).  Denote the nonnegative stabilizer stiffness on mode
\(j\) by \(s_{n,j}\).

\begin{proposition}[Pointwise-decoupling obstruction]
\label{prop:v16v58-pointwise}
Suppose there is a number \(\rho_*>0\) such that
\[
 \rho_0-a_2\lambda_j+s_{n,j}\geq\rho_*
 \qquad (j\leq J_n)                                   \tag{58.7}
\]
for every \(n\).  Then each fixed \(j\) satisfies, for all sufficiently
large \(n\),
\[
 s_{n,j}\geq a_2\lambda_j-(\rho_0-\rho_*).             \tag{58.8}
\]
Consequently \(s_{n,j}\to0\) for every fixed \(j\) is incompatible
with (58.7) whenever the retained spectrum contains a mode with
\(a_2\lambda_j>\rho_0-\rho_*\).
\end{proposition}

\begin{proof}
For fixed \(j\), the condition \(J_n\to\infty\) implies \(j\leq J_n\)
for every sufficiently large \(n\).  Rearranging (58.7) then gives
(58.8).  If its right-hand side is positive, it contradicts
\(s_{n,j}\to0\).
\end{proof}

The proposition is a Hessian-level statement.  Pointwise vanishing of
\(s_{n,j}\) is stronger than the first-variation source decoupling
(57.22), so the two must not be identified.  It nevertheless rules
out a common proposed shortcut: moving the stabilizer entirely into
the ultraviolet while retaining a uniform positive quadratic form on
every eventually fixed conformal mode.

\subsection{The order of the cutoff and stabilizer limits}

\begin{corollary}[Noncommuting coercive limits]
\label{cor:v16v58-order}
For the envelope (58.1), the following two procedures cannot share a
positive lower bound independent of both parameters:
\begin{enumerate}
\item take \(\Lambda\to\infty\) at a fixed stabilizer coefficient
      \(\epsilon>0\);
\item then take \(\epsilon\downarrow0\).
\end{enumerate}
Indeed,
\[
 \inf_{\lambda\geq0}q_{\epsilon,\infty}(\lambda)
 =\rho_0-\frac{a_2^2}{4\epsilon b_4}
 \longrightarrow-\infty
 \quad\text{as }\epsilon\downarrow0.                  \tag{58.9}
\]
The same failure holds for every diagonal
\((\epsilon_n,\Lambda_n)\to(0,\infty)\) by Theorem
\ref{thm:v16v58-expanding-no-go}.
\end{corollary}

At fixed \(\epsilon\), a positive floor is possible only when
\(\rho_0>a_2^2/(4\epsilon b_4)\).  Thus a fixed stabilizer may be used
to construct a regulated measure and take a cutoff limit, but the
resulting gap estimate degenerates before the stabilizer is removed.
Tightness uniform in \(\epsilon\), if true for another reason, must be
proved independently; it cannot be imported from this quadratic
floor.

\subsection{Constraint-first finite-dimensional reduction}

The preceding no-go acts only on the unstable conformal directions.
This motivates an exact constraint reduction before applying the
probability compiler.  Let \(H_n\) be a finite-dimensional Euclidean
space with an orthogonal decomposition
\[
 H_n=K_n\oplus C_n,                                   \tag{58.10}
\]
where \(C_n\) contains the declared conformal directions and \(K_n\)
is the constrained physical slice.  Let \(P_n\) be the orthogonal
projection onto \(K_n\), and let \(A_n\) be a symmetric operator.

\begin{theorem}[Restricted Gaussian gap]
\label{thm:v16v58-restricted-gaussian}
Assume
\[
 \langle u,A_nu\rangle\geq\rho_K|u|^2
 \qquad (u\in K_n)                                    \tag{58.11}
\]
with \(\rho_K>0\) independent of \(n\).  Then
\[
 d\gamma_n^K(u)
 =Z_n^{-1}\exp\!\left(-\frac12\langle u,A_nu\rangle\right)
 \,d{\rm vol}_{K_n}(u)                                \tag{58.12}
\]
is a probability measure on \(K_n\) and satisfies
\[
 \operatorname{Var}_{\gamma_n^K}(f)
 \leq\frac1{\rho_K}\int_{K_n}|\nabla_{K_n}f|^2
 \,d\gamma_n^K                                       \tag{58.13}
\]
for every smooth \(f\).  No lower bound on \(A_n\) over \(C_n\) is
required.
\end{theorem}

\begin{proof}
The compression \(A_n^K=P_nA_nP_n|_{K_n}\) is symmetric and bounded
below by \(\rho_K I\), so (58.12) is the centered Gaussian on \(K_n\)
with covariance \((A_n^K)^{-1}\).  Diagonalize \(A_n^K\).  The
one-dimensional Gaussian Poincar\'e inequality in each eigenvector
direction has constant at most \(1/\rho_K\); tensorization gives
(58.13).  Equivalently, the Bakry--\'Emery Hessian bound on \(K_n\) is
\(\nabla_{K_n}^2(\langle u,A_nu\rangle/2)=A_n^K
\succeq\rho_KI\).
\end{proof}

\begin{proposition}[Restricted nonlinear corridor]
\label{prop:v16v58-restricted-nonlinear}
On \(K_n\), suppose
\[
 u\!\cdot\!\nabla_{K_n}U_{{\rm core},n}(u)
 \geq\rho_{K,n}|u|^2-C_{K,n}                          \tag{58.14}
\]
and the restricted interactions satisfy
\[
 \|\nabla_{K_n}I_{i,n}(u)\|
 \leq L^K_{1,i,n}|u|+L^K_{0,i,n}.                     \tag{58.15}
\]
If the analogues of (57.19)--(57.20), with no conformal stabilizer,
hold on \(K_n\), then the radial return (57.21) holds on \(K_n\).
\end{proposition}

\begin{proof}
Apply Lemma~\ref{lem:v16v57-aggregate} to the intrinsic gradients on
the Euclidean slice \(K_n\) and add (58.14).
\end{proof}

This is the only one of the three routes considered here that evades
the V57--V58 spectral obstruction at the level of the stated
hypotheses: the bad directions are absent from the integration domain.
It is not yet a construction of the gravitational constraint.  A
valid continuum route must additionally supply:
\begin{enumerate}
\item a geometrically defined constraint map and a regular slice;
\item its Faddeev--Popov or coarea Jacobian, with a sign or absolute
      continuity statement sufficient for a positive law;
\item compatibility of the slice with shell transfer and reflection;
\item reconstruction of the full Einstein Ward identity, including
      the constraint reaction and gauge identities.
\end{enumerate}

\subsection{Decision table forced by the no-go}

\begin{center}
\begin{tabular}{p{0.24\linewidth}p{0.31\linewidth}p{0.33\linewidth}}
\hline
Route & Rigorous gain & Remaining obstruction\\
\hline
Vanishing \(R^2\) coefficient
& source may decouple under V71 bounds
& uniform conformal floor fails by (58.3)\\
Moving ultraviolet stabilizer
& can control newly introduced high modes
& cannot vanish on every fixed unstable mode by (58.8)\\
Fixed stabilizer, two limits
& may construct a regulated cutoff limit
& coercive constant diverges negatively as it is removed\\
Constraint-first slice
& removes the bad sector before the gap estimate
& slice Jacobian, positivity, transfer, and Ward recovery unproved\\
\hline
\end{tabular}
\end{center}

\begin{assessment}[What V58 closes]
\label{ass:v16v58-status}
V58 proves the cutoff-exact and diagonal versions of the conformal
stabilizer obstruction.  It also proves that an arbitrary moving
spectral stiffness cannot both vanish on every fixed mode and retain a
uniform floor on an eventually retained unstable mode.  The
fixed-stabilizer order of limits therefore supplies no uniform gap.

The constraint-first route survives this obstruction, and its exact
finite-dimensional Gaussian and nonlinear radial compilers are now
stated.  The essential upstream bottleneck is no longer a choice of
vanishing coefficient; it is construction of a regular positive
constraint slice whose Jacobian, transfer properties, and Ward
identity can be controlled.
\end{assessment}

\begin{decision}[V59 coarea slice and Jacobian load]
\label{dec:v16v58-next}
Build the finite-dimensional nonlinear constraint layer.  Starting
from a submersion \(C_n(x)=0\), derive the exact coarea density on the
slice, compute its logarithmic gradient and Hessian loads, and state
checkable conditions under which they fit the V57 slope--force and
V53 Poincar\'e compilers.  Track the reaction term in the ambient
integration-by-parts identity so that constraint removal is not
mistaken for the unconstrained Einstein equation.
\end{decision}

\begin{firewall}[Claim boundary after V58]
\label{fire:v16v58}
The spectral statements concern declared quadratic conformal
envelopes and modewise stiffness.  They do not prove that the full
gravitational Hessian is diagonal in this basis.  The restricted
Gaussian theorem does not construct a diffeomorphism gauge slice,
solve the Hamiltonian constraint, prove determinant positivity, or
recover the Einstein equation.
\end{firewall}

\section*{Version V58 append-only record}
\addcontentsline{toc}{section}{Version V58 append-only record}
The complete V57 source was retained byte for byte before this
continuation.  It had \(593477\) bytes, \(17686\) lines, and SHA--256
\[
 \texttt{F9A1E42488E74D75249C3E2A81BFF16773198C289F8D4CC10277E4C55606BC91}.
\]
V58 proved the expanding-window and moving-stiffness no-go theorems,
made the order-of-limits failure explicit, and isolated a
constraint-first slice as the remaining real-positive coercivity
route within the present alternatives.


\section{V59: nonlinear coarea slice, Jacobian loads, and constraint reaction}
\label{sec:v16v59}

\subsection{The positive conditional law}

Let \(\Omega\subset\mathbb R^m\) be open and let
\[
 C:\Omega\longrightarrow\mathbb R^r,\qquad r<m,       \tag{59.1}
\]
be a \(C^3\) submersion in a neighbourhood of
\[
 M=C^{-1}(0).                                         \tag{59.2}
\]
Thus \(M\) is a smooth submanifold of dimension \(d=m-r\).  Put
\[
 A_C(x)=DC(x)DC(x)^{\mathsf T},\qquad
 J_C(x)=\sqrt{\det A_C(x)}.                           \tag{59.3}
\]
The submersion hypothesis makes \(J_C>0\) on \(M\).

\begin{proposition}[Coarea-conditioned Gibbs density]
\label{prop:v16v59-coarea}
Let \(U:\Omega\to\mathbb R\) be measurable and suppose
\[
 0<Z_M:=\int_M e^{-U(x)}J_C(x)^{-1}\,d{\rm vol}_M(x)
 <\infty.                                             \tag{59.4}
\]
Then the normalized zero-level conditional law is
\[
 d\nu_M(x)
 =Z_M^{-1}e^{-\Phi_M(x)}\,d{\rm vol}_M(x),\qquad
 \Phi_M=U|_M+\log J_C|_M.                             \tag{59.5}
\]
Equivalently, for compactly supported continuous \(f\),
\[
 \int_\Omega f(x)e^{-U(x)}\delta_0(C(x))\,dx
 =\int_M f(x)e^{-U(x)}J_C(x)^{-1}\,d{\rm vol}_M(x).
                                                               \tag{59.6}
\]
\end{proposition}

\begin{proof}
The coarea formula states
\[
 \int_\Omega g(x)J_C(x)\,dx
 =\int_{\mathbb R^r}\int_{C^{-1}(y)}
   g(x)\,d{\rm vol}_{C^{-1}(y)}(x)\,dy.               \tag{59.7}
\]
Apply it to
\(g=f e^{-U}J_C^{-1}\) and evaluate the resulting level-set
disintegration at \(y=0\).  This gives (59.6); normalization gives
(59.5).
\end{proof}

The factor in (59.5) is positive.  This proves positivity of the
finite-dimensional law obtained by conditioning Lebesgue measure.
It does not prove that this law equals a diffeomorphism-orbit quotient:
that identification additionally requires an orbit chart, treatment
of residual copies, and the appropriate gauge determinant.

\subsection{Exact logarithmic Jacobian derivatives}

For \(v\in\mathbb R^m\), write
\[
 B_v=D(DC)[v]=D^2C[v,\cdot].                           \tag{59.8}
\]
Then
\[
 DA_C[v]=B_vDC^{\mathsf T}+DCB_v^{\mathsf T}.         \tag{59.9}
\]

\begin{lemma}[Jacobian gradient and Hessian]
\label{lem:v16v59-jacobian}
At every point where \(C\) is a submersion,
\[
 D\log J_C[v]
 =\frac12\operatorname{tr}\!\left(A_C^{-1}DA_C[v]\right),
                                                               \tag{59.10}
\]
and
\[
 D^2\log J_C[v,w]
 =\frac12\operatorname{tr}\!\left(
 A_C^{-1}D^2A_C[v,w]
 -A_C^{-1}DA_C[w]A_C^{-1}DA_C[v]\right).             \tag{59.11}
\]
\end{lemma}

\begin{proof}
Since \(\log J_C=\frac12\log\det A_C\), Jacobi's determinant
formula gives (59.10).  Differentiate it in direction \(w\), using
\[
 D(A_C^{-1})[w]=-A_C^{-1}DA_C[w]A_C^{-1},             \tag{59.12}
\]
and cyclicity of the trace to obtain (59.11).
\end{proof}

Assume on a tubular neighbourhood of \(M\) that
\[
 \sigma_{\min}(DC)\geq s>0,\quad
 \|DC\|_{\rm op}\leq M_1,\quad
 \|D^2C\|_{\rm op}\leq M_2,\quad
 \|D^3C\|_{\rm op}\leq M_3.                           \tag{59.13}
\]
Here the multilinear operator norms are Euclidean.

\begin{proposition}[Checkable Jacobian loads]
\label{prop:v16v59-jacobian-load}
Under (59.13), for every unit \(v\),
\[
 |D\log J_C[v]|\leq
 G_J:=\frac{rM_1M_2}{s^2},                            \tag{59.14}
\]
and
\[
 |D^2\log J_C[v,v]|
 \leq H_J:=
 \frac{r(M_1M_3+M_2^2)}{s^2}
 +\frac{2rM_1^2M_2^2}{s^4}.                          \tag{59.15}
\]
\end{proposition}

\begin{proof}
The spectral bound in (59.13) gives
\(\|A_C^{-1}\|_{\rm op}\leq s^{-2}\).  From (59.9),
\[
 \|DA_C[v]\|_{\rm op}\leq2M_1M_2.                    \tag{59.16}
\]
Since \(|\operatorname{tr}T|\leq r\|T\|_{\rm op}\),
(59.10) yields (59.14).  Differentiating (59.9) once more gives
\[
 \|D^2A_C[v,v]\|_{\rm op}
 \leq2M_1M_3+2M_2^2.                                 \tag{59.17}
\]
Apply the trace bound to the first term in (59.11).  For the second,
conjugate \(DA_C[v]\) by \(A_C^{-1/2}\); it is symmetric, so
\[
 \operatorname{tr}
 \bigl(A_C^{-1}DA_C[v]A_C^{-1}DA_C[v]\bigr)
 \leq r s^{-4}(2M_1M_2)^2.                           \tag{59.18}
\]
Combining the two half-weighted estimates proves (59.15).
\end{proof}

\subsection{Intrinsic curvature and the Poincar\'e load}

Let \(P_x\) denote orthogonal projection onto \(T_xM\), and let
\({\rm II}\) be the second fundamental form with convention
\[
 {\rm II}(v,w)=(D_v\widetilde w)^\perp.               \tag{59.19}
\]
Differentiating \(C(\gamma(t))=0\) twice along a unit-speed curve in
\(M\) shows
\[
 DC\,{\rm II}(v,v)=-D^2C[v,v].
\]
Since the inverse of \(DC\) on the normal space has norm at most
\(s^{-1}\),
\[
 \|{\rm II}\|_{\rm op}\leq\kappa:=\frac{M_2}{s}.      \tag{59.20}
\]

\begin{lemma}[Restriction Hessian]
\label{lem:v16v59-restriction-hessian}
For any \(C^2\) ambient function \(F\) and tangent vectors \(v,w\),
\[
 {\rm Hess}_M(F|_M)(v,w)
 =D^2F(v,w)+\langle\nabla F,{\rm II}(v,w)\rangle.     \tag{59.21}
\]
\end{lemma}

\begin{proof}
Use the Gauss decomposition
\(D_v\widetilde w=\nabla_v^M w+{\rm II}(v,w)\) in the definitions of
the ambient and intrinsic Hessians.
\end{proof}

Consequently the intrinsic Jacobian Hessian obeys
\[
 {\rm Hess}_M(\log J_C)(v,v)
 \geq-H_J-G_J\kappa
 \qquad (|v|=1,\ v\in T_xM).                          \tag{59.22}
\]

\begin{theorem}[Coarea Bakry--\'Emery criterion]
\label{thm:v16v59-coarea-gap}
Assume \(M\) is complete and without boundary.  Suppose for every unit
tangent \(v\),
\begin{align}
 D^2U(v,v)&\geq\rho_U,                                \tag{59.23}\\
 |(\nabla U)^\perp|&\leq G_U,                         \tag{59.24}\\
 {\rm Ric}_M(v,v)&\geq-K_{\rm Ric}.                   \tag{59.25}
\end{align}
If
\[
 \rho_M:=
 \rho_U-\kappa G_U-H_J-G_J\kappa-K_{\rm Ric}>0,       \tag{59.26}
\]
then the coarea law (59.5) satisfies
\[
 \operatorname{Var}_{\nu_M}(f)
 \leq\frac1{\rho_M}\int_M|\nabla_Mf|^2\,d\nu_M.       \tag{59.27}
\]
\end{theorem}

\begin{proof}
Lemma~\ref{lem:v16v59-restriction-hessian}, (59.20), and (59.24)
give
\[
 {\rm Hess}_M(U|_M)(v,v)\geq\rho_U-\kappa G_U.         \tag{59.28}
\]
Add (59.22) and the Ricci bound (59.25).  The result is
\[
 {\rm Ric}_M+{\rm Hess}_M\Phi_M\succeq\rho_Mg_M.
\]
The Bakry--\'Emery Poincar\'e criterion on the complete weighted
manifold then gives (59.27).
\end{proof}

The theorem shows exactly where a nearly singular slice fails: the
loads contain \(s^{-2}\), \(s^{-3}\), and \(s^{-4}\).  A formal
constraint with no uniform lower singular-value bound does not feed a
uniform gap.

\subsection{Radial alternative without global convexity}

Suppose \(M\) is star-shaped with respect to the chosen slice
coordinates \(u\), so the intrinsic radial vector \(u\) is defined.
If
\[
 \|\nabla_M\log J_C(u)\|
 \leq L^J_1|u|+L^J_0,                                 \tag{59.29}
\]
then
\[
 u\!\cdot\!\nabla_M\log J_C(u)
 \geq-(L^J_1+\tau_J)|u|^2
      -\frac{(L^J_0)^2}{4\tau_J}                     \tag{59.30}
\]
for every \(\tau_J>0\).  This is Lemma
\ref{lem:v16v57-aggregate} with one interaction.  Therefore the coarea
Jacobian enters the V57 corridor as one more declared slope--force
pair.  The uniform bound (59.14) is the special case
\(L^J_1=0\), \(L^J_0=G_J\).

\subsection{The exact constrained integration-by-parts identity}

The positive coarea law does not satisfy the unconstrained ambient
Euler--Lagrange identity.  Let
\[
 {\bf H}(x)=\sum_{\alpha=1}^d{\rm II}(e_\alpha,e_\alpha)             \tag{59.31}
\]
be the mean-curvature vector for an orthonormal tangent frame.  For a
constant ambient vector \(v\),
\[
 \operatorname{div}_M(Pv)=\langle{\bf H},v\rangle.    \tag{59.32}
\]

\begin{theorem}[Constraint score and reaction]
\label{thm:v16v59-reaction}
For every smooth compactly supported \(f\) on \(M\) and every constant
ambient vector \(v\),
\[
 \mathbb E_{\nu_M}\langle\nabla_Mf,v\rangle
 =
 \mathbb E_{\nu_M}\!\left[
 f\left\langle
 P\nabla U+\nabla_M\log J_C-{\bf H},v
 \right\rangle\right].                               \tag{59.33}
\]
In particular,
\[
 \mathbb E_{\nu_M}
 \left[P\nabla U+\nabla_M\log J_C-{\bf H}\right]=0.   \tag{59.34}
\]
Equivalently, the ambient force satisfies
\[
 \mathbb E_{\nu_M}\nabla U
 =\mathbb E_{\nu_M}{\cal R}_C,\qquad
 {\cal R}_C:=(I-P)\nabla U+{\bf H}-\nabla_M\log J_C.  \tag{59.35}
\]
\end{theorem}

\begin{proof}
Apply weighted Stokes' theorem to the tangent vector field \(fPv\):
\[
 \int_M\operatorname{div}_M(fPv)e^{-\Phi_M}\,d{\rm vol}_M
 =
 \int_M f\langle\nabla_M\Phi_M,Pv\rangle
 e^{-\Phi_M}\,d{\rm vol}_M.                           \tag{59.36}
\]
Expand the divergence using (59.32), substitute
\(\nabla_M\Phi_M=P\nabla U+\nabla_M\log J_C\), and rearrange to obtain
(59.33).  Set \(f=1\) by approximation to obtain (59.34).  Finally
write
\(\nabla U=P\nabla U+(I-P)\nabla U\) and use (59.34), which gives
(59.35).
\end{proof}

The vector \({\cal R}_C\) is the constraint reaction.  Recovering an
ambient Einstein equation requires showing that its physical pairing
is either the intended gauge/constraint identity or vanishes after
projection onto admissible tests.  Simply deleting the conformal
coordinate and then asserting
\(\mathbb E\nabla U=0\) would omit all three terms in (59.35).

\begin{assessment}[What V59 closes]
\label{ass:v16v59-status}
V59 constructs the exact positive coarea-conditioned law for a regular
finite-dimensional constraint, derives its logarithmic Jacobian
gradient and Hessian, and gives explicit loads in terms of the first
three derivatives and smallest singular value of the constraint map.
It supplies both a Bakry--\'Emery gap criterion and a radial
slope--force alternative.

The ambient integration-by-parts identity now carries the complete
constraint reaction: normal force, mean curvature, and logarithmic
coarea Jacobian.  This prevents the constraint-first escape from being
misreported as an unconstrained Einstein equation.  A gravitational
constraint map with uniform \(s,M_1,M_2,M_3\), compatible transfer,
and a physical reaction cancellation has not been constructed.
\end{assessment}

\begin{decision}[V60 gravitational slice audit]
\label{dec:v16v59-next}
Perform the mandatory YM/NS companion audit, then test candidate
metric constraints against the V59 certificate.  Start with a
finite-mode constant-volume conformal slice, compute its coarea
Jacobian and reaction explicitly, and determine whether it removes
only the global volume mode or the full high-frequency conformal
instability.  Use the failure, if any, to select the next upstream
constraint.
\end{decision}

\begin{firewall}[Claim boundary after V59]
\label{fire:v16v59}
The coarea density is a conditional Euclidean law, not yet a
Faddeev--Popov quotient.  The uniform derivative hypotheses have not
been verified for metric constraints, and the constrained reaction
has not been identified with a continuum diffeomorphism or Einstein
Ward identity.
\end{firewall}

\section*{Version V59 append-only record}
\addcontentsline{toc}{section}{Version V59 append-only record}
The complete V58 source was retained byte for byte before this
continuation.  It had \(604802\) bytes, \(17987\) lines, and SHA--256
\[
 \texttt{32B7017C375B63C1D3855D44AFB65EC49D86F0F603C2FB2B5E122BDC78BBB029}.
\]
V59 built the nonlinear coarea layer, quantified its Jacobian gap
loads, and derived the exact constrained score and ambient reaction
identity required before a constraint-first gravitational route can
claim an Einstein limit.


\section{V60: companion audit and the constant-volume conformal slice}
\label{sec:v16v60}

\subsection{Mandatory YM/NS companion audit}

\begin{audit}[V60 companion checkpoint]
\label{aud:v16v60-companions}
The newest active numeric companion notes inspected before the metric
constraint calculation were
\[
\begin{array}{c|r|r|l}
\hbox{programme}&\hbox{version}&\hbox{bytes}&\hbox{SHA--256}\\ \hline
\mathrm{YM}&55&477422&
\texttt{E0B3464E673F16B30DA54C46965C93FF5F47FBD421577AAE9DBC35319255F921}\\
\mathrm{NS}&26&278283&
\texttt{82C887F8C2C69E4F28FAD7C3DC8DE5B9099CB5A4BF431EBA1FFF3E91935978AD}.
\end{array}                                             \tag{60.1}
\]
Their line counts were respectively \(12998\) and \(5319\).
\end{audit}

YM V55 exactly certifies eight previously unobstructed mixed
third-axis Wilson charts and localizes the remaining certificate
defects to four centres.  Its failed fixed-splitting screens are
carefully separated from physical pencil negativity, and its next
move is a rational third-axis bridge.  NS V26 remains unchanged and
computes rank-one Oseen-kernel derivative norms, with only a small
second-derivative improvement.

\begin{theorem}[V60 import decision]
\label{thm:v16v60-audit-decision}
No companion theorem supplies the missing gravitational constraint
certificate of V59.  The YM distinction between failure of a
sufficient certificate and failure of the underlying operator is
imported only as a methodological guard: failure of a coarea
singular-value or gap bound must not be reported as nonexistence of a
physical constraint slice.
\end{theorem}

\begin{proof}
YM V55 concerns exact finite Wilson-pencil charts and does not define a
metric constraint, coarea Jacobian, or Einstein reaction.  NS V26
concerns rank-one Oseen tensors and has different variables and
operators.  Neither verifies any hypothesis in (59.13),
(59.23)--(59.26), or (59.35).  The stated guard is logical and imports
no numerical or operator estimate.
\end{proof}

\subsection{Finite-mode volume constraint}

Let \((X,\bar g)\) be a compact \(d\)-dimensional Riemannian manifold
of volume \(V\), and let
\[
 \phi_0=V^{-1/2},\qquad
 \int_X\phi_j\,d{\rm vol}_{\bar g}=0,\qquad
 \int_X\phi_i\phi_j\,d{\rm vol}_{\bar g}=\delta_{ij}
                                                               \tag{60.2}
\]
for \(1\leq i,j\leq N\).  Put
\[
 \sigma_a=\sum_{j=0}^Na_j\phi_j,\qquad
 g_a=e^{2\sigma_a}\bar g,                              \tag{60.3}
\]
and impose the single constraint
\[
 C_N(a)=\int_Xe^{d\sigma_a}\,d{\rm vol}_{\bar g}-V=0. \tag{60.4}
\]

\begin{proposition}[Local geometry of the volume slice]
\label{prop:v16v60-volume-geometry}
At \(a=0\),
\[
 DC_N(0)=d\sqrt V\,e_0^{\mathsf T},\qquad
 D^2C_N(0)=d^2I_{N+1}.                                \tag{60.5}
\]
Hence \(M_N=C_N^{-1}(0)\) is a regular hypersurface near zero with
\[
 T_0M_N=\operatorname{span}\{e_1,\ldots,e_N\}.         \tag{60.6}
\]
Writing it as \(a_0=\chi(a_1,\ldots,a_N)\), one has
\[
 D\chi(0)=0,\qquad
 D^2\chi(0)=-\frac d{\sqrt V}I_N.                     \tag{60.7}
\]
\end{proposition}

\begin{proof}
Differentiation of (60.4) gives
\[
 \partial_iC_N(a)=d\int_X\phi_i e^{d\sigma_a}\,d{\rm vol}_{\bar g},
\quad
 \partial_{ij}C_N(a)
 =d^2\int_X\phi_i\phi_j e^{d\sigma_a}\,d{\rm vol}_{\bar g}.       \tag{60.8}
\]
Use (60.2) at \(a=0\) to obtain (60.5).  The implicit-function
theorem gives (60.6)--(60.7), the latter by twice differentiating
\(C_N(\chi(a'),a')=0\).
\end{proof}

\subsection{Coarea Jacobian and reaction at the background}

Because (60.4) is scalar,
\[
 J_N(a)=|\nabla C_N(a)|.                              \tag{60.9}
\]

\begin{proposition}[Exact base-point Jacobian load]
\label{prop:v16v60-base-jacobian}
At the background \(a=0\),
\[
 J_N(0)=d\sqrt V,\qquad
 \nabla_{M_N}\log J_N(0)=0,                           \tag{60.10}
\]
and on tangent coordinates \(1\leq j,k\leq N\),
\[
 {\rm Hess}_{M_N}\log J_N(0)[e_j,e_k]
 =\frac{d^2}{V}\delta_{jk}.                           \tag{60.11}
\]
With the second-fundamental-form convention of (59.19),
\[
 {\rm II}_0(e_j,e_k)
 =-\frac d{\sqrt V}\delta_{jk}e_0,\qquad
 {\bf H}(0)=-\frac{Nd}{\sqrt V}e_0.                  \tag{60.12}
\]
\end{proposition}

\begin{proof}
Let \(S=|\nabla C_N|^2\), so \(\log J_N=\frac12\log S\).
At zero, (60.5) implies
\[
 S=d^2V,\qquad
 \partial_jS=0\quad(j\geq1),\qquad
 \partial_{jk}S=4d^4\delta_{jk}\quad(j,k\geq1).       \tag{60.13}
\]
Indeed, the two contributions to
\(\frac12\partial_{jk}S\) are
\(\sum_iC_{ij}C_{ik}=d^4\delta_{jk}\) and
\(\sum_iC_iC_{ijk}=d^4\delta_{jk}\), because
\[
 C_{0jk}(0)=d^3\int_X\phi_0\phi_j\phi_k\,d{\rm vol}_{\bar g}
 =\frac{d^3}{\sqrt V}\delta_{jk}.                     \tag{60.14}
\]
Thus the ambient tangent Hessian of \(\log J_N\) is
\(2d^2\delta_{jk}/V\).  Its ambient normal derivative is
\[
 D\log J_N(0)[e_0]=\frac d{\sqrt V}.                  \tag{60.15}
\]
Equation (60.7) gives the second fundamental form in (60.12).
The restriction formula (59.21) subtracts
\(d^2\delta_{jk}/V\) from the ambient tangent Hessian, proving
(60.11).  Tracing (60.12) gives the mean curvature.
\end{proof}

At the background, the V59 constraint reaction therefore becomes
\[
 {\cal R}_{C_N}(0)
 =
 \left(\partial_0U(0)-\frac{Nd}{\sqrt V}\right)e_0.   \tag{60.16}
\]
It is purely normal at that one point.  The dimension-dependent
mean-curvature term in (60.16) is not negligible as \(N\to\infty\);
it must be combined with the coarea normalization or a declared
counterterm before a continuum reaction can be discussed.

\subsection{Exact volume-preserving high-frequency paths}

The single constraint removes only the constant conformal coordinate.
For a mean-zero normalized mode \(\phi\), define
\[
 c_\phi(t)
 =-\frac1d\log\!\left(
 \frac1V\int_Xe^{dt\phi}\,d{\rm vol}_{\bar g}\right),
\qquad
 \sigma_\phi(t)=t\phi+c_\phi(t).                      \tag{60.17}
\]

\begin{lemma}[Exact constrained path]
\label{lem:v16v60-exact-path}
The metric \(g_t=e^{2\sigma_\phi(t)}\bar g\) has volume \(V\) for all
\(t\) for which (60.17) is defined, and
\[
 c_\phi(0)=c_\phi'(0)=0,\qquad
 c_\phi''(0)=-\frac dV.                               \tag{60.18}
\]
\end{lemma}

\begin{proof}
By construction,
\[
 \int_Xe^{d\sigma_\phi(t)}\,d{\rm vol}_{\bar g}
 =e^{dc_\phi(t)}\int_Xe^{dt\phi}\,d{\rm vol}_{\bar g}=V.
\]
The derivatives in (60.18) follow from
\(\int\phi=0\) and \(\int\phi^2=1\).
\end{proof}

\begin{theorem}[Constant volume does not remove conformal instability]
\label{thm:v16v60-volume-fails}
Let \(d=4\), let \(\bar g\) be flat, and let
\(-\bar\Delta\phi=\lambda\phi\) with
\(\int\phi=0\), \(\int\phi^2=1\).  For the Euclidean Einstein--Hilbert
conformal functional in the convention imported in frozen V39,
\[
 S_{\rm EH}^E[e^{2\sigma}\bar g]
 =-\frac1{16\pi G}
   \int_X6e^{2\sigma}|\bar\nabla\sigma|^2
   \,d{\rm vol}_{\bar g}
   +S_\Lambda[e^{2\sigma}\bar g],                    \tag{60.19}
\]
the exact volume-preserving path (60.17) satisfies
\[
 \left.\frac{d^2}{dt^2}\right|_{t=0}
 S_{\rm EH}^E[g_t]
 =-\frac{12}{16\pi G}\lambda.                        \tag{60.20}
\]
The cosmological contribution has zero derivative at every order.
Hence the constrained Hessian is unbounded below when the retained
mean-zero spectrum has \(\lambda\to\infty\).
\end{theorem}

\begin{proof}
The cosmological functional is constant by
Lemma~\ref{lem:v16v60-exact-path} and Proposition
\ref{prop:v16v57-volume-no-go}.  At \(t=0\),
\(\nabla\sigma_\phi(0)=0\) and
\(\partial_t\nabla\sigma_\phi(0)=\nabla\phi\), since \(c_\phi\) is
spatially constant.  Therefore
\[
 \left.\frac{d^2}{dt^2}\right|_{0}
 \int_X6e^{2\sigma_\phi(t)}
 |\nabla\sigma_\phi(t)|^2\,d{\rm vol}_{\bar g}
 =12\int_X|\nabla\phi|^2\,d{\rm vol}_{\bar g}
 =12\lambda.                                         \tag{60.21}
\]
Multiplication by the prefactor in (60.19) proves (60.20).
\end{proof}

The theorem is stronger than saying that the global constraint has
too few equations.  It exhibits exact paths inside the nonlinear
constraint surface on which the same high-frequency negative
quadratic form survives.

\subsection{Upstream selection: local unimodularity}

The next candidate constrains the determinant locally rather than its
integral:
\[
 d{\rm vol}_{g}(x)=d{\rm vol}_{\bar g}(x)
 \quad\text{at every retained cell or mode}.          \tag{60.22}
\]
Infinitesimally, (60.22) imposes
\[
 \operatorname{tr}_{\bar g}h=0,                       \tag{60.23}
\]
so it removes every pure conformal variation
\(h=2\phi\bar g\), not merely the constant one.  This avoids
Theorem~\ref{thm:v16v60-volume-fails} at the tangent-space level.
It changes the field equation: variations are trace-free, so a
continuum Einstein equation can be recovered only by deriving its
trace as an integration constant from the Bianchi identity and
stress conservation.

\begin{assessment}[What V60 closes]
\label{ass:v16v60-status}
The companion audit supplies no missing metric estimate.  V60 computes
the regular finite-mode constant-volume slice, its base coarea
Jacobian Hessian, and its full base reaction.  It then constructs
exact nonlinear volume-preserving conformal paths and proves that
their Einstein--Hilbert Hessian remains unbounded below.

The single global volume constraint is therefore rejected as the
coercivity mechanism.  Local unimodularity is selected upstream
because its tangent space removes all pure conformal modes.  Its
coarea scaling, trace-free field equation, Bianchi recovery, and
compatibility with matter anomalies remain to be proved.
\end{assessment}

\begin{decision}[V61 unimodular equation and trace recovery]
\label{dec:v16v60-next}
Formulate the local determinant constraint at finite cutoff and prove
the exact trace-free Euler--Lagrange equation.  In the smooth
continuum, combine it with the contracted Bianchi identity and
covariant stress conservation to recover the Einstein equation with
a spacetime-constant integration parameter.  Track precisely how a
trace anomaly or failed stress conservation modifies that conclusion,
and separate this classical recovery from positivity of the quantum
measure.
\end{decision}

\begin{firewall}[Claim boundary after V60]
\label{fire:v16v60}
The calculation is for finite conformal modes about a flat background.
It neither proves that local unimodularity yields a positive
gravitational measure nor controls its coarea Jacobian, ghost sector,
global slice copies, reflection positivity, or quantum trace Ward
identity.
\end{firewall}

\section*{Version V60 append-only record}
\addcontentsline{toc}{section}{Version V60 append-only record}
The complete V59 source was retained byte for byte before this
continuation.  It had \(616301\) bytes, \(18343\) lines, and SHA--256
\[
 \texttt{69F96DBB04BC90D42CA818A5FC572FDA9A88D671753F18B47E9085B21619CF8D}.
\]
V60 recorded the mandatory companion audit, computed the
constant-volume conformal slice and reaction, rejected it through an
exact high-frequency constrained path, and selected local
unimodularity for the next upstream test.


\section{V61: local unimodularity and recovery of the Einstein trace}
\label{sec:v16v61}

\subsection{A regular finite-cutoff determinant slice}

At a Euclidean cutoff, let the metric variable at each retained cell
\(y\in\mathcal Y_n\) be a positive-definite symmetric
\(d\times d\) matrix \(g_y\).  Fix positive-definite reference matrices
\(\bar g_y\) and impose
\[
 C_y(g)=\log\det g_y-\log\det\bar g_y=0
 \qquad(y\in\mathcal Y_n).                            \tag{61.1}
\]
The ambient matrix space carries the product Frobenius metric.

\begin{proposition}[Exact finite-cutoff coarea Jacobian]
\label{prop:v16v61-finite-jacobian}
The differential and gradient of (61.1) are
\[
 DC_y(g)[h]=\operatorname{tr}(g_y^{-1}h_y),
 \qquad
 \nabla C_y(g)=(0,\ldots,g_y^{-1},\ldots,0).          \tag{61.2}
\]
The constraint Gram matrix is diagonal, and hence
\[
 J_{C_n}(g)
 =\prod_{y\in\mathcal Y_n}\|g_y^{-1}\|_F.             \tag{61.3}
\]
On the constraint surface,
\[
 \sigma_{\min}(DC_n)
 =\min_y\|g_y^{-1}\|_F
 \geq
 \min_y\left[
 \sqrt d\,(\det\bar g_y)^{-1/d}\right]>0.             \tag{61.4}
\]
Thus the local determinant constraints form a regular submersion at
every positive-definite constrained metric.
\end{proposition}

\begin{proof}
Jacobi's formula gives
\[
 D\log\det g_y[h_y]=\operatorname{tr}(g_y^{-1}h_y),
\]
which is the Frobenius pairing with \(g_y^{-1}\).  Distinct
constraints act on disjoint matrix blocks, so
\[
 DC_nDC_n^{\mathsf T}
 =\operatorname{diag}
 \bigl(\|g_y^{-1}\|_F^2:y\in\mathcal Y_n\bigr),
\]
proving (61.3) and the equality in (61.4).  If
\(\lambda_1,\ldots,\lambda_d\) are the eigenvalues of \(g_y\), the
arithmetic--geometric mean inequality gives
\[
 \|g_y^{-1}\|_F^2
 =\sum_{i=1}^d\lambda_i^{-2}
 \geq d\left(\prod_{i=1}^d\lambda_i^{-2}\right)^{1/d}
 =d(\det\bar g_y)^{-2/d}.                             \tag{61.5}
\]
\end{proof}

Unlike the single volume constraint of V60, this slice has one
constraint at every retained cell.  Its tangent space is
\[
 T_gM_n
 =\left\{h:\operatorname{tr}(g_y^{-1}h_y)=0
 \text{ for every }y\right\}.                         \tag{61.6}
\]
Every local pure conformal vector \(h_y=2\phi_yg_y\) has trace
\(2d\phi_y\), and therefore belongs to (61.6) only when
\(\phi_y=0\) at every cell.

The lower singular-value certificate (61.4) is automatic, but the
upper V59 derivative loads are not.  Fixed determinant permits
arbitrarily anisotropic matrices, for which
\(\|g_y^{-1}\|_F\), \(D^2C_y\), and \(D^3C_y\) can diverge.  A compact
shape bound or a separate anisotropy Lyapunov estimate is still
needed for a uniform coarea gap.

\subsection{Finite-cutoff trace-free critical equation}

Let \(S_n(g)\) be differentiable and write its Frobenius gradient at
cell \(y\) as \(E_y(g)\):
\[
 DS_n(g)[h]=\sum_y\langle E_y(g),h_y\rangle_F.         \tag{61.7}
\]

\begin{theorem}[Constrained multiplier equation]
\label{thm:v16v61-finite-euler}
A constrained critical point of \(S_n\) on (61.1) satisfies
\[
 E_y(g)=p_y g_y^{-1}\qquad(y\in\mathcal Y_n)           \tag{61.8}
\]
for scalars \(p_y\).  Equivalently,
\[
 E_y-\frac{\operatorname{tr}(g_yE_y)}d\,g_y^{-1}=0.   \tag{61.9}
\]
\end{theorem}

\begin{proof}
By (61.6), the derivative (61.7) vanishes on the kernel of the
surjective constraint differential.  Its annihilator is the span of
the mutually orthogonal constraint normals \(g_y^{-1}\), which gives
(61.8).  Contracting (61.8) with \(g_y\) gives
\(p_y=\operatorname{tr}(g_yE_y)/d\), proving (61.9).
\end{proof}

At finite cutoff the multipliers \(p_y\) are cell dependent.  Their
constancy is not a consequence of the algebraic constraint.  It can
emerge only from a continuum differential identity.

\subsection{Continuum trace-free equation}

Let \(M\) now be a connected smooth \(d\)-manifold with \(d>2\).
Assume that the renormalized metric variation has the form
\[
 \delta S[g;h]
 =\frac12\int_M
 \bigl(G^{\mu\nu}-\kappa T^{\mu\nu}\bigr)h_{\mu\nu}
 \,d{\rm vol}_g                                      \tag{61.10}
\]
for every compactly supported smooth \(h\), where
\[
 G_{\mu\nu}=R_{\mu\nu}-\frac12Rg_{\mu\nu}.            \tag{61.11}
\]
The cosmological variation is pure trace and therefore disappears
from determinant-preserving variations.

\begin{theorem}[Trace-free Euler--Lagrange equation]
\label{thm:v16v61-tracefree}
If (61.10) vanishes for every compactly supported variation satisfying
\(g^{\mu\nu}h_{\mu\nu}=0\), then
\[
 R_{\mu\nu}-\frac1dRg_{\mu\nu}
 =
 \kappa\left(T_{\mu\nu}-\frac1dTg_{\mu\nu}\right).
                                                               \tag{61.12}
\]
\end{theorem}

\begin{proof}
Set \(E^{\mu\nu}=G^{\mu\nu}-\kappa T^{\mu\nu}\).  Every compactly
supported trace-free tensor is an admissible test, so the fundamental
lemma of the calculus of variations gives
\[
 E_{\mu\nu}-\frac1dE^\alpha{}_\alpha g_{\mu\nu}=0.    \tag{61.13}
\]
The trace-free part of \(G_{\mu\nu}\) equals the trace-free part of
\(R_{\mu\nu}\), and (61.13) is exactly (61.12).
\end{proof}

\subsection{Bianchi recovery and its failure mode}

Define the stress nonconservation current
\[
 J_\nu:=\nabla^\mu T_{\mu\nu}.                        \tag{61.14}
\]

\begin{theorem}[Integration parameter and reaction current]
\label{thm:v16v61-bianchi}
Every smooth solution of (61.12) obeys
\[
 \partial_\nu\left(
 \frac{d-2}{2}R+\kappa T\right)
 =d\kappa J_\nu.                                      \tag{61.15}
\]
If \(J_\nu=0\), connectedness implies the existence of a spacetime
constant \(\Lambda_{\rm int}\) such that
\[
 \frac{d-2}{2}R+\kappa T=d\Lambda_{\rm int},          \tag{61.16}
\]
and (61.12) is then equivalent to
\[
 G_{\mu\nu}+\Lambda_{\rm int}g_{\mu\nu}
 =\kappa T_{\mu\nu}.                                  \tag{61.17}
\]
If \(J\neq0\), a scalar \(\Lambda_{\rm eff}\) satisfying the full
equation can exist only when
\[
 \partial_\nu\Lambda_{\rm eff}=\kappa J_\nu.          \tag{61.18}
\]
In particular \(J\) must be an exact one-form, and the integration
parameter is not constant.
\end{theorem}

\begin{proof}
Take the divergence of (61.12) and use the contracted Bianchi identity
\(\nabla^\mu R_{\mu\nu}=\frac12\partial_\nu R\).  This gives
\[
 \left(\frac12-\frac1d\right)\partial_\nu R
 =\kappa\left(J_\nu-\frac1d\partial_\nu T\right).
\]
Multiplication by \(d\) and rearrangement proves (61.15).  If \(J=0\),
the scalar on the left is constant on connected \(M\), giving
(61.16).  Subtract \(\kappa T_{\mu\nu}\) from (61.12) and use (61.16);
the result is (61.17).  Conversely, the trace-free part of (61.17)
is (61.12).  Taking the divergence of a proposed equation with
\(\Lambda_{\rm eff}\) and using \(\nabla^\mu G_{\mu\nu}=0\) proves
(61.18).
\end{proof}

Equation (61.15) is the continuum counterpart of the finite
multipliers \(p_y\): conservation is the compatibility condition that
collapses their continuum trace reaction to one constant.

\subsection{Trace anomaly ledger}

Suppose a quantum matter sector has a renormalized trace
\[
 T=T_{\rm cl}+{\cal A}[g,A,\psi],                     \tag{61.19}
\]
where \({\cal A}\) is the trace anomaly.  If the renormalized stress
tensor is still covariantly conserved, then
\[
 \frac{d-2}{2}R+\kappa(T_{\rm cl}+{\cal A})
 =d\Lambda_{\rm int}.                                 \tag{61.20}
\]
Thus a trace anomaly changes the scalar curvature relation but does
not by itself turn \(\Lambda_{\rm int}\) into a spacetime-dependent
field.  If the regularization also produces a diffeomorphism anomaly
or an uncancelled cutoff defect \(J_\nu\), then (61.15), rather than
(61.20), is authoritative.

\begin{corollary}[Ward-defect stability]
\label{cor:v16v61-ward-defect}
Let a cutoff sequence satisfy the trace-free equation with currents
\(J^{(n)}\).  If
\[
 J^{(n)}\longrightarrow0
\quad\text{in }\mathcal D'(T^*M),                     \tag{61.21}
\]
and
\[
 F_n:=\frac{d-2}{2}R_n+\kappa T_n
\longrightarrow F
\quad\text{in }\mathcal D'(M),                        \tag{61.22}
\]
then \(\partial_\nu F=0\) distributionally.  On each connected
component, \(F\) is a constant distribution.
\end{corollary}

\begin{proof}
Equation (61.15) gives
\(\partial_\nu F_n=d\kappa J_\nu^{(n)}\).  Pass to the distributional
limit using (61.21)--(61.22).  A distribution with zero derivative on
a connected manifold is constant.
\end{proof}

This is the correct weak target for the continuum programme: the
cutoff conservation defect must vanish strongly enough to pass the
derivative identity, while the curvature--trace combination must be
tight as a distribution.  Pointwise convergence of either term is
unnecessary.

\subsection{What unimodularity does and does not solve}

The local determinant constraint removes the pure conformal tangent
direction and retains the trace-free Einstein equation.  Conservation
then reconstructs the missing trace with an integration constant.
The following independent items remain:
\begin{enumerate}
\item control of matrix anisotropy so that the V59 coarea Jacobian
      loads are uniform;
\item a diffeomorphism gauge slice and its determinant;
\item positivity and reflection transfer for the combined constrained
      law;
\item the distributional Ward-defect convergence (61.21);
\item identification of all anomaly contributions in (61.19).
\end{enumerate}

\begin{assessment}[What V61 closes]
\label{ass:v16v61-status}
V61 proves that the cellwise determinant constraint is a regular
finite-cutoff submersion with an exact positive coarea Jacobian and an
automatic lower singular-value bound.  Its constrained critical
equation is precisely trace-free.  In the smooth continuum, Bianchi
plus stress conservation recovers the full Einstein equation with a
constant integration parameter; a trace anomaly modifies the scalar
relation, while a conservation defect produces the exact source
\(d\kappa J\).

This resolves the classical equation-level objection to removing the
conformal mode.  It does not construct a positive quantum measure:
anisotropy, diffeomorphism gauge directions, the gauge determinant,
reflection positivity, and the cutoff Ward limit remain open.
\end{assessment}

\begin{decision}[V62 linearized unimodular gauge complex]
\label{dec:v16v61-next}
Combine local unimodularity with a diffeomorphism gauge at the
linearized flat-background level.  Prove the finite-mode orthogonal
decomposition into transverse-traceless, gauge, and constrained scalar
sectors; compute the Faddeev--Popov symbol and its zero modes; and
identify the exact positive quadratic form on the physical sector.
Then state which bounds must survive nonlinear reconstruction before
the V59 positive-law compiler can be invoked.
\end{decision}

\begin{firewall}[Claim boundary after V61]
\label{fire:v16v61}
The Bianchi recovery is a classical or distributional implication
under an assumed trace-free equation and stress Ward identity.  It is
not a construction of the cutoff measure, a proof of anomaly
cancellation, or a derivation of the Standard Model--Einstein
continuum.
\end{firewall}

\section*{Version V61 append-only record}
\addcontentsline{toc}{section}{Version V61 append-only record}
The complete V60 source was retained byte for byte before this
continuation.  It had \(627340\) bytes, \(18654\) lines, and SHA--256
\[
 \texttt{ABEB4CB06AA8212A9F4FA398E40671B0228337C63A014A3291CD4946AD3276D5}.
\]
V61 proved regularity of the finite local-determinant slice, derived
its trace-free Euler--Lagrange equation, recovered the Einstein trace
from Bianchi and stress conservation, and isolated the exact
distributional Ward-defect target.


\section{V62: the linearized unimodular gauge complex}
\label{sec:v16v62}

\subsection{Modewise orthogonal decomposition}

Work on a flat \(d\)-torus with \(d>2\), and fix a nonzero Fourier
momentum \(k\).  Write \(n=k/|k|\) and
\[
 P_{\mu\nu}=\delta_{\mu\nu}-n_\mu n_\nu.              \tag{62.1}
\]
The linearized local determinant constraint is
\[
 \operatorname{tr}h=0.                               \tag{62.2}
\]
Let \({\rm Sym}_0^2(\mathbb R^d)\) denote the real trace-free symmetric
tensors at this momentum.

\begin{theorem}[Unimodular York decomposition at one mode]
\label{thm:v16v62-decomposition}
Every \(h\in{\rm Sym}_0^2(\mathbb R^d)\) has a unique orthogonal
decomposition
\[
 h=t+n\odot v+sS,\qquad
 S=n\otimes n-\frac1{d-1}P,                           \tag{62.3}
\]
where
\[
 n^\mu t_{\mu\nu}=0,\quad \operatorname{tr}t=0,\quad
 n\cdot v=0,\quad
 (n\odot v)_{\mu\nu}=n_\mu v_\nu+v_\mu n_\nu.         \tag{62.4}
\]
The coefficients are
\[
 s=n^\mu n^\nu h_{\mu\nu},\qquad
 v=P(hn),\qquad
 t=h-n\odot v-sS.                                    \tag{62.5}
\]
\end{theorem}

\begin{proof}
The tensors \(t\), \(n\odot v\), and \(S\) are pairwise orthogonal:
transversality kills the contractions with \(t\), while
\(S(n,v)=0\).  Moreover
\[
 |n\odot v|^2=2|v|^2,\qquad
 |S|^2=\frac d{d-1}.                                  \tag{62.6}
\]
For a given \(h\), (62.5) produces a transverse \(v\), and the
remaining \(t\) is transverse because
\[
 hn=v+sn,\qquad
 (n\odot v)n=v,\qquad Sn=n.
\]
It is trace-free because all three other terms in (62.3) are
trace-free.  Orthogonality proves uniqueness.
\end{proof}

The three summands are respectively transverse-traceless, residual
volume-preserving gauge, and the constrained scalar.  The scalar is
not a pure conformal tensor.

\subsection{Residual gauge and its Faddeev--Popov symbol}

A linearized diffeomorphism acts by
\[
 (\delta_\xi h)_{\mu\nu}
 =i(k_\mu\xi_\nu+k_\nu\xi_\mu).                       \tag{62.7}
\]
It preserves (62.2) exactly when \(k\cdot\xi=0\).  Thus the residual
gauge parameter has \(d-1\) transverse components.  Define the
transverse divergence gauge
\[
 F_\nu(h)=P_\nu{}^\rho k^\mu h_{\mu\rho}=0.           \tag{62.8}
\]

\begin{proposition}[Residual Faddeev--Popov determinant]
\label{prop:v16v62-fp}
On \(k^\perp\),
\[
 D_\xi F(\delta_\xi h)=i|k|^2\xi,\qquad
 |\det D_\xi F|=|k|^{2(d-1)}.                         \tag{62.9}
\]
The gauge condition removes exactly the \(n\odot v\) summand in
(62.3).  Its only Fourier zero modes occur at \(k=0\).
\end{proposition}

\begin{proof}
Contract (62.7) with \(k^\mu\) and project:
\[
 P_\nu{}^\rho k^\mu(\delta_\xi h)_{\mu\rho}
 =iP_\nu{}^\rho\bigl(|k|^2\xi_\rho+k_\rho(k\cdot\xi)\bigr)
 =i|k|^2\xi_\nu.
\]
This proves (62.9).  By (62.3),
\[
 F(h)=|k|P(hn)=|k|v,
\]
so \(F=0\) if and only if \(v=0\).
\end{proof}

The determinant in (62.9) is strictly positive in magnitude for each
nonzero mode.  This is a linearized local statement; it does not
settle the sign of a nonlinear gauge determinant across Gribov
horizons.

\subsection{Quadratic Einstein form on the slice}

In the positive physical normalization \(c_G>0\), the flat-background
Euclidean Einstein--Hilbert quadratic form restricted to
\(\operatorname{tr}h=0\) can be written modewise as
\[
 Q_k(h)
 =c_G|k|^2\left(
 \frac12|h|^2-|n^\mu h_{\mu\cdot}|^2\right).          \tag{62.10}
\]
This is the trace-free specialization of the Fierz--Pauli form.  The
normalization is left explicit because earlier action conventions
differ by an overall positive factor.

\begin{theorem}[Exact physical-sector floor]
\label{thm:v16v62-floor}
For the decomposition (62.3),
\[
 Q_k(h)
 =c_G|k|^2\left[
 \frac12|t|^2+\frac{d-2}{2(d-1)}s^2\right].           \tag{62.11}
\]
The residual gauge summand has zero quadratic form.  On the gauge
slice \(F=0\),
\[
 Q_k(h)\geq
 c_G\frac{d-2}{2d}|k|^2|h|^2.                        \tag{62.12}
\]
\end{theorem}

\begin{proof}
Orthogonality and (62.6) give
\[
 |h|^2=|t|^2+2|v|^2+\frac d{d-1}s^2.                 \tag{62.13}
\]
Also
\[
 hn=v+sn,\qquad |hn|^2=|v|^2+s^2.                   \tag{62.14}
\]
Substitute (62.13)--(62.14) into (62.10); the \(v\)-terms cancel and
(62.11) follows.  On \(F=0\), \(v=0\).  The TT coefficient relative
to its norm is \(c_G|k|^2/2\), while the scalar coefficient relative
to \(|sS|^2=ds^2/(d-1)\) is
\(c_G|k|^2(d-2)/(2d)\).  The latter is the smaller one, proving
(62.12).
\end{proof}

The pure conformal instability of V39 and V60 is absent because
\(h=2\phi\delta\) violates (62.2).  Its removal does not require a
vanishing fourth-order stabilizer.  The surviving scalar \(sS\) has a
strictly positive coefficient for \(d>2\).

\subsection{Infrared zero modes and shape anisotropy}

At \(k=0\), both the Faddeev--Popov symbol (62.9) and the quadratic
form (62.10) vanish.  Constant trace-free tensors change the flat
torus shape while preserving its volume.  They are the linear
counterpart of the uncontrolled fixed-determinant anisotropy noted
after (61.6).

\begin{proposition}[Nonzero-mode cutoff floor]
\label{prop:v16v62-cutoff-floor}
Suppose the Fourier cutoff excludes \(k=0\) and the torus has
\[
 |k|^2\geq\lambda_{\rm IR}>0                          \tag{62.15}
\]
on every retained mode.  Then on the unimodular residual-gauge slice,
\[
 Q(h)=\sum_{k\ne0}Q_k(h_k)
 \geq c_G\frac{d-2}{2d}\lambda_{\rm IR}
       \sum_{k\ne0}|h_k|^2.                           \tag{62.16}
\]
\end{proposition}

\begin{proof}
Sum (62.12) and use (62.15).
\end{proof}

There is no corresponding conclusion when the constant shape modes
are integrated over an unbounded fundamental domain.  A valid
finite-volume construction must fix them, quotient them by the
mapping-class group with a controlled domain, or supply a separate
shape potential whose continuum role is declared.

\subsection{Nonlinear certificate required by V59}

The linear calculation identifies the exact quantities that a
nonlinear shell construction must bound.  Let
\({\cal F}_n(g)\) be a residual volume-preserving gauge map on the
local determinant slice.  A sufficient finite-cutoff certificate
consists of:
\begin{align}
 &\sigma_{\min}\bigl(D_\xi{\cal F}_n(g)\bigr)
   \geq c_{\rm FP}>0
   \quad\text{off declared zero modes},               \tag{62.17}\\
 &\left|\det D_\xi{\cal F}_n(g)\right|>0
   \quad\text{with a fixed sign on each chart},        \tag{62.18}\\
 &{\rm Hess}_{\rm slice}S_n
   +{\rm Hess}_{\rm slice}\log J_{C_n}
   -{\rm Hess}_{\rm slice}\log|\det D_\xi{\cal F}_n|
   \succeq\rho_{\rm phys}I,                            \tag{62.19}\\
 &\|\nabla_{\rm slice}I_{i,n}\|
   \leq L_{1,i,n}|z|+L_{0,i,n},                        \tag{62.20}
\end{align}
with the V57 margin positive uniformly in the shell and cutoff.
The sign in (62.19) follows because the positive gauge determinant
multiplies the density and therefore subtracts
\(\log|\det D_\xi{\cal F}_n|\) from the effective potential.

At flat background, (62.17)--(62.19) reduce modewise to
(62.9), (62.12), and removal of \(k=0\).  Away from that point they
are new analytic hypotheses.  In particular, positivity at one
background does not exclude a Gribov horizon.

\begin{assessment}[What V62 closes]
\label{ass:v16v62-status}
V62 gives the exact finite-mode decomposition of a trace-free metric
perturbation into TT, residual gauge, and constrained scalar sectors.
The residual Faddeev--Popov symbol is \(|k|^2\) on each transverse
gauge component, and the Einstein quadratic form is strictly positive
on the nonzero-mode gauge slice, including the constrained scalar.

The remaining linear obstruction is the constant volume-preserving
shape sector.  Nonlinearly, uniform Faddeev--Popov invertibility,
determinant sign, coarea loads, and a physical Hessian or radial
margin have not been proved.
\end{assessment}

\begin{decision}[V63 local nonlinear gauge chart]
\label{dec:v16v62-next}
Use a quantitative inverse-function theorem to construct a local
residual gauge chart around a background satisfying the spectral
floor (62.17).  Derive an explicit radius that preserves
Faddeev--Popov invertibility and determinant sign, and bound the first
two derivatives of its logarithmic determinant.  Feed those bounds
into the V59 Jacobian and V57 slope--force budgets without asserting a
global gauge slice.
\end{decision}

\begin{firewall}[Claim boundary after V62]
\label{fire:v16v62}
The result is linearized about a flat torus and omits the zero modes.
It does not prove a global nonlinear gauge fixing, absence of Gribov
copies, reflection positivity, nonlinear Einstein coercivity, or the
Standard Model--Einstein continuum.
\end{firewall}

\section*{Version V62 append-only record}
\addcontentsline{toc}{section}{Version V62 append-only record}
The complete V61 source was retained byte for byte before this
continuation.  It had \(638881\) bytes, \(18976\) lines, and SHA--256
\[
 \texttt{E598DB939016355AE7E61EE3797C7BDAD8522755A1879B765F2AD705017B99B3}.
\]
V62 proved the linearized unimodular York decomposition, computed the
residual Faddeev--Popov symbol and zero modes, and established the
strict nonzero-mode Einstein quadratic floor on the constrained gauge
slice.


\section{V63: a quantitative local gauge chart and determinant loads}
\label{sec:v16v63}

\subsection{A contraction-radius inverse theorem}

Let \(X\subset\mathbb R^m\) be a parameter domain for metrics already
restricted by the local determinant constraint.  Let
\[
 \Psi:X\times\mathbb R^p\longrightarrow\mathbb R^p               \tag{63.1}
\]
be a \(C^2\) residual gauge map, where \(\xi\in\mathbb R^p\) is the
volume-preserving gauge parameter.  Fix
\[
 \Psi(x_0,0)=0,\qquad
 M_0:=D_\xi\Psi(x_0,0),\qquad
 \|M_0^{-1}\|_{\rm op}\leq\beta.                     \tag{63.2}
\]

\begin{theorem}[Explicit local gauge radius]
\label{thm:v16v63-radius}
Fix \(q\in(0,1)\) and \(r>0\).  Suppose that on
\[
 {\cal D}_{\delta,r}
 =\{(x,\xi):|x-x_0|\leq\delta,\ |\xi|\leq r\}         \tag{63.3}
\]
one has
\[
 \|D_\xi\Psi(x,\xi)-M_0\|_{\rm op}\leq\frac q\beta,
 \qquad
 \|D_x\Psi(x,0)\|_{\rm op}\leq K_x,                  \tag{63.4}
\]
and
\[
 \delta\leq\frac{(1-q)r}{\beta K_x}.                 \tag{63.5}
\]
Then every \(x\) with \(|x-x_0|\leq\delta\) has a unique
\(\xi_*(x)\in\overline B_r(0)\) satisfying
\[
 \Psi(x,\xi_*(x))=0.                                  \tag{63.6}
\]
Moreover
\[
 |\xi_*(x)|\leq\frac{\beta K_x}{1-q}|x-x_0|,          \tag{63.7}
\]
\[
 \|D_\xi\Psi(x,\xi_*(x))^{-1}\|_{\rm op}
 \leq B:=\frac{\beta}{1-q},                           \tag{63.8}
\]
and, wherever the derivative is evaluated,
\[
 D\xi_*(x)
 =-D_\xi\Psi(x,\xi_*(x))^{-1}
   D_x\Psi(x,\xi_*(x)).                               \tag{63.9}
\]
\end{theorem}

\begin{proof}
For fixed \(x\), define
\[
 T_x(\xi)=\xi-M_0^{-1}\Psi(x,\xi).                    \tag{63.10}
\]
The first estimate in (63.4) gives
\[
 \|D_\xi T_x\|_{\rm op}
 \leq\|M_0^{-1}\|_{\rm op}
     \|M_0-D_\xi\Psi\|_{\rm op}\leq q.                \tag{63.11}
\]
Since \(\Psi(x_0,0)=0\), the second estimate in (63.4) gives
\[
 |T_x(0)|
 =|M_0^{-1}\Psi(x,0)|
 \leq\beta K_x|x-x_0|
 \leq(1-q)r.                                         \tag{63.12}
\]
Thus \(T_x\) maps \(\overline B_r\) to itself and is a contraction.
Its unique fixed point satisfies (63.6).  Applying (63.11) to the
fixed-point equation gives
\[
 |\xi_*|\leq |T_x(0)|+q|\xi_*|,
\]
which proves (63.7).  Finally,
\[
 D_\xi\Psi
 =M_0\left[I+M_0^{-1}(D_\xi\Psi-M_0)\right],
\]
and the bracket has inverse norm at most \((1-q)^{-1}\) by the
Neumann series.  This proves (63.8).  Differentiating (63.6) proves
(63.9).
\end{proof}

\subsection{Orientation and determinant bounds}

Put
\[
 M(x)=D_\xi\Psi(x,\xi_*(x)).                           \tag{63.13}
\]

\begin{corollary}[No local Gribov horizon]
\label{cor:v16v63-sign}
On the connected ball \(|x-x_0|\leq\delta\),
\[
 \det M(x)\neq0,\qquad
 \operatorname{sign}\det M(x)
 =\operatorname{sign}\det M_0.                       \tag{63.14}
\]
Also
\[
 \sigma_{\min}(M(x))\geq B^{-1}
 =\frac{1-q}{\beta},\qquad
 |\det M(x)|\geq B^{-p}.                              \tag{63.15}
\]
\end{corollary}

\begin{proof}
Invertibility follows from (63.8).  A continuous nonzero determinant
has constant sign on a connected set.  The least singular-value bound
is the reciprocal of the inverse norm, and the determinant is the
product of the \(p\) singular values.
\end{proof}

This excludes a gauge horizon only inside the certified chart.  It
does not exclude a second gauge parameter outside \(\overline B_r\)
that reaches the same gauge slice.

\subsection{First and second logarithmic determinant loads}

Assume \(M\) is \(C^2\) on the chart and, for unit directions \(v,w\)
in \(X\),
\[
 \|DM(x)[v]\|_{\rm op}\leq K_1,\qquad
 \|D^2M(x)[v,w]\|_{\rm op}\leq K_2.                  \tag{63.16}
\]

\begin{proposition}[Faddeev--Popov log loads]
\label{prop:v16v63-logdet}
On the certified chart,
\[
 |D\log|\det M|[v]|
 \leq G_{\rm FP}:=pBK_1,                              \tag{63.17}
\]
and
\[
 |D^2\log|\det M|[v,w]|
 \leq H_{\rm FP}:=p(BK_2+B^2K_1^2).                 \tag{63.18}
\]
\end{proposition}

\begin{proof}
The fixed sign in (63.14) permits differentiation of
\(\log|\det M|\).  Jacobi's formula gives
\[
 D\log|\det M|[v]
 =\operatorname{tr}(M^{-1}DM[v]).                    \tag{63.19}
\]
Use \(|\operatorname{tr}A|\leq p\|A\|_{\rm op}\),
(63.8), and (63.16) to obtain (63.17).  A second derivative gives
\[
 D^2\log|\det M|[v,w]
 =\operatorname{tr}\!\left(
 M^{-1}D^2M[v,w]
 -M^{-1}DM[w]M^{-1}DM[v]\right).                    \tag{63.20}
\]
The same norm and trace bounds prove (63.18).
\end{proof}

The constants \(K_1,K_2\) are composite chart derivatives.  They can
be computed from derivatives of \(\Psi\) using (63.9).  For example,
if
\[
 \|D_x\Psi\|_{\rm op}\leq A_x,                        \tag{63.21}
\]
then
\[
 \|D\xi_*\|_{\rm op}\leq BA_x.                        \tag{63.22}
\]
Consequently, if
\[
 \|D_xD_\xi\Psi\|_{\rm op}\leq A_{x\xi},\qquad
 \|D_\xi^2\Psi\|_{\rm op}\leq A_{\xi\xi},             \tag{63.23}
\]
then the chain rule gives the checkable first load
\[
 K_1\leq A_{x\xi}+A_{\xi\xi}BA_x.                    \tag{63.24}
\]
An analogous bound for \(K_2\) uses third derivatives of \(\Psi\)
and the twice-differentiated equation (63.6).  It must be supplied by
the reconstruction atlas rather than inferred from (63.4).

\subsection{Insertion into the positive constrained law}

Within one gauge chart, suppose the gauge-orbit change of variables is
valid and the determinant orientation is fixed.  The constrained
gauge-fixed density has the formal local form
\[
 d\nu_{\rm gf}(x)
 =Z^{-1}e^{-U(x)}
 J_{C_n}(x)^{-1}|\det M(x)|\,d{\rm vol}_{\rm slice}(x).
                                                               \tag{63.25}
\]
Its effective potential is
\[
 \Phi_{\rm gf}
 =U+\log J_{C_n}-\log|\det M|.                        \tag{63.26}
\]

\begin{corollary}[Gap and radial ledger]
\label{cor:v16v63-ledger}
The Faddeev--Popov contribution in (63.26) has gradient load
\[
 \|\nabla\{-\log|\det M|\}\|\leq G_{\rm FP},          \tag{63.27}
\]
so it enters the V57 radial compiler as
\[
 L^{\rm FP}_1=0,\qquad L^{\rm FP}_0=G_{\rm FP}.       \tag{63.28}
\]
Its worst lower-Hessian contribution is
\[
 {\rm Hess}\{-\log|\det M|\}\succeq-H_{\rm FP}I.      \tag{63.29}
\]
Therefore a Bakry--\'Emery floor \(\rho_{\rm bare}\) before the gauge
determinant remains positive if
\[
 \rho_{\rm bare}>H_{\rm FP},                          \tag{63.30}
\]
while a radial floor uses the Young loss
\[
 G_{\rm FP}|x|
 \leq\tau_{\rm FP}|x|^2
 +\frac{G_{\rm FP}^2}{4\tau_{\rm FP}}.               \tag{63.31}
\]
\end{corollary}

\begin{proof}
Equations (63.27) and (63.29) are the sign-reversed absolute bounds
from Proposition~\ref{prop:v16v63-logdet}.  The remaining statements
are the V59 Hessian ledger and V57 slope--force compiler.
\end{proof}

\subsection{Flat-background specialization and infrared collapse}

For V62's flat residual gauge map, the nonzero-mode reference operator
is
\[
 M_0(k)=|k|^2I_{k^\perp}.                             \tag{63.32}
\]
If \(|k|^2\geq\lambda_{\rm IR}\), then
\[
 \beta=\lambda_{\rm IR}^{-1}.                         \tag{63.33}
\]
Thus (63.4) requires
\[
 \|D_\xi\Psi-M_0\|_{\rm op}
 \leq q\lambda_{\rm IR},                              \tag{63.34}
\]
and the metric radius allowed by (63.5) scales at most like
\[
 \delta\leq
 \frac{(1-q)r\lambda_{\rm IR}}{K_x}.                 \tag{63.35}
\]

\begin{proposition}[Thermodynamic-radius obstruction]
\label{prop:v16v63-ir}
If a sequence of tori has
\(\lambda_{\rm IR}^{(n)}\to0\) while \(r,K_x,q\) stay fixed, the
certified radius in (63.35) tends to zero.  Hence this particular
global Fourier-chart proof is not uniform in the thermodynamic limit.
\end{proposition}

\begin{proof}
The right-hand side of (63.35) is proportional to
\(\lambda_{\rm IR}^{(n)}\).
\end{proof}

This is a limitation of the certificate, not a proof that local gauge
charts disappear at large volume.  A volume-uniform route must use
local elliptic estimates, weighted spaces, or an infrared block
treated separately from the ultraviolet gauge inverse.

\begin{assessment}[What V63 closes]
\label{ass:v16v63-status}
V63 constructs an explicit nonlinear residual-gauge chart under a
checkable perturbation bound, preserves Faddeev--Popov invertibility
and orientation, and derives the exact first and second logarithmic
determinant loads entering the V57 and V59 compilers.

At flat background the certificate depends on
\(\lambda_{\rm IR}^{-1}\); its admissible metric radius collapses as
the volume grows.  The result is therefore a valid local finite-volume
chart, not a volume-uniform gauge construction or a global exclusion
of Gribov copies.
\end{assessment}

\begin{decision}[V64 infrared/ultraviolet gauge split]
\label{dec:v16v63-next}
Separate the residual gauge space at a fixed physical threshold
\(\lambda_*>0\).  Apply V63 only to the ultraviolet inverse, and treat
the finite infrared block by an explicit compact chart or pinned
coordinates.  Prove a Schur-complement invertibility theorem with
quantified mixed blocks and determine whether its constants can remain
uniform as the volume grows.
\end{decision}

\begin{firewall}[Claim boundary after V63]
\label{fire:v16v63}
The local chart assumes derivative bounds for the chosen nonlinear
gauge map and a valid orbit change of variables.  It does not prove
global uniqueness, a volume-uniform slice, determinant reflection
positivity, or the full continuum measure.
\end{firewall}

\section*{Version V63 append-only record}
\addcontentsline{toc}{section}{Version V63 append-only record}
The complete V62 source was retained byte for byte before this
continuation.  It had \(648057\) bytes, \(19241\) lines, and SHA--256
\[
 \texttt{D08207459BA96D60E74F82B0F565F64E509F153FCCE041E7344B351626D09D48}.
\]
V63 proved a quantitative local residual-gauge chart, controlled its
Faddeev--Popov determinant and logarithmic loads, and isolated the
infrared collapse of the global Fourier inverse-function radius.


\section{V64: infrared--ultraviolet Schur control and determinant extensivity}
\label{sec:v16v64}

\subsection{A quantified block inverse}

Let the residual gauge space split as
\[
 {\cal G}_n={\cal G}_{<,n}\oplus{\cal G}_{>,n},       \tag{64.1}
\]
where the second summand contains modes above a fixed physical
threshold.  In this decomposition write the Faddeev--Popov derivative
as
\[
 M_n=
 \begin{pmatrix}
 A_n&B_n\\
 C_n&D_n
 \end{pmatrix}.                                      \tag{64.2}
\]

\begin{theorem}[Infrared--ultraviolet Schur certificate]
\label{thm:v16v64-schur}
Assume
\[
 \|A_n^{-1}\|\leq\alpha,\qquad
 \|D_n^{-1}\|\leq\beta,\qquad
 \|B_n\|\leq b,\qquad
 \|C_n\|\leq c,                                      \tag{64.3}
\]
and
\[
 \theta:=\alpha b\beta c<1.                           \tag{64.4}
\]
Then \(M_n\) is invertible.  Its Schur complement
\[
 S_n=A_n-B_nD_n^{-1}C_n                               \tag{64.5}
\]
satisfies
\[
 \|S_n^{-1}\|\leq\frac{\alpha}{1-\theta}.             \tag{64.6}
\]
The inverse blocks obey
\begin{align}
 \|(M_n^{-1})_{<<}\|
 &\leq\frac{\alpha}{1-\theta}, \nonumber\\
 \|(M_n^{-1})_{<>}\|
 &\leq\frac{\alpha b\beta}{1-\theta}, \nonumber\\
 \|(M_n^{-1})_{><}\|
 &\leq\frac{\beta c\alpha}{1-\theta},                 \tag{64.7}\\
 \|(M_n^{-1})_{>>}\|
 &\leq\frac{\beta}{1-\theta}.
\end{align}
If the blocks are real and their determinant orientations are fixed,
then
\[
 \operatorname{sign}\det M_n
 =\operatorname{sign}\det A_n\,
  \operatorname{sign}\det D_n.                       \tag{64.8}
\]
\end{theorem}

\begin{proof}
Factor
\[
 S_n=A_n\left(I-A_n^{-1}B_nD_n^{-1}C_n\right).       \tag{64.9}
\]
The operator subtracted from the identity has norm at most
\(\theta<1\).  The Neumann series proves (64.6).  The standard block
inverse identity is
\[
 M_n^{-1}=
 \begin{pmatrix}
 S_n^{-1}&-S_n^{-1}B_nD_n^{-1}\\
 -D_n^{-1}C_nS_n^{-1}
 &D_n^{-1}+D_n^{-1}C_nS_n^{-1}B_nD_n^{-1}
 \end{pmatrix}.                                      \tag{64.10}
\]
The first three bounds in (64.7) follow immediately.  The final block
is bounded by
\[
 \beta+\frac{\beta c\alpha b\beta}{1-\theta}
 =\frac{\beta}{1-\theta}.
\]
Finally \(\det M_n=\det D_n\det S_n\).  The path
\(A_n-tB_nD_n^{-1}C_n\), \(0\leq t\leq1\), stays invertible by
(64.4), so \(\det S_n\) has the sign of \(\det A_n\), proving
(64.8).
\end{proof}

For the flat ultraviolet symbol \(D_n(k)=|k|^2I\) on
\(|k|^2\geq\lambda_*\), one can take
\[
 \beta=\lambda_*^{-1},                                \tag{64.11}
\]
independently of the total volume.  Thus the global lowest eigenvalue
has been removed from the ultraviolet inverse.  Uniformity now rests
on the infrared inverse \(\alpha\) and the mixed-product margin
\(\theta\).

\subsection{The low block is not uniformly finite}

Consider a flat \(d\)-torus of side length \(L\).  Its Fourier
eigenvalues are
\[
 \lambda_m=\frac{4\pi^2}{L^2}|m|^2,\qquad
 m\in\mathbb Z^d.                                    \tag{64.12}
\]
Let \(N_<(L,\lambda_*)\) be the number of nonzero modes with
\(\lambda_m<\lambda_*\).

\begin{proposition}[Infrared dimension growth]
\label{prop:v16v64-ir-count}
For every fixed \(\lambda_*>0\),
\[
 N_<(L,\lambda_*)
 \geq
 \left(
 2\left\lfloor
 \frac{L\sqrt{\lambda_*}}{4\pi\sqrt d}
 \right\rfloor+1
 \right)^d-1                                         \tag{64.13}
\]
for all sufficiently large \(L\).  In particular
\[
 N_<(L,\lambda_*)\longrightarrow\infty
 \quad\text{at least on the order of }L^d.            \tag{64.14}
\]
\end{proposition}

\begin{proof}
Every lattice vector satisfying
\[
 |m_i|\leq
 \frac{L\sqrt{\lambda_*}}{4\pi\sqrt d}
\quad(1\leq i\leq d)
\]
obeys
\[
 |m|^2\leq\frac{L^2\lambda_*}{16\pi^2},
\qquad
 \lambda_m\leq\frac{\lambda_*}{4}<\lambda_*.
\]
Counting these vectors and removing \(m=0\) gives (64.13).  The
right-hand side grows as a positive constant times \(L^d\).
\end{proof}

Therefore the proposal to handle the finite infrared block explicitly
is finite only at each fixed volume.  Its dimension is not uniform in
the thermodynamic limit.  Pinning a fixed number of global modes
cannot repair the V63 radius by itself.

\subsection{Why the dimension factor in V63 is too expensive}

Proposition~\ref{prop:v16v63-logdet} used
\[
 |\operatorname{tr}T|\leq p\|T\|_{\rm op}.            \tag{64.15}
\]
If \(p\) contains the infrared count (64.14), its constants can be
extensive even when every operator norm is uniform.  This does not
show that the logarithmic determinant density diverges; it shows that
the crude finite-dimensional trace bound cannot establish a
volume-uniform slope or Hessian load.

The correct replacement is a trace-ideal estimate, matching the
orientation and response architecture already isolated in V51.

\begin{theorem}[Dimension-free determinant-response compiler]
\label{thm:v16v64-trace-ideal}
Let \(M(x)\) be invertible on a finite-dimensional Hilbert space, with
dimension allowed to vary.  For a unit parameter direction \(v\), put
\[
 R_v=M^{-1}DM[v].                                     \tag{64.16}
\]
If
\[
 \|R_v\|_{\mathfrak S_1}\leq T_1,                    \tag{64.17}
\]
then
\[
 |D\log|\det M|[v]|\leq T_1.                          \tag{64.18}
\]
If, for unit \(v,w\),
\[
 \|M^{-1}D^2M[v,w]\|_{\mathfrak S_1}\leq T_2,
\qquad
 \|R_v\|_{\mathfrak S_2},
 \|R_w\|_{\mathfrak S_2}\leq T_{\rm HS},              \tag{64.19}
\]
then
\[
 |D^2\log|\det M|[v,w]|
 \leq T_2+T_{\rm HS}^2.                               \tag{64.20}
\]
All constants are independent of the dimension if the stated
Schatten bounds are.
\end{theorem}

\begin{proof}
Jacobi's formula (63.19) and the trace inequality
\(|\operatorname{Tr}A|\leq\|A\|_{\mathfrak S_1}\)
give (64.18).  Equation (63.20) gives two terms.  The first is bounded
by \(T_2\).  For the second, Schatten H\"older gives
\[
 |\operatorname{Tr}(R_wR_v)|
 \leq\|R_wR_v\|_{\mathfrak S_1}
 \leq\|R_w\|_{\mathfrak S_2}
      \|R_v\|_{\mathfrak S_2}
 \leq T_{\rm HS}^2.                                  \tag{64.21}
\]
\end{proof}

\begin{corollary}[Oriented two-factor loader]
\label{cor:v16v64-two-factor}
Suppose
\[
 R_v=A_vB_v,\qquad
 \|A_v\|_{\mathfrak S_2}\leq a,\qquad
 \|B_v\|_{\mathfrak S_2}\leq b.                      \tag{64.22}
\]
Then
\[
 \|R_v\|_{\mathfrak S_1}\leq ab                     \tag{64.23}
\]
and (64.18) holds with \(T_1=ab\).
\end{corollary}

\begin{proof}
This is the Schatten product inequality
\(\|A_vB_v\|_{\mathfrak S_1}
\leq\|A_v\|_{\mathfrak S_2}\|B_v\|_{\mathfrak S_2}\).
\end{proof}

The factor order in (64.22) matters.  V51's covariance-response
analysis showed that an unoriented Gram identity does not by itself
identify the response \(M^{-1}DM\).  Here the inverse
Faddeev--Popov factor must be on the correct side of the derivative
before the trace-ideal budget is claimed.

\subsection{Local density rather than global trace}

On a translation-covariant cutoff, a total logarithmic determinant is
expected to be extensive.  A continuum construction should therefore
distinguish
\[
 \log|\det M_n|
 \quad\text{from}\quad
 |\Lambda_n|^{-1}\log|\det M_n|,                      \tag{64.24}
\]
where \(|\Lambda_n|\) is the physical volume or number of cells.
Uniform control of the normalized quantity does not give a uniform
global gradient in the unnormalized Euclidean norm.  The V57 radial
compiler must use the same volume normalization as the soft
coordinate norm and action density.  Otherwise an extensive
determinant force is incorrectly compared with an intensive
quadratic floor.

\begin{assessment}[What V64 closes]
\label{ass:v16v64-status}
V64 proves a sharp Schur-complement certificate for an
infrared--ultraviolet Faddeev--Popov split.  A fixed ultraviolet
threshold removes the global spectral inverse from that block, but
the number of modes below the threshold grows proportionally to
volume.  The proposed finite low-block treatment is therefore not
uniform.

The dimension factor in the V63 determinant loads is replaced by
trace-class and Hilbert--Schmidt response bounds.  This creates a
precise bridge to V51's oriented response machinery and exposes the
next owner: a local or normalized Schatten estimate for
\(M^{-1}DM\), together with consistent volume scaling.
\end{assessment}

\begin{decision}[V65 companion audit and local response owner]
\label{dec:v16v64-next}
Perform the mandatory YM/NS companion audit.  Then formulate a
translation-local Faddeev--Popov resolvent response whose
volume-normalized trace is controlled by a diagonal kernel bound.
Separate an intensive determinant-density theorem from the stronger
global Schatten bound, and determine which one the soft radial
compiler actually needs under its normalized coordinate convention.
\end{decision}

\begin{firewall}[Claim boundary after V64]
\label{fire:v16v64}
The Schur theorem is conditional on a uniformly invertible infrared
block and small mixed product.  No such nonlinear gravitational block
has been constructed.  The trace-ideal compiler states sufficient
bounds but does not prove them, and an intensive log-determinant limit
does not by itself imply reflection positivity or a global
Poincar\'e gap.
\end{firewall}

\section*{Version V64 append-only record}
\addcontentsline{toc}{section}{Version V64 append-only record}
The complete V63 source was retained byte for byte before this
continuation.  It had \(657964\) bytes, \(19563\) lines, and SHA--256
\[
 \texttt{BDE7BC6C0C405ED42EB75B483526F83BCB2903E67EE69E47463BF3D7F47240F1}.
\]
V64 proved the infrared--ultraviolet Schur certificate, showed that
the low-mode count grows with volume, and replaced dimension-dependent
determinant loads by oriented trace-ideal response criteria.


\section{V65: companion audit and intensive determinant response}
\label{sec:v16v65}

\subsection{Mandatory YM/NS companion audit}

\begin{audit}[V65 companion checkpoint]
\label{aud:v16v65-companions}
The newest active numeric YM note at the checkpoint was V56:
\[
\begin{array}{c|r|r|l}
\hbox{file}&\hbox{lines}&\hbox{bytes}&\hbox{SHA--256}\\ \hline
\texttt{Renewal\_Yang\_Mills\_Mass\_Gap\_Research\_Notes\_V56.tex}
&13172&483563&
\texttt{9A728CAE9E3651DA670139B8B6D5F35E29E4A7243E1870DD0E4793F2521958F5}\\
\texttt{Renewal\_NS\_Stability\_Research\_Notes\_V26.tex}
&5319&278283&
\texttt{82C887F8C2C69E4F28FAD7C3DC8DE5B9099CB5A4BF431EBA1FFF3E91935978AD}.
\end{array}                                             \tag{65.1}
\]
\end{audit}

YM V56 uses exact rational intermediate charts and overlap inequalities
to cross the four V54--V55 certificate defects, enlarging its direct
plus region to a mixed third-coordinate slab.  It retains a direct
minus statement only on the inverted slab and makes no interacting
measure claim.  NS V26 remains unchanged.

\begin{theorem}[V65 import decision]
\label{thm:v16v65-audit-decision}
No companion estimate controls the gravitational
Faddeev--Popov resolvent response.  The YM result reinforces one
applicable construction rule: when one global chart loses a
certificate, an overlapping atlas with exact transition coverage can
be valid even though no single chart covers the whole region.
\end{theorem}

\begin{proof}
Neither companion note contains \(M^{-1}DM\), a metric gauge
determinant, a coarea density, or a Schatten estimate for the
gravitational response.  YM V56's conclusion follows from multiple
local certificates and explicit overlaps; the stated rule imports
that logical pattern without importing its Wilson constants.
\end{proof}

\subsection{Normalized trace ideals}

Let \({\cal H}_N\) have dimension \(p_N\), and let \(N\) denote the
number of physical cells.  Internal multiplicity may make
\(p_N=rN\).  Define
\[
 \|A\|_{1,N}:=\frac1N\|A\|_{\mathfrak S_1},\qquad
 \|A\|_{2,N}:=\frac1{\sqrt N}\|A\|_{\mathfrak S_2}.   \tag{65.2}
\]

\begin{theorem}[Intensive log-determinant compiler]
\label{thm:v16v65-normalized}
Let \(M_N(x)\) be invertible with fixed determinant sign on a chart,
and put
\[
 \ell_N(x)=\frac1N\log|\det M_N(x)|,\qquad
 R_{N,v}=M_N^{-1}DM_N[v].                             \tag{65.3}
\]
If, for unit directions in the chosen parameter norm,
\[
 \|R_{N,v}\|_{1,N}\leq t_1,                           \tag{65.4}
\]
then
\[
 |D\ell_N[v]|\leq t_1.                                \tag{65.5}
\]
If also
\[
 \|M_N^{-1}D^2M_N[v,w]\|_{1,N}\leq t_2,\qquad
 \|R_{N,v}\|_{2,N},\|R_{N,w}\|_{2,N}\leq t_{\rm HS}, \tag{65.6}
\]
then
\[
 |D^2\ell_N[v,w]|\leq t_2+t_{\rm HS}^2.              \tag{65.7}
\]
These estimates are uniform in \(N\) when their normalized inputs
are.
\end{theorem}

\begin{proof}
Divide the two identities (63.19)--(63.20) by \(N\).  Then
\[
 \frac1N|\operatorname{Tr}A|\leq\|A\|_{1,N}
\]
and
\[
 \frac1N\|AB\|_{\mathfrak S_1}
 \leq\frac1N\|A\|_{\mathfrak S_2}\|B\|_{\mathfrak S_2}
 =\|A\|_{2,N}\|B\|_{2,N}.
\]
These give (65.5) and (65.7).
\end{proof}

\subsection{A diagonal-kernel owner}

Let the response act on lattice fields with internal fibre
\(\mathbb C^r\), and write its kernel as \(R_{N,v}(x,y)\).  Then
\[
 D\ell_N[v]
 =\frac1N\sum_{x\in\Lambda_N}
   \operatorname{tr}_{\mathbb C^r}R_{N,v}(x,x).       \tag{65.8}
\]
Suppose the parameter direction is a cell field
\(v=(v_y)\), with normalized norm
\[
 \|v\|_{2,N}^2=\frac1N\sum_y|v_y|^2.                 \tag{65.9}
\]

\begin{proposition}[Local diagonal sensitivity]
\label{prop:v16v65-diagonal}
Assume there is a nonnegative kernel \(K_N(x,y)\) such that
\[
 \left|
 \operatorname{tr}_{\mathbb C^r}R_{N,v}(x,x)
 \right|
 \leq\sum_yK_N(x,y)|v_y|,                             \tag{65.10}
\]
and
\[
 \sup_y\sum_xK_N(x,y)\leq k_1.                        \tag{65.11}
\]
Then
\[
 |D\ell_N[v]|\leq k_1\|v\|_{2,N}.                    \tag{65.12}
\]
\end{proposition}

\begin{proof}
Sum (65.10), interchange the finite sums, and use (65.11):
\[
 |D\ell_N[v]|
 \leq\frac{k_1}{N}\sum_y|v_y|
 \leq k_1
 \left(\frac1N\sum_y|v_y|^2\right)^{1/2}.
\]
\end{proof}

This is an intensive theorem.  It requires only the diagonal trace of
the oriented response, not a global trace-class norm.  A translation
local kernel with uniformly summable tails is a natural sufficient
owner for (65.11).

\subsection{Second derivative through an averaged Schur kernel}

Write the second logarithmic response as
\[
 {\cal S}_{N;v,w}
 =M_N^{-1}D^2M_N[v,w]-R_{N,w}R_{N,v}.                \tag{65.13}
\]
Assume
\[
 \left|
 \operatorname{tr}_{\mathbb C^r}
 {\cal S}_{N;v,w}(x,x)\right|
 \leq\sum_{y,z}K_N^{(2)}(x;y,z)|v_y||w_z|.            \tag{65.14}
\]
Set
\[
 \overline K_N^{(2)}(y,z)
 =\frac1N\sum_xK_N^{(2)}(x;y,z),\qquad
 T_N(y,z)=N\overline K_N^{(2)}(y,z).                 \tag{65.15}
\]

\begin{proposition}[Intensive Hessian from a local kernel]
\label{prop:v16v65-hessian-kernel}
If
\[
 \sup_y\sum_zT_N(y,z)\leq k_2,\qquad
 \sup_z\sum_yT_N(y,z)\leq k_2,                        \tag{65.16}
\]
then
\[
 |D^2\ell_N[v,w]|
 \leq k_2\|v\|_{2,N}\|w\|_{2,N}.                     \tag{65.17}
\]
\end{proposition}

\begin{proof}
Equations (65.13)--(65.15) give
\[
 |D^2\ell_N[v,w]|
 \leq\frac1N\sum_{y,z}T_N(y,z)|v_y||w_z|
 =\langle |v|,T_N|w|\rangle_{2,N}.                   \tag{65.18}
\]
The two Schur sums in (65.16) imply
\(\|T_N\|_{\ell^2\to\ell^2}\leq k_2\).  Cauchy--Schwarz in the
normalized inner product proves (65.17).
\end{proof}

\subsection{Which estimate the probability compiler needs}

Let the total gauge determinant potential be
\[
 V_N^{\rm FP}(x)=-\log|\det M_N(x)|=-N\ell_N(x).      \tag{65.19}
\]
For the normalized inner product in (65.9), define the normalized
gradient by
\[
 D\ell_N[v]=\langle\nabla_N\ell_N,v\rangle_{2,N}.     \tag{65.20}
\]
Then
\[
 \nabla_{\rm Eucl}V_N^{\rm FP}
 =-\nabla_N\ell_N.                                   \tag{65.21}
\]
Thus (65.12) implies
\[
 \|\nabla_{\rm Eucl}V_N^{\rm FP}\|_{\rm Eucl}
 \leq\sqrt N\,k_1.                                   \tag{65.22}
\]
Inserting this as a pure additive force in V57 gives, for
\(\tau>0\),
\[
 x\!\cdot\!\nabla V_N^{\rm FP}
 \geq-\tau|x|_{\rm Eucl}^2-\frac{Nk_1^2}{4\tau}.     \tag{65.23}
\]
The additive Lyapunov constant is extensive.

By contrast, (65.17) says
\[
 {\rm Hess}_{\rm Eucl}V_N^{\rm FP}\succeq-k_2I       \tag{65.24}
\]
with no factor \(N\).  Indeed multiply the density Hessian estimate by
\(N\) and use
\(\|v\|_{2,N}^2=N^{-1}|v|_{\rm Eucl}^2\).

\begin{theorem}[Compiler selection]
\label{thm:v16v65-selection}
An intensive first-derivative estimate such as (65.12) is insufficient
by itself for V56's dimension-uniform global radial-to-local-core
argument, because it produces the extensive constant (65.23).
An intensive Hessian estimate such as (65.17) is sufficient to make a
dimension-uniform Bakry--\'Emery comparison with a bare quadratic
floor, provided
\[
 \rho_{\rm bare}>k_2+\hbox{all other Hessian losses}. \tag{65.25}
\]
If global convexity is unavailable, a volume-uniform proof must replace
the single global Lyapunov core by a local, block, or tensorized
criterion that permits extensive additive constants.
\end{theorem}

\begin{proof}
Equation (65.23) has constant \(Nk_1^2/(4\tau)\), whereas the V56
global local-core radius derived from such a constant grows with
\(N\); its fixed-ball oscillation estimate is therefore not uniform
without additional structure.  Equation (65.24), on the other hand,
is an operator lower bound independent of \(N\).  Adding Hessians gives
(65.25), and the Bakry--\'Emery theorem yields a uniform gap.  The last
statement records the remaining alternative when the premise of the
Hessian route fails.
\end{proof}

\begin{assessment}[What V65 closes]
\label{ass:v16v65-status}
The companion audit records YM's successful overlapping-chart repair
but imports no gravitational estimate.  V65 separates an intensive
determinant density from a global trace-class response, derives
first- and second-derivative diagonal-kernel owners, and tracks their
exact volume scaling.

The result selects the next probability route.  An intensive
Faddeev--Popov Hessian load can enter a uniform convexity margin
directly.  An intensive force bound alone produces an extensive
Lyapunov constant and therefore requires a local or block Poincar\'e
compiler rather than the global V56 core argument.
\end{assessment}

\begin{decision}[V66 block conditional-gap compiler]
\label{dec:v16v65-next}
Construct a finite-volume block Poincar\'e theorem allowing extensive
additive free-energy terms.  Assume uniform one-block conditional gaps
and quantify cross-block Hessian influence by an operator matrix.
Prove a volume-independent global gap under a strict influence margin,
then assign the local Faddeev--Popov, vacuum, curvature, and matter
responses to its entries.
\end{decision}

\begin{firewall}[Claim boundary after V65]
\label{fire:v16v65}
No diagonal sensitivity kernel has yet been proved for the nonlinear
gravitational Faddeev--Popov operator.  The normalized bounds do not
imply a global trace-class limit, and the compiler-selection theorem
does not supply the missing block conditional gaps or influence
margin.
\end{firewall}

\section*{Version V65 append-only record}
\addcontentsline{toc}{section}{Version V65 append-only record}
The complete V64 source was retained byte for byte before this
continuation.  It had \(667632\) bytes, \(19867\) lines, and SHA--256
\[
 \texttt{F31E793A05B8937E61CE0ADBFC0F0565600AEB27E2AAB3A2ADD84B662FA6348D}.
\]
V65 recorded the required companion audit, proved intensive
diagonal-kernel determinant-response estimates, and showed why the
Hessian and radial compilers have different volume requirements.


\section{V66: a block conditional-gap compiler for extensive actions}
\label{sec:v16v66}

\subsection{Product blocks and conditional generators}

Let
\[
 x=(x_1,\ldots,x_N)\in
 \mathbb R^{m_1}\times\cdots\times\mathbb R^{m_N},
\qquad
 d\mu(x)=Z^{-1}e^{-U(x)}\,dx,                         \tag{66.1}
\]
where \(U\) is \(C^2\) and the usual integrations by parts are
justified.  For fixed \(x_{\widehat i}\), let \(\mu_i(\,\cdot\mid
x_{\widehat i})\) be the conditional law of block \(x_i\).
Assume the uniform one-block Poincar\'e inequalities
\[
 \operatorname{Var}_{\mu_i(\cdot\mid x_{\widehat i})}(f)
 \leq\frac1{\rho_i}
 \int|\nabla_if|^2\,d\mu_i(\cdot\mid x_{\widehat i}),
 \qquad \rho_i>0.                                    \tag{66.2}
\]
The constants may depend on the block type but not on the total
number of blocks or the exterior configuration.

For \(i\ne j\), suppose
\[
 \|\nabla_i\nabla_jU(x)\|_{\rm op}\leq\kappa_{ij},
 \qquad
 \kappa_{ij}=\kappa_{ji}\geq0.                        \tag{66.3}
\]
Define the symmetric influence matrix
\[
 {\cal A}_{ii}=\rho_i,\qquad
 {\cal A}_{ij}=-\kappa_{ij}\quad(i\ne j).             \tag{66.4}
\]

\subsection{The conditional one-form estimate}

\begin{lemma}[Conditional Witten estimate]
\label{lem:v16v66-witten}
Let
\[
 L_i=-\Delta_i+\nabla_iU\!\cdot\!\nabla_i             \tag{66.5}
\]
be the nonnegative conditional generator.  Under (66.2), every smooth
scalar \(h\) satisfies, after conditioning on \(x_{\widehat i}\),
\[
 \int\left(
 \|\nabla_i^2h\|_{\rm HS}^2+
 \left\langle\nabla_ih,
 (\nabla_i^2U)\nabla_ih\right\rangle
 \right)d\mu_i
 \geq\rho_i\int|\nabla_ih|^2\,d\mu_i.                \tag{66.6}
\]
\end{lemma}

\begin{proof}
The conditional Bochner identity is
\[
 \int(L_ih)^2\,d\mu_i
 =\int\left(
 \|\nabla_i^2h\|_{\rm HS}^2+
 \langle\nabla_ih,(\nabla_i^2U)\nabla_ih\rangle
 \right)d\mu_i.                                      \tag{66.7}
\]
The conditional mean of \(L_ih\) is zero.  By the spectral theorem,
the Poincar\'e gap (66.2) implies
\[
 \int(L_ih)^2\,d\mu_i
 \geq\rho_i\int hL_ih\,d\mu_i
 =\rho_i\int|\nabla_ih|^2\,d\mu_i.                  \tag{66.8}
\]
Combine (66.7)--(66.8).
\end{proof}

\subsection{Global block theorem}

\begin{theorem}[Block conditional-gap criterion]
\label{thm:v16v66-block-gap}
If
\[
 {\cal A}\succeq\rho_*I_N
 \quad\text{for some }\rho_*>0,                       \tag{66.9}
\]
then
\[
 \operatorname{Var}_\mu(f)
 \leq\frac1{\rho_*}\int\sum_{i=1}^N|\nabla_if|^2\,d\mu
                                                               \tag{66.10}
\]
for every smooth \(f\).  The constant is independent of \(N\).
\end{theorem}

\begin{proof}
It is enough first to consider centered \(f\) in the range of the
global generator
\[
 L=-\Delta+\nabla U\!\cdot\!\nabla.
\]
Choose \(h\) with \(Lh=f\).  The global Bochner identity gives
\[
 \int(Lh)^2\,d\mu
 =
 \int\left(
 \|\nabla^2h\|_{\rm HS}^2+
 \langle\nabla h,(\nabla^2U)\nabla h\rangle
 \right)d\mu.                                        \tag{66.11}
\]
For each \(i\), retain the diagonal derivative and Hessian terms and
apply Lemma~\ref{lem:v16v66-witten}, then integrate the exterior
variables.  The discarded mixed second-derivative squares are
nonnegative.  With \(a_i=|\nabla_ih|\), (66.3) yields
\[
 \int(Lh)^2\,d\mu
 \geq
 \int\left[
 \sum_i\rho_i a_i^2-
 \sum_{i\ne j}\kappa_{ij}a_ia_j
 \right]d\mu
 =
 \int a^{\mathsf T}{\cal A}a\,d\mu.                  \tag{66.12}
\]
By (66.9),
\[
 \int(Lh)^2\,d\mu
 \geq\rho_*\int|\nabla h|^2\,d\mu.                   \tag{66.13}
\]
Integration by parts and Cauchy--Schwarz give
\[
 \operatorname{Var}_\mu(f)
 =\int fLh\,d\mu
 =\int\nabla f\!\cdot\!\nabla h\,d\mu
 \leq
 \left(\int|\nabla f|^2d\mu\right)^{1/2}
 \left(\int|\nabla h|^2d\mu\right)^{1/2}.             \tag{66.14}
\]
Since \(\int(Lh)^2=\operatorname{Var}_\mu(f)\), insert (66.13) into
(66.14) and cancel the square root of the variance.  This proves
(66.10) on the generator range.  Standard resolvent approximation,
using \(h_\varepsilon=(L+\varepsilon)^{-1}f\), extends the inequality
to the Dirichlet-form domain.
\end{proof}

\begin{corollary}[Row-sum margin]
\label{cor:v16v66-row-sum}
It is sufficient that
\[
 \rho_*:=
 \inf_i\left(\rho_i-\sum_{j\ne i}\kappa_{ij}\right)>0.             \tag{66.15}
\]
\end{corollary}

\begin{proof}
For \(a\in\mathbb R^N\),
\[
 2|a_ia_j|\leq a_i^2+a_j^2.
\]
Therefore
\[
 a^{\mathsf T}{\cal A}a
 \geq\sum_i
 \left(\rho_i-\sum_{j\ne i}\kappa_{ij}\right)a_i^2
 \geq\rho_*|a|^2.
\]
\end{proof}

The theorem permits \(U\) and its additive Lyapunov constants to be
extensive.  Only the conditional gaps and the influence row sums must
remain uniform.

\subsection{One-block gaps from local radial control}

Let the conditional potential in block \(i\) be
\[
 U_i(x_i;x_{\widehat i})
 =U_{i,{\rm core}}(x_i)+I_i(x_i;x_{\widehat i}).      \tag{66.16}
\]
Suppose, uniformly in the exterior configuration,
\[
 x_i\!\cdot\!\nabla_iU_i
 \geq\rho_{i,0}|x_i|^2-C_{i,0},                      \tag{66.17}
\]
and that on a fixed connected ball the conditional density has the
bounded-oscillation property required by V56.  Then V56 gives a
conditional gap \(\rho_i>0\).  The constants \(C_{i,0}\) may be
nonzero for every block; they do not add in (66.2), because the
one-block theorem is applied before the global assembly.

This is precisely where the intensive first-derivative estimate of
V65 can be useful: if the diagonal-kernel force seen by any fixed
block is uniformly bounded, it contributes a uniform local Young
constant rather than the extensive global constant (65.23).

\subsection{Physical influence ledger}

Decompose the constrained gauge-fixed potential into
\[
 U=U^{\rm grav}_0+V^{\rm FP}+V^{\rm vac}+V^R
   +V^{\rm gauge}+V^{\rm ferm}+V^{\rm Higgs}.         \tag{66.18}
\]
For distinct blocks define
\[
 \kappa_{ij}\leq
 \kappa^{0}_{ij}
 +\kappa^{\rm FP}_{ij}
 +\kappa^{\rm vac}_{ij}
 +\kappa^{R}_{ij}
 +\kappa^{\rm gauge}_{ij}
 +\kappa^{\rm ferm}_{ij}
 +\kappa^{\rm Higgs}_{ij},                            \tag{66.19}
\]
where each term is a uniform upper bound on the corresponding mixed
block Hessian norm.

The earlier owners enter as follows:
\begin{enumerate}
\item the V65 second diagonal-kernel response supplies
      \(\kappa^{\rm FP}_{ij}\) when its bilinear kernel is retained
      before taking the Schur row sum;
\item V46's determinant Hessian and reconstruction responses supply
      \(\kappa^{\rm vac}_{ij}\);
\item V47--V54's bulk, face, joint, and collar response arrays supply
      \(\kappa^R_{ij}\);
\item V51's oriented covariance response supplies the matter-mediated
      metric Hessian entries only after its Hilbert--Schmidt tails are
      verified.
\end{enumerate}
Finite interaction range is sufficient but not necessary.  Summable
off-diagonal decay,
\[
 \sup_i\sum_{j\ne i}\kappa_{ij}<\infty,               \tag{66.20}
\]
is the actual owner required by (66.15).

\begin{proposition}[Locality-to-gap compiler]
\label{prop:v16v66-locality}
Assume every conditional block has gap
\(\rho_i\geq\rho_{\rm loc}>0\) and
\[
 \sup_i\sum_{j\ne i}\kappa_{ij}\leq\kappa_{\rm row}
 <\rho_{\rm loc}.                                    \tag{66.21}
\]
Then
\[
 \operatorname{Var}_\mu(f)
 \leq\frac1{\rho_{\rm loc}-\kappa_{\rm row}}
 \int|\nabla f|^2\,d\mu.                             \tag{66.22}
\]
\end{proposition}

\begin{proof}
Apply Corollary~\ref{cor:v16v66-row-sum} with
\(\rho_*=\rho_{\rm loc}-\kappa_{\rm row}\), then Theorem
\ref{thm:v16v66-block-gap}.
\end{proof}

\subsection{The new bottleneck}

The compiler changes the gravitational question from a global
dimension-dependent Lyapunov estimate to two local statements:
\begin{align}
 &\inf_{i,x_{\widehat i}}
 \operatorname{gap}\bigl(\mu_i(\cdot\mid x_{\widehat i})\bigr)>0,
                                                               \tag{66.23}\\
 &\sup_i\sum_{j\ne i}
 \|\nabla_i\nabla_jU\|_{\rm op}
 <\inf_i\rho_i.                                       \tag{66.24}
\end{align}
For a massless elliptic Faddeev--Popov inverse, locality of the
differential \(DM\) does not automatically imply summability of
\(M^{-1}DM\).  The next version must calculate the decay exponent of
the resolvent-response kernel and determine whether derivatives
improve it enough for (66.20).

\begin{assessment}[What V66 closes]
\label{ass:v16v66-status}
V66 proves a volume-independent block Poincar\'e theorem from uniform
conditional gaps and a positive influence matrix.  The proof uses a
conditional Witten estimate and the global Bochner identity, so it
allows nonconvex conditional potentials and extensive additive
constants.

All current gravitational, gauge-determinant, vacuum, curvature, and
matter responses now have a common target: one-block gap ownership and
a summable mixed-Hessian row.  Neither target has yet been verified
for the physical constrained cutoff law.
\end{assessment}

\begin{decision}[V67 massless resolvent decay]
\label{dec:v16v66-next}
Compute a model lattice Faddeev--Popov response
\(M^{-1}DM\) when \(M\) is a massless Laplacian on mean-zero fields and
\(DM\) is a local first- or second-order coefficient perturbation.
Derive the spatial decay and row-sum threshold in dimension four.
If the row is logarithmically or power divergent, introduce a
derivative cancellation or neutral-block decomposition before
returning to the nonlinear metric operator.
\end{decision}

\begin{firewall}[Claim boundary after V66]
\label{fire:v16v66}
The block theorem is a sufficient analytic criterion.  It does not
prove the conditional laws are positive, establish their gaps, or
verify the influence bounds.  It does not by itself imply
reflection positivity, clustering in physical distance, or the
Standard Model--Einstein continuum.
\end{firewall}

\section*{Version V66 append-only record}
\addcontentsline{toc}{section}{Version V66 append-only record}
The complete V65 source was retained byte for byte before this
continuation.  It had \(677451\) bytes, \(20170\) lines, and SHA--256
\[
 \texttt{EDC59EEFB70561DB53B8601908B5D0E299CF36A2F59E2AFE95AA54B5E89AC202}.
\]
V66 proved the block conditional-gap criterion, its row-sum corollary,
and the locality ledger that admits extensive determinant and action
terms while retaining a volume-independent global Poincar\'e gap.


\section{V67: exact massless graph-resolvent response}
\label{sec:v16v67}

\subsection{Weighted graph Faddeev--Popov model}

Let \(G_N=(V_N,E_N)\) be a finite connected graph with an oriented
incidence operator
\[
 D:\mathbb R^{V_N}\longrightarrow\mathbb R^{E_N}.     \tag{67.1}
\]
For positive edge conductances \(a=(a_e)\), put
\[
 A=\operatorname{diag}(a_e),\qquad
 L(a)=D^{\mathsf T}AD.                                \tag{67.2}
\]
The constants form the kernel of \(L(a)\).  Let
\[
 {\cal H}_0=\left\{u:\sum_{x\in V_N}u_x=0\right\},
\qquad
 G(a)=L(a)|_{{\cal H}_0}^{-1}.                        \tag{67.3}
\]
The determinant below is the determinant on \({\cal H}_0\):
\[
 {\cal F}(a)=-\log\det{}_{{\cal H}_0}L(a).            \tag{67.4}
\]

For an edge \(e\), let
\[
 b_e=D^{\mathsf T}{\bf e}_e,\qquad
 q_{ef}=b_e^{\mathsf T}G(a)b_f.                       \tag{67.5}
\]

\subsection{Exact first and second variations}

\begin{theorem}[Effective-resistance response]
\label{thm:v16v67-response}
For every pair of edges,
\[
 \partial_{a_e}{\cal F}(a)=-q_{ee},\qquad
 \partial_{a_e}\partial_{a_f}{\cal F}(a)=q_{ef}^2.   \tag{67.6}
\]
Consequently \({\cal F}\) is convex in the conductance variables.
\end{theorem}

\begin{proof}
Since
\[
 \partial_{a_e}L=b_eb_e^{\mathsf T},                  \tag{67.7}
\]
Jacobi's formula gives
\[
 \partial_{a_e}\log\det L
 =\operatorname{Tr}(G b_eb_e^{\mathsf T})
 =b_e^{\mathsf T}Gb_e=q_{ee}.                        \tag{67.8}
\]
Also
\[
 \partial_{a_f}G=-G b_fb_f^{\mathsf T}G.             \tag{67.9}
\]
Differentiate (67.8) and insert (67.9):
\[
 \partial_{a_f}q_{ee}
 =-(b_e^{\mathsf T}Gb_f)(b_f^{\mathsf T}Gb_e)
 =-q_{ef}^2.
\]
The signs in (67.6) follow from (67.4).  The Hessian matrix
\((q_{ef}^2)\) is positive semidefinite by the Schur product theorem,
because \(Q=(q_{ef})=DGD^{\mathsf T}\succeq0\).
\end{proof}

The massless inverse \(G\) itself has no volume-uniform operator norm,
but the differentiated response in (67.6) has a stronger exact
structure.

\subsection{Projection identity and uniform row sums}

\begin{theorem}[Conductance projection]
\label{thm:v16v67-projection}
The edge operator
\[
 P=A^{1/2}DGD^{\mathsf T}A^{1/2}                     \tag{67.10}
\]
is an orthogonal projection.  If
\[
 a_e\geq a_->0\qquad(e\in E_N),                       \tag{67.11}
\]
then
\[
 0\leq q_{ee}\leq a_e^{-1}\leq a_-^{-1},             \tag{67.12}
\]
and
\[
 \sum_{f\in E_N}q_{ef}^2
 \leq\frac1{a_ea_-}\,P_{ee}
 \leq a_-^{-2}.                                      \tag{67.13}
\]
The bounds are independent of the graph size and its spectral gap.
\end{theorem}

\begin{proof}
Symmetry of \(P\) is clear.  On \({\cal H}_0\),
\(GL=LG=I\), and \(D\) annihilates constants.  Therefore
\[
 P^2
 =A^{1/2}DG D^{\mathsf T}ADG D^{\mathsf T}A^{1/2}
 =A^{1/2}DG L G D^{\mathsf T}A^{1/2}
 =P.                                                  \tag{67.14}
\]
Thus \(0\preceq P\preceq I\), and
\[
 P_{ef}=\sqrt{a_ea_f}\,q_{ef}.                       \tag{67.15}
\]
The diagonal identity gives
\(q_{ee}=P_{ee}/a_e\), proving (67.12).  Since a projection satisfies
\[
 \sum_fP_{ef}^2=(P^2)_{ee}=P_{ee},                   \tag{67.16}
\]
equation (67.15) gives
\[
 \sum_fq_{ef}^2
 =\frac1{a_e}\sum_f\frac{P_{ef}^2}{a_f}
 \leq\frac{P_{ee}}{a_ea_-},
\]
which is (67.13).
\end{proof}

\begin{corollary}[Volume-uniform block influence]
\label{cor:v16v67-influence}
Treat each conductance as a scalar block.  The Faddeev--Popov
potential (67.4) has
\[
 |\partial_{a_e}{\cal F}|\leq a_-^{-1},\qquad
 \sup_e\sum_f
 |\partial_{a_e}\partial_{a_f}{\cal F}|
 \leq a_-^{-2}.                                      \tag{67.17}
\]
Its full Hessian is nonnegative, so it consumes no convexity floor in
conductance coordinates.
\end{corollary}

\begin{proof}
Combine (67.6), (67.12), and (67.13).  Convexity was proved in
Theorem~\ref{thm:v16v67-response}.
\end{proof}

This is stronger than the spatial-decay estimate requested in V66.
It gives the complete row sum without estimating the Green kernel
pointwise.

\subsection{Derivative cancellation and spatial reading}

For a translation-invariant lattice Laplacian in \(d>2\), its Green
kernel has the continuum-scale behaviour
\[
 G(x,y)\sim |x-y|^{2-d}.                              \tag{67.18}
\]
The edge response \(q_{ef}\) differentiates the Green kernel once at
each endpoint, so its formal off-diagonal decay is
\[
 |q_{ef}|\sim{\rm dist}(e,f)^{-d},\qquad
 q_{ef}^2\sim{\rm dist}(e,f)^{-2d}.                  \tag{67.19}
\]
The latter is summable.  Equations (67.13)--(67.16) are the exact
finite-graph substitute for this heuristic and remain valid in every
dimension and at every volume.

The cancellation depends on the divergence-form derivative
\(\partial_{a_e}L=b_eb_e^{\mathsf T}\).  A zeroth-order local
perturbation would expose \(G\) without the two endpoint derivatives
and would not inherit the projection identity.

\subsection{Matrix and nonlinear limits of the model}

The exact result applies to a scalar, symmetric, positive
Faddeev--Popov operator that is linear in positive conductances.
For a nonlinear metric-dependent operator \(M(g)\),
\[
 D^2\{-\log\det M\}[v,w]
 =
 \operatorname{Tr}\bigl(
 M^{-1}DM[w]M^{-1}DM[v]\bigr)
 -\operatorname{Tr}\bigl(M^{-1}D^2M[v,w]\bigr).
                                                               \tag{67.20}
\]
The first term is nonnegative when \(M\) is positive and \(v=w\), but
the second has no automatic sign.  Likewise, a nonsymmetric
Faddeev--Popov operator need not admit the projection (67.10).

\begin{proposition}[Linear-coordinate safe corridor]
\label{prop:v16v67-safe}
Suppose a gauge operator can be parameterized on a chart as
\[
 M(a)=D^{\mathsf T}\operatorname{diag}(a_e)D
 \quad\text{with }a_e\geq a_->0,                     \tag{67.21}
\]
and the metric-to-conductance map \(a=a(g)\) is affine on the
reconstruction block.  Then the pulled-back
\(-\log\det M(a(g))\) is convex on that block.
\end{proposition}

\begin{proof}
The Hessian of \({\cal F}\circ a\) is
\[
 D^2({\cal F}\circ a)[v,v]
 =D^2{\cal F}[Da[v],Da[v]]
 +D{\cal F}[D^2a[v,v]].
\]
The second term vanishes for affine \(a\), and the first is
nonnegative by Theorem~\ref{thm:v16v67-response}.
\end{proof}

The affine condition is restrictive but identifies the precise
upstream design advantage of conductance coordinates: the dangerous
second term in (67.20) disappears.

\begin{assessment}[What V67 closes]
\label{ass:v16v67-status}
V67 solves the massless resolvent-response model exactly.  For a
weighted graph Laplacian, first determinant forces are uniformly
local in size, the mixed Hessian row sum is at most \(a_-^{-2}\), and
the negative log determinant is convex.  These statements are
independent of volume and of the vanishing spectral gap.

The mechanism is the conductance projection, equivalently two
derivatives on the massless Green kernel.  A nonlinear,
matrix-valued, or nonsymmetric gravitational Faddeev--Popov operator
has not been reduced to this model; its \(D^2M\) remainder is the next
owner.
\end{assessment}

\begin{decision}[V68 nonlinear pullback remainder]
\label{dec:v16v67-next}
Derive a quantitative pullback theorem for
\({\cal F}(a(g))\).  Bound the negative term
\(D{\cal F}[D^2a]\) by the uniform first force (67.12), retain the
positive Gram Hessian, and obtain cellwise slope and influence
budgets in terms of \(Da\) and \(D^2a\).  Determine the small-curvature
condition on the metric-to-conductance map that preserves the V66
block margin.
\end{decision}

\begin{firewall}[Claim boundary after V67]
\label{fire:v16v67}
The graph theorem is an exact model result, not a derivation of the
gravitational Faddeev--Popov operator.  It assumes positivity,
symmetry, scalar conductances, and removal of the constant kernel.
It proves neither a nonlinear gauge atlas nor reflection positivity.
\end{firewall}

\section*{Version V67 append-only record}
\addcontentsline{toc}{section}{Version V67 append-only record}
The complete V66 source was retained byte for byte before this
continuation.  It had \(687735\) bytes, \(20491\) lines, and SHA--256
\[
 \texttt{617EBC37F99F9A1402E4828C5F9B4C78E6A7DAC65ACC98188815D2068EDA940A}.
\]
V67 proved the exact graph-Laplacian determinant response, its
projection identity, and volume-uniform force and influence bounds for
the massless gauge model.


\section{V68: nonlinear conductance pullback and curvature margin}
\label{sec:v16v68}

\subsection{Metric blocks and conductance sensitivities}

Retain the connected weighted graph of V67.  Let
\[
 g=(g_1,\ldots,g_N)\longmapsto
 a(g)=(a_e(g):e\in E_N)                               \tag{68.1}
\]
be a \(C^2\) reconstruction map from metric blocks to positive
conductances satisfying
\[
 a_e(g)\geq a_->0.                                   \tag{68.2}
\]
Define
\[
 V^{\rm FP}(g)
 =-\log\det{}_{{\cal H}_0}
 \left[D^{\mathsf T}\operatorname{diag}(a_e(g))D\right].
                                                               \tag{68.3}
\]
For block directions \(v_i\) and \(w_j\), set
\[
 A_{ei}
 :=\sup_g\sup_{|v_i|=1}|D_i a_e(g)[v_i]|,            \tag{68.4}
\]
\[
 B_{eij}
 :=\sup_g\sup_{\substack{|v_i|=1\\|w_j|=1}}
 |D_iD_j a_e(g)[v_i,w_j]|.                           \tag{68.5}
\]
These constants retain the support pattern of the reconstruction.

\subsection{Exact pullback Hessian}

\begin{theorem}[Conductance pullback identity]
\label{thm:v16v68-pullback}
For arbitrary directions \(v,w\),
\[
 DV^{\rm FP}(g)[v]
 =-\sum_e q_{ee}\,Da_e(g)[v],                        \tag{68.6}
\]
and
\[
 D^2V^{\rm FP}(g)[v,w]
 =
 \sum_{e,f}q_{ef}^2\,Da_e(g)[v]Da_f(g)[w]
 -\sum_eq_{ee}\,D^2a_e(g)[v,w].                     \tag{68.7}
\]
The first term in (68.7) is positive semidefinite.  Therefore
\[
 D^2V^{\rm FP}(g)[v,v]
 \geq-\sum_eq_{ee}|D^2a_e(g)[v,v]|.                 \tag{68.8}
\]
\end{theorem}

\begin{proof}
Apply the ordinary first- and second-order chain rules to
\({\cal F}(a(g))\), then insert the conductance derivatives (67.6).
The first term in (68.7) is the quadratic form of the positive
semidefinite matrix \((q_{ef}^2)\), proving (68.8).
\end{proof}

Thus all loss of convexity comes from the curvature of the
metric-to-conductance map.  The massless inverse does not create an
additional negative term.

\subsection{Local force and mixed-influence bounds}

Define the sensitivity sums
\[
 A_{\rm out}:=\sup_i\sum_eA_{ei},\qquad
 A_{\rm in}:=\sup_e\sum_iA_{ei},                     \tag{68.9}
\]
and
\[
 B_{\rm row}:=
 \sup_i\sum_j\sum_eB_{eij}.                           \tag{68.10}
\]

\begin{proposition}[Uniform local force]
\label{prop:v16v68-force}
For every block \(i\),
\[
 \|\nabla_iV^{\rm FP}(g)\|
 \leq\frac1{a_-}\sum_eA_{ei}
 \leq\frac{A_{\rm out}}{a_-}.                        \tag{68.11}
\]
\end{proposition}

\begin{proof}
Use (68.6), \(q_{ee}\leq a_-^{-1}\) from (67.12), and take the
supremum over unit block directions.
\end{proof}

\begin{proposition}[Uniform mixed-Hessian row]
\label{prop:v16v68-row}
Let
\[
 \kappa_{ij}^{\rm FP}
 :=\sup_g\|\nabla_i\nabla_jV^{\rm FP}(g)\|_{\rm op}.
                                                               \tag{68.12}
\]
Then
\[
 \sup_i\sum_j\kappa_{ij}^{\rm FP}
 \leq
 \frac{A_{\rm out}A_{\rm in}}{a_-^2}
 +\frac{B_{\rm row}}{a_-}.                           \tag{68.13}
\]
\end{proposition}

\begin{proof}
For unit \(v_i,w_j\), equation (68.7) and the definitions give
\[
 |D_iD_jV^{\rm FP}[v_i,w_j]|
 \leq
 \sum_{e,f}q_{ef}^2A_{ei}A_{fj}
 +\sum_eq_{ee}B_{eij}.                               \tag{68.14}
\]
Sum over \(j\).  For the Gram term,
\[
 \sum_j\sum_{e,f}q_{ef}^2A_{ei}A_{fj}
 \leq
 A_{\rm in}\sum_eA_{ei}\sum_fq_{ef}^2
 \leq\frac{A_{\rm in}A_{\rm out}}{a_-^2},            \tag{68.15}
\]
where (67.13) was used in the last step.  The second term is at most
\[
 \frac1{a_-}\sum_{j,e}B_{eij}
 \leq\frac{B_{\rm row}}{a_-}.                        \tag{68.16}
\]
Take the supremum over directions and \(i\).
\end{proof}

The quantities in (68.9)--(68.10) are uniform for a finite-range
reconstruction with uniformly bounded local degree and derivatives.
They need not be uniform for a global spectral parameterization.

\subsection{A sharper lower-form estimate}

The row estimate charges the positive Gram term as an absolute mixed
influence.  For a global Hessian comparison one can retain its sign.
Define the nonnegative block-curvature form
\[
 {\cal B}(v)
 :=\sum_eq_{ee}|D^2a_e(g)[v,v]|.                     \tag{68.17}
\]

\begin{corollary}[Pullback form-smallness]
\label{cor:v16v68-form}
If
\[
 {\cal B}(v)\leq\eta_a|v|^2                          \tag{68.18}
\]
uniformly on the chart, then
\[
 {\rm Hess}\,V^{\rm FP}\succeq-\eta_aI.              \tag{68.19}
\]
It is sufficient that
\[
 \sum_e|D^2a_e(g)[v,v]|\leq a_-\eta_a|v|^2.          \tag{68.20}
\]
\end{corollary}

\begin{proof}
Equation (68.8) gives the first conclusion.  The bound
\(q_{ee}\leq a_-^{-1}\) shows that (68.20) implies (68.18).
\end{proof}

When \(a(g)\) is affine, \(\eta_a=0\) and
\(V^{\rm FP}\) is convex, recovering Proposition
\ref{prop:v16v67-safe}.  For nonlinear \(a\), (68.20) is the exact
small-curvature substitute.

\subsection{Insertion into the block gap}

Let
\[
 \rho_{\rm loc}:=\inf_i\rho_i,\qquad
 \kappa_{\rm other}:=
 \sup_i\sum_{j\ne i}
 \left(
 \kappa^0_{ij}+\kappa^{\rm vac}_{ij}+\kappa^R_{ij}
 +\kappa^{\rm gauge}_{ij}
 +\kappa^{\rm ferm}_{ij}+\kappa^{\rm Higgs}_{ij}
 \right).                                            \tag{68.21}
\]

\begin{theorem}[Nonlinear conductance corridor]
\label{thm:v16v68-corridor}
Assume the one-block conditional gaps in (66.2) satisfy
\(\rho_i\geq\rho_{\rm loc}>0\).  If
\[
 \rho_{\rm loc}
 >
 \kappa_{\rm other}
 +\frac{A_{\rm out}A_{\rm in}}{a_-^2}
 +\frac{B_{\rm row}}{a_-},                           \tag{68.22}
\]
then the full constrained measure has Poincar\'e constant at most
\[
 \left[
 \rho_{\rm loc}
 -\kappa_{\rm other}
 -\frac{A_{\rm out}A_{\rm in}}{a_-^2}
 -\frac{B_{\rm row}}{a_-}
 \right]^{-1}.                                       \tag{68.23}
\]
\end{theorem}

\begin{proof}
Proposition~\ref{prop:v16v68-row} bounds the complete
Faddeev--Popov influence row.  Add the other rows in (68.21) and apply
Proposition~\ref{prop:v16v66-locality}.
\end{proof}

This is a conservative sufficient condition because it discards the
positive matrix structure of the Gram term when forming absolute block
influences.  Under a global Hessian route, (68.19) charges only the
map-curvature loss \(\eta_a\).

\subsection{Coordinate-design consequence}

Suppose each edge conductance depends only on a bounded neighbourhood
\({\cal N}(e)\) of metric blocks and
\[
 |{\cal N}(e)|\leq d_a,\qquad
 \#\{e:i\in{\cal N}(e)\}\leq d_g,                    \tag{68.24}
\]
while
\[
 A_{ei}\leq A_0{\bf1}_{i\in{\cal N}(e)},\qquad
 B_{eij}\leq B_0
 {\bf1}_{i,j\in{\cal N}(e)}.                         \tag{68.25}
\]
Then
\[
 A_{\rm out}\leq d_gA_0,\qquad
 A_{\rm in}\leq d_aA_0,\qquad
 B_{\rm row}\leq d_gd_aB_0.                          \tag{68.26}
\]
Hence (68.22) becomes the explicit local condition
\[
 \rho_{\rm loc}>
 \kappa_{\rm other}
 +\frac{d_gd_aA_0^2}{a_-^2}
 +\frac{d_gd_aB_0}{a_-}.                             \tag{68.27}
\]

\begin{proof}
Each sum in (68.9)--(68.10) contains at most the incidences counted in
(68.24), and every nonzero summand is bounded by (68.25).  Substitute
the resulting bounds into (68.22).
\end{proof}

Thus locality and bounded reconstruction curvature are enough to turn
the exact graph projection into a volume-uniform influence
certificate.  The remaining gap is structural: the actual residual
metric gauge operator must be represented by a positive
divergence-form graph operator with such local conductances.

\begin{assessment}[What V68 closes]
\label{ass:v16v68-status}
V68 derives the exact nonlinear pullback of the graph determinant.
Its Hessian splits into a positive resolvent Gram term and a sole
sign-indefinite reconstruction-curvature term.  Explicit local force,
influence-row, and form-smallness bounds are obtained, leading to a
volume-independent block-gap corridor.

The certificate is now checkable from a local metric-to-conductance
map.  It has not been shown that the gravitational residual
Faddeev--Popov operator is scalar, symmetric, positive, or affine
enough to satisfy it.
\end{assessment}

\begin{decision}[V69 vector-valued conductances]
\label{dec:v16v68-next}
Extend the projection identity to a vector gauge field with
positive-definite matrix conductances on edges.  Derive dimension and
ellipticity dependence for the first force and mixed-Hessian row, and
identify precisely which nonsymmetric lower-order terms destroy the
projection.  This will decide whether the actual linearized
volume-preserving diffeomorphism operator can enter the V68 corridor.
\end{decision}

\begin{firewall}[Claim boundary after V68]
\label{fire:v16v68}
The corridor assumes uniform one-block conditional gaps and all
non-Faddeev--Popov influence rows.  It does not verify the physical
gauge operator representation, nonlinear gauge atlas, determinant
orientation, Ward limit, or reflection positivity.
\end{firewall}

\section*{Version V68 append-only record}
\addcontentsline{toc}{section}{Version V68 append-only record}
The complete V67 source was retained byte for byte before this
continuation.  It had \(696081\) bytes, \(20752\) lines, and SHA--256
\[
 \texttt{F30B0CF8FB559A8175B55DB50AACE51C2191EA5D3F2EDE9F0607F344D91B4568}.
\]
V68 pulled the massless graph determinant through a nonlinear local
metric reconstruction and quantified the exact map-curvature margin
required by the block Poincar\'e compiler.


\section{V69: matrix conductances and the vector gauge determinant}
\label{sec:v16v69}

\subsection{Vector-valued weighted graph operator}

Let the connected graph and incidence map of V67 be fixed, and let
the gauge fibre have dimension \(r_g\).  Define
\[
 {\bf D}=D\otimes I_{r_g}:
 \mathbb R^{V_N}\otimes\mathbb R^{r_g}
 \longrightarrow
 \mathbb R^{E_N}\otimes\mathbb R^{r_g}.               \tag{69.1}
\]
At every edge let \(A_e\) be a real symmetric positive-definite
\(r_g\times r_g\) matrix, and put
\[
 {\bf A}=\bigoplus_{e\in E_N}A_e,\qquad
 {\bf L}={\bf D}^{\mathsf T}{\bf A}{\bf D}.           \tag{69.2}
\]
On the orthogonal complement \({\cal H}_0\) of the
\(r_g\)-dimensional constant kernel, write
\[
 {\bf G}={\bf L}|_{{\cal H}_0}^{-1},\qquad
 {\bf Q}={\bf D}{\bf G}{\bf D}^{\mathsf T}.           \tag{69.3}
\]
Assume the uniform ellipticity bounds
\[
 a_-I_{r_g}\preceq A_e\preceq a_+I_{r_g}.            \tag{69.4}
\]

\subsection{Projection and block row identity}

\begin{theorem}[Matrix conductance projection]
\label{thm:v16v69-projection}
The edge-fibre operator
\[
 {\bf P}={\bf A}^{1/2}{\bf Q}{\bf A}^{1/2}            \tag{69.5}
\]
is an orthogonal projection.  Its edge blocks satisfy
\[
 \sum_f\|P_{ef}\|_{\rm HS}^2
 =\operatorname{Tr}P_{ee}\leq r_g.                   \tag{69.6}
\]
Consequently
\[
 0\preceq Q_{ee}\preceq A_e^{-1},\qquad
 \operatorname{Tr}Q_{ee}\leq\frac{r_g}{a_-},         \tag{69.7}
\]
and
\[
 \sum_f\|Q_{ef}\|_{\rm HS}^2
 \leq\frac{r_g}{a_-^2}.                              \tag{69.8}
\]
\end{theorem}

\begin{proof}
The proof of (67.14) is unchanged:
\[
 {\bf P}^2
 ={\bf A}^{1/2}{\bf D}{\bf G}
 {\bf L}{\bf G}{\bf D}^{\mathsf T}{\bf A}^{1/2}
 ={\bf P}.
\]
It is symmetric, hence an orthogonal projection.  Taking the trace of
the \((e,e)\) block of \({\bf P}^2={\bf P}\) gives
\[
 \sum_f\operatorname{Tr}(P_{ef}P_{fe})
 =\operatorname{Tr}P_{ee}.
\]
Since \(P_{fe}=P_{ef}^{\mathsf T}\), this is (69.6), and
\(0\preceq P_{ee}\preceq I_{r_g}\) gives its final bound.
Equation (69.5) yields
\[
 Q_{ef}=A_e^{-1/2}P_{ef}A_f^{-1/2}.                  \tag{69.9}
\]
The diagonal inequality in (69.7) follows by congruence.  Finally,
\[
 \|Q_{ef}\|_{\rm HS}
 \leq a_-^{-1}\|P_{ef}\|_{\rm HS},
\]
and summing its square proves (69.8).
\end{proof}

\subsection{Matrix determinant response}

Let
\[
 {\cal F}({\bf A})=-\log\det{}_{{\cal H}_0}{\bf L}.   \tag{69.10}
\]
For an edge-matrix variation
\({\bf H}=\bigoplus_eH_e\),
\[
 D{\cal F}[{\bf H}]
 =-\operatorname{Tr}({\bf Q}{\bf H})
 =-\sum_e\operatorname{Tr}(Q_{ee}H_e).               \tag{69.11}
\]

\begin{theorem}[Vector force and Hessian row]
\label{thm:v16v69-response}
For variations supported at one edge,
\[
 |D_e{\cal F}[H_e]|
 \leq\frac{r_g}{a_-}\|H_e\|_{\rm op}.                \tag{69.12}
\]
For edge-supported variations \(H_e,K_f\),
\[
 D_eD_f{\cal F}[H_e,K_f]
 =\operatorname{Tr}
 \left(Q_{ef}K_fQ_{fe}H_e\right).                    \tag{69.13}
\]
In particular,
\[
 \sum_f
 |D_eD_f{\cal F}[H_e,K_f]|
 \leq
 \frac{r_g}{a_-^2}
 \|H_e\|_{\rm op}\sup_f\|K_f\|_{\rm op}.             \tag{69.14}
\]
The full Hessian of \({\cal F}\) is positive semidefinite.
\end{theorem}

\begin{proof}
Equation (69.12) follows from (69.7) and
\[
 |\operatorname{Tr}(Q_{ee}H_e)|
 \leq\operatorname{Tr}(Q_{ee})\|H_e\|_{\rm op}.
\]
Differentiate \({\bf G}\) as in (67.9) to obtain (69.13).  Then
\[
 |\operatorname{Tr}(Q_{ef}K_fQ_{fe}H_e)|
 \leq
 \|H_e\|_{\rm op}\|K_f\|_{\rm op}
 \|Q_{ef}\|_{\rm HS}^2,
\]
and (69.8) proves (69.14).  For a general variation,
\[
 D^2{\cal F}[{\bf H},{\bf H}]
 =\operatorname{Tr}\left[
 ({\bf G}^{1/2}{\bf D}^{\mathsf T}{\bf H}
 {\bf D}{\bf G}^{1/2})^2\right]\geq0.                \tag{69.15}
\]
\end{proof}

The internal fibre costs only the explicit factor \(r_g\).  There is
still no dependence on the number of vertices, edges, or the
massless spectral gap.

\subsection{Nonnegative lower-order terms}

Consider
\[
 {\bf M}={\bf D}^{\mathsf T}{\bf A}{\bf D}+W,         \tag{69.16}
\]
where \(W=W^{\mathsf T}\succeq0\) on \({\cal H}_0\), and let
\({\bf G}_W={\bf M}^{-1}\).  Define
\[
 {\bf P}_W
 ={\bf A}^{1/2}{\bf D}{\bf G}_W
 {\bf D}^{\mathsf T}{\bf A}^{1/2}.                   \tag{69.17}
\]

\begin{proposition}[Projection becomes a contraction]
\label{prop:v16v69-contraction}
One has
\[
 0\preceq{\bf P}_W\preceq I,\qquad
 {\bf P}_W^2\preceq{\bf P}_W.                        \tag{69.18}
\]
Consequently the row estimates (69.6)--(69.8) remain valid with
\({\bf Q}\) replaced by
\({\bf D}{\bf G}_W{\bf D}^{\mathsf T}\).
\end{proposition}

\begin{proof}
Since
\({\bf M}\succeq{\bf D}^{\mathsf T}{\bf A}{\bf D}\),
the variational characterization of the inverse on the range gives
\({\bf P}_W\preceq{\bf P}\preceq I\), where \({\bf P}\) is (69.5).
Positivity is immediate.  A positive contraction satisfies
\({\bf P}_W^2\preceq{\bf P}_W\).  Hence
\[
 \sum_f\|(P_W)_{ef}\|_{\rm HS}^2
 =\operatorname{Tr}(P_W^2)_{ee}
 \leq\operatorname{Tr}(P_W)_{ee}\leq r_g,
\]
and the proof of (69.7)--(69.8) applies verbatim.
\end{proof}

Thus a symmetric nonnegative curvature or gauge-pinning potential does
not destroy the volume-uniform response mechanism.

\subsection{Sign-indefinite symmetric perturbations}

Suppose instead that \(W=W^{\mathsf T}\) satisfies the relative form
bound
\[
 W\succeq-\eta\,{\bf D}^{\mathsf T}{\bf A}{\bf D},
 \qquad 0\leq\eta<1.                                  \tag{69.19}
\]
Then
\[
 {\bf M}\succeq(1-\eta)
 {\bf D}^{\mathsf T}{\bf A}{\bf D}.                  \tag{69.20}
\]

\begin{proposition}[Relative negative lower-order load]
\label{prop:v16v69-relative}
Under (69.19),
\[
 0\preceq{\bf P}_W\preceq\frac1{1-\eta}I.             \tag{69.21}
\]
The force bound (69.12) is multiplied by
\((1-\eta)^{-1}\), and the absolute Hessian row bound (69.14) is
multiplied by \((1-\eta)^{-2}\).
\end{proposition}

\begin{proof}
The inverse form inequality from (69.20) gives (69.21).  A positive
operator bounded by \(cI\) satisfies
\({\bf P}_W^2\preceq c{\bf P}_W\preceq c^2I\).
Thus the diagonal trace is at most \(cr_g\), and the block-square row
is at most \(c^2r_g\).  Insert \(c=(1-\eta)^{-1}\) in the proofs of
(69.12) and (69.14).
\end{proof}

This is the precise stability corridor for a symmetric lower-order
term: it may be negative, but only with a relative bound strictly
below one.

\subsection{What nonsymmetry destroys}

Let
\[
 {\bf M}={\bf L}+K,\qquad K\ne K^{\mathsf T}.          \tag{69.22}
\]
Then \({\bf M}^{-1}\) need not be symmetric, and
\[
 {\bf A}^{1/2}{\bf D}{\bf M}^{-1}
 {\bf D}^{\mathsf T}{\bf A}^{1/2}                    \tag{69.23}
\]
is neither an orthogonal projection nor a positive contraction in
general.  Moreover,
\[
 \operatorname{Tr}
 ({\bf M}^{-1}\delta{\bf M}
  {\bf M}^{-1}\delta{\bf M})                          \tag{69.24}
\]
has no fixed sign.  Therefore both the projection row identity and
the Gram convexity used in V67--V68 fail.

\begin{proposition}[Small nonsymmetric relative perturbation]
\label{prop:v16v69-nonsymmetric}
On \({\cal H}_0\), suppose
\[
 {\bf M}={\bf L}^{1/2}(I+E){\bf L}^{1/2},
 \qquad \|E\|_{\rm op}\leq\eta<1.                    \tag{69.25}
\]
Then \({\bf M}\) is invertible,
\[
 \|{\bf L}^{1/2}{\bf M}^{-1}{\bf L}^{1/2}\|
 \leq(1-\eta)^{-1},                                  \tag{69.26}
\]
and, for real \(E\),
\[
 \det{}_{{\cal H}_0}{\bf M}>0.                       \tag{69.27}
\]
However no convexity conclusion follows without symmetry.
\end{proposition}

\begin{proof}
The Neumann series gives
\((I+E)^{-1}\) with norm at most \((1-\eta)^{-1}\), proving
invertibility and (69.26).  Every real eigenvalue of \(I+E\) is
positive because its distance from \(1\) is less than one; nonreal
eigenvalues occur in conjugate pairs with positive product.  Hence
\(\det(I+E)>0\), and \(\det{\bf L}>0\) on \({\cal H}_0\), proving
(69.27).  The sign failure in (69.24) for nonsymmetric matrices
precludes the last conclusion.
\end{proof}

The relative smallness in (69.25) must be uniform in volume.  A
first-order drift compared with a massless second-order operator often
fails this at the lowest modes, so (69.25) is a real hypothesis, not a
formal order-counting argument.

\subsection{Reading for the gravitational gauge operator}

At a flat background, V62's residual volume-preserving
Faddeev--Popov symbol is \(|k|^2I_{k^\perp}\), matching the principal
part of (69.2) after removal of constant modes.  On a curved
background, the candidate operator contains:
\begin{enumerate}
\item a vector divergence-form principal part, compatible with
      matrix conductances;
\item curvature and projection terms whose symmetry and relative sign
      must be computed in the chosen gauge;
\item possible first-order nonsymmetric terms, which must either be
      removed by a weighted inner product or satisfy (69.25).
\end{enumerate}
Only after this decomposition is explicit can the factor
\((1-\eta)^{-2}r_g/a_-^2\) be entered in the V68/V66 influence ledger.

\begin{assessment}[What V69 closes]
\label{ass:v16v69-status}
V69 extends the massless projection mechanism to vector gauge fields
with positive matrix conductances.  It proves force and Hessian-row
bounds with explicit fibre factor \(r_g\), shows that symmetric
nonnegative lower-order terms preserve the result, and quantifies
sign-indefinite symmetric terms by a relative bound.

Nonsymmetric terms destroy both projection and Gram positivity.
A small relative perturbation preserves invertibility and determinant
orientation but not convexity.  The actual curved residual
diffeomorphism operator has not yet been decomposed into these pieces.
\end{assessment}

\begin{decision}[V70 companion audit and curved residual operator]
\label{dec:v16v69-next}
Perform the mandatory YM/NS companion audit.  Then compute the
continuum linearization of a concrete residual
volume-preserving gauge condition about a smooth unimodular
background.  Identify its vector Laplacian, Ricci term, projection,
formal adjoint, and kernel, and decide whether a weighted inner product
makes the operator symmetric with a relative lower-order bound.
\end{decision}

\begin{firewall}[Claim boundary after V69]
\label{fire:v16v69}
The matrix graph model does not yet identify the nonlinear
gravitational Faddeev--Popov operator.  The relative perturbation
bounds, determinant sign, gauge atlas, one-block gaps, reflection
positivity, and continuum Ward limits remain unproved.
\end{firewall}

\section*{Version V69 append-only record}
\addcontentsline{toc}{section}{Version V69 append-only record}
The complete V68 source was retained byte for byte before this
continuation.  It had \(705371\) bytes, \(21059\) lines, and SHA--256
\[
 \texttt{E56B4C333B5EA5109C0F2FF8BDE631FF801EC4CC4C5439D1A21F0F1446511B25}.
\]
V69 extended the conductance projection and determinant-response
bounds to vector fibres and classified symmetric, sign-indefinite,
and nonsymmetric lower-order perturbations.


\section{V70: companion audit and the curved unimodular residual operator}
\label{sec:v16v70}

\subsection{Mandatory YM/NS companion audit}

\begin{audit}[V70 companion checkpoint]
\label{aud:v16v70-companions}
The newest companion files inspected were
\[
\begin{array}{c|r|r|l}
\hbox{file}&\hbox{lines}&\hbox{bytes}&\hbox{SHA--256}\\ \hline
\texttt{Renewal\_Yang\_Mills\_Mass\_Gap\_Research\_Notes\_V56.tex}
&13172&483563&
\texttt{9A728CAE9E3651DA670139B8B6D5F35E29E4A7243E1870DD0E4793F2521958F5}\\
\texttt{Renewal\_NS\_Stability\_Research\_Notes\_V26.tex}
&5319&278283&
\texttt{82C887F8C2C69E4F28FAD7C3DC8DE5B9099CB5A4BF431EBA1FFF3E91935978AD}.
\end{array}                                             \tag{70.1}
\]
\end{audit}

YM V56 remains at its exact overlapping-chart repair of four
third-axis certificate defects; NS V26 remains at its rank-one Oseen
kernel norms.  Neither contains a curved metric gauge operator,
Korn inequality, Hodge projection, or metric Faddeev--Popov
determinant.

\begin{theorem}[V70 import decision]
\label{thm:v16v70-audit-decision}
No companion estimate enters the curved residual-gauge calculation.
YM's certified-atlas logic remains relevant only to the later problem
of covering a nonlinear gauge region by oriented local charts.
\end{theorem}

\begin{proof}
The companion variables and operators do not verify any analytic
hypothesis below.  The atlas observation is a construction pattern,
not a transferred theorem.
\end{proof}

\subsection{Unimodular tangent and residual diffeomorphisms}

Let \((X,g)\) be a compact smooth Riemannian \(d\)-manifold without
boundary.  Define
\[
 {\cal S}_0
 =\{h\in C^\infty({\rm Sym}^2T^*X):
   \operatorname{tr}_g h=0\},                         \tag{70.2}
\]
and
\[
 {\cal V}_T
 =\{\xi\in C^\infty(T^*X):\delta\xi=0\},              \tag{70.3}
\]
where \(\delta\xi=-\nabla^\mu\xi_\mu\).  The symmetrized derivative
and its formal adjoint are
\[
 (\delta^*\xi)_{\mu\nu}
 =\frac12(\nabla_\mu\xi_\nu+\nabla_\nu\xi_\mu),
 \qquad
 (\delta h)_\nu=-\nabla^\mu h_{\mu\nu}.               \tag{70.4}
\]
For \(\xi\in{\cal V}_T\),
\[
 \operatorname{tr}_g(2\delta^*\xi)
 =2\nabla^\mu\xi_\mu=0,                               \tag{70.5}
\]
so \(2\delta^*\xi={\cal L}_{\xi^\sharp}g\) is tangent to the local
determinant slice.

Let \(P_T\) be the \(L^2(g)\)-orthogonal Hodge projection onto the
closure of \({\cal V}_T\).  Use the residual gauge condition
\[
 {\cal F}_g(h)=P_T\delta h=0,\qquad h\in{\cal S}_0.   \tag{70.6}
\]

\subsection{Linearized Faddeev--Popov operator}

\begin{theorem}[Exact residual operator]
\label{thm:v16v70-operator}
The derivative of (70.6) along a residual infinitesimal
diffeomorphism is
\[
 {\cal M}_g\xi
 =2P_T\delta\delta^*\xi,\qquad \xi\in{\cal V}_T.      \tag{70.7}
\]
It is a self-adjoint nonnegative operator on \({\cal V}_T\), and
\[
 \langle\xi,{\cal M}_g\xi\rangle_{L^2}
 =2\|\delta^*\xi\|_{L^2}^2
 =\frac12\|{\cal L}_{\xi^\sharp}g\|_{L^2}^2.          \tag{70.8}
\]
Its kernel is exactly the space of Killing one-forms:
\[
 \ker{\cal M}_g
 =\{\xi:\delta^*\xi=0\}.                              \tag{70.9}
\]
\end{theorem}

\begin{proof}
The gauge-orbit tangent is \(h=2\delta^*\xi\), so differentiating
(70.6) gives (70.7).  For \(\xi,\zeta\in{\cal V}_T\),
\[
 \langle\zeta,{\cal M}_g\xi\rangle
 =2\langle\zeta,\delta\delta^*\xi\rangle
 =2\langle\delta^*\zeta,\delta^*\xi\rangle,           \tag{70.10}
\]
because \(P_T\zeta=\zeta\) and \(\delta\) is the adjoint of
\(\delta^*\).  This proves self-adjointness and (70.8).
The right side of (70.8) vanishes exactly when
\(\delta^*\xi=0\), which is the Killing equation.
\end{proof}

Thus the Ricci term that appears after commuting covariant
derivatives does not create a negative determinant operator.  The
factorization (70.7) is the authoritative sign statement.

\subsection{Rough-Laplacian and Ricci form}

Use the positive rough Laplacian
\[
 \nabla^*\nabla=-\nabla^\mu\nabla_\mu.                \tag{70.11}
\]

\begin{proposition}[Curvature decomposition]
\label{prop:v16v70-curvature}
On divergence-free one-forms,
\[
 2\delta\delta^*\xi
 =\nabla^*\nabla\xi-\operatorname{Ric}(\xi),          \tag{70.12}
\]
and therefore
\[
 {\cal M}_g
 =P_T(\nabla^*\nabla-\operatorname{Ric})P_T
 \quad\text{on }{\cal V}_T.                           \tag{70.13}
\]
\end{proposition}

\begin{proof}
From (70.4),
\[
 (2\delta\delta^*\xi)_\nu
 =-\nabla^\mu\nabla_\mu\xi_\nu
  -\nabla^\mu\nabla_\nu\xi_\mu.                      \tag{70.14}
\]
The commutator identity for a one-form gives
\[
 \nabla^\mu\nabla_\nu\xi_\mu
 =\nabla_\nu\nabla^\mu\xi_\mu
  +\operatorname{Ric}_{\nu\rho}\xi^\rho.             \tag{70.15}
\]
The first term vanishes when \(\delta\xi=0\).  Equations
(70.11), (70.14), and (70.15) yield (70.12); compression by \(P_T\)
gives (70.13).
\end{proof}

On a positive-Ricci background, the second term in (70.12) is
negative if read separately.  Identity (70.8) shows why applying the
relative-lower-order test of V69 directly to \(-\operatorname{Ric}\)
can be unnecessarily pessimistic: the full operator is a square.

\subsection{Korn gap and zero-mode removal}

Let \({\cal K}_g=\ker\delta^*\) and work on
\[
 {\cal V}_T^\perp
 ={\cal V}_T\cap{\cal K}_g^\perp.                    \tag{70.16}
\]

\begin{theorem}[Residual invertibility is a Korn inequality]
\label{thm:v16v70-korn}
The following statements are equivalent for \(\lambda_K>0\):
\begin{align}
 2\|\delta^*\xi\|_{L^2}^2
 &\geq\lambda_K\|\xi\|_{L^2}^2
 \quad(\xi\in{\cal V}_T^\perp),                       \tag{70.17}\\
 {\cal M}_g|_{{\cal V}_T^\perp}
 &\succeq\lambda_KI.                                 \tag{70.18}
\end{align}
When they hold,
\[
 \|{\cal M}_g^{-1}\|_{L^2\to L^2}
 \leq\lambda_K^{-1}.                                 \tag{70.19}
\]
\end{theorem}

\begin{proof}
Equivalence follows from the quadratic identity (70.8) and the
min--max characterization of the bottom positive eigenvalue.
The inverse bound is the spectral theorem.
\end{proof}

On a flat torus, constant translations are precisely the obvious
Killing zero modes, and for a nonzero divergence-free Fourier mode
(70.17) has \(\lambda_K=|k|^2\).  Consequently the global inverse
again loses volume uniformity.  V67--V69 show, however, that a
factorized determinant response may remain uniform even when
\(\lambda_K\to0\).

\subsection{Abstract factor projection}

The graph incidence structure was not essential to the projection
identity.

\begin{proposition}[General square-factor projection]
\label{prop:v16v70-general-factor}
Let \(B:{\cal H}\to{\cal K}\) be injective on the declared
zero-mode complement, and let \(A\succ0\) on \({\cal K}\).  Put
\[
 M=B^*AB,\qquad
 P=A^{1/2}BM^{-1}B^*A^{1/2}.                         \tag{70.20}
\]
Then \(P\) is the orthogonal projection onto
\(\overline{\operatorname{Ran}(A^{1/2}B)}\).
\end{proposition}

\begin{proof}
It is symmetric, and
\[
 P^2
 =A^{1/2}B M^{-1}(B^*AB)M^{-1}B^*A^{1/2}
 =P.
\]
Its range is visibly contained in
\(\overline{\operatorname{Ran}(A^{1/2}B)}\), and it acts as the
identity on that range.
\end{proof}

At a fixed background, (70.7) fits (70.20) with
\[
 B=\sqrt2\,\delta^*|_{{\cal V}_T^\perp},
 \qquad A=I.                                         \tag{70.21}
\]
Thus its fixed-background factor response has the same projection
mechanism as the vector graph model.  The new difficulty is that a
metric variation changes \(B\), the \(L^2\) inner products, the Hodge
projection, and the zero-mode subspace simultaneously.

\subsection{Moving-subspace ledger}

A nonlinear determinant calculation must choose identifications
\[
 U_g:{\cal V}_{T,g}^\perp\longrightarrow{\cal H}_{\rm ref}       \tag{70.22}
\]
with a fixed reference Hilbert space.  If \(U_g\) is unitary for the
metric-dependent \(L^2\) products, then
\[
 \widetilde{\cal M}_g
 =U_g{\cal M}_gU_g^{-1}                               \tag{70.23}
\]
remains self-adjoint and nonnegative.  Differentiating it produces
\[
 D\widetilde{\cal M}
 =(DU){\cal M}U^{-1}
 +U(D{\cal M})U^{-1}
 -U{\cal M}U^{-1}(DU)U^{-1}.                         \tag{70.24}
\]
The first and third terms form a commutator.  Their trace against
\(\widetilde{\cal M}^{-1}\) cancels formally:
\[
 \operatorname{Tr}\left[
 \widetilde{\cal M}^{-1}
 [\, (DU)U^{-1},\widetilde{\cal M}\,]\right]=0,       \tag{70.25}
\]
whenever cyclicity is justified.  At second order, trace-ideal
control is still required.  A nonunitary identification can introduce
spurious nonsymmetry and lose the V69 square structure.

\begin{assessment}[What V70 closes]
\label{ass:v16v70-status}
The companion audit supplies no curved-gauge estimate.  V70 computes
the residual volume-preserving Faddeev--Popov operator exactly:
it is the Hodge compression of \(2\delta\delta^*\), is self-adjoint
and nonnegative, and has only Killing fields in its kernel.  Its
invertibility is precisely a Korn inequality.

The apparent negative Ricci potential is controlled by the exact
square factorization.  The V69 projection extends to this factor at a
fixed background.  Metric variation still moves the factor, inner
product, Hodge projection, and Killing complement; those response
terms are not yet bounded.
\end{assessment}

\begin{decision}[V71 moving-factor determinant response]
\label{dec:v16v70-next}
Differentiate a finite-dimensional family
\(M(g)=B(g)^*A(g)B(g)\) after unitary trivialization.  Separate the
conductance response \(DA\), the factor response \(DB\), and the
commutator from the moving Hilbert structure.  Prove exact
first-variation cancellation for the commutator and quantify the
remaining logarithmic determinant loads using a right inverse of
\(B\), without assuming a volume-uniform Korn gap where derivative
cancellation suffices.
\end{decision}

\begin{firewall}[Claim boundary after V70]
\label{fire:v16v70}
The calculation is the linearization of one residual gauge condition
at a smooth background.  It does not construct a nonlinear gauge
atlas, prove uniform Korn constants, control moving Hodge projections,
establish determinant reflection positivity, or complete the
continuum measure.
\end{firewall}

\section*{Version V70 append-only record}
\addcontentsline{toc}{section}{Version V70 append-only record}
The complete V69 source was retained byte for byte before this
continuation.  It had \(716306\) bytes, \(21405\) lines, and SHA--256
\[
 \texttt{DE1953B67826D73A1E838EFA07C5545974644BA62A9859DAD8B8F6C58AF19BB1}.
\]
V70 recorded the mandatory companion audit, derived the curved
unimodular residual operator and Korn gap, and lifted the graph
projection to its general square-factor form.


\section{V71: moving-factor determinant response}
\label{sec:v16v71}

\subsection{Finite-dimensional square family}

Let \({\cal H}\) and \({\cal K}\) be fixed finite-dimensional real
Hilbert spaces.  Let
\[
 B(x):{\cal H}\to{\cal K},\qquad
 A(x)=A(x)^*\succ0,                                   \tag{71.1}
\]
be \(C^2\) families, and assume \(B(x)\) is injective throughout the
chart.  Define
\[
 M(x)=B(x)^*A(x)B(x),\qquad
 {\cal F}(x)=-\log\det M(x).                          \tag{71.2}
\]
Put
\[
 R=M^{-1}B^*A,\qquad
 Q=BM^{-1}B^*.                                       \tag{71.3}
\]
Then
\[
 RB=I_{\cal H}.                                      \tag{71.4}
\]

\subsection{Exact first response}

\begin{theorem}[Conductance and factor split]
\label{thm:v16v71-first}
For every parameter direction \(v\),
\[
 D\log\det M[v]
 =
 \operatorname{Tr}_{\cal K}\bigl(Q\,DA[v]\bigr)
 +2\operatorname{Tr}_{\cal H}\bigl(R\,DB[v]\bigr).   \tag{71.5}
\]
Equivalently,
\[
 D{\cal F}[v]
 =-\operatorname{Tr}(Q\,DA[v])
  -2\operatorname{Tr}(R\,DB[v]).                     \tag{71.6}
\]
\end{theorem}

\begin{proof}
Differentiate (71.2):
\[
 DM[v]
 =DB[v]^*AB+B^*DA[v]B+B^*ADB[v].                    \tag{71.7}
\]
Jacobi's formula gives
\(\operatorname{Tr}(M^{-1}DM[v])\).  Cyclicity turns the middle term
into \(\operatorname{Tr}(QDA[v])\).  The last term is
\(\operatorname{Tr}(RDB[v])\).  The first term is its transpose and
has the same real trace.  This proves (71.5)--(71.6).
\end{proof}

The response orientation is explicit: the right inverse \(R\) must
stand to the left of \(DB\).  A norm bound on \(DB\) and a separate
bound on \(M^{-1}\) would unnecessarily insert the inverse Korn gap.

\subsection{Derivative cancellation inside the factor}

\begin{proposition}[Right-factor cancellation]
\label{prop:v16v71-right-factor}
Suppose
\[
 DB[v]=B S_v+N_v                                     \tag{71.8}
\]
for an endomorphism \(S_v\) of \({\cal H}\) and a defect
\(N_v:{\cal H}\to{\cal K}\).  Then
\[
 RDB[v]=S_v+RN_v,                                    \tag{71.9}
\]
and
\[
 D\log\det M[v]
 =\operatorname{Tr}(QDA[v])
  +2\operatorname{Tr}(S_v)
  +2\operatorname{Tr}(RN_v).                         \tag{71.10}
\]
In particular, when \(N_v=0\), the factor response contains no inverse
Korn constant.
\end{proposition}

\begin{proof}
Multiply (71.8) by \(R\) and use (71.4), then substitute the result
in (71.5).
\end{proof}

The same mechanism includes a change of domain frame:
\(B\mapsto BU^{-1}\) differentiates as \(DB=-B\,U^{-1}DU\,U^{-1}\)
at fixed geometric factor, so it is entirely of the \(BS_v\) form.

\subsection{Unitary trivialization creates no determinant load}

Let \(U(x)\) be an orthogonal family on \({\cal H}\), and set
\[
 \widetilde M(x)=U(x)M(x)U(x)^{-1}.                  \tag{71.11}
\]

\begin{theorem}[Exact commutator cancellation]
\label{thm:v16v71-unitary}
At every finite cutoff,
\[
 \det\widetilde M(x)=\det M(x).                       \tag{71.12}
\]
Consequently all derivatives of their logarithmic determinants agree.
At first order, the additional term is the commutator
\[
 [K_v,\widetilde M],\qquad K_v=DU[v]U^{-1},           \tag{71.13}
\]
and
\[
 \operatorname{Tr}\left(
 \widetilde M^{-1}[K_v,\widetilde M]\right)=0.        \tag{71.14}
\]
\end{theorem}

\begin{proof}
Similarity invariance of the determinant proves (71.12), hence the
all-order statement.  Differentiating (71.11) gives (71.13) in
addition to the conjugated geometric derivative.  Cyclicity gives
\[
 \operatorname{Tr}(\widetilde M^{-1}K_v\widetilde M)
 -\operatorname{Tr}(K_v)=0,
\]
which is (71.14).
\end{proof}

Thus the moving \(L^2\) Hilbert structure should be trivialized
unitarily.  Treating its commutator terms separately by absolute
operator norms would create a spurious determinant force.  At an
infinite-dimensional limit, this cancellation must be preserved by
the chosen regularized determinant; V71 proves it only at the finite
cutoffs.

\subsection{Whitened second response}

Define
\[
 X_v=M^{-1/2}DM[v]M^{-1/2}.                           \tag{71.15}
\]
It is self-adjoint for real \(v\).

\begin{theorem}[Exact second-response identity]
\label{thm:v16v71-second}
For parameter directions \(v,w\),
\[
 D^2\log\det M[v,w]
 =
 \operatorname{Tr}\left(
 M^{-1/2}D^2M[v,w]M^{-1/2}\right)
 -\operatorname{Tr}(X_vX_w).                         \tag{71.16}
\]
Consequently,
\[
 |D^2\log\det M[v,w]|
 \leq
 \|M^{-1/2}D^2M[v,w]M^{-1/2}\|_{\mathfrak S_1}
 +\|X_v\|_{\mathfrak S_2}\|X_w\|_{\mathfrak S_2}.    \tag{71.17}
\]
\end{theorem}

\begin{proof}
Differentiate
\(\operatorname{Tr}(M^{-1}DM[v])\) in direction \(w\), using
\(D(M^{-1})[w]=-M^{-1}DM[w]M^{-1}\).  Cyclicity rewrites the quadratic
term as \(\operatorname{Tr}(X_wX_v)\), proving (71.16).
Trace norm and Schatten H\"older give (71.17).
\end{proof}

The normalized versions of the two norms in (71.17) are the exact
V65 owners.  They can remain volume uniform even when
\(\|M^{-1}\|\) diverges, provided derivative cancellation is retained
inside the whitened composites.

\subsection{Whitened factor orientation}

Put
\[
 C=A^{1/2}B,\qquad
 W=CM^{-1/2}.                                        \tag{71.18}
\]
Then
\[
 W^*W=I_{\cal H},                                    \tag{71.19}
\]
so \(W\) is an isometry.  For a direction \(v\), define
\[
 Z_v=W^*DC[v]M^{-1/2}.                               \tag{71.20}
\]

\begin{proposition}[Gram whitening]
\label{prop:v16v71-whitening}
One has
\[
 X_v=Z_v+Z_v^*,\qquad
 D\log\det M[v]=2\operatorname{Tr}Z_v,               \tag{71.21}
\]
and
\[
 \|X_v\|_{\mathfrak S_2}\leq2\|Z_v\|_{\mathfrak S_2}.             \tag{71.22}
\]
\end{proposition}

\begin{proof}
Since \(M=C^*C\),
\[
 DM[v]=DC[v]^*C+C^*DC[v].
\]
Conjugate by \(M^{-1/2}\) and use (71.18)--(71.20) to obtain the first
identity.  Taking the trace gives the second; the triangle inequality
gives (71.22).
\end{proof}

This is the square-factor analogue of V51's oriented
energy-response amplitude.  The inverse appears only in the composite
\(DC\,M^{-1/2}\), and the isometry \(W\) costs no norm.

\subsection{A response corridor without a uniform Korn inverse}

Let \(N\) be the number of physical cells and use the normalized
Schatten norms (65.2).

\begin{theorem}[Moving-factor determinant corridor]
\label{thm:v16v71-corridor}
Suppose, uniformly in the cutoff and unit local parameter directions,
\begin{align}
 \|Q\,DA[v]\|_{1,N}&\leq a_1,                         \tag{71.23}\\
 \|S_v\|_{1,N}+\|RN_v\|_{1,N}&\leq b_1,              \tag{71.24}\\
 \|M^{-1/2}D^2M[v,w]M^{-1/2}\|_{1,N}&\leq c_2,       \tag{71.25}\\
 \|Z_v\|_{2,N}&\leq z_2.                              \tag{71.26}
\end{align}
Then
\[
 \frac1N|D\log\det M[v]|
 \leq a_1+2b_1,                                      \tag{71.27}
\]
and
\[
 \frac1N|D^2\log\det M[v,w]|
 \leq c_2+4z_2^2.                                    \tag{71.28}
\]
No separate bound on \(\|M^{-1}\|\) occurs.
\end{theorem}

\begin{proof}
Apply the normalized trace inequality to (71.10) for (71.27).
Equations (71.17) and (71.22), divided by \(N\), give (71.28).
\end{proof}

The theorem does not make a diverging inverse harmless by assertion.
It replaces that inverse by four oriented composites that must be
proved finite.  In particular, a factor defect \(N_v\) with a
component in the lowest Korn modes can still make \(RN_v\) diverge.

\subsection{Kernel-stratum obstruction}

The preceding derivatives require a fixed domain dimension and an
injective \(B\).  For the curved residual operator, this means that
the dimension of the Killing space must remain fixed on the chart.

\begin{proposition}[Kernel jumps prevent a smooth log determinant]
\label{prop:v16v71-kernel}
If a continuous family \(M(t)\succeq0\) has an eigenvalue
\(\lambda(t)>0\) for \(t\ne0\) with \(\lambda(t)\to0\) as \(t\to0\),
then
\[
 \log\det M(t)\longrightarrow-\infty                 \tag{71.29}
\]
unless that eigenvalue is removed from the determinant on the entire
family.  Removing it only at \(t=0\) changes the determinant dimension
and does not define a continuous logarithm.
\end{proposition}

\begin{proof}
The determinant is the product of its positive eigenvalues.  If all
other eigenvalues remain bounded above and below, the factor
\(\lambda(t)\) forces (71.29); if another eigenvalue also degenerates,
the same conclusion holds more strongly.  A determinant on a
different-dimensional complement is a different function.
\end{proof}

Hence a nonlinear gauge atlas must be stratified by stabilizer type,
or must remove a fixed symmetry space declared throughout the chart.
Generic metrics and highly symmetric backgrounds cannot be joined by
one smooth determinant chart without an additional treatment of the
emerging Killing modes.

\begin{assessment}[What V71 closes]
\label{ass:v16v71-status}
V71 differentiates a general moving square factor and separates
conductance, factor, and Hilbert-trivialization responses.  Unitary
commutators cancel exactly at finite cutoff.  A right-factor
decomposition \(DB=BS+N\) isolates the only part that needs the
potentially large right inverse, and whitened trace-ideal bounds give
a determinant corridor with no separate Korn-gap constant.

The remaining geometric owners are \(RN_v\),
\(M^{-1/2}D^2MM^{-1/2}\), and \(Z_v\).  Kernel-dimension jumps at
symmetric metrics prevent a single smooth determinant chart.
\end{assessment}

\begin{decision}[V72 variation of the Killing operator]
\label{dec:v16v71-next}
For \(B_g=\sqrt2\,\delta_g^*\) after unitary trivialization, compute
\(DB_g[h]\xi\) in local coordinates.  Split the terms containing
\(\nabla\xi\) into left/right factor responses and isolate the
zeroth-order defect proportional to \(\nabla h\,\xi\).  Determine its
pairing with low Korn modes and state the Sobolev or neutrality
condition needed to control \(R_gN_h\).
\end{decision}

\begin{firewall}[Claim boundary after V71]
\label{fire:v16v71}
The moving-factor theorem is finite-dimensional and conditional on
fixed kernel rank and trace-ideal bounds.  It does not establish those
bounds for \(\delta_g^*\), construct a nonlinear gauge atlas, or prove
the positive continuum law and Ward identities.
\end{firewall}

\section*{Version V71 append-only record}
\addcontentsline{toc}{section}{Version V71 append-only record}
The complete V70 source was retained byte for byte before this
continuation.  It had \(726909\) bytes, \(21722\) lines, and SHA--256
\[
 \texttt{BB7C56AF95033A5ED76FAA0C60D9AA12A5B1B993C8B2093F15F8DCA51894B94A}.
\]
V71 proved the moving-square-factor response identities, exact
unitary commutator cancellation, a trace-ideal determinant corridor,
and the fixed-kernel-stratum requirement.


\section{V72: variation of the Killing factor and the low-mode defect}
\label{sec:v16v72}

\subsection{Exact geometric variation}

Represent the residual gauge parameter as a vector field \(\xi\), and
write the square factor from V70 as
\[
 B_g\xi=\sqrt2\,\delta_g^*(\xi^\flat)
 =\frac1{\sqrt2}{\cal L}_\xi g.                       \tag{72.1}
\]
Before the unitary Hilbert-space trivialization of V71, a metric
variation \(h\) gives the exact derivative
\[
 D_gB[h]\xi=\frac1{\sqrt2}{\cal L}_\xi h.             \tag{72.2}
\]

\begin{proposition}[Local-coordinate factor variation]
\label{prop:v16v72-variation}
In covariant components,
\[
 (D_gB[h]\xi)_{\mu\nu}
 =\frac1{\sqrt2}\left(
 \xi^\rho\nabla_\rho h_{\mu\nu}
 +h_{\rho\nu}\nabla_\mu\xi^\rho
 +h_{\mu\rho}\nabla_\nu\xi^\rho\right).              \tag{72.3}
\]
Equivalently,
\[
 D_gB[h]={\cal P}_h+N_h,                              \tag{72.4}
\]
where
\begin{align}
 ({\cal P}_h\xi)_{\mu\nu}
 &=\frac1{\sqrt2}\left(
 h_{\rho\nu}\nabla_\mu\xi^\rho
 +h_{\mu\rho}\nabla_\nu\xi^\rho\right),              \tag{72.5}\\
 (N_h\xi)_{\mu\nu}
 &=\frac1{\sqrt2}\xi^\rho\nabla_\rho h_{\mu\nu}.      \tag{72.6}
\end{align}
\end{proposition}

\begin{proof}
The identity
\[
 ({\cal L}_\xi h)_{\mu\nu}
 =\xi^\rho\nabla_\rho h_{\mu\nu}
  +h_{\rho\nu}\nabla_\mu\xi^\rho
  +h_{\mu\rho}\nabla_\nu\xi^\rho
\]
is the covariant formula for the Lie derivative of a covariant
two-tensor.  Insert it in (72.2).
\end{proof}

The unitary trivialization terms omitted from (72.2) are commutators
at the level of \(M=B^*B\); their determinant response cancels by
Theorem~\ref{thm:v16v71-unitary}.  The geometric terms
(72.5)--(72.6) remain.

\subsection{Principal and zeroth-order estimates}

\begin{proposition}[Raw Sobolev bounds]
\label{prop:v16v72-raw}
There is a dimensional constant \(c_d\) such that
\[
 \|{\cal P}_h\xi\|_{L^2}
 \leq c_d\|h\|_{L^\infty}\|\nabla\xi\|_{L^2},         \tag{72.7}
\]
and
\[
 \|N_h\xi\|_{L^2}
 \leq c_d\|\nabla h\|_{L^\infty}\|\xi\|_{L^2}.        \tag{72.8}
\]
\end{proposition}

\begin{proof}
Each component in (72.5) is one bounded coefficient \(h\) times one
component of \(\nabla\xi\); each component in (72.6) is one bounded
coefficient \(\nabla h\) times one component of \(\xi\).  Pointwise
Cauchy--Schwarz followed by integration gives the estimates.
\end{proof}

Let \(K_{\rm Ric}\geq0\) satisfy
\[
 \operatorname{Ric}\succeq-K_{\rm Ric}g.             \tag{72.9}
\]
On the Killing-orthogonal residual space, let the Korn gap be
\(\lambda_K\) as in (70.17).

\begin{lemma}[Korn control of the principal response]
\label{lem:v16v72-korn-control}
For residual \(\xi\perp{\cal K}_g\),
\[
 \|\nabla\xi\|_{L^2}^2
 \leq
 2\left(1+\frac{K_{\rm Ric}}{\lambda_K}\right)
 \|\delta^*\xi\|_{L^2}^2.                            \tag{72.10}
\]
\end{lemma}

\begin{proof}
The integrated form of (70.12) gives
\[
 2\|\delta^*\xi\|_{L^2}^2
 =\|\nabla\xi\|_{L^2}^2
  -\int\operatorname{Ric}(\xi,\xi)\,d{\rm vol}_g.
                                                               \tag{72.11}
\]
Hence
\[
 \|\nabla\xi\|_{L^2}^2
 \leq2\|\delta^*\xi\|_{L^2}^2
     K_{\rm Ric}\|\xi\|_{L^2}^2.
\]
The Korn inequality
\(2\|\delta^*\xi\|^2\geq\lambda_K\|\xi\|^2\)
gives (72.10).
\end{proof}

Even the derivative-bearing response (72.5) therefore has a
gap-independent factor bound only when either \(K_{\rm Ric}=0\) or a
stronger volume-uniform Korn estimate is available.  The zeroth-order
term (72.6) is more severe after application of the right inverse.

\subsection{Exact right-inverse form of the defect}

At a fixed finite cutoff let
\[
 M=B^*B,\qquad R=M^{-1}B^*,\qquad
 W=BM^{-1/2}.                                        \tag{72.12}
\]
Then
\[
 R=M^{-1/2}W^*.                                      \tag{72.13}
\]

\begin{theorem}[Neutrality criterion]
\label{thm:v16v72-neutrality}
For an operator \(N:{\cal H}\to{\cal K}\), the following are
equivalent:
\begin{enumerate}
\item there is an endomorphism \(S\) of \({\cal H}\) such that
      \[
      B^*N=MS;                                        \tag{72.14}
      \]
\item the projected defect satisfies
      \[
      W^*N=M^{1/2}S.                                  \tag{72.15}
      \]
\end{enumerate}
When either holds,
\[
 RN=S.                                                \tag{72.16}
\]
Thus a uniform bound on \(S\) controls the factor response without a
separate Korn inverse.
\end{theorem}

\begin{proof}
Since \(W^*=M^{-1/2}B^*\), multiplying (72.14) by \(M^{-1/2}\)
gives (72.15), and conversely.  Multiplying (72.14) by \(M^{-1}\)
gives (72.16).
\end{proof}

Condition (72.14) says that the adjoint defect vanishes to the same
spectral order as \(M\) on its low modes.  Mere orthogonality to the
exact kernel gives only one point and is not enough for a family whose
small positive eigenvalues approach zero.

\begin{corollary}[Approximate neutrality]
\label{cor:v16v72-approx}
If
\[
 B^*N=MS+E,                                           \tag{72.17}
\]
then
\[
 RN=S+M^{-1}E.                                       \tag{72.18}
\]
Consequently a volume-uniform response requires an oriented bound on
\(M^{-1}E\), not merely \(\|E\|\to0\).
\end{corollary}

\begin{proof}
Multiply (72.17) by \(M^{-1}\).
\end{proof}

\subsection{Fourier resonance obstruction}

The need for (72.14) already appears in the scalar one-dimensional
factor model.  Let \(X_L=\mathbb R/L\mathbb Z\),
\[
 B=\partial_x,\qquad M=B^*B=-\partial_x^2            \tag{72.19}
\]
on mean-zero functions, and define
\[
 N_h\xi=(\partial_xh)\xi.                             \tag{72.20}
\]
This is the scalar analogue of (72.6).

\begin{proposition}[High--high to low blow-up]
\label{prop:v16v72-resonance}
There are normalized Fourier modes \(h_L,\xi_L\) with
\(\|\partial_xh_L\|_{L^\infty}\) bounded and output frequency
\[
 q_L=\frac{2\pi}{L}                                   \tag{72.21}
\]
such that
\[
 \|RN_{h_L}\xi_L\|_{L^2}
 =\frac{|p_L|}{|q_L|}
 \longrightarrow\infty                               \tag{72.22}
\]
as \(L\to\infty\), with \(p_L\) bounded away from zero and infinity.
\end{proposition}

\begin{proof}
Use normalized complex Fourier modes.  Choose an integer \(m_L\) so
that
\[
 p_L=\frac{2\pi m_L}{L}\longrightarrow1,
\]
set
\[
 h_L(x)=e^{ip_Lx},\qquad
 \xi_L(x)=e^{i(q_L-p_L)x}.
\]
Then
\[
 N_{h_L}\xi_L=ip_Le^{iq_Lx}.                         \tag{72.23}
\]
On the nonzero \(q_L\)-mode,
\[
 R=M^{-1}B^*
 =(-\partial_x^2)^{-1}(-\partial_x)
\]
has multiplier \(-i/q_L\).  Hence the output amplitude is
\(|p_L/q_L|\), proving (72.22).  Meanwhile
\(\|\partial_xh_L\|_\infty=|p_L|\) stays bounded.
\end{proof}

This is not a gravitational counterexample; it is an exact warning
against estimating the connection-variation defect solely by
\(\|\nabla h\|_\infty\).  Products of two non-low inputs can feed the
lowest nonzero output mode and expose the right inverse.

\subsection{A checkable low-mode condition}

Let \((u_j,\lambda_j)\) be an orthonormal eigenbasis of
\(M\) on the zero-mode complement, and put
\[
 w_j=\lambda_j^{-1/2}Bu_j.                            \tag{72.24}
\]
Then \((w_j)\) is orthonormal in \(\operatorname{Ran}B\), and
\[
 RN\xi
 =\sum_j\lambda_j^{-1/2}
 \langle w_j,N\xi\rangle\,u_j.                       \tag{72.25}
\]

\begin{proposition}[Spectral neutrality owner]
\label{prop:v16v72-spectral}
If an operator \(S\) and a constant \(\varepsilon\geq0\) satisfy
\[
 \sum_j
 \left|
 \lambda_j^{-1/2}\langle w_j,N\xi\rangle
 -\langle u_j,S\xi\rangle
 \right|^2
 \leq\varepsilon^2\|\xi\|^2,                         \tag{72.26}
\]
then
\[
 \|RN-S\|_{\rm op}\leq\varepsilon.                   \tag{72.27}
\]
\end{proposition}

\begin{proof}
Subtract the eigen-expansion of \(S\xi\) from (72.25).  Parseval's
identity turns (72.26) into
\(\|(RN-S)\xi\|^2\leq\varepsilon^2\|\xi\|^2\).
\end{proof}

For \(N=N_h\), (72.26) is the concrete low-mode cancellation that the
reconstruction and gauge choice must prove.  A Sobolev estimate with a
volume-uniform negative-order norm can imply it; an ordinary local
\(L^\infty\) coefficient bound cannot, by Proposition
\ref{prop:v16v72-resonance}.

\begin{assessment}[What V72 closes]
\label{ass:v16v72-status}
V72 computes the exact metric variation of the Killing factor and
separates derivative-bearing terms from the connection defect
\(\xi\cdot\nabla h\).  It proves the exact neutrality condition
\(B^*N=MS\) that removes the inverse Korn gap and gives its spectral
form.

A Fourier model proves that bounded local coefficients alone do not
control the defect uniformly in volume because high--high
interactions can feed the lowest output frequency.  The physical
metric reconstruction has not yet been shown to satisfy the required
neutrality or negative-order Sobolev bound.
\end{assessment}

\begin{decision}[V73 neutral blocks and conserved metric response]
\label{dec:v16v72-next}
Go upstream to the source of the connection defect.  Test whether
shell-local metric variations with vanishing block average, together
with the divergence-free gauge parameter and discrete summation by
parts, force an extra output derivative in \(B^*N_h\).  Formulate a
neutral-block theorem and identify the boundary flux that obstructs
exact divisibility by \(M\).
\end{decision}

\begin{firewall}[Claim boundary after V72]
\label{fire:v16v72}
The Fourier resonance is a scalar model and does not prove failure of
the gravitational response.  The neutrality criteria are sufficient
owners, not verified properties of the current metric shell map.
No nonlinear gauge atlas, determinant limit, or continuum measure is
constructed.
\end{firewall}

\section*{Version V72 append-only record}
\addcontentsline{toc}{section}{Version V72 append-only record}
The complete V71 source was retained byte for byte before this
continuation.  It had \(737588\) bytes, \(22058\) lines, and SHA--256
\[
 \texttt{4F2063B7108D0D42764C8635BFD10CE5AA9A3445596D266FFAD4414FAE52BDAC}.
\]
V72 computed the Killing-factor variation, isolated the low-mode
connection defect, and proved exact and approximate neutrality
criteria together with a resonance obstruction to naive local bounds.


\section{V73: neutral blocks and the boundary-flux owner}
\label{sec:v16v73}

\subsection{Discrete massless right inverse}

Let \(\Lambda_N=(\mathbb Z/L\mathbb Z)^d\) be a periodic lattice with
\(N=L^d\) sites.  Let
\[
 D:\ell^2(\Lambda_N)\longrightarrow
 \ell^2(\Lambda_N;\mathbb R^d)                        \tag{73.1}
\]
be the forward difference, and let
\[
 M=D^*D                                                   \tag{73.2}
\]
on mean-zero scalar fields.  Its right inverse from edge fields is
\[
 R=M^{-1}D^*.                                        \tag{73.3}
\]
For a vector field \(w\), use the unitary Fourier convention
\[
 \widehat w(q)
 =N^{-1/2}\sum_{x\in\Lambda_N}e^{-iq\cdot x}w(x).     \tag{73.4}
\]
The difference symbol is
\[
 s_\mu(q)=e^{iq_\mu}-1,\qquad
 \lambda(q)=|s(q)|^2.                                \tag{73.5}
\]
For \(q\ne0\),
\[
 \widehat{Rw}(q)
 =\frac{\overline{s(q)}\cdot\widehat w(q)}
 {|s(q)|^2}.                                         \tag{73.6}
\]

\subsection{Neutral support gives the missing derivative}

\begin{theorem}[Neutral-block right-inverse estimate]
\label{thm:v16v73-neutral-block}
Let \(Q\subset\Lambda_N\) have graph diameter at most \(D_Q\), and let
\(w\) be supported in \(Q\).  Suppose the componentwise neutrality
condition
\[
 \sum_{x\in Q}w(x)=0.                                 \tag{73.7}
\]
Then
\[
 \|Rw\|_{\ell^2(\Lambda_N)}
 \leq C_dD_Q\|w\|_{\ell^1(Q)},                        \tag{73.8}
\]
where \(C_d\) is independent of \(L\), \(N\), and the lowest nonzero
eigenvalue of \(M\).
\end{theorem}

\begin{proof}
Fix \(x_Q\in Q\).  By (73.7),
\[
 \widehat w(q)
 =N^{-1/2}\sum_{x\in Q}
 \bigl(e^{-iq\cdot x}-e^{-iq\cdot x_Q}\bigr)w(x).
                                                               \tag{73.9}
\]
Hence
\[
 |\widehat w(q)|
 \leq N^{-1/2}|q|D_Q\|w\|_{\ell^1(Q)}.               \tag{73.10}
\]
For \(q_\mu\in[-\pi,\pi]\),
\[
 |e^{iq_\mu}-1|
 =2|\sin(q_\mu/2)|
 \geq\frac2\pi|q_\mu|,
\]
so
\[
 |q|\leq\frac\pi2|s(q)|.                              \tag{73.11}
\]
Equations (73.6), (73.10), and (73.11) give, for every nonzero mode,
\[
 |\widehat{Rw}(q)|
 \leq\frac{\pi D_Q}{2\sqrt N}\|w\|_{\ell^1(Q)}.       \tag{73.12}
\]
There are at most \(N-1\) such modes.  Parseval's identity yields
(73.8) with \(C_d=\pi/2\); allowing equivalent vector and graph
norms only changes the dimensional constant.
\end{proof}

The theorem replaces the dangerous multiplier \(|s(q)|^{-1}\) by
the zero of \(\widehat w(q)\) at \(q=0\).  A fixed block diameter then
makes the estimate uniform in total volume.

\subsection{Neutrality from discrete Stokes}

Let \(Q\) now be a finite lattice block with internal oriented edges.
For a scalar \(h\) and an edge field \(\xi\), define the model
connection response
\[
 w_e=(D_Qh)_e\,\xi_e.                                 \tag{73.13}
\]
The componentwise sum of \(w\) is the discrete pairing
\[
 \sum_{e\in E(Q)}(D_Qh)_e\xi_e
 =\langle D_Qh,\xi\rangle_{E(Q)}.                    \tag{73.14}
\]

\begin{proposition}[Boundary-flux identity]
\label{prop:v16v73-flux}
There is an exact discrete Green formula
\[
 \langle D_Qh,\xi\rangle_{E(Q)}
 =
 \langle h,D_Q^*\xi\rangle_Q
 +{\cal B}_{\partial Q}(h,\xi),                       \tag{73.15}
\]
where \({\cal B}_{\partial Q}\) is the signed flux over edges crossing
\(\partial Q\).  Therefore
\[
 D_Q^*\xi=0,\qquad
 {\cal B}_{\partial Q}(h,\xi)=0                       \tag{73.16}
\]
imply the neutrality (73.7) for the response (73.13).
\end{proposition}

\begin{proof}
Expand the left side as the sum over oriented edges of
\((h({\rm head}\,e)-h({\rm tail}\,e))\xi_e\).  Regroup all terms by
vertices.  Contributions from internal edges give
\(\langle h,D_Q^*\xi\rangle_Q\); the unmatched crossing-edge terms
are exactly \({\cal B}_{\partial Q}\).  This proves (73.15), and
(73.16) gives the last assertion.
\end{proof}

The divergence-free condition alone is sufficient on a periodic
block.  On a shell block with boundary, the flux is an independent
owner.

\subsection{Why mean-zero metric variation is not enough}

\begin{proposition}[Average-zero does not imply response neutrality]
\label{prop:v16v73-average-no-go}
The condition
\[
 \sum_{x\in Q}h(x)=0                                  \tag{73.17}
\]
does not imply
\[
 \sum_{e\in E(Q)}(D_Qh)_e\xi_e=0.                    \tag{73.18}
\]
\end{proposition}

\begin{proof}
Take a one-dimensional block with vertices \(0,1,2\), internal edges
\((0,1)\) and \((1,2)\), and
\[
 h=(-1,0,1),\qquad \xi=(1,0).
\]
Then \(\sum_xh(x)=0\), while
\[
 \sum_e(D_Qh)_e\xi_e=(0-(-1))\cdot1+(1-0)\cdot0=1.
\]
\end{proof}

Thus the correct cancellation is neutrality of the connection
response or vanishing of its discrete Stokes flux, not merely a
vanishing block average of \(h\).

\subsection{A neutral-block decomposition}

Suppose
\[
 w=\sum_{\alpha\in{\cal I}_N}w_\alpha,                \tag{73.19}
\]
where each \(w_\alpha\) is supported in a block \(Q_\alpha\) of
diameter at most \(D_0\) and satisfies
\[
 \sum_xw_\alpha(x)=0.                                 \tag{73.20}
\]

\begin{proposition}[Global response from neutral blocks]
\label{prop:v16v73-decomposition}
One has
\[
 \|Rw\|_{\ell^2}
 \leq C_dD_0\sum_\alpha
 \|w_\alpha\|_{\ell^1}.                              \tag{73.21}
\]
If, in addition, the Fourier responses have uniformly bounded
overlap,
\[
 \sum_\alpha
 |\widehat{Rw_\alpha}(q)|^2
 \leq m_0
 \left|\sum_\alpha\widehat{Rw_\alpha}(q)\right|^2
 \quad\text{after a fixed colouring into }m_0
 \text{ noninteracting classes},                     \tag{73.22}
\]
then the sharper square-sum estimate
\[
 \|Rw\|_{\ell^2}^2
 \leq C_d^2m_0D_0^2
 \sum_\alpha\|w_\alpha\|_{\ell^1}^2                 \tag{73.23}
\]
holds.
\end{proposition}

\begin{proof}
The triangle inequality and Theorem
\ref{thm:v16v73-neutral-block} give (73.21).  Under the stated
colouring/overlap hypothesis, apply the pointwise Fourier estimate
(73.12) within each class, use Cauchy--Schwarz over the at most
\(m_0\) classes, and sum modes by Parseval.  This gives (73.23).
\end{proof}

The hypothesis (73.22) is deliberately explicit: arbitrary neutral
blocks can add coherently, so a square-sum bound does not follow from
neutrality alone.

\subsection{Insertion into the moving-factor defect}

Let \(N_h\xi\) be the discrete counterpart of (72.6), partitioned
into shell blocks.  If every block satisfies the divergence and flux
conditions (73.16), Theorem~\ref{thm:v16v73-neutral-block} supplies
the missing negative-order bound in the scalar factor model.  In the
tensor gauge problem the required statement is componentwise:
\[
 \sum_{x\in Q_\alpha}
 (N_h\xi)_{\mu\nu}(x)=0
 \quad\text{for every retained tensor component}.    \tag{73.24}
\]
The V48 collar traces are the natural existing owner for
\({\cal B}_{\partial Q}\), but they currently bound its size rather
than prove its cancellation.

\begin{theorem}[Neutrality/flux alternative]
\label{thm:v16v73-alternative}
For a fixed-diameter shell block, the low-mode factor response is
volume uniform by either of the following sufficient routes:
\begin{enumerate}
\item exact componentwise neutrality (73.24);
\item a boundary-flux correction \(c_{\partial Q}\) such that
      \(N_h\xi-c_{\partial Q}\) is neutral and
      \(Rc_{\partial Q}\) has a separately summable owner.
\end{enumerate}
A bound on the uncorrected boundary flux with no negative-order
response estimate is insufficient.
\end{theorem}

\begin{proof}
The first route is Theorem~\ref{thm:v16v73-neutral-block} applied to
each component.  For the second, use linearity
\[
 RN_h\xi=R(N_h\xi-c_{\partial Q})+Rc_{\partial Q}.
\]
The first term is controlled by the neutral-block theorem and the
second by its declared owner.
\end{proof}

\begin{assessment}[What V73 closes]
\label{ass:v16v73-status}
V73 proves that fixed-diameter block neutrality supplies exactly the
low-frequency zero needed to control the massless right inverse
uniformly in volume.  Discrete Stokes identifies the owner:
divergence-free response plus vanishing boundary flux.  A mean-zero
metric variation alone does not imply the required cancellation.

For shell reconstructions, the unresolved term is now a concrete
collar flux.  The existing V48 trace machinery bounds such terms but
does not make them neutral; a correction or summable negative-order
flux response is still required.
\end{assessment}

\begin{decision}[V74 tensor shell flux correction]
\label{dec:v16v73-next}
Lift the scalar neutral-block theorem to the tensor connection
variation.  Write the discrete adjoint of the symmetrized gradient,
derive the face flux for \(N_h\xi\), and construct the minimal
block-constant correction that cancels its zero Fourier mode.  Bound
the corrected bulk response and the remaining face response using the
V48 collar owner, with scaling explicit in block diameter and face
area.
\end{decision}

\begin{firewall}[Claim boundary after V73]
\label{fire:v16v73}
The neutral-block theorem is proved for a scalar periodic lattice
factor.  Its tensor lift, shell flux correction, collar summability,
and compatibility with the nonlinear metric gauge remain open.  It
does not construct the continuum determinant or measure.
\end{firewall}

\section*{Version V73 append-only record}
\addcontentsline{toc}{section}{Version V73 append-only record}
The complete V72 source was retained byte for byte before this
continuation.  It had \(747706\) bytes, \(22393\) lines, and SHA--256
\[
 \texttt{9AC1A00E6B4E269493845C663D92812DD231A34BF7058A512197B1A2F2FEAC8E}.
\]
V73 proved the neutral-block massless right-inverse estimate,
identified the exact boundary-flux obstruction, and selected the V48
collar traces as the next physical owner.


\section{V74: tensor shell correction and the four-dimensional flux bound}
\label{sec:v16v74}

\subsection{Tensor symmetrized-gradient model}

On the periodic lattice \(\Lambda_N\), let
\[
 B:\ell^2(\Lambda_N;\mathbb R^d)
 \longrightarrow
 \ell^2(\Lambda_N;{\rm Sym}^2\mathbb R^d)             \tag{74.1}
\]
be a translation-invariant discrete symmetrized gradient.  Assume its
nonzero Fourier symbol satisfies
\[
 B(q)^*B(q)\succeq c_B|s(q)|^2I_d,\qquad c_B>0,       \tag{74.2}
\]
where \(s(q)\) is (73.5).  Remove the constant vector kernel and put
\[
 M=B^*B,\qquad R=M^{-1}B^*.                           \tag{74.3}
\]
Then
\[
 \|R(q)\|_{\rm op}
 \leq c_B^{-1/2}|s(q)|^{-1}.                          \tag{74.4}
\]

\begin{proof}
The least singular value of \(B(q)\) is at least
\(\sqrt{c_B}|s(q)|\).  The Moore--Penrose right inverse
\((B(q)^*B(q))^{-1}B(q)^*\) therefore has norm at most the reciprocal.
\end{proof}

For the standard flat symmetrized gradient, (74.2) follows directly
from
\[
 |\operatorname{sym}(s\otimes\xi)|^2
 =\frac12\left(|s|^2|\xi|^2+|s\cdot\xi|^2\right)
 \geq\frac12|s|^2|\xi|^2.                            \tag{74.5}
\]

\subsection{A dimension-dependent localized bound}

Define
\[
 {\cal G}_{d,N}
 :=\frac1N\sum_{q\ne0}\frac1{|s(q)|^2}.              \tag{74.6}
\]

\begin{lemma}[Uniform massless lattice integral]
\label{lem:v16v74-green}
If \(d>2\), then
\[
 {\cal G}_d:=\sup_N{\cal G}_{d,N}<\infty.             \tag{74.7}
\]
For \(d=2\), the same quantity grows logarithmically with the side
length, and for \(d=1\) it grows linearly.
\end{lemma}

\begin{proof}
On the Brillouin zone,
\[
 |s(q)|\geq\frac2\pi|q|
\]
as in (73.11).  Partition the discrete momenta into dyadic shells
\(2^{-j-1}<|q|\leq2^{-j}\).  The normalized number of points in such a
shell is at most \(C_d2^{-jd}\), while
\(|q|^{-2}\leq C2^{2j}\).  The shell contribution is therefore at
most \(C_d2^{-j(d-2)}\), whose sum converges exactly when \(d>2\).
The innermost shell gives the stated logarithmic and linear growth in
dimensions two and one.  Away from \(q=0\), the summand is uniformly
bounded.
\end{proof}

\begin{theorem}[Localized tensor response]
\label{thm:v16v74-localized}
Let \(w\) be a symmetric-tensor field supported in a block
\(Q\subset\Lambda_N\).  If \(d>2\), then
\[
 \|Rw\|_{\ell^2}
 \leq c_B^{-1/2}{\cal G}_d^{1/2}
 \|w\|_{\ell^1(Q)}.                                  \tag{74.8}
\]
No neutrality assumption is required.
\end{theorem}

\begin{proof}
The unitary Fourier transform gives
\[
 |\widehat w(q)|
 \leq N^{-1/2}\|w\|_{\ell^1(Q)}.                     \tag{74.9}
\]
By (74.4),
\[
 |\widehat{Rw}(q)|^2
 \leq
 \frac1{c_BN|s(q)|^2}\|w\|_{\ell^1(Q)}^2.
\]
Sum over \(q\ne0\), use Parseval and (74.7), and take the square root.
\end{proof}

This theorem resolves the exact infrared issue for a single
fixed-support response in four dimensions.  Neutrality remains useful
because it gives the stronger dimension-independent estimate V73 and
better spatial decay, but it is not necessary for an \(L^2\) response
when \(d=4\).

\subsection{Minimal block-constant correction}

Let
\[
 m_Q=\sum_{x\in Q}w(x),\qquad
 c_Q(x)=\frac{m_Q}{|Q|}{\bf1}_Q(x),\qquad
 w_Q^0=w-c_Q.                                        \tag{74.10}
\]

\begin{proposition}[Neutral bulk plus flux correction]
\label{prop:v16v74-correction}
The field \(w_Q^0\) is neutral:
\[
 \sum_xw_Q^0(x)=0.                                   \tag{74.11}
\]
Among tensor fields supported in \(Q\) with sum \(m_Q\), \(c_Q\)
has the least \(\ell^2\) norm, namely
\[
 \|c_Q\|_{\ell^2}=|Q|^{-1/2}|m_Q|.                   \tag{74.12}
\]
Moreover, for \(d>2\),
\[
 \|Rw\|_{\ell^2}
 \leq
 C_dD_Q\|w_Q^0\|_{\ell^1}
 +c_B^{-1/2}{\cal G}_d^{1/2}|m_Q|.                  \tag{74.13}
\]
\end{proposition}

\begin{proof}
Equation (74.11) follows from (74.10).  For any supported \(c\) with
\(\sum_Qc=m_Q\), Cauchy--Schwarz gives
\[
 |m_Q|^2\leq |Q|\|c\|_{\ell^2}^2,
\]
with equality exactly for the constant choice (74.10), proving
(74.12).  Apply the neutral-block theorem componentwise to \(w_Q^0\).
Apply Theorem~\ref{thm:v16v74-localized} to \(c_Q\), noting that
\[
 \|c_Q\|_{\ell^1}=|m_Q|.
\]
Add the two responses.
\end{proof}

The correction does not claim that \(c_Q\) is in the range of \(B\).
Its response is controlled by the integrability of the
four-dimensional massless Green multiplier.

\subsection{The tensor connection flux}

For the continuum connection term from (72.6),
\[
 (N_h\xi)_{\mu\nu}
 =\frac1{\sqrt2}\xi^\rho\nabla_\rho h_{\mu\nu},       \tag{74.14}
\]
let \(Q\) be a smooth block.  Define its tensor mean
\[
 (m_Q)_{\mu\nu}
 =\int_Q(N_h\xi)_{\mu\nu}\,d{\rm vol}_g.              \tag{74.15}
\]

\begin{theorem}[Exact tensor flux identity]
\label{thm:v16v74-flux}
One has
\[
 (m_Q)_{\mu\nu}
 =
 \frac1{\sqrt2}\int_{\partial Q}
 h_{\mu\nu}\,\xi^\rho n_\rho\,dS_g
 -\frac1{\sqrt2}\int_Q
 h_{\mu\nu}\nabla_\rho\xi^\rho\,d{\rm vol}_g.         \tag{74.16}
\]
For a residual volume-preserving parameter,
\(\nabla_\rho\xi^\rho=0\), so
\[
 (m_Q)_{\mu\nu}
 =
 \frac1{\sqrt2}\int_{\partial Q}
 h_{\mu\nu}(\xi\cdot n)\,dS_g.                        \tag{74.17}
\]
\end{theorem}

\begin{proof}
Apply the divergence theorem componentwise to
\(\nabla_\rho(h_{\mu\nu}\xi^\rho)\) and rearrange.
\end{proof}

Thus V73's abstract obstruction is exactly the normal gauge flux
weighted by the metric response.

\begin{corollary}[Collar trace owner]
\label{cor:v16v74-collar}
For divergence-free \(\xi\),
\[
 |m_Q|
 \leq\frac1{\sqrt2}
 \|h\|_{L^2(\partial Q)}
 \|\xi\cdot n\|_{L^2(\partial Q)}.                    \tag{74.18}
\]
If
\[
 \|h\|_{L^2(\partial Q)}
 \leq C_h^{\partial},\qquad
 \|\xi\cdot n\|_{L^2(\partial Q)}
 \leq C_\xi^{\partial},                              \tag{74.19}
\]
then the correction term in (74.13) is at most
\[
 \frac{{\cal G}_d^{1/2}}{\sqrt{2c_B}}\,
 C_h^{\partial}C_\xi^{\partial}.                     \tag{74.20}
\]
\end{corollary}

\begin{proof}
Cauchy--Schwarz on the face proves (74.18); insert (74.19) into the
second term of (74.13).
\end{proof}

The V48 collar machinery provides the natural type of
\(L^2(\partial Q)\) trace estimate for \(h\).  A corresponding trace
owner for the residual gauge parameter \(\xi\cdot n\) is still needed.

\subsection{Block scaling}

If \(Q\) is a regular block of side length \(\ell\), its volume and
face area scale as
\[
 |Q|\asymp\ell^d,\qquad |\partial Q|\asymp\ell^{d-1}. \tag{74.21}
\]
Uniform pointwise trace bounds would give
\[
 |m_Q|\lesssim
 \ell^{d-1}\|h\|_{L^\infty(\partial Q)}
 \|\xi\cdot n\|_{L^\infty(\partial Q)}.              \tag{74.22}
\]
For a fixed microscopic block size this is a local constant.  For
growing blocks it is not intensive; the \(L^2\) trace allocation in
(74.18) must be combined with the shell normalization before the
V66 influence row is claimed.

\begin{theorem}[Four-dimensional shell response criterion]
\label{thm:v16v74-criterion}
In \(d=4\), a family of tensor connection responses admits a
volume-uniform one-block right-inverse bound if:
\begin{enumerate}
\item every response is supported in a block of uniformly bounded
      diameter;
\item its neutral part has uniformly bounded block \(\ell^1\) norm;
\item its face product in (74.18) is uniformly bounded.
\end{enumerate}
The bound is the right side of (74.13) with (74.20).
\end{theorem}

\begin{proof}
Use Proposition~\ref{prop:v16v74-correction}; Lemma
\ref{lem:v16v74-green} applies because \(4>2\), and
Corollary~\ref{cor:v16v74-collar} controls the nonneutral correction.
\end{proof}

\begin{assessment}[What V74 closes]
\label{ass:v16v74-status}
V74 lifts the massless right-inverse estimate to a tensor
symmetrized-gradient factor.  In dimensions above two, a localized
nonneutral response is already \(L^2\)-controlled because
\(|q|^{-2}\) is locally integrable.  The minimal block-constant
correction separates the neutral bulk from its mean.

For the actual connection variation, that mean is exactly a boundary
flux.  V48 can own the metric trace factor, while a gauge-parameter
normal-trace estimate and the shell scaling remain missing.  The
result controls one block; summability of all block-to-block
influences is not yet proved.
\end{assessment}

\begin{decision}[V75 companion audit and face-row summability]
\label{dec:v16v74-next}
Perform the mandatory YM/NS companion audit.  Then estimate the
spatial tail of the corrected tensor response and its squared Gram
kernel.  Prove a face-to-block row-sum criterion in four dimensions,
including the overlap multiplicity of shell collars, and decide
whether the V48 trace constants scale compatibly.
\end{decision}

\begin{firewall}[Claim boundary after V74]
\label{fire:v16v74}
The tensor result is a flat translation-invariant model with an
assumed symbol Korn bound.  It does not verify the nonlinear curved
factor, the gauge-parameter face trace, the all-block influence row,
determinant orientation, or the continuum measure.
\end{firewall}

\section*{Version V74 append-only record}
\addcontentsline{toc}{section}{Version V74 append-only record}
The complete V73 source was retained byte for byte before this
continuation.  It had \(757417\) bytes, \(22696\) lines, and SHA--256
\[
 \texttt{7D0307F101563576C8C28D00FAB536FE4B917B2EF9FF563DEA0D84F0AB70E38A}.
\]
V74 derived the tensor shell correction, the four-dimensional
localized massless response bound, and the exact boundary-flux owner
for the connection variation.


\section{V75: companion audit and collective face-flux summability}
\label{sec:v16v75}

\subsection{Mandatory YM/NS companion audit}

\begin{audit}[V75 companion checkpoint]
\label{aud:v16v75-companions}
The newest active numeric companions inspected were
\[
\begin{array}{c|r|r|l}
\hbox{file}&\hbox{lines}&\hbox{bytes}&\hbox{SHA--256}\\ \hline
\texttt{Renewal\_Yang\_Mills\_Mass\_Gap\_Research\_Notes\_V58.tex}
&13540&497016&
\texttt{1F01247F8B5C7DB818E3A873219D3203316A547892DD786B1B375435ED58193D}\\
\texttt{Renewal\_NS\_Stability\_Research\_Notes\_V26.tex}
&5319&278283&
\texttt{82C887F8C2C69E4F28FAD7C3DC8DE5B9099CB5A4BF431EBA1FFF3E91935978AD}.
\end{array}                                             \tag{75.1}
\]
\end{audit}

YM V58 exactly constructs three new charts and, together with five
inherited charts, covers a fourth-orientation seed cube.  It retains
the correct firewall: the inverted cube supplies the other direct
pencil sign, while a full mixed quarter cell and interacting measure
remain open.  NS V26 remains unchanged.

\begin{theorem}[V75 import decision]
\label{thm:v16v75-audit-decision}
No companion theorem supplies a gauge-parameter face trace or a
gravitational flux row sum.  YM's exact adjacent-chart assembly is
consistent with pairing shared shell faces, but none of its Wilson
floors enters the estimates below.
\end{theorem}

\begin{proof}
The companions contain no symmetrized-gradient right inverse,
connection flux, or collar trace for residual diffeomorphisms.
Therefore their hypotheses do not imply any V74 owner.
\end{proof}

\subsection{Single-response control is not collective control}

Theorem~\ref{thm:v16v74-localized} proves a uniform \(L^2\) bound for
one localized response in four dimensions.  A translated family of
nonneutral responses can still have a divergent collective Gram
operator.

Use the scalar factor
\[
 B=D,\qquad M=D^*D,\qquad R=M^{-1}D^*                \tag{75.2}
\]
on the periodic lattice.  Let \(w_0\) be a fixed compactly supported
edge field and \(w_x=\tau_xw_0\).

\begin{proposition}[Nonneutral translated-family obstruction]
\label{prop:v16v75-collective-no-go}
Let
\[
 m=\widehat w_0(0)=\sum_yw_0(y).                      \tag{75.3}
\]
If there is a sequence of lowest nonzero momenta \(q_N\to0\) for
which
\[
 |\widehat q_N\cdot m|\geq c_m|m|>0,                 \tag{75.4}
\]
then the convolution Gram matrix
\[
 {\cal K}_N(x,y)=\langle Rw_x,Rw_y\rangle_{\ell^2}    \tag{75.5}
\]
satisfies
\[
 \|{\cal K}_N\|_{\ell^2\to\ell^2}
 \geq c\,|q_N|^{-2}\longrightarrow\infty.            \tag{75.6}
\]
Consequently its absolute row sum cannot be bounded uniformly.
\end{proposition}

\begin{proof}
Translation invariance diagonalizes (75.5).  Its Fourier multiplier
is
\[
 |\widehat{Rw_0}(q)|^2
 =
 \frac{|\overline{s(q)}\cdot\widehat w_0(q)|^2}
 {|s(q)|^4}.                                         \tag{75.7}
\]
Because \(w_0\) is finitely supported,
\(\widehat w_0(q)=m+O(|q|)\), up to the harmless Fourier
normalization common to every translated column.  Along (75.4),
the numerator in (75.7) is bounded below by
\(c|q_N|^2|m|^2\), while the denominator is at most
\(C|q_N|^4\).  The operator norm of a convolution is the supremum of
the absolute multiplier, proving (75.6).  An absolute row bound
dominates the convolution operator norm.
\end{proof}

In \(d=4\), each column \(Rw_x\) can have a uniform \(L^2\) norm even
while the collective operator diverges like the square of the system
diameter.  Therefore the unpaired face correction in V74 is
insufficient for the V66 all-block theorem.

\subsection{Neutral responses remove the infrared pole}

\begin{proposition}[Neutral translated-family bound]
\label{prop:v16v75-neutral-family}
If
\[
 \sum_yw_0(y)=0,                                     \tag{75.8}
\]
then
\[
 \sup_{N,q\ne0}|\widehat{Rw_0}(q)|
 \leq C_dD_0\|w_0\|_{\ell^1},                        \tag{75.9}
\]
where \(D_0\) is the support diameter.  Hence the Gram convolution
operator satisfies
\[
 \|{\cal K}_N\|_{\ell^2\to\ell^2}
 \leq C_d^2D_0^2\|w_0\|_{\ell^1}^2                  \tag{75.10}
\]
uniformly in volume.
\end{proposition}

\begin{proof}
The Fourier cancellation proof (73.9)--(73.12), without summing the
modes, gives (75.9).  The Gram multiplier is the square of the
response multiplier, so (75.10) follows.
\end{proof}

The result controls the spectral norm required by the positive-matrix
condition (66.9).  Absolute spatial row summability may require
additional smoothness or moment cancellation; it is not inferred from
(75.9).

\subsection{Pairing an internal shell face}

Let \(Q_i,Q_j\) be adjacent equal-volume blocks sharing a face \(f\).
Let \(F_f\) be the tensor flux assigned with opposite orientations to
the two blocks.  Define the paired block source
\[
 c_f(x)
 =\frac{F_f}{|Q_i|}{\bf1}_{Q_i}(x)
  -\frac{F_f}{|Q_j|}{\bf1}_{Q_j}(x).                 \tag{75.11}
\]

\begin{theorem}[Face-pair neutrality]
\label{thm:v16v75-face-pair}
The source (75.11) is neutral:
\[
 \sum_xc_f(x)=0.                                     \tag{75.12}
\]
Its support diameter is at most twice the block diameter, and
\[
 \|Rc_f\|_{\ell^2}
 \leq2C_dD_0|F_f|.                                   \tag{75.13}
\]
For all translations of the same oriented face type, the collective
Gram operator has norm at most
\[
 4C_d^2D_0^2
 \sup_f|F_f|^2.                                      \tag{75.14}
\]
\end{theorem}

\begin{proof}
Equal block volumes and the opposite signs in (75.11) prove
(75.12).  Also
\[
 \|c_f\|_{\ell^1}=2|F_f|.
\]
Apply Theorem~\ref{thm:v16v73-neutral-block} to obtain (75.13), and
Proposition~\ref{prop:v16v75-neutral-family} to its translated
family for (75.14).
\end{proof}

Internal face pairing converts the dangerous monopole response into a
block dipole.  On a periodic tessellation every face is internal.  On
a manifold with physical boundary, the unpaired exterior faces need a
boundary condition or a separate boundary determinant.

\subsection{Overlap multiplicity}

Suppose the block adjacency graph has degree at most \(d_{\rm face}\).
Its faces can be partitioned into at most
\[
 m_{\rm col}\leq2d_{\rm face}-1                       \tag{75.15}
\]
matchings by a greedy edge colouring.  Within one matching the paired
sources have disjoint block supports.

\begin{proposition}[All-face spectral row]
\label{prop:v16v75-all-face}
Assume each face family of one translation type satisfies (75.14),
and at most \(m_{\rm type}\) geometric face types occur.  Then the
collective paired-face response operator obeys
\[
 \|{\cal R}_{\rm face}\|_{\ell^2\to\ell^2}
 \leq
 2C_dD_0
 \sqrt{m_{\rm col}m_{\rm type}}\,
 F_*,
 \qquad
 F_*:=\sup_f\|F_f\|_{\rm op}.                         \tag{75.16}
\]
Its Gram operator is bounded by the square of the right side.
\end{proposition}

\begin{proof}
Split the faces by matching and geometric type.  Within each class,
translation/Floquet diagonalization gives the bound (75.14), while
disjoint support prevents an additional volume factor.  Cauchy--
Schwarz across at most \(m_{\rm col}m_{\rm type}\) classes gives
(75.16).  Squaring bounds the Gram operator.
\end{proof}

This is a spectral influence bound.  It can enter the matrix condition
(66.9) directly when the paired-face Gram is retained as one operator,
instead of replacing it by entrywise absolute values.

\subsection{Trace scaling on a regular block}

For a block \(Q\) of physical side \(\ell\), the standard trace
inequality has the scale
\[
 \|u\|_{L^2(\partial Q)}^2
 \leq C_{\rm tr}
 \left(
 \ell^{-1}\|u\|_{L^2(Q)}^2
 +\ell\|\nabla u\|_{L^2(Q)}^2
 \right).                                            \tag{75.17}
\]
If \(\xi\) has zero block mean, Poincar\'e gives
\[
 \|\xi\|_{L^2(Q)}
 \leq C_P\ell\|\nabla\xi\|_{L^2(Q)},                 \tag{75.18}
\]
and therefore
\[
 \|\xi\cdot n\|_{L^2(\partial Q)}
 \leq C_{\xi}\ell^{1/2}
 \|\nabla\xi\|_{L^2(Q)}.                             \tag{75.19}
\]

\begin{proof}
Insert (75.18) into (75.17) and take square roots; normal projection
does not increase the trace norm.
\end{proof}

Combining (74.18) and (75.19) gives
\[
 |F_f|
 \leq
 C\ell^{1/2}
 \|h\|_{L^2(f)}
 \|\nabla\xi\|_{L^2(Q_i\cup Q_j)}.                   \tag{75.20}
\]
Thus a V48-type metric trace bound
\[
 \|h\|_{L^2(f)}
 \leq C_h\ell^{-1/2}\|h\|_{{\cal X}(Q_i\cup Q_j)}    \tag{75.21}
\]
would make the flux operator intensive:
\[
 |F_f|
 \leq CC_h
 \|h\|_{\cal X}\|\nabla\xi\|_{L^2}.                  \tag{75.22}
\]
The precise V48 collar norm must be matched to (75.21); its existence
cannot be inferred from dimensional scaling alone.

\begin{assessment}[What V75 closes]
\label{ass:v16v75-status}
The companion audit records YM's fourth-orientation seed cube but
imports no gauge estimate.  V75 proves that a collection of
individually controlled nonneutral face responses can still have a
divergent Gram operator.  Exact pairing of each internal face converts
the flux monopole into a neutral block dipole and yields a
volume-uniform collective spectral bound with explicit colouring
multiplicity.

The remaining owner is now sharply scaled: the metric collar trace
must contribute the reciprocal \(\ell^{1/2}\) factor needed to match
the gauge trace in (75.19).  This has not yet been verified against
the actual V48 reconstruction constants.
\end{assessment}

\begin{decision}[V76 V48 scaling match]
\label{dec:v16v75-next}
Read the exact V48 collar inequalities in the active prefix and map
their mesh powers to (75.21).  Derive the resulting face-flux
influence constant without changing norms.  If the powers fail, move
upstream to a face-averaged reconstruction coordinate or a
constraint-compatible normal gauge condition rather than hiding the
loss in an unspecified constant.
\end{decision}

\begin{firewall}[Claim boundary after V75]
\label{fire:v16v75}
The paired-face theorem is a flat periodic block model.  It does not
prove the curved shell traces, handle physical boundary faces, control
kernel-stratum transitions, or establish the determinant and
continuum measure.
\end{firewall}

\section*{Version V75 append-only record}
\addcontentsline{toc}{section}{Version V75 append-only record}
The complete V74 source was retained byte for byte before this
continuation.  It had \(766866\) bytes, \(23005\) lines, and SHA--256
\[
 \texttt{0F01EBBE3F979D3AD2F06203403B456A8B27719682AAFD4999DB9B2729B0713B}.
\]
V75 recorded the mandatory companion audit, proved the collective
nonneutral obstruction, and constructed the neutral paired-face
spectral influence bound with explicit trace scaling.


\section{V76: exact collar-norm match for the paired gauge flux}
\label{sec:v16v76}

\subsection{The V48 owner does not have the required norm}

The exact smooth-face reconstruction estimate proved in V48 is
\[
 \|h|_\Sigma\|_{L^1(\Sigma)}
 \leq C_{\rm tr}\left(
 \ell^{-1}\|h\|_{L^1({\cal C}_\ell)}
 +\|\nabla_nh\|_{L^1({\cal C}_\ell)}
 \right).                                            \tag{76.1}
\]
In the notation of (48.18)--(48.19),
\[
 \mathfrak F_{n,q}
 \leq C_{\rm tr}\left(
 \ell_{n,q}^{-1}\mathfrak W^{(0)}_{n,q}
 +\mathfrak W^{(1)}_{n,q}\right).                   \tag{76.2}
\]
The tensor flux in V74 instead satisfies
\[
 |m_Q|
 \leq2^{-1/2}
 \|h\|_{L^2(\partial Q)}
 \|\xi\cdot n\|_{L^2(\partial Q)}.                   \tag{76.3}
\]

\begin{proposition}[\(L^1\) trace does not load the \(L^2\) flux]
\label{prop:v16v76-mismatch}
There is no constant \(C\), independent of \(h\), such that
\[
 \|h\|_{L^2(\Sigma)}
 \leq C\|h\|_{L^1(\Sigma)}                            \tag{76.4}
\]
on a nonatomic face of positive finite measure.  Consequently
(76.1) alone does not imply the metric factor required in (76.3).
\end{proposition}

\begin{proof}
Choose measurable sets \(E_m\subset\Sigma\) with
\(|E_m|=m^{-1}\), and set \(h_m=m{\bf1}_{E_m}\) in one fixed tensor
component.  Then
\[
 \|h_m\|_{L^1}=1,\qquad
 \|h_m\|_{L^2}=m^{1/2}\longrightarrow\infty.
\]
Smooth approximations preserve the conclusion.
\end{proof}

An \(L^1\!\times L^\infty\) estimate could use V48, but it would
require a uniform \(L^\infty\) face trace of the residual gauge
parameter.  No such owner exists in the active notes, and in four
dimensions it would require substantially higher Sobolev regularity.
The \(L^2\) route is therefore the correct upstream repair.

\subsection{An explicit \(L^2\) collar theorem}

Let
\[
 {\cal C}_\ell=[0,\ell]\times\Sigma,\qquad
 g=dt^2+\gamma_t,\qquad
 dA_t=J(t,x)dA_0,                                    \tag{76.5}
\]
with
\[
 J(t,x)\geq j_->0.                                   \tag{76.6}
\]
As in V48, use parallel transport along the normal geodesics.

\begin{theorem}[Uniform \(L^2\) collar trace]
\label{thm:v16v76-L2-trace}
For every tensor field \(u\in W^{1,2}({\cal C}_\ell)\),
\[
 \|u|_\Sigma\|_{L^2(\Sigma)}^2
 \leq\frac2{j_-}\left(
 \ell^{-1}\|u\|_{L^2({\cal C}_\ell)}^2
 +\ell\|\nabla_nu\|_{L^2({\cal C}_\ell)}^2
 \right).                                            \tag{76.7}
\]
Hence
\[
 \|u|_\Sigma\|_{L^2(\Sigma)}
 \leq\sqrt{\frac2{j_-}}\left(
 \ell^{-1/2}\|u\|_{L^2({\cal C}_\ell)}
 +\ell^{1/2}\|\nabla_nu\|_{L^2({\cal C}_\ell)}
 \right).                                            \tag{76.8}
\]
\end{theorem}

\begin{proof}
For almost every \(x\) and \(t\in[0,\ell]\),
\[
 |u(0,x)|
 \leq |u(t,x)|+\int_0^t|\nabla_nu(s,x)|\,ds.
\]
Using \((a+b)^2\leq2a^2+2b^2\) and Cauchy--Schwarz in \(s\),
\[
 |u(0,x)|^2
 \leq2|u(t,x)|^2
 +2\ell\int_0^\ell|\nabla_nu(s,x)|^2\,ds.            \tag{76.9}
\]
Average in \(t\), integrate in \(x\), and use
\(dt\,dA_0\leq j_-^{-1}dV_g\).  This proves (76.7).
Taking square roots and using
\(\sqrt{a+b}\leq\sqrt a+\sqrt b\) gives (76.8).
\end{proof}

The factor \(\ell^{-1/2}\) is necessary for a field constant in the
normal direction.  Thus the new powers are sharp at the level of
dimensional scaling.

\subsection{Discrete \(L^2\) version}

For a collar with layers \(k=0,\ldots,L\), lattice spacing \(a\), and
\(\ell_a=(L+1)a\), let \(U_{k,x}\) be tensor valued.

\begin{theorem}[Discrete \(L^2\) collar trace]
\label{thm:v16v76-discrete}
With uniformly comparable cell and face weights,
\[
 a^{d-1}\sum_x|U_{0,x}|^2
 \leq C_w\left[
 \frac2{\ell_a}a^d\sum_{k=0}^L\sum_x|U_{k,x}|^2
 +2\ell_a a^d\sum_{k=0}^{L-1}\sum_x
 |\nabla_n^aU_{k,x}|^2
 \right].                                            \tag{76.10}
\]
\end{theorem}

\begin{proof}
For each \(k,x\),
\[
 |U_{0,x}|^2
 \leq2|U_{k,x}|^2
 +2k\sum_{j=0}^{k-1}|U_{j+1,x}-U_{j,x}|^2.           \tag{76.11}
\]
Average over \(k=0,\ldots,L\), use \(k\leq L+1\), sum in \(x\), and
replace differences by \(a\nabla_n^aU\).  Multiplication by
\(a^{d-1}\) gives (76.10); uniform weight comparison costs \(C_w\).
\end{proof}

\subsection{Metric and gauge trace pairing}

For a shell collar \({\cal C}_{n,q}\) of width \(\ell_{n,q}\), define
the Hilbert reconstruction owners
\[
 {\cal H}^{(j)}_{n,q}
 :=
 \sup_{\substack{s\in G_n\\\|a\|_S=1}}
 \|\nabla^jh_{n,a}(s)\|_{L^2({\cal C}_{n,q})},
 \qquad j=0,1.                                       \tag{76.12}
\]
Theorem~\ref{thm:v16v76-L2-trace} gives
\[
 \|h_{n,a}|_{\Sigma_{n,q}}\|_{L^2}
 \leq C_2\left(
 \ell_{n,q}^{-1/2}{\cal H}^{(0)}_{n,q}
 +\ell_{n,q}^{1/2}{\cal H}^{(1)}_{n,q}
 \right),                                            \tag{76.13}
\]
where \(C_2=\sqrt{2/j_0}\) under the V48 uniform geometry
\(j_-\geq j_0>0\).

Assume the residual gauge parameter has zero mean on the union of the
two face-adjacent blocks and a local Poincar\'e estimate
\[
 \|\xi\|_{L^2(Q_i\cup Q_j)}
 \leq C_P\ell_{n,q}\|\nabla\xi\|_{L^2(Q_i\cup Q_j)}. \tag{76.14}
\]
Applying (76.8) to \(\xi\) gives
\[
 \|\xi\cdot n\|_{L^2(\Sigma_{n,q})}
 \leq C_\xi\ell_{n,q}^{1/2}
 \|\nabla\xi\|_{L^2(Q_i\cup Q_j)},                   \tag{76.15}
\]
with \(C_\xi=C_2(C_P+1)\).

\begin{theorem}[Exact paired face-flux load]
\label{thm:v16v76-flux-load}
Under (76.12)--(76.15), the flux operator of a unit soft metric
direction satisfies
\[
 |F_{n,q}|
 \leq C_{\rm flux}
 \left(
 {\cal H}^{(0)}_{n,q}
 +\ell_{n,q}{\cal H}^{(1)}_{n,q}
 \right)
 \|\nabla\xi\|_{L^2(Q_i\cup Q_j)},                   \tag{76.16}
\]
where
\[
 C_{\rm flux}=2^{-1/2}C_2C_\xi.                      \tag{76.17}
\]
\end{theorem}

\begin{proof}
Insert (76.13) and (76.15) into the exact
\(L^2\!\times L^2\) flux estimate (74.18).  The factors
\(\ell^{-1/2}\) and \(\ell^{1/2}\) cancel in the zeroth-order term,
while the derivative term gains \(\ell\).
\end{proof}

Define the matched face owner
\[
 \Gamma_{n,q}^{\rm face}
 :={\cal H}^{(0)}_{n,q}
 +\ell_{n,q}{\cal H}^{(1)}_{n,q}.                   \tag{76.18}
\]
There is no negative power of the collar width in (76.18).

\subsection{Collective influence constant}

Let \(D_{n,q}\) be the lattice or physical block diameter in the norm
used by the neutral-pair estimate, and let
\[
 m_{n,q}=m_{\rm col}m_{\rm type}                     \tag{76.19}
\]
be the face colouring/type multiplicity of V75.

\begin{corollary}[Paired-face Gram budget]
\label{cor:v16v76-gram}
The paired-face response operator obeys
\[
 \|{\cal R}_{n,q}^{\rm face}\|
 \leq
 2C_dC_{\rm flux}D_{n,q}\sqrt{m_{n,q}}\,
 \Gamma_{n,q}^{\rm face},                            \tag{76.20}
\]
and its Gram contribution to the influence operator is at most
\[
 \kappa_{n,q}^{\rm face}
 :=
 4C_d^2C_{\rm flux}^2D_{n,q}^2m_{n,q}
 \left(\Gamma_{n,q}^{\rm face}\right)^2.             \tag{76.21}
\]
\end{corollary}

\begin{proof}
Use Theorem~\ref{thm:v16v76-flux-load} as the face amplitude
\(F_*\) in Proposition~\ref{prop:v16v75-all-face}, then square.
\end{proof}

A volume-uniform V66 insertion is therefore certified by
\[
 \sup_n\sup_i
 \sum_{q:\,i\sim q}
 D_{n,q}^2m_{n,q}
 \left(
 {\cal H}^{(0)}_{n,q}
 +\ell_{n,q}{\cal H}^{(1)}_{n,q}
 \right)^2
 <\infty,                                            \tag{76.22}
\]
and the resulting supremum must be strictly below the remaining
conditional-gap margin after all other interactions are included.
The incidence notation \(i\sim q\) means that block \(i\) meets the
face shell \(q\).

\subsection{Relation to the V49 Hilbert owners}

The quantities in (76.12) are precisely the one-shell operator norms
of the V49 Hilbert reconstruction maps \(R_n^{(0)}\) and
\(R_n^{(1)}\).  Trace-norm convergence in (49.15) implies operator-norm
tail control, but V49's weighted \(L^1\) array (49.20) does not
automatically imply the square row (76.22).  The new sufficient tail
owner is
\[
 \lim_{Q\to\infty}\sup_n
 \sum_{q\geq Q}
 D_{n,q}^2m_{n,q}
 \left(
 \|P_qR_n^{(0)}\|_{\rm op}
 +\ell_{n,q}\|P_qR_n^{(1)}\|_{\rm op}
 \right)^2=0.                                       \tag{76.23}
\]

\begin{proposition}[Trace-class sufficient condition]
\label{prop:v16v76-traceclass}
If
\[
 \sup_{n,q}D_{n,q}^2m_{n,q}\leq C_D                 \tag{76.24}
\]
and
\[
 R_n^{(0)}\to R^{(0)},\qquad
 \ell_{n,\bullet}R_n^{(1)}\to\widetilde R^{(1)}
 \quad\text{in trace norm with uniform shell localization},    \tag{76.25}
\]
then (76.23) holds.
\end{proposition}

\begin{proof}
For nonnegative numbers \(a_q\),
\[
 \sum_{q\geq Q}a_q^2
 \leq\left(\sum_{q\geq Q}a_q\right)^2.               \tag{76.26}
\]
Operator norms are bounded by trace norms.  The uniform shell
localization theorem used in V49 makes the trace-norm tails of both
families in (76.25) vanish.  Apply (76.26), the triangle inequality,
and (76.24).
\end{proof}

The weighted convergence of
\(\ell_{n,\bullet}R_n^{(1)}\) is a new requirement.  It is weaker than
the \(\ell^{-1}\)-weighted derivative load appearing in some V48
\(L^1\) terms, but it has not been proved by the physical
reconstruction.

\begin{assessment}[What V76 closes]
\label{ass:v16v76-status}
V76 checks the actual V48 theorem and records the norm mismatch
instead of hiding it.  The V48 \(L^1\) collar owner cannot load the
\(L^2\) gauge flux.  A new continuum and discrete \(L^2\) collar
theorem supplies the sharp powers \(\ell^{-1/2}\) and
\(\ell^{1/2}\); after pairing with the gauge trace, they combine into
the intensive owner
\({\cal H}^{(0)}+\ell{\cal H}^{(1)}\).

The resulting all-face Gram budget and a trace-class sufficient tail
condition are explicit.  The remaining missing premise is physical
weighted Hilbert reconstruction convergence together with local
zero-mean gauge control.
\end{assessment}

\begin{decision}[V77 gauge block mean and Poincar\'e ownership]
\label{dec:v16v76-next}
Construct the local zero-mean decomposition used in (76.14) without
removing physical residual gauge modes.  Separate block translations,
prove a stable local Korn--Poincar\'e estimate modulo those constants,
and track the finite-dimensional block-mean interaction through the
infrared Schur complement of V64.
\end{decision}

\begin{firewall}[Claim boundary after V76]
\label{fire:v16v76}
The \(L^2\) collar and flux theorems are analytic compilers.  They do
not prove the required reconstruction convergence, local gauge
Poincar\'e estimate, nonlinear factor response, determinant
orientation, or the Standard Model--Einstein continuum.
\end{firewall}

\section*{Version V76 append-only record}
\addcontentsline{toc}{section}{Version V76 append-only record}
The complete V75 source was retained byte for byte before this
continuation.  It had \(777425\) bytes, \(23322\) lines, and SHA--256
\[
 \texttt{C04AB6D2F61A4D6EE21A52B577D199CF15244A732BE50CD29D941C4D632DC1F6}.
\]
V76 matched the exact V48 collar powers, proved the needed
\(L^2\) trace theorem, and derived the volume-uniform paired-face Gram
budget and its Hilbert reconstruction owner.


\section{V77: local Korn decomposition and the coarse rigid block}
\label{sec:v16v77}

\subsection{Why block mean removal is insufficient}

Let \(Q\subset\mathbb R^d\) be a cube of side \(\ell\), with centre
\(x_Q\).  For a vector field \(\xi\), write
\[
 e(\xi)=\frac12(\nabla\xi+\nabla\xi^{\mathsf T}).     \tag{77.1}
\]
The local Faddeev--Popov energy in V70 controls \(e(\xi)\), not the
full gradient.

\begin{proposition}[Rotational obstruction]
\label{prop:v16v77-rotation}
There is no constant \(C\) such that
\[
 \|\nabla\xi\|_{L^2(Q)}
 \leq C\|e(\xi)\|_{L^2(Q)}                            \tag{77.2}
\]
for every zero-mean \(\xi\in H^1(Q;\mathbb R^d)\).
\end{proposition}

\begin{proof}
Let \(\Omega\ne0\) be skew symmetric and set
\[
 \xi(x)=\Omega(x-x_Q).
\]
The cube symmetry gives \(\int_Q\xi=0\), while
\[
 e(\xi)=\frac12(\Omega+\Omega^{\mathsf T})=0
\]
and \(\nabla\xi=\Omega\ne0\).
\end{proof}

Thus the zero-mean hypothesis used for the gauge trace in V76 is
enough for an ordinary Poincar\'e estimate in terms of
\(\nabla\xi\), but it does not convert that gradient into the
symmetrized-gradient gauge energy.  Local rotations must be retained
with translations in the infrared block.

\subsection{Canonical rigid-motion projection}

Let
\[
 {\cal R}(Q)
 =\{x\mapsto a+\Omega(x-x_Q):
 a\in\mathbb R^d,\ \Omega^{\mathsf T}=-\Omega\}.      \tag{77.3}
\]
Denote by \(\Pi_Q^{\rm rig}\) the \(L^2(Q)\)-orthogonal projection onto
\({\cal R}(Q)\), and write
\[
 \xi=r_Q+\xi_Q^\perp,\qquad
 r_Q=\Pi_Q^{\rm rig}\xi.                              \tag{77.4}
\]

\begin{theorem}[Scale-stable local Korn--Poincar\'e estimate]
\label{thm:v16v77-local-korn}
There is a constant \(C_K(d)\), independent of \(\ell\), such that
\[
 \|\nabla\xi_Q^\perp\|_{L^2(Q)}
 \leq C_K\|e(\xi)\|_{L^2(Q)},                         \tag{77.5}
\]
and
\[
 \|\xi_Q^\perp\|_{L^2(Q)}
 \leq C_K\ell\|e(\xi)\|_{L^2(Q)}.                    \tag{77.6}
\]
Consequently,
\[
 \|\xi_Q^\perp\|_{L^2(\partial Q)}
 \leq C_{\partial K}\ell^{1/2}
 \|e(\xi)\|_{L^2(Q)}.                                \tag{77.7}
\]
\end{theorem}

\begin{proof}
Scale \(Q\) to the unit cube.  On the unit cube, Korn's second
inequality states
\[
 \inf_{r\in{\cal R}(Q)}
 \|\xi-r\|_{H^1}\leq C\|e(\xi)\|_{L^2}.              \tag{77.8}
\]
For completeness, if (77.8) failed on the orthogonal complement of
\({\cal R}(Q)\), there would be a sequence with unit \(H^1\) norm and
\(e(\xi_n)\to0\).  Rellich compactness and weak \(H^1\) convergence
give a nonzero limit satisfying \(e(\xi)=0\).  The kernel of \(e\) on
a connected cube is exactly \({\cal R}(Q)\), contradicting
orthogonality.  This proves (77.8).

The \(H^1\) estimate on the orthogonal complement gives (77.5) and
the unit-scale version of (77.6).  Rescaling derivatives and volume
to side \(\ell\) yields the factor \(\ell\) in (77.6).  Finally the
scaled trace inequality
\[
 \|u\|_{L^2(\partial Q)}^2
 \leq C\left(
 \ell^{-1}\|u\|_{L^2(Q)}^2+
 \ell\|\nabla u\|_{L^2(Q)}^2\right)
\]
combined with (77.5)--(77.6) proves (77.7).
\end{proof}

\subsection{Explicit trace of the rigid part}

Write
\[
 r_Q(x)=a_Q+\Omega_Q(x-x_Q).                          \tag{77.9}
\]

\begin{proposition}[Rigid face scaling]
\label{prop:v16v77-rigid-trace}
There is a dimensional constant \(C_R\) such that
\[
 \|r_Q\|_{L^2(\partial Q)}
 \leq C_R\left(
 \ell^{(d-1)/2}|a_Q|
 +\ell^{(d+1)/2}|\Omega_Q|\right).                   \tag{77.10}
\]
Therefore
\[
 \|\xi\cdot n\|_{L^2(\partial Q)}
 \leq C_{\partial K}\ell^{1/2}\|e(\xi)\|_{L^2(Q)}
 +C_R\left(
 \ell^{(d-1)/2}|a_Q|
 +\ell^{(d+1)/2}|\Omega_Q|\right).                   \tag{77.11}
\]
\end{proposition}

\begin{proof}
On \(\partial Q\), \(|x-x_Q|\leq C_d\ell\).  Integrate the pointwise
bound
\[
 |r_Q(x)|\leq|a_Q|+C_d\ell|\Omega_Q|
\]
over a face measure comparable to \(\ell^{d-1}\), and take square
roots.  Add (77.7) to obtain (77.11).
\end{proof}

The first term in (77.11) is the fluctuation trace used by V76.
The last two are genuine coarse variables; deleting them would remove
residual gauge directions rather than estimate them.

\subsection{Fine/coarse quadratic splitting}

At a finite cutoff, collect all block rigid coordinates
\[
 r=(a_Q,\Omega_Q)_Q                                  \tag{77.12}
\]
and all orthogonal fluctuations into \(z\).  Write the quadratic
residual gauge operator as
\[
 {\cal M}_n=
 \begin{pmatrix}
 A_n&B_n\\
 B_n^*&F_n
 \end{pmatrix}.                                      \tag{77.13}
\]
The local Korn theorem supplies the fine-block floor
\[
 F_n\succeq\gamma_f\ell^{-2}I                        \tag{77.14}
\]
when the fine norm is the block \(L^2\) norm and the mesh geometry is
uniform.

\begin{theorem}[Rigid-mode Schur reduction]
\label{thm:v16v77-schur}
Assume (77.14).  Define
\[
 S_n=A_n-B_nF_n^{-1}B_n^*.                            \tag{77.15}
\]
Then
\[
 \det{\cal M}_n=\det F_n\,\det S_n,                  \tag{77.16}
\]
on any fixed zero-mode complement, and
\[
 \langle r,S_nr\rangle
 =\inf_z
 \left\langle
 \binom rz,
 {\cal M}_n\binom rz
 \right\rangle.                                      \tag{77.17}
\]
In particular \(S_n\succeq0\) whenever
\({\cal M}_n\succeq0\).
\end{theorem}

\begin{proof}
Block Gaussian elimination gives (77.16).  Completing the square,
\[
 \left\langle
 \binom rz,
 {\cal M}_n\binom rz
 \right\rangle
 =
 \langle r,S_nr\rangle
 +\left\langle
 z+F_n^{-1}B_n^*r,\,
 F_n(z+F_n^{-1}B_n^*r)
 \right\rangle.                                     \tag{77.18}
\]
Taking the infimum proves (77.17), and positivity follows.
\end{proof}

The determinant response separates into a uniformly local fine factor
and the effective coarse rigid operator:
\[
 -\log\det{\cal M}_n
 =-\log\det F_n-\log\det S_n.                        \tag{77.19}
\]
This is the correct location for the block translations and rotations.

\subsection{Infrared status of the coarse operator}

If all blocks carry the same translation \(a_Q=a\) and the same
rotation compatible with the periodic geometry, the assembled vector
field is a global rigid motion and the gauge energy vanishes.  These
global modes must be removed once.  After their removal, a slowly
varying family \(r_Q\) still has energy of difference type; its lowest
positive eigenvalue decays with the system diameter.

\begin{proposition}[No volume-uniform coarse inverse from local Korn]
\label{prop:v16v77-coarse-no-go}
The fine estimate (77.14) alone gives no volume-independent bound on
\(\|S_n^{-1}\|\).  In particular, for a periodic chain of block
translations with energy
\[
 \sum_Q|a_{Q+1}-a_Q|^2,                               \tag{77.20}
\]
the first positive eigenvalue is
\[
 4\sin^2(\pi/N_{\rm block})
 \longrightarrow0.                                   \tag{77.21}
\]
\end{proposition}

\begin{proof}
The translation energy (77.20) is a possible Schur form consistent
with a uniform fine floor.  Its Fourier eigenvalues are
\(4\sin^2(\pi k/N_{\rm block})\); the first positive one is
(77.21).  Therefore no conclusion about a uniform coarse inverse
follows from (77.14).
\end{proof}

The outcome matches V64: the coarse block is massless.  It must be
handled through a conductance projection and differentiated response,
not through a uniform inverse norm.

\subsection{Corrected gauge trace ledger}

Combining (76.13) with (77.11) splits the face flux into
\[
 F_{n,q}=F_{n,q}^{\rm fine}+F_{n,q}^{\rm rig}.        \tag{77.22}
\]
The fine term has the V76 owner
\[
 \|F_{n,q}^{\rm fine}\|
 \leq C_{\rm flux}
 \left({\cal H}^{(0)}_{n,q}
 +\ell_{n,q}{\cal H}^{(1)}_{n,q}\right)
 \|e(\xi)\|_{L^2(Q_i\cup Q_j)}.                      \tag{77.23}
\]
The rigid term is explicitly finite dimensional:
\[
 \|F_{n,q}^{\rm rig}\|
 \leq C
 \left(
 \ell^{-1/2}{\cal H}^{(0)}+\ell^{1/2}{\cal H}^{(1)}
 \right)
 \left(
 \ell^{(d-1)/2}|a_Q|
 +\ell^{(d+1)/2}|\Omega_Q|\right).                   \tag{77.24}
\]
It must be assigned to the response of \(S_n\), not charged to the
fine Korn constant.

\begin{assessment}[What V77 closes]
\label{ass:v16v77-status}
V77 corrects the local gauge decomposition by retaining rotations as
well as translations.  It proves a scale-stable Korn--Poincar\'e and
face-trace estimate on the orthogonal fluctuation space, then derives
the exact Schur determinant and energy reduction onto coarse rigid
coordinates.

The fine face flux retains the V76 intensive owner.  The rigid flux
belongs to a massless coarse Schur operator whose inverse is not
volume uniform.  Its differentiated determinant response, rather than
its inverse norm, is the next owner.
\end{assessment}

\begin{decision}[V78 coarse rigid conductance symbol]
\label{dec:v16v77-next}
Compute the Schur operator for piecewise rigid translations and
rotations on a regular block lattice.  Identify the discrete
translation and microrotation symbols, remove only the true global
rigid kernel, and test whether the operator factors through a
positive matrix conductance as required by V69.  Preserve any
hourglass or uncoupled rotation modes as exact obstructions.
\end{decision}

\begin{firewall}[Claim boundary after V77]
\label{fire:v16v77}
The local Korn theorem is rigorous, but the coarse Schur operator has
not been derived from the actual metric discretization.  Its
conductance factorization, rigid face response, conditional gap,
determinant orientation, and continuum transfer remain unproved.
\end{firewall}

\section*{Version V77 append-only record}
\addcontentsline{toc}{section}{Version V77 append-only record}
The complete V76 source was retained byte for byte before this
continuation.  It had \(788359\) bytes, \(23678\) lines, and SHA--256
\[
 \texttt{D73D58A8D3C4F0959A1377385E1AA960B97026D8E4572E7E9049E9F1D0521D31}.
\]
V77 proved the local rigid-motion Korn decomposition, the corrected
gauge face ledger, and the fine/coarse Schur reduction that isolates
the massless rigid block.


\section{V78: coarse rigid symbol and the rotational hourglass}
\label{sec:v16v78}

\subsection{Piecewise rigid face mismatch}

Partition a flat periodic lattice into equal cubes of side \(\ell\).
At block \(Q\), let
\[
 r_Q(x)=a_Q+\Omega_Q(x-x_Q),\qquad
 \Omega_Q^{\mathsf T}=-\Omega_Q.                     \tag{78.1}
\]
For the neighbour \(Q+\mu\) in coordinate direction \(e_\mu\), the
shared face centre is
\[
 x_{Q,\mu}=x_Q+\frac\ell2e_\mu
 =x_{Q+\mu}-\frac\ell2e_\mu.                         \tag{78.2}
\]
The centre mismatch is
\[
 J_{Q,\mu}
 :=r_{Q+\mu}(x_{Q,\mu})-r_Q(x_{Q,\mu})
 =D_\mu a_Q-\frac\ell2
 (\Omega_{Q+\mu}+\Omega_Q)e_\mu,                     \tag{78.3}
\]
where \(D_\mu a_Q=a_{Q+\mu}-a_Q\).

Consider first the centre-sampled energy
\[
 E_0(a,\Omega)=\sum_{Q,\mu}|J_{Q,\mu}|^2.            \tag{78.4}
\]

\subsection{Exact Fourier symbol}

For block momentum \(k\), put
\[
 s_\mu(k)=e^{ik_\mu}-1,\qquad
 c_\mu(k)=e^{ik_\mu}+1.                              \tag{78.5}
\]

\begin{proposition}[Centre-sampled symbol]
\label{prop:v16v78-symbol}
The Fourier form of (78.4) is
\[
 E_0(k;a,\Omega)
 =
 \sum_{\mu=1}^d
 \left|
 s_\mu(k)a-\frac\ell2c_\mu(k)\Omega e_\mu
 \right|^2.                                          \tag{78.6}
\]
\end{proposition}

\begin{proof}
On a Fourier mode,
\[
 D_\mu a_Q=s_\mu(k)a\,e^{ik\cdot Q},\qquad
 \Omega_{Q+\mu}+\Omega_Q
 =c_\mu(k)\Omega\,e^{ik\cdot Q}.
\]
Insert these identities in (78.3) and sum the squared block
amplitudes.
\end{proof}

\subsection{An exact hourglass obstruction}

\begin{theorem}[Checkerboard rotational kernel]
\label{thm:v16v78-hourglass}
Assume the block lattice admits
\[
 k_\pi=(\pi,\ldots,\pi).                              \tag{78.7}
\]
Then for every nonzero skew matrix \(\Omega\),
\[
 E_0(k_\pi;0,\Omega)=0.                               \tag{78.8}
\]
Thus the centre-sampled face energy has a rotational kernel unrelated
to a global rigid motion.
\end{theorem}

\begin{proof}
At \(k_\pi\),
\[
 s_\mu=-2,\qquad c_\mu=0
\]
for every \(\mu\).  Setting \(a=0\) makes every term in (78.6) vanish,
independently of \(\Omega\).  The mode alternates sign from block to
block and is therefore not a spatially constant rigid motion.
\end{proof}

This is the coarse analogue of a finite-element hourglass mode.  A
Schur operator derived only from face-centre matching cannot satisfy
the V69 conductance certificate.

\subsection{A rotational-difference repair}

Add
\[
 E_\gamma(a,\Omega)
 =
 E_0(a,\Omega)
 +\gamma\ell^2
 \sum_{Q,\mu}|D_\mu\Omega_Q|_F^2,
 \qquad \gamma>0.                                    \tag{78.9}
\]
Its Fourier symbol is
\[
 E_\gamma(k;a,\Omega)
 =
 \sum_\mu
 \left|s_\mu a-\frac\ell2c_\mu\Omega e_\mu\right|^2
 +\gamma\ell^2|s(k)|^2|\Omega|_F^2.                 \tag{78.10}
\]

\begin{theorem}[Kernel classification after rotational repair]
\label{thm:v16v78-kernel}
On a periodic block lattice,
\[
 E_\gamma(k;a,\Omega)=0
\]
if and only if
\[
 k=0,\qquad \Omega=0,                                \tag{78.11}
\]
with \(a\) arbitrary.  Thus the only kernel is the space of global
translations.
\end{theorem}

\begin{proof}
If \(k\ne0\), then \(|s(k)|>0\).  The final term in (78.10) forces
\(\Omega=0\).  The remaining terms are
\(\sum_\mu|s_\mu a|^2=|s(k)|^2|a|^2\), so \(a=0\).
If \(k=0\), then \(s_\mu=0\) and \(c_\mu=2\).  Equation (78.10)
reduces to
\[
 \ell^2\sum_\mu|\Omega e_\mu|^2
 =\ell^2|\Omega|_F^2,
\]
which vanishes exactly when \(\Omega=0\); \(a\) is free.
\end{proof}

\begin{corollary}[Coarse Korn symbol]
\label{cor:v16v78-coarse-korn}
There is \(c_{\gamma,d}>0\) such that
\[
 E_\gamma(k;a,\Omega)
 \geq c_{\gamma,d}\left[
 |s(k)|^2|a|^2+
 \ell^2(1+|s(k)|^2)|\Omega|_F^2
 \right]                                             \tag{78.12}
\]
for every \(k\), after omitting the \(k=0\) translation kernel.
\end{corollary}

\begin{proof}
Divide the two sides by the positive quadratic norm on the right.
The resulting Rayleigh quotient is continuous on the compact
Brillouin zone and unit sphere in \((a,\ell\Omega)\).  By Theorem
\ref{thm:v16v78-kernel}, it never vanishes on the declared complement,
so it has a positive minimum \(c_{\gamma,d}\).
\end{proof}

The translation branch remains massless with symbol
\(|s(k)|^2\).  The rotation branch is gapped at \(k=0\) and controlled
at the edge of the Brillouin zone.

\subsection{Positive square-factor representation}

Define the local factor
\[
 {\cal B}_\gamma(a,\Omega)
 =
 \left(
 \left(
 D_\mu a-\frac\ell2(T_\mu+I)\Omega e_\mu
 \right)_\mu,\,
 \left(\sqrt\gamma\,\ell D_\mu\Omega\right)_\mu
 \right),                                            \tag{78.13}
\]
where \(T_\mu\) is the block shift.  Then
\[
 E_\gamma(a,\Omega)=\|{\cal B}_\gamma(a,\Omega)\|^2,
\qquad
 S_\gamma={\cal B}_\gamma^*{\cal B}_\gamma.          \tag{78.14}
\]

\begin{theorem}[Coarse conductance eligibility]
\label{thm:v16v78-eligibility}
After removal of global translations, \(S_\gamma\) is positive and
injective.  If its face and rotation-difference squares are weighted
by uniformly positive-definite matrix conductances, then the general
factor projection of Proposition
\ref{prop:v16v70-general-factor} applies.  The determinant response
therefore has a volume-uniform projection owner at fixed factor, with
the fibre dimension
\[
 r_{\rm rig}=d+\frac{d(d-1)}2=\frac{d(d+1)}2.         \tag{78.15}
\]
\end{theorem}

\begin{proof}
Equation (78.14) and Theorem~\ref{thm:v16v78-kernel} give positivity
and injectivity on the quotient.  Positive matrix weights replace the
identity in (78.14) by an operator \(A\succ0\), yielding
\(S_\gamma={\cal B}_\gamma^*A{\cal B}_\gamma\).
Proposition~\ref{prop:v16v70-general-factor} then gives the
orthogonal projection.  The dimension in (78.15) counts translations
and skew matrices per block.
\end{proof}

The theorem proves eligibility of the repaired model, not that the
actual Schur complement (77.15) contains the added
rotation-difference term with \(\gamma>0\).

\subsection{Bounded-domain rigid kernel}

On a nonperiodic Euclidean domain, a global rigid rotation has
\[
 \Omega_Q=\Omega,\qquad a_Q=a+\Omega x_Q.             \tag{78.16}
\]
Substitution in (78.3) gives zero centre mismatch.  The
rotation-difference term also vanishes.  Thus the bounded-domain
kernel consists of all global rigid motions, as expected.  It must be
removed or pinned once; it is not an hourglass defect.

\subsection{Physical reading}

The artificial-looking term in (78.9) has a natural possible source.
Matching the full affine field across a face, rather than sampling
only the face centre, detects tangential variation of
\(\Omega_{Q+\mu}-\Omega_Q\).  The next version tests whether the
integrated face mismatch produces (78.9) automatically and computes
its coefficient from the face second moment.

\begin{assessment}[What V78 closes]
\label{ass:v16v78-status}
V78 computes the exact translation--microrotation symbol of a
centre-sampled piecewise-rigid face energy and finds a checkerboard
rotational kernel.  Adding a positive rotation-difference term removes
every spurious zero mode, leaving only true global translations on the
torus, and yields a positive square factor eligible for the V69
determinant projection.

The repaired coefficient has not been derived from the actual Schur
operator.  It must arise from full-face reconstruction or be declared
as an auxiliary stabilization, with the latter reintroducing the
source-decoupling obligations of V57.
\end{assessment}

\begin{decision}[V79 full-face moment calculation]
\label{dec:v16v78-next}
Integrate the squared mismatch of two affine rigid fields over their
entire shared face.  Prove orthogonal separation of the face-centre
jump and tangential rotation difference, compute the exact second
moment of a cube face, and determine which rotation components are
controlled after summing all face orientations.  Reject centre
quadrature if it is the sole source of the hourglass.
\end{decision}

\begin{firewall}[Claim boundary after V78]
\label{fire:v16v78}
The coarse energy is a regular-lattice model, not the computed Schur
complement of the nonlinear metric gauge.  Its weights, moving-factor
response, collar reconstruction, conditional gap, and continuum
transfer remain unproved.
\end{firewall}

\section*{Version V78 append-only record}
\addcontentsline{toc}{section}{Version V78 append-only record}
The complete V77 source was retained byte for byte before this
continuation.  It had \(798219\) bytes, \(23999\) lines, and SHA--256
\[
 \texttt{45BB7C910B249819D51B0DCBDF174DCB1AC2F69D042947503283AF13E133ED1B}.
\]
V78 derived the coarse rigid symbol, proved the checkerboard
rotational obstruction, and identified the positive
rotation-difference repair required for conductance eligibility.


\section{V79: full-face moments remove the rotational hourglass}
\label{sec:v16v79}

\subsection{Affine mismatch on one face}

Let two cubes \(Q_-\) and \(Q_+\) of side \(\ell\) share the face
orthogonal to \(e_\mu\), with face centre \(x_f\).  Their centres are
\[
 x_-=x_f-\frac\ell2e_\mu,\qquad
 x_+=x_f+\frac\ell2e_\mu.                            \tag{79.1}
\]
Let
\[
 r_\pm(x)=a_\pm+\Omega_\pm(x-x_\pm),
 \qquad \Omega_\pm^{\mathsf T}=-\Omega_\pm.           \tag{79.2}
\]
Write every face point as
\[
 x=x_f+y,\qquad y\cdot e_\mu=0,\qquad
 |y_\nu|\leq\frac\ell2\quad(\nu\ne\mu).              \tag{79.3}
\]

\begin{lemma}[Centre/tangential decomposition]
\label{lem:v16v79-decomposition}
The face mismatch is
\[
 r_+(x)-r_-(x)
 =J_\mu+\Delta_\mu\Omega\,y,                          \tag{79.4}
\]
where
\[
 J_\mu
 =a_+-a_--\frac\ell2(\Omega_++\Omega_-)e_\mu,
 \qquad
 \Delta_\mu\Omega=\Omega_+-\Omega_-.                 \tag{79.5}
\]
\end{lemma}

\begin{proof}
Use
\[
 x-x_+=y-\frac\ell2e_\mu,\qquad
 x-x_-=y+\frac\ell2e_\mu
\]
in (79.2) and collect the constant and \(y\)-dependent terms.
\end{proof}

\subsection{Exact cube-face second moment}

Let \(P_\mu=I-e_\mu\otimes e_\mu\) be tangential projection.

\begin{lemma}[Face moment tensor]
\label{lem:v16v79-moment}
For the centred \((d-1)\)-dimensional cube face,
\[
 \int_fy\,dS=0,\qquad
 \int_fy\otimes y\,dS
 =\frac{\ell^{d+1}}{12}P_\mu.                        \tag{79.6}
\]
\end{lemma}

\begin{proof}
Odd integrals vanish by reflection symmetry.  Mixed second moments
also vanish.  For a tangential coordinate \(\nu\ne\mu\),
\[
 \int_fy_\nu^2\,dS
 =
 \left(\int_{-\ell/2}^{\ell/2}t^2\,dt\right)
 \ell^{d-2}
 =\frac{\ell^{d+1}}{12}.
\]
There is no normal coordinate on the face, proving (79.6).
\end{proof}

\begin{theorem}[Exact full-face rigid energy]
\label{thm:v16v79-face-energy}
The integrated mismatch is
\[
 \int_f|r_+-r_-|^2\,dS
 =
 \ell^{d-1}|J_\mu|^2
 +\frac{\ell^{d+1}}{12}
 \|\Delta_\mu\Omega\,P_\mu\|_F^2.                   \tag{79.7}
\]
\end{theorem}

\begin{proof}
Insert (79.4).  The cross term vanishes by the first identity in
(79.6).  For the quadratic term,
\[
 \int_f|\Delta_\mu\Omega\,y|^2\,dS
 =
 \operatorname{Tr}\left[
 \Delta_\mu\Omega^*\Delta_\mu\Omega
 \int_fy\otimes y\,dS\right],
\]
and the second identity in (79.6) gives (79.7).
\end{proof}

The first term is the centre-sampled energy of V78.  The second is the
missing positive rotational-difference term, with exact coefficient
\(1/12\).

\subsection{Fourier symbol and kernel}

After division by the common face area \(\ell^{d-1}\), the periodic
full-face energy has symbol
\[
 E_{\rm face}(k;a,\Omega)
 =
 \sum_{\mu=1}^d\left[
 \left|
 s_\mu a-\frac\ell2c_\mu\Omega e_\mu
 \right|^2
 +\frac{\ell^2}{12}|s_\mu|^2
 \|\Omega P_\mu\|_F^2
 \right].                                            \tag{79.8}
\]

\begin{theorem}[No full-face hourglass]
\label{thm:v16v79-no-hourglass}
For \(d\geq2\),
\[
 E_{\rm face}(k;a,\Omega)=0
\]
if and only if
\[
 k=0,\qquad\Omega=0,                                 \tag{79.9}
\]
with \(a\) arbitrary.  Hence the only periodic kernel is the true
global translation space.
\end{theorem}

\begin{proof}
Suppose first that \(k\ne0\), and choose \(\mu\) with
\(s_\mu\ne0\).  The second term for that \(\mu\) in (79.8) implies
\[
 \Omega P_\mu=0.                                     \tag{79.10}
\]
Thus \(\Omega e_\nu=0\) for every \(\nu\ne\mu\).  Skew symmetry also
forces every component involving \(e_\mu\) to vanish, so
\(\Omega=0\).  The first terms in (79.8) then give
\(|s(k)|^2|a|^2=0\), hence \(a=0\).

If \(k=0\), then \(s_\mu=0\) and \(c_\mu=2\).  The first terms reduce
to
\[
 \ell^2\sum_\mu|\Omega e_\mu|^2
 =\ell^2|\Omega|_F^2,
\]
so \(\Omega=0\), while \(a\) remains arbitrary.
\end{proof}

\begin{corollary}[Full-face coarse Korn bound]
\label{cor:v16v79-korn}
There is \(c_d>0\) such that, after removing the \(k=0\) translation
kernel,
\[
 E_{\rm face}(k;a,\Omega)
 \geq c_d\left[
 |s(k)|^2|a|^2+
 \ell^2(1+|s(k)|^2)|\Omega|_F^2
 \right].                                            \tag{79.11}
\]
\end{corollary}

\begin{proof}
Use the compactness argument from Corollary
\ref{cor:v16v78-coarse-korn}, with the kernel classification in
Theorem~\ref{thm:v16v79-no-hourglass}.
\end{proof}

\subsection{The checkerboard calculation}

At \(k_\pi=(\pi,\ldots,\pi)\) and \(a=0\), the centre terms vanish,
but the full-face contribution is
\[
 E_{\rm face}(k_\pi;0,\Omega)
 =
 \frac{\ell^2}{3}
 \sum_{\mu=1}^d\|\Omega P_\mu\|_F^2.                 \tag{79.12}
\]
Since
\[
 \sum_{\mu=1}^d\|\Omega P_\mu\|_F^2
 =
 \sum_\mu\left(
 \|\Omega\|_F^2-|\Omega e_\mu|^2\right)
 =(d-1)\|\Omega\|_F^2,                               \tag{79.13}
\]
one obtains
\[
 E_{\rm face}(k_\pi;0,\Omega)
 =\frac{d-1}{3}\ell^2\|\Omega\|_F^2>0               \tag{79.14}
\]
for \(\Omega\ne0\).  Thus centre quadrature is exactly the source of
the V78 hourglass in this model.

\subsection{Weighted face metric}

Let the face mismatch be measured by a positive matrix weight
\(W_{Q,\mu}(y)\) satisfying
\[
 w_-I\preceq W_{Q,\mu}(y)\preceq w_+I.               \tag{79.15}
\]
Define
\[
 E_W=\sum_{Q,\mu}\int_{f_{Q,\mu}}
 \langle r_{Q+\mu}-r_Q,\,
 W_{Q,\mu}(r_{Q+\mu}-r_Q)\rangle\,dS.                \tag{79.16}
\]

\begin{proposition}[Elliptic full-face corridor]
\label{prop:v16v79-weighted}
One has
\[
 w_-E_{\rm face}\leq E_W\leq w_+E_{\rm face}.        \tag{79.17}
\]
Consequently the kernel remains the true rigid kernel and the coarse
Korn constant in (79.11) is multiplied by at least \(w_-\).
\end{proposition}

\begin{proof}
Apply (79.15) pointwise to every face integrand and sum.  The lower
bound implies that \(E_W=0\) only when \(E_{\rm face}=0\).
\end{proof}

The energy (79.16) is a positive square factor with local matrix
conductances.  Therefore V69--V71 apply at a fixed reconstruction
provided the face weights and factor variations satisfy their
ellipticity and response bounds.

\subsection{Schur-complement criterion}

\begin{theorem}[Physical origin criterion]
\label{thm:v16v79-schur-origin}
Let \(S_n\) be the rigid Schur complement (77.15).  If there is a
uniform \(c_S>0\) such that
\[
 \langle r,S_nr\rangle
 \geq c_SE_W(r)                                      \tag{79.18}
\]
for all coarse rigid fields modulo the global kernel, then:
\begin{enumerate}
\item \(S_n\) has no hourglass modes;
\item it satisfies the coarse Korn bound
      \(c_Sw_-c_d\) times the right side of (79.11);
\item its determinant is eligible for the square-factor projection
      response of V70.
\end{enumerate}
\end{theorem}

\begin{proof}
Combine (79.18), Proposition~\ref{prop:v16v79-weighted}, and
Corollary~\ref{cor:v16v79-korn}.  A positive quadratic form has a
square root, and the explicit local face factor in (79.16) supplies
the conductance representation for the lower comparison.
\end{proof}

The comparison (79.18) is the remaining physical finite-element
owner.  It cannot be replaced by pointwise face-centre agreement.

\begin{assessment}[What V79 closes]
\label{ass:v16v79-status}
V79 proves that full-face integration automatically supplies the
rotation-difference term introduced provisionally in V78.  The exact
cube-face second moment is \(\ell^{d+1}/12\), and the resulting
Fourier symbol has no checkerboard or other spurious kernel.

Thus no auxiliary rotational stabilizer is needed in the coarse rigid
model if the actual Schur energy dominates a uniformly elliptic
full-face mismatch.  That domination, and the variation of its face
weights under nonlinear metric reconstruction, remain unproved.
\end{assessment}

\begin{decision}[V80 companion audit and quadrature certificate]
\label{dec:v16v79-next}
Perform the mandatory YM/NS companion audit.  Then characterize the
positive face quadrature rules that preserve (79.7): compute the
weighted centroid and covariance matrix, prove that positive-definite
tangential covariance is exactly the condition excluding rotational
hourglasses, and give a minimal symmetric rule suitable for the
finite cutoff.
\end{decision}

\begin{firewall}[Claim boundary after V79]
\label{fire:v16v79}
The calculation concerns affine piecewise-rigid fields on regular
flat cubes.  It does not prove the Schur domination (79.18), nonlinear
weight response, curved-face analogue, determinant orientation,
conditional gaps, or continuum transfer.
\end{firewall}

\section*{Version V79 append-only record}
\addcontentsline{toc}{section}{Version V79 append-only record}
The complete V78 source was retained byte for byte before this
continuation.  It had \(806970\) bytes, \(24277\) lines, and SHA--256
\[
 \texttt{8CF19B580235E7EB1DF23003F2092A9A85E7EA1E7246DF3AF670093711D53B8E}.
\]
V79 computed the exact full-face second moment, removed the rotational
hourglass without an auxiliary stabilizer, and isolated the Schur
domination criterion for the physical coarse operator.


\section{V80: companion audit and a positive face-quadrature certificate}
\label{sec:v16v80}

\subsection{Mandatory YM/NS companion audit}

\begin{audit}[V80 companion checkpoint]
\label{aud:v16v80-companions}
The newest active numeric companion notes inspected were
\[
\begin{array}{c|r|r|l}
\hbox{file}&\hbox{lines}&\hbox{bytes}&\hbox{SHA--256}\\ \hline
\texttt{Renewal\_Yang\_Mills\_Mass\_Gap\_Research\_Notes\_V60.tex}
&14002&513804&
\texttt{4D77E2498272E2A824F380CD0F7A7A14255A475CCE8AEE61D15FA9DCA8894A74}\\
\texttt{Renewal\_NS\_Stability\_Research\_Notes\_V26.tex}
&5319&278283&
\texttt{82C887F8C2C69E4F28FAD7C3DC8DE5B9099CB5A4BF431EBA1FFF3E91935978AD}.
\end{array}                                             \tag{80.1}
\]
\end{audit}

YM V60 now proves both direct Wilson response signs on paired
first-quarter-face plates, while correctly withholding a
quarter-cell or interacting-measure claim.  NS V26 remains unchanged.

\begin{theorem}[V80 import decision]
\label{thm:v16v80-audit-decision}
No companion result supplies a face moment or metric gauge
quadrature.  YM's use of compatible local charts again supports an
atlas strategy, but none of its exact floors enters the quadrature
certificate below.
\end{theorem}

\begin{proof}
The companion operators contain no piecewise-rigid face mismatch,
tangential covariance, or Faddeev--Popov factor.  Their hypotheses are
therefore disjoint from those used below.
\end{proof}

\subsection{Moment decomposition of a positive rule}

Let a flat face \(f\) have tangential vector space \(T_f\) of
dimension \(r=d-1\).  A positive quadrature rule consists of points
\[
 y_\alpha\in T_f,\qquad w_\alpha>0,\qquad
 1\leq\alpha\leq m.                                  \tag{80.2}
\]
Define its total weight, centroid, and centred covariance by
\[
 W=\sum_\alpha w_\alpha,\qquad
 \bar y=W^{-1}\sum_\alpha w_\alpha y_\alpha,          \tag{80.3}
\]
\[
 C=\sum_\alpha w_\alpha
 (y_\alpha-\bar y)\otimes(y_\alpha-\bar y).          \tag{80.4}
\]

\begin{theorem}[Exact quadrature square completion]
\label{thm:v16v80-square}
For every vector \(J\) and linear map \(L:T_f\to\mathbb R^d\),
\[
 \sum_\alpha w_\alpha|J+Ly_\alpha|^2
 =
 W|J+L\bar y|^2+\operatorname{Tr}(L^*LC).            \tag{80.5}
\]
\end{theorem}

\begin{proof}
Write \(y_\alpha=\bar y+z_\alpha\), where
\(\sum_\alpha w_\alpha z_\alpha=0\).  Expand the square.  The cross
term vanishes, the constant term is the first term in (80.5), and
\[
 \sum_\alpha w_\alpha|Lz_\alpha|^2
 =\operatorname{Tr}\left(
 L^*L\sum_\alpha w_\alpha z_\alpha\otimes z_\alpha
 \right).
\]
\end{proof}

\subsection{Exactness and coercivity}

For a cube face of side \(\ell\), the continuum moments from V79 are
\[
 W_f=\ell^{d-1},\qquad
 \bar y_f=0,\qquad
 C_f=\frac{\ell^{d+1}}{12}P_f.                       \tag{80.6}
\]

\begin{theorem}[Affine-square exactness criterion]
\label{thm:v16v80-exactness}
A positive quadrature rule reproduces
\[
 \int_f|J+Ly|^2\,dS                                  \tag{80.7}
\]
for every \(J\) and \(L\) if and only if
\[
 W=\ell^{d-1},\qquad
 \bar y=0,\qquad
 C=\frac{\ell^{d+1}}{12}P_f.                         \tag{80.8}
\]
\end{theorem}

\begin{proof}
Sufficiency follows from Theorem~\ref{thm:v16v80-square} and
(80.6).  For necessity, first take \(L=0\) to identify \(W\).
Compare the term linear in \(J\) for arbitrary \(J,L\) to obtain
\(\bar y=0\).  With those fixed, equality for all \(L\) identifies
the quadratic form \(\operatorname{Tr}(L^*LC)\), hence \(C=C_f\).
\end{proof}

\begin{proposition}[Tangential affine coercivity]
\label{prop:v16v80-covariance}
The quadrature controls every tangential affine part:
\[
 \sum_\alpha w_\alpha|Ly_\alpha|^2
 \geq c_{\rm cov}\ell^{d+1}|L|_{\rm HS}^2            \tag{80.9}
\]
for all \(L:T_f\to\mathbb R^d\) if and only if
\[
 C\succeq c_{\rm cov}\ell^{d+1}P_f.                  \tag{80.10}
\]
\end{proposition}

\begin{proof}
By (80.5) with \(J=0\) and \(\bar y=0\), the left side is
\(\operatorname{Tr}(L^*LC)\).  The operator inequality (80.10)
implies (80.9).  Conversely, choose \(Lz=\langle z,v\rangle u\) for
unit \(u\in\mathbb R^d\) and arbitrary tangent \(v\).  Then (80.9)
becomes
\(\langle v,Cv\rangle\geq c_{\rm cov}\ell^{d+1}|v|^2\),
which is (80.10).
\end{proof}

Positive-definite tangential covariance is therefore exactly the
certificate for all affine tangential mismatches.  It is sufficient
for excluding rotational hourglasses.  If it fails, a particular
coupled skew sector may still be controlled by other face
orientations, so failure of one covariance certificate is not by
itself a global hourglass theorem.

\subsection{Point-count lower bound}

\begin{proposition}[Minimum number of positive points]
\label{prop:v16v80-minimum}
If \(C\) is positive definite on the \(r\)-dimensional tangent space,
then
\[
 m\geq r+1=d.                                        \tag{80.11}
\]
\end{proposition}

\begin{proof}
The centred vectors
\[
 z_\alpha=y_\alpha-\bar y
\]
satisfy the nontrivial positive relation
\(\sum_\alpha w_\alpha z_\alpha=0\).  Hence their span has dimension
at most \(m-1\).  Since
\[
 \operatorname{Ran}C
 \subseteq\operatorname{span}\{z_\alpha\},
\]
positive definiteness of \(C\) requires \(r\leq m-1\).
\end{proof}

Thus one centre point is maximally deficient: its covariance is zero.

\subsection{Minimal physical four-dimensional rule}

In spacetime dimension \(d=4\), the face is three dimensional.  In
orthonormal tangential coordinates define four points
\[
\begin{aligned}
 y_1&=\frac{\ell}{\sqrt{12}}(1,1,1),&
 y_2&=\frac{\ell}{\sqrt{12}}(1,-1,-1),\\
 y_3&=\frac{\ell}{\sqrt{12}}(-1,1,-1),&
 y_4&=\frac{\ell}{\sqrt{12}}(-1,-1,1),
\end{aligned}                                        \tag{80.12}
\]
with weights
\[
 w_\alpha=\frac{\ell^3}{4}.                          \tag{80.13}
\]
Every coordinate has magnitude \(\ell/\sqrt{12}<\ell/2\), so all
points lie in the face.

\begin{theorem}[Exact tetrahedral face rule]
\label{thm:v16v80-tetrahedron}
The four-point rule (80.12)--(80.13) satisfies
\[
 W=\ell^3,\qquad
 \bar y=0,\qquad
 C=\frac{\ell^5}{12}I_3.                             \tag{80.14}
\]
It therefore integrates every affine mismatch square exactly and
uses the minimum possible number of positive points.
\end{theorem}

\begin{proof}
The four sign vectors in (80.12) sum to zero.  Their off-diagonal
coordinate products cancel, while every diagonal square equals
\(\ell^2/12\).  Hence
\[
 \frac14\sum_{\alpha=1}^4y_\alpha y_\alpha^{\mathsf T}
 =\frac{\ell^2}{12}I_3.
\]
Multiplication by total weight \(\ell^3\) gives (80.14).
Theorem~\ref{thm:v16v80-exactness} gives exactness, and Proposition
\ref{prop:v16v80-minimum} proves minimality.
\end{proof}

\subsection{General positive rule and hourglass exclusion}

For every face orientation \(\mu\), let the rule have centroid zero
and
\[
 C_\mu\succeq c_{\rm cov}\ell^{d+1}P_\mu,
 \qquad c_{\rm cov}>0.                               \tag{80.15}
\]
Then the quadrature analogue of (79.8) obeys
\[
 E_{\rm quad}(k;a,\Omega)
 \geq
 W_-\sum_\mu
 \left|s_\mu a-\frac\ell2c_\mu\Omega e_\mu\right|^2
 +c_{\rm cov}\ell^{d+1}
 \sum_\mu|s_\mu|^2\|\Omega P_\mu\|_F^2,              \tag{80.16}
\]
where \(W_-=\inf_\mu W_\mu>0\).

\begin{corollary}[Quadrature kernel certificate]
\label{cor:v16v80-kernel}
Under (80.15), the periodic quadrature energy has only the global
translation kernel.
\end{corollary}

\begin{proof}
The proof of Theorem~\ref{thm:v16v79-no-hourglass} uses only
strict positivity of the centre and tangential covariance terms.
Equation (80.16) supplies both.
\end{proof}

\subsection{Cutoff implementation ledger}

The tetrahedral rule is exact for piecewise-rigid affine fields, but a
nonlinear reconstructed mismatch is not affine along the face.  A
valid cutoff use must split
\[
 u(y)=J+Ly+\varepsilon(y)                             \tag{80.17}
\]
and bound the quadrature remainder.  If
\[
 \sum_\alpha w_\alpha|\varepsilon(y_\alpha)|^2
 \leq\eta_{\rm quad}
 \left[
 W|J|^2+c_{\rm cov}\ell^{d+1}|L|^2
 \right],\qquad \eta_{\rm quad}<1,                   \tag{80.18}
\]
then
\[
 \sum_\alpha w_\alpha|u(y_\alpha)|^2
 \geq
 (1-\sqrt{\eta_{\rm quad}})^2
 \left[
 W|J|^2+c_{\rm cov}\ell^{d+1}|L|^2
 \right].                                            \tag{80.19}
\]

\begin{proof}
Let \(v_\alpha=J+Ly_\alpha\), and use the weighted direct-sum norm.
The affine covariance certificate gives
\[
 \|v\|_w^2
 \geq W|J|^2+c_{\rm cov}\ell^{d+1}|L|^2.
\]
Equation (80.18) therefore implies
\(\|\varepsilon\|_w\leq
\sqrt{\eta_{\rm quad}}\|v\|_w\).  The reverse triangle inequality
gives
\[
 \|v+\varepsilon\|_w
 \geq(1-\sqrt{\eta_{\rm quad}})\|v\|_w.
\]
Square and apply the affine lower bound once more to obtain (80.19).
\end{proof}

The exact requirement is a strict remainder margin relative to the
moment floor.

\begin{assessment}[What V80 closes]
\label{ass:v16v80-status}
The companion audit imports no gauge estimate.  V80 gives the exact
moment characterization of positive face quadrature, proves that
positive tangential covariance is equivalent to coercivity for every
affine tangential mismatch, and establishes the minimum point count.

In four dimensions an explicit four-point tetrahedral rule reproduces
the full-face affine energy exactly, including the \(1/12\) second
moment, and therefore removes the centre-quadrature hourglass with no
auxiliary source.  Curved faces and nonlinear quadrature remainders
remain to be controlled.
\end{assessment}

\begin{decision}[V81 curved-face covariance stability]
\label{dec:v16v80-next}
Transport the tetrahedral rule through a curved face chart.  Quantify
how bilipschitz distortion, density weights, and centroid drift change
the covariance floor.  Prove an explicit perturbation radius that
preserves positive covariance and then load the nonlinear mismatch
remainder from tangential \(W^{2,2}\) reconstruction norms.
\end{decision}

\begin{firewall}[Claim boundary after V80]
\label{fire:v16v80}
The quadrature theorem is exact for flat affine faces.  It does not
prove curved-chart covariance, nonlinear remainder bounds, Schur
domination, moving determinant response, conditional gaps, or the
continuum measure.
\end{firewall}

\section*{Version V80 append-only record}
\addcontentsline{toc}{section}{Version V80 append-only record}
The complete V79 source was retained byte for byte before this
continuation.  It had \(815863\) bytes, \(24582\) lines, and SHA--256
\[
 \texttt{45F3024932D72CEFB2B70CFB82B96DDE734DE20293D3B9B248E10147BE2DDCF3}.
\]
V80 recorded the mandatory companion audit and proved the exact
positive face-quadrature certificate, including the minimal
four-dimensional tetrahedral rule.


\section{V81: curved-face covariance and nonlinear quadrature remainder}
\label{sec:v16v81}

\subsection{A perturbation theorem for weighted covariance}

Let \(T\) be an \(r\)-dimensional Euclidean tangent space.  A reference
positive rule \((y_\alpha,w_\alpha)\) has
\[
 W_0=\sum_\alpha w_\alpha=\ell^r,\qquad
 \sum_\alpha w_\alpha y_\alpha=0,\qquad
 C_0=\sum_\alpha w_\alpha y_\alpha\otimes y_\alpha
 \succeq c_0\ell^{r+2}I_T.                           \tag{81.1}
\]
Let a perturbed rule be
\[
 y_\alpha'=Ty_\alpha+e_\alpha,\qquad
 w_\alpha'>0,                                        \tag{81.2}
\]
where
\[
 \sigma_{\min}(T)\geq\sigma_->0,\qquad
 |e_\alpha|\leq\varepsilon_y\ell,\qquad
 w_\alpha'\geq(1-\varepsilon_w)w_\alpha              \tag{81.3}
\]
with \(0\leq\varepsilon_w<1\).

Denote the perturbed weighted centroid by \(\bar y'\) and covariance
by
\[
 C'=\sum_\alpha w_\alpha'
 (y_\alpha'-\bar y')\otimes(y_\alpha'-\bar y').       \tag{81.4}
\]

\begin{theorem}[Covariance stability without centring by hand]
\label{thm:v16v81-covariance}
If
\[
 \varepsilon_y<\sqrt{c_0}\,\sigma_-,                 \tag{81.5}
\]
then
\[
 C'\succeq
 (1-\varepsilon_w)
 \left(\sqrt{c_0}\,\sigma_--\varepsilon_y\right)^2
 \ell^{r+2}I_T.                                      \tag{81.6}
\]
\end{theorem}

\begin{proof}
For a unit \(v\in T\), weighted variance has the variational form
\[
 v^{\mathsf T}C'v
 =\inf_{c\in\mathbb R}
 \sum_\alpha w_\alpha'
 \left(v\cdot y_\alpha'-c\right)^2.                  \tag{81.7}
\]
The weight bound gives at least \(1-\varepsilon_w\) times the same
infimum with weights \(w_\alpha\).  In the weighted
\(\ell^2(w)\) quotient by constants,
\[
 \operatorname{dist}
 \bigl((v\cdot Ty_\alpha)_\alpha,\mathbb R{\bf1}\bigr)
 =
 \left(v^{\mathsf T}TC_0T^{\mathsf T}v\right)^{1/2}
 \geq\sqrt{c_0}\,\sigma_-\ell^{(r+2)/2}.             \tag{81.8}
\]
The perturbation vector
\((v\cdot e_\alpha)_\alpha\) has weighted norm at most
\[
 \varepsilon_y\ell W_0^{1/2}
 =\varepsilon_y\ell^{(r+2)/2}.                       \tag{81.9}
\]
Distance to a closed subspace is one-Lipschitz, so subtracting
(81.9) from (81.8), squaring, and multiplying by
\(1-\varepsilon_w\) proves (81.6).
\end{proof}

The centroid drift is already included in the minimization (81.7);
no separate recentering estimate is needed.

\subsection{Curved chart loading}

Let a curved \(r\)-face be parameterized near its centre by
\[
 \Phi(y)=x_f+Ey+\rho(y),                              \tag{81.10}
\]
where \(E:T\to T_{x_f}f\) is an isometry and
\[
 |\rho(y)|\leq K_\Phi|y|^2,\qquad
 |J_\Phi(y)-1|\leq K_J\ell                           \tag{81.11}
\]
on the reference face.  Transport the quadrature points to a common
tangent frame and use sampled density weights
\[
 y_\alpha'=E y_\alpha+\rho(y_\alpha),\qquad
 w_\alpha'=w_\alpha J_\Phi(y_\alpha).                \tag{81.12}
\]

\begin{corollary}[Curved tetrahedral covariance floor]
\label{cor:v16v81-curved}
For the four-dimensional tetrahedral rule of V80,
\[
 c_0=\frac1{12},\qquad r=3.                           \tag{81.13}
\]
There is a dimensional \(C_\rho\) such that if
\[
 C_\rho K_\Phi\ell<\frac1{\sqrt{12}},
 \qquad K_J\ell<1,                                   \tag{81.14}
\]
then
\[
 C'\succeq
 (1-K_J\ell)
 \left(\frac1{\sqrt{12}}-C_\rho K_\Phi\ell\right)^2
 \ell^5 I_3.                                         \tag{81.15}
\]
\end{corollary}

\begin{proof}
At the tetrahedral nodes, \(|y_\alpha|\leq C\ell\), so
(81.11) gives
\[
 |\rho(y_\alpha)|\leq C_\rho K_\Phi\ell^2.
\]
Thus Theorem~\ref{thm:v16v81-covariance} applies with
\(\sigma_-=1\), \(\varepsilon_y=C_\rho K_\Phi\ell\), and
\(\varepsilon_w=K_J\ell\).
\end{proof}

The covariance loss is linear in the dimensionless chart curvature
\(K_\Phi\ell\) before squaring.  A collapsing face size therefore
helps this geometric part, provided the reconstruction derivatives do
not grow faster.

\subsection{Bramble--Hilbert point remainder}

Let \(\widehat f\) be the unit reference face and let
\(\Pi_1\) be the \(H^2\)-stable projection onto affine
vector-valued functions.  In physical dimension four the face
dimension is \(r=3\), and
\[
 H^2(\widehat f)\hookrightarrow C^0(\widehat f).      \tag{81.16}
\]

\begin{theorem}[Four-dimensional nodal affine remainder]
\label{thm:v16v81-remainder}
Let \(u\in H^2(f;\mathbb R^m)\), let
\[
 u_{\rm aff}=\Pi_1u=J+Ly,\qquad
 \varepsilon=u-u_{\rm aff}.                          \tag{81.17}
\]
For any fixed finite reference quadrature transported by a uniformly
bilipschitz face chart,
\[
 \sum_\alpha w_\alpha'
 |\varepsilon(y_\alpha')|^2
 \leq C_{\rm BH}\ell^4
 \|\nabla_T^2u\|_{L^2(f)}^2.                         \tag{81.18}
\]
The constant depends on the reference rule and chart bounds, not on
\(\ell\).
\end{theorem}

\begin{proof}
On the unit reference face, point evaluation is continuous on
\(H^2\) by (81.16).  On the kernel-orthogonal complement of affine
functions, the Bramble--Hilbert lemma gives
\[
 \|\widehat u-\Pi_1\widehat u\|_{H^2}
 \leq C\|D^2\widehat u\|_{L^2}.
\]
Evaluation at finitely many nodes and summation with bounded positive
weights therefore gives
\[
 \sum_\alpha \widehat w_\alpha
 |\widehat\varepsilon(\widehat y_\alpha)|^2
 \leq C\|D^2\widehat u\|_{L^2(\widehat f)}^2.         \tag{81.19}
\]
Under \(y=\ell\widehat y\),
\[
 \|D^2\widehat u\|_{L^2(\widehat f)}^2
 =\ell^{4-r}\|\nabla_T^2u\|_{L^2(f)}^2,
\]
while the physical quadrature weights contribute \(\ell^r\).
Their product is \(\ell^4\), proving (81.18).  Uniform chart
distortion only changes the constant.
\end{proof}

The restriction to a three-dimensional face matters:
\(H^2\) point evaluation is continuous because \(2>3/2\).  The same
argument would require higher Sobolev order on faces of dimension at
least four.

\subsection{Nonlinear mismatch lower bound}

Let the curved positive rule have total weight \(W'\) and covariance
floor
\[
 C'\succeq c_{\rm curv}\ell^{r+2}I.                  \tag{81.20}
\]
For \(u=J+Ly+\varepsilon\), define
\[
 E_{\rm aff}
 :=W'|J+L\bar y'|^2+
 \operatorname{Tr}(L^*LC').                          \tag{81.21}
\]

\begin{theorem}[Quadrature remainder corridor]
\label{thm:v16v81-corridor}
For every \(\tau\in(0,1)\),
\[
 \sum_\alpha w_\alpha'|u(y_\alpha')|^2
 \geq
 (1-\tau)E_{\rm aff}
 -(\tau^{-1}-1)C_{\rm BH}\ell^4
 \|\nabla_T^2u\|_{L^2(f)}^2.                         \tag{81.22}
\]
If
\[
 C_{\rm BH}\ell^4\|\nabla_T^2u\|_{L^2(f)}^2
 \leq\eta E_{\rm aff},\qquad
 \eta<\tau,                                           \tag{81.23}
\]
then the right side has the positive floor
\[
 \left[
 (1-\tau)-(\tau^{-1}-1)\eta
 \right]E_{\rm aff}.                                 \tag{81.24}
\]
\end{theorem}

\begin{proof}
In the weighted nodal Hilbert space, apply
\[
 |v+e|^2\geq(1-\tau)|v|^2-(\tau^{-1}-1)|e|^2
\]
with \(v=u_{\rm aff}\), then use (81.18) and the exact moment
identity (80.5).  Under (81.23), factor out \(E_{\rm aff}\);
the coefficient in (81.24) is positive precisely under its stated
inequality.
\end{proof}

\subsection{Reconstruction owner}

For a reconstructed face mismatch \(u_{n,q}(a)\) in a unit soft
direction, define
\[
 {\cal T}^{(2)}_{n,q}
 :=\sup_{\|a\|_S=1}
 \|\nabla_T^2u_{n,q}(a)\|_{L^2(f_{n,q})}.             \tag{81.25}
\]
The negative quadrature remainder entering the coarse Schur floor is
\[
 {\cal E}^{\rm quad}_{n,q}
 =
 C_{\rm BH}\ell_{n,q}^4
 \left({\cal T}^{(2)}_{n,q}\right)^2.                \tag{81.26}
\]
This is a face \(W^{2,2}\) owner.  V54 introduced
\(W^{2,2}\) reconstruction control for the EGHY second variation, but
its bulk-to-face trace and shell tail must be checked in the current
tangential norm; V81 does not identify them automatically.

\begin{assessment}[What V81 closes]
\label{ass:v16v81-status}
V81 proves quantitative stability of the face covariance under
bilipschitz point distortion and positive density perturbation,
including centroid drift.  The four-dimensional tetrahedral rule
retains a positive floor whenever \(K_\Phi\ell\) and \(K_J\ell\) stay
inside an explicit radius.

A scaled Bramble--Hilbert theorem controls the nonlinear nodal
remainder by \(\ell^4\) times a tangential \(W^{2,2}\) face norm.
The remaining owner is a physical reconstruction tail for that norm
and a proof that the rigid Schur energy dominates the curved
quadrature energy.
\end{assessment}

\begin{decision}[V82 Schur domination by face mismatch]
\label{dec:v16v81-next}
Prove a two-block extension theorem: among fine corrections that make
neighbouring traces compatible, the symmetrized-gradient energy is
bounded below by \(\ell^{-1}\) times the squared face mismatch modulo
one common rigid motion.  Use it to derive the Schur comparison
(79.18), with the curved quadrature remainder (81.26) subtracted
explicitly.
\end{decision}

\begin{firewall}[Claim boundary after V81]
\label{fire:v16v81}
The curved covariance theorem assumes a controlled face chart and
positive sampled density.  The nonlinear remainder bound does not
prove its reconstruction owner, the Schur domination, determinant
response, conditional gaps, or continuum transfer.
\end{firewall}

\section*{Version V81 append-only record}
\addcontentsline{toc}{section}{Version V81 append-only record}
The complete V80 source was retained byte for byte before this
continuation.  It had \(826528\) bytes, \(24915\) lines, and SHA--256
\[
 \texttt{D2FE96D7F848B31600A8B0667ACA4D1EDCD69F942E74723EA85BFC1A6ADE52B4}.
\]
V81 proved curved-face covariance stability and the
four-dimensional \(W^{2,2}\) nonlinear quadrature-remainder corridor.


\section{V82: Schur domination by compatible face extension}
\label{sec:v16v82}

\subsection{Two neighbouring blocks}

Let \(Q_-\) and \(Q_+\) be adjacent \(d\)-dimensional cubes of side
\(\ell\), with shared face \(f\).  On each block decompose
\[
 u_\pm=r_\pm+z_\pm,                                   \tag{82.1}
\]
where \(r_\pm\in{\cal R}(Q_\pm)\) are rigid motions and
\[
 z_\pm\perp{\cal R}(Q_\pm)
 \quad\text{in }L^2(Q_\pm).                           \tag{82.2}
\]
Assume trace compatibility,
\[
 u_-|_f=u_+|_f.                                      \tag{82.3}
\]
Then the rigid mismatch
\[
 j_f=(r_+-r_-)|_f                                    \tag{82.4}
\]
satisfies
\[
 j_f=-(z_+-z_-)|_f.                                  \tag{82.5}
\]

\begin{theorem}[Two-block extension lower bound]
\label{thm:v16v82-two-block}
There is \(c_{\rm ext}(d)>0\), independent of \(\ell\), such that
\[
 \|e(z_-)\|_{L^2(Q_-)}^2+
 \|e(z_+)\|_{L^2(Q_+)}^2
 \geq
 c_{\rm ext}\ell^{-1}\|j_f\|_{L^2(f)}^2.             \tag{82.6}
\]
\end{theorem}

\begin{proof}
The local Korn--Poincar\'e estimates (77.5)--(77.6), followed by the
scaled trace inequality, give
\[
 \|z_\pm\|_{L^2(f)}^2
 \leq C\ell\|e(z_\pm)\|_{L^2(Q_\pm)}^2.              \tag{82.7}
\]
Using (82.5),
\[
 \|j_f\|_{L^2(f)}^2
 \leq2\|z_-\|_{L^2(f)}^2+2\|z_+\|_{L^2(f)}^2.
\]
Insert (82.7) and rearrange.
\end{proof}

The factor \(\ell^{-1}\) is forced by scaling: bulk
symmetrized-gradient energy has dimension
\(\ell^{d-2}\), while a squared face displacement has dimension
\(\ell^{d-1}\).

\subsection{All internal faces}

Let \({\cal Q}\) be a shape-regular cubical partition with uniformly
bounded face degree.  For each block use the decomposition
\[
 u_Q=r_Q+z_Q,\qquad z_Q\perp{\cal R}(Q),              \tag{82.8}
\]
and assume that \(u\) has compatible traces on every internal face.

\begin{theorem}[Global face-extension inequality]
\label{thm:v16v82-global}
There is \(c_{\rm ext}>0\), depending only on dimension and shape
regularity, such that
\[
 \sum_Q\|e(z_Q)\|_{L^2(Q)}^2
 \geq
 c_{\rm ext}\ell^{-1}
 \sum_{f=Q\mid Q'}\|r_{Q'}-r_Q\|_{L^2(f)}^2.         \tag{82.9}
\]
\end{theorem}

\begin{proof}
For every internal face,
\[
 (r_{Q'}-r_Q)|_f=-(z_{Q'}-z_Q)|_f.
\]
Sum the squared triangle bound over faces.  Each block trace appears
at most its uniformly bounded face degree.  The trace estimate
(82.7), summed over blocks, gives
\[
 \sum_f\|r_{Q'}-r_Q\|_{L^2(f)}^2
 \leq C\ell\sum_Q\|e(z_Q)\|_{L^2(Q)}^2.
\]
Rearrange.
\end{proof}

\subsection{Schur complement of the compatible energy}

Let the fine/coarse quadratic form be
\[
 {\cal E}_n(r,z)
 =2\sum_Q\|\delta^*z_Q\|_{L^2(Q)}^2,                 \tag{82.10}
\]
on the affine space of fine corrections satisfying the compatibility
constraints with \(r\).  Since \(e(z)=\delta^*z\), define the coarse
Schur energy
\[
 \langle r,S_nr\rangle
 :=\inf_{z\ {\rm compatible}}{\cal E}_n(r,z).         \tag{82.11}
\]

\begin{theorem}[Flat full-face Schur domination]
\label{thm:v16v82-flat-schur}
For a flat cubical partition,
\[
 \langle r,S_nr\rangle
 \geq
 2c_{\rm ext}\ell^{-1}
 \sum_{f=Q\mid Q'}
 \int_f|r_{Q'}-r_Q|^2\,dS.                           \tag{82.12}
\]
Consequently the Schur operator has only the true global rigid kernel
and dominates the V79 full-face conductance factor.
\end{theorem}

\begin{proof}
Apply Theorem~\ref{thm:v16v82-global} to every admissible \(z\) in
(82.11), multiply by two, and take the infimum.  The kernel statement
then follows from Theorem~\ref{thm:v16v79-no-hourglass}, with global
rigid motions interpreted according to the boundary geometry.
\end{proof}

This proves the comparison (79.18) for the declared conforming flat
model, with its necessary factor \(\ell^{-1}\).  It does not yet prove
that the programme's actual reconstructed gauge fields use this
conforming extension space.

\subsection{Uniformly curved blocks}

Assume each block chart has metric and volume density uniformly
equivalent to the flat reference and satisfies local Korn and trace
constants
\[
 C_K\leq C_K^*,\qquad C_{\rm tr}\leq C_{\rm tr}^*.   \tag{82.13}
\]
Let the physical gauge energy obey
\[
 {\cal E}_{g,n}(z)
 \geq c_g\sum_Q\|e_g(z_Q)\|_{L^2(Q,g)}^2             \tag{82.14}
\]
with \(c_g>0\).

\begin{corollary}[Curved Schur domination]
\label{cor:v16v82-curved-schur}
Under (82.13)--(82.14), there is
\(c_{\rm ext}^g>0\), independent of the number of blocks, such that
\[
 \langle r,S_{g,n}r\rangle
 \geq c_{\rm ext}^g\ell^{-1}
 \sum_f\|j_f\|_{L^2(f,g)}^2.                          \tag{82.15}
\]
\end{corollary}

\begin{proof}
Repeat the proof of Theorem~\ref{thm:v16v82-global} with the uniform
curved Korn and trace constants, then insert (82.14) and take the
fine-field infimum.
\end{proof}

\subsection{Affine part and nonlinear remainder}

In a curved face chart, split the coarse mismatch as
\[
 j_f=u_{\rm aff}+\varepsilon,\qquad
 u_{\rm aff}=J+Ly.                                   \tag{82.16}
\]
For every \(\tau\in(0,1)\),
\[
 \|j_f\|_{L^2(f)}^2
 \geq
 (1-\tau)\|u_{\rm aff}\|_{L^2(f)}^2
 -(\tau^{-1}-1)\|\varepsilon\|_{L^2(f)}^2.           \tag{82.17}
\]
The continuum Bramble--Hilbert estimate gives
\[
 \|\varepsilon\|_{L^2(f)}^2
 \leq C_{\rm BH}^{L^2}\ell^4
 \|\nabla_T^2j_f\|_{L^2(f)}^2.                       \tag{82.18}
\]

\begin{theorem}[Curved affine Schur corridor]
\label{thm:v16v82-affine-corridor}
Let the curved affine moment form satisfy
\[
 \|u_{\rm aff}\|_{L^2(f,g)}^2
 \geq
 c_{\rm mom}\left[
 \ell^{d-1}|J|^2+\ell^{d+1}|L|^2
 \right].                                            \tag{82.19}
\]
Then
\[
\begin{split}
 \langle r,S_{g,n}r\rangle
 \geq c_{\rm ext}^g\ell^{-1}\sum_f\bigg\{
 &(1-\tau)c_{\rm mom}
 \left[\ell^{d-1}|J_f|^2+\ell^{d+1}|L_f|^2\right]\\
 &-(\tau^{-1}-1)C_{\rm BH}^{L^2}\ell^4
 \|\nabla_T^2j_f\|_{L^2(f)}^2
 \bigg\}.                                            \tag{82.20}
\end{split}
\]
\end{theorem}

\begin{proof}
Insert (82.17)--(82.19) into (82.15), face by face.
\end{proof}

The negative remainder in the Schur energy has the scaled form
\[
 {\cal R}^{\rm Schur}_{f}
 =
 C\ell^3\|\nabla_T^2j_f\|_{L^2(f)}^2.                \tag{82.21}
\]
The power is \(\ell^3\), not \(\ell^4\), because extension from a
face to a bulk energy costs \(\ell^{-1}\).

\subsection{A strict domination margin}

\begin{corollary}[No-hourglass curved corridor]
\label{cor:v16v82-margin}
If, for some \(\eta_{\rm Schur}<1\),
\[
 \sum_f
 C_{\rm BH}^{L^2}\ell^4
 \|\nabla_T^2j_f\|_{L^2(f)}^2
 \leq
 \eta_{\rm Schur}
 \sum_f c_{\rm mom}
 \left[\ell^{d-1}|J_f|^2+\ell^{d+1}|L_f|^2\right],   \tag{82.22}
\]
then choosing
\[
 \eta_{\rm Schur}<\tau<1                             \tag{82.23}
\]
gives
\[
\begin{split}
 \langle r,S_{g,n}r\rangle
 \geq c_{\rm ext}^g\ell^{-1}
 (1-\tau)\left(1-\frac{\eta_{\rm Schur}}{\tau}\right)
 \sum_f c_{\rm mom}
 \left[\ell^{d-1}|J_f|^2+\ell^{d+1}|L_f|^2\right].
                                                               \tag{82.24}
\end{split}
\]
\end{corollary}

\begin{proof}
The negative coefficient in (82.20) is
\(\tau^{-1}-1=(1-\tau)/\tau\).  Apply (82.22), factor the positive
affine form, and use (82.23).
\end{proof}

This is the exact Schur counterpart of the V81 quadrature remainder
margin.

\subsection{Determinant consequence}

Under (82.24), the coarse factor has the local full-face square
structure of V79 with a strictly positive residual coefficient.
After removal of the true global rigid kernel:
\begin{enumerate}
\item the coarse determinant is positive at each fixed background;
\item fixed-factor conductance responses have the projection owner of
      V69--V70;
\item moving-factor responses still require the V71--V76 trace-ideal
      and collar bounds.
\end{enumerate}

\begin{assessment}[What V82 closes]
\label{ass:v16v82-status}
V82 proves that compatible fine corrections pay at least
\(\ell^{-1}\) times the squared coarse rigid face mismatch.  This
establishes the previously missing Schur domination for a conforming
flat cubical model and extends it to uniformly curved blocks under
explicit Korn, trace, and energy-equivalence constants.

The curved nonlinear mismatch subtracts the exact
\(\ell^3\|\nabla_T^2j_f\|_2^2\) bulk-energy remainder.  A strict
Bramble--Hilbert margin preserves the full-face no-hourglass floor.
The actual reconstruction has not yet been verified to use the
compatible extension space or satisfy that margin.
\end{assessment}

\begin{decision}[V83 finite-cutoff positivity certificate]
\label{dec:v16v82-next}
Assemble the local determinant constraint, residual gauge slice,
fine Korn factor, coarse full-face Schur factor, and their two
Jacobians into one finite-cutoff density.  State and prove a single
certificate theorem for positivity and a volume-uniform Poincar\'e
gap using V66, with every unresolved physical input appearing as a
named numerical margin rather than an implicit assumption.
\end{decision}

\begin{firewall}[Claim boundary after V82]
\label{fire:v16v82}
The extension theorem concerns a declared conforming block model.
It does not prove that the programme's metric reconstruction satisfies
the compatibility, curved Korn, Bramble--Hilbert, determinant
response, conditional-gap, reflection, or continuum hypotheses.
\end{firewall}

\section*{Version V82 append-only record}
\addcontentsline{toc}{section}{Version V82 append-only record}
The complete V81 source was retained byte for byte before this
continuation.  It had \(836065\) bytes, \(25222\) lines, and SHA--256
\[
 \texttt{1A81BBD11E45B59675FFC972D8F92840C11602061E0086A3B3B571285832E7C4}.
\]
V82 proved the compatible face-extension inequality and the flat and
curved Schur domination corridors, including the nonlinear
face-remainder margin.


\section{V83: a single finite-cutoff positive-law and gap certificate}
\label{sec:v16v83}

\subsection{The constrained gauge-fixed density}

Let \({\cal X}_n\) be the finite-dimensional cutoff field space,
including metric, gauge, Higgs, and any bosonized or determinant
matter variables retained at this stage.  Let
\[
 C_n(x)=0                                             \tag{83.1}
\]
be the local determinant constraint and let
\[
 {\cal F}_n(x)=0                                      \tag{83.2}
\]
be the residual volume-preserving gauge condition on an oriented
chart.  Denote
\[
 J_{C,n}=\sqrt{\det(DC_nDC_n^{\mathsf T})},\qquad
 \Delta_{{\rm FP},n}=\det D_\xi{\cal F}_n.            \tag{83.3}
\]
Let \(W_{{\rm mat},n}\) be the complete matter weight after the
declared finite-cutoff fermion operation.  On the physical slice
\({\cal P}_n\), define
\[
 d\mu_n
 =
 Z_n^{-1}
 e^{-S_{{\rm bos},n}}
 W_{{\rm mat},n}
 J_{C,n}^{-1}
 \Delta_{{\rm FP},n}\,
 d{\rm vol}_{{\cal P}_n}.                            \tag{83.4}
\]

Whenever every factor is positive, the effective potential is
\[
 U_n
 =S_{{\rm bos},n}
 -\log W_{{\rm mat},n}
 +\log J_{C,n}
 -\log\Delta_{{\rm FP},n}.                           \tag{83.5}
\]

\subsection{Named finite-cutoff certificate}

\begin{definition}[Certificate \({\sf C}_{\rm SM-E,n}\)]
\label{def:v16v83-certificate}
The \(n\)-th cutoff satisfies \({\sf C}_{\rm SM-E,n}\) with margins
\[
 (s_C,\delta_{\rm FP},a_-,m_{\rm cov},m_{\rm Schur},
 \rho_{\rm loc},\kappa_{\rm row})                    \tag{83.6}
\]
if all of the following hold.
\begin{enumerate}
\item \textbf{Constraint regularity:}
      \[
      \sigma_{\min}(DC_n)\geq s_C>0.                 \tag{83.7}
      \]
\item \textbf{Gauge orientation and fixed kernel stratum:}
      after removing one declared symmetry space of fixed dimension,
      \[
      \Delta_{{\rm FP},n}\geq\delta_{\rm FP}>0        \tag{83.8}
      \]
      throughout every chart, and oriented overlaps agree.
\item \textbf{Positive matter weight and integrability:}
      \[
      W_{{\rm mat},n}>0
      \quad\mu_n\text{-a.e.},\qquad
      0<Z_n<\infty.                                  \tag{83.9}
      \]
\item \textbf{Fine gauge factor:} the local Korn complement has a
      uniform conditional factor bound and the V76 face owner
      satisfies its weighted Hilbert tail.
\item \textbf{Coarse face covariance:}
      \[
      C_{f,n}\succeq m_{\rm cov}\ell_f^{d+1}P_f
      \quad\text{for every face}.                    \tag{83.10}
      \]
\item \textbf{Schur remainder margin:} the ratio in (82.22) obeys
      \[
      \eta_{{\rm Schur},n}\leq1-m_{\rm Schur},
      \qquad m_{\rm Schur}>0.                        \tag{83.11}
      \]
\item \textbf{One-block conditional gaps:} for a fixed block
      decomposition \(x=(x_i)_{i\in I_n}\),
      \[
      \operatorname{gap}
      \bigl(\mu_{n,i}(\cdot\mid x_{\widehat i})\bigr)
      \geq\rho_{\rm loc}>0                           \tag{83.12}
      \]
      for every exterior configuration in the chart.
\item \textbf{Mixed influence:} with
      \[
      \kappa_{ij,n}
      :=\sup_x\|\nabla_i\nabla_jU_n(x)\|_{\rm op},
      \quad i\ne j,                                  \tag{83.13}
      \]
      one has
      \[
      \sup_i\sum_{j\ne i}\kappa_{ij,n}
      \leq\kappa_{\rm row}<\rho_{\rm loc}.            \tag{83.14}
      \]
\end{enumerate}
\end{definition}

The mixed row in (83.14) includes, without omission,
\[
 \kappa^{\rm FP}+\kappa^{C}
 +\kappa^{\rm vac}+\kappa^R
 +\kappa^{\rm gauge}+\kappa^{\rm Higgs}
 +\kappa^{\rm ferm}.                                 \tag{83.15}
\]
The owners are respectively V68--V76, V59--V61, V46/V57,
V47--V54/V76, and the still-required Standard Model response bounds.

\subsection{Positive-law theorem}

\begin{theorem}[Finite-cutoff probability and uniform gap]
\label{thm:v16v83-positive-gap}
Assume \({\sf C}_{\rm SM-E,n}\) for every \(n\), with all margins in
(83.6) independent of \(n\).  Then (83.4) is a positive probability
measure and
\[
 \operatorname{Var}_{\mu_n}(f)
 \leq
 \frac1{\rho_*}
 \int|\nabla_{\cal P}f|^2\,d\mu_n,\qquad
 \rho_*:=\rho_{\rm loc}-\kappa_{\rm row}>0.           \tag{83.16}
\]
The Poincar\'e constant is independent of the number of blocks.
\end{theorem}

\begin{proof}
Constraint regularity gives the positive coarea density
\(J_{C,n}^{-1}d{\rm vol}_{{\cal P}_n}\) by V59.
Fixed determinant orientation, (83.8), and (83.9) make every factor in
(83.4) positive.  Integrability normalizes it to a probability.

The fine Korn and coarse Schur margins ensure that the declared gauge
coordinates and determinant factors are regular on the chart; their
derivative loads are included in (83.13)--(83.15).  Conditions
(83.12)--(83.14) are exactly the hypotheses of Corollary
\ref{cor:v16v66-row-sum}.  Applying Theorem
\ref{thm:v16v66-block-gap} gives (83.16).
\end{proof}

The theorem does not manufacture any certificate item.  It proves
that the list is sufficient and that no dimension factor is hidden in
the final gap.

\subsection{A sharper influence-matrix form}

Define
\[
 {\cal A}_{n,ii}=\rho_{i,n},\qquad
 {\cal A}_{n,ij}=-\kappa_{ij,n}\quad(i\ne j).         \tag{83.17}
\]

\begin{theorem}[Anisotropic block variance bound]
\label{thm:v16v83-anisotropic}
Under the conditional hypotheses of V66, assume
\[
 {\cal A}_n\succ0.                                   \tag{83.18}
\]
For
\[
 b_i(f)=\left(\int|\nabla_if|^2\,d\mu_n\right)^{1/2},             \tag{83.19}
\]
one has
\[
 \operatorname{Var}_{\mu_n}(f)
 \leq b(f)^{\mathsf T}{\cal A}_n^{-1}b(f).           \tag{83.20}
\]
\end{theorem}

\begin{proof}
Let \(Lh=f-\mathbb E f\) as in the proof of Theorem
\ref{thm:v16v66-block-gap}, and set
\[
 a_i=\left(\int|\nabla_ih|^2\,d\mu_n\right)^{1/2}.
\]
The conditional Bochner argument (66.11)--(66.12), followed by
Cauchy--Schwarz in the mixed block terms, gives
\[
 \operatorname{Var}_{\mu_n}(f)
 =\int(Lh)^2\,d\mu_n
 \geq a^{\mathsf T}{\cal A}_na.                     \tag{83.21}
\]
Integration by parts gives
\[
 \operatorname{Var}_{\mu_n}(f)
 =\int\nabla f\cdot\nabla h\,d\mu_n
 \leq b(f)^{\mathsf T}a.                             \tag{83.22}
\]
By Cauchy--Schwarz in the
\({\cal A}_n\)-inner product,
\[
 b^{\mathsf T}a
 \leq
 (b^{\mathsf T}{\cal A}_n^{-1}b)^{1/2}
 (a^{\mathsf T}{\cal A}_na)^{1/2}.
\]
Combine this with (83.21)--(83.22) and cancel the variance square
root to obtain (83.20).
\end{proof}

This form can retain a positive full-face Gram operator rather than
charging it entrywise to the absolute row.  It is therefore the
preferred finite-cutoff target when the response factorization is
known.

\subsection{Chart gluing requirement}

Suppose oriented charts \(({\cal U}_\alpha,\mu_\alpha)\) cover the
physical slice.  A partition of unity inserted into a density changes
its logarithmic derivatives and can destroy (83.14).  The chart
densities must instead agree on overlaps as measures:
\[
 e^{-U_\alpha}d{\rm vol}_\alpha
 =
 e^{-U_\beta}d{\rm vol}_\beta
 \quad\text{on }{\cal U}_\alpha\cap{\cal U}_\beta.    \tag{83.23}
\]

\begin{proposition}[Overlap cocycle]
\label{prop:v16v83-overlap}
Let \(T_{\alpha\beta}\) be the coordinate transition on an overlap.
Condition (83.23) is equivalent to
\[
 U_\beta\circ T_{\alpha\beta}
 =
 U_\alpha-\log|\det DT_{\alpha\beta}|.               \tag{83.24}
\]
If the coarea and Faddeev--Popov factors are both derived from the
same ambient orbit change of variables, their Jacobians satisfy this
cocycle automatically at finite cutoff.
\end{proposition}

\begin{proof}
Apply the finite-dimensional change-of-variables formula to the right
side of (83.23).  Equality of densities is precisely (83.24).
Jacobians of composed coordinate maps multiply, so their logarithms
obey the cocycle identity.
\end{proof}

This is the measure-theoretic version of the exact overlap discipline
seen in the companion YM atlas.  Determinant sign must still remain
fixed within each oriented chart.

\subsection{Current evidence ledger}

The active notes now prove the following parts of the certificate
under declared reconstruction hypotheses:
\begin{enumerate}
\item local determinant submersion and its exact coarea Jacobian
      (V59--V61);
\item flat and curved-background residual square factors and their
      kernel description (V62, V70);
\item vector conductance and moving-factor determinant compilers
      (V67--V76);
\item local Korn, full-face quadrature, and conforming Schur
      domination (V77--V82);
\item the global block-gap implication (V66 and V83).
\end{enumerate}
The unproved inputs include the positive chiral Standard Model matter
weight, physical one-block gaps, nonlinear atlas coverage and
orientation, the reconstruction tails in V76/V81, the complete
influence row, reflection positivity, tightness, Ward convergence,
and continuum identification.

\begin{assessment}[What V83 closes]
\label{ass:v16v83-status}
V83 assembles the constraint, gauge, fine/coarse determinant, matter,
and interaction layers into one finite-cutoff density and proves a
single sufficient theorem for positivity and a volume-independent
Poincar\'e gap.  The sharper anisotropic form retains operator
structure instead of forcing every response through an absolute row.

The theorem also identifies the exact overlap cocycle needed for a
multi-chart gauge atlas.  It remains conditional because several
physical certificate entries, especially matter positivity and local
conditional gaps, have not been proved.
\end{assessment}

\begin{decision}[V84 localized covariance influence]
\label{dec:v16v83-next}
Strengthen the block theorem from a global variance bound to a
localized covariance estimate.  Add a sweeping-gradient or
resolvent hypothesis that is actually sufficient, derive decay from a
finite-range strict influence matrix, and state clearly why the
Poincar\'e certificate alone does not prove physical clustering.
\end{decision}

\begin{firewall}[Claim boundary after V83]
\label{fire:v16v83}
Certificate \({\sf C}_{\rm SM-E,n}\) is not verified for the physical
cutoff theory.  The theorem proves only its stated implications and
does not establish reflection positivity, exponential clustering,
Hamiltonian reconstruction, or the Standard Model--Einstein
continuum.
\end{firewall}

\section*{Version V83 append-only record}
\addcontentsline{toc}{section}{Version V83 append-only record}
The complete V82 source was retained byte for byte before this
continuation.  It had \(845764\) bytes, \(25546\) lines, and SHA--256
\[
 \texttt{250AE0774D6EA2972BB8C25C9EEAC4BD497C3E02725ECA2272A371FE36F01F76}.
\]
V83 assembled the finite-cutoff positive-law certificate, proved its
uniform block gap and anisotropic variance bound, and derived the
gauge-chart overlap cocycle.


\section{V84: localized covariance from the block Witten resolvent}
\label{sec:v16v84}

\subsection{The gradient intertwining equation}

Retain the block Gibbs law and notation of V66:
\[
 L=-\Delta+\nabla U\cdot\nabla,\qquad
 {\cal A}_{ii}=\rho_i,\qquad
 {\cal A}_{ij}=-\kappa_{ij}.                          \tag{84.1}
\]
Let \(g\) be centered and let
\[
 h=L^{-1}g.                                          \tag{84.2}
\]
Write
\[
 u_i=\nabla_i h,\qquad v_i=\nabla_i g.               \tag{84.3}
\]

\begin{lemma}[Block gradient intertwining]
\label{lem:v16v84-intertwining}
For every block \(i\),
\[
 v_i
 =
 (L+\nabla_i^2U)u_i
 +\sum_{j\ne i}(\nabla_i\nabla_jU)u_j.               \tag{84.4}
\]
\end{lemma}

\begin{proof}
Differentiate \(Lh=g\) in block \(i\).  Derivatives commute with the
Euclidean Laplacian, while
\[
 \nabla_i(\nabla U\cdot\nabla h)
 =
 \nabla U\cdot\nabla(\nabla_i h)
 +\sum_j(\nabla_i\nabla_jU)\nabla_jh.
\]
This is (84.4).
\end{proof}

\subsection{Diagonal one-form coercivity}

\begin{lemma}[Integrated diagonal block floor]
\label{lem:v16v84-diagonal}
Under the conditional Poincar\'e gap (66.2),
\[
 \left\langle
 u_i,(L+\nabla_i^2U)u_i
 \right\rangle_{L^2(\mu)}
 \geq\rho_i\|u_i\|_{L^2(\mu)}^2.                    \tag{84.5}
\]
\end{lemma}

\begin{proof}
Because \(u_i=\nabla_i h\),
\[
\begin{split}
 \langle u_i,Lu_i\rangle
 +\langle u_i,(\nabla_i^2U)u_i\rangle
 =
 \int\bigg[
 &\|\nabla_i^2h\|_{\rm HS}^2
 +\langle\nabla_ih,(\nabla_i^2U)\nabla_ih\rangle\\
 &+\sum_{k\ne i}\|\nabla_k\nabla_ih\|_{\rm HS}^2
 \bigg]d\mu.                                         \tag{84.6}
\end{split}
\]
Condition on \(x_{\widehat i}\).  Lemma
\ref{lem:v16v66-witten} bounds the first line below by
\(\rho_i\int|\nabla_ih|^2\), and the second line is nonnegative.
Integrate the exterior variables.
\end{proof}

\subsection{Componentwise sweeping estimate}

Define
\[
 a_i=\|u_i\|_{L^2(\mu)},\qquad
 b_i=\|v_i\|_{L^2(\mu)}.                             \tag{84.7}
\]

\begin{theorem}[Influence-matrix gradient resolvent]
\label{thm:v16v84-sweeping}
Assume
\[
 \|\,\nabla_i\nabla_jU\,\|_{\rm op}\leq\kappa_{ij}
 \quad(i\ne j)                                       \tag{84.8}
\]
and that
\[
 {\cal A}=D-K,\qquad
 D=\operatorname{diag}(\rho_i),\quad
 K_{ij}=\kappa_{ij}{\bf1}_{i\ne j},                  \tag{84.9}
\]
is a nonsingular \(M\)-matrix.  Equivalently, it is sufficient that
\[
 \theta:=\|D^{-1}K\|_{\ell^\infty\to\ell^\infty}<1. \tag{84.10}
\]
Then
\[
 a\leq{\cal A}^{-1}b
 \quad\text{componentwise}.                          \tag{84.11}
\]
\end{theorem}

\begin{proof}
Take the \(L^2(\mu)\) inner product of (84.4) with \(u_i\).
Lemma~\ref{lem:v16v84-diagonal}, Cauchy--Schwarz, and (84.8) give
\[
 \rho_i a_i^2
 \leq a_ib_i+\sum_{j\ne i}\kappa_{ij}a_ia_j.         \tag{84.12}
\]
If \(a_i>0\), divide by \(a_i\); if \(a_i=0\), the resulting
inequality is automatic.  Thus
\[
 {\cal A}a\leq b                                     \tag{84.13}
\]
componentwise.  Under (84.10),
\[
 {\cal A}^{-1}
 =(I-D^{-1}K)^{-1}D^{-1}
 =\sum_{m=0}^\infty(D^{-1}K)^mD^{-1}                \tag{84.14}
\]
has nonnegative entries.  Multiplying (84.13) by it preserves the
componentwise order and proves (84.11).
\end{proof}

This estimate is stronger than the scalar Poincar\'e conclusion of
V66 and uses the same conditional inputs.  Its extra conclusion comes
from retaining the differentiated Poisson equation block by block.

\subsection{Localized covariance}

For a smooth observable \(f\), put
\[
 b_i(f)=\|\nabla_if\|_{L^2(\mu)}.                    \tag{84.15}
\]

\begin{theorem}[Block covariance matrix bound]
\label{thm:v16v84-covariance}
Under the hypotheses of Theorem~\ref{thm:v16v84-sweeping},
\[
 |\operatorname{Cov}_\mu(f,g)|
 \leq
 b(f)^{\mathsf T}{\cal A}^{-1}b(g).                 \tag{84.16}
\]
\end{theorem}

\begin{proof}
Replace \(g\) by its centered version and let \(h=L^{-1}g\).  Then
\[
 \operatorname{Cov}_\mu(f,g)
 =\int fLh\,d\mu
 =\sum_i\int\nabla_if\cdot\nabla_ih\,d\mu.           \tag{84.17}
\]
Cauchy--Schwarz blockwise gives
\[
 |\operatorname{Cov}_\mu(f,g)|
 \leq b(f)^{\mathsf T}a(h).
\]
Apply the componentwise estimate (84.11), with
\(b(h\hbox{-source})=b(g)\), and use nonnegativity of
\({\cal A}^{-1}\).
\end{proof}

Unlike (83.20), equation (84.16) retains the locations of the
gradient supports.

\subsection{Finite-range exponential decay}

Let the blocks carry a graph distance \({\rm dist}(i,j)\).  Assume
\[
 \kappa_{ij}=0
 \quad\text{when }{\rm dist}(i,j)>R_0,                \tag{84.18}
\]
and
\[
 \rho_i\geq\rho_{\min}>0,\qquad
 \sup_i\frac1{\rho_i}\sum_j\kappa_{ij}\leq\theta<1.  \tag{84.19}
\]

\begin{proposition}[Decay of the inverse influence matrix]
\label{prop:v16v84-inverse-decay}
One has
\[
 0\leq({\cal A}^{-1})_{ij}
 \leq
 \frac1{\rho_{\min}(1-\theta)}
 \theta^{\lceil{\rm dist}(i,j)/R_0\rceil}.            \tag{84.20}
\]
\end{proposition}

\begin{proof}
Use the Neumann expansion (84.14) with
\[
 T=D^{-1}K.
\]
Finite range implies
\((T^m)_{ij}=0\) whenever
\(m<\lceil{\rm dist}(i,j)/R_0\rceil\).  Positivity and the row bound
\(\|T\|_{\infty\to\infty}\leq\theta\) give
\[
 (T^m)_{ij}\leq\theta^m.
\]
Also \(D^{-1}_{jj}\leq\rho_{\min}^{-1}\).  Sum the geometric tail.
\end{proof}

\begin{corollary}[Finite-cutoff block clustering]
\label{cor:v16v84-clustering}
If \(\nabla_if=0\) outside \(S_f\) and
\(\nabla_jg=0\) outside \(S_g\), then
\[
\begin{split}
 |\operatorname{Cov}_\mu(f,g)|
 \leq&
 \frac{
 \theta^{\lceil{\rm dist}(S_f,S_g)/R_0\rceil}
 }{\rho_{\min}(1-\theta)}
 \left(\sum_{i\in S_f}b_i(f)\right)
 \left(\sum_{j\in S_g}b_j(g)\right).
                                                               \tag{84.21}
\end{split}
\]
\end{corollary}

\begin{proof}
Insert (84.20) in (84.16) and restrict the sums to the two gradient
supports.
\end{proof}

This is exponential decay in block-graph distance at a fixed cutoff.
To become a physical correlation length, \(R_0\), the block diameter,
and \(1-\theta\) must have compatible continuum scaling.

\subsection{Why a bare Poincar\'e gap is insufficient}

\begin{proposition}[Uniform gap with polynomial covariance]
\label{prop:v16v84-gap-no-cluster}
There is a sequence of centered Gaussian measures on periodic
one-dimensional lattices with uniformly bounded Poincar\'e constants
but covariances that have a polynomial, rather than exponential,
spatial tail.
\end{proposition}

\begin{proof}
Let
\[
 a_r=(1+|r|)^{-2}\quad(r\in\mathbb Z)
\]
and periodize it on the \(N\)-cycle to obtain a convolution operator
\(A_N\).  Define the covariance
\[
 C_N=I+\varepsilon A_NA_N^*,\qquad \varepsilon>0.    \tag{84.22}
\]
It is positive and
\[
 I\preceq C_N\preceq
 \left(1+\varepsilon\|a\|_{\ell^1(\mathbb Z)}^2\right)I          \tag{84.23}
\]
uniformly in \(N\).  A centered Gaussian with covariance \(C_N\) has
Poincar\'e constant \(\|C_N\|\), so the constants are uniform.

For distances much smaller than \(N\), the off-diagonal covariance
contains
\[
 \varepsilon(A_NA_N^*)_{0r}
 =\varepsilon\sum_s a_s a_{s-r}
 \geq\varepsilon a_0a_r
 =\varepsilon(1+r)^{-2}.                             \tag{84.24}
\]
Thus no uniform exponential upper bound can hold.
\end{proof}

The counterexample separates the V83 global gap from the localized
influence result proved here.

\subsection{Programme insertion}

For the finite-cutoff SM--Einstein law, (84.21) follows if the named
certificate of V83 is strengthened only by:
\begin{enumerate}
\item a block interaction range \(R_0\) in the chosen variables, or an
      exponentially weighted analogue of (84.18);
\item the strict normalized influence margin \(\theta<1\);
\item differentiable local observables with controlled block
      gradients.
\end{enumerate}
The paired-face Gram should be retained as a positive operator when
possible.  Replacing it by long-range entrywise absolute bounds can
destroy (84.18) even though the factor response is harmless.

\begin{assessment}[What V84 closes]
\label{ass:v16v84-status}
V84 upgrades the conditional-gap theorem to a componentwise
Witten-resolvent estimate and a localized covariance bound.  A
finite-range strict influence matrix yields exponential covariance
decay in block distance with an explicit rate.  No separate
sweeping-gradient assumption is needed beyond the V66 conditional
gaps and mixed-Hessian bounds.

A Gaussian counterexample proves that the global Poincar\'e gap alone
does not imply clustering.  Physical clustering still requires
uniform continuum scaling of the block range and influence margin.
\end{assessment}

\begin{decision}[V85 companion audit and physical scaling]
\label{dec:v16v84-next}
Perform the mandatory YM/NS companion audit.  Then convert (84.21)
from block distance to physical distance for block size
\(\ell_n\), allowing an interaction range \(R_{0,n}\) and margin
\(\theta_n\).  Derive the necessary and sufficient asymptotic
condition for a nonzero physical inverse correlation length, and
separate that Euclidean covariance result from an
Osterwalder--Schrader Hamiltonian mass gap.
\end{decision}

\begin{firewall}[Claim boundary after V84]
\label{fire:v16v84}
The covariance theorem assumes a positive finite-cutoff law, uniform
conditional gaps, and controlled mixed Hessians.  Those physical
inputs remain unverified.  Block covariance decay is not reflection
positivity, Hamiltonian reconstruction, or a Standard Model--Einstein
continuum theorem.
\end{firewall}

\section*{Version V84 append-only record}
\addcontentsline{toc}{section}{Version V84 append-only record}
The complete V83 source was retained byte for byte before this
continuation.  It had \(856515\) bytes, \(25861\) lines, and SHA--256
\[
 \texttt{A76800462CEFD63C3115AF46EE802D44F558124436F77F906C506577E623CE36}.
\]
V84 proved the localized block covariance and finite-range decay
theorems and demonstrated that a uniform global gap alone does not
establish clustering.


\section{V85: companion audit and the physical clustering corridor}
\label{sec:v16v85}

\subsection{Mandatory companion audit}

The newest Yang--Mills and Navier--Stokes notes in the shared folder
were inspected before choosing the present direction.  The audited
files are
\[
\begin{array}{lll}
\text{YM V62:}&14424\text{ lines},&
\texttt{0110BCA73FEB2EBCF99A7E3946291D21A9BEB4FBD97FABACD7475A3D2622A5BD},\\
\text{NS V26:}&5319\text{ lines},&
\texttt{82C887F8C2C69E4F28FAD7C3DC8DE5B9099CB5A4BF431EBA1FFF3E91935978AD}.
\end{array}                                                   \tag{85.1}
\]

YM V62 proves exact positivity of both direct signs on paired second
quarter faces.  Its fourth atlas still has only depth \(t\) in the
third coordinate.  Thus it provides a stronger finite-cutoff symbol
input for the gauge part of the present programme, but it does not
provide a full-torus floor, a limiting measure, reflection
positivity, clustering, or a Hamiltonian mass gap.  NS V26 computes
pointwise rank-one norms of the first two Oseen-kernel derivatives.
Its channel-sensitive treatment reinforces the V84 decision to keep
the block matrix and observable-gradient profile, but no
Navier--Stokes statement is imported into the Euclidean field
theory.

\begin{theorem}[Audit decision]
\label{thm:v16v85-audit}
The V84 influence-matrix route remains the shortest noncircular route
from a finite-cutoff positive law to physical clustering.  YM V62 may
eventually strengthen the gauge entry of the V83 certificate, while
NS V26 supplies no hypothesis of the covariance theorem.  Neither
companion note closes a continuum or Osterwalder--Schrader input.
\end{theorem}

\begin{proof}
The YM claim boundary explicitly leaves the third coordinate, the
full torus, and every continuum conclusion open.  The NS estimates
act on the Oseen kernel in the classical three-dimensional flow
problem, not on the Hessian or Gibbs generator of the SM--Einstein
law.  The asserted uses and nonuses follow from the hypotheses of
Theorem~\ref{thm:v16v84-covariance}.
\end{proof}

\subsection{From block distance to physical distance}

At cutoff \(n\), let \(\ell_n>0\) be the physical length assigned to
one edge of the block graph.  We use the induced physical graph
metric
\[
 d_n^{\rm phys}(S,T)
 :=\ell_n d_n^{\rm blk}(S,T).                         \tag{85.2}
\]
For a quasi-isometric geometric realization this definition may be
replaced by Euclidean distance at the price of the fixed
quasi-isometry constants.  Let \(R_{0,n}\geq1\) be the interaction
range in block units and set
\[
\begin{split}
 a_n&:=\ell_nR_{0,n},\\
 \lambda_n&:=-\log\theta_n,\\
 m_n&:=\frac{\lambda_n}{a_n},\\
 G_n(f)&:=\sum_{i\in S_f}
       \|\nabla_i f\|_{L^2(\mu_n)}.                   \tag{85.3}
\end{split}
\]
Here \(a_n\) is the physical range of one influence step and \(m_n\)
is the inverse correlation length certified by the V84 estimate.

\begin{theorem}[Exact physical form of the block estimate]
\label{thm:v16v85-physical}
Suppose the hypotheses of Corollary~\ref{cor:v16v84-clustering} hold
at cutoff \(n\), with parameters
\(\rho_{\min,n}>0\), \(0<\theta_n<1\), and \(R_{0,n}\).
For observables supported on block sets \(S_f,S_g\), put
\[
 r_n=d_n^{\rm phys}(S_f,S_g).
\]
Then
\[
 |\operatorname{Cov}_{\mu_n}(f,g)|
 \leq
 A_n(f,g)e^{-m_nr_n},                                \tag{85.4}
\]
where
\[
 A_n(f,g)
 :=
 \frac{G_n(f)G_n(g)}
 {\rho_{\min,n}(1-\theta_n)}.                        \tag{85.5}
\]
\end{theorem}

\begin{proof}
Since \(r_n=\ell_nd_n^{\rm blk}(S_f,S_g)\),
\[
\begin{split}
 \theta_n^{
 \lceil d_n^{\rm blk}(S_f,S_g)/R_{0,n}\rceil}
 &\leq
 \theta_n^{d_n^{\rm blk}(S_f,S_g)/R_{0,n}}\\
 &=\exp\left\{
 -\frac{-\log\theta_n}{\ell_nR_{0,n}}\,r_n
 \right\}
 =e^{-m_nr_n}.                                      \tag{85.6}
\end{split}
\]
Insert this in (84.21).  Notice that the ceiling only improves the
estimate; no continuum approximation has been used.
\end{proof}

\subsection{Necessary and sufficient scaling for the certified rate}

The rate and the prefactor are logically distinct.  The following
elementary statement records the exact distinction and prevents a
false inference when \(\theta_n\) approaches one.

\begin{proposition}[Uniform exponential majorant criterion]
\label{prop:v16v85-majorant}
Let \(A_n>0\) and \(m_n>0\).  There exist constants \(C<\infty\) and
\(m>0\) such that
\[
 A_ne^{-m_nr}\leq Ce^{-mr}
 \qquad(n\geq1,\ r\geq0)                             \tag{85.7}
\]
if and only if
\[
 \sup_n A_n<\infty,
 \qquad
 \inf_n m_n>0.                                      \tag{85.8}
\]
\end{proposition}

\begin{proof}
If (85.8) holds, choose
\(C=\sup_nA_n\) and \(m=\inf_nm_n\).
Conversely, setting \(r=0\) in (85.7) gives \(A_n\leq C\).
If \(m_n<m\) for some \(n\), division by \(e^{-mr}\) gives
\(A_ne^{(m-m_n)r}\leq C\), which is impossible as
\(r\to\infty\).  Hence \(m_n\geq m\) for every \(n\).
\end{proof}

\begin{corollary}[The physical-rate corridor]
\label{cor:v16v85-rate-corridor}
For the family of upper bounds furnished by V84, a nonzero uniform
physical inverse correlation length is equivalent to
\[
 \liminf_{n\to\infty}
 \frac{-\log\theta_n}{\ell_nR_{0,n}}>0.              \tag{85.9}
\]
For a fixed pair or a normalized class of observables, a uniform
clustering amplitude additionally requires
\[
 \sup_n
 \frac{G_n(f)G_n(g)}
 {\rho_{\min,n}(1-\theta_n)}
 <\infty.                                           \tag{85.10}
\]
\end{corollary}

\begin{proof}
Apply Proposition~\ref{prop:v16v85-majorant} to (85.4), discarding
finitely many cutoffs if the limit inferior is used.  The first
condition in (85.8) is precisely (85.10), and the second is (85.9).
\end{proof}

The \emph{equivalence} in
Corollary~\ref{cor:v16v85-rate-corridor} concerns the exponential
majorants produced by the influence estimate.  Failure of (85.9)
does not prove that the underlying measures fail to cluster; it
proves that this certificate supplies no positive uniform physical
rate.

\subsection{Near-critical influence matrices}

\begin{lemma}[Conversion from influence margin to mass scale]
\label{lem:v16v85-log-margin}
For \(0<\theta<1\),
\[
 1-\theta\leq-\log\theta\leq\frac{1-\theta}{\theta}. \tag{85.11}
\]
Consequently, if \(\theta_n\to1\), then
\[
 \frac{-\log\theta_n}{\ell_nR_{0,n}}
 \sim
 \frac{1-\theta_n}{\ell_nR_{0,n}}.                  \tag{85.12}
\]
In this regime (85.9) is equivalent to
\[
 \liminf_{n\to\infty}
 \frac{1-\theta_n}{\ell_nR_{0,n}}>0.                \tag{85.13}
\]
\end{lemma}

\begin{proof}
The lower bound in (85.11) follows from
\(-\log\theta=\int_\theta^1s^{-1}\,ds\geq1-\theta\);
the upper bound follows by replacing \(s^{-1}\) by
\(\theta^{-1}\).  Division by \(\ell_nR_{0,n}\) and
\(\theta_n\to1\) prove the remaining assertions.
\end{proof}

There are therefore three distinct cases.
\begin{enumerate}
\item If \(a_n=\ell_nR_{0,n}\to0\) and
      \(\sup_n\theta_n<1\), then \(m_n\to\infty\).
\item If \(a_n\to0\), \(\theta_n\to1\), and
      \(1-\theta_n\asymp a_n\), then the certified physical mass
      scale remains finite and positive.
\item If \(1-\theta_n=o(a_n)\), then \(m_n\to0\), so this route loses
      every positive physical rate.
\end{enumerate}
Even the second case has a prefactor problem:
\((1-\theta_n)^{-1}\asymp a_n^{-1}\).  It yields uniform clustering
only if the conditional gap and the renormalized observable
gradients compensate as required by (85.10).

\subsection{Observable normalization is part of the continuum input}

For a class \({\cal O}_n\) of cutoff observables define
\[
 {\cal G}_n:=\sup_{f\in{\cal O}_n}G_n(f).             \tag{85.14}
\]

\begin{theorem}[Uniform clustering for a normalized observable class]
\label{thm:v16v85-class}
Assume
\[
\begin{split}
 &\inf_n\frac{-\log\theta_n}{\ell_nR_{0,n}}
 \geq m_*>0,\\
 &\sup_n
 \frac{{\cal G}_n^2}
 {\rho_{\min,n}(1-\theta_n)}
 \leq C_*<\infty.                                   \tag{85.15}
\end{split}
\]
Then for \(f,g\in{\cal O}_n\),
\[
 |\operatorname{Cov}_{\mu_n}(f,g)|
 \leq C_*e^{-m_*d_n^{\rm phys}(S_f,S_g)}.           \tag{85.16}
\]
\end{theorem}

\begin{proof}
Equations (85.4)--(85.5) and (85.15) give the claim directly.
\end{proof}

The norm \(G_n\) depends on the cutoff coordinate system.  Thus a
field-strength renormalization cannot be appended after the
clustering proof: it must enter the definition of
\({\cal O}_n\), and (85.15) must be verified for the renormalized
observables.

\subsection{Programme insertion}

The V83 certificate should be augmented, at every cutoff, by the
two dimensionless/physical ledgers
\[
\begin{array}{c|c}
\text{rate ledger}&
m_n=(-\log\theta_n)/(\ell_nR_{0,n})\\
\text{amplitude ledger}&
{\cal G}_n^2/[\rho_{\min,n}(1-\theta_n)].
\end{array}                                         \tag{85.17}
\]
The first measures propagation of influence across physical
distance; the second measures the normalization of the chosen
observable sector.  Both are required before taking a clustering
statement through the cutoff limit.

\begin{assessment}[What V85 closes]
\label{ass:v16v85-status}
V85 converts the V84 block estimate exactly into physical units.
For that certificate, (85.9) is the necessary and sufficient
positive-rate condition, while (85.10) is the independent amplitude
condition.  The near-critical rule is
\(1-\theta_n\gtrsim\ell_nR_{0,n}\), together with compensation of
the resulting prefactor by the conditional gap and observable
normalization.

No Osterwalder--Schrader conclusion follows from these inequalities
alone.  Reflection positivity, existence of the limiting Schwinger
functions, and a spectral reconstruction remain separate inputs.
\end{assessment}

\begin{decision}[V86: the precise Osterwalder--Schrader handoff]
\label{dec:v16v85-next}
State and prove the positive-spectral-measure lemma that converts
exponential Euclidean-time decay into a lower spectral-support bound
in one observable channel.  Then formulate the additional density
and uniformity conditions needed to turn channel gaps into a
Hamiltonian mass gap, without identifying either conclusion with
the finite-cutoff Poincar\'e gap.
\end{decision}

\begin{firewall}[Claim boundary after V85]
\label{fire:v16v85}
The physical clustering theorem remains conditional on the
finite-cutoff law, the V83 local certificate, finite-range
influences, and the uniform rate and amplitude ledgers.  It neither
constructs a limiting SM--Einstein measure nor proves reflection
positivity or a Hamiltonian mass gap.
\end{firewall}

\section*{Version V85 append-only record}
\addcontentsline{toc}{section}{Version V85 append-only record}
The complete V84 source was retained byte for byte before this
continuation.  It had \(866750\) bytes, \(26201\) lines, and SHA--256
\[
 \texttt{733D9F12F65C0E001EAE1B5EEDBAB7F24848C435ACDC25F5C88C2AD793A31505}.
\]
V85 performed the required companion audit and established the exact
rate and amplitude corridors for converting block clustering to
physical clustering.


\section{V86: the Osterwalder--Schrader spectral handoff}
\label{sec:v16v86}

\subsection{A positive Laplace transform determines its lower support}

The step from Euclidean-time decay to an energy gap uses positivity
of a spectral measure.  It is not a consequence of an absolute
cross-covariance estimate by itself.

\begin{lemma}[Lower support of a positive Laplace transform]
\label{lem:v16v86-laplace}
Let \(\nu\) be a nonzero finite positive Borel measure on
\([0,\infty)\), put
\[
 F(t)=\int_{[0,\infty)}e^{-tE}\,d\nu(E),
 \qquad
 \beta=\inf\operatorname{supp}\nu.                  \tag{86.1}
\]
Then
\[
 \lim_{t\to\infty}-\frac1t\log F(t)=\beta.           \tag{86.2}
\]
For \(m\geq0\), the following are equivalent:
\[
\begin{split}
 &\operatorname{supp}\nu\subset[m,\infty),\\
 &F(t)\leq\nu([0,\infty))e^{-mt}
   \quad(t\geq0).                                   \tag{86.3}
\end{split}
\]
Moreover, if for some \(C<\infty\) and \(t_0\geq0\),
\[
 F(t)\leq Ce^{-mt}\quad(t\geq t_0),                 \tag{86.4}
\]
then \(\operatorname{supp}\nu\subset[m,\infty)\).
\end{lemma}

\begin{proof}
The support is contained in \([\beta,\infty)\), hence
\[
 F(t)\leq\nu([0,\infty))e^{-\beta t}.               \tag{86.5}
\]
For every \(\varepsilon>0\), the definition of the lower support
edge gives
\[
 q_\varepsilon:=
 \nu([\beta,\beta+\varepsilon])>0.
\]
Therefore
\[
 F(t)\geq q_\varepsilon
 e^{-(\beta+\varepsilon)t}.                         \tag{86.6}
\]
Taking \(-t^{-1}\log\), then the lower and upper limits, and finally
\(\varepsilon\downarrow0\), proves (86.2).

The first implication in (86.3) follows by integrating
\(e^{-tE}\leq e^{-mt}\).  Conversely, if
\(\nu([0,m-\varepsilon])>0\) for some \(\varepsilon>0\), then
\[
 F(t)\geq
 \nu([0,m-\varepsilon])e^{-(m-\varepsilon)t},
\]
which contradicts either the second line of (86.3) or (86.4) as
\(t\to\infty\).  Taking the union over positive rational
\(\varepsilon\) shows that \(\nu([0,m))=0\).
\end{proof}

\subsection{One observable channel after OS reconstruction}

Assume for this subsection that a continuum Schwinger functional has
already been constructed on a reflection-positive observable
algebra and that the Osterwalder--Schrader reconstruction gives
\[
 ({\cal H},\Omega,H),\qquad H\geq0,\qquad H\Omega=0. \tag{86.7}
\]
For a reconstructed observable \(F\), let
\[
 \psi_F=(F-\langle\Omega,F\Omega\rangle)\Omega.      \tag{86.8}
\]
Write \(P_H\) for the projection-valued spectral measure of \(H\) and
define
\[
 \nu_F(B)=\|P_H(B)\psi_F\|^2.                       \tag{86.9}
\]
Then the spectral theorem gives the positive autocorrelation
\[
 C_F(t)
 :=\langle\psi_F,e^{-tH}\psi_F\rangle
 =\int_{[0,\infty)}e^{-tE}\,d\nu_F(E).              \tag{86.10}
\]

\begin{theorem}[Euclidean decay gives an observable-channel gap]
\label{thm:v16v86-channel}
If, for some \(m>0\), \(C_F(t)\leq C_Fe^{-mt}\) for all sufficiently
large \(t\), then
\[
 P_H([0,m))\psi_F=0.                                \tag{86.11}
\]
Equivalently, the restriction of \(H\) to the cyclic spectral
subspace generated by \(\psi_F\) has no spectrum below \(m\).
\end{theorem}

\begin{proof}
Apply Lemma~\ref{lem:v16v86-laplace} to the positive measure
\(\nu_F\).  It gives \(\nu_F([0,m))=0\), which is precisely
\[
 \|P_H([0,m))\psi_F\|^2=0.
\]
Every vector in the cyclic spectral subspace is a limit of
\(\varphi(H)\psi_F\).  Spectral projections commute with
\(\varphi(H)\), so (86.11) holds throughout that subspace.
\end{proof}

This theorem must be applied to an autocorrelation, or to every
linear combination in a family.  A signed cross-correlation can
decay through cancellation while its two spectral channels still
contain low energy.

\subsection{Passage of the cutoff estimate to one channel}

Let \(F_n\) be normalized cutoff approximants to \(F\), and let
\(C_{F,n}(t)\) be their connected Euclidean-time two-point
functions in physical units.

\begin{proposition}[Uniform cutoff decay survives pointwise passage]
\label{prop:v16v86-limit-channel}
Suppose
\[
 \lim_{n\to\infty}C_{F,n}(t)=C_F(t)
 \quad\text{for every }t\geq t_0,                   \tag{86.12}
\]
and, with constants independent of \(n\),
\[
 |C_{F,n}(t)|\leq C_Fe^{-mt}
 \quad(n\geq1,\ t\geq t_0).                         \tag{86.13}
\]
Then
\[
 |C_F(t)|\leq C_Fe^{-mt}
 \quad(t\geq t_0).                                  \tag{86.14}
\]
If the limiting functional is reflection positive and (86.10)
identifies this limit with the reconstructed autocorrelation, then
(86.11) follows.
\end{proposition}

\begin{proof}
Take the limit \(n\to\infty\) in (86.13) for each fixed \(t\).
Reflection positivity and OS reconstruction supply the positive
spectral representation (86.10), after which
Theorem~\ref{thm:v16v86-channel} applies.
\end{proof}

The uniform rate and amplitude conditions (85.15) are exactly what
is needed to obtain a bound of the form (86.13) from V85, once the
physical supports of the two insertions are separated by Euclidean
time \(t\).  Pointwise convergence in (86.12) is still a separate
continuum-construction hypothesis.

\subsection{From channel gaps to a Hamiltonian gap}

Let
\[
 {\cal H}_0:=\Omega^\perp.                           \tag{86.15}
\]
This definition presumes that the vacuum sector to be removed is
the one-dimensional span of \(\Omega\).  If the zero-energy space is
larger, uniqueness must first be proved or the entire kernel of
\(H\) must be removed.

\begin{theorem}[Dense-channel criterion for a full energy gap]
\label{thm:v16v86-dense-gap}
Let \({\cal D}\subset{\cal H}_0\) be dense.  Suppose there is one
\(m>0\) such that
\[
 P_H([0,m))\psi=0
 \qquad(\psi\in{\cal D}).                            \tag{86.16}
\]
Then
\[
 \operatorname{spec}(H|_{{\cal H}_0})
 \subset[m,\infty).                                 \tag{86.17}
\]
\end{theorem}

\begin{proof}
The spectral projection \(Q=P_H([0,m))\) is bounded.  It vanishes on
the dense set \({\cal D}\), hence by continuity it vanishes on all of
\({\cal H}_0\).  The spectral theorem then gives (86.17).
\end{proof}

\begin{corollary}[Uniform local-observable handoff]
\label{cor:v16v86-local-gap}
Suppose the centered vectors generated from a gauge-invariant local
observable algebra form a dense linear subspace
\({\cal D}\subset{\cal H}_0\).  If every
\(\psi\in{\cal D}\) has
\[
 \langle\psi,e^{-tH}\psi\rangle
 \leq C_\psi e^{-mt}                                \tag{86.18}
\]
for all sufficiently large \(t\), with the same \(m>0\) but an
allowed vector-dependent \(C_\psi\), then \(H|_{{\cal H}_0}\) has
gap at least \(m\).
\end{corollary}

\begin{proof}
Theorem~\ref{thm:v16v86-channel} gives (86.16) for each
\(\psi\in{\cal D}\).  Apply
Theorem~\ref{thm:v16v86-dense-gap}.
\end{proof}

The common rate is essential.  Positive rates \(m_\psi\) that tend
to zero along a dense family do not exclude spectrum accumulating at
zero.  Density is also essential: decay in one Wilson-loop, matter,
or curvature channel proves only the corresponding channel gap.

\subsection{The two generators must not be identified}

At finite cutoff, V66--V85 use the reversible diffusion generator
\[
 L_n=-\Delta+\nabla U_n\cdot\nabla
 \quad\text{on }L^2(\mu_n).                          \tag{86.19}
\]
After OS reconstruction, \(H\) is the generator of physical
Euclidean-time translations on the reconstructed Hilbert space.
The Poincar\'e estimate controls \(L_n^{-1}\) on centered
configuration functions.  The mass-gap statement controls the
spectral projection of \(H\) on physical states.  They are different
operators on different Hilbert spaces.

\begin{proposition}[Required bridge between the two gaps]
\label{prop:v16v86-two-generators}
A lower bound on the spectrum of \(L_n\) implies no bound on
\(\operatorname{spec}(H)\) from the results proved so far.  The
valid bridge in the present programme is the chain
\[
\begin{split}
\text{conditional gaps for }L_n
&\Longrightarrow \text{localized covariance at cutoff }n\\
&\Longrightarrow \text{uniform physical Euclidean-time decay}\\
&\Longrightarrow \text{positive OS spectral support}.
                                                               \tag{86.20}
\end{split}
\]
Each arrow uses, respectively, V84, V85 plus continuum convergence,
and Theorem~\ref{thm:v16v86-channel} plus reflection positivity.
\end{proposition}

\begin{proof}
The first statement is a claim-boundary assertion: no unitary
intertwiner or spectral identity between \(L_n\) and \(H\) has been
assumed or proved.  The three implications in (86.20) are exactly
Theorem~\ref{thm:v16v84-covariance},
Theorem~\ref{thm:v16v85-class} with
Proposition~\ref{prop:v16v86-limit-channel}, and
Theorem~\ref{thm:v16v86-channel}.
\end{proof}

\subsection{OS handoff ledger for the SM--Einstein programme}

The spectral conclusion requires the following entries in addition
to the V83 finite-cutoff certificate:
\begin{enumerate}
\item a common continuum gauge-invariant observable algebra and
      normalized cutoff approximants;
\item convergence of its Schwinger functions, including the
      connected time-separated two-point functions;
\item reflection positivity of the limit and the regularity needed
      for OS reconstruction;
\item the uniform rate and amplitude bounds of V85 on a
      linear family whose reconstructed centered vectors are dense;
\item vacuum uniqueness, or a precise replacement of
      \(\Omega^\perp\) by the orthogonal complement of
      \(\ker H\).
\end{enumerate}
For the gravitational sector, the third entry is a major open
bottleneck: gauge fixing, ghosts, and the conformal direction cannot
be discarded before positivity is proved on the physical observable
algebra.

\begin{assessment}[What V86 closes]
\label{ass:v16v86-status}
V86 proves the exact positive-measure theorem behind the OS handoff.
Uniform Euclidean-time decay gives a gap in one reconstructed
observable channel.  A common rate on a dense linear family gives a
full Hamiltonian energy gap, provided the limiting
reflection-positive theory and unique vacuum already exist.

No such limiting SM--Einstein Schwinger functional or reflection
positivity theorem has yet been constructed.  The finite-cutoff
Poincar\'e gap is not itself the Hamiltonian gap.
\end{assessment}

\begin{decision}[V87: reflection positivity under cutoff passage]
\label{dec:v16v86-next}
Go upstream to the missing continuum handoff.  Prove a closure
theorem showing that reflection positivity passes to a limiting
Schwinger functional under convergence on the positive-time
observable algebra.  Separate this algebraic closure from tightness,
existence, Euclidean invariance, and the gravitational
gauge-quotient problem.
\end{decision}

\begin{firewall}[Claim boundary after V86]
\label{fire:v16v86}
V86 is a conditional spectral theorem.  It proves neither the
existence nor reflection positivity of the SM--Einstein continuum
functional, does not establish density of the required physical
observable channels, and does not claim a full Standard
Model--Einstein mass gap.
\end{firewall}

\section*{Version V86 append-only record}
\addcontentsline{toc}{section}{Version V86 append-only record}
The complete V85 source was retained byte for byte before this
continuation.  It had \(877730\) bytes, \(26521\) lines, and SHA--256
\[
 \texttt{1A3EBE7AF1670E4FB69AA4E6B88675D15733EEF71E3EEAE7C873F149AFC240C3}.
\]
V86 proved the positive Laplace-transform support lemma, the
observable-channel OS gap theorem, and the dense-channel criterion
for a full Hamiltonian energy gap.


\section{V87: closure of reflection positivity under cutoff passage}
\label{sec:v16v87}

\subsection{A common reflected observable algebra}

Let \({\cal A}\) be a unital complex algebra of continuum cylinder
observables and let \({\cal A}_+\subset{\cal A}\) be the subalgebra
supported at nonnegative Euclidean time.  Let
\(\Theta:{\cal A}\to{\cal A}\) be the antilinear time-reflection
operation, including complex conjugation, so that
\[
 \Theta(cF)=\overline c\,\Theta F.                   \tag{87.1}
\]
At cutoff \(n\), use corresponding objects
\(({\cal A}_n,{\cal A}_{n,+},\Theta_n)\) and a linear functional
\({\cal S}_n\).  The comparison must be made through maps
\[
 \iota_n:{\cal A}_+\longrightarrow{\cal A}_{n,+}    \tag{87.2}
\]
that preserve products and the reflected products needed below.
This prevents a change of observable normalization from being hidden
inside the limiting argument.

\begin{definition}[Pulled-back OS form]
\label{def:v16v87-os-form}
For \(F,G\in{\cal A}_+\), define
\[
 Q_n(F,G)
 :=
 {\cal S}_n\bigl((\Theta_n\iota_nF)(\iota_nG)\bigr).
                                                               \tag{87.3}
\]
The cutoff functional is reflection positive on the embedded
algebra when every finite matrix
\[
 \bigl(Q_n(F_\alpha,F_\beta)\bigr)_{\alpha,\beta=1}^N
 \quad\text{is positive semidefinite}.              \tag{87.4}
\]
\end{definition}

\subsection{Algebraic closure}

\begin{theorem}[Reflection positivity is closed under pointwise
form convergence]
\label{thm:v16v87-rp-closure}
Suppose (87.4) holds for every \(n\), and suppose that for every
\(F,G\in{\cal A}_+\) the limit
\[
 Q(F,G):=\lim_{n\to\infty}Q_n(F,G)                  \tag{87.5}
\]
exists.  Then \(Q\) is positive semidefinite:
\[
 \sum_{\alpha,\beta=1}^N
 \overline{c_\alpha}c_\beta
 Q(F_\alpha,F_\beta)\geq0                           \tag{87.6}
\]
for all finite families \(F_\alpha\in{\cal A}_+\) and coefficients
\(c_\alpha\in\mathbb C\).  If a limiting functional
\({\cal S}\) satisfies
\[
 Q(F,G)={\cal S}((\Theta F)G),                      \tag{87.7}
\]
then \({\cal S}\) is reflection positive on \({\cal A}_+\).
\end{theorem}

\begin{proof}
For fixed \(N\), \(F_\alpha\), and \(c_\alpha\), cutoff reflection
positivity gives
\[
 q_n:=
 \sum_{\alpha,\beta=1}^N
 \overline{c_\alpha}c_\beta
 Q_n(F_\alpha,F_\beta)\geq0.                        \tag{87.8}
\]
The sum is finite, so (87.5) implies
\[
 q_n\longrightarrow
 \sum_{\alpha,\beta=1}^N
 \overline{c_\alpha}c_\beta
 Q(F_\alpha,F_\beta).
\]
A limit of nonnegative real numbers is nonnegative, proving
(87.6).  Equation (87.7) is exactly the OS reflection-positivity
inequality for \({\cal S}\).
\end{proof}

Entrywise convergence of every finite reflected Gram matrix is
therefore enough.  Uniform positivity floors are not required for
this closure step; zero eigenvalues may appear in the limit and are
removed by the OS null-space quotient.

\begin{corollary}[Closure of the null-space quotient form]
\label{cor:v16v87-null}
Under Theorem~\ref{thm:v16v87-rp-closure}, the set
\[
 {\cal N}:=\{F\in{\cal A}_+:Q(F,F)=0\}              \tag{87.9}
\]
is a linear subspace, and
\[
 \langle[F],[G]\rangle_{\rm OS}:=Q(F,G)             \tag{87.10}
\]
is a well-defined inner product on
\({\cal A}_+/{\cal N}\).
\end{corollary}

\begin{proof}
Positivity of \(Q\) implies the Cauchy--Schwarz inequality
\[
 |Q(F,G)|^2\leq Q(F,F)Q(G,G).                       \tag{87.11}
\]
Indeed, positivity of \(Q(F+zG,F+zG)\) for all \(z\in\mathbb C\)
gives the usual quadratic discriminant argument.  Hence
\(Q(F,F)=0\) implies \(Q(F,G)=0\) for every \(G\).
This proves linearity of \({\cal N}\), independence of
representatives in (87.10), and positive definiteness on the
quotient.
\end{proof}

\subsection{Weak convergence of probability laws}

The algebraic theorem does not itself produce the limit (87.5).
For honest positive cutoff measures, weak convergence supplies it
for bounded continuous cylinder observables.

\begin{proposition}[Weak-limit reflection positivity]
\label{prop:v16v87-weak-rp}
Let \(X\) be a Hausdorff space of generalized fields with a
continuous reflection \(\vartheta:X\to X\).  Let
\(\mu_n\) be reflection-positive probability measures on \(X\) and
suppose
\[
 \mu_n\Rightarrow\mu                              \tag{87.12}
\]
weakly.  Let \({\cal A}_+\) consist of bounded continuous cylinder
functions supported at nonnegative time, with
\[
 (\Theta F)(\phi)=\overline{F(\vartheta\phi)}.       \tag{87.13}
\]
Then \(\mu\) is reflection positive on \({\cal A}_+\).
\end{proposition}

\begin{proof}
For a finite family \(F_\alpha\in{\cal A}_+\) and coefficients
\(c_\alpha\), set \(F=\sum_\alpha c_\alpha F_\alpha\).  The function
\[
 \phi\longmapsto
 \overline{F(\vartheta\phi)}F(\phi)                 \tag{87.14}
\]
is bounded and continuous.  Weak convergence gives
\[
\begin{split}
 \int\overline{F(\vartheta\phi)}F(\phi)\,d\mu(\phi)
 &=
 \lim_{n\to\infty}
 \int\overline{F(\vartheta\phi)}F(\phi)\,d\mu_n(\phi)
 \geq0.                                             \tag{87.15}
\end{split}
\]
Expanding the left side proves all finite OS matrix inequalities.
\end{proof}

\begin{proposition}[Extension to unbounded cylinder polynomials]
\label{prop:v16v87-ui}
In the setting of
Proposition~\ref{prop:v16v87-weak-rp}, let \(F\) be a continuous
cylinder polynomial for which the random variables
\[
 Y_n(\phi)=\overline{F(\vartheta\phi)}F(\phi)        \tag{87.16}
\]
are uniformly integrable with respect to \(\mu_n\), and suppose
\(Y(\phi)=\overline{F(\vartheta\phi)}F(\phi)\) is
\(\mu\)-integrable.  Then
\[
 \int Y\,d\mu=\lim_{n\to\infty}\int Y_n\,d\mu_n
 \geq0.                                             \tag{87.17}
\]
The same assertion holds for every finite linear combination of the
chosen polynomials, and hence reflection positivity extends to their
linear span.
\end{proposition}

\begin{proof}
For \(M>0\), choose a bounded continuous truncation
\(\chi_M:\mathbb C\to\mathbb C\) that equals the identity on
\(\{|z|\leq M\}\) and has modulus at most \(2M\).  Weak convergence
passes the expectation of \(\chi_M(Y_n)\) to that of
\(\chi_M(Y)\).  Uniform integrability makes
\[
 \sup_n\int|Y_n-\chi_M(Y_n)|\,d\mu_n\longrightarrow0
 \quad(M\to\infty).                                 \tag{87.18}
\]
The same truncation argument, together with Fatou applied to tail
indicators and the assumed integrability of \(Y\), passes the
remaining limit.  Thus the expectations converge.  Each cutoff
expectation is nonnegative by reflection positivity, so the limit is
nonnegative.  Apply this argument to each finite linear combination
before expanding its quadratic form.
\end{proof}

\subsection{Other OS axioms closed by the same mechanism}

\begin{proposition}[Normalization, reality, and symmetry]
\label{prop:v16v87-other-closures}
Suppose \({\cal S}_n(A)\to{\cal S}(A)\) on a common algebra.
Then each of the following properties passes to the limit whenever
it is expressed on that algebra:
\[
\begin{array}{ll}
\text{normalization:}&{\cal S}_n(1)=1,\\
\text{reality:}&{\cal S}_n(A^*)=\overline{{\cal S}_n(A)},\\
\text{group invariance:}&
{\cal S}_n(\tau_gA)={\cal S}_n(A).
\end{array}                                         \tag{87.19}
\]
\end{proposition}

\begin{proof}
Take the scalar limit of the corresponding equality.  For group
invariance, both \(A\) and \(\tau_gA\) must belong to the common
convergence algebra.
\end{proof}

This proposition does not create full Euclidean invariance from a
lattice subgroup.  One must also prove convergence of approximate
rotations/translations and continuity of the limiting group action.
Likewise, reflection positivity does not imply OS regularity or
temperedness.

\subsection{Gauge-invariant and gravitational observables}

\begin{theorem}[Physical-subalgebra version]
\label{thm:v16v87-physical-algebra}
Let \({\cal A}^{\rm phys}_+\) be a common algebra of
gauge-invariant positive-time observables.  If its cutoff embeddings
\(\iota_n\) intertwine reflection and every cutoff pulled-back OS
form (87.3) is positive, then pointwise convergence of these forms
produces a positive limiting OS form on
\({\cal A}^{\rm phys}_+\), irrespective of whether an auxiliary
gauge-fixed representation is positive on a larger algebra.
\end{theorem}

\begin{proof}
Apply Theorem~\ref{thm:v16v87-rp-closure} after restricting every
form to \({\cal A}^{\rm phys}_+\).  The proof uses only the restricted
quadratic forms and therefore makes no assertion about auxiliary
variables outside that algebra.
\end{proof}

The theorem identifies a viable target but does not verify its
hypotheses.  In the SM--Einstein programme, a common physical algebra
must handle internal gauge symmetry, diffeomorphism symmetry, moving
reflection hypersurfaces, and metric-dependent localization.
Ordinary gauge-fixed ghost expressions do not automatically define
positive probability measures, and positivity of their physical
restriction remains to be proved.

\subsection{Separation from tightness and existence}

\begin{proposition}[Reflection positivity is not a compactness
theorem]
\label{prop:v16v87-rp-no-tightness}
The reflection-positivity inequalities alone do not imply tightness
of a family of cutoff laws.
\end{proposition}

\begin{proof}
Let \(\mu_n=\delta_{\phi_n}\), where \(\phi_n\) is a
reflection-invariant constant field of amplitude \(n\).
For every positive-time observable \(F\),
\[
 \int\overline{F(\vartheta\phi)}F(\phi)\,d\mu_n(\phi)
 =|F(\phi_n)|^2\geq0.                               \tag{87.20}
\]
Thus every \(\mu_n\) is reflection positive.  On any field topology
in which the constant amplitudes \(n\) escape every compact set, the
family \(\{\delta_{\phi_n}\}\) is not tight.
\end{proof}

Consequently the order of proof must be:
\[
\text{uniform moment/regularity bounds}
\Longrightarrow\text{tightness and subsequential existence}
\Longrightarrow\text{convergence of reflected forms}
\Longrightarrow\text{limiting reflection positivity}.          \tag{87.21}
\]

\begin{assessment}[What V87 closes]
\label{ass:v16v87-status}
V87 proves that reflection positivity is closed under pointwise
convergence of reflected Gram forms.  Weak convergence of positive
measures supplies this convergence for bounded continuous cylinder
observables, and uniform integrability extends it to cylinder
polynomials.  The result remains valid on a common gauge-invariant
physical subalgebra.

The theorem is a closure mechanism, not an existence theorem.
Tightness, convergence, approximate-to-continuous Euclidean
invariance, and construction of the gravitational physical algebra
remain open.
\end{assessment}

\begin{decision}[V88: a concrete tightness criterion]
\label{dec:v16v87-next}
Go one step farther upstream.  Place the cutoff bosonic fields in a
fixed negative Sobolev space on a compact physical volume and prove
a tightness criterion from a uniform stronger negative-Sobolev
moment bound and a compact embedding.  State separately what is
needed for fermionic and ghost sectors, which are not probability
measures on ordinary distribution spaces.
\end{decision}

\begin{firewall}[Claim boundary after V87]
\label{fire:v16v87}
No cutoff SM--Einstein law has been shown reflection positive, tight,
or convergent.  V87 only proves preservation of reflection
positivity after the required common-algebra convergence is
available.
\end{firewall}

\section*{Version V87 append-only record}
\addcontentsline{toc}{section}{Version V87 append-only record}
The complete V86 source was retained byte for byte before this
continuation.  It had \(889227\) bytes, \(26847\) lines, and SHA--256
\[
 \texttt{D17A71BCF424DAFFC9A734E3713C94DD747267E0C70BFE9E647E060CCA17F11D}.
\]
V87 proved closure of reflection positivity on a common continuum
observable algebra and separated that closure from tightness and
existence.


\section{V88: negative-Sobolev tightness for the bosonic cutoff laws}
\label{sec:v16v88}

\subsection{A fixed distribution space}

Let \(M\) be a compact \(d\)-dimensional Riemannian manifold used as
the physical finite volume, and fix a positive elliptic reference
operator
\[
 \Lambda=(1-\Delta_{\bar g})^{1/2}.                  \tag{88.1}
\]
For \(r\in\mathbb R\), define
\[
 \|u\|_{H^r(M)}^2
 :=\|\Lambda^ru\|_{L^2(M)}^2.                       \tag{88.2}
\]
All cutoff interpolation maps must take their fields into these
same reference spaces.  Changing the norm with \(n\) would remove
the compactness content of the estimates below.

\begin{lemma}[Compact negative-Sobolev embedding]
\label{lem:v16v88-compact}
If \(0\leq t<s\), the inclusion
\[
 H^{-t}(M)\hookrightarrow H^{-s}(M)                 \tag{88.3}
\]
is compact.
\end{lemma}

\begin{proof}
Let \(e_k\) be an orthonormal eigenbasis of
\(-\Delta_{\bar g}\), with eigenvalues
\(0\leq\lambda_0\leq\lambda_1\leq\cdots\).
If \(\|u\|_{H^{-t}}\leq1\), then for the projection
\(P_Nu=\sum_{k\leq N}\langle u,e_k\rangle e_k\),
\[
\begin{split}
 \|(I-P_N)u\|_{H^{-s}}^2
 &=
 \sum_{k>N}(1+\lambda_k)^{-s}
 |\langle u,e_k\rangle|^2\\
 &\leq
 (1+\lambda_{N+1})^{-(s-t)}
 \|u\|_{H^{-t}}^2.                                  \tag{88.4}
\end{split}
\]
The right side tends to zero uniformly on the unit ball, whereas
each \(P_N\) has finite rank.  Thus the inclusion is the operator-norm
limit of finite-rank maps and is compact.
\end{proof}

\subsection{The moment-to-tightness theorem}

Let \(J_n\) denote the chosen interpolation of a cutoff bosonic field
into distributions on \(M\), and let
\[
 \widehat\mu_n=(J_n)_\#\mu_n                         \tag{88.5}
\]
be its pushforward law.

\begin{theorem}[Uniform stronger-norm moment implies tightness]
\label{thm:v16v88-tightness}
Fix \(p>0\) and \(0\leq t<s\).  Suppose
\[
 \sup_n
 \int\|u\|_{H^{-t}}^p\,d\widehat\mu_n(u)
 \leq C_p<\infty.                                   \tag{88.6}
\]
Then \(\{\widehat\mu_n\}\) is tight as a family of probability
measures on \(H^{-s}(M)\).  Consequently every sequence has a weakly
convergent subsequence in \(H^{-s}(M)\).
\end{theorem}

\begin{proof}
Let \(B_R\) be the closed radius-\(R\) ball of \(H^{-t}(M)\), and let
\(K_R\) be its closure in \(H^{-s}(M)\).  By
Lemma~\ref{lem:v16v88-compact}, \(K_R\) is compact in \(H^{-s}(M)\).
Markov's inequality gives
\[
 \widehat\mu_n(H^{-s}\setminus K_R)
 \leq
 \widehat\mu_n(\|u\|_{H^{-t}}>R)
 \leq C_pR^{-p}.                                    \tag{88.7}
\]
Given \(\varepsilon>0\), choose
\(R>(C_p/\varepsilon)^{1/p}\).  This is tightness.
The target \(H^{-s}(M)\) is a separable Hilbert space and hence
Polish; relative sequential compactness of tight probability
families follows from Prokhorov's theorem.
\end{proof}

\begin{corollary}[Several bosonic fields]
\label{cor:v16v88-product}
Let \(u_n^{(a)}\), \(1\leq a\leq q\), take values in
\(H^{-s_a}(M;E_a)\), where each \(E_a\) is a fixed
finite-rank tensor or internal bundle.  If for some
\(t_a<s_a\) and \(p_a>0\),
\[
 \sup_n
 {\mathbb E}_{\mu_n}
 \|u_n^{(a)}\|_{H^{-t_a}}^{p_a}<\infty
 \quad(1\leq a\leq q),                              \tag{88.8}
\]
then their joint laws are tight in
\[
 \prod_{a=1}^qH^{-s_a}(M;E_a).                      \tag{88.9}
\]
\end{corollary}

\begin{proof}
For each \(a\), choose a compact set \(K_a\) whose complement has
probability at most \(\varepsilon/q\), uniformly in \(n\), by
Theorem~\ref{thm:v16v88-tightness}.  The finite product
\(\prod_aK_a\) is compact, and the union bound makes the probability
of its complement at most \(\varepsilon\).
\end{proof}

\subsection{A spectral sufficient condition}

The moment hypothesis (88.6) can be checked mode by mode.  Write
\[
 u_{n,k}=\langle J_nu_n,e_k\rangle.                 \tag{88.10}
\]

\begin{proposition}[Polynomial mode growth]
\label{prop:v16v88-mode}
Suppose there are constants \(C,q\), independent of \(n,k\), such
that every retained cutoff mode obeys
\[
 {\mathbb E}_{\mu_n}|u_{n,k}|^2
 \leq C(1+\lambda_k)^q,                             \tag{88.11}
\]
and omitted modes have coefficient zero.  If
\[
 t>q+\frac d2,                                      \tag{88.12}
\]
then
\[
 \sup_n{\mathbb E}_{\mu_n}
 \|J_nu_n\|_{H^{-t}}^2<\infty.                      \tag{88.13}
\]
Hence the laws are tight in \(H^{-s}(M)\) for every \(s>t\).
\end{proposition}

\begin{proof}
Tonelli's theorem and (88.11) give
\[
\begin{split}
 {\mathbb E}\|J_nu_n\|_{H^{-t}}^2
 &=
 \sum_k(1+\lambda_k)^{-t}
 {\mathbb E}|u_{n,k}|^2\\
 &\leq
 C\sum_k(1+\lambda_k)^{-(t-q)}.                     \tag{88.14}
\end{split}
\]
Weyl's eigenvalue bound implies
\(\lambda_k\geq c k^{2/d}\) for all sufficiently large \(k\).
The final series is therefore bounded by a constant times
\[
 1+\sum_{k\geq1}k^{-2(t-q)/d},
\]
which converges exactly under (88.12).  Apply
Theorem~\ref{thm:v16v88-tightness} with \(p=2\).
\end{proof}

The exponent \(q\) must be computed after field-strength
renormalization.  A bound on unrenormalized lattice coefficients does
not by itself control the continuum interpolation.

\subsection{Cylinder convergence and moments}

\begin{proposition}[Consequences of an \(H^{-s}\) weak limit]
\label{prop:v16v88-cylinder}
Suppose \(\widehat\mu_{n_j}\Rightarrow\mu\) in \(H^{-s}(M)\).  For
test functions \(f_1,\ldots,f_r\in H^s(M)\) and every bounded
continuous \(B:\mathbb R^r\to\mathbb C\),
\[
\begin{split}
 &\lim_{j\to\infty}
 \int B(\langle u,f_1\rangle,\ldots,\langle u,f_r\rangle)
 \,d\widehat\mu_{n_j}(u)\\
 &\qquad=
 \int B(\langle u,f_1\rangle,\ldots,\langle u,f_r\rangle)
 \,d\mu(u).                                         \tag{88.15}
\end{split}
\]
If a cylinder polynomial \(P\) of degree \(r_0\) has a uniformly
integrable family
\[
 P(\langle u,f_1\rangle,\ldots,\langle u,f_r\rangle)
 \quad\text{under }\widehat\mu_{n_j},                \tag{88.16}
\]
then (88.15) also holds with \(B=P\).
\end{proposition}

\begin{proof}
Each dual pairing \(u\mapsto\langle u,f_a\rangle\) is continuous on
\(H^{-s}\), so the first claim is the definition of weak
convergence.  The second follows by bounded continuous truncation
and uniform integrability, as in
Proposition~\ref{prop:v16v87-ui}.
\end{proof}

This is the convergence input required by V87 for reflected cylinder
forms.  For polynomial Schwinger functions, tightness alone is
insufficient; the uniform-integrability ledger must be retained.

\subsection{Local-volume version}

Let \(M_1\Subset M_2\Subset\cdots\) exhaust a noncompact spacetime.
For compatible restrictions, equip
\[
 H^{-s}_{\rm loc}:=
 \{u:u|_{M_j}\in H^{-s}(M_j)\text{ for every }j\}    \tag{88.17}
\]
with its projective-limit topology.

\begin{proposition}[Local tightness by diagonal compact sets]
\label{prop:v16v88-local}
If for every \(j\) there are \(t_j<s\), \(p_j>0\), and \(C_j<\infty\)
such that
\[
 \sup_n{\mathbb E}
 \|J_nu_n\|_{H^{-t_j}(M_j)}^{p_j}\leq C_j,           \tag{88.18}
\]
then the laws are tight in \(H^{-s}_{\rm loc}\).
\end{proposition}

\begin{proof}
Given \(\varepsilon>0\), apply
Theorem~\ref{thm:v16v88-tightness} on \(M_j\) to choose compact
\(K_j\subset H^{-s}(M_j)\) with
\[
 \sup_n{\mathbb P}(u_n|_{M_j}\notin K_j)
 \leq\varepsilon2^{-j}.                             \tag{88.19}
\]
The set of compatible distributions whose \(j\)-th restriction lies
in \(K_j\) is a closed subset of the compact product
\(\prod_jK_j\), hence compact in the projective-limit topology.
The union bound makes its complement have probability at most
\(\varepsilon\).
\end{proof}

\subsection{What the theorem does not cover}

\paragraph{Connections and metrics.}
The linear-space theorem applies to gauge potentials or metric
perturbations only after choosing compatible interpolation maps into
fixed bundles and controlling the gauge/coordinate representatives.
Weak convergence in \(H^{-s}\) does not preserve pointwise
nondegeneracy or a fixed metric signature.  One must either obtain
stronger geometric control or construct the limit directly through
gauge- and diffeomorphism-invariant observables.

\paragraph{Fermions and ghosts.}
Berezin integrals are not probability measures on an ordinary
distribution space.  Theorem~\ref{thm:v16v88-tightness} therefore
applies only to a positive bosonic marginal.  Fermionic Schwinger
functions require uniform bounds on the coefficient kernels of the
Grassmann generating functional and graded reflection positivity.
Faddeev--Popov ghosts are auxiliary graded variables and cannot be
declared tight by (88.6).  If they are integrated out, positivity and
uniform integrability of the resulting physical bosonic
observables must be proved.

\subsection{Insertion into the finite-cutoff certificate}

For each renormalized bosonic field species \(a\), append the entry
\[
 {\mathfrak T}_{a,n}
 :=
 {\mathbb E}_{\mu_n}
 \|J_nu_n^{(a)}\|_{H^{-t_a}}^{p_a}.                 \tag{88.20}
\]
The continuum compactness gate is
\[
 \sup_n{\mathfrak T}_{a,n}<\infty,
 \qquad t_a<s_a,                                    \tag{88.21}
\]
together with uniform integrability for every unbounded cylinder
observable used in the limiting Schwinger functions.

\begin{assessment}[What V88 closes]
\label{ass:v16v88-status}
V88 gives a complete compactness mechanism for a finite family of
renormalized bosonic cutoff fields: a uniform moment in a stronger
negative Sobolev norm gives tightness in a weaker one.  A polynomial
mode-variance bound provides a concrete sufficient condition, and a
local version handles volume exhaustion.

No estimate in the current SM--Einstein certificate yet supplies the
uniform quantities (88.20).  Fermionic and ghost sectors require a
separate graded-functional treatment, and distributional metric
limits require additional geometric control.
\end{assessment}

\begin{decision}[V89: obtain moments from conditional coercivity]
\label{dec:v16v88-next}
Return to the finite-cutoff law and identify a noncircular route from
the named block coercivity margins to (88.20).  Prove a block
sub-Gaussian moment estimate under a uniform conditional
log-Sobolev inequality and a controlled mean, then sum it in the
negative-Sobolev norm.  Record explicitly why a Poincar\'e inequality
alone does not provide the same exponential concentration.
\end{decision}

\begin{firewall}[Claim boundary after V88]
\label{fire:v16v88}
V88 proves a conditional tightness theorem.  It does not establish
the required uniform Sobolev moments for the SM--Einstein fields,
construct a continuum law, or solve the fermionic, ghost, gauge, or
metric-signature problems.
\end{firewall}

\section*{Version V88 append-only record}
\addcontentsline{toc}{section}{Version V88 append-only record}
The complete V87 source was retained byte for byte before this
continuation.  It had \(901076\) bytes, \(27178\) lines, and SHA--256
\[
 \texttt{72F0293C17B0B4457CCF9ED6DF48CC49F09F1AC0488BA1E28D90637EC4A4162F}.
\]
V88 proved the negative-Sobolev tightness criterion and isolated the
uniform bosonic moment estimates required from the cutoff action.


\section{V89: log-Sobolev concentration and Sobolev moments}
\label{sec:v16v89}

\subsection{Conditional entropy and the missing mixing input}

For a nonnegative integrable random variable \(Y\), write
\[
 \operatorname{Ent}_\mu(Y)
 :=
 \int Y\log\frac{Y}{\int Y\,d\mu}\,d\mu.             \tag{89.1}
\]
Let \(\mu_i(\,\cdot\,|x_{\widehat i})\) denote the conditional law of
block \(i\).

\begin{definition}[Conditional LSI and entropy tensorization]
\label{def:v16v89-conditional}
The conditional log-Sobolev floors are numbers
\(\rho_i^{\rm LS}>0\) for which
\[
 \operatorname{Ent}_{\mu_i}(F^2)
 \leq
 \frac{2}{\rho_i^{\rm LS}}
 \int|\nabla_iF|^2\,d\mu_i                          \tag{89.2}
\]
for every exterior configuration and every smooth \(F\).
An approximate tensorization constant \(C_{\rm AT}\) satisfies
\[
 \operatorname{Ent}_\mu(F^2)
 \leq
 C_{\rm AT}\sum_i
 \int\operatorname{Ent}_{\mu_i(\cdot|x_{\widehat i})}(F^2)
 \,d\mu(x_{\widehat i}).                            \tag{89.3}
\]
\end{definition}

Conditional inequalities alone do not make (89.3) automatic.
Equation (89.3) is the entropy-level mixing input that must be
proved from the interaction or entered as a separate certificate.

\begin{theorem}[Conditional-to-global log-Sobolev inequality]
\label{thm:v16v89-global-lsi}
If (89.2)--(89.3) hold and
\[
 \rho_{\min}^{\rm LS}:=\inf_i\rho_i^{\rm LS}>0,
\]
then
\[
 \operatorname{Ent}_\mu(F^2)
 \leq
 \frac{2}{\alpha}
 \int|\nabla F|^2\,d\mu,
 \qquad
 \alpha:=\frac{\rho_{\min}^{\rm LS}}{C_{\rm AT}}.   \tag{89.4}
\]
\end{theorem}

\begin{proof}
Insert (89.2) into (89.3), integrate the exterior variables, and
sum:
\[
\begin{split}
 \operatorname{Ent}_\mu(F^2)
 &\leq
 2C_{\rm AT}\sum_i\frac1{\rho_i^{\rm LS}}
 \int|\nabla_iF|^2\,d\mu\\
 &\leq
 \frac{2C_{\rm AT}}{\rho_{\min}^{\rm LS}}
 \int|\nabla F|^2\,d\mu.
                                                               \tag{89.5}
\end{split}
\]
This is (89.4).
\end{proof}

The scalar influence condition of V84 is a plausible input to an
independent proof of (89.3), but V84 proves a covariance resolvent,
not entropy tensorization.  We therefore do not reuse
\(\theta<1\) as if (89.3) had already been obtained.

\subsection{Herbst concentration with constants}

\begin{theorem}[Sub-Gaussian Laplace bound]
\label{thm:v16v89-herbst}
Assume the global log-Sobolev inequality (89.4).  If a locally
Lipschitz \(X\) obeys
\[
 |\nabla X|\leq L\quad\mu\text{-almost everywhere}, \tag{89.6}
\]
then for every \(\lambda\in\mathbb R\),
\[
 \log{\mathbb E}_\mu
 e^{\lambda(X-{\mathbb E}_\mu X)}
 \leq\frac{\lambda^2L^2}{2\alpha}.                  \tag{89.7}
\]
Consequently,
\[
 \mu(|X-{\mathbb E}X|\geq r)
 \leq2\exp\left(-\frac{\alpha r^2}{2L^2}\right).
                                                               \tag{89.8}
\]
\end{theorem}

\begin{proof}
First take \(\lambda>0\) and apply (89.4) to
\(F=e^{\lambda X/2}\).  With
\(\psi(\lambda)=\log{\mathbb E}e^{\lambda X}\), one obtains
\[
 \lambda\psi'(\lambda)-\psi(\lambda)
 \leq\frac{\lambda^2L^2}{2\alpha}.                  \tag{89.9}
\]
Since the left side divided by \(\lambda^2\) is
\(\frac{d}{d\lambda}(\psi(\lambda)/\lambda)\),
integration from \(0\) to \(\lambda\), using
\(\lim_{\lambda\to0}\psi(\lambda)/\lambda={\mathbb E}X\), gives
(89.7).  Apply the same argument to \(-X\) for negative
\(\lambda\).  Chernoff's bound and minimization at
\(\lambda=\alpha r/L^2\) give each one-sided tail; their sum is
(89.8).
\end{proof}

\begin{corollary}[All centered moments]
\label{cor:v16v89-moments}
For every \(p>0\),
\[
 {\mathbb E}|X-{\mathbb E}X|^p
 \leq
 2\Gamma\left(\frac p2+1\right)
 \left(\frac{2L^2}{\alpha}\right)^{p/2}.             \tag{89.10}
\]
\end{corollary}

\begin{proof}
Use the layer-cake formula and (89.8):
\[
\begin{split}
 {\mathbb E}|X-{\mathbb E}X|^p
 &=
 p\int_0^\infty r^{p-1}
 \mu(|X-{\mathbb E}X|\geq r)\,dr\\
 &\leq
 2p\int_0^\infty r^{p-1}
 e^{-\alpha r^2/(2L^2)}\,dr,
                                                               \tag{89.11}
\end{split}
\]
which evaluates to (89.10).
\end{proof}

\subsection{Application to the interpolation norm}

At cutoff \(n\), give the finite-dimensional configuration space its
block Euclidean norm and let
\[
 J_n:\mathbb R^{N_n}\longrightarrow H^{-t}(M;E)     \tag{89.12}
\]
be the linear interpolation and field-renormalization map.  Define
\[
 \sigma_{t,n}:=
 \|J_n\|_{\mathbb R^{N_n}\to H^{-t}}.                \tag{89.13}
\]

\begin{lemma}[Lipschitz constant of the Sobolev norm]
\label{lem:v16v89-norm-lipschitz}
The function
\[
 X_n(x)=\|J_nx\|_{H^{-t}}                           \tag{89.14}
\]
is \(\sigma_{t,n}\)-Lipschitz.
\end{lemma}

\begin{proof}
The reverse triangle inequality and (89.13) give
\[
 |X_n(x)-X_n(y)|
 \leq\|J_n(x-y)\|_{H^{-t}}
 \leq\sigma_{t,n}|x-y|.                             \tag{89.15}
\]
Rademacher's theorem then gives
\(|\nabla X_n|\leq\sigma_{t,n}\) almost everywhere.
\end{proof}

\begin{theorem}[Uniform log-Sobolev corridor for tightness]
\label{thm:v16v89-lsi-tight}
Let \(\alpha_n\) be global log-Sobolev constants for \(\mu_n\).
Assume
\[
 \sup_n\frac{\sigma_{t,n}^2}{\alpha_n}\leq V_t<\infty,
 \qquad
 \sup_n{\mathbb E}_{\mu_n}
 \|J_nx\|_{H^{-t}}\leq M_t<\infty.                  \tag{89.16}
\]
Then for every \(p\geq1\),
\[
\begin{split}
 \sup_n{\mathbb E}_{\mu_n}\|J_nx\|_{H^{-t}}^p
 \leq
 2^{p-1}\left[
 M_t^p+
 2\Gamma\left(\frac p2+1\right)(2V_t)^{p/2}
 \right].                                          \tag{89.17}
\end{split}
\]
Consequently the interpolated laws are tight in \(H^{-s}\) for every
\(s>t\).
\end{theorem}

\begin{proof}
Apply Corollary~\ref{cor:v16v89-moments} to \(X_n\), using
Lemma~\ref{lem:v16v89-norm-lipschitz}.  Then
\[
 X_n^p
 \leq2^{p-1}\left(
 |X_n-{\mathbb E}X_n|^p+|{\mathbb E}X_n|^p
 \right).
\]
Equations (89.16) and (89.10) give (89.17).  Apply
Theorem~\ref{thm:v16v88-tightness}.
\end{proof}

\begin{corollary}[Conditional certificate form]
\label{cor:v16v89-conditional-tight}
If cutoff \(n\) satisfies (89.2)--(89.3), it is sufficient in
(89.16) to verify
\[
 \sup_n
 \frac{\sigma_{t,n}^2C_{{\rm AT},n}}
 {\rho_{\min,n}^{\rm LS}}<\infty,
 \qquad
 \sup_n{\mathbb E}\|J_nx\|_{H^{-t}}<\infty.         \tag{89.18}
\]
\end{corollary}

\begin{proof}
Theorem~\ref{thm:v16v89-global-lsi} gives
\(\alpha_n=\rho_{\min,n}^{\rm LS}/C_{{\rm AT},n}\).
Substitution in (89.16) proves the claim.
\end{proof}

\subsection{The mean is an independent infrared quantity}

The concentration estimate controls fluctuations about
\({\mathbb E}X_n\).  It does not control the location of that mean.
For example, translating a fixed log-Sobolev measure by a vector of
length \(n\) leaves its log-Sobolev constant unchanged but makes
\({\mathbb E}\|x\|\) diverge.  Thus the second condition in (89.16)
cannot be deleted.

\begin{proposition}[One sufficient mean bound]
\label{prop:v16v89-mean}
If for some orthonormal basis \(e_k\) of the target Hilbert space,
\[
 \sum_k\sup_n
 {\mathbb E}|\langle J_nx,e_k\rangle|^2
 <\infty,                                           \tag{89.19}
\]
then the mean condition in (89.16) holds.
\end{proposition}

\begin{proof}
Jensen and Tonelli give
\[
\begin{split}
 \left({\mathbb E}\|J_nx\|_{H^{-t}}\right)^2
 &\leq{\mathbb E}\|J_nx\|_{H^{-t}}^2\\
 &=\sum_k{\mathbb E}|\langle J_nx,e_k\rangle|^2
 \leq\sum_k\sup_m
 {\mathbb E}|\langle J_mx,e_k\rangle|^2.
                                                               \tag{89.20}
\end{split}
\]
The final quantity is finite and independent of \(n\).
\end{proof}

\subsection{Why the Poincar\'e gap is weaker}

\begin{proposition}[Poincar\'e does not imply sub-Gaussian
concentration]
\label{prop:v16v89-poincare-counter}
There are probability measures with a positive Poincar\'e gap for
which the conclusion (89.7) fails for a one-Lipschitz observable.
\end{proposition}

\begin{proof}
Let
\[
 d\mu(x)=\frac12e^{-|x|}\,dx.                       \tag{89.21}
\]
The one-dimensional Cheeger inequality for this density gives
\[
 \operatorname{Var}_\mu(f)
 \leq4\int|f'(x)|^2\,d\mu(x),                       \tag{89.22}
\]
so \(\mu\) has a positive Poincar\'e gap.  For the one-Lipschitz
observable \(X(x)=x\),
\[
 {\mathbb E}e^{\lambda X}
 =\frac1{1-\lambda^2}
 \quad\text{only for }|\lambda|<1,                  \tag{89.23}
\]
and the moment-generating function is infinite for
\(|\lambda|\geq1\).  A bound of the form (89.7) would make it finite
for every real \(\lambda\), a contradiction.
\end{proof}

Thus the V83--V84 Poincar\'e certificate supplies covariance control
but not the sub-Gaussian concentration used in
Theorem~\ref{thm:v16v89-lsi-tight}.

\subsection{New entries in the continuum certificate}

For each cutoff and bosonic species, record
\[
\begin{array}{c|c}
\text{conditional entropy floor}&\rho_{\min,n}^{\rm LS}\\
\text{entropy mixing cost}&C_{{\rm AT},n}\\
\text{interpolation scale}&\sigma_{t,n}\\
\text{infrared location}&
M_{t,n}={\mathbb E}\|J_nx\|_{H^{-t}}.
\end{array}                                         \tag{89.24}
\]
The quantitative gate is
\[
 \sup_n
 \frac{\sigma_{t,n}^2C_{{\rm AT},n}}
 {\rho_{\min,n}^{\rm LS}}<\infty,
 \qquad
 \sup_nM_{t,n}<\infty.                              \tag{89.25}
\]
This gate is sufficient for the V88 tightness entry and for uniform
integrability of every fixed polynomial in finitely many
continuous linear observables.

\begin{assessment}[What V89 closes]
\label{ass:v16v89-status}
V89 supplies a rigorous route from conditional log-Sobolev
inequalities and approximate entropy tensorization to global
sub-Gaussian concentration.  Uniform control of the interpolation
Lipschitz scale and of the mean Sobolev norm then produces every
polynomial moment required by V88.

The existing influence and Poincar\'e certificate does not yet prove
the entropy tensorization constant or the mean bound.  These are new
ultraviolet-mixing and infrared-location inputs, respectively.
\end{assessment}

\begin{decision}[V90: companion audit and entropy-mixing options]
\label{dec:v16v89-next}
Perform the mandatory YM/NS audit.  Then compare two noncircular ways
to prove the missing entropy gate: uniform full-Hessian convexity
through the coarse/fine Schur floor, and approximate tensorization
on the physical observable algebra.  Select the branch compatible
with gauge and gravitational zero modes, and continue in V91.
\end{decision}

\begin{firewall}[Claim boundary after V89]
\label{fire:v16v89}
No global log-Sobolev inequality, entropy tensorization estimate, or
uniform Sobolev mean has been proved for the SM--Einstein cutoff law.
V89 establishes the implication from those quantitative inputs to
bosonic tightness.
\end{firewall}

\section*{Version V89 append-only record}
\addcontentsline{toc}{section}{Version V89 append-only record}
The complete V88 source was retained byte for byte before this
continuation.  It had \(912138\) bytes, \(27504\) lines, and SHA--256
\[
 \texttt{E5E5CE48BC82BE42B4E068018F7B24058D18CD728BD73F6A475F4B53393DC015}.
\]
V89 proved the log-Sobolev concentration corridor leading to the
negative-Sobolev moments of V88.


\section{V90: companion audit and the entropy-mixing branch}
\label{sec:v16v90}

\subsection{Mandatory companion audit}

The newest active companion files are
\[
\begin{array}{lll}
\text{YM V64:}&14831\text{ lines},&
\texttt{69B9A7675E3F46414962EF6CD74AAD10E17E1EEA975C543B78419885D2B34808},\\
\text{NS V26:}&5319\text{ lines},&
\texttt{82C887F8C2C69E4F28FAD7C3DC8DE5B9099CB5A4BF431EBA1FFF3E91935978AD}.
\end{array}                                                   \tag{90.1}
\]
YM V63 found obstructions to a uniform full-radius splitting in the
fourth third-coordinate layer.  V64 repaired them with
centre-specific rational bridges and proves both direct signs on the
paired slabs through depth \(2t\).  The terminal layer and cap, the
full torus, and all continuum conclusions remain open.  This is a
useful structural warning: local positive charts and bridge
certificates need not assemble into one globally convex coordinate
chart.  NS V26 is unchanged and contributes no entropy inequality.

\begin{theorem}[V90 audit decision]
\label{thm:v16v90-audit}
The full-Hessian route may be used on an individual gauge-fixed
chart, but it cannot be adopted as the global entropy mechanism for
the SM--Einstein programme.  The primary branch will instead seek
approximate entropy tensorization for a common physical observable
algebra through gauge-compatible block resampling.
\end{theorem}

\begin{proof}
The YM bridge construction directly exhibits the need for multiple
local charts even at the symbol level.  The no-go result below shows
independently that strict global Hessian convexity is incompatible
with compact periodic gauge directions.  Conditional resampling on
the physical algebra does not require coercivity along a gauge
orbit.
\end{proof}

\subsection{Branch A: what a full Schur floor would prove}

Consider a \(C^2\) finite-cutoff action in Euclidean fine/coarse
coordinates with Hessian
\[
 {\mathcal H}=
 \begin{pmatrix}
 H_{ff}&H_{fc}\\
 H_{cf}&H_{cc}
 \end{pmatrix},
 \qquad
 S_c=H_{cc}-H_{cf}H_{ff}^{-1}H_{fc}.                \tag{90.2}
\]

\begin{lemma}[Quantitative Schur-to-Hessian floor]
\label{lem:v16v90-schur-floor}
Assume, pointwise in configuration space,
\[
 H_{ff}\succeq aI,\qquad
 S_c\succeq bI,\qquad
 \|H_{ff}^{-1}H_{fc}\|\leq\gamma,                   \tag{90.3}
\]
where \(a,b>0\).  Then
\[
 {\mathcal H}\succeq\lambda I,\qquad
 \lambda=
 \frac{\min\{a,b\}}{\max\{2,1+2\gamma^2\}}.         \tag{90.4}
\]
\end{lemma}

\begin{proof}
Put \(B=H_{ff}^{-1}H_{fc}\) and
\(w=v_f+Bv_c\).  Completion of the square gives
\[
\begin{split}
 \left\langle
 \binom{v_f}{v_c},
 {\mathcal H}\binom{v_f}{v_c}
 \right\rangle
 =
 \langle w,H_{ff}w\rangle+\langle v_c,S_cv_c\rangle
 \geq a|w|^2+b|v_c|^2.                             \tag{90.5}
\end{split}
\]
Moreover,
\[
\begin{split}
 |v_f|^2+|v_c|^2
 &=|w-Bv_c|^2+|v_c|^2\\
 &\leq2|w|^2+(1+2\gamma^2)|v_c|^2\\
 &\leq
 \max\{2,1+2\gamma^2\}(|w|^2+|v_c|^2).
                                                               \tag{90.6}
\end{split}
\]
Equations (90.5)--(90.6) prove (90.4).
\end{proof}

\begin{theorem}[Chartwise Bakry--\'Emery route]
\label{thm:v16v90-bakry}
If a probability density on \(\mathbb R^N\) has the form
\[
 d\mu=Z^{-1}e^{-U(x)}\,dx
\]
and \(\nabla^2U\succeq\lambda I\) globally for some \(\lambda>0\),
then
\[
 \operatorname{Ent}_\mu(F^2)
 \leq\frac2\lambda\int|\nabla F|^2\,d\mu.           \tag{90.7}
\]
Thus (90.3) would give the V89 global LSI with the value (90.4).
\end{theorem}

\begin{proof}
For the diffusion generator
\({\mathcal L}=\Delta-\nabla U\cdot\nabla\), the carré du champ is
\(\Gamma(f)=|\nabla f|^2\), and direct differentiation gives
\[
 \Gamma_2(f)
 =
 \|\nabla^2f\|_{\rm HS}^2+
 \langle\nabla f,(\nabla^2U)\nabla f\rangle
 \geq\lambda\Gamma(f).                              \tag{90.8}
\]
Let \(P_t\) be its Markov semigroup.  The standard entropy
dissipation identity and the gradient estimate obtained by
differentiating
\(P_t|\nabla f|^2- e^{2\lambda t}|\nabla P_tf|^2\)
give
\[
 \operatorname{Ent}_\mu(F^2)
 =
 \int_0^\infty
 \int\frac{|\nabla P_t(F^2)|^2}{P_t(F^2)}\,d\mu\,dt
 \leq
 \frac2\lambda\int|\nabla F|^2\,d\mu.               \tag{90.9}
\]
This is (90.7).
\end{proof}

\subsection{Why branch A cannot be global}

\begin{proposition}[Compact-direction convexity obstruction]
\label{prop:v16v90-compact-no-go}
Let \(N\) be a closed Riemannian manifold of positive dimension.
There is no \(C^2\) function \(U:N\to\mathbb R\) satisfying
\[
 \operatorname{Hess}U\succeq\lambda g
 \quad\text{for any }\lambda>0.                     \tag{90.10}
\]
\end{proposition}

\begin{proof}
Taking the metric trace in (90.10) gives
\[
 \Delta U\geq\lambda\dim N.
\]
Integration over \(N\) contradicts
\[
 \int_N\Delta U\,d{\rm vol}=0
\]
by the divergence theorem.
\end{proof}

Periodic gauge holonomies contain precisely such compact
directions.  Before quotienting, exact gauge and diffeomorphism
orbits also produce null directions.  After quotienting, Gribov
charts and orbit-space singularities prevent the immediate use of
one global Euclidean Hessian.  Lemma~\ref{lem:v16v90-schur-floor}
remains useful inside a chart and for the local conditional floors;
it is not the global entropy proof.

\subsection{Branch B: entropy loss under physical block resampling}

Let \(E_i\) be conditional expectation obtained by resampling block
\(i\) while fixing its exterior.  For a density \(h\geq0\) with
\(\int h\,d\mu=1\), define
\[
 {\cal E}_i(h)
 :=
 \int\left[
 E_i(h\log h)-(E_ih)\log(E_ih)
 \right]d\mu.                                       \tag{90.11}
\]
This is the integrated conditional entropy of \(h\) in block \(i\).

\begin{lemma}[Exact entropy loss of one heat-bath update]
\label{lem:v16v90-one-update}
For every block \(i\),
\[
 \operatorname{Ent}_\mu(h)
 -\operatorname{Ent}_\mu(E_ih)
 ={\cal E}_i(h).                                    \tag{90.12}
\]
\end{lemma}

\begin{proof}
Because \(E_i\) preserves the integral of \(h\),
\[
\begin{split}
 \operatorname{Ent}_\mu(h)
 -\operatorname{Ent}_\mu(E_ih)
 &=
 \int h\log h\,d\mu
 -\int(E_ih)\log(E_ih)\,d\mu\\
 &=
 \int\left[
 E_i(h\log h)-(E_ih)\log(E_ih)
 \right]d\mu.
                                                               \tag{90.13}
\end{split}
\]
\end{proof}

Choose a finite sweep \(i_1,\ldots,i_m\), set
\[
 h_0=h,\qquad h_k=E_{i_k}h_{k-1}.                  \tag{90.14}
\]

\begin{theorem}[Sweep criterion for approximate tensorization]
\label{thm:v16v90-sweep}
Let \({\cal H}_{\rm phys}\) be a class of densities stable under the
updates \(E_i\).  Suppose there are
\(0\leq\eta<1\) and \(C_{\rm mix}<\infty\) such that every
\(h\in{\cal H}_{\rm phys}\) satisfies
\[
 \operatorname{Ent}_\mu(h_m)
 \leq\eta\operatorname{Ent}_\mu(h)                  \tag{90.15}
\]
and
\[
 \sum_{k=1}^m{\cal E}_{i_k}(h_{k-1})
 \leq
 C_{\rm mix}\sum_i{\cal E}_i(h).                   \tag{90.16}
\]
Then
\[
 \operatorname{Ent}_\mu(h)
 \leq
 \frac{C_{\rm mix}}{1-\eta}
 \sum_i{\cal E}_i(h)
 \qquad(h\in{\cal H}_{\rm phys}).                   \tag{90.17}
\]
\end{theorem}

\begin{proof}
Telescoping Lemma~\ref{lem:v16v90-one-update} gives
\[
\begin{split}
 \operatorname{Ent}_\mu(h)-\operatorname{Ent}_\mu(h_m)
 =
 \sum_{k=1}^m{\cal E}_{i_k}(h_{k-1}).               \tag{90.18}
\end{split}
\]
The left side is at least
\((1-\eta)\operatorname{Ent}_\mu(h)\) by (90.15), and the right side
is bounded by (90.16).  Division by \(1-\eta\) proves (90.17).
\end{proof}

For \(h=F^2/\int F^2d\mu\), equation (90.17) is (89.3) on the chosen
physical function class, with
\[
 C_{\rm AT}=\frac{C_{\rm mix}}{1-\eta}.             \tag{90.19}
\]
Combining it with the conditional LSIs (89.2) proves the restricted
global LSI needed for concentration of physical observables.

\subsection{Gauge compatibility of the sweep}

\begin{proposition}[Descent to the physical algebra]
\label{prop:v16v90-gauge-sweep}
Let a symmetry group \({\cal G}\) act on configurations and suppose
each update satisfies
\[
 E_i(F\circ g)=(E_iF)\circ g
 \quad(g\in{\cal G}).                               \tag{90.20}
\]
Then the invariant function class is stable under every \(E_i\).
Consequently Theorem~\ref{thm:v16v90-sweep} may be proved entirely
on invariant densities.
\end{proposition}

\begin{proof}
If \(F\circ g=F\), then (90.20) gives
\[
 (E_iF)\circ g=E_i(F\circ g)=E_iF.
\]
Products, positive squares, and normalization preserve invariance,
so every density generated by the invariant observable algebra
remains in that class throughout the sweep.
\end{proof}

For lattice gauge theory, resampling one link without its boundary
transformation data need not satisfy (90.20).  Blocks must therefore
include a gauge-covariant boundary sector or be defined directly on
orbit variables.  For gravity, an analogous construction must track
diffeomorphisms of the block boundary.  This is now the concrete
content of the entropy-mixing bottleneck.

\subsection{Selected quantitative ledger}

At cutoff \(n\), the selected branch records
\[
\begin{array}{c|c}
\text{sweep contraction}&\eta_n<1\\
\text{intermediate-entropy stability}&C_{{\rm mix},n}\\
\text{conditional LSI floor}&\rho_{\min,n}^{\rm LS}\\
\text{physical interpolation scale}&\sigma_{t,n}.
\end{array}                                         \tag{90.21}
\]
By V89 and (90.19), the ultraviolet tightness ratio becomes
\[
 \frac{\sigma_{t,n}^2C_{{\rm mix},n}}
 {(1-\eta_n)\rho_{\min,n}^{\rm LS}}.                \tag{90.22}
\]
Uniform boundedness of (90.22), together with the mean bound, is the
selected entropy gate.

\begin{assessment}[What V90 closes]
\label{ass:v16v90-status}
V90 proves that a positive fine block and positive coarse Schur
complement yield a quantitative full Hessian floor and hence a
chartwise Bakry--\'Emery LSI.  A compact-direction theorem shows why
this cannot be the global gauge mechanism.

The selected branch is the exact entropy-loss identity for
gauge-compatible block resampling.  Sweep contraction and
intermediate-entropy stability imply approximate tensorization with
constant \(C_{\rm mix}/(1-\eta)\).  These two sweep estimates remain
to be proved for the physical SM--Einstein blocks.
\end{assessment}

\begin{decision}[V91: physical observable coordinates]
\label{dec:v16v90-next}
Continue after the audit by replacing raw gauge and metric
coordinates with a countable separating family of bounded physical
cylinder observables.  Prove tightness of their joint laws in a
compact product, construct the limiting positive functional, and
state the additional consistency needed for it to represent a
distributional field theory rather than only an abstract observable
law.
\end{decision}

\begin{firewall}[Claim boundary after V90]
\label{fire:v16v90}
Neither sweep contraction nor entropy stability has been established
for the SM--Einstein law.  The global LSI, bosonic tightness,
reflection-positive limit, and full continuum theory therefore
remain conditional.
\end{firewall}

\section*{Version V90 append-only record}
\addcontentsline{toc}{section}{Version V90 append-only record}
The complete V89 source was retained byte for byte before this
continuation.  It had \(923213\) bytes, \(27877\) lines, and SHA--256
\[
 \texttt{CA7FC92EC97B6D862FE75A06C92273F009D0E207C581EC2377F725BC6473A600}.
\]
V90 performed the required companion audit, proved the two
entropy-mixing mechanisms, and selected gauge-compatible physical
block resampling as the global branch.


\section{V91: compact physical-observable coordinates}
\label{sec:v16v91}

\subsection{A countable physical coordinate family}

Let
\[
 {\cal O}_1,{\cal O}_2,\ldots                         \tag{91.1}
\]
be a countable family of gauge- and diffeomorphism-invariant
observables chosen from a common continuum label set.  Examples of
the intended types are traced holonomies on a countable loop family,
smeared invariant composites on a countable dense test-function
family, and relational curvature observables.  Their existence and
separation properties are hypotheses, not conclusions of this
section.

For every \(k\), choose a compact metric value space \(K_k\).
Already bounded observables may take values directly in \(K_k\).
For a real unbounded observable \(Q_k\), use the compactifying map
\[
 b(q)=\frac{q}{\sqrt{1+q^2}}\in(-1,1)               \tag{91.2}
\]
and take \(K_k=[-1,1]\).  At cutoff \(n\), let
\[
 Y_n=(Y_{1,n},Y_{2,n},\ldots)\in
 K:=\prod_{k=1}^\infty K_k                          \tag{91.3}
\]
be the corresponding physical observable vector, and denote its law
by
\[
 \nu_n=(Y_n)_\#\mu_n.                               \tag{91.4}
\]

\subsection{Automatic compactness of bounded coordinates}

\begin{lemma}[The countable compact product]
\label{lem:v16v91-product}
The space \(K\) is compact and metrizable.  If each metric \(d_k\)
has diameter at most one, a compatible metric is
\[
 d_K(x,y)=\sum_{k=1}^\infty2^{-k}d_k(x_k,y_k).       \tag{91.5}
\]
\end{lemma}

\begin{proof}
The series in (91.5) converges uniformly and defines the product
topology.  Given a sequence in \(K\), compactness of \(K_1\) gives a
subsequence converging in the first coordinate; iterating and taking
the diagonal subsequence gives convergence in every coordinate.
The tail of (91.5) is uniformly small, so coordinatewise convergence
implies convergence in \(d_K\).  Thus \(K\) is sequentially compact.
As a metric space it is compact.
\end{proof}

\begin{theorem}[Physical-coordinate subsequential law]
\label{thm:v16v91-coordinate-limit}
Every sequence \(\{\nu_n\}\) has a subsequence
\(\nu_{n_j}\) converging weakly to a probability measure \(\nu\) on
\(K\).  For every bounded continuous cylinder function
\[
 F(x)=f(x_{k_1},\ldots,x_{k_r}),                    \tag{91.6}
\]
one has
\[
 \lim_{j\to\infty}{\mathbb E}_{\mu_{n_j}}F(Y_{n_j})
 =
 \int_KF\,d\nu.                                     \tag{91.7}
\]
\end{theorem}

\begin{proof}
All probability laws on the compact metric space \(K\) form a tight
family.  Prokhorov's theorem gives a weakly convergent subsequence and
a probability limit.  Equation (91.7) is the definition of weak
convergence.
\end{proof}

This construction needs no raw gauge-fixed field tightness.  Its
price is that it initially produces only an abstract law of bounded
physical coordinates.

\subsection{The limiting positive functional}

Let \({\cal C}_{\rm cyl}(K)\) be the algebra of continuous cylinder
functions on \(K\), and define
\[
 {\cal S}(F):=\int_KF\,d\nu.                         \tag{91.8}
\]

\begin{theorem}[Positive normalized physical functional]
\label{thm:v16v91-positive-functional}
The functional \({\cal S}\) is linear, normalized, and positive:
\[
 {\cal S}(1)=1,\qquad
 {\cal S}(\overline FF)\geq0.                       \tag{91.9}
\]
Its finite-coordinate marginals are mutually consistent.  Moreover,
\({\cal C}_{\rm cyl}(K)\) is uniformly dense in \(C(K)\).
\end{theorem}

\begin{proof}
The first two claims follow from integration against the probability
measure \(\nu\).  Every finite-coordinate marginal is the pushforward
of the same measure \(\nu\), so marginalization gives consistency.
The cylinder algebra contains constants, is closed under complex
conjugation, and separates points of \(K\), because two different
points differ in some coordinate.  The complex Stone--Weierstrass
theorem therefore gives uniform density in \(C(K)\).
\end{proof}

\subsection{Reflection positivity in coordinate space}

Assume time reflection induces a continuous involution
\(\vartheta_K:K\to K\) and preserves a positive-time cylinder
subalgebra \({\cal C}_+\).  Define
\[
 (\Theta F)(x)=\overline{F(\vartheta_Kx)}.           \tag{91.10}
\]

\begin{theorem}[Coordinate-limit reflection positivity]
\label{thm:v16v91-coordinate-rp}
Suppose every cutoff coordinate law satisfies
\[
 \int_K(\Theta F)(x)F(x)\,d\nu_n(x)\geq0
 \qquad(F\in{\cal C}_+).                            \tag{91.11}
\]
Then every weak subsequential limit \(\nu\) satisfies the same
inequality.
\end{theorem}

\begin{proof}
For fixed \(F\in{\cal C}_+\), the integrand
\(\overline{F(\vartheta_Kx)}F(x)\) is bounded and continuous.
Weak convergence and (91.11) give
\[
 \int_K\overline{F(\vartheta_Kx)}F(x)\,d\nu(x)
 =
 \lim_{j\to\infty}
 \int_K\overline{F(\vartheta_Kx)}F(x)\,d\nu_{n_j}(x)
 \geq0.                                             \tag{91.12}
\]
Apply the same argument to every finite linear combination to obtain
the full reflected Gram-matrix condition.
\end{proof}

The hypothesis (91.11) is still a substantive finite-cutoff
condition.  Positivity of the ordinary probability law \(\nu_n\)
does not imply reflection positivity.

\subsection{No escape through compactification endpoints}

The compactifying transform (91.2) may hide divergence by placing
limiting mass at \(1\) or \(-1\).

\begin{lemma}[Endpoint mass and uniform tails]
\label{lem:v16v91-endpoint}
Let \(Q_n\) be real random variables and
\(Y_n=b(Q_n)\).  If
\[
 \lim_{R\to\infty}\sup_n{\mathbb P}(|Q_n|>R)=0,     \tag{91.13}
\]
then every weak limit of \(Y_n\) on \([-1,1]\) gives zero mass to
\(\{-1,1\}\).
\end{lemma}

\begin{proof}
For \(0<a<1\),
\[
 |b(q)|\geq a
 \quad\Longleftrightarrow\quad
 |q|\geq\frac{a}{\sqrt{1-a^2}}.                    \tag{91.14}
\]
The closed-set Portmanteau bound gives
\[
 \nu(\{|y|\geq a\})
 \leq\limsup_j{\mathbb P}(|Y_{n_j}|\geq a).         \tag{91.15}
\]
As \(a\uparrow1\), the threshold on the right of (91.14) tends to
infinity, so (91.13) makes the right side tend to zero.
Continuity from above gives
\(\nu(\{-1,1\})=0\).
\end{proof}

\begin{theorem}[Recovery of countably many finite observables]
\label{thm:v16v91-recovery}
Suppose (91.13) holds for every compactified coordinate \(Q_{k,n}\).
Then every subsequential coordinate limit \(\nu\) is supported on
\[
 K^\circ:=
 \prod_{k:\,Q_k\ {\rm unbounded}}(-1,1)
 \times
 \prod_{k:\,{\cal O}_k\ {\rm bounded}}K_k.          \tag{91.16}
\]
Consequently the inverse maps \(b^{-1}\) define finite real random
variables \(Q_k\) simultaneously, \(\nu\)-almost surely.
\end{theorem}

\begin{proof}
By Lemma~\ref{lem:v16v91-endpoint}, the event that coordinate \(k\)
is an endpoint has \(\nu\)-measure zero.  There are only countably
many coordinates, so the union of all endpoint events has measure
zero.  On the complement, every \(b^{-1}(Y_k)\) is finite.
\end{proof}

Separate coordinatewise tightness is enough for existence in the
countable product topology.  It is not enough for Sobolev
regularity, locality, or a continuum energy bound.

\subsection{Algebraic relations survive when they are closed}

\begin{proposition}[Passage of exact bounded relations]
\label{prop:v16v91-relations}
Let \(R:K_{k_1}\times\cdots\times K_{k_r}\to\mathbb C\) be continuous.
If
\[
 R(Y_{k_1,n},\ldots,Y_{k_r,n})=0
 \quad\mu_n\text{-almost surely for every }n,        \tag{91.17}
\]
then
\[
 R(Y_{k_1},\ldots,Y_{k_r})=0
 \quad\nu\text{-almost surely}.                     \tag{91.18}
\]
\end{proposition}

\begin{proof}
For \(\varepsilon>0\), the set
\[
 C_\varepsilon=\{x:|R(x_{k_1},\ldots,x_{k_r})|
 \geq\varepsilon\}
\]
is closed.  It has \(\nu_n\)-measure zero, so the Portmanteau theorem
gives
\[
 \nu(C_\varepsilon)\leq
 \limsup_j\nu_{n_j}(C_\varepsilon)=0.
\]
Take the countable union over positive rational \(\varepsilon\).
\end{proof}

This preserves finite Mandelstam-type relations, involution
relations, and other continuous identities already exact at cutoff.
Relations involving derivatives, shrinking loops, or unbounded
products require quantitative continuity or uniform integrability.

\subsection{What is needed to recover a field theory}

The coordinate law becomes more than an abstract compactification
only after the following are proved.
\begin{enumerate}
\item \emph{Separation:} the chosen invariant observables separate
      physical configurations, or their generalized completion,
      modulo null sets.
\item \emph{Label continuity:} when test functions, loops, or
      relational supports converge, the corresponding random
      coordinates converge in probability or in a controlled
      \(L^p\) norm.
\item \emph{Locality:} disjoint-support algebraic relations and the
      net of local subalgebras survive the limit.
\item \emph{No escape and moments:} endpoint mass vanishes and every
      unbounded Schwinger function used later is uniformly
      integrable.
\item \emph{Symmetry:} the countable rational Euclidean action
      extends continuously to the full group; reflection acts
      continuously as assumed in (91.10).
\end{enumerate}
For linear distribution-valued fields, these entries are usually
encoded by continuity of the characteristic functional on the test
function space.  For Wilson loops and relational gravitational
observables, an analogous continuity theorem must be proved.

\subsection{Insertion into the programme}

The physical-coordinate route and the raw-field route of V88 are
complementary:
\[
\begin{array}{c|c|c}
\text{route}&\text{compactness input}&\text{output}\\
\hline
\text{raw fields}&H^{-t}\text{ moment}&
\text{distribution-valued subsequential law}\\
\text{physical coordinates}&\text{boundedness}&
\text{abstract invariant coordinate law}.
\end{array}                                         \tag{91.19}
\]
The second route is compatible with gauge zero modes but requires
the five reconstruction entries above.  The first route has stronger
geometric content but requires gauge-compatible Sobolev bounds.

\begin{assessment}[What V91 closes]
\label{ass:v16v91-status}
V91 constructs subsequential probability laws for any countable
family of bounded physical observables and proves positivity,
consistency, preservation of continuous algebraic relations, and
conditional preservation of reflection positivity.  It also gives
an exact no-escape criterion for compactified unbounded observables.

This is an abstract observable continuum.  It is not yet a
distributional local field theory because separation, label
continuity, locality, full Euclidean symmetry, and unbounded moment
control remain open.
\end{assessment}

\begin{decision}[V92: extend rational Euclidean symmetry]
\label{dec:v16v91-next}
Prove that invariance under a countable dense subgroup of Euclidean
motions extends to the full connected Euclidean group when the
observable labels are stochastically continuous.  Formulate the
uniform cutoff modulus that passes this continuity to the limit, and
separate fixed-background Euclidean symmetry from dynamical-metric
diffeomorphism invariance.
\end{decision}

\begin{firewall}[Claim boundary after V91]
\label{fire:v16v91}
No separating physical coordinate family with the required
continuity, locality, reflection positivity, or no-escape bounds has
yet been verified for the full SM--Einstein cutoff law.  V91 proves
the compact coordinate mechanism conditional on those choices.
\end{firewall}

\section*{Version V91 append-only record}
\addcontentsline{toc}{section}{Version V91 append-only record}
The complete V90 source was retained byte for byte before this
continuation.  It had \(934630\) bytes, \(28241\) lines, and SHA--256
\[
 \texttt{5FD27E89BB9FDA0EEFA5517CCD6379CB8CBAA359B51371CE48EB9282FF3AC5B5}.
\]
V91 constructed compact subsequential laws for countable bounded
physical-observable coordinates and identified the conditions needed
to recover a local continuum theory.


\section{V92: from rational to full Euclidean symmetry}
\label{sec:v16v92}

\subsection{The elementary dense-subgroup closure}

Let \(G\) be a second-countable topological group and
\(G_0\subset G\) a countable dense subgroup.  Let
\(\tau_g\) act by algebra automorphisms on an observable algebra
\({\cal A}\).

\begin{lemma}[Dense-subgroup invariance]
\label{lem:v16v92-dense}
Let \({\cal S}:{\cal A}\to\mathbb C\) satisfy
\[
 {\cal S}(\tau_hF)={\cal S}(F)
 \quad(h\in G_0).                                   \tag{92.1}
\]
If for every \(F\in{\cal A}\) the map
\[
 g\longmapsto{\cal S}(\tau_gF)                      \tag{92.2}
\]
is continuous on \(G\), then \({\cal S}\) is \(G\)-invariant.
\end{lemma}

\begin{proof}
For \(g\in G\), choose \(h_j\in G_0\) with \(h_j\to g\).  Then
\[
 {\cal S}(\tau_gF)
 =
 \lim_{j\to\infty}{\cal S}(\tau_{h_j}F)
 =
 {\cal S}(F).                                       \tag{92.3}
\]
\end{proof}

Thus the missing content is continuity, not the cardinality of the
symmetry group.

\subsection{Uniform label continuity at cutoff}

Let \(({\cal L},d_{\cal L})\) be a separable metric label space.
A label \(\xi\) may encode a test function, loop, frame, or
relational support.  Let \(Y_{\xi,n}\) be a bounded physical cutoff
observable, with
\[
 |Y_{\xi,n}|\leq1.                                  \tag{92.4}
\]

\begin{definition}[Uniform stochastic modulus]
\label{def:v16v92-modulus}
A function \(\omega:[0,\infty)\to[0,\infty)\) is a uniform
\(L^1\) label modulus when
\[
 \omega(r)\longrightarrow0\quad(r\downarrow0)
                                                               \tag{92.5}
\]
and
\[
 \sup_n{\mathbb E}_{\mu_n}
 |Y_{\xi,n}-Y_{\zeta,n}|
 \leq\omega(d_{\cal L}(\xi,\zeta))                  \tag{92.6}
\]
for all \(\xi,\zeta\in{\cal L}\).
\end{definition}

\begin{theorem}[Continuity passes to a coordinate-law limit]
\label{thm:v16v92-modulus-limit}
Let \({\cal L}_0\subset{\cal L}\) be countable and dense.  Suppose
the joint laws of
\(\{Y_{\xi,n}:\xi\in{\cal L}_0\}\) converge along a subsequence to
\(\{Y_\xi:\xi\in{\cal L}_0\}\), and assume (92.6).  Then
\[
 {\mathbb E}|Y_\xi-Y_\zeta|
 \leq\omega(d_{\cal L}(\xi,\zeta))
 \quad(\xi,\zeta\in{\cal L}_0).                     \tag{92.7}
\]
There is a unique extension
\[
 \xi\longmapsto Y_\xi\in L^1(\nu)                  \tag{92.8}
\]
to all of \({\cal L}\), and the extension obeys the same modulus.
\end{theorem}

\begin{proof}
For \(\xi,\zeta\in{\cal L}_0\), the function
\((y_\xi,y_\zeta)\mapsto|y_\xi-y_\zeta|\) is bounded and continuous.
Finite-coordinate weak convergence therefore passes (92.6) to the
limit and gives (92.7).

For \(\xi\in{\cal L}\), choose
\(\xi_j\in{\cal L}_0\) with \(\xi_j\to\xi\).  Equation (92.7) makes
\(\{Y_{\xi_j}\}\) Cauchy in the complete space \(L^1(\nu)\); define
\(Y_\xi\) as its limit.  Two approximating sequences give the same
limit by interlacing them and using (92.7).  Passing to the limit in
(92.7) proves the modulus for arbitrary labels.
\end{proof}

\begin{corollary}[Continuity of cylinder expectations]
\label{cor:v16v92-cylinder-cont}
Let \(f:\mathbb C^r\to\mathbb C\) be bounded and Lipschitz.
Then
\[
 (\xi_1,\ldots,\xi_r)\longmapsto
 {\mathbb E}_\nu f(Y_{\xi_1},\ldots,Y_{\xi_r})       \tag{92.9}
\]
is continuous.
\end{corollary}

\begin{proof}
If \(L_f\) is the Lipschitz constant, then the difference of two
expectations is at most
\[
 L_f\sum_{a=1}^r
 {\mathbb E}|Y_{\xi_a}-Y_{\zeta_a}|
 \leq
 L_f\sum_{a=1}^r\omega(d_{\cal L}(\xi_a,\zeta_a)).
                                                               \tag{92.10}
\]
\end{proof}

\subsection{Varying lattice symmetry groups}

At cutoff \(n\), the exact symmetry group \(G_n\) is generally only a
lattice subgroup.  It need not contain one fixed dense subgroup.
We therefore state an approximation theorem.

\begin{theorem}[Approximate cutoff symmetry implies exact limiting
symmetry]
\label{thm:v16v92-approx-symmetry}
Suppose \(G\) acts continuously on \({\cal L}\).  For every
\(g\in G\), assume there are \(g_n\in G_n\) such that
\[
 \delta_n(g,\xi):=
 d_{\cal L}(g_n\xi,g\xi)\longrightarrow0            \tag{92.11}
\]
for each label \(\xi\).  Assume the cutoff laws are exactly invariant:
\[
 (Y_{\xi_1,n},\ldots,Y_{\xi_r,n})
 \stackrel{\rm law}{=}
 (Y_{g_n\xi_1,n},\ldots,Y_{g_n\xi_r,n})             \tag{92.12}
\]
for all finite label families.  If (92.6) holds and the joint laws on
a dense label set converge as in
Theorem~\ref{thm:v16v92-modulus-limit}, then the extended limiting
process is \(G\)-invariant:
\[
 (Y_{\xi_1},\ldots,Y_{\xi_r})
 \stackrel{\rm law}{=}
 (Y_{g\xi_1},\ldots,Y_{g\xi_r}).                    \tag{92.13}
\]
\end{theorem}

\begin{proof}
Fix labels \(\xi_a\).  Choose
\(\zeta_a^{(m)}\in{\cal L}_0\) tending to \(\xi_a\), and choose
\(\eta_a^{(m)}\in{\cal L}_0\) so that
\[
 d_{\cal L}(\eta_a^{(m)},g\zeta_a^{(m)})\longrightarrow0.
                                                               \tag{92.14}
\]
For a bounded Lipschitz \(f\), cutoff invariance gives
\[
 {\mathbb E}f(Y_{\zeta_1^{(m)},n},\ldots,Y_{\zeta_r^{(m)},n})
 =
 {\mathbb E}f(Y_{g_n\zeta_1^{(m)},n},\ldots,
                  Y_{g_n\zeta_r^{(m)},n}).           \tag{92.15}
\]
Compare the right side with the cutoff vector labelled by
\(\eta_a^{(m)}\).  Equation (92.6) bounds the difference of the two
expectations by
\[
 L_f\sum_{a=1}^r
 \omega\!\left(
 d_{\cal L}(g_n\zeta_a^{(m)},\eta_a^{(m)})
 \right).                                           \tag{92.16}
\]
For fixed \(m\), take \(n=n_j\to\infty\).  Both endpoint label
families lie in \({\cal L}_0\), while (92.11) turns (92.16) into the
same expression with \(g\zeta_a^{(m)}\) in place of
\(g_n\zeta_a^{(m)}\).  Now let \(m\to\infty\).  Equations
(92.14), (92.7), and continuity of the \(G\)-action show that the two
limiting expectations are those of the vectors labelled by
\((\xi_a)\) and \((g\xi_a)\), and their difference is zero.
Bounded Lipschitz functions determine probability laws on finite
products of compact metric spaces, proving (92.13).
\end{proof}

\subsection{A usable test-function modulus}

For a distribution-valued field \(\Phi_n\), take labels
\(\xi=f\) in a normed test-function space \({\cal T}\), and consider
a one-Lipschitz compactification \(b\):
\[
 Y_{f,n}=b(\langle\Phi_n,f\rangle).                 \tag{92.17}
\]

\begin{proposition}[Sobolev moment implies a label modulus]
\label{prop:v16v92-sobolev-modulus}
If
\[
 \sup_n{\mathbb E}\|\Phi_n\|_{H^{-s}}\leq M_s,
                                                               \tag{92.18}
\]
then
\[
 \sup_n{\mathbb E}|Y_{f,n}-Y_{h,n}|
 \leq M_s\|f-h\|_{H^s}.                             \tag{92.19}
\]
\end{proposition}

\begin{proof}
Because \(b\) is one-Lipschitz,
\[
\begin{split}
 |Y_{f,n}-Y_{h,n}|
 &\leq|\langle\Phi_n,f-h\rangle|\\
 &\leq\|\Phi_n\|_{H^{-s}}\|f-h\|_{H^s}.
                                                               \tag{92.20}
\end{split}
\]
Take expectations and use (92.18).
\end{proof}

Thus the V88--V89 Sobolev corridor supplies precisely the continuity
needed to extend a rational test-function construction.  For loop
observables, an analogous estimate must control the response to
deformation of the loop; boundedness alone gives no such modulus.

\subsection{Full Euclidean invariance of the limiting functional}

Take \(G=E(d)\), or its orientation-preserving component if
reflection is treated separately.  Let \({\cal A}_{\rm cyl}\) be
generated by the extended variables \(Y_\xi\), and define
\[
 \tau_gY_\xi=Y_{g\xi}.                              \tag{92.21}
\]

\begin{corollary}[Euclidean closure]
\label{cor:v16v92-euclidean}
Under Theorem~\ref{thm:v16v92-approx-symmetry}, the limiting
functional satisfies
\[
 {\cal S}(\tau_gF)={\cal S}(F)
 \quad(g\in E(d),\ F\in{\cal A}_{\rm cyl})          \tag{92.22}
\]
for bounded continuous cylinder \(F\).
\end{corollary}

\begin{proof}
Equation (92.13) gives equality of all finite-coordinate
expectations.  These are exactly (92.22).
\end{proof}

\subsection{Background Euclidean symmetry and gravitational symmetry}

The preceding theorem concerns a fixed label space carrying an
\(E(d)\)-action.  In a dynamical-metric theory it does not establish
diffeomorphism invariance.  The two requirements are:
\[
\begin{array}{c|c}
\text{OS/relativistic reconstruction}&
\text{Euclidean covariance of Schwinger functions}\\
\text{gravitational gauge principle}&
\text{diffeomorphism-invariant physical observables}.
\end{array}                                         \tag{92.23}
\]
A relational observable may carry both an asymptotic Euclidean label
and internal dynamical reference data.  Its cutoff approximation must
obey the uniform modulus with respect to both.  Averaging a
coordinate-local field over diffeomorphisms is not a substitute
unless the average exists and preserves locality, reflection
positivity, and the no-escape bounds.

\subsection{New continuity ledger}

For each observable type, record a label metric and a cutoff-uniform
modulus:
\[
 \Omega_{\cal O}(r)
 :=
 \sup_n\sup_{d_{\cal L}(\xi,\zeta)\leq r}
 {\mathbb E}|Y_{\xi,n}-Y_{\zeta,n}|.                \tag{92.24}
\]
The symmetry-continuity gate is
\[
 \Omega_{\cal O}(r)\longrightarrow0
 \quad(r\downarrow0),                               \tag{92.25}
\]
together with the approximation errors (92.11) for the regulator
symmetry groups.

\begin{assessment}[What V92 closes]
\label{ass:v16v92-status}
V92 proves that a cutoff-uniform stochastic label modulus extends the
countable coordinate law to all labels and turns approximate lattice
symmetry into exact limiting Euclidean symmetry.  A uniform
negative-Sobolev first moment supplies the modulus for smeared
distributional fields.

No such modulus has yet been proved for Wilson loops, relational
gravitational observables, or the full interacting cutoff law.
Euclidean symmetry and diffeomorphism invariance remain distinct
requirements.
\end{assessment}

\begin{decision}[V93: locality under the observable limit]
\label{dec:v16v92-next}
Prove a closed-support formulation of locality for bounded physical
coordinates.  Show that exact cutoff commutation/factorization
relations pass to the limit when support labels converge with the
V92 modulus, and distinguish algebraic locality from clustering and
the Markov property.
\end{decision}

\begin{firewall}[Claim boundary after V92]
\label{fire:v16v92}
The full Euclidean and diffeomorphism symmetries of an
SM--Einstein continuum have not been established.  V92 supplies only
the continuity theorem that would pass verified cutoff symmetries to
a subsequential physical-observable limit.
\end{firewall}

\section*{Version V92 append-only record}
\addcontentsline{toc}{section}{Version V92 append-only record}
The complete V91 source was retained byte for byte before this
continuation.  It had \(946512\) bytes, \(28575\) lines, and SHA--256
\[
 \texttt{D8468E1B5C70FFFBBC594A81EE500A3310EA1471F8966A5245011F72F46F409C}.
\]
V92 proved the uniform-modulus passage from countable rational labels
and lattice symmetries to continuous labels and full Euclidean
invariance.


\section{V93: locality, independence, and the Markov distinction}
\label{sec:v16v93}

\subsection{Local observable nets}

Let \({\cal R}\) be a countable base of bounded regions in Euclidean
spacetime.  For each \(O\in{\cal R}\), let
\({\cal A}_n(O)\) be the cutoff algebra generated by bounded physical
observables whose labels are supported in \(O\).  Assume isotony:
\[
 O_1\subset O_2
 \quad\Longrightarrow\quad
 {\cal A}_n(O_1)\subset{\cal A}_n(O_2).              \tag{93.1}
\]
Choose the V91 coordinate family so that the region attached to each
generator is part of its label, and let \({\cal A}(O)\) be the
limiting cylinder algebra generated by the same labels.

\begin{proposition}[Isotony survives the coordinate construction]
\label{prop:v16v93-isotony}
The limiting algebras obey
\[
 O_1\subset O_2
 \quad\Longrightarrow\quad
 {\cal A}(O_1)\subset{\cal A}(O_2).                 \tag{93.2}
\]
\end{proposition}

\begin{proof}
Every generator with label supported in \(O_1\) is also a generator
allowed in \(O_2\).  Taking finite sums and products preserves the
inclusion.
\end{proof}

This is an algebraic organization statement.  It supplies no decay
or conditional independence.

\subsection{Closed local identities at positive separation}

Let the label metric of V92 control both the observable and its
support.  Assume that if compact supports \(A,B\) obey
\({\rm dist}(A,B)>0\), they can be approximated by labels from the
countable dense set while retaining half that separation.

\begin{theorem}[Passage of separated bounded identities]
\label{thm:v16v93-separated}
Let \(\xi,\zeta\) have supports at distance \(\delta>0\), and let
\(\xi_j,\zeta_j\) be dense cutoff-compatible labels converging to
them with
\[
 {\rm dist}(\operatorname{supp}\xi_j,
             \operatorname{supp}\zeta_j)\geq\delta/2.           \tag{93.3}
\]
Suppose \(R:K_\xi\times K_\zeta\to\mathbb C\) is Lipschitz and the
cutoff identity
\[
 R(Y_{\xi_j,n},Y_{\zeta_j,n})=0
 \quad\mu_n\text{-almost surely}                    \tag{93.4}
\]
holds for every \(j,n\).  Under the V92 uniform modulus,
\[
 R(Y_\xi,Y_\zeta)=0
 \quad\nu\text{-almost surely}.                     \tag{93.5}
\]
\end{theorem}

\begin{proof}
For fixed \(j\), Proposition~\ref{prop:v16v91-relations} passes
(93.4) to the coordinate limit:
\[
 R(Y_{\xi_j},Y_{\zeta_j})=0
 \quad\nu\text{-almost surely}.                     \tag{93.6}
\]
The V92 construction gives
\[
 Y_{\xi_j}\to Y_\xi,\qquad
 Y_{\zeta_j}\to Y_\zeta
 \quad\text{in }L^1(\nu).                           \tag{93.7}
\]
The Lipschitz property of \(R\) makes
\[
 R(Y_{\xi_j},Y_{\zeta_j})
 \longrightarrow R(Y_\xi,Y_\zeta)
 \quad\text{in }L^1(\nu).                           \tag{93.8}
\]
The left side is zero for every \(j\), so the limit is zero in
\(L^1\), proving (93.5).
\end{proof}

The positive-separation assumption avoids contact terms.  Extending
an identity to touching supports requires distributional control of
the diagonal singularity and is not supplied by stochastic
continuity alone.

\subsection{Schwinger permutation relations}

For bosonic bounded coordinates, multiplication is already
commutative.  The useful cutoff statement is therefore symmetry of
the labelled Schwinger functions.  For parity-homogeneous graded
observables, the target relation is
\[
\begin{split}
 {\cal S}_n(
 A\,F_{\xi,n}F_{\zeta,n}\,B)
 =
 (-1)^{p_\xi p_\zeta}
 {\cal S}_n(
 A\,F_{\zeta,n}F_{\xi,n}\,B),                       \tag{93.9}
\end{split}
\]
when the two supports are separated and \(A,B\) denote other ordered
insertions.

\begin{proposition}[Permutation relations pass to the limit]
\label{prop:v16v93-permutation}
If all insertions in (93.9) are bounded, their joint cutoff laws
converge, and (93.9) holds for a dense separated label family, then
the limiting Schwinger functional obeys the same relation for every
positive-separated pair reached by the V92 modulus.
\end{proposition}

\begin{proof}
For dense labels, both sides of (93.9) are bounded continuous
cylinder expectations, so finite-coordinate weak convergence passes
the equality to the limit.  Approximate general separated labels as
in (93.3).  Repeated use of
\[
 \left|\prod_{a=1}^rz_a-\prod_{a=1}^rw_a\right|
 \leq\sum_{a=1}^r|z_a-w_a|                         \tag{93.10}
\]
for \(|z_a|,|w_a|\leq1\), followed by the V92 \(L^1\) modulus, passes
each limiting expectation to the desired labels.
\end{proof}

For fermionic insertions, (93.9) belongs to a graded generating
functional and cannot be encoded by an ordinary commutative
probability-coordinate algebra alone.

\subsection{Exact independence is stronger than locality}

\begin{theorem}[Independence is closed under joint weak convergence]
\label{thm:v16v93-independence}
Let \(X_n\) and \(Z_n\) take values in compact metric spaces and
suppose their joint laws converge to that of \((X,Z)\).  If
\[
 {\mathbb E}[f(X_n)g(Z_n)]
 =
 {\mathbb E}f(X_n)\,{\mathbb E}g(Z_n)               \tag{93.11}
\]
for every bounded continuous \(f,g\), then \(X\) and \(Z\) are
independent.
\end{theorem}

\begin{proof}
Weak convergence passes all three expectations in (93.11) to the
limit, giving
\[
 {\mathbb E}[f(X)g(Z)]
 =
 {\mathbb E}f(X)\,{\mathbb E}g(Z).                  \tag{93.12}
\]
Finite linear combinations of product functions are uniformly dense
in the continuous functions on the compact product by
Stone--Weierstrass.  Hence the joint law equals the product of its
marginals.
\end{proof}

An interacting Euclidean field is generally not independent on
arbitrary disjoint regions.  Algebraic locality means compatible
local algebras and, after OS reconstruction, graded commutation at
spacelike separation.  It does not mean (93.11).

\subsection{Clustering is an asymptotic estimate}

\begin{proposition}[Uniform clustering passes to the coordinate
limit]
\label{prop:v16v93-clustering-limit}
Let \(F_n,G_n\) be bounded physical cylinder observables whose joint
expectations converge, and suppose
\[
 |\operatorname{Cov}_{\mu_n}(F_n,G_n)|
 \leq C_{F,G}e^{-m\,d_n^{\rm phys}(F_n,G_n)}        \tag{93.13}
\]
with \(C_{F,G},m>0\) independent of \(n\).  If the physical support
distance tends to \(d(F,G)\), then
\[
 |\operatorname{Cov}_\nu(F,G)|
 \leq C_{F,G}e^{-m\,d(F,G)}.                        \tag{93.14}
\]
\end{proposition}

\begin{proof}
Finite-coordinate weak convergence passes the expectations of
\(F_nG_n\), \(F_n\), and \(G_n\) to the corresponding limiting
expectations.  Take the limit superior in (93.13) and use convergence
of the distances.
\end{proof}

This is the V85-to-continuum step for bounded coordinates.  It is
neither exact independence nor the Markov property.

\subsection{Conditional Markov structure is not weakly closed}

For random variables \(X,Y,Z\), the Markov relation across a boundary
variable \(Y\) is
\[
 X\ \perp\!\!\!\perp\ Z\mid Y.                     \tag{93.15}
\]

\begin{proposition}[A collapsing-boundary counterexample]
\label{prop:v16v93-markov-counter}
Conditional independence need not survive joint weak convergence.
\end{proposition}

\begin{proof}
Let \(U\) be a symmetric random sign and put
\[
 X_n=U,\qquad Y_n=U/n,\qquad Z_n=U.                 \tag{93.16}
\]
Given \(Y_n\), both \(X_n\) and \(Z_n\) are deterministic, so they
are conditionally independent.  Jointly,
\[
 (X_n,Y_n,Z_n)\Longrightarrow(U,0,U).               \tag{93.17}
\]
Conditioning on the constant limiting variable \(Y=0\) changes
nothing, while \(X=Z=U\) are not independent.  Thus the limiting
triple fails (93.15).
\end{proof}

The example shows that a boundary coordinate can lose the information
that mediated conditional independence.

\begin{theorem}[Kernel criterion for passage of a Markov relation]
\label{thm:v16v93-kernel}
Let \(X_n,Y_n,Z_n\) take values in compact metric spaces and suppose
\[
 {\mathbb P}_n(dx,dy,dz)
 =
 a_n(dx|y)b_n(dz|y)\lambda_n(dy).                   \tag{93.18}
\]
Assume \(\lambda_n\Rightarrow\lambda\), and that there are Feller
kernels \(a,b\) such that for every bounded continuous \(f,g\),
\[
\begin{split}
 \sup_y\left|
 \int f\,da_n(\cdot|y)-\int f\,da(\cdot|y)
 \right|&\longrightarrow0,\\
 \sup_y\left|
 \int g\,db_n(\cdot|y)-\int g\,db(\cdot|y)
 \right|&\longrightarrow0.                         \tag{93.19}
\end{split}
\]
Then the joint laws converge to
\[
 {\mathbb P}(dx,dy,dz)=a(dx|y)b(dz|y)\lambda(dy),   \tag{93.20}
\]
which satisfies \(X\perp\!\!\!\perp Z\mid Y\).
\end{theorem}

\begin{proof}
For a product test function \(f(x)h(y)g(z)\), (93.18) gives
\[
 \int h(y)a_nf(y)b_ng(y)\,d\lambda_n(y),            \tag{93.21}
\]
where \(a_nf(y)=\int f\,da_n(\cdot|y)\), and similarly for
\(b_ng\).  Uniform convergence in (93.19), the Feller continuity of
\(af,bg\), and weak convergence of \(\lambda_n\) pass (93.21) to
\[
 \int h(y)af(y)bg(y)\,d\lambda(y).                  \tag{93.22}
\]
Finite sums of product test functions are dense in the continuous
functions on the compact triple product.  Hence (93.20) is the weak
limit.  Its product-kernel form is exactly conditional
independence.
\end{proof}

For an SM--Einstein domain Markov property, the boundary variable
must contain the complete gauge-covariant and gravitational boundary
sector, and its conditional kernels must satisfy a stability theorem
of the form (93.19).

\subsection{Relation to OS locality}

The standard OS reconstruction of local relativistic fields uses a
package: Euclidean covariance, reflection positivity, permutation or
graded symmetry, regularity, and analytic continuation.  V87, V92,
and Proposition~\ref{prop:v16v93-permutation} address closure of
three parts of that package.  They do not by themselves prove
microcausality, because the regularity and analytic-continuation
hypotheses have not been verified.

\begin{assessment}[What V93 closes]
\label{ass:v16v93-status}
V93 proves passage of bounded separated identities, graded
permutation relations, exact independence, and uniform clustering
under their respective hypotheses.  It demonstrates that
conditional Markov structure is not weakly closed and gives a
uniform-kernel criterion that preserves it.

Local net structure, clustering, independence, the Markov property,
and reconstructed Lorentzian locality are distinct.  The full
SM--Einstein programme still lacks the graded fermionic functional,
boundary-kernel stability, and OS regularity.
\end{assessment}

\begin{decision}[V94: uniqueness of the continuum limit]
\label{dec:v16v93-next}
Address subsequence dependence.  Prove that convergence of a
measure-determining physical cylinder algebra to one set of values
forces uniqueness of the coordinate-law limit and convergence of the
whole cutoff sequence.  Then formulate a boundary-influence
criterion that can make the limiting values independent of finite
volume and boundary condition.
\end{decision}

\begin{firewall}[Claim boundary after V93]
\label{fire:v16v93}
No locality or Markov theorem has been verified for the actual
SM--Einstein cutoff measures.  V93 supplies the closure criteria and
prevents factorization, clustering, and locality from being
identified with one another.
\end{firewall}

\section*{Version V93 append-only record}
\addcontentsline{toc}{section}{Version V93 append-only record}
The complete V92 source was retained byte for byte before this
continuation.  It had \(957609\) bytes, \(28909\) lines, and SHA--256
\[
 \texttt{778D73A8C21F504398D39DFFD5E6F05C05007F9B06F5D461D510D0261B4036FB}.
\]
V93 proved distinct closure theorems for local identities,
clustering, independence, and stable conditional kernels.


\section{V94: uniqueness of the physical-observable limit}
\label{sec:v16v94}

\subsection{A measure-determining cylinder algebra}

Retain the compact countable product
\[
 K=\prod_{k=1}^\infty K_k                           \tag{94.1}
\]
of V91, and let \({\cal C}_{\rm cyl}(K)\) denote its continuous
cylinder algebra.

\begin{lemma}[Cylinder functions determine probability measures]
\label{lem:v16v94-determining}
If probability measures \(\nu,\widetilde\nu\) on \(K\) satisfy
\[
 \int F\,d\nu=\int F\,d\widetilde\nu
 \quad(F\in{\cal C}_{\rm cyl}(K)),                  \tag{94.2}
\]
then \(\nu=\widetilde\nu\).
\end{lemma}

\begin{proof}
By Theorem~\ref{thm:v16v91-positive-functional}, the cylinder algebra
is uniformly dense in \(C(K)\).  If \(G\in C(K)\), choose cylinder
functions \(F_j\) with \(\|F_j-G\|_\infty\to0\).  Equation (94.2)
passes to \(G\), so the two measures agree on every continuous
function.  The Riesz representation theorem gives
\(\nu=\widetilde\nu\).
\end{proof}

\begin{theorem}[Uniqueness upgrades subsequential to full convergence]
\label{thm:v16v94-full-convergence}
Let \(\nu_n\) be probability measures on \(K\).  Suppose that for
every \(F\) in a uniformly dense subalgebra
\({\cal D}\subset{\cal C}_{\rm cyl}(K)\), the scalar limit
\[
 L(F):=\lim_{n\to\infty}\int F\,d\nu_n              \tag{94.3}
\]
exists.  Then there is a unique probability measure \(\nu\) on
\(K\) with
\[
 L(F)=\int F\,d\nu,
\]
and
\[
 \nu_n\Rightarrow\nu                               \tag{94.4}
\]
without extraction of a subsequence.
\end{theorem}

\begin{proof}
Compactness of \(K\) makes \(\{\nu_n\}\) tight.  Any subsequence has
a further subsequence converging weakly to some \(\nu'\).
For \(F\in{\cal D}\), weak convergence and (94.3) give
\[
 \int F\,d\nu'=L(F).                                \tag{94.5}
\]
Uniform density and Lemma~\ref{lem:v16v94-determining} show that all
such subsequential limits coincide; call the common measure \(\nu\).
If (94.4) failed, some subsequence would remain outside a weak
neighborhood of \(\nu\), but compactness would give it a further
subsequence converging to a different limit.  This contradiction
proves full convergence.
\end{proof}

\begin{corollary}[Finite-dimensional distributions suffice]
\label{cor:v16v94-fdds}
If every finite-coordinate marginal of \(\nu_n\) converges and the
limits are consistent, then \(\nu_n\Rightarrow\nu\) on \(K\).
\end{corollary}

\begin{proof}
The marginal convergence gives (94.3) for every continuous cylinder
function.  Apply Theorem~\ref{thm:v16v94-full-convergence}.
\end{proof}

\subsection{Boundary influence through the localized covariance}

Let \(\Lambda\) be a finite block region, \(\tau\) a boundary
condition, and
\[
 d\mu_\Lambda^\tau(x)
 =
 Z_\tau^{-1}e^{-U_\Lambda(x;\tau)}\,dx.             \tag{94.6}
\]
Suppose two boundary conditions are joined by a \(C^1\) path
\(\tau_s\), \(0\leq s\leq1\), and put
\[
 W_s(x)=\partial_sU_\Lambda(x;\tau_s).              \tag{94.7}
\]

\begin{lemma}[Boundary interpolation identity]
\label{lem:v16v94-interpolation}
If \(F\) is independent of \(s\) and differentiation under the
integral is valid, then
\[
 \frac d{ds}{\mathbb E}_{\mu_\Lambda^{\tau_s}}F
 =
 -\operatorname{Cov}_{\mu_\Lambda^{\tau_s}}(F,W_s).
                                                               \tag{94.8}
\]
\end{lemma}

\begin{proof}
Differentiate the numerator and denominator in
\[
 {\mathbb E}_{\mu_\Lambda^{\tau_s}}F
 =
 \frac{\int Fe^{-U_\Lambda(x;\tau_s)}dx}
 {\int e^{-U_\Lambda(x;\tau_s)}dx}.
\]
The numerator contributes \(-{\mathbb E}(FW_s)\), and the derivative
of the reciprocal partition function contributes
\({\mathbb E}F\,{\mathbb E}W_s\).  Their sum is (94.8).
\end{proof}

Let \(S_F\) be the block-gradient support of \(F\), and suppose the
gradient support of every \(W_s\) lies in a boundary layer
\(T_\Lambda\).

\begin{theorem}[Exponential boundary-insensitivity bound]
\label{thm:v16v94-boundary}
Assume uniformly along the path \(\tau_s\) that the V84 hypotheses
hold with
\[
 \rho_{\min,s}\geq\rho_*>0,\qquad
 \theta_s\leq\theta_*<1,\qquad
 R_{0,s}\leq R_* .                                  \tag{94.9}
\]
Define
\[
 G_s(H)=\sum_i
 \|\nabla_iH\|_{L^2(\mu_\Lambda^{\tau_s})}
\]
and assume
\[
 \int_0^1G_s(W_s)\,ds\leq B_{\partial,\Lambda}.     \tag{94.10}
\]
Then
\[
\begin{split}
 \left|
 {\mathbb E}_{\mu_\Lambda^{\tau_1}}F
 -
 {\mathbb E}_{\mu_\Lambda^{\tau_0}}F
 \right|
 \leq&
 \frac{
 \sup_sG_s(F)\,B_{\partial,\Lambda}
 }{\rho_*(1-\theta_*)}\\
 &\times
 \theta_*^{
 \lceil{\rm dist}(S_F,T_\Lambda)/R_*\rceil}.
                                                               \tag{94.11}
\end{split}
\]
\end{theorem}

\begin{proof}
By Lemma~\ref{lem:v16v94-interpolation} and the V84 localized
covariance estimate,
\[
\begin{split}
 \left|\frac d{ds}{\mathbb E}_{\mu_\Lambda^{\tau_s}}F\right|
 &\leq
 \frac{G_s(F)G_s(W_s)}{\rho_*(1-\theta_*)}
 \theta_*^{
 \lceil{\rm dist}(S_F,T_\Lambda)/R_*\rceil}.
                                                               \tag{94.12}
\end{split}
\]
Integrate from \(s=0\) to \(s=1\) and use (94.10).
\end{proof}

If \(B_{\partial,\Lambda}\) grows polynomially in the boundary size,
the exponential factor still wins when a fixed local observable is
moved arbitrarily far from the boundary.  If it grows exponentially,
its rate must be compared explicitly with the influence rate.

\subsection{Uniqueness of an infinite-volume Gibbs state}

\begin{theorem}[DLR uniqueness from vanishing boundary oscillation]
\label{thm:v16v94-dlr}
Fix a cutoff and suppose infinite-volume Gibbs measures satisfy the
DLR equation.  Assume that for every bounded local \(F\), supported
in \(S_F\), there are regions \(\Lambda_j\uparrow\Gamma\) such that
\[
 \varepsilon_j(F):=
 \sup_{\tau,\tau'}
 \left|
 {\mathbb E}_{\mu_{\Lambda_j}^{\tau}}F
 -
 {\mathbb E}_{\mu_{\Lambda_j}^{\tau'}}F
 \right|
 \longrightarrow0.                                 \tag{94.13}
\]
Then there is at most one infinite-volume Gibbs measure.
\end{theorem}

\begin{proof}
Let \(\mu\) and \(\widetilde\mu\) be two DLR measures.  Their
conditional expectations of \(F\), given the exterior of
\(\Lambda_j\), are the finite-volume expectations with the realized
exterior boundary conditions.  Therefore
\[
\begin{split}
 |\mu(F)-\widetilde\mu(F)|
 &\leq
 \iint
 \left|
 {\mathbb E}_{\mu_{\Lambda_j}^{\tau}}F
 -
 {\mathbb E}_{\mu_{\Lambda_j}^{\tau'}}F
 \right|
 d\pi_j(\tau,\tau')\\
 &\leq\varepsilon_j(F),                             \tag{94.14}
\end{split}
\]
where \(\pi_j\) is any coupling of the two exterior marginals.
Letting \(j\to\infty\) gives equality on every bounded local
observable.  Such observables determine the Gibbs measure.
\end{proof}

\begin{corollary}[A sufficient influence corridor for DLR uniqueness]
\label{cor:v16v94-dlr-influence}
If all boundary pairs admit interpolation paths satisfying
Theorem~\ref{thm:v16v94-boundary} and
\[
 \sup_{\tau,\tau'}
 \sup_sG_s(F)\,B_{\partial,\Lambda_j}
 \theta_*^{
 \lceil{\rm dist}(S_F,T_{\Lambda_j})/R_*\rceil}
 \longrightarrow0,                                 \tag{94.15}
\]
then (94.13) holds and the infinite-volume Gibbs state is unique.
\end{corollary}

\begin{proof}
Theorem~\ref{thm:v16v94-boundary} bounds the supremum in (94.13) by
the expression in (94.15) divided by
\(\rho_*(1-\theta_*)\).  Apply
Theorem~\ref{thm:v16v94-dlr}.
\end{proof}

\subsection{Two kinds of uniqueness}

The DLR theorem removes dependence on volume boundary conditions at
a fixed regulator.  It does not remove dependence on the
renormalization trajectory or on ultraviolet subsequences.  The
complete uniqueness ledger is
\[
\begin{array}{c|c}
\text{thermodynamic uniqueness}&
\text{vanishing boundary oscillation (94.13)}\\
\text{continuum uniqueness}&
\text{unique limits on a determining algebra (94.3)}.
\end{array}                                         \tag{94.16}
\]
Both are needed if volume and cutoff are removed simultaneously.

\begin{proposition}[Inheritance by a unique limit]
\label{prop:v16v94-inheritance}
If the full sequence \(\nu_n\Rightarrow\nu\) by
Theorem~\ref{thm:v16v94-full-convergence}, then every property proved
by the closure theorems V87, V92, and V93 holds for this unique
\(\nu\), provided their uniform hypotheses hold for the full
sequence.
\end{proposition}

\begin{proof}
Those closure proofs use only convergence of bounded cylinder
expectations and their stated uniform estimates.  Full convergence
supplies the same limits without choosing a subsequence.
\end{proof}

\subsection{Programme insertion}

For each determining cylinder observable \(F\), record:
\[
\begin{array}{c|c}
\text{UV Cauchy modulus}&
\sup_{m,n\geq N}|{\mathbb E}_{\nu_m}F-
                    {\mathbb E}_{\nu_n}F|\\
\text{boundary oscillation}&
\sup_{\tau,\tau'}|{\mathbb E}_{\mu_\Lambda^\tau}F-
                    {\mathbb E}_{\mu_\Lambda^{\tau'}}F|.
\end{array}                                         \tag{94.17}
\]
The first must vanish as \(N\to\infty\); the second must vanish as
the boundary recedes.  Compactness then upgrades these scalar
ledgers to uniqueness of the corresponding probability law.

\begin{assessment}[What V94 closes]
\label{ass:v16v94-status}
V94 proves that convergence on a measure-determining physical
cylinder algebra yields a unique continuum coordinate law and full,
rather than subsequential, convergence.  The localized covariance
theorem gives an explicit boundary-insensitivity estimate, which
implies DLR uniqueness when its boundary prefactor loses to the
influence decay.

Neither the UV Cauchy ledger nor the uniform boundary path estimate
has been verified for the SM--Einstein cutoff laws.  Thermodynamic
and continuum uniqueness remain separate open gates.
\end{assessment}

\begin{decision}[V95: companion audit and Cauchy comparison]
\label{dec:v16v94-next}
Perform the mandatory YM/NS audit.  Then derive a same-space
interpolation identity comparing two nearby cutoff actions on a
common refined configuration space.  Use the localized covariance
bound to obtain a quantitative Cauchy estimate for physical
observables, and identify the renormalization-counterterm mismatch
that must tend to zero.
\end{decision}

\begin{firewall}[Claim boundary after V94]
\label{fire:v16v94}
No unique SM--Einstein continuum law has been constructed.  V94
proves the determining-algebra and boundary-influence criteria that
would establish uniqueness once the required uniform cutoff
comparisons are available.
\end{firewall}

\section*{Version V94 append-only record}
\addcontentsline{toc}{section}{Version V94 append-only record}
The complete V93 source was retained byte for byte before this
continuation.  It had \(969154\) bytes, \(29243\) lines, and SHA--256
\[
 \texttt{ADE1270FA74AACAF50DF765393414D0613FB956D7B0E4E3659F3EB542CB711A2}.
\]
V94 proved uniqueness from a determining algebra and derived an
explicit boundary-influence route to DLR uniqueness.


\section{V95: companion audit and adjacent-cutoff comparison}
\label{sec:v16v95}

\subsection{Mandatory companion audit}

The newest companion files remain
\[
\begin{array}{lll}
\text{YM V64:}&14831\text{ lines},&
\texttt{69B9A7675E3F46414962EF6CD74AAD10E17E1EEA975C543B78419885D2B34808},\\
\text{NS V26:}&5319\text{ lines},&
\texttt{82C887F8C2C69E4F28FAD7C3DC8DE5B9099CB5A4BF431EBA1FFF3E91935978AD}.
\end{array}                                                   \tag{95.1}
\]
YM V64 still supplies a centre-specific bridge repair rather than a
global continuum action comparison; its terminal third-coordinate
layer and continuum gates remain open.  NS V26 remains a
channel-sensitive kernel calculation with no renormalization-group
input.  Neither changes the V94 decision.

\begin{theorem}[V95 audit decision]
\label{thm:v16v95-audit}
No cutoff-Cauchy estimate can be imported from the companion
programmes.  The next required object is a common-refinement action
interpolation whose defect is measured through the localized
covariance resolvent.
\end{theorem}

\begin{proof}
The audited YM theorem compares symbol charts at one regulator, and
the NS theorem compares tensor norms in a classical kernel.  Neither
compares two probability laws at adjacent ultraviolet cutoffs.
\end{proof}

\subsection{The same-space interpolation identity}

Let \(\lambda\) be a common finite-dimensional reference measure and
let
\[
 d\mu_s=Z_s^{-1}e^{-U_s}\,d\lambda,
 \qquad 0\leq s\leq1,                               \tag{95.2}
\]
where \(U_s\) is \(C^1\) in \(s\).  Put
\[
 \dot U_s=\partial_sU_s.                            \tag{95.3}
\]

\begin{lemma}[Action interpolation]
\label{lem:v16v95-action}
For an \(s\)-independent integrable observable \(F\), if
differentiation under the integral is valid, then
\[
 \frac d{ds}{\mathbb E}_{\mu_s}F
 =
 -\operatorname{Cov}_{\mu_s}(F,\dot U_s).           \tag{95.4}
\]
\end{lemma}

\begin{proof}
This is the numerator/partition-function calculation of
Lemma~\ref{lem:v16v94-interpolation}, with the boundary path replaced
by an action path.
\end{proof}

For each \(s\), define
\[
 b_{i,s}(H)=\|\nabla_iH\|_{L^2(\mu_s)}.              \tag{95.5}
\]

\begin{theorem}[Localized comparison of two Gibbs laws]
\label{thm:v16v95-local-comparison}
Assume every \(\mu_s\) satisfies the V84 hypotheses with
\[
 \rho_{\min,s}\geq\rho_*>0,\qquad
 \theta_s\leq\theta_*<1,\qquad
 R_{0,s}\leq R_* .                                  \tag{95.6}
\]
If the gradient support of \(F\) is \(S_F\), then
\[
\begin{split}
 |{\mathbb E}_{\mu_1}F-{\mathbb E}_{\mu_0}F|
 \leq&
 \frac1{\rho_*(1-\theta_*)}
 \int_0^1
 \left(\sum_{i\in S_F}b_{i,s}(F)\right)\\
 &\quad\times
 \left(
 \sum_j
 \theta_*^{
 \lceil{\rm dist}(S_F,j)/R_*\rceil}
 b_{j,s}(\dot U_s)
 \right)ds.                                        \tag{95.7}
\end{split}
\]
\end{theorem}

\begin{proof}
Equations (95.4) and (84.16) give
\[
 \left|\frac d{ds}{\mathbb E}_{\mu_s}F\right|
 \leq
 \sum_{i\in S_F}\sum_j
 b_{i,s}(F)({\cal A}_s^{-1})_{ij}
 b_{j,s}(\dot U_s).                                 \tag{95.8}
\]
As in Theorem~\ref{thm:v16v94-boundary}, (95.6) and
Proposition~\ref{prop:v16v84-inverse-decay} imply
\[
 ({\cal A}_s^{-1})_{ij}
 \leq
 \frac{
 \theta_*^{\lceil{\rm dist}(i,j)/R_*\rceil}
 }{\rho_*(1-\theta_*)}.                             \tag{95.9}
\]
Since \(i\in S_F\),
\({\rm dist}(i,j)\geq{\rm dist}(S_F,j)\), and
\(\theta_*\in(0,1)\) reverses this into the required upper bound.
Integrate (95.8) over \(s\).
\end{proof}

Only the gradient of the action mismatch appears.  An
\(x\)-independent normalization counterterm drops out exactly.

\subsection{Physical form of the defect}

Let \(\ell\) be the common-refinement block scale and set
\[
 m_*=\frac{-\log\theta_*}{\ell R_*}.                \tag{95.10}
\]
Define the observable-weighted action defect
\[
\begin{split}
 {\mathfrak D}(S_F;U_\bullet)
 :=
 \int_0^1\sum_j
 e^{-m_*d^{\rm phys}(S_F,j)}
 b_{j,s}(\dot U_s)\,ds.                             \tag{95.11}
\end{split}
\]

\begin{corollary}[Physical weighted comparison]
\label{cor:v16v95-physical-comparison}
Under Theorem~\ref{thm:v16v95-local-comparison},
\[
 |{\mathbb E}_{\mu_1}F-{\mathbb E}_{\mu_0}F|
 \leq
 \frac{\sup_sG_s(F)}
 {\rho_*(1-\theta_*)}
 {\mathfrak D}(S_F;U_\bullet),                      \tag{95.12}
\]
where \(G_s(F)=\sum_{i\in S_F}b_{i,s}(F)\).
\end{corollary}

\begin{proof}
The ceiling estimate (85.6), with the common parameters in (95.10),
gives
\[
 \theta_*^{\lceil{\rm dist}(S_F,j)/R_*\rceil}
 \leq e^{-m_*d^{\rm phys}(S_F,j)}.
\]
Insert this in (95.7) and take the supremum of \(G_s(F)\).
\end{proof}

\subsection{Putting adjacent cutoffs on one space}

Let \(x_n\) denote the coarse configuration.  Introduce auxiliary
fine variables \(\eta_n\) through a probability kernel
\(q_n(d\eta_n|x_n)\), and lift the coarse law to
\[
 d\widetilde\mu_n(x_n,\eta_n)
 =
 d\mu_n(x_n)q_n(d\eta_n|x_n).                       \tag{95.13}
\]
Then
\[
 {\mathbb E}_{\widetilde\mu_n}\widehat F_n
 =
 {\mathbb E}_{\mu_n}F_n                             \tag{95.14}
\]
whenever \(\widehat F_n(x_n,\eta_n)=F_n(x_n)\).
Write the fine law \(\mu_{n+1}\) and
\(\widetilde\mu_n\) relative to a common reference measure on this
refined space, with actions \(U_{n+1}\) and
\(\widetilde U_n\).

\begin{theorem}[Adjacent-cutoff observable estimate]
\label{thm:v16v95-adjacent}
Let \(F_{n+1}\) be the fine approximant and
\(\widehat F_n\) the lifted coarse approximant.  Suppose
\[
 \varepsilon_n^{\rm obs}(F)
 :=
 {\mathbb E}_{\mu_{n+1}}
 |F_{n+1}-\widehat F_n|<\infty.                     \tag{95.15}
\]
Interpolate
\[
 U_{n,s}=(1-s)\widetilde U_n+sU_{n+1}.              \tag{95.16}
\]
If this path satisfies Corollary~\ref{cor:v16v95-physical-comparison},
then
\[
\begin{split}
 |{\mathbb E}_{\mu_{n+1}}F_{n+1}
  -{\mathbb E}_{\mu_n}F_n|
 \leq&
 \varepsilon_n^{\rm obs}(F)\\
 &+
 \frac{\sup_sG_{n,s}(\widehat F_n)}
 {\rho_{*,n}(1-\theta_{*,n})}
 {\mathfrak D}_n(S_F).                              \tag{95.17}
\end{split}
\]
\end{theorem}

\begin{proof}
Add and subtract
\({\mathbb E}_{\mu_{n+1}}\widehat F_n\).  The first difference is
bounded by (95.15).  Equation (95.14) identifies the other endpoint,
and Corollary~\ref{cor:v16v95-physical-comparison} bounds the
difference between the two action laws.
\end{proof}

\subsection{The summable Cauchy gate}

Define the right side of (95.17) by
\[
 \varepsilon_n(F).                                  \tag{95.18}
\]

\begin{theorem}[Summable adjacent errors give a unique scalar limit]
\label{thm:v16v95-summable}
If
\[
 \sum_{n=1}^\infty\varepsilon_n(F)<\infty,          \tag{95.19}
\]
then \(\{{\mathbb E}_{\mu_n}F_n\}\) is Cauchy and hence convergent.
If (95.19) holds on a measure-determining cylinder algebra, the
physical coordinate laws converge to a unique continuum law by V94.
\end{theorem}

\begin{proof}
For \(m>n\), telescope (95.17):
\[
 |{\mathbb E}_{\mu_m}F_m-
   {\mathbb E}_{\mu_n}F_n|
 \leq\sum_{k=n}^{m-1}\varepsilon_k(F).              \tag{95.20}
\]
The tail tends to zero uniformly in \(m>n\) by (95.19).  The second
claim is Theorem~\ref{thm:v16v94-full-convergence}.
\end{proof}

Merely proving \(\varepsilon_n(F)\to0\) is insufficient.  For
example, increments \(1/n\) vanish but are not summable.

\subsection{Counterterm mismatch ledger}

Suppose along the common-space interpolation,
\[
 \dot U_{n,s}
 =
 \sum_{a=1}^A\dot g_{a,n}(s){\cal O}_{a,n}
 +{\cal R}_{n,s}+c_n(s),                            \tag{95.21}
\]
where \(c_n(s)\) is configuration independent.

\begin{proposition}[Linear defect decomposition]
\label{prop:v16v95-defect-decomp}
The weighted defect obeys
\[
\begin{split}
 {\mathfrak D}_n(S_F)
 \leq&
 \sum_{a=1}^A
 \int_0^1|\dot g_{a,n}(s)|\,ds\,
 {\mathfrak N}_{a,n}(S_F)\\
 &+
 {\mathfrak R}_n(S_F),                              \tag{95.22}
\end{split}
\]
where
\[
\begin{split}
 {\mathfrak N}_{a,n}(S)
 &:=
 \sup_s\sum_j e^{-m_{*,n}d^{\rm phys}(S,j)}
 b_{j,n,s}({\cal O}_{a,n}),\\
 {\mathfrak R}_n(S)
 &:=
 \int_0^1\sum_j e^{-m_{*,n}d^{\rm phys}(S,j)}
 b_{j,n,s}({\cal R}_{n,s})\,ds.                    \tag{95.23}
\end{split}
\]
The scalar term \(c_n(s)\) contributes zero.
\end{proposition}

\begin{proof}
Differentiate (95.21) in each block, use the triangle inequality in
\(L^2(\mu_{n,s})\), multiply by the nonnegative exponential weight,
sum over \(j\), and integrate over \(s\).
\(\nabla_jc_n(s)=0\) for every \(j\).
\end{proof}

Relevant and marginal couplings must be tuned so that their weighted
increments in (95.22), together with the observable discretization
error, are summable.  The remainder must satisfy the same condition.
Calling \({\cal R}_{n,s}\) irrelevant does not prove this estimate.

\subsection{Programme insertion}

For every determining physical observable \(F\), add the adjacent
cutoff ledger
\[
 \varepsilon_n^{\rm obs}(F)
 +
 \frac{\sup_sG_{n,s}(\widehat F_n)}
 {\rho_{*,n}(1-\theta_{*,n})}
 \left[
 \sum_a\|\dot g_{a,n}\|_{L^1_s}{\mathfrak N}_{a,n}(S_F)
 +{\mathfrak R}_n(S_F)
 \right].                                          \tag{95.24}
\]
Its sum over \(n\) must be finite.

\begin{assessment}[What V95 closes]
\label{ass:v16v95-status}
V95 derives an exact adjacent-cutoff comparison on a common refined
space.  The action mismatch enters only through an exponentially
localized block-gradient norm; additive normalization constants
cancel.  Summability of this defect and the observable interpolation
error yields convergence on the determining algebra and hence a
unique coordinate-law limit.

No common-refinement lift, uniformly admissible interpolation path,
or summable counterterm ledger has yet been verified for the full
SM--Einstein action.
\end{assessment}

\begin{decision}[V96: path stability and reference-measure changes]
\label{dec:v16v95-next}
Continue after the audit.  Generalize the comparison to adjacent
cutoffs whose Gaussian or Haar reference measures differ.  Derive
the Radon--Nikodym action correction when absolute continuity holds,
and prove a singularity obstruction showing when no same-space
interpolation is possible.  This decides whether the V95 route can
be used directly or must be replaced by a coupling comparison.
\end{decision}

\begin{firewall}[Claim boundary after V95]
\label{fire:v16v95}
The cutoff sequence has not been proved Cauchy.  V95 supplies the
comparison identity and the summable defect criterion; all
renormalization and common-space hypotheses remain open.
\end{firewall}

\section*{Version V95 append-only record}
\addcontentsline{toc}{section}{Version V95 append-only record}
The complete V94 source was retained byte for byte before this
continuation.  It had \(980075\) bytes, \(29579\) lines, and SHA--256
\[
 \texttt{F217364EE0068BA9A257A8D0C264D54BF4FA0E3E83E5C36E370BAF679B325BF1}.
\]
V95 performed the required companion audit and established the
localized, summable adjacent-cutoff comparison criterion.


\section{V96: reference-measure changes and singular cutoff lifts}
\label{sec:v16v96}

\subsection{Absorbing an absolutely continuous reference change}

Let
\[
 d\mu_a=Z_a^{-1}e^{-U_a}\,d\lambda_a,
 \qquad a=0,1,                                      \tag{96.1}
\]
be two laws on the same measurable configuration space.

\begin{lemma}[Radon--Nikodym action correction]
\label{lem:v16v96-rn}
Suppose \(\lambda_1\ll\lambda_0\) and
\[
 r=\frac{d\lambda_1}{d\lambda_0}>0
 \quad\lambda_0\text{-almost everywhere}.           \tag{96.2}
\]
Then \(\mu_1\) may be written relative to \(\lambda_0\) as
\[
 d\mu_1=Z_1^{-1}e^{-\overline U_1}\,d\lambda_0,
 \qquad
 \overline U_1=U_1-\log r.                          \tag{96.3}
\]
Consequently the effective mismatch for the V95 interpolation is
\[
 \Delta U_{\rm eff}
 =
 U_1-U_0-\log r.                                    \tag{96.4}
\]
\end{lemma}

\begin{proof}
Substitute \(d\lambda_1=r\,d\lambda_0\) in (96.1):
\[
 e^{-U_1}d\lambda_1
 =
 e^{-U_1+\log r}d\lambda_0
 =
 e^{-\overline U_1}d\lambda_0.
\]
The normalizing integral is unchanged.
\end{proof}

The reference Jacobian is part of the action defect.  Omitting it
compares the wrong probability laws.

\subsection{Finite-dimensional Gaussian references}

Let \(\gamma_{C_a}\) be a centered nondegenerate Gaussian on
\(\mathbb R^N\) with covariance \(C_a\succ0\).

\begin{proposition}[Gaussian correction]
\label{prop:v16v96-gaussian}
The Gaussian references are equivalent and
\[
\begin{split}
 \log\frac{d\gamma_{C_1}}{d\gamma_{C_0}}(x)
 &=
 \frac12\log\frac{\det C_0}{\det C_1}
 -\frac12\langle x,(C_1^{-1}-C_0^{-1})x\rangle.
                                                               \tag{96.5}
\end{split}
\]
Thus, up to an irrelevant scalar,
\[
 \Delta U_{\rm eff}(x)
 =
 U_1(x)-U_0(x)
 +\frac12\langle x,(C_1^{-1}-C_0^{-1})x\rangle.     \tag{96.6}
\]
\end{proposition}

\begin{proof}
Divide the two Gaussian densities
\[
 (2\pi)^{-N/2}(\det C_a)^{-1/2}
 \exp\left(-\frac12\langle x,C_a^{-1}x\rangle\right).
\]
This gives (96.5), and Lemma~\ref{lem:v16v96-rn} gives (96.6).
\end{proof}

The quadratic correction is generally unbounded.  Its contribution
must be estimated through the weighted gradient norm (95.11) or
through a direct Hessian argument; a bounded-density perturbation
lemma alone cannot control it.

\subsection{Haar references and refinement variables}

\begin{proposition}[Product-Haar refinement]
\label{prop:v16v96-haar}
Let a compact gauge group \(G\) carry normalized Haar measure \(dg\).
If the fine edge set is a disjoint union
\[
 E_{n+1}=E_n\sqcup E_n^{\rm new},                   \tag{96.7}
\]
then the lifted coarse reference
\[
 \left(\prod_{e\in E_n}dg_e\right)
 \left(\prod_{e\in E_n^{\rm new}}dg_e\right)        \tag{96.8}
\]
is exactly the fine product-Haar reference.  No Radon--Nikodym
correction is present; all mismatch lies in the lifted action.
\end{proposition}

\begin{proof}
Both sides of (96.8) are the normalized product of one Haar measure
for every edge in \(E_{n+1}\).  Equality follows from the product
construction.
\end{proof}

For subdivision maps in which a coarse holonomy is a product of fine
holonomies, one must choose a conditional Haar kernel whose
pushforward gives the coarse Haar law.  Deterministically fixing the
new edges does not have this property and creates a singular lift.

\subsection{Stability of conditional gaps under bounded corrections}

\begin{lemma}[Holley--Stroock comparison for a Poincar\'e gap]
\label{lem:v16v96-hs}
Suppose \(\mu\) satisfies
\[
 \operatorname{Var}_\mu(f)
 \leq\frac1\rho\int|\nabla f|^2\,d\mu.              \tag{96.9}
\]
Let
\[
 d\nu=\frac{e^{-V}}{\int e^{-V}d\mu}\,d\mu
\]
and assume
\[
 \operatorname{osc}V:=\sup V-\inf V\leq\omega<\infty.
                                                               \tag{96.10}
\]
Then
\[
 \operatorname{Var}_\nu(f)
 \leq\frac{e^\omega}{\rho}
 \int|\nabla f|^2\,d\nu.                            \tag{96.11}
\]
\end{lemma}

\begin{proof}
Let \(r=d\nu/d\mu\), \(m=\inf r\), and \(M=\sup r\).
Then \(M/m\leq e^\omega\).  Using the variational definition of
variance,
\[
\begin{split}
 \operatorname{Var}_\nu(f)
 &\leq\int(f-{\mathbb E}_\mu f)^2\,d\nu\\
 &\leq M\operatorname{Var}_\mu(f)
 \leq\frac M\rho\int|\nabla f|^2\,d\mu
 \leq\frac M{m\rho}\int|\nabla f|^2\,d\nu.
                                                               \tag{96.12}
\end{split}
\]
\end{proof}

\begin{corollary}[Pathwise conditional-gap floor]
\label{cor:v16v96-path-gaps}
Suppose the conditional block law at \(s=0\) has gap \(\rho_i(0)\),
and the change of its conditional potential along the V95 path has
blockwise oscillation at most \(\omega_i\).  Then
\[
 \rho_i(s)\geq e^{-\omega_i}\rho_i(0)
 \quad(0\leq s\leq1).                               \tag{96.13}
\]
\end{corollary}

\begin{proof}
Apply Lemma~\ref{lem:v16v96-hs} to each conditional law, with the
exterior configuration fixed, and take the uniform oscillation
bound.
\end{proof}

The mixed-Hessian influence entries must still be recomputed along
the path.  Equation (96.13) controls only the diagonal conditional
floor.  If the reference correction is quadratic and unbounded,
\(\omega_i=\infty\), so direct Gaussian or Schur analysis is needed.

\subsection{A singularity obstruction}

\begin{theorem}[No finite action correction across singular
references]
\label{thm:v16v96-singular}
If \(\lambda_1\) is not absolutely continuous with respect to
\(\lambda_0\), there is no finite measurable function \(V\) and
constant \(Z\in(0,\infty)\) such that
\[
 d\lambda_1=Z^{-1}e^{-V}\,d\lambda_0.               \tag{96.14}
\]
In particular, the straight action interpolation of V95 cannot
compare the two references on the \(\lambda_0\) space.
\end{theorem}

\begin{proof}
Every measure of the form on the right of (96.14) is absolutely
continuous with respect to \(\lambda_0\).  This contradicts the
hypothesis.
\end{proof}

\begin{corollary}[Deterministic new-mode lifts are singular]
\label{cor:v16v96-deterministic}
Suppose the fine law has a density in a new variable
\(\eta\in\mathbb R^r\), whereas the lifted coarse law fixes
\(\eta=\eta_0\) almost surely.  Then the two laws are mutually
singular in the \(\eta\) coordinate, and no V95 action interpolation
exists between them.
\end{corollary}

\begin{proof}
The hyperplane \(\{\eta=\eta_0\}\) has full lifted-coarse measure and
zero fine measure, because the fine \(\eta\)-marginal is absolutely
continuous with respect to \(r\)-dimensional Lebesgue measure.
\end{proof}

This is why the auxiliary kernel \(q_n\) in (95.13) must have
appropriate full support.

\subsection{Coupling comparison when densities are unavailable}

Let \(X=(X_i)\) and \(Y=(Y_i)\) live in possibly different
coordinate realizations with block metrics \(d_i\).  Let \(\pi\) be
a coupling of their laws and put
\[
 a_i={\mathbb E}_\pi d_i(X_i,Y_i).                  \tag{96.15}
\]

\begin{theorem}[Coupled influence comparison]
\label{thm:v16v96-coupling}
Suppose a nonnegative matrix \(C\) and a defect vector
\(\delta\geq0\) satisfy
\[
 a\leq\delta+Ca                                    \tag{96.16}
\]
componentwise, and \(I-C\) is a nonsingular \(M\)-matrix; it is
sufficient that
\[
 \|C\|_{\ell^\infty\to\ell^\infty}<1.              \tag{96.17}
\]
Then
\[
 a\leq(I-C)^{-1}\delta.                             \tag{96.18}
\]
If an observable pair \(F_X,F_Y\) obeys
\[
 |F_X(X)-F_Y(Y)|
 \leq\varepsilon^{\rm obs}
 +\sum_iL_i d_i(X_i,Y_i),                           \tag{96.19}
\]
then
\[
 |{\mathbb E}F_X-{\mathbb E}F_Y|
 \leq
 \varepsilon^{\rm obs}
 +L^{\mathsf T}(I-C)^{-1}\delta.                    \tag{96.20}
\]
\end{theorem}

\begin{proof}
Under (96.17),
\[
 (I-C)^{-1}=\sum_{k=0}^\infty C^k
\]
has nonnegative entries.  Multiply (96.16) by it to obtain (96.18).
Taking \(\pi\)-expectations in (96.19), then using (96.15) and
(96.18), proves (96.20).
\end{proof}

In a heat-bath construction, \(C_{ij}\) measures how a disagreement
in exterior block \(j\) changes the coupled conditional update in
block \(i\), while \(\delta_i\) measures the intrinsic mismatch
between the two cutoff specifications.  Constructing a coupling
that satisfies (96.16) is a separate task, but it does not require
mutual absolute continuity.

\subsection{Choice criterion for adjacent cutoffs}

The comparison route is now determined by:
\[
\begin{array}{c|c}
\text{equivalent common-space laws}&
\Delta U_{\rm eff}=U_1-U_0-\log(d\lambda_1/d\lambda_0)\\
\text{singular or incompatible lifts}&
\text{coupling defect }L^{\mathsf T}(I-C)^{-1}\delta.
\end{array}                                         \tag{96.21}
\]
In the first row, the V95 weighted gradient defect must be summable.
In the second, the coupling bound (96.20) must be summable along the
cutoff sequence.

\begin{assessment}[What V96 closes]
\label{ass:v16v96-status}
V96 gives the exact Radon--Nikodym correction for changes of
reference measure, including the Gaussian quadratic term and the
product-Haar refinement.  It proves that deterministic lifts of
continuous new modes are singular and cannot be repaired by a
finite action interpolation.  For singular comparisons it supplies
an \(M\)-matrix coupling alternative.

The full SM--Einstein regulator has not yet been equipped with either
a summably matched full-support lift or a summable coupled
specification.
\end{assessment}

\begin{decision}[V97: projective Gaussian lift design]
\label{dec:v16v96-next}
Construct a canonical full-support auxiliary Gaussian lift for a
linear coarse-graining map.  Compute its Radon--Nikodym correction
against the fine Gaussian reference by Schur complements, and
identify the covariance choice that makes the reference correction
vanish.  This isolates the interacting counterterm mismatch from an
avoidable reference mismatch.
\end{decision}

\begin{firewall}[Claim boundary after V96]
\label{fire:v16v96}
No summable adjacent-cutoff comparison has been established for the
actual theory.  V96 decides which comparison mechanisms are
mathematically admissible and exposes singular lifts that would make
the V95 interpolation invalid.
\end{firewall}

\section*{Version V96 append-only record}
\addcontentsline{toc}{section}{Version V96 append-only record}
The complete V95 source was retained byte for byte before this
continuation.  It had \(990982\) bytes, \(29945\) lines, and SHA--256
\[
 \texttt{A1821264A5325962D4E0D51D0B75D96EEF7E91C0AC8C3CB40C7B9232ECB52741}.
\]
V96 incorporated reference-measure changes into the cutoff action
defect and supplied a coupling alternative for singular lifts.


\section{V97: the canonical projective Gaussian lift}
\label{sec:v16v97}

\subsection{Coarse and fluctuation coordinates}

Let \(Y=\mathbb R^N\), \(X=\mathbb R^m\), and let
\[
 P:Y\to X,\qquad \operatorname{rank}P=m.             \tag{97.1}
\]
Let the fine Gaussian reference be \(\gamma_C\), where
\(C\succ0\), and define
\[
 C_c=PCP^{\mathsf T}\succ0,\qquad
 R=CP^{\mathsf T}C_c^{-1}.                          \tag{97.2}
\]
Then \(PR=I_X\).  Put
\[
 \eta=y-RPy,\qquad
 C_\eta=C-CP^{\mathsf T}C_c^{-1}PC.                 \tag{97.3}
\]

\begin{lemma}[Gaussian coarse/fluctuation factorization]
\label{lem:v16v97-factor}
If \(y\sim\gamma_C\), then \(x=Py\) and \(\eta=y-Rx\) satisfy
\[
 x\sim\gamma_{C_c},\qquad
 \eta\in\ker P,\qquad
 \operatorname{Cov}(\eta)=C_\eta,\qquad
 \operatorname{Cov}(\eta,x)=0.                     \tag{97.4}
\]
The restriction of \(C_\eta\) to \(\ker P\) is positive definite,
and \(x,\eta\) are independent.
\end{lemma}

\begin{proof}
The covariance of \(x\) is \(PCP^{\mathsf T}=C_c\), and
\[
 P\eta=Py-PRPy=0.
\]
Moreover,
\[
\begin{split}
 \operatorname{Cov}(\eta,x)
 &=
 CP^{\mathsf T}-R C_c=0,\\
 \operatorname{Cov}(\eta)
 &=
 (I-RP)C(I-RP)^{\mathsf T}
 =C_\eta.                                          \tag{97.5}
\end{split}
\]
To identify the degeneracy, write
\[
 C_\eta=C^{1/2}(I-Q)C^{1/2},\qquad
 Q=C^{1/2}P^{\mathsf T}C_c^{-1}PC^{1/2}.
\]
The matrix \(Q\) is the orthogonal projection onto
\(\operatorname{ran}(C^{1/2}P^{\mathsf T})\).  Hence
\(\ker C_\eta=\operatorname{ran}P^{\mathsf T}\), which intersects
\(\ker P=(\operatorname{ran}P^{\mathsf T})^\perp\) only at zero.
Thus the restriction to \(\ker P\) is positive definite.  Finally,
uncorrelated jointly Gaussian vectors are independent.
\end{proof}

\begin{theorem}[Exact projective reference lift]
\label{thm:v16v97-exact-lift}
Let \(\gamma_{C_\eta}^{\ker P}\) denote the nondegenerate Gaussian on
\(\ker P\) with covariance (97.3).  The map
\[
 T:X\times\ker P\to Y,\qquad T(x,\eta)=Rx+\eta      \tag{97.6}
\]
is a linear isomorphism and
\[
 T_\#\left(
 \gamma_{C_c}\otimes\gamma_{C_\eta}^{\ker P}
 \right)=\gamma_C.                                 \tag{97.7}
\]
Thus the coarse Gaussian reference lifted with its canonical
conditional fluctuation is exactly the fine Gaussian reference.
\end{theorem}

\begin{proof}
For \(y\in Y\), the inverse is
\[
 x=Py,\qquad\eta=y-RPy,
\]
so \(T\) is an isomorphism.  Lemma~\ref{lem:v16v97-factor} says that
these inverse coordinates of \(\gamma_C\) have the product law in
(97.7).
\end{proof}

\subsection{Optimality of the conditional-mean lift}

\begin{proposition}[Least-square linear reconstruction]
\label{prop:v16v97-optimal}
Among all right inverses \(S:X\to Y\) satisfying \(PS=I_X\), the map
\(R\) in (97.2) minimizes
\[
 {\mathbb E}_{\gamma_C}|y-SPy|^2.                  \tag{97.8}
\]
More precisely,
\[
\begin{split}
 {\mathbb E}|y-Sx|^2
 =
 {\mathbb E}|\eta|^2+
 \operatorname{tr}\left[
 (R-S)C_c(R-S)^{\mathsf T}
 \right].                                          \tag{97.9}
\end{split}
\]
\end{proposition}

\begin{proof}
Since \(y=Rx+\eta\),
\[
 y-Sx=\eta+(R-S)x.
\]
The two summands are uncorrelated by (97.4).  Expanding the squared
norm and taking expectation gives (97.9); its second term is
nonnegative and vanishes at \(S=R\).
\end{proof}

\subsection{Cost of a noncanonical full-support lift}

Let \(D\) be any positive covariance on \(\ker P\), and lift the
coarse reference with
\[
 \gamma_{C_c}\otimes\gamma_D^{\ker P}.              \tag{97.10}
\]

\begin{proposition}[Pure fluctuation Radon--Nikodym correction]
\label{prop:v16v97-noncanonical}
Relative to the coordinates \((x,\eta)\), the fine and lifted
references differ only in the fluctuation variable.  Up to an
additive scalar, the effective fine-minus-lift action correction is
\[
 \Delta U_{\rm ref}(\eta)
 =
 \frac12\left\langle
 \eta,(C_\eta^{-1}-D^{-1})\eta
 \right\rangle_{\ker P}.                            \tag{97.11}
\]
It vanishes identically when \(D=C_\eta\).
\end{proposition}

\begin{proof}
The coarse factor \(\gamma_{C_c}\) is common to both product
references.  Apply Proposition~\ref{prop:v16v96-gaussian} on
\(\ker P\) to \(\gamma_D\) and \(\gamma_{C_\eta}\).
\end{proof}

The inverses and determinants here are taken on \(\ker P\); the
ambient matrix \(C_\eta\) has kernel
\(\operatorname{ran}P^{\mathsf T}\) in the corresponding oblique
decomposition.

\subsection{Isolation of the interacting mismatch}

Let the fine law be
\[
 d\mu_f(y)=Z_f^{-1}e^{-V_f(y)}\,d\gamma_C(y),        \tag{97.12}
\]
and the coarse law
\[
 d\mu_c(x)=Z_c^{-1}e^{-V_c(x)}\,d\gamma_{C_c}(x).   \tag{97.13}
\]
Lift \(\mu_c\) with the canonical Gaussian fluctuation and use the
coordinates (97.6).

\begin{corollary}[Reference-free interacting defect]
\label{cor:v16v97-interacting}
The V95 effective action mismatch is, up to an additive scalar,
\[
 \Delta U_{\rm int}(x,\eta)
 =
 V_f(Rx+\eta)-V_c(x).                               \tag{97.14}
\]
No Gaussian determinant or quadratic correction remains.
\end{corollary}

\begin{proof}
Theorem~\ref{thm:v16v97-exact-lift} makes the two reference measures
identical.  Apply Lemma~\ref{lem:v16v96-rn} with Radon--Nikodym
derivative one.
\end{proof}

\begin{proposition}[Exact Wilsonian marginal identity]
\label{prop:v16v97-wilson}
The coarse action reproduces the fine \(x=Py\) marginal exactly if
and only if, up to an \(x\)-independent constant,
\[
 e^{-V_c(x)}
 =
 \int_{\ker P}e^{-V_f(Rx+\eta)}
 \,d\gamma_{C_\eta}^{\ker P}(\eta).                 \tag{97.15}
\]
\end{proposition}

\begin{proof}
Use the product factorization (97.7) in the fine partition integral
and integrate the fluctuation coordinate.  The resulting density
with respect to \(\gamma_{C_c}\) is the right side of (97.15).
Equality of normalized coarse densities is equivalent to equality
up to a scalar factor.
\end{proof}

Even under (97.15), the product lift of \(\mu_c\) is generally not
the full fine law: the interacting conditional distribution of
\(\eta\) given \(x\) is proportional to
\(e^{-V_f(Rx+\eta)}d\gamma_{C_\eta}(\eta)\), rather than the bare
Gaussian.  Equation (97.14) measures precisely this remaining
coupling.

\subsection{Compatibility across several refinements}

Let \(P_{n+1,n}:Y_{n+1}\to Y_n\) be full-rank coarse maps with
\[
 P_{n+1,n-1}=P_{n,n-1}P_{n+1,n}.                   \tag{97.16}
\]
If fine Gaussian covariances obey
\[
 C_n=P_{n+1,n}C_{n+1}P_{n+1,n}^{\mathsf T},        \tag{97.17}
\]
then Theorem~\ref{thm:v16v97-exact-lift} applies at every step and
the Gaussian marginals form a projective family.

\begin{proof}
Equation (97.17) is exactly the covariance pushforward condition.
Centered Gaussian laws are determined by their covariances, so
\[
 (P_{n+1,n})_\#\gamma_{C_{n+1}}=\gamma_{C_n}.
\]
The composition condition (97.16) gives consistency over multiple
steps.
\end{proof}

\subsection{Programme insertion}

For every linearized bosonic species, the regulator should record
\[
 P_{n+1,n},\quad C_{n+1},\quad
 C_n=P_{n+1,n}C_{n+1}P_{n+1,n}^{\mathsf T},\quad
 C_{\eta,n}.                                       \tag{97.18}
\]
The auxiliary kernel in V95 should be the canonical
\(\gamma_{C_{\eta,n}}\).  The summable defect ledger then tests only
\[
 V_{n+1}(R_nx+\eta)-V_n(x),                         \tag{97.19}
\]
plus the observable interpolation error.

\begin{assessment}[What V97 closes]
\label{ass:v16v97-status}
V97 constructs the unique conditional-mean Gaussian decomposition
associated with a full-rank coarse map.  Its canonical fluctuation
kernel lifts the coarse Gaussian reference exactly to the fine one,
so the adjacent-cutoff defect contains only the interacting action
mismatch.  Noncanonical full-support lifts add the explicit
quadratic correction (97.11).

The SM--Einstein covariances and coarse maps have not yet been shown
to form the projective family (97.17), and the interacting defect
(97.19) has not been bounded.
\end{assessment}

\begin{decision}[V98: a cumulant remainder for the coarse action]
\label{dec:v16v97-next}
Use the exact marginal formula (97.15) to derive a second-order
cumulant expansion with a rigorous remainder under an exponential
moment bound.  Separate the local relevant/marginal counterterms from
the fluctuation remainder that must enter the V95 summability ledger.
\end{decision}

\begin{firewall}[Claim boundary after V97]
\label{fire:v16v97}
No interacting renormalization step has been closed.  V97 removes an
avoidable Gaussian reference mismatch and leaves the true
fine-to-coarse interaction defect exposed.
\end{firewall}

\section*{Version V97 append-only record}
\addcontentsline{toc}{section}{Version V97 append-only record}
The complete V96 source was retained byte for byte before this
continuation.  It had \(1001548\) bytes, \(30280\) lines, and
SHA--256
\[
 \texttt{CD4A77B2CDCE69DA2A04B6AD0B9DDEF5B2A2D3A44B3AA01320FAB16A786A34D1}.
\]
V97 constructed the exact projective Gaussian reference lift and
isolated the interacting Wilsonian defect.


\section{V98: cumulant control of one Wilsonian shell}
\label{sec:v16v98}

\subsection{The exact shell effective action}

Use the V97 canonical coordinates and let
\(\mathbb E_\eta\) denote expectation under
\(\gamma_{C_\eta}^{\ker P}\).  Suppose the fine interaction has the
form
\[
 V_f(Rx+\eta)=V_0(x)+\varepsilon Z(x,\eta).          \tag{98.1}
\]
The exact coarse interaction is
\[
 W_\varepsilon(x)
 =
 V_0(x)-\log\mathbb E_\eta
 e^{-\varepsilon Z(x,\eta)}.                        \tag{98.2}
\]

\subsection{A scalar cumulant remainder}

\begin{theorem}[Second-order logarithmic moment expansion]
\label{thm:v16v98-cumulant}
Let \(Z\) be real, put
\[
 m=\mathbb EZ,\qquad X=Z-m,\qquad
 \sigma^2=\mathbb EX^2,                             \tag{98.3}
\]
and suppose for some \(a>0\), \(M<\infty\),
\[
 \mathbb E e^{a|X|}\leq M.                          \tag{98.4}
\]
For \(|\varepsilon|\leq a/2\),
\[
 -\log\mathbb Ee^{-\varepsilon Z}
 =
 \varepsilon m-\frac{\varepsilon^2}{2}\sigma^2
 +R_3(\varepsilon),                                 \tag{98.5}
\]
where
\[
 |R_3(\varepsilon)|
 \leq
 \frac43K_{a,M}|\varepsilon|^3,
 \qquad
 K_{a,M}:=
 M\left(\frac6{ae}\right)^3.                        \tag{98.6}
\]
\end{theorem}

\begin{proof}
Set
\[
 \psi(u)=\log\mathbb Ee^{-uX}.
\]
Then
\[
 \psi(0)=\psi'(0)=0,\qquad\psi''(0)=\sigma^2,
\]
and, under the tilted probability law
\[
 d\mathbb P_u=
 \frac{e^{-uX}}{\mathbb Ee^{-uX}}\,d\mathbb P,
\]
one has
\[
 \psi'''(u)
 =
 -\mathbb E_u
 (X-\mathbb E_uX)^3.                                \tag{98.7}
\]
Jensen gives \(\mathbb Ee^{-uX}\geq1\).  For
\(|u|\leq a/2\),
\[
\begin{split}
 \mathbb E_u|X|^3
 &\leq
 \mathbb E\left[|X|^3e^{a|X|/2}\right]\\
 &\leq
 \sup_{r\geq0}\left(r^3e^{-ar/2}\right)
 \mathbb Ee^{a|X|}
 =K_{a,M}.                                         \tag{98.8}
\end{split}
\]
Since
\[
 |x-y|^3\leq4(|x|^3+|y|^3)
\]
and
\(|\mathbb E_uX|^3\leq\mathbb E_u|X|^3\),
\[
 |\psi'''(u)|\leq8K_{a,M}.                          \tag{98.9}
\]
Taylor's theorem with integral remainder gives
\[
 \left|\psi(\varepsilon)
 -\frac{\varepsilon^2}{2}\sigma^2\right|
 \leq\frac{8K_{a,M}}6|\varepsilon|^3.               \tag{98.10}
\]
Finally,
\[
 -\log\mathbb Ee^{-\varepsilon Z}
 =\varepsilon m-\psi(\varepsilon),
\]
which proves (98.5)--(98.6).
\end{proof}

\begin{corollary}[Pointwise shell expansion]
\label{cor:v16v98-pointwise}
If (98.4) holds for the centered shell variable
\[
 X_x=Z(x,\eta)-\mathbb E_\eta Z(x,\eta)
\]
uniformly in \(x\), then
\[
\begin{split}
 W_\varepsilon(x)
 =
 V_0(x)
 +\varepsilon\mathbb E_\eta Z(x,\eta)
 -\frac{\varepsilon^2}{2}
 \operatorname{Var}_\eta Z(x,\eta)
 +R_{\varepsilon}(x),                              \tag{98.11}\\
 \sup_x|R_\varepsilon(x)|
 \leq\frac43K_{a,M}|\varepsilon|^3.                 \tag{98.12}
\end{split}
\]
\end{corollary}

\begin{proof}
Apply Theorem~\ref{thm:v16v98-cumulant} for each fixed \(x\).
\end{proof}

\subsection{The action norm needed by V95}

Pointwise control of an effective action is not the V95 defect
control, which uses block gradients.  For a fixed observable support
\(S\), define the seminorm
\[
 \|A\|_{{\mathcal B}_{S,n}}
 :=
 \sup_{0\leq s\leq1}
 \sum_j e^{-m_{*,n}d^{\rm phys}(S,j)}
 \|\nabla_jA\|_{L^2(\mu_{n,s})}.                    \tag{98.13}
\]

\begin{lemma}[Banach-valued Taylor remainder]
\label{lem:v16v98-banach}
Let \(\varepsilon\mapsto W_\varepsilon\) be three times continuously
differentiable as a map into the completion of
\({\mathcal B}_{S,n}\).  If
\[
 \sup_{|u|\leq\varepsilon_0}
 \|\partial_u^3W_u\|_{{\mathcal B}_{S,n}}
 \leq M_{3,n}(S),                                   \tag{98.14}
\]
then for \(|\varepsilon|\leq\varepsilon_0\),
\[
\left\|
 W_\varepsilon-W_0-\varepsilon W'_0
 -\frac{\varepsilon^2}{2}W''_0
 \right\|_{{\mathcal B}_{S,n}}
 \leq\frac{|\varepsilon|^3}{6}M_{3,n}(S).           \tag{98.15}
\]
\end{lemma}

\begin{proof}
The Bochner integral form of Taylor's theorem is
\[
 W_\varepsilon-W_0-\varepsilon W'_0
 -\frac{\varepsilon^2}{2}W''_0
 =
 \frac12\int_0^\varepsilon
 (\varepsilon-u)^2W'''_u\,du.
\]
Take the seminorm, use (98.14), and integrate.
\end{proof}

For (98.2), differentiation at zero gives
\[
 W'_0=\mathbb E_\eta Z,\qquad
 W''_0=-\operatorname{Var}_\eta Z.                 \tag{98.16}
\]
Establishing (98.14) requires exponential moment bounds jointly with
the block gradients of \(Z\).  The scalar estimate (98.12) does not
imply it.

\subsection{Relevant, marginal, and remainder sectors}

Let \(\Pi_{\rm rm}\) be a fixed finite-rank projection onto the chosen
relevant and marginal local counterterm space, and put
\[
 K_1=\mathbb E_\eta Z,\qquad
 K_2=-\frac12\operatorname{Var}_\eta Z.             \tag{98.17}
\]
Choose the candidate coarse action
\[
 V_c
 =
 V_0+\Pi_{\rm rm}
 \left(\varepsilon K_1+\varepsilon^2K_2\right).     \tag{98.18}
\]

\begin{proposition}[Exact counterterm/remainder decomposition]
\label{prop:v16v98-decomposition}
The exact marginal mismatch is
\[
\begin{split}
 W_\varepsilon-V_c
 =
 (I-\Pi_{\rm rm})
 \left(\varepsilon K_1+\varepsilon^2K_2\right)
 +R_\varepsilon.                                   \tag{98.19}
\end{split}
\]
If Lemma~\ref{lem:v16v98-banach} applies, then
\[
\begin{split}
 \|W_\varepsilon-V_c\|_{{\mathcal B}_{S,n}}
 \leq&
 |\varepsilon|
 \|(I-\Pi_{\rm rm})K_1\|_{{\mathcal B}_{S,n}}\\
 &+
 \varepsilon^2
 \|(I-\Pi_{\rm rm})K_2\|_{{\mathcal B}_{S,n}}
 +\frac{|\varepsilon|^3}{6}M_{3,n}(S).             \tag{98.20}
\end{split}
\]
\end{proposition}

\begin{proof}
Insert (98.11) and (98.18) and collect the projected and
complementary terms.  The seminorm bound is the triangle inequality
and (98.15).
\end{proof}

The projection must be fixed by symmetry and scaling, not chosen
after seeing the desired estimate.  In particular, all
gauge-invariant and diffeomorphism-compatible counterterms allowed by
the regulator must be included before claiming that the complement
is irrelevant.

\subsection{Comparison on the coarse marginal space}

For a coarse-measurable observable \(F(x)\), the fine expectation
equals expectation under the exact marginal action
\(W_\varepsilon\) by Proposition~\ref{prop:v16v97-wilson}.

\begin{theorem}[One-shell coarse-observable error]
\label{thm:v16v98-one-shell}
Suppose the interpolation from \(V_c\) to \(W_\varepsilon\) satisfies
the uniform V95 influence hypotheses.  Then
\[
\begin{split}
 |{\mathbb E}_{\mu_f}F(Py)
  -{\mathbb E}_{\mu_c}F(x)|
 \leq
 \frac{\sup_sG_s(F)}
 {\rho_*(1-\theta_*)}
 \|W_\varepsilon-V_c\|_{{\mathcal B}_{S_F,n}}.
                                                               \tag{98.21}
\end{split}
\]
Under Proposition~\ref{prop:v16v98-decomposition}, the right side is
bounded by (98.20) times the displayed prefactor.
\end{theorem}

\begin{proof}
Push the fine law forward by \(P\), obtaining the exact action
\(W_\varepsilon\).  Apply
Corollary~\ref{cor:v16v95-physical-comparison} on the coarse space to
the path joining \(V_c\) and \(W_\varepsilon\).
\end{proof}

\subsection{The true locality bottleneck}

The variance term is
\[
 \operatorname{Var}_\eta Z
 =
 \mathbb E_\eta Z^2-(\mathbb E_\eta Z)^2.           \tag{98.22}
\]
Even when \(Z\) is a sum of local terms, this contains covariances
between distinct shell locations.  To prove (98.20) summable one
needs decay of the shell covariance and derivative bounds uniform in
the coarse field \(x\).  Calling the Gaussian shell ultraviolet does
not establish those bounds when its covariance has long tails or
when gauge/gravitational constraints couple its blocks.

\begin{assessment}[What V98 closes]
\label{ass:v16v98-status}
V98 gives a rigorous second-order cumulant expansion with an explicit
cubic scalar remainder and a Banach-valued remainder theorem in the
exact V95 weighted-gradient seminorm.  It separates the
relevant/marginal counterterms from the complementary first and
second cumulants and the cubic shell remainder.

Uniform exponential-gradient moments, locality of the second
cumulant, and summability of (98.20) have not been proved for the
SM--Einstein shell.
\end{assessment}

\begin{decision}[V99: consolidate the continuum gates]
\label{dec:v16v98-next}
As the present cycle approaches V100, consolidate V77--V98 into one
dependency theorem.  State the shortest verified implication chain
from the finite-cutoff symbol and block certificates to a unique
reflection-positive continuum and Hamiltonian gap, listing every
unproved gate without replacing it by a conclusion.
\end{decision}

\begin{firewall}[Claim boundary after V98]
\label{fire:v16v98}
No renormalization step or interacting continuum limit has been
proved.  V98 identifies the precise cumulant and weighted-gradient
estimates required to make the V95 adjacent-cutoff series converge.
\end{firewall}

\section*{Version V98 append-only record}
\addcontentsline{toc}{section}{Version V98 append-only record}
The complete V97 source was retained byte for byte before this
continuation.  It had \(1010846\) bytes, \(30573\) lines, and
SHA--256
\[
 \texttt{64E640176B29857A5EE436E6DA16097C7B24B136AB61C89211189C467ED6156F}.
\]
V98 proved the controlled one-shell cumulant expansion and the exact
counterterm/remainder decomposition.


\section{V99: consolidated continuum dependency theorem}
\label{sec:v16v99}

\subsection{The eight gates}

The results V77--V98 organize the remaining construction into eight
logically separate gates.

\begin{enumerate}
\item \textbf{Finite-cutoff physical law.}
      The gauge-fixed or physical-coordinate cutoff functional is
      normalized and positive on the intended observable algebra.
      Fine modes have a uniform floor, coarse rigid modes are
      controlled by full-face moments, and the Schur margin remains
      positive.
\item \textbf{Conditional mixing.}
      Conditional Poincar\'e floors and mixed-Hessian influences give
      the V84 \(M\)-matrix resolvent.  Conditional log-Sobolev floors
      and the V90 physical heat-bath sweep give the entropy
      tensorization needed for moments.
\item \textbf{Uniform physical scaling.}
      The rate and amplitude ledgers obey
      \[
      \inf_n\frac{-\log\theta_n}{\ell_nR_{0,n}}>0,
      \qquad
      \sup_n
      \frac{{\cal G}_n^2}
      {\rho_{\min,n}(1-\theta_n)}<\infty.            \tag{99.1}
      \]
\item \textbf{Cutoff compactness.}
      Renormalized bosonic fields satisfy the V88 negative-Sobolev
      moment bound, or a separating bounded physical-coordinate
      family satisfies the V91 no-escape conditions.  Fermionic
      kernels satisfy the corresponding graded bounds.
\item \textbf{Renormalization Cauchy property.}
      Canonical Gaussian/Haar lifts put adjacent references on a
      common space, and the V95 observable plus weighted action
      defects are summable on a measure-determining algebra.
\item \textbf{Physical continuum structure.}
      The limit has the V92 uniform label continuity, full Euclidean
      symmetry, the V93 local net and graded permutation relations,
      and the V94 uniqueness ledgers.  Gauge and diffeomorphism
      identities survive on a separating physical algebra.
\item \textbf{Reflection-positive OS limit.}
      Reflected cutoff Gram forms are positive on the common
      positive-time physical algebra and converge.  The limiting
      functional has the regularity required for OS reconstruction
      and a unique vacuum.
\item \textbf{Dense-channel spectral gap.}
      Uniform Euclidean-time clustering holds with one rate \(m>0\)
      on a linear observable family whose centered reconstructed
      vectors are dense in the nonvacuum Hilbert space.
\end{enumerate}

\subsection{The conditional composition theorem}

\begin{theorem}[Finite cutoff to a unique gapped OS continuum]
\label{thm:v16v99-composition}
Assume Gates \(1\)--\(8\) for a cofinal sequence of finite-volume
cutoffs, with all convergence, scaling, and summability bounds
uniform in the order in which volume and ultraviolet regulators are
removed.  Then:
\begin{enumerate}
\item the physical cylinder functionals converge without subsequence
      dependence to a unique normalized positive functional
      \({\cal S}\);
\item \({\cal S}\) is Euclidean invariant and reflection positive on
      the common physical algebra and obeys its verified local and
      graded relations;
\item its connected bounded physical correlations cluster
      exponentially at rate at least the \(m\) in Gates \(3\) and
      \(8\);
\item OS reconstruction produces
      \(({\cal H},\Omega,H)\) with
      \[
      \operatorname{spec}(H|_{\Omega^\perp})
      \subset[m,\infty).                            \tag{99.2}
      \]
\end{enumerate}
\end{theorem}

\begin{proof}
Gate \(4\) gives subsequential existence by V88 or V91.  Gate \(5\)
gives convergence of every member of a measure-determining algebra
by Theorem~\ref{thm:v16v95-summable}; Gate \(6\) and
Theorem~\ref{thm:v16v94-full-convergence} give uniqueness and
full-sequence convergence.  The closure theorems V87, V92, and V93
pass reflection positivity, Euclidean invariance, and the stated
bounded local relations to \({\cal S}\).

Gates \(1\)--\(3\), through V84--V85, give a cutoff-uniform physical
clustering estimate.  Gate \(5\) passes its bounded correlation
functions to the limit, as in
Proposition~\ref{prop:v16v93-clustering-limit}.  Gate \(7\) permits
OS reconstruction and identifies time-separated autocorrelations
with positive spectral Laplace transforms.  Gate \(8\) and
Corollary~\ref{cor:v16v86-local-gap} then give (99.2).
\end{proof}

\subsection{Uniformity in the double limit}

\begin{proposition}[Commutation criterion for volume and cutoff
limits]
\label{prop:v16v99-double}
Let \(a_{n,L}(F)\) be the expectation of a bounded determining
observable at cutoff \(n\) and volume \(L\).  Suppose
\[
\begin{split}
 \sup_L|a_{n+1,L}(F)-a_{n,L}(F)|
 &\leq\varepsilon_n(F),
 \qquad\sum_n\varepsilon_n(F)<\infty,\\
 \sup_n\sup_{\tau,\tau'}
 |a_{n,L}^{\tau}(F)-a_{n,L}^{\tau'}(F)|
 &\leq\delta_L(F),
 \qquad\delta_L(F)\to0.                             \tag{99.3}
\end{split}
\]
If one reference-boundary thermodynamic limit exists for every
\(n\), then the iterated limits exist, are boundary independent, and
agree:
\[
 \lim_{n\to\infty}\lim_{L\to\infty}a_{n,L}(F)
 =
 \lim_{L\to\infty}\lim_{n\to\infty}a_{n,L}(F).      \tag{99.4}
\]
\end{proposition}

\begin{proof}
The first line of (99.3) makes the cutoff sequence uniformly Cauchy
in \(L\).  Hence \(a_{n,L}\) converges as \(n\to\infty\) uniformly in
\(L\).  The second line removes the boundary condition uniformly in
\(n\).  The elementary theorem on interchange of a uniform limit
with a pointwise limit gives (99.4).
\end{proof}

\subsection{Verified inputs and open gates}

The following pieces are proved in this cycle:
\[
\begin{array}{c|c}
\text{piece}&\text{location}\\
\hline
\text{rigid coarse decomposition and face-moment repair}&V77\text{--}V82\\
\text{finite-cutoff certificate architecture}&V83\\
\text{localized covariance and physical rate}&V84\text{--}V85\\
\text{positive spectral-support handoff}&V86\\
\text{closure of reflection positivity}&V87\\
\text{tightness, entropy, and physical coordinates}&V88\text{--}V91\\
\text{symmetry, locality distinctions, and uniqueness}&V92\text{--}V94\\
\text{adjacent-cutoff and shell comparison machinery}&V95\text{--}V98.
\end{array}                                         \tag{99.5}
\]
These are implication theorems and local structural results.  The
following programme-specific inputs remain unproved:
\begin{enumerate}
\item the complete YM full-torus sign/symbol certificate;
\item positivity of the full chiral fermion, gauge, Higgs, and
      gravitational cutoff physical functional;
\item control of the Einstein conformal direction and a global
      gauge/diffeomorphism-compatible physical observable algebra;
\item the entropy-sweep constants and uniform Sobolev means;
\item projective regulator covariances and summable interacting
      cumulant defects;
\item cutoff reflection positivity, OS regularity, vacuum
      uniqueness, and density of uniformly gapped channels.
\end{enumerate}

\begin{assessment}[Position at V99]
\label{ass:v16v99-status}
V99 closes the logical architecture: once the named
programme-specific gates are supplied with the recorded uniformity,
the already proved implication chain reaches a unique
reflection-positive continuum with exponential clustering and a
Hamiltonian gap.  The theorem does not supply those missing
estimates.

The nearest upstream analytic target is the summable shell
weighted-gradient bound (98.20), while the nearest foundational
target is positivity/reflection positivity of the full physical
cutoff functional.  Neither should be bypassed by a formal continuum
claim.
\end{assessment}

\begin{decision}[V100: audit, checkpoint, and rollover]
\label{dec:v16v99-next}
Perform the mandatory companion audit.  Record a compact theorem
inventory and the exact unresolved frontier, freeze the completed
cycle as \texttt{V100\_f16}, and restart at V1 with a concise context
and result summary.  The restarted V1 must also continue
substantively by deriving a usable sufficient condition for the
Banach-valued third-cumulant bound in (98.14).
\end{decision}

\begin{firewall}[Claim boundary after V99]
\label{fire:v16v99}
The conditional composition theorem is not a proof of the full
SM--Einstein continuum, because several Gates \(1\)--\(8\) remain
unverified.  The full programme remains open.
\end{firewall}

\section*{Version V99 append-only record}
\addcontentsline{toc}{section}{Version V99 append-only record}
The complete V98 source was retained byte for byte before this
continuation.  It had \(1020017\) bytes, \(30904\) lines, and
SHA--256
\[
 \texttt{76D8EDBEDB9757FC8A8E0DBE4BA721E1BCBBC452CFBCCF2FD7996309A30826DC}.
\]
V99 consolidated the proved implication chain and the remaining
programme-specific continuum gates.


\section{V100: companion audit and sixteenth-cycle checkpoint}
\label{sec:v16v100}

\subsection{Mandatory companion audit}

The V100 audit inspected
\[
\begin{array}{lll}
\text{YM V65:}&15119\text{ lines},&
\texttt{5543FA7C0DFDE051056F156EFCB588EFE636C80F66FC640448AEAFE48A4FE165},\\
\text{NS V26:}&5319\text{ lines},&
\texttt{82C887F8C2C69E4F28FAD7C3DC8DE5B9099CB5A4BF431EBA1FFF3E91935978AD}.
\end{array}                                                   \tag{100.1}
\]
YM V65 completes the mixed/fourth inversion pair: both direct signs
now hold on four complete transverse quarter cells, with exact common
floor \(\alpha_{57}>0\).  The four sign cells
\((+++),(---),(+-+),(-+-)\) remain untreated, so the full torus and
all continuum conclusions remain open.  NS V26 is unchanged.

\begin{theorem}[V100 audit decision]
\label{thm:v16v100-audit}
YM V65 strengthens the prospective gauge-symbol entry of Gate \(1\)
in V99 from paired slabs to four complete quarter cells.  It does not
alter the continuum, entropy, renormalization, reflection-positivity,
or spectral gates.  No NS result is imported.
\end{theorem}

\begin{proof}
The YM V65 firewall expressly leaves four transverse sign cells and
every continuum conclusion open.  The NS note contains no statement
about the cutoff probability or Schwinger functional.
\end{proof}

\subsection{Sixteenth-cycle result inventory}

\begin{theorem}[Results established in V1--V100 of this cycle]
\label{thm:v16v100-inventory}
The cycle establishes the following reusable mathematical blocks.
\begin{enumerate}
\item Rigid-motion Korn decomposition and a full-face moment
      discretization that removes the rotational hourglass mode.
\item A finite-cutoff positive-law certificate with explicit
      fine/coarse Schur, gauge, matter, conditional-gap, and
      influence margins.
\item A block Witten-resolvent covariance theorem and exponential
      decay under finite-range strict influence.
\item The exact physical correlation rate
      \(m_n=-\log\theta_n/(\ell_nR_{0,n})\) and its independent
      amplitude condition.
\item The positive spectral-measure theorem converting uniform
      Euclidean-time autocorrelation decay into channel gaps and,
      on a dense family, a Hamiltonian gap.
\item Closure of reflection positivity, Euclidean symmetry, bounded
      local identities, clustering, and determining-algebra
      uniqueness under the recorded convergence hypotheses.
\item Negative-Sobolev and physical-coordinate compactness routes,
      together with the conditional-LSI/entropy-sweep concentration
      mechanism.
\item Localized boundary and adjacent-cutoff comparison theorems,
      including reference-measure corrections, singular-lift
      coupling, the canonical projective Gaussian lift, and a
      controlled one-shell cumulant expansion.
\end{enumerate}
\end{theorem}

\begin{proof}
The items are respectively V77--V83, V84--V86, V87--V94, and
V95--V98, consolidated by Theorem~\ref{thm:v16v99-composition}.
\end{proof}

\subsection{Exact unresolved frontier}

The full SM--Einstein continuum has not been reached.  The shortest
current frontier consists of:
\begin{enumerate}
\item completion of the remaining YM sign cells and insertion into a
      positive full physical cutoff law;
\item a gauge- and diffeomorphism-compatible physical observable
      algebra, including the graded chiral fermion sector and control
      of the Einstein conformal direction;
\item uniform entropy-sweep, Sobolev-mean, and physical clustering
      ledgers;
\item projective regulator matching and a summable bound for the
      shell remainder in the V98 weighted-gradient seminorm;
\item finite-cutoff reflection positivity, limiting OS regularity,
      vacuum uniqueness, and a dense common-rate channel family.
\end{enumerate}

\begin{assessment}[Sixteenth-cycle checkpoint]
\label{ass:v16v100-status}
The cycle has converted a collection of local symbol and block
estimates into a rigorous conditional continuum architecture and has
isolated the remaining estimates without circularly identifying a
Poincar\'e gap, clustering rate, and Hamiltonian mass gap.

The next analytic attack is the derivative version of the shell
cumulant bound.  The next foundational attack is construction of the
reflection-positive physical cutoff functional.  Both remain open.
\end{assessment}

\begin{decision}[Freeze and restart]
\label{dec:v16v100-restart}
Freeze this file as \texttt{V100\_f16}.  Begin the seventeenth cycle
at V1 with context, a concise result summary, the unresolved frontier,
and a substantive theorem giving sufficient exponential-gradient
moments for the Banach-valued third-cumulant bound.
\end{decision}

\begin{firewall}[Claim boundary at the sixteenth-cycle freeze]
\label{fire:v16v100}
The checkpoint records rigorous implication and comparison theorems,
not a construction of the full Standard Model--Einstein continuum.
All frontier items above remain to be proved.
\end{firewall}

\section*{Version V100 append-only record}
\addcontentsline{toc}{section}{Version V100 append-only record}
The complete V99 source was retained byte for byte before this
continuation.  It had \(1028941\) bytes, \(31116\) lines, and
SHA--256
\[
 \texttt{E7CF77D21D1DF9CB9BAE4E479FC5E94729F53FC1D3A2842C9816CAA95A692D56}.
\]
V100 performed the mandatory companion audit and fixed the exact
inventory and frontier for the sixteenth-cycle freeze.

\end{document}
